\documentclass[12pt,a4paper]{article}

\usepackage{mathtools}
\usepackage{mathrsfs}
\usepackage{float}
\usepackage{caption}
\usepackage{tikz}
\usepackage{tikz-cd}
\usetikzlibrary{arrows,positioning,calc,decorations.pathmorphing}

% Macros-v2_203: v3.077 source-current bindings for unconditionality source integration. No mathematical macro changes.
%% Macros-v2_201: downstream full-carrier geometry update, five current root bindings, no mathematical macro change.
%% Macros-v2_200: Paper III-F v1.280 current-source binding for framework v3.074; mathematical macro definitions unchanged.
%% Macros-v2_198: current project v3.072 III-D and bibliography binding; mathematical macros unchanged.
%% Current project v3.070 source-reference bindings; previous macro bodies unchanged.
% Current II-D reference root for PapersMinimal v3.063; mathematical macros otherwise unchanged.
% Macros-v2_188: re-based current paper identities and A--G family; no change to geometry.
% Macros-v2_184: active Paper III-E v1.200/III-F v1.267 binding for framework v3.031; mathematical macros unchanged.
% Macros-v2_183: updated shared cross-reference bindings for Paper1B/III-E/III-F v3.030; mathematical macros unchanged.
% Macros-v2_182: active Paper II-F/V-A cross-references for framework v3.029; definitions unchanged.
% Paper-series v2.381: new I-B torsion/modular results, bibliography and current source roots.
% Paper-series v2.380: new exceptional bibliography and synchronized paper references.
% Paper-series v2.379 current 31-paper macro/aux registry.
% Paper-series v2.378: current 31-paper AUX/bibliography registry updated for Paper I-B v1.2 literature additions.
% Framework v2.371: current-source reference registry refreshed; old theorem labels retained.
% Framework v2.370: current 30-paper source/AUX root rebinding; native mathematical macros unchanged.
% Framework v2.368: synchronized merged 30-root registry; mathematical notation unchanged.
% Framework v2.358: active 30-paper source registry synchronized; existing mathematical macros unchanged.
% Framework v2.357: selective new source-owned certificates; active reference roots updated for five revised papers.
% Framework v2.356: active 30-root registry synchronized to the 30 papers in this release.
% Framework v2.275: exact positive dense-twin word cubature, finite repeatability and eighth-moment classification, nested sampling frames, joint action and POVM-jet exports; preserves v2.274.
% Framework v2.274: filename and root-source synchronization.
% Framework v2.273: restored missing Problem 3 theorems from v2.268; exact exceptional frames, twin geometry, convolution repeatability and projector/occurrence firewalls; merged on v2.272 without overwriting independent branches.
% Framework v2.272: merged v2.270 covariant native-readout mathematics with v2.271 Problem 4 finite-action Hopf hypercharge/holonomy corrections; no universal bulk-gauge uniqueness claim.
% Independent branches: v2.270 readout/guard exports and v2.270 K16 action-consumed charge closure; all prior scoped assumptions retained.
% Framework v2.269: synchronized 30 active AUX roots; mathematics and bibliography preserved from Macros-v2_161.
% Framework v2.266: Paper II-C converse-derived outcome resources; II-E, III-C, IV-A/B, V-A and VI-B downstream mathematical closures; 30-root registry synchronized.
% Framework v2.265: II-C coordinate-stage constitution and exact first-symbol host inheritance; prior mathematical domains retained.
% Framework v2.264: paper II-B reverse/audit and IV-A observable-minimal exchange; synchronized external references to 30 active roots.
% Framework v2.263: Paper I-C adds bridge-controlled root-carry parity and source-relative old-cap cancellation; other root mathematical bodies retained; active 30-root references synchronized.
% Framework v2.262: integrates Paper II-F radius, Casimir and readout results; active 30-root references synchronized.
% Framework v2.261: merged Paper II-B GP2-level/geometric repairs; preserves v2.260 II-D repairs; active references synchronized.
% Framework v2.260: Paper II-D signed RS, finite-symbol typing, and conditional curvature/Jordan repairs; references synchronized.
% Framework v2.259: conditional native SU(2) atlas and explicit later discharge; v2.258 mathematical bodies retained elsewhere.
% Framework v2.258: preserves parallel II-A/dynamics update and adds complete readout/frequency closures.
% Framework v2.257: exact minimal complete position readout and arbitrary-preparation frequency law.
% Framework v2.256: exact native inertia and sharp bounded ordered-source closure; preserves concurrent v2.255 work.
% Framework v2.255: preserves parallel detector integration and adds exact finite readout closures.
% Framework v2.254: Pasterski bibliography and synchronized 30-paper external-reference registry.
% Framework v2.253: typed raw coordinate defect in I-A and synchronized VI-A import; active references synchronized.
% Framework v2.252: synchronized source registry; optional III-G shell-frame calculation.
% Framework v2.251: forward common-action execution, joint source jets and scoped coupled control.
% Framework v2.251: exact full-heat variations, fermion source jets and explicit physical completion frontier.
% Framework v2.251: certified interacting source jets, pole stability and scoped completion proofs.
% Framework v2.251: synchronized active source registry; mathematical notation unchanged.
% Framework v2.251: synchronized active source registry; mathematical notation unchanged.
% Framework v2.251: synchronized active source registry; mathematical notation unchanged.
% Framework v2.251: synchronized source registry; no mathematical notation changes.
% Framework v2.244: synchronized 30-paper AUX registry and rooted phase-source action-order tests.
% Framework v2.242: signed rooted phase lift, complete transverse return and exact family cubic.
% Framework v2.241: proved rooted conormal geometry, conditional cap realization and scoped action tests.
% Framework v2.239: all 30 active source roots synchronized; no mathematical macro changes.
% Framework v2.238: Paper IV A/B/C split and 30-root reference registry.
% Framework v2.237: Paper VII-D technical split and 28-root reference registry.
% Framework v2.236: Paper VI-C technical-volume reference added; all current root bindings advanced.
% Framework v2.235: 26-root AUX and bibliography bindings synchronized; mathematical macros unchanged.
% BEGIN INLINE SOURCE: Paper00-TeXBuildCompatibility-v1_0.tex
% Rendering-only repairs for inherited long appendices and bitmap text fonts.
% Integer counters through 26 retain the original labels; later values use AA,
% AB, ... . This file changes no action coefficient, state or theorem premise.
\RequirePackage{alphalph}
\makeatletter
\let\BHPrimitiveAlph\@Alph
\newalphalph{\BHExtendedAlph}{\BHPrimitiveAlph}{26}
\let\@Alph\BHExtendedAlph
\AtBeginDocument{%
  \@ifpackageloaded{microtype}{\microtypesetup{expansion=false}}{}%
}
\makeatother

% END INLINE SOURCE: Paper00-TeXBuildCompatibility-v1_0.tex

% Framework v2.164: current reference/bibliography profile; mathematical definitions retained.
% Framework v2.162: finite metric, functional Hilbert and scoped Ward integration.
% Framework v2.161: inherits Macros-v2_57.tex exactly except the current external-reference binding.
% Macros v2.57 / framework v2.156: current reference registry only; mathematical definitions unchanged.
% Framework v2.143 / Macros v2.55: current external-reference bindings; mathematical notation unchanged.
% Framework v2.110 / Macros v2.54: current external-reference registry and bibliography; mathematical notation inherited.
% Framework v2.67 / Macros-v2_53: adds Born-action-cell and joint-EH readout symbols; absolute unit realization is metrological; MATCH_{J/scal}^{EH} remains the sole G24 normalization bridge.
% Framework v2.35 / Macros-v2_52: grouped paper numbering, active AUX registry, and Born-v6_20 bibliography binding synchronized; notation unchanged.
% Framework v2.26: active-source AUX and bibliography bindings synchronized; notation unchanged.
% Framework v2.09: reviewed early-reciprocity integration; scoped physical imports.
% Framework v2.08: Born--Euler mathematics and conditional complete-action potential criterion.
% Framework v2.07: active AUX registry includes Paper I-B v1.2; shared bibliography Born-v6_18.
% Macros-v2_46 / framework v2.04: mathematical atlas and downstream realization correction.
% Macros-v2_46 / framework v2.03: scoped structure integration and current dependency bindings.
% Macros-v2_46 / framework v2.01: shared source bindings; prior revision comments are historical.
% v2.42 / framework v1.99: Synchronized RG/AS owner and bibliography bindings.
% v2.41 / framework v1.99: active manuscript AUX bindings and geometric-closure bibliography synchronized.
% Release-v1.80 centralized source synchronization for scalar-circle/comparison-band integration.
% Release v1.79: Hurwitz hierarchy and kinetic-normalized action response integration; physical mass matching remains conditional.
% Release v1.78: approved M8 closure-batch integration; physical and analytic domains retained.
% Current v1.78 registry: all active TeX roots are synchronized below.
% Earlier release comments are historical scope records, not current status.
% Release-v1.77 synchronization: native same-owner Noether closure, generator--flow/history typing, analytic-certification firewalls, scalaron/K16, LF/QCD, quantum-Ward, and RAR boundary-owner transport are registered; shared mathematical notation is unchanged.
% Release-v1.76 synchronization: registry points to the v1.76 active paper roots.
% Release-v1.76 synchronization: registers the local A+Phi+g Noether/constraint closure and dynamics-frontier propagation; shared mathematical notation is unchanged.
% Release-v1.75 synchronization: registry points to the v1.75 active paper roots.
% Release-v1.75 synchronization: integrates the joint-Dirac Einstein bulk real form and the local odd-scalar heat-coefficient action on explicit analytic first-jet/interior domains; v1.74 Yang--Mills closure is retained. Fermionic reality, pure-metric curvature-squared/exact-nonlocal terms, noncanonical boundaries, physical vacuum selection, and actual solution-history tokens remain open.
% Release-v1.74 synchronization: registry points to the v1.74 active paper roots and retains the v1.73 theorem labels.
% Release-v1.74 synchronization: canonical Macros-v2_34.tex; generated local nonlinear Yang--Mills action, constrained connection-field phase space, source-free Gauss momentum map, and Legendre closure are integrated on their proved domains. The remaining non-gauge ACT_cl sectors and actual solution-history selection remain open.
% Release-v1.71 synchronization: closes the former JOINT_x/JOINT_5
% connection-occurrence gap by canonical structure-group extension of the
% nonflat Born connection, faithful associated representation, and the local
% cotangent--Clifford operator constructor.  Actual fluctuation members,
% actions, vacua, solutions, calibrations, and analytic QFT domains remain
% separately typed.  The active-AUX registry is advanced once.
% Release-v1.71 canonical macro and active-AUX registry.
% ======================================================================
% Macros-v2_39.tex -- canonical publication macro file
% Release v1.71 retains the v1.70 null-ray/QCD boundaries and additionally
% removes the duplicate structural connection actuality premise.  It does not
% claim global bundle topology, select a gauge/Higgs member, or promote the
% finite construction to an interacting continuum theory.
% Release v1.68 constructs the finite light-front total-fraction operators and
% parton number/momentum operator measures, proves their simplex and Hausdorff-limit
% theorems, and separates positive cut measures from scheme-renormalized
% light-ray distributions.  It also separates finite real Schur poles, stable
% norm-resolvent/Riesz pole-residue limits, and continued resonance
% pole/residue limits.  The Paper-VII-C scattering package is claim-indexed:
% local QFT and spectral claims import no S-matrix or infrared premise, while
% asymptotic amplitudes, LSZ, S-matrix poles, completeness, and cross sections
% import only their named massive or massless/long-range clauses.
% Release v1.67 factorizes BND48_H into a same-action Hessian/Schur
% descendant and still-open typed boundary-trace, critical-lift, and
% refinement-limit conditions.  It proves the finite-collar Dirichlet-to-
% Neumann law sqrt(y)tanh(L sqrt(y)), its nonuniform critical limit, and the
% finite-return ultra-infrared turnover.  The Paper-I-A Hardy constant mode is
% an exact geometric near-hit, not yet the physical metric boundary mode.
% The current local f(R) Hessian remains massive and therefore still refutes
% a RAR prediction on the minimal action branch.
% Release v1.66 derives the transverse depth-six normalizer bridge from the
% already generated C6 Reynolds image and Paper-I-A binary-dihedral normalizer.
% The bridge is the unique least 2D3-stable scalar-containing realization of
% the faithful quotient 2D3/C6=C2, has trace defect six and hence K=48, but is
% transverse to the nested active carrier.  Exact K=44 and K=52 contraction,
% semigroup, and energy-midpoint identities are registered as structural
% readouts rather than projector events.  The active E48 carrier gap remains
% unchanged.  A same-action boundary-symbol/Calderon intertwiner is isolated
% as the necessary and still unproved owner bridge to the static DtN sort;
% representation equivalence alone neither derives RAR nor creates a K=48
% action junction, pole, occupation, or cosmological transition.
% Release v1.65 derives the native E_48-centered carrier mass gap from the
% GP1--GP3-generated active spectrum, distinguishes its carrier-free interval
% from a particle located at E_48, and withdraws the unsupported promotion to
% a positive continuum Yang--Mills mass gap.  The latter remains a valid
% conditional transfer theorem on an enlarged continuum branch, but it is not
% an active framework prediction or an unfinished native-gap proof.
% Release v1.64 synchronizes Derivation Batches 14--22.  It registers the
% GP1--GP3 heat-profile orbit, same-Hessian electroweak response and pole
% ownership, K16/K48 action-response corrections, perturbation state and
% non-Gaussianity typing, finite-vacuum/gap/QCD limit repairs, factorized
% position-area ownership, the actual static metric response, and the ordered
% Wilson/string-tension limits.  Strict continuum convergence, positive
% physical gap, surviving higher cumulants, and positive tension remain open.
% Release v1.61 closes TIMEWIDTH16 as an independent direct-production datum.
% JUNC_{2|7} continues to own the finite intrinsic microscopic profile/width,
% while the actual same-action background/production history determines its
% physical-time pullback, support, and duration.  The current Paper-I-A--VII-C AUX
% registry is advanced once; shared mathematical notation is unchanged.
% Release v1.60 derives the final-octet primitive-remnant depletion, canonical
% Hopf-current representative, exact Hopf descent to dOmega/(4pi), and
% closed-FLRW a^{-3} neutral-current shape.  It retypes K=48 as the onset
% of the final one-unit geometric transfer and records cDM/RAR as distinct
% physical projections of that common neutral parent while preserving
% occurrence, Hilbert-owner, stress-matching, and static-kernel premises.
% Release v1.59 transports the v1.58 exceptional-shell/K16 closure through
% Papers III-F and VI-B--VII-C.  It replaces Paper VI-B's obsolete mixed-Hessian
% output typing by the third-variation pair vertex and Bogoliubov occupation
% support, derives the post-seed Higgs-pair feedback vertex and its exact
% alpha--delta kernel factors, propagates the unique K16/K48 action-junction
% split, identifies the QCD delta count with the upstream 2O shell atom only
% as a common-scalar numerical compatibility, and narrows THERM16 to dynamical
% approach/beta selection because the finite Gibbs/KMS target is already
% unique on a supplied simple observable algebra once its Hamiltonian and
% inverse temperature are supplied; direct-sum central weights remain typed.
% Release v1.58 registers the exceptional nearest-shell closure, the derived
% transverse K=16 action event, the unique interior action junction, the
% type-correct ordered pair vertex, the thin-layer squeeze law, and the
% radiation-like/matter-like kinematic split.  REG_0, transfer magnitude,
% physical layer width, and thermalization remain explicitly conditional.
% Release v1.57 registers the active-depth-six obstruction, the native
% K=48 carrier desert, the typed 7|9 owner reduction, and the dynamical
% reheating handoff across Papers III-B--VI-B.
% Release v1.56 registers the least-support n_2 inflation dressing of the
% persistent n_7 Laurent response, the analytic global-convexity and positive
% endpoint-measure proofs, the derived inflation core and K=16 coefficient
% freeze, and the no-particle dark-sector split into static RAR plus the
% conditional homogeneous effective-curvature response at reflected K=48.
% There is no n_6 action, scalaron-DM mass, or second Einstein--Hilbert term.
% Release v1.55 registers the protected-octet resolved/residual reflection
% q_sp+Z_W=8, its j<->8-j and K<->64-K scale covariance, and the derived
% rank-one-exit quadratic grade C_1^2/U=C_2 in Papers III-C, III-E, and VIII.
% The reflection is a typed trace/scale theorem, not a time reversal or a
% physical event-occurrence theorem; reheating transfer remains conditional.
% Release v1.54 registers the exact positive-part Laurent indices, the
% active-domain/terminal-endpoint action split, the no-second-EH firewall, and
% the q_DM=6/K_DM=48 scalaron readout typing across Papers III-D/E/F, VI,
% VIII, and IX.
% Release v1.53 registers the common nonnegative Laurent obstacle projection
% for n_1^+ and n_7^+, discharges the three former exit/freeze assumptions,
% and advances the Paper-III-D/E/F, VIII, and IX consumers.  Shared notation
% and bibliography ownership are unchanged.
% Release v1.51 registers the PQTI derivation batch and its updated
% Paper-III-A/C/D/E/F, VIII, and IX consumers.  Shared notation and
% bibliography ownership are unchanged.
% Release v1.50 advances the Paper-III-B AUX target to v1.40 after the
% rigorous audit repair; shared notation and bibliography ownership are unchanged.
% Release v1.49 advances the Paper-II-E/F AUX target to v3.14 after the
% minimal audit repair; shared notation and bibliography ownership are unchanged.
% Release v1.46 synchronizes Paper II-A's minimal frozen constructor-origin
% theorem and its five dependent AUX targets.  Shared notation is unchanged.
% Release v1.45 updates the Paper-II-B AUX target to v2.23 after the
% GP3-rigidity and appendix-placement repair; shared notation is unchanged.
% Release v1.44 advances the centralized Paper-I-A--VII-C AUX registry after the
% Paper-I-A v5.16 protected-completion scope repair.  Every current
% paper imports this one file; no paper-level opt-in or \externaldocument list
% is permitted.  The bibliography remains centralized here through Born-v6_13.
% Shared bibliography, theorem infrastructure, colors, notation alignment,
% and current Paper-I-A--VII-C external-reference registration.
% This file centralizes common imports, bibliography sorting, colors, theorem
% infrastructure, and the notation shared by the complete series.
% v2.9 (release v1.43) made the Paper-I-A--VII-C AUX registry automatic and
% centralized the bibliography and all external-reference ownership.
% v2.8 release-registry refresh (v1.33) registers the rooted generated-parent
% cascade and the POL/REG/free-law/11+4/Born-group closure revisions.  No
% shared mathematical notation or Paper-II-A fundamental principle is changed.
% v2.8 release-registry refresh (v1.32) registers the universal generated-
% continuation revisions in Papers II-A, III-A, III-C, IV, and VIII.  No
% shared mathematical notation or Paper-II-A fundamental principle is changed.
% v2.8 release-registry refresh (v1.30) registers the CRP4/projective-state
% source revisions and their direct registry-only consumers.  No shared
% mathematical notation or Paper-II-A axiom is changed.
% v2.7 release-registry refresh (v1.27) updates only the current external
% Paper-I-A--VII-C basenames after the placement/assumption synchronization; no
% shared mathematical notation or Paper-II-A axiom is changed.
% v2.7 synchronizes the geometric defect-repair revision: Paper-II-A v3.42,
% Paper-III-C v1.26, Paper-III-D v1.27, Paper-III-E v1.32, and Paper-III-F
% v1.34.  It registers the common trace-defect/spectral-filtration theorem,
% its owner-typed response pullback, and the reduced GRED and TERM interfaces.
% No shared mathematical notation is changed.
% v2.6 synchronizes Paper-III-E v1.31's cone-type repair, minimal-response
% classification, and explicit classical action-role boundary, together with
% Paper-III-F v1.32's downstream response-atlas synchronization.  No shared
% mathematical notation is changed.
% v2.5 synchronizes the Paper-III-D v1.26 faithful-trace-duality proof,
% exact ordered-owner joint locus, and typed alternative-presentation
% bookkeeping repair, together with its Paper-III-E and closure-frontier
% references.  No shared mathematical notation is changed.
% v2.4 synchronizes the Paper-III-C v1.25 orbit/process/owner repair, its
% downstream Papers III-B--F, V, VI, VIII, and IX revisions, and Born-v6_13.
% v2.3 synchronizes the Paper-II-A v3.41 witness/support/principal inheritance
% layer, the Paper-III-A v1.24 joint-realizability/readout firewall, all dependent
% host-rate revisions, and the current Paper-I-A--VII-B source registry.
% v2.2 synchronizes the generated active-realization, determinant-depth,
% dimension-weighted hierarchy, and typed cosmological-event interfaces.
% v2.1 synchronizes the exceptional-support interface revisions: Paper I-A
% v5.9, Paper II-A v3.30, Paper II-E/F v3.7, and Paper IV v4.16.
% v2.0 introduces the stage-safe G0 occurrence notation, synchronizes the
% active Paper-I-A--IV source registry (current rigorous Paper-III closure:
% III-A v1.17, III-B v1.22, III-C--E v1.17, III-F v1.14, IV v4.14), and makes
% external-reference self-skip depend on logical paper role rather than a
% versioned job name.
% ======================================================================

% ======================================================================
% macros.tex — Shared macros for the Born Geometry series
% ======================================================================

% ======================================================================
% 1. PACKAGES, BIBLIOGRAPHY, COLORS, AND UNIT DEFINITIONS
% ======================================================================
% Core packages shared across the complete Paper-I-A--VII-C series.  Loading them here makes
% the macro file a genuine series preamble while remaining compatible with existing
% paper preambles that load the same packages before \input.
\usepackage{amsmath}
\usepackage{amssymb}
\usepackage{amsthm}
\usepackage{mathrsfs}
\usepackage{array}
\usepackage{booktabs}
\usepackage{xcolor}
\makeatletter
\@ifpackageloaded{geometry}{}{\usepackage{geometry}\geometry{margin=1in}}
\makeatother
\usepackage{mdframed}

% Shared support packages required by the macros below and by most papers.
\usepackage{ragged2e}
\usepackage{longtable}
\usepackage{tabularx}
\usepackage{enumitem}
\usepackage[most]{tcolorbox}

% Series-wide numeric bibliography convention using the portable BibTeX path.
% All Papers I-A--VII-C use the ordinary \cite interface.
\usepackage[numbers,sort&compress]{natbib}
\providecommand{\printbibliography}{%
  \bibliographystyle{plainnat}%
  \bibliography{Born-v6_60}%
}

% Shared colors.  \providecolor avoids clobbering a paper-specific override.
\providecolor{postulatecolor}{RGB}{230,242,255}
\providecolor{compare}{RGB}{230,245,230}
\providecolor{definitioncolor}{RGB}{230,242,255}

\providecolor{openriskcolor}{RGB}{255,235,235}
\providecolor{bridgecolor}{RGB}{240,240,255}
\providecolor{gold}{RGB}{255,215,0}

% ======================================================================
% 2. CUSTOM ENVIRONMENTS
% ======================================================================
% Common table column types used in all papers.  Guarded so that a paper may
% define its own versions before loading this file.
\makeatletter
\@ifundefined{NC@find@L}{%
  \newcolumntype{L}[1]{>{\RaggedRight\arraybackslash}p{#1}}%
}{}
\@ifundefined{NC@find@C}{%
  \newcolumntype{C}[1]{>{\Centering\arraybackslash}p{#1}}%
}{}
\@ifundefined{NC@find@Y}{%
  \newcolumntype{Y}{>{\RaggedRight\arraybackslash}X}%
}{}
\makeatother

% ======================================================================
% 3. THEOREM STYLES AND ENVIRONMENTS
% ======================================================================
% Custom underlined theorem style (useful for remarks or notes)
\newtheoremstyle{underlined}
  {3pt}   % Space above
  {3pt}   % Space below
  {}      % Body font (empty means normal/roman text, typical for remarks)
  {}      % Indent amount
  {}      % Theorem head font (can be \itshape, \bfseries, etc.)
  {.}     % Punctuation after theorem head
  { }     % Space after theorem head
  {\underline{\thmname{#1}\thmnumber{ #2}}\thmnote{ (#3)}} % Theorem head formatting

% Standard plain theorem style environments.  Each environment is guarded so
% the shared macro file can be loaded after a paper-specific theorem setup.
\theoremstyle{plain}
\makeatletter
\@ifundefined{theorem}{\newtheorem{theorem}{Theorem}[section]}{}
\@ifundefined{lemma}{\newtheorem{lemma}{Lemma}[section]}{}
\@ifundefined{proposition}{\newtheorem{proposition}{Proposition}[section]}{}
\@ifundefined{corollary}{\newtheorem{corollary}{Corollary}[section]}{}
\makeatother

% Definition style environments (typically non-italicized body)
\theoremstyle{definition}
\makeatletter
\@ifundefined{definition}{\newtheorem{definition}{Definition}[section]}{}
\@ifundefined{postulate}{\newtheorem{postulate}{Postulate}[section]}{}
\@ifundefined{closure}{\newtheorem{closure}{Closure}[section]}{}
\makeatother

% Applying the custom underlined style
\theoremstyle{underlined}
\makeatletter
\@ifundefined{remark}{\newtheorem{remark}{Remark}[section]}{}
\@ifundefined{convention}{\newtheorem{convention}{Convention}[section]}{}
\@ifundefined{principle}{\newtheorem{principle}{Principle}[section]}{}
\makeatother

% Framed proof environment.
% Every \begin{proof} ... \end{proof} is automatically framed/shaded.
% This avoids mdframed/trivlist nesting, which can drop the top border after
% theorem-like environments.
\makeatletter
\renewenvironment{proof}[1][\proofname]{%
  \par\addvspace{6pt}%
  \noindent
  \begin{tcolorbox}[%
    enhanced,
    breakable,
    colback=black!4,
    colframe=black!35,
    boxrule=0.35pt,
    arc=1pt,
    left=6pt,
    right=6pt,
    top=4pt,
    bottom=4pt,
    boxsep=0pt,
    before skip=0pt,
    after skip=6pt,
    parbox=false
  ]%
  \pushQED{\qed}%
  \normalfont\itshape #1\@addpunct{.}\enspace
  \normalfont\ignorespaces
}{%
  \unskip\nobreak\hfill\popQED
  \end{tcolorbox}%
  \par\addvspace{2pt}%
}
\makeatother

\mdfdefinestyle{postulatebox}{
  backgroundcolor=yellow!10,
  linecolor=black,
  linewidth=0.4pt,
  skipabove=\baselineskip,
  skipbelow=\baselineskip,
  nobreak=true
}

% ======================================================================
% PROOF-READABILITY HELPERS
% ======================================================================

\newcommand{\ProofPlan}[1]{%
  \par\smallskip
  \noindent\textbf{Proof plan. }#1%
  \par\smallskip
}

\newlist{proofsteps}{enumerate}{2}
\setlist[proofsteps,1]{
  label=\textbf{(\alph*)},
  ref=\alph*,
  leftmargin=2.2em,
  itemsep=0.45\baselineskip,
  topsep=0.35\baselineskip
}
\setlist[proofsteps,2]{
  label=\textbf{(\roman*)},
  ref=\roman*,
  leftmargin=2.0em,
  itemsep=0.25\baselineskip,
  topsep=0.2\baselineskip
}


% ======================================================================
% 4. NUMBER SYSTEMS AND ALPHABETS
% ======================================================================
\providecommand{\RR}{\mathbb{R}} % Real numbers
\providecommand{\PP}{\mathbb{P}} % Projective space / Probability
\providecommand{\CC}{\mathbb{C}} % Complex numbers
\providecommand{\ZZ}{\mathbb{Z}} % Integers
\providecommand{\NN}{\mathbb{N}} % Natural numbers
\providecommand{\HH}{\mathbb{H}} % Quaternions / Hyperbolic space
\providecommand{\LFRayGeom}{\mathsf{LF\!\mbox{-}RAY\!\mbox{-}GEOM}}
\providecommand{\LFRep}{\mathsf{LF\!\mbox{-}REP}}
\providecommand{\LFState}{\mathsf{LF\!\mbox{-}STATE}}
\providecommand{\LFLim}{\mathsf{LF\!\mbox{-}LIM}}
\providecommand{\LFScat}{\mathsf{LF\!\mbox{-}SCAT}}
% Number systems / projective spaces
\providecommand{\CP}{\mathbb{CP}}
\providecommand{\RP}{\mathbb{RP}}

% ======================================================================
% 5. GROUPS AND LIE THEORY
% ======================================================================
\providecommand{\SO}{\ensuremath{\mathrm{SO}}}     % Special Orthogonal group
\providecommand{\SU}{\ensuremath{\mathrm{SU}}}     % Special Unitary group
\providecommand{\Spin}{\ensuremath{\mathrm{Spin}}} % Spin group
\providecommand{\CS}{\ensuremath{\mathrm{CS}}}     % Chern-Simons (capitalized)
\providecommand{\cs}{\ensuremath{\mathrm{cs}}}     % Chern-Simons (lowercase)

% ======================================================================
% 6. COMMON MANIFOLDS AND TOPOLOGY
% ======================================================================
\providecommand{\Sone}{S^1}                 % Standard circle

% ======================================================================
% 7. MATH OPERATORS (Declared via \DeclareMathOperator)
% ======================================================================
\DeclareMathOperator{\End}{End}   % Endomorphisms
\DeclareMathOperator{\Aut}{Aut}   % Automorphisms
\DeclareMathOperator{\spec}{spec} % Lowercase spectrum
\DeclareMathOperator{\tr}{tr}     % Lowercase trace (if you already have \Tr)
\DeclareMathOperator{\supp}{supp} % Support
\DeclareMathOperator{\TrA}{Tr_{A}} % Partial trace over subsystem A
\DeclareMathOperator{\TrB}{Tr_{B}} % Partial trace over subsystem B
% Lie/algebra operators
\DeclareMathOperator{\Ad}{Ad}
\DeclareMathOperator{\Sym}{Sym}
\DeclareMathOperator{\su}{su}
% ======================================================================
% 8. STANDARD OPERATOR-LIKE SYMBOLS (Safe definitions via \providecommand)
% ======================================================================
\providecommand{\vol}{\operatorname{vol}}   % Volume form
\providecommand{\Vol}{\operatorname{Vol}}   % Total volume
\providecommand{\Ric}{\operatorname{Ric}}   % Ricci tensor
\providecommand{\Rm}{\operatorname{Rm}}     % Riemann tensor
\providecommand{\Tr}{\operatorname{Tr}}     % Uppercase trace
\providecommand{\Spec}{\operatorname{Spec}} % Uppercase spectrum
\providecommand{\diag}{\operatorname{diag}} % Diagonal matrix
\providecommand{\rank}{\operatorname{rank}} % Rank
\providecommand{\scal}{\mathcal R} % Scalar curvature
\providecommand{\Cl}{\operatorname{Cl}}     % Clifford algebra
\providecommand{\diam}{\operatorname{diam}} % Diameter

% ======================================================================
% 9. FUNCTIONS, TRIGONOMETRY, AND CALCULUS
% ======================================================================

% Codifferential naming conventions
% (Avoid \ddag: commonly defined by physics.sty and symbol packages)
\providecommand{\codiff}{\mathrm{d}^{\dagger}} % Codifferential operator

% Quantum mechanics bra-ket notation (needed for basis tables)
\providecommand{\ket}[1]{\lvert #1 \rangle}
\providecommand{\bra}[1]{\langle #1 \rvert}
\providecommand{\braket}[2]{\langle #1 \mid #2 \rangle}

% ======================================================================
% 10. PHYSICAL SCALES AND CONSTANTS (Paper 1c / Born Geometry)
% ======================================================================
\providecommand{\Rp}{\ensuremath{R_{\mathrm P}}}     % matched common EH curvature scale (not defined from observed G)
\providecommand{\hBorn}{\ensuremath{h_{\mathrm B}}} % physical representative 2\pi\hbar of one primitive Born scalar action cell
\providecommand{\UEHjoint}{\ensuremath{U_{J}}} % joint-Dirac EH curvature-coordinate scale 8 kappa_J
\providecommand{\Gread}{\ensuremath{G_{\rm scal}}} % scalar-action coordinate c_op^3/(h_B U); physical matching explicit
\providecommand{\GreadJoint}{\ensuremath{G_{J,\mathrm{raw}}}} % raw ordinary-trace coordinate = G_scal/96
\providecommand{\Rds}{R_{\mathrm{dS}}}    % de Sitter Curvature
\providecommand{\Mpl}{M_{\mathrm{Pl}}}    % Planck mass
\providecommand{\lp}{\ell_{\mathrm{Pl}}}  % Planck length
\providecommand{\Pmax}{P_{\mathrm{max}}}  % Maximum momentum
\providecommand{\Mplbar}{\bar{M}_{\mathrm{Pl}}} % Reduced Planck mass
\providecommand{\eff}{\mathrm{eff}}       % Subscript for effective quantities
\providecommand{\Cosmo}{\ensuremath{\Lambda}}                % Cosmological constant (use \Cosmo in source)
\providecommand{\LCDM}{\ensuremath{\Cosmo\mathrm{CDM}}}  % \LambdaCDM (typesets with \Cosmo)
\providecommand{\Lcut}{\ensuremath{\Lambda_{\mathrm{cut}}}}  % Cutoff scale (brace when used as script)
\providecommand{\Rcs}{\ensuremath{R_{\mathrm{cs}}}}      % Configuration-space radius
\providecommand{\CosmoEff}{\ensuremath{\Cosmo_{\eff}}}   % Effective cosmological constant

% Disambiguation: \Rmom is the momentum-space (S^3) radius. 
% Reserve plain R for the 4D curvature scalar in f(R).
\providecommand{\Rmom}{\ensuremath{R_{\mathrm{mom}}}} 
\providecommand{\Punit}{P_{0}} % Momentum unit set by \Rmom (= \hbar/\Rmom)
\providecommand{\Rvec}{\ensuremath{R_{\mathrm{vec}}}} % invariant vector-space radius in the reduced round sector
% Specific numeric values and expressions
\providecommand{\Vlogcap}{L_{\mathrm{Hopf}}^{(1)}}
\providecommand{\Hopf}{\operatorname{Hopf}}
% Scalar phase (S^1) sector + universal cover notation (Paper I-A)
\providecommand{\tper}{T_\star}      % Time period/scale
\providecommand{\Eper}{E_\star}      % Energy period/scale
\providecommand{\ttilde}{\tilde t}   % Universal cover time
\providecommand{\Etilde}{\tilde E}   % Universal cover energy

% ======================================================================
% 10A. SERIES-WIDE LOGICAL / NOTATIONAL DISAMBIGUATIONS
% ======================================================================
% Use these symbols to avoid collisions among: (i) the Paper--III
% no-insertion/geometric-closure class, (ii) Paper--I-A trimer diagonals,
% (iii) Paper--V-A/V-C family-space diagonals, (iv) scalar-index running,
% and (v) protected shell-clock labels.
\providecommand{\Rprim}{\ensuremath{\mathsf R_{\mathrm{prim}}^{\Sigma}}} % primitive finite structural readouts
\providecommand{\Rloop}{\ensuremath{\mathsf R_{\circ}^{\Sigma}}} % closed-loop readout role
\providecommand{\Rbdry}{\ensuremath{\mathsf R_{\partial}^{\Sigma}}} % boundary-conditioned readout role
\providecommand{\Rgen}{\ensuremath{\mathsf R_{\mathrm{fin}}^{\Sigma}}} % finite generated closure of primitive roles
\providecommand{\PhysOcc}{\operatorname{PhysOcc}}
\providecommand{\Actual}{\operatorname{Actual}}
\providecommand{\GeomLic}{\operatorname{GeomLic}}
\providecommand{\InvDist}{\operatorname{InvDist}}
\providecommand{\Dcl}{\ensuremath{\mathsf{D}_1}}      % unmarked positive Dirac-even closed spectral-readout form class
\providecommand{\Dheat}{\ensuremath{\mathsf D_{\mathrm{heat}}}} % one-parameter heat-semigroup subclass
\providecommand{\DtriI}{\ensuremath{D^{\mathrm{tri}}_1}} % Paper--I-A cyclic-trimer diagonal
\providecommand{\DtriII}{\ensuremath{D^{\mathrm{tri}}_2}}
\providecommand{\DfamI}{\ensuremath{D^{\mathrm{fam}}_1}} % Family-space diagonal selector
\providecommand{\DfamII}{\ensuremath{D^{\mathrm{fam}}_2}}
\providecommand{\alpharun}{\ensuremath{\alpha_{\mathrm{run}}}} % scalar spectral-index running
\providecommand{\Ksh}{\ensuremath{K_{\mathrm{sh}}}} % protected shell-clock coordinate
\providecommand{\ddef}{\ensuremath{\delta_{\mathrm{def}}}} % universal defect scalar after Paper--III realization

% ======================================================================
% 10B. SERIES-WIDE LOGICAL STATUS TAGS
% ======================================================================
% Upright small-caps wrapper for audit/status labels.  It prevents LaTeX from
% trying to use unavailable small-caps italic font shapes inside theorem bodies
% such as TU/lmr/m/scit.
\providecommand{\AuditLabel}[1]{}
\renewcommand{\AuditLabel}[1]{{\normalfont\scshape #1}}

% These tags are bookkeeping labels for consistency audits. They do not
% create new mathematical hypotheses; they record how a datum is being used
% in the proof chain.
\providecommand{\StatusPrimitive}{}
\renewcommand{\StatusPrimitive}{\AuditLabel{Primitive}}
\providecommand{\StatusDerived}{}
\renewcommand{\StatusDerived}{\AuditLabel{Derived}}
\providecommand{\StatusImported}{}
\renewcommand{\StatusImported}{\AuditLabel{Imported}}
\providecommand{\StatusConditional}{}
\renewcommand{\StatusConditional}{\AuditLabel{Conditional}}
\providecommand{\StatusBridge}{}
\renewcommand{\StatusBridge}{\AuditLabel{Bridge-derived}}
\providecommand{\StatusComparison}{}
\renewcommand{\StatusComparison}{\AuditLabel{Comparison-only}}

% ======================================================================
% 10C. SERIES-WIDE CLAIM-WORDING AUDIT LABELS
% ======================================================================
% These are textual audit labels used to keep modal language consistent
% across the series.  They introduce no new mathematical assumptions.
\providecommand{\ClaimForced}{}
\renewcommand{\ClaimForced}{\AuditLabel{forced}}
\providecommand{\ClaimUnique}{}
\renewcommand{\ClaimUnique}{\AuditLabel{unique}}
\providecommand{\ClaimCanonical}{}
\renewcommand{\ClaimCanonical}{\AuditLabel{canonical}}
\providecommand{\ClaimMinimal}{}
\renewcommand{\ClaimMinimal}{\AuditLabel{minimal}}
\providecommand{\ClaimClosed}{}
\renewcommand{\ClaimClosed}{\AuditLabel{closed}}
\providecommand{\ClaimPrediction}{}
\renewcommand{\ClaimPrediction}{\AuditLabel{prediction}}

% ======================================================================
% 10D. SERIES-WIDE THEOREM-CHAIN AUDIT LABELS
% ======================================================================
% These labels mark the direction of load-bearing theorem chains. They are
% audit annotations only; they introduce no new mathematical assumptions.
\providecommand{\ChainInput}{}
\renewcommand{\ChainInput}{\AuditLabel{Chain input}}
\providecommand{\ChainOutput}{}
\renewcommand{\ChainOutput}{\AuditLabel{Chain output}}
\providecommand{\ChainGuard}{}
\renewcommand{\ChainGuard}{\AuditLabel{Chain guard}}
\providecommand{\ChainNoBackflow}{}
\renewcommand{\ChainNoBackflow}{\AuditLabel{No backflow}}

% ======================================================================
% 10E. SERIES-WIDE LOCAL THEOREM-AUDIT LABELS
% ======================================================================
% These labels mark a statement-level dependency audit around individual
% theorem/definition nodes. They are bookkeeping labels only.
\providecommand{\LocalAudit}{}
\renewcommand{\LocalAudit}{\AuditLabel{Local audit}}
\providecommand{\UsesOnly}{}
\renewcommand{\UsesOnly}{\AuditLabel{Uses only}}
\providecommand{\Exports}{}
\renewcommand{\Exports}{\AuditLabel{Exports}}
\providecommand{\DoesNotUse}{}
\renewcommand{\DoesNotUse}{\AuditLabel{Does not use}}

% ======================================================================
% 10F. SERIES-WIDE NUMERICAL-ROW AUDIT LABELS
% ======================================================================
% These labels classify individual rows in numerical, benchmark, status,
% and comparison tables.  They are audit annotations only and do not change
% any theorem, value, formula, or comparison number.
\providecommand{\RowAudit}{}
\renewcommand{\RowAudit}{\AuditLabel{Numerical-row audit}}
\providecommand{\RowTheorem}{}
\renewcommand{\RowTheorem}{\AuditLabel{Theorem-fixed row}}
\providecommand{\RowBridge}{}
\renewcommand{\RowBridge}{\AuditLabel{Bridge-fixed row}}
\providecommand{\RowRepresentative}{}
\renewcommand{\RowRepresentative}{\AuditLabel{Representative row}}
\providecommand{\RowComparison}{}
\renewcommand{\RowComparison}{\AuditLabel{Comparison row}}
\providecommand{\RowCalibrated}{}
\renewcommand{\RowCalibrated}{\AuditLabel{Calibrated row}}
\providecommand{\RowEstimate}{}
\renewcommand{\RowEstimate}{\AuditLabel{Estimate row}}
\providecommand{\RowOpen}{}
\renewcommand{\RowOpen}{\AuditLabel{Open row}}
\providecommand{\RowNoSelector}{}
\renewcommand{\RowNoSelector}{\AuditLabel{No selector}}

% ======================================================================
% 10G. TABLE HELPERS
% ======================================================================
% Table helpers used by Paper V-C.  The unused readout-label preset has been
% removed; only active table helpers remain.
\providecommand{\gpsection}[1]{%
  \addlinespace[0.55em]
  \multicolumn{3}{@{}l}{\textsc{#1}}\\[-0.25em]
  \midrule
}
\providecommand{\mathcell}[1]{%
  \(\begin{gathered}#1\end{gathered}\)%
}

\newcommand{\DerBy}[1]{%
  \begin{tabular}[t]{@{}l@{}}#1\end{tabular}%
}

\newcommand{\TableNotes}[1]{%
  \par\medskip
  \begin{minipage}{0.94\linewidth}
  \footnotesize\emph{Notes.} #1
  \end{minipage}\par
}
% ======================================================================
% 10H. PAPER VI-A / QCD NOTATION
% ======================================================================
% Shared symbols used by Paper VI-A and downstream QCD/Yang--Mills handoffs.
\providecommand{\Pqcd}{\mathfrak P_{\rm VIII}^{\rm QCD}}
\providecommand{\Iqcd}{\mathfrak I_{\rm QCD}^{(7)}}
\providecommand{\muQCD}{\mu_{\rm QCD}^{\rm int}}
\providecommand{\LambdaQCD}{\Lambda_{\rm QCD}}
\providecommand{\LambdaBH}{\Lambda_{\rm BH}^{\rm QCD}}
\providecommand{\HQCD}{\mathcal H_{\rm QCD}^{\rm BH}}
\providecommand{\HQCDfin}{H_{\rm QCD}^{\rm fin}}
\providecommand{\Hraw}{\mathcal H_{\rm raw}}
\providecommand{\HphysQCD}{\mathcal H_{\rm phys}^{\rm QCD}}
\providecommand{\Vcol}{V_C}
\providecommand{\Ps}{P_{\rm sing}}
\providecommand{\PiOct}{\Pi_8}
\providecommand{\MSbar}{\overline{\rm MS}}
\providecommand{\YM}{\mathrm{YM}}
\providecommand{\Ktri}{K_{\triangle}}
\providecommand{\KYM}{K_{\YM}^{\rm BH}}
\providecommand{\aBH}{a_{\rm BH}}
\providecommand{\epsBH}{\varepsilon_{\rm BH}}
\newcommand{\RowConstructive}{\textnormal{\textsc{Constructive}}}
\newcommand{\RowObligation}{\textnormal{\textsc{Bridge obligation}}}
\newcommand{\RowTarget}{\textnormal{\textsc{Target row}}}
\newcommand{\RowValidation}{\textnormal{\textsc{Validation row}}}
\providecommand{\slashed}[1]{\not{#1}}

% ======================================================================
% 10I. PAPER VI-B / CONTINUOUS-QFT NOTATION
% ======================================================================
% Shared symbols used by Paper VI-B and any downstream continuous-QFT or
% Yang--Mills handoff paper.  All definitions are guarded so individual
% papers may override them locally before loading this file if needed.
\providecommand{\A}{\mathfrak A}
\providecommand{\Hphys}{\mathcal H_{\rm phys}}
\providecommand{\QCD}{\mathrm{QCD}}
\providecommand{\BH}{\mathrm{BH}}
\providecommand{\loc}{\mathrm{loc}}

\providecommand{\RBH}{\mathcal R_{\rm BH}}
\providecommand{\DG}{\mathfrak D_G}
\providecommand{\PW}{\mathcal{P\!W}} % Peter--Weyl spectral-symbol subspace
\providecommand{\Toe}{\operatorname{Toe}} % Toeplitz compression/representative
\providecommand{\Hol}{\operatorname{Hol}} % Holonomy

% ======================================================================
% 11. LATE-LOADED PACKAGES (Must be at the end of the preamble)
% ======================================================================
% Load late packages only if a paper has not already loaded them, avoiding
% option clashes in paper-specific preambles.
\makeatletter
\@ifpackageloaded{hyperref}{}{\usepackage{hyperref}}
\makeatother
% --- Series-wide disambiguations (Pass 3) ---
\providecommand{\diracdelta}{\delta_{\rm D}} % Dirac delta distribution (delta reserved for hierarchy)
\providecommand{\Hg}{H_{\mathrm g}} % Hodge harmonic projector (keep H for Hubble)
\providecommand{\Gr}{G_{\mathrm r}} % Green operator (keep G for Newton)
\providecommand{\FD}{\mathrm{D}} % functional/variational derivative symbol (delta reserved)
\providecommand{\Id}{\mathbf{1}} % identity/Kronecker tensor (avoid \delta_{ij})


% BEGIN INLINE SOURCE: Paper00-ActiveExternalReferences-v2_130.tex
% Framework v2.191: current reference binding; existing definitions retained.
% Framework v2.190: current22-root registry with native clock-evolution proofs.
% Framework v2.189: exact current22-root external registry, including retained v2.188 hardening owners.
% Framework v2.187: current registry synchronization only; source release v2.186 preserved.
% Framework v2.186: authoritative current 22-paper registry.
% Framework v2.182: all current 22 logical paper AUX bindings.
% Framework v2.170: current I-B v1.43 external AUX binding.
% Framework v2.167: exact current 22-root registry after CP--BL2 integration.
% Framework v2.166: exact current 22-root compact-release AUX registry.
% Framework v2.165: authoritative current root bindings.
% Framework v2.164: exact current 22 manuscript AUX bindings.
% Framework v2.162: finite metric, functional Hilbert and scoped Ward integration.
% Framework v2.161: current 22 manuscript AUX bindings; source only.
% Framework v2.156 / registry v2.113: synchronized current 22 manuscript roots; source-only registry.
% Framework v2.143 / registry v2.111: active manuscript AUX basenames; no generated AUX or PDF included.
% ======================================================================
% Central external-reference imports, numbers only, no PDF jumping
% ======================================================================

% xr-hyper is loaded centrally so that papers may retain their existing
% \usepackage{xr-hyper} lines without redefining \externaldocument after this
% registry is installed.  If plain xr was loaded earlier, retain it.
\makeatletter
\@ifpackageloaded{xr-hyper}{}{%
  \@ifpackageloaded{xr}{}{\usepackage{xr-hyper}}%
}
\makeatother

% Make external-document registration idempotent.  The registry identifies a
% logical paper by its Paper<number><optional letter> prefix.  Prefix-qualified
% imports remain independent.  Current paper roots contain no local
% \externaldocument calls; this guard also makes accidental duplicates harmless.
\ExplSyntaxOn
\prop_new:N \g_bh_external_refs_prop
\cs_new_eq:NN \BHOriginalExternalDocument \externaldocument

\cs_new_protected:Npn \bh_external_document:nnnn #1#2#3#4
 {
  \tl_set:Nn \l_tmpa_tl {#3}
  \regex_replace_once:nnN
   { \A (?: xref/ )? (Paper[0-9]+[A-G]?) .* \Z }
   { \1 }
   \l_tmpa_tl
  \tl_set:Nx \l_tmpb_tl { #1 | \tl_use:N \l_tmpa_tl }
  \prop_if_in:NVTF \g_bh_external_refs_prop \l_tmpb_tl
   {
    \PackageInfo{BornHopf}
     {Skipping~duplicate~external~refs~for~#3}
   }
   {
    \prop_gput:NVn \g_bh_external_refs_prop \l_tmpb_tl {#3}
    \IfNoValueTF{#4}
     { \BHOriginalExternalDocument[#1][#2]{#3} }
     { \BHOriginalExternalDocument[#1][#2]{#3}[#4] }
   }
 }

\RenewDocumentCommand{\externaldocument}{O{} O{} m o}
 { \bh_external_document:nnnn {#1}{#2}{#3}{#4} }

\NewDocumentCommand{\BHExternalRefs}{m}
 {
  % Do not import the aux file of the logical paper currently being compiled,
  % even when the job and registry use different version suffixes.
  \tl_set:Nx \l_tmpa_tl { \c_sys_jobname_str }
  \regex_replace_once:nnN
   { \A (Paper[0-9]+[A-G]?) .* \Z }
   { \1 }
   \l_tmpa_tl
  \tl_set:Nn \l_tmpb_tl { #1 }
  \regex_replace_once:nnN
   { \A (Paper[0-9]+[A-G]?) .* \Z }
   { \1 }
   \l_tmpb_tl
  \str_if_eq:eeTF { \tl_use:N \l_tmpa_tl } { \tl_use:N \l_tmpb_tl }
   {
    \PackageInfo{BornHopf}
     {Skipping~external~refs~for~#1~because~this~is~the~current~document}
   }
   {
    \IfFileExists{xref/#1.aux}
     { \externaldocument[][nocite]{xref/#1}[#1.pdf] }
     {
      \IfFileExists{#1.aux}
       { \externaldocument[][nocite]{#1} }
       {
        \PackageInfo{BornHopf}
         {External~aux~file~#1.aux~not~found;~no~labels~imported}
       }
     }
   }
 }
\ExplSyntaxOff

% External references are automatic.  This file is the sole owner of both the
% bibliography and the current Paper-I-A--VII-D AUX registry (30 roots including I-C).  Missing auxiliaries
% are informational on the first pass and are populated by the source-only
% two-pass series build.  Basenames deliberately omit temporary download/copy
% suffixes such as ``(1)'' or ``-1''.
  % Paper I-A: round S3 compendium.
  \BHExternalRefs{Paper1A-Geometry-v5_181}

  % Paper I-B: complete exceptional geometry and finite spectral structures.
  \BHExternalRefs{Paper1B-ExceptionalGeometry-v1_27}
  % Paper I-D: technical companion of I-A through I-C.
  \BHExternalRefs{Paper1D-TechnicalMaterial-v1_69}
  % Paper I-C: reciprocal M8 compendium.
  \BHExternalRefs{Paper1C-M8Geometry-v1_146}

\ifdefined\BHMathOnly
  % Papers I-A and I-B import only the mathematical atlas labels.
\else
  % Paper II: five physical-mathematical main papers plus one technical proof volume.
  \BHExternalRefs{Paper2A-GeometricConstitution-v3_227}
  \BHExternalRefs{Paper2B-MomentumCarrier-v2_147}
  \BHExternalRefs{Paper2C-ReciprocalLocalization-v2_177}
  \BHExternalRefs{Paper2D-SymmetrySpectra-v2_145}
  % The former II-F export labels remain in II-E; current II-F is the technical volume.
  \BHExternalRefs{Paper2E-KinematicExports-v3_140}
  % Paper II-F: complete technical proof and calculation owner for II-A--II-E.
  \BHExternalRefs{Paper2F-TechnicalMaterial-v1_71}

  % Paper III: Dynamics and Intrinsic Geometric Constitution (six components).
  \BHExternalRefs{Paper3A-GeneratedDynamics-v1_181}
  \BHExternalRefs{Paper3B-SpectralAction-v1_192}
  \BHExternalRefs{Paper3C-ProtectedHierarchy-v1_196}
  \BHExternalRefs{Paper3D-RunningResponse-v1_168}
  \BHExternalRefs{Paper3E-GravityAction-v1_204}
  \BHExternalRefs{Paper3F-CompleteDynamics-v1_281}
  % Paper III-G: complete preserved technical proof owner.
  \BHExternalRefs{Paper3G-TechnicalMaterial-v1_58}

  % Paper IV-A: generated finite quantum mechanics.
  \BHExternalRefs{Paper4A-FiniteQuantumMechanics-v4_217}
  % Paper IV-B: complete finite quantum processes and physical response.
  \BHExternalRefs{Paper4B-QuantumResponse-v1_55}
  % Paper IV-C: extended proofs and finite technical calculations.
  \BHExternalRefs{Paper4C-TechnicalMaterial-v1_50}

  % Paper V-A: The Standard Model.
  \BHExternalRefs{Paper5A-FiniteMatterGeometry-v3_219}

  % Paper V-B: Finite Spectrum and Flavor Closure.
  \BHExternalRefs{Paper5B-TerminalSpectrum-v5_202}

  % Paper V-C: Finite Interaction Boundaries of the Standard Model.
  \BHExternalRefs{Paper5C-InteractionResponse-v2_185}
  % Paper V-D: preserved technical proof and calculation owner.
  \BHExternalRefs{Paper5D-TechnicalMaterial-v1_58}

  % Paper VI-A: Gravity and Closed Cosmology.
  \BHExternalRefs{Paper6A-GravityCosmology-v2_268}

  % Paper VI-B: Geometric Dark Sector and Cosmological Readouts.
  \BHExternalRefs{Paper6B-CosmologyReadouts-v1_215}

  % Paper VI-C: complete technical proof and comparison owner.
  \BHExternalRefs{Paper6C-TechnicalMaterial-v1_53}

  % Paper VII-A: Finite QCD Kernel and Nonperturbative Closure.
\BHExternalRefs{Paper7A-QCDKernel-v1_185}

  % Paper VII-B: Hadronic, Thermal, and Nuclear QCD Readouts.
  \BHExternalRefs{Paper7B-QCDReadouts-v1_169}

  % Paper VII-C: Continuous Quantum Field Theory.
  \BHExternalRefs{Paper7C-ContinuousQFT-v2_208}

  % Paper VII-D: complete technical material for VII-A through VII-C.
  \BHExternalRefs{Paper7D-TechnicalMaterial-v1_60}

\fi

% END INLINE SOURCE: Paper00-ActiveExternalReferences-v2_130.tex


\usepackage{tocloft}
\setlength{\cftsecnumwidth}{3.4em}
\setlength{\cftsubsecnumwidth}{4.3em}
\setlength{\cftsubsubsecnumwidth}{5.3em}
\hypersetup{hidelinks}


\usepackage{xr-hyper}

\title{\textbf{Born--Hopf Geometry III-F: Complete Local Dynamics, Sources, and Boundary Response}\\[0.5em]
\large Full metric work, physical source variations, and causal response}

\author{Peter Kugelmann\thanks{\texttt{pkugelmann@proton.me}}}
\date{Version 1.281 --- 20 November 2026 \\ Framework v3.077}
\setlength{\emergencystretch}{8em}
\begin{document}
\maketitle
\begin{abstract}
Paper III-F supplies the conditional complete local dynamical
and source-action realization: cotangent localization, the joint spectral
connection, metric/Noether variation, finite charged and neutral source
constructions, Schur and boundary terms, causal Cauchy evolution, and exact
response tests. Bare heat coefficients, local finite kernels, and positive
Schur providers are kept distinct from physically matched source work.
Crucial Lorentzian action, Euler, boundary, FLRW, and transverse-response
statements remain here; optional calculations are in Paper~III-G.
\end{abstract}
\tableofcontents
\clearpage

\section{Physical response target and constituted dynamical inputs}
\label{paper3f:sec:organized-scope}
\subsection{Introduction and Paper-III series allocation}
\label{p3f:sec:series-allocation}
Paper~III-F is the locality and backreaction closure of the Paper~III
sequence.  It imports the fixed momentum seed and reciprocal smooth-role
interfaces, the declared spectral-action class, the protected
Einstein--Hilbert normalization branch of Paper~III-E, and the current smooth
metric response \(\mathfrak f_{\rm sp}\).  It then proves the local doubled
geometry, the explicit polarization-dependent variational asymmetry, the
common-action source equation, and a conditional Born-spherical gauge--gravity
reduction.  No downstream particle, cosmological-epoch, or observational
assignment is used as a premise.

\subparagraph{Inherited geometric and finite-role prefix.}
The current upstream chain is
\[
\begin{gathered}
\text{Paper~II-B}\Longrightarrow
 \mathcal H_p\cong\mathcal H_1,\quad m_p=1,\\
\text{II-C/III-A/GP2}\Longrightarrow
 \text{smooth reciprocal role and Born involution},\\
\text{III-F evaluation duality}\Longrightarrow
 \text{scale-free cotangent target }T^*U,\\
+\text{generated self-adjoint operator owner}
 \xRightarrow{\text{III-B v1.73}}
 (\mathrm{CLOOP}_{\rm heat}^{\rm der}),\\
 \text{one-profile additive class }\Dcl,\\
\substack{\mathsf C_{\rm IR}^{\rm der}\\
 +(\mathrm{GRAV\!-\!ASYM})}
\xRightarrow{\text{III-E}}
 F_g(u)=e^{-u},\quad f_{2,g}=(4\pi)^{-2},\quad
 \Lambda_g^2=96\pi^2\Lambda_{\rm EH}^2,\\
\substack{\mathrm{REG}_{\sigma}\in\mathfrak A_{\rm REG}^{+}
 +(\mathrm{CRP4})\\
 +(\mathrm{TERM}_f)+(\mathrm{MRESP}_{\min})}
\xRightarrow{\text{III-E}}
 \mathfrak f_{\rm sp}(\scal),\quad
 \sigma_{\rm act}(\scal),\\
 w_{\rm nl}=\frac18\,8^{n_{\rm act}},\quad
 R_+=U=\Lambda_{\rm EH}^2,\\
 C=U\alpha_{\rm hier}(\sigma_{\rm fin})^{-7},\quad
 \mathcal K_p=6/\Rmom^2,\\
 U=(3/\pi)\mathcal K_p\quad\text{only under }(\mathrm{CELL}_{K\to U}),\\
 \Rp:=U\quad\text{(legacy alias)}.
\end{gathered}
\]
On the displayed heat branch, \((\mathrm{GRAV}_{\circ})\),
\((\mathrm{GRAV\!-\!HEAT})\), \((\mathrm{CORN})\), and
\((\mathrm{BD}_{\rm EH})\) are derived certifications rather than
additional premises.  The first two certify the generated process law; the
last two certify the protected scalar constituent and the Einstein
specialization of the natural scalar-density Born pullback.
Here \(\mathsf C_{\rm IR}^{\rm der}\) abbreviates
\((\mathrm{ND}_{\rm sp})\), the G0-occurring Hurwitz-even active orbit family,
and the orbit-invariant identity \(P_{\rm term}=P_0\).  Paper~II-C removes the nonzero
protected completion scale and relative chiral phase from the normalized even
algebra exactly; equivariant descent, finite functional calculus, and GP2
co-realization inherit the same G0 witness without choosing a diagonal, so
\((\mathrm{IR}_{O(3)})\) is only a certification of the
resulting scalar Haar projector.
The smooth reciprocal role spaces are not finite Hilbert carriers.  The term
\emph{carrier} below is reserved for actual finite representation carriers such
as \(\mathcal H_p\); local manifolds, bases, fibres, and phase spaces are named
as such.

\subparagraph{Geometric primacy G0.}
The legacy localization tag \((\mathrm{LOC}^{\rm ev}_{xp})\) is now a
derived certification of the Paper-III-A generator-germ theorem and adds no
G0 root.  The occurring action parent and its derived variational domain
\((\mathrm{VDOM}_x^{\min})\), the derived joint datum
\((\mathrm{JOINT}_x^{\rm der})\), the generated stationary-family role
\((\mathrm{STATFAM}_x^{\rm der})\), and one claim-local actual
solution-history token \((\mathrm{SOLHIST}_{U}^{\rm phys})\) are kept
distinct.  The residual physical seven-dimensional gauge reduction
\((\mathrm{GRED}_7^{\rm phys})\), optional comparison
\((\mathrm{BULK}_8)\), coefficient routing
\((\mathrm{MATCH}_{8\to7})\), sharp frontier
\(\mathsf{SF}_{\Lambda_F}\), and \((\mathrm{SOLHIST}_{U}^{\rm phys})\)
are G0 bridge, occurrence, or readout certificates whenever they are used
physically.  The joint datum and stationary family are generated roles and
add no G0 occurrence.  Their realization must be
intrinsic to the constituted geometry and result-independent; it cannot
retroactively select or alter that geometry.  The quasi-local estimates
\((\mathrm{QLOC})\) and claim-specific continuation/reconstruction tags
remain separate analytic hypotheses.  G0 types their physical use but derives
neither their regularity nor their continuation assumptions.
The generated gravitational heat role and its legacy
\((\mathrm{GRAV}_{\circ})\), \((\mathrm{GRAV\!-\!HEAT})\), and
\((\mathrm{EH}_x)\) labels reuse the upstream root witness; they are not new
entries in this G0 ledger.
The generator/action convention \((\mathrm{CAL}_{\rm GRED})\) is a
calibration: it fixes the numerical representative of the gauge coefficient
but adds no new round-fibre geometry.  The primitive-turn duration is likewise
an explicit calibration,
\[
 \boxed{(\mathrm{CAL}_t):\qquad T_t>0},
\]
and enters only the numerical representative
\(c_{\rm op}=2\pi R_x/T_t\).  Among the macroscopic branches below, the
minimal variational domain is derived once the occurring action signature is
frozen.  The former nonzero \((\mathrm{JOINT}_x^{\rm occ})\) input is
discharged by Theorem~\ref{p3f:thm:generated-fibre-spin-joint-operator}; the
renormalized zero-derivative member, any numerical nonhorizontal
\((\mathrm{REG}_{\sigma})\) boundary, \((\mathrm{TERM}_f)\), and
\((\mathrm{GRED}_7^{\rm phys})\) contain the genuine remaining occurrence,
relation, or boundary data.  The local GRED base, principal bundle,
homogeneous Born-invariant background connection, its faithful fibre-spin
extensions, and their joint operator image are derived below and are not
members of that residual family.  Later theorem consequences do not
retroactively remove the earlier-stage status of any datum that remains.
Paper~III-E's \((\mathrm{CRP4})\) is the certified
action/contact response port of
Theorem~\ref{p3b:thm:v2211-action-contact-certificate}, not a consequence
of naming the action-density dual by a determinant symbol.  Its actual
normalized coefficient and connection defect must be evaluated before the
contact determinant-root action is transported to the physical action line.
The historical \((\mathrm{OWN}_{\rm resp/det})\) name refers to that same
certificate.  First and second variations of the occurring action are
generated in its own coefficient dual without this identification.
On the certified CRP4/PQTI branch the printed
\(\widehat\gamma_\sigma=16\sqrt{\alpha_\sigma}\) follows from those
identified owners; it is not an unconditional consequence of the
coefficient compiler or a freely chosen repair coefficient.

\begin{definition}[Generated localization certification, polarization, common-action, and reduction hypotheses]
\label{p3:def:lat-pol-csa-gred-hypotheses}
\label{p3:def:glxp-pol-csa-gred-hypotheses}
The following records are used below.  Item~(i) is a derived certification of
the current upstream prefix; the remaining listed hypotheses are additional
only when their corresponding branches are invoked.
\begin{enumerate}[label=(\roman*)]
\item \emph{Derived evaluation-preserving reciprocal localization
\((\mathrm{LOC}^{\rm ev}_{xp})\), legacy \((\mathrm{GL}_{xp})\).}
These labels no longer name additional data.  Paper~III-A,
Theorem~\ref{p3a:thm:canonical-generator-germ-localization}, combines the
Paper-II-C vector first-symbol/transvection dual with the scalar
succession/energy dual and proves that the already occurring reciprocal
generator role co-realizes, on every smooth four-dimensional
reciprocal-position chart \(U\), the scale-free cotangent generator bundle
\(T^*U\).  The generated localization is full-rank, preserves the
nondegenerate evaluation pairing and the scalar/vector provenance, is
compatible with restriction and cotangent-lift overlaps, and uses no metric
identification between tangent and cotangent representations.  In the derived
\((\mathrm{CAUS})\) domain the same local role also carries Paper~III-A's
exact positive polar.  The legacy tags certify this generated theorem and add
no G0 root.  Paper~II-C,
Corollary~\ref{p2c:cor:compact-transvection-cotangent-germ}, closes the exact
first-order compact-transvection/cotangent-germ bridge, and Paper~II-D,
Corollary~\ref{p2d:cor:finite-spatial-momentum-principal-symbol}, realizes the
same cotangent component functions as principal symbols of finite
reciprocal-position transport operators.  These results do not identify the
global compact carrier itself with the noncompact cotangent bundle.  Any
remaining \((\mathrm{PAIR}_{xp})\) label therefore denotes only a stronger
global or all-resolution finite-band operator comparison.  Numerical
representatives still use the displayed reciprocal/time calibrations when
dimensionful claims are made.

\item \emph{Derived minimal metric variational domain
\((\mathrm{VDOM}_x^{\min})\), legacy \((\mathrm{POL})\).}
For an occurring minimal local action, freeze its configuration functor before
stationarity is imposed:
\[
 \mathcal C_{\rm loc}(U)
 =\operatorname{Met}(U)\times
   \prod_{i=1}^{N}\Gamma_c(E_i),
\]
where the \(E_i\to U\) are exactly the generated retained local field bundles.
The compact reciprocal-momentum carrier, its round geometry, scale, and
connection are coefficient/representation objects and are not maps from
\(U\) into an additional moduli factor.  Hence
\[
 T\mathcal C_{\rm loc}(U)
 =\Gamma_c(\operatorname{Sym}^2T^*U)
  \oplus\bigoplus_{i=1}^{N}\Gamma_c(E_i).
\]
The varied fields are therefore \(g_x\) and all retained local section-valued
arguments.  In the Levi--Civita formulation, \(\nabla^{g_x}\) and its dual
connection co-vary functorially with \(g_x\); they are not independent fields.
Variations are compactly supported in the stated unit-localizer domain and
respect the declared boundary conditions.

The legacy label \((\mathrm{POL})\) denotes this already typed tangent domain
and adds no selector.  A Palatini, metric-affine, bimetric, or varying compact-
representation theory adds a new configuration factor and is a different
sector parent.  The present statement does not select
\(\mathcal C_{\rm loc}\) from all possible theories; it derives the variation
roster from the arguments of the occurring minimal action and forbids both an
inserted nonargument variation and an unmarked omission of a genuine field
argument.

\item \emph{Derived effective-position natural scalar-density handoff
\(\mathsf{MET}_{x}^{\rm der}\), legacy \((\mathrm{EH}_{x})\).}
Paper~III-E proves that, for every diffeomorphism-natural local scalar
\(\ell(g)\) built only from the metric and its Levi--Civita curvature data,
with \(I_{s,\ell}(g_s):=\int\ell(g_s)\,d\mu_{g_s}\), the Born-orbit average
\[
 S_{{\rm B},\ell}(g_x,g_p)
 :=\tfrac12 I_{x,\ell}(g_x)+\tfrac12 I_{p,\ell}(g_p)
\]
obeys
\[
 \boxed{S^{\rm eff}_{x,\ell}[g]
 :=j_{\rm B}^{*}S_{{\rm B},\ell}[g]
 =I_{x,\ell}[g],\qquad j_{\rm B}(g)=(g,\beta_*g).}
\]
Consequently the position variation inherits one full Born-diagonal copy for
\(\ell=1\), \(\ell=\scal\), every pure-metric heat coefficient, and any
declared local \(\ell=f(\scal)\).  There is no additional routing parameter:
the two generated half-weighted representatives pull back to
\(\tfrac12+\tfrac12=1\).  The legacy label \((\mathrm{EH}_x)\) records the
\(\ell=\scal\) specialization and is a derived certificate, not a premise.
Terms containing independent bundle curvature, matter fields, source
insertions, or fermions are not covered by this scalar-density theorem and
retain the position-role factor \(1/2\) unless their own generated
identification is proved.  Holding the compact momentum geometry fixed is a
typing consequence of \((\mathrm{VDOM}_x^{\min})\), because that
representation object is not an argument of the occurring local action.

\item \emph{Zero-derivative and two-derivative Einstein alternatives.}
The zero-derivative variation is parameterized by the generic local-action family
\[
 \boxed{\mathfrak A_{\rm REN}
 :=\{\mathrm{REN}_0(\Lambda_{\rm ren}):\Lambda_{\rm ren}\in\mathbb R\}},
 \qquad
 \mathcal E^{(0)}_{\mu\nu}=\Lambda_{\rm ren}g_{\mu\nu},
\]
with source member
\[
 \boxed{(\mathrm{REN}_0^0):\quad\Lambda_{\rm ren}=0.}
\]
The bare protected \(1/8\) constituent does not select a member of
\(\mathfrak A_{\rm REN}\).  In addition, one of the following mutually
explicit higher-order controls is declared:
\[
\begin{aligned}
(\mathcal E_{\ge4}=0)&:\quad\text{the higher-curvature metric variation vanishes}\\
&\qquad\text{exactly on the domain},\\
(\mathrm{IR}_{2\partial}(\varepsilon_*))&:\quad
\text{its declared norm is at most \(\varepsilon_*\)}\\
&\qquad\text{relative to the EH variation on a stated field/domain class.}
\end{aligned}
\]
The exact tag
\((\mathrm{EH}_{2\partial}(\Lambda_{\rm ren}))\) means
\((\mathrm{REN}_0(\Lambda_{\rm ren}))+(\mathcal E_{\ge4}=0)\); its
vacuum-term-free specialization is
\((\mathrm{EH}_{2\partial}^{0})=(\mathrm{REN}_0^0)+(\mathcal E_{\ge4}=0)\).
The generated heat expansion fixes the protected support/routing coefficient
of its bare zero-derivative scalar term to \(1/8\), exactly as for the linear
Einstein term.  On the generic local \(f(R)\) atlas this does not fix the
renormalized subtraction, sign, magnitude, or background identification
represented by \(\Lambda_{\rm ren}\).  This is a generic action interface,
not an unresolved input of every physical branch.  On the generated native
cyclic terminal branch, Corollary~\ref{p3f:cor:v2118-native-cyclic-ren0-firewall}
shows that the inverse-Euler endpoint already fixes the complete zero-order
intercept and no independent \(\mathrm{REN}_0\) member is composed on top of
it.  Numerical subtraction on one generic background does not imply
\((\mathrm{REN}_0^0)\) for neighboring variations, and a derivative expansion
does not imply either higher-order alternative.  Norm, domain, boundary
conditions, and error size must accompany the effective branch.
Only genuinely topological densities or boundary terms whose variation
vanishes on the declared compact-support/boundary domain may be quotiented
without a remainder.  Generic \(R^2\), \(R_{\mu\nu}R^{\mu\nu}\), and
equivalent non-topological curvature invariants remain in
\(\mathcal E_{\ge4}\).

\item \emph{Derived joint spectral datum
\((\mathrm{JOINT}_x^{\rm der})\),\linebreak[4]
generated stationary family
\((\mathrm{STATFAM}_x^{\rm der})\),\linebreak[4]
and actual solution-history token
\((\mathrm{SOLHIST}_{U}^{\rm phys})\); legacy \((\mathrm{STAT}_x)\) and
\((\mathrm{CSA})\).}
Theorems~\ref{p3f:thm:unique-born-invariant-gred-connection}
and~\ref{p3f:thm:generated-fibre-spin-joint-operator} construct a nonflat
principal background, extend it to both simple factors of
\(\Spin(4)_{\rm fib}\), represent it faithfully on the generated carrier,
and insert it into one local Dirac family.  The metric, the nonzero connection
anchor, and every later generated gauge or odd-field slot consequently belong
to one rooted operator family.  They use the same declared \(\Dcl\) profile
and one endomorphism-valued curvature.  This theorem output is denoted
\((\mathrm{JOINT}_x^{\rm der})\); the historical notation
\((\mathrm{JOINT}_x^{\rm occ})\) names no additional premise in the current
framework.

The conclusion is stronger than matrix-algebra availability and weaker than
a selected physical field configuration.  Nonzeroness is forced by the
explicit curvature of \(\omega_{\rm B}\), faithful infinitesimal
representation prevents its loss, and the cotangent--Clifford contraction is
the licensed operator bridge.  The resulting affine gauge/odd family is
generated as a whole; no particular fluctuation, scalar section, numerical
coefficient, action, or stationary solution is thereby selected.  The
same-connection, same-curvature, and same-profile relations follow by
associated-bundle and functional-calculus closure and are not further
choices.  Parity-even derivative-free ordinary-trace kinetic closure is
likewise derived below and is not a second ``common-normalization'' premise.

For every constituent whose Lorentzian variational role is generated, the
differential/equalizer/quotient word of
Theorem~\ref{p3f:thm:v173-stationary-solution-family} generates the complete
stationary-solution family.  In particular it applies directly to
\((\mathrm{ACT}_{{\rm J,EH+YM+\Phi},U_0}^{\rm loc,der})\); a broader claim consumes only its relevant residual \((\mathrm{ACT}_{C}^{\rm res})\) sectors defined below.  This role is denoted \((\mathrm{STATFAM}_x^{\rm der})\).  A particular physical metric/field tuple
on the spacetime domain \(U\) is separately recorded as
\[
 (\mathrm{SOLHIST}_{U}^{\rm phys}):\qquad
 [q_{\rm world}]\in\mathfrak M_{\rm stat,U}([S_L]).
\]
Because \(q_{\rm world}\) is already a section over all of \(U\), it is the
covariant on-shell history on that domain; no second history token is required.
This does not replace an initial-state, quantum-state, boundary, global
extension, or preparation premise.  The raw bosonic position trace enters with
the III-B position-role coefficient \(1/2\).  The retained fermionic quadratic
form is assigned that same position-role coefficient before a uniform canonical
field rescaling; a single nonzero scalar multiplying the entire classical
fermionic action can be removed by one uniform field rescaling, provided every
fermionic interaction is rescaled with it.  Relative species weights,
sector-dependent profiles, anomalies, measures, or quantum determinants cannot
be removed in that way and are additional insertions.  The notation
\((\mathrm{JOINT}_x)\) henceforth abbreviates the derived ancestry together
with its relation closure.  The legacy \((\mathrm{STAT}_x)\) denotes the old
compound of the now-derived family and an actual member; the legacy
\((\mathrm{CSA})\) adds no independent kinetic weight or profile.

\item \emph{Residual seven-dimensional gauge reduction
\((\mathrm{GRED}_7^{\rm phys})\) and
generator calibration \((\mathrm{CAL}_{\rm GRED})\).}
Paper~III-B now supplies the theorem-level generated compact readout role
\[
 SU(2)_{\rm diag}\longrightarrow X_7^{\rm cmp}
 \longrightarrow K_{\rm rel},
 \qquad
 X_7^{\rm cmp}=\mathcal M_8/U(1)_{\rm diag},
\]
from the already constituted reciprocal product by a conservative graph lift;
the full product and its factor tests remain co-realized.  When this theorem
is composed with Paper~III-A's generated chart
\(U\subset K_x\) and Paper~III-B's diffeomorphism
\(j_x:K_x\to K_{\rm rel}\),
Theorems~\ref{p3f:thm:canonical-local-gred-pullback} and
\ref{p3f:thm:unique-born-invariant-gred-connection} derive
\[
 X_{7,U}^{\rm can}\to U,\qquad
 \omega_{\rm B}=\tfrac12(\theta_x+\theta_p).
\]
Thus the local base, principal carrier, gluing, and Born-invariant homogeneous
principal background are not free data.  The surviving physical realization
family may be written schematically as
\[
 \boxed{\mathfrak A_{\rm GRED}^{\rm phys}
 =\{(g_F,\nabla_{\rm iso}^{\rm fluc},\mathcal R_{\rm act})
      \text{ on }X_{7,U}^{\rm can}
      \text{ and compatible with }\omega_{\rm B}\}.}
\]
Here \(F=S^3_{\rm diag}\) is already fixed and is not a free member.
The datum \(\nabla_{\rm iso}^{\rm fluc}\) denotes any physically occurring
connection on an enlarged fibre-isometry bundle, including independent
left/right fluctuations not generated by the diagonal principal background;
\(\mathcal R_{\rm act}\) denotes adoption of the corresponding reduction in
the physical action.  Once \(g_F\) is supplied, the
normalized fibre probability measure is the dependent datum
\(d\mu_F:=\operatorname{Vol}(F)^{-1}d\operatorname{vol}_F\), not a second
selector.  The residual interface \((\mathrm{GRED}_7^{\rm phys})\) declares
the current physical member and the already normalized four-dimensional
Einstein coefficient \(\mathcal C_{\rm EH}\).  The separate calibration fixes
the generator basis and metric
\(\Gamma_{AB}=\kappa r_{\rm fib}^2\delta_{AB}\), the Yang--Mills convention,
and, when used, compact-cell saturation
\(r_{\rm fib}\Lambda_{\rm cut}=1\).  No Born projector or scalar-circle volume
is part of this seven-dimensional reduction.  The equality \(g_L=g_R\) is
geometric on the symmetric round vector fibre; the absolute coefficient
requires the displayed reduction and calibration.  The tags
\((\mathrm{CORN})\) and \((\mathrm{BD}_{\rm EH})\) remain theorem-status
certifications and add no hypotheses.  For compatibility with earlier
formula labels, later occurrences of \((\mathrm{GRED}_7)\) abbreviate this
residual \((\mathrm{GRED}_7^{\rm phys})\) on the already derived carrier and
background; they do not restore the discharged data.

For the premise-minimal common-heat gauge chain used through Paper~VI-B, the
former constitutive connection occurrence is replaced by the derived
certificate \((\mathrm{JOINT}_x^{\rm der})\); its same-datum factor
connections, curvatures, profiles, and ordinary-trace kinetic relations are
derived.  The
remaining member of \(\mathfrak A_{\rm GRED}^{\rm phys}\) is retained only
for the optional diagonal-fibre left--right normalization or enlarged
isometry-connection fluctuation branch.  It is not an additional active
structural root of the premise-minimal chain.

\item \emph{Optional eight-dimensional Born comparison
\((\mathrm{BULK}_8)\).}
The compact comparison space
\[
 Y_8=(S_x^3\times S_p^3)\times(S_t^1\times S_E^1)
\]
has equal auxiliary metrics inside each reciprocal pair, normalized comparison
measures, Born involution
\(\mathcal B(x,p,t,E)=(p,x,E,t)\), and a specified nonnegative Laplace-type
operator \(L_8\) satisfying \([\mathcal B,L_8]=0\).  Put
\(P_+=(1+\mathcal B)/2\).  The scalar-circle metrics here are auxiliary
comparison data, not physical time/energy moduli.  The equivariant heat
expansion is taken on its standard smooth compact domain; fixed-set terms are
kept separate as in Paper~III-G, Appendix~\ref{p3f:app:born-even-fixed-set}.

\item \emph{Optional coefficient-routing bridge
\((\mathrm{MATCH}_{8\to7})\).}
Only when the eight-dimensional projected bulk diagnostic is used to calibrate
the seven-dimensional coefficient, a result-independent normalized-volume
map is supplied such that
\[
 \widehat{\mathcal C}_{8,{\rm bulk}}^{(+)}=\mathcal C_{\rm EH}.
\]
The Born-even factor is applied before this match and is not multiplied again
by the physical scalar rank or the corner weight.  This bridge is not a
consequence of Born exchange or of \((\mathrm{GRED}_7)\), and no generated
signature in the present source selects its normalized-volume routing map.

\item \emph{Quasi-local profile hypothesis \((\mathrm{QLOC})\).}
Let \(q\) be a supplied bounded smooth analytic coarse profile,
\(Q_\Lambda=q(D^2/\Lambda^2)\), and for
\(a\in C_c^\infty(U)\) put
\(A_{a,\Lambda}:=Q_\Lambda M_aQ_\Lambda\), with localization support
\(\operatorname{supp}a\).  There is \(\Lambda_0>0\) such that for every
\(N\in\mathbb N\) there is a constant \(C_N\), independent of
\(a,b,\Lambda\), for which all \(\Lambda\ge\Lambda_0\) and all separated
supports obey
\begin{equation}
\label{p3:eq:quantified-qloc}
 \bigl\|[A_{a,\Lambda},A_{b,\Lambda}]\bigr\|
 \le C_N\|a\|_\infty\|b\|_\infty
 \bigl(1+\Lambda d_g(\operatorname{supp}a,
                         \operatorname{supp}b)\bigr)^{-N}.
\end{equation}
Indeed, if \(E=\operatorname{supp}a\) and \(F=\operatorname{supp}b\) are
disjoint, the exact algebraic identity
\[
 [A_{a,\Lambda},A_{b,\Lambda}]
 =Q_\Lambda M_aQ_\Lambda^2M_bQ_\Lambda
  -Q_\Lambda M_bQ_\Lambda^2M_aQ_\Lambda
\]
shows that both terms factor through the separated-support operators
\[
 \mathbf1_E Q_\Lambda^2\mathbf1_F,
 \qquad
 \mathbf1_F Q_\Lambda^2\mathbf1_E.
\]
Thus \((\mathrm{QLOC})\) is reduced to an off-diagonal estimate for
\(Q_\Lambda^2\), rather than assumed directly for every commutator.
The norm is the operator norm on the declared local Hilbert realization,
\(d_g\) is the Riemannian distance of the elliptic representative used in the
spectral calculus, and the constants may depend on its bounded-geometry class
and on \(q\), but not on the two observables.  A sharp spectral projector does
not automatically have this property.
\end{enumerate}
\end{definition}

\begin{theorem}[Configuration-tangent derivation of the minimal variational domain]
\label{p3f:thm:minimal-configuration-tangent-vdom}
On the occurring minimal local-action parent of
Definition~\ref{p3:def:lat-pol-csa-gred-hypotheses}, the complete admissible
first variation is a section of \(T\mathcal C_{\rm loc}(U)\).  Therefore
\((\mathrm{VDOM}_x^{\min})\) is uniquely determined: it varies all and only
the local section-valued action arguments, co-varies dependent Levi--Civita
data, and contains no independent compact-momentum representation variation.
No \((\mathrm{POL})\) selector remains on this branch.
\end{theorem}

\begin{proof}
The tangent functor of a finite product is the product/direct sum of the
tangent functors.  The tangent of \(\operatorname{Met}(U)\) is
\(\Gamma_c(\operatorname{Sym}^2T^*U)\) on the compactly supported variation
domain, and each linear field space is its own tangent.  The compact momentum
representation is absent from \(\mathcal C_{\rm loc}(U)\), so its variation
cannot occur in the tangent without inserting a new parent factor.  Conversely
omitting one displayed field direction would require an independently
generated constraint.  Functorial dependence of the Levi--Civita connection
on \(g_x\) gives its co-variation.  This proves uniqueness inside the frozen
minimal action signature and preserves the boundary to enlarged theories.
\end{proof}

\begin{theorem}[Same-datum JOINT relation closure]
\label{p3f:thm:joint-same-datum-relation-closure}
Let \(P\to U\) be one generated principal bundle, let
\(\rho\) be its already generated representation on the retained local
carrier, and let \(A\) be a generated principal connection.  On a local spin
chart define the covariantly represented one-form module
\[
 \Omega^1_{\rm cov}(U;\rho)
 :=c_g\!\left(\Omega^1(U)\otimes\rho_*(\mathfrak g)\right).
\]
Then \(c_g(\rho_*(A))\in\Omega^1_{\rm cov}(U;\rho)\), and the twisted
operator \(D_{g,A}:=D_g+c_g(\rho_*(A))\) is one position-sector Dirac datum.
Every
associated summand receives the functorially induced connection
\(d+\rho_*(A)\), its curvature is \(\rho_*(F_A)\), and every declared profile
\(q(D_{g,A}^2/\Lambda^2)\) is one functional-calculus operation on that single
datum.  Hence, inside this same-owner branch, common connection, curvature,
and profile are relation-closure consequences; they are not three additional
selectors.

If an optional finite noncommutative spectral triple is supplied and its
Connes one-form module equals the displayed covariant module, then the same
term also lies in \(\Omega_D^1(\mathcal A)\).  That optional equality is not
needed for geometric covariant representability.  This theorem alone does not
generate \(A\), the profile function \(q\), a particular fluctuation member,
or adoption of the resulting functional as a physical action; the canonical
nonzero choice used by \((\mathrm{JOINT}_x^{\rm der})\) is generated separately
in Theorem~\ref{p3f:thm:generated-fibre-spin-joint-operator}.
\end{theorem}

\begin{proof}
Formation of an associated bundle is functorial in a principal connection:
in every local trivialization the covariant derivative is
\(d+\rho_*(A)\), and
\[
 (d+\rho_*(A))^2
 =\rho_*(dA+A\wedge A)=\rho_*(F_A).
\]
Direct sums and tensor representations use the same \(A\) through the
corresponding generated representation maps; inserting an independent
connection would change the single-datum input.  Once \(D_{g,A}\) is fixed,
the first-order commutator identity
\([D_g,f]=c_g(df)\) puts every local expression
\(A_\mu^a\rho_*(T_a)[D_g,x^\mu]\) in the displayed covariant operator module.
Thus the represented connection is obtained by the already generated
cotangent, representation, multiplication, and Clifford-contraction
operations, without assuming an optional finite spectral triple.  Once
\(D_{g,A}\) is fixed,
functional calculus gives one operator \(q(D_{g,A}^2/\Lambda^2)\), and its
restrictions to generated invariant summands are dependent restrictions of
that same operator.  These are exact relation closures of the occurring
datum.  Paper~II-A,
Theorem~\ref{p2a:thm:geometric-defect-repair-admission}, admits them by exact
transport, while its zero-member and declared-output firewalls preserve every
boundary listed in the final paragraph.
\end{proof}

\begin{remark}[No reopened momentum selector]
The former \(\mathsf P_{\rm sli}\) route is historical only.  Every use of
\(m_p=1\) or \(\mathcal H_1\) below imports the current Paper-II-B seed theorem.
Any finite position-side level remains a separate position interface and is not
set to one by this paper.
\end{remark}

\label{p3:part:dynamical-asymmetry-locality}


\paragraph{Technical proof ownership and constituted scope.}
Every extended auxiliary derivation, explicit matrix verification,
and optional comparison moved during this reorganization is
preserved in Paper~III-G, Section~\ref{p3g:sec:sources}.
The physical momentum carrier is inherited from Paper~II-B:
$\mathcal H_p\cong\mathcal H_1$ and $m_p=1$, without later
scalar, action, regulator, or source data reselecting it.

\section{Doubled local geometry, position localization, and causal polar}
\label{paper3f:sec:organized-local}
\subsection{The Dynamical Asymmetry}
\label{p3:sec:asymmetry}

Born reciprocity is the microscopic kinematic involution derived upstream by
the premise-minimal Paper-II-C/GP2 closure.  On the occurring minimal local
action, macroscopic asymmetry follows from the derived tangent domain
\((\mathrm{VDOM}_x^{\min})\): the reciprocal-position metric enters the
fluctuated operator and is varied, whereas the compact momentum radius remains
representation data outside the configuration functor.  The local cotangent target and its canonical gluing are derived below from
scale-free evaluation duality, while Paper~III-A's generator-germ theorem
already supplies the physical co-realization of the pre-existing reciprocal
generator role on that target.  No separate \((\mathrm{GL}_{xp})\)
occurrence remains; the later variational statements are conditional only on
the occurrence of the displayed minimal action parent and its derived
\((\mathrm{VDOM}_x^{\min})\).

\subsubsection{Local doubled geometry and infrared polarizations}
\label{p3:sec:local-born}

The theorem-level statement used here is local rather than global.  On a smooth
reciprocal-position chart, nondegenerate evaluation between operation vectors
and generator readouts fixes the algebraic dual bundle type and its coordinate
covariance without a cone or scale calibration.  Paper~III-A's generated-role theorem co-realizes the pre-existing physical
reciprocal generator role on that derived cotangent target.  Only in the
derived \((\mathrm{CAUS})\) domain is Paper~III-A's exact crossed positive
polar of the generated clocked-operation cone imported.  Two geometrically
distinct objects must be kept separate.  The physical local phase space is the
cotangent bundle \(M_U=T^*U\), with its canonical tautological one-form,
symplectic form, and vertical polarization.  The doubled Born algebra lives
canonically on the exact generalized bundle
\[
        E_U:=TU\oplus T^*U
\]
over the base chart.  A connection may split the tangent exact sequence of
\(T^*U\) and thereby identify its tangent bundle with a pullback of \(E_U\), but
that identification is auxiliary and is not needed to define the Born data.
This distinction removes any hidden assumption of a canonical or integrable
horizontal distribution on \(T(T^*U)\).

\begin{definition}[Local phase space and doubled bundle]
\label{p3:def:local_doubled_patch}
Let \(U\) be a local reciprocal-position chart.  Theorem~\ref{p3:thm:local_reciprocal_atlas} derives \(T^*U\) as the
coordinate-natural local generator target and its cotangent-lift gluing on
overlaps; Paper~III-A,
Theorem~\ref{p3a:thm:canonical-generator-germ-localization}, supplies its
physical co-realization.  Define unconditionally on this generated local role
\[
        M_U:=T^*U,
        \qquad
        \pi:M_U\to U.
\]
The local phase space carries the canonical exact sequence
\[
0\longrightarrow \pi^*T^*U
 \xrightarrow{\ \iota\ }
TM_U
 \xrightarrow{\ d\pi\ }
\pi^*TU
\longrightarrow0.
\]
Its vertical subbundle \(\ker d\pi\cong\pi^*T^*U\), tautological one-form
\(\vartheta\), and symplectic form \(\omega_{\rm can}=d\vartheta\) are canonical.
A linear connection on \(TU\) gives a splitting of this exact sequence, but no
such horizontal splitting is canonical and its horizontal distribution need
not be integrable.

Independently, define the exact doubled bundle over the base
\[
        E_U:=TU\oplus T^*U,
        \qquad
        F:=TU\oplus0,
        \qquad
        G:=0\oplus T^*U.
\]
The summands \(F\) and \(G\) are canonical subbundles of \(E_U\).  On overlaps
they glue by
\[
        X+\xi
        \longmapsto
        \phi_*X+(\phi^{-1})^*\xi,
\]
the direct sum of the tangent and cotangent representations.  A chosen
connection identifies \(TM_U\) with \(\pi^*E_U\), but all Born-algebraic
statements below are made intrinsically on \(E_U\), not on a presumed canonical
horizontal split of \(TM_U\).
\end{definition}

\begin{lemma}[Scale-free coordinate-natural evaluation dual]
\label{p3:lem:no-metric-reciprocal-duality}
Let \(U\) be a smooth reciprocal-position chart and let
\(V=T_xU\) at a point \(x\in U\).  Suppose local operation increments transform
as vectors in \(V\) and local generator readouts pair linearly and
nondegenerately with them through evaluation
\[
        \langle p,v\rangle=p(v)
\]
on the vector part.  If an independent scalar operation line is included, use
the crossed scalar/vector evaluation
\[
        \langle(E,p),(\tau,\xi)\rangle=E\tau-p(\xi)
\]
on the full local operation space.  Then the vector generator representation is
canonically \(V^*\).  Under a coordinate change \(A:V\to V'\), the generator
covector transforms by
\[
        p\longmapsto(A^{-1})^*p.
\]
No identification of this dual representation with \(V\) is canonical before
a metric or other nondegenerate bilinear form is supplied.  This conclusion
fixes neither a positive cone, its aperture, an energy normalization,
\(c_{\rm op}\), nor \((\mathrm{CAL}_{xp})\).
\end{lemma}

\begin{proof}
Nondegenerate linear evaluation identifies the generator readout space with
the algebraic dual of the local operation space.  If \(v'=Av\), invariance of evaluation
requires
\[
        p'(v')=p(v)
        \qquad\text{for all }v,
\]
which uniquely gives \(p'=(A^{-1})^*p\).  Thus the coordinate-natural bundle
assembled from the generator fibres is the cotangent bundle.  A map
\(V^*\to V\) would instead require a nondegenerate bilinear form, a frame, or
another intertwiner not contained in the bare evaluation-dual construction.
Hence the
cotangent transformation law is derived without inserted metric data.  No
order or positivity structure was used, so no physical cone statement follows
from this lemma alone.
\end{proof}

\begin{lemma}[Uniqueness of cotangent-lift transition maps]
\label{p3:lem:cotangent-lift-uniqueness}
Let \(\phi:U\to U'\) be a reciprocal-position coordinate transition.  Let
\(\Phi\) be any \(C^1\) map
\[
        \Phi:T^*U\longrightarrow T^*U'
\]
covering \(\phi\).  If \(\Phi\) preserves the tautological one-form, then it
is necessarily the
cotangent lift
\[
        \Phi=T^*(\phi^{-1}),
        \qquad
        (x,p)\longmapsto
        \bigl(\phi(x),(D\phi_x^{-1})^*p\bigr).
\]
These maps satisfy the atlas cocycle and preserve the canonical symplectic
form.  Fibrewise linearity is therefore a conclusion, not a premise.  An
affine momentum shift or an independent \(O(d,d)\), \(B\)-field, or
\(\beta\)-field twist requires extra transition data.
\end{lemma}

\begin{proof}
Let \(\vartheta_U\) and \(\vartheta_{U'}\) be the tautological one-forms.  For
\(\Phi(x,p)=(\phi(x),p')\), the condition
\(\Phi^*\vartheta_{U'}=\vartheta_U\) gives, for every tangent vector \(v\),
\[
        p'(D\phi_x\,d\pi(v))=p(d\pi(v)).
\]
Therefore \(p'=(D\phi_x^{-1})^*p\), proving uniqueness.  Functoriality of the
cotangent lift gives the cocycle law, and
\[
        \Phi^*d\vartheta_{U'}=d\vartheta_U
\]
proves symplectic invariance.  A shift \(p'\mapsto p'+\alpha_{\phi(x)}\) changes
the tautological form by the pullback of a one-form \(\alpha\); more general
twists likewise require an independently supplied tensor.  They are therefore
not transitions of the uninserted cotangent atlas.
\end{proof}

\begin{theorem}[Scale-free cotangent target and calibrated physical reciprocal atlas]
\label{p3:thm:local_reciprocal_atlas}
Let \(x\) lie in a smooth four-dimensional reciprocal-position manifold.
After shrinking an arbitrary chart about \(x\), choose a coordinate-ball
neighborhood \(U\ni x\).  If local operation increments transform in
\(TU\), then \(U\) has a canonical \emph{local generator target} with the
following properties:
\begin{enumerate}[label=(\roman*),leftmargin=2.2em]
\item \(U\) is a smooth four-dimensional coordinate-ball patch,
\[
        U\cong B^4\subset\RR^4;
\]
\item scale-free evaluation duality in
      Lemma~\ref{p3:lem:no-metric-reciprocal-duality} identifies the local
      generator representation with \(T^*U\), not with a second copy of
      \(TU\);
\item reciprocal-position coordinate changes act on the local generator target
      by the cotangent representation, and any \(C^1\) lift covering the base
      transition and preserving the tautological form is the ordinary
      cotangent lift (hence is fibrewise linear);
\item consequently the derived local target carries the tautological one-form
\[
        \vartheta=p_\mu dx^\mu
\]
      and canonical symplectic form
\[
        \omega_{\rm can}=d\vartheta=dp_\mu\wedge dx^\mu,
\]
      and these structures obey the atlas cocycle on overlaps.
\end{enumerate}
Paper~III-A,
Theorem~\ref{p3a:thm:canonical-generator-germ-localization}, proves that the
already constituted physical reciprocal generator role is co-realized on this
derived target without a second G0 occurrence.  Hence the physical local phase
space is
\[
        M_U=T^*U,
\]
with canonical vertical polarization and exact tangent sequence
\[
0\longrightarrow \pi^*T^*U
 \longrightarrow TM_U
 \xrightarrow{\ d\pi\ }
\pi^*TU
\longrightarrow0.
\]
After a local affine connection is chosen one may split this sequence as
\[
        TM_U\cong\pi^*(TU\oplus T^*U),
\]
but neither this horizontal splitting nor its integrability is canonical.
The exact doubled base bundle is intrinsically
\[
        E_U=TU\oplus T^*U.
\]

If the derived \((\mathrm{CAUS})\) domain of
Definition~\ref{p3:def:causal-wick-response-hypotheses} is also imported, then
Paper~III-A's calibrated physical positive polar is realized on the same
cotangent target:
\[
 (\mathcal C_x^+)^\vee
 =\{(E,p):E\ge c_{\rm op}\|p\|_{h^{-1}}\}.
\]
This last statement, unlike the cotangent type and canonical forms, consumes
the derived clocked-Hopf doctrine and its characteristic/succession
provenance through \((\mathrm{CAUS})\).  A numerical dimensionful
representative additionally consumes \((\mathrm{CAL}_{xp})\) and
\((\mathrm{CAL}_t)\); no propagation-family occurrence is
added.

The theorem determines the effective local phase-space target, while the
Paper-III-A generator-germ theorem supplies its generated physical
co-realization.  It does not identify the
global compact momentum carrier with a noncompact cotangent bundle, does not
choose a pointwise intertwiner between the independently covariant compact
vector factors before localization, and the scale-free part does not fix the
reciprocal calibration \((\mathrm{CAL}_{xp})\).  The exact first-order
compact-transvection/cotangent-germ bridge is already supplied upstream by
Paper~II-C; only a stronger global/all-resolution finite-band operator
intertwiner remains outside the present theorem.  The fixed Paper-II-B momentum seed and
its derived label \(m_p=1\) are unchanged.
In particular, the local evaluation dual and a local principal symbol
\(\sigma_1(D)(x,\xi)\) use the cotangent variable \(\xi\) directly and require
no finite-band pairing premise.  Paper~II-C,
Theorem~\ref{p2c:thm:first-symbol-pairing} and
Corollary~\ref{p2c:cor:compact-transvection-cotangent-germ}, supply the exact
first-symbol and compact-transvection/cotangent-germ pairings.  Paper~II-D,
Corollary~\ref{p2d:cor:finite-spatial-momentum-principal-symbol}, supplies the
finite reciprocal-position principal-symbol realization.  A universal
all-resolution map into the frozen compact carrier remains a distinct and, on
the faithful-injective reading, excluded stronger object.
\end{theorem}

\begin{proof}
Choose a chart \(\psi:W\to O\subset\RR^4\) about \(x\).  Since \(O\) is
open, it contains a Euclidean ball centered at \(\psi(x)\); its inverse image
is the required \(U\).  Thus the ball statement is obtained by shrinking and
is not asserted for an arbitrary preassigned chart domain.
Lemma~\ref{p3:lem:no-metric-reciprocal-duality} gives the
coordinate-natural cotangent transformation law directly from evaluation
duality.  Hence the
local generator target is \(T^*U\) without assuming a metric identification of
tangent and cotangent representations.

For a position-coordinate transition \(\phi:U\to U'\),
Lemma~\ref{p3:lem:cotangent-lift-uniqueness} shows that any \(C^1\) lift
covering the base transition and preserving the tautological one-form is
\[
        T^*(\phi^{-1}).
\]
Functoriality supplies the cocycle, and preservation of \(\vartheta\) implies
preservation of \(d\vartheta\).  Thus the cotangent atlas, tautological form,
symplectic form, and vertical polarization are theorem consequences of the
local dual target, not additional localization hypotheses.

Paper~III-A,
Theorem~\ref{p3a:thm:canonical-generator-germ-localization}, now proves the
previously external identification: the already constituted reciprocal
generator role has these cotangent generator germs as its chartwise local
co-realization.  Therefore \(M_U=T^*U\) is the physical local phase space
without another occurrence premise.  Its canonical exact tangent sequence
gives the vertical subbundle; a chosen connection may split it, while the
exact doubled bundle \(TU\oplus T^*U\) exists independently of that choice.
No step equates the global compact carrier with the local noncompact cotangent
bundle.

Finally, within the derived \((\mathrm{CAUS})\) domain, Paper~III-A,
Theorem~\ref{p3:thm:gp1-operation-generator-polar-dual}, computes the positive
polar of the generated clocked-operation cone.  The generated cotangent co-realization therefore carries that physical
polar on the same target.  This implication does not retroactively make the
scale-free bundle construction depend on the cone calibration.
\end{proof}

\begin{remark}[Momentum realization chain and operator-equality firewall]
\label{p3f:rem:momentum-realization-chain}
On the spatial part of a reciprocal-position chart, the current framework now
has three separately typed realizations of one local generator role:
\[
 \widehat P_i(a)\in\mathscr K_a
 \xleftrightarrow[\text{first-order germ}]{\ \mathfrak B_a\ }
 e_b^i\in T_b^*S_B^3,
 \qquad
 \sigma_{\rm phys}(\Pi_i^{(a)})(x,p)=p(e_i(x)).
\]
The first relation is the exact Paper-II-C bundle isomorphism; the second is
the Paper-II-D principal-symbol statement on the reciprocal-position scalar
band.  They do not assert an operator equality
$\widehat P_i=\Pi_i^{(a)}$, a common Hilbert carrier, or a flow
$m_p=1\to\infty$.  The large-radius commuting momentum variables are obtained
only through the later Paper-VII-C trace-moment limit of
$\Pi_i^{(a)}$.
\end{remark}

\begin{corollary}[Homogeneous covector scaling and conditional physical-polar boundary]
\label{p3:cor:homogeneous-cotangent-inverse-scaling}
Let the reciprocal-position spatial metric on a homogeneous round branch be
\[
        h(\widetilde t)=a(\widetilde t)^2\gamma,
\]
where \(\gamma\) is a fixed unit-round reference metric.  Then for every
fixed comoving covector \(p\),
\[
        \|p\|_{h^{-1}}
        =\frac{1}{a(\widetilde t)}\,
          \|p\|_{\gamma^{-1}}.
\]
If the complete \((\mathrm{CAUS})\) domain holds, the
boundary of the physical local GP1 polar cone consequently obeys
\[
        E_{\partial}(\widetilde t)
        =c_{\rm op}\|p\|_{h^{-1}}
        =\frac{c_{\rm op}}{a(\widetilde t)}
          \|p\|_{\gamma^{-1}}.
\]
Thus expansion of the position metric produces the inverse scaling of the
corresponding local cotangent generator norm.  Only under \((\mathrm{CAUS})\) is the boundary equality a physical
local reciprocal-generator statement.  The statement is kinematic: it neither derives an equation for
\(a(\widetilde t)\), proves
\(\partial_{\widetilde t}a>0\), identifies the compact scalar-circle
circumference with \(a\), nor asserts that every mode lies on the polar
boundary.
\end{corollary}

\begin{proof}
The inverse metric is
\[
        h^{-1}(\widetilde t)=a(\widetilde t)^{-2}\gamma^{-1},
\]
so the norm of a fixed covector scales by \(a^{-1}\).  Substitution into the
boundary equality of Paper~III-A within the complete \((\mathrm{CAUS})\) domain,
Theorem~\ref{p3:thm:gp1-operation-generator-polar-dual}, gives the second
display.  The final qualifications follow because the polar theorem supplies no
equation for \(a\), the sign of its
derivative, a scalar-circle/metric identification, or a mode-on-shell
condition.
\end{proof}

\begin{definition}[Sharp scalar spectral-frontier occurrence]
\label{p3:def:sharp-scalar-resolution-frontier}
On the homogeneous reciprocal-position branch
\[
 h(\widetilde t)=a(\widetilde t)^2\gamma,
\]
the interface
\[
 (\mathrm{OCC}_{x,\Lambda_F}^{\rm sharp,full})
\]
asserts that, for an independently declared physical frontier scale
\(\Lambda_F>0\), with resolution length
\[
 \boxed{\ell_F:=\Lambda_F^{-1}},
\]
the intrinsic sharp scalar
projector
\[
 \Pi_{\Lambda_F}(\widetilde t)
 :=\mathbf1_{[0,\Lambda_F^2]}
   \!\left(\Delta_{h(\widetilde t)}\right)
\]
is generated as the complete instantaneous scalar test frontier, together
with the primitive frontier \(F_{x,0}=E_0^x\oplus E_1^x\), and that GP2
cumulatively retains every such generated test.  This is an operational
frontier occurrence; it is not implied merely by placing a common profile in
the spectral action.  This paper denotes the occurrence
\(\mathsf{SF}_{\Lambda_F}\) so that the independently frozen scale remains
visible at every use.  The identification
\[
 (\mathrm{ID}_{F/A}):\qquad \Lambda_F=\Lambda_{\rm cut}
\]
is a separate matching condition and is not used in the frontier theorem.
\end{definition}

\begin{theorem}[Exact homogeneous scalar spectral-resolution staircase]
\label{p3:thm:exact-spectral-resolution-staircase}
Assume Definition~\ref{p3:def:sharp-scalar-resolution-frontier}, let
\(a(\widetilde s)>0\) be an arbitrary scale history, and fix
\(\widetilde t_0\).  No continuity hypothesis is needed.  The instantaneous
resolution is
\[
 \boxed{
 m_x^{\rm inst}(\widetilde t)
 =\max\!\left\{1,
 \left\lfloor
 \sqrt{1+\bigl(\Lambda_Fa(\widetilde t)\bigr)^2}-1
 \right\rfloor\right\},}
\]
and the primitive history gives the integer record first:
\[
 \boxed{
 m_x(\widetilde t)
 =\sup_{\widetilde s\in[\widetilde t_0,\widetilde t]}
 m_x^{\rm inst}(\widetilde s)\in\mathbb N\cup\{+\infty\}.}
\]
Equivalently,
\[
 m_x(\widetilde t)
 =\max\!\left\{1,\sup\!\left\{
 k\in\mathbb N:\ \exists\widetilde s\le\widetilde t,
 \ a(\widetilde s)\ge
 \frac{\sqrt{k(k+2)}}{\Lambda_F}
 \right\}\right\}.
\]
Put
\[
 A_*(\widetilde t):=
 \sup_{\widetilde s\in[\widetilde t_0,\widetilde t]}
 a(\widetilde s),
 \qquad
 \Phi_F(A):=\sqrt{1+(\Lambda_FA)^2}-1.
\]
If \(A_*=+\infty\), then \(m_x=+\infty\).  If \(A_*<\infty\) and either
\(A_*\) is attained or \(\Phi_F(A_*)\notin\mathbb N\), then
\[
 m_x(\widetilde t)
 =\max\!\left\{1,\left\lfloor
 \Phi_F\bigl(A_*(\widetilde t)\bigr)
 \right\rfloor\right\}.
\]
If \(A_*<\infty\) is not attained and \(\Phi_F(A_*)=k\in\mathbb N\), the
exact threshold correction is
\[
 \boxed{m_x(\widetilde t)=\max\{1,k-1\}.}
\]
Continuity on the compact past interval is merely one sufficient condition
for attainment; it is not needed for the integer-record theorem.
If \(a\) is nondecreasing in the inherited scalar-succession orientation,
then \(A_*(\widetilde t)=a(\widetilde t)\) because the endpoint belongs to the
interval, and the instantaneous and retained formulas coincide without any
continuity assumption.
\end{theorem}

\begin{proof}
Constant homothety gives
\[
 \Delta_{a^2\gamma}=a^{-2}\Delta_\gamma,
 \qquad
 \lambda_k(\widetilde t)=\frac{k(k+2)}{a(\widetilde t)^2}.
\]
Hence degree \(k\) lies in the sharp range exactly when
\[
 k(k+2)\le \bigl(\Lambda_Fa\bigr)^2
 \iff
 k\le\sqrt{1+\bigl(\Lambda_Fa\bigr)^2}-1.
\]
Paper~II-C, Theorem~\ref{p2c:thm:exact-reciprocal-scalar-bands} and
Proposition~\ref{p2c:prop:least-generated-position-resolution}, identifies
the complete band through degree \(m\) with \(\mathcal P_x(m)\) and reads its
resolution from the highest present harmonic.  The primitive \(E_1^x\) seed
supplies the outer maximum with one.
All homothetic projectors use the same eigenspaces, so cumulative retention is
their union and its largest degree is exactly the displayed integer supremum.
If the scale supremum is infinite, every finite threshold is crossed.  If it
is finite and attained, monotonicity of \(\Phi_F\) gives the floor at that
maximum.  If it is finite and not attained, the values of \(\Phi_F(a)\)
approach \(\Phi_F(A_*)\) strictly from below.  Their floors still reach
\(\lfloor\Phi_F(A_*)\rfloor\) unless that limit is an integer \(k\), in which
case every floor is at most \(k-1\) and values arbitrarily close to the
supremum realize \(k-1\).  This proves the threshold correction.  The
frontier occurrence is the datum and the band/resolution is derived from it;
no level selector is introduced.
\end{proof}

\begin{corollary}[Sharp-frontier discharge of the scalar generated boundary]
\label{p3f:cor:sharp-frontier-discharges-generated-boundary}
On the occurrence \(\mathsf{SF}_{\Lambda_F}\), define
\[
 m_F(\widetilde t):=m_x^{\rm inst}(\widetilde t),
 \qquad
 F_{\Lambda_F}^{\rm inst}(\widetilde t)
 :=F_{x,0}+\operatorname{Ran}\Pi_{\Lambda_F}(\widetilde t).
\]
Then the complete instantaneous retained frontier has exactly the Paper~II-C
scalar band
\[
 \boxed{
 F_{\Lambda_F}^{\rm inst}(\widetilde t)
 =\mathcal P_x\!\bigl(m_F(\widetilde t)\bigr),}
\]
so the registered scalar generated-boundary condition is discharged at that
stage:
\[
 \boxed{
 \mathsf{SF}_{\Lambda_F}
 \Longrightarrow
 \mathsf P_{\rm gbd}^{\rm scalar}
 \!\bigl(x,m_F(\widetilde t)\bigr).}
\]
With cumulative GP2 retention, the same statement holds for the retained
record level \(m_x(\widetilde t)\).  Chronologically this is a
\textsc{derived-later} discharge: Paper~II-C leaves the scalar generated
boundary unresolved for a generic occurring frontier, whereas the independently
occurring sharp spectral frontier here generates the complete harmonic band.
This is branch-specific; it does not turn an arbitrary intrinsic grammar into
a generated boundary, and it does not identify \(\Lambda_F\) with the smooth
action scale.
\end{corollary}

\begin{proof}
The spectral part satisfies
\(\lambda_k(\widetilde t)\le\Lambda_F^2\) exactly through the largest
spectrally admitted degree.  Definition~\ref{p3:def:sharp-scalar-resolution-frontier}
also retains the primitive frontier \(F_{x,0}=E_0^x\oplus E_1^x\).  Therefore
the combined instantaneous frontier contains exactly the complete degrees
\(0\le k\le m_F(\widetilde t)\), where the outer maximum with one in
Theorem~\ref{p3:thm:exact-spectral-resolution-staircase} is precisely the
primitive-frontier contribution.  Paper~II-C identifies that direct sum with
\(\mathcal P_x(m_F(\widetilde t))\).  Hence the retained boundary is
generated rather than separately inserted.  Cumulative retention replaces
the instantaneous level by the record level.
The generic Paper~II-A warning remains unchanged because no analogous
frontier occurrence has been supplied for an arbitrary constructor grammar.
\end{proof}

\begin{corollary}[Sharp-frontier generation of the finite position history]
\label{p3f:cor:sharp-frontier-generates-position-history}
On the occurrence \(\mathsf{SF}_{\Lambda_F}\), define the cumulatively retained
finite frontier
\[
 F_x(\widetilde t)
 :=\operatorname{Ret}_{\widetilde s\in[\widetilde t_0,\widetilde t]}
 \left[
   F_{x,0}+\operatorname{Ran}\Pi_{\Lambda_F}(\widetilde s)
 \right],
\]
where \(\operatorname{Ret}\) denotes the GP2 cumulative-retention closure
already included in Definition~\ref{p3:def:sharp-scalar-resolution-frontier}.
Then
\[
 \boxed{
 \mathsf{SF}_{\Lambda_F}
 \Longrightarrow
 (\mathrm{HIST}_x^{\rm fin})}
\]
on this branch.  More explicitly,
\[
 F_x(\widetilde s)\subseteq F_x(\widetilde t)
 \qquad
 (\widetilde s\le \widetilde t),
\]
and Paper~II-C's least-resolution functor gives
\[
 \boxed{
 m_x(\widetilde t)=m_x^{\rm res}\!\bigl(F_x(\widetilde t)\bigr).}
\]
Hence no second G0 history root is introduced on the sharp-frontier branch:
the independently occurring frontier already generates the retained history.
Chronologically this is a \textsc{derived-later} discharge of the generic
Paper~II-C/Paper~III-A history interface.  It does not assert nonconstant
growth: if the sharp spectral ranges are constant, the derived history is
constant; if they change, a permanently constant history would omit generated
retained tests.
\end{corollary}

\begin{proof}
Definition~\ref{p3:def:sharp-scalar-resolution-frontier} constitutes, at every
\(\widetilde s\), the complete instantaneous scalar frontier together with
cumulative GP2 retention.  Corollary~\ref{p3f:cor:sharp-frontier-discharges-generated-boundary}
identifies every instantaneous frontier with the complete Paper~II-C band
\(\mathcal P_x(m_F(\widetilde s))\).  Taking the already licensed retention
closure therefore produces a finite frontier-valued history in
\(\mathcal A_x^{\rm alg}\), and the retained-inclusion relation is immediate
from the definition of \(\operatorname{Ret}\).  Paper~II-C then assigns its
least resolution by \(m_x^{\rm res}\).  Omitting this retained history while
keeping \(\mathsf{SF}_{\Lambda_F}\) would omit generated retained tests rather
than define a different admissible frontier realization.
\end{proof}

\begin{corollary}[Exact scalar entry thresholds and test count]
\label{p3:cor:spectral-resolution-thresholds}
For \(k\ge2\), degree \(k\) first enters when
\[
 \boxed{a\ge a_k:=\frac{\sqrt{k(k+2)}}{\Lambda_F}.}
\]
At resolution \(m\), the number of resolved scalar tests is
\[
 \boxed{
 \dim\mathcal P_x(m)
 =\frac{(m+1)(m+2)(2m+3)}6.}
\]
Thus nondecreasing \(a\) makes the instantaneous resolution nondecreasing and
equal to the cumulatively retained one; unbounded \(a\) makes
\(m_x\to\infty\); and bounded \(a\) crosses only finitely many thresholds.
\end{corollary}

\begin{proof}
The threshold is the equality case
\(k(k+2)=\Lambda_F^2a^2\).
Theorem~\ref{p2c:thm:exact-reciprocal-scalar-bands} gives
\(\dim E_k^x=(k+1)^2\), so
\[
 \dim\mathcal P_x(m)
 =\sum_{k=0}^{m}(k+1)^2
 =\frac{(m+1)(m+2)(2m+3)}6.
\]
The remaining statements follow directly from the threshold sequence and the
staircase theorem.
\end{proof}

\begin{corollary}[Derived primitive plateau and exact finite-band activation]
\label{p3f:cor:v259-derived-plateau-finite-band-activation}
Let
\[
 a_L:=\frac{\sqrt{L(L+2)}}{\Lambda_F},\qquad L\ge2.
\]
For the instantaneous sharp frontier,
\[
 a(\widetilde t)<a_2=\frac{\sqrt8}{\Lambda_F}
 \quad\Longrightarrow\quad
 \boxed{m_x^{\rm inst}(\widetilde t)=1}.
\]
Hence, on a nondecreasing sharp-frontier history that begins below the first
nonprimitive threshold, the retained level is also
\[
 \boxed{m_x(\widetilde t)=1\qquad(\widetilde t<\tau_2)}.
\]
Thus the early level-one position-resolution plateau is a consequence of the
occurring sharp frontier and the subthreshold scale domain; it is not an
additional x-side level selector or initial-condition premise.

More generally, whenever the instantaneous scale satisfies
\[
 a(\widetilde t)\ge a_L,
\]
one has
\[
 \boxed{\mathcal P_x(L)\subseteq
 F_{x,0}+\operatorname{Ran}\Pi_{\Lambda_F}(\widetilde t)}.
\]
With cumulative GP2 retention, if the threshold has been crossed at any
\(\widetilde s\le\widetilde t\), then
\[
 \boxed{\mathcal P_x(L)\subseteq F_x(\widetilde t)}.
\]
In particular, on a nondecreasing history the single condition
\(a(\widetilde t)\ge a_L\) permanently activates the complete finite band
\(\mathcal P_x(L)\).  A downstream statement whose only x-side domain
requirement is the presence of that complete band therefore needs no second
band-occurrence premise after this threshold.

This result does not select the scale history, a state, an action member, a
source, a coupling, a calibration, or an analytic continuum limit.  It also
does not identify the equal labels \(m_x=m_p=1\) with equal carriers or equal
physical radii.
\end{corollary}

\begin{proof}
For \(L=2\), Theorem~\ref{p3:thm:exact-spectral-resolution-staircase} gives
\(\sqrt{1+(\Lambda_Fa)^2}-1<2\) exactly when \(a<a_2\); the outer maximum
with one therefore fixes the instantaneous level at one.  On a nondecreasing
history no earlier higher threshold can have been crossed before \(\tau_2\),
so the retained record agrees with the instantaneous level.  For general
\(L\), the defining spectral inequality
\(L(L+2)/a^2\le\Lambda_F^2\) is equivalent to \(a\ge a_L\), and the sharp
projector therefore contains every degree through \(L\).  Corollary~\ref{p3f:cor:sharp-frontier-discharges-generated-boundary}
identifies that range with the complete band, while
Corollary~\ref{p3f:cor:sharp-frontier-generates-position-history} supplies its
persistent retained inclusion.
\end{proof}

\begin{theorem}[Exact spectral-frontier event process and physical-time obstruction]
\label{p3f:thm:spectral-frontier-event-process}
Assume \(\mathsf{SF}_{\Lambda_F}\) and let the parent scale history
\(a:[\widetilde t_0,\widetilde t_1]\to(0,\infty)\) be continuous and
nondecreasing in the inherited scalar-succession orientation.  Put
\(m_0=m_x(\widetilde t_0)\).  For every later reached level \(k>m_0\), let
\[
 \tau_k:=\inf\left\{\widetilde t\ge\widetilde t_0:
 a(\widetilde t)\ge
 a_k\right\},
 \qquad
 a_k=\frac{\sqrt{k(k+2)}}{\Lambda_F}.
\]
Then the first-entry events are generated by the scale/frontier parent and
satisfy
\[
 a(\tau_k)=a_k,
 \qquad
 \tau_k<\tau_{k+1},
 \qquad
 m_x(\tau_k^-)=k-1,
 \qquad
 m_x(\tau_k)=k.
\]
At each event the retained frontier advances by the unique complete-band
inclusion
\[
 \boxed{
 \mathcal P_x(k)
 =\mathcal P_x(k-1)\oplus E_k^x,
 \qquad
 \iota_{k\leftarrow k-1}:
 \mathcal P_x(k-1)\hookrightarrow\mathcal P_x(k).}
\]
Consequently the event-indexed descendant is the exact successor law
\[
 \boxed{
 S_n(k)=k+n,
 \qquad
 \iota_{k+n\leftarrow k}
 =\iota_{k+n\leftarrow k+n-1}\cdots
  \iota_{k+1\leftarrow k}}
\]
wherever the indicated events occur.  This closes the local
position-frontier successor obligation in event order and introduces neither
a level selector nor a second history occurrence.

If \(a\) is strictly increasing through any one of these entry times, there
is no translation-homogeneous deterministic physical-time semigroup
\(T_h\) on \(m_x\) alone with
\[
 m_x(\widetilde t+h)=T_h(m_x(\widetilde t))
\]
for all admissible \(\widetilde t,h\ge0\).  On every finite nonterminal
level, however, the generated threshold residual
\[
 \boxed{
 r_F(\widetilde t)
 :=a_{m_x(\widetilde t)+1}-a(\widetilde t)>0}
\]
makes \((m_x,r_F)\) fibre-separating because
\[
 a(\widetilde t)
 =a_{m_x(\widetilde t)+1}-r_F(\widetilde t).
\]
Thus any already generated exact scale process transports exactly to this
completed state.  The sharp frontier does not itself select the scale-factor
equation, its initial condition, a physical-time event rate, or a particular
cosmological solution.
\end{theorem}

\begin{proof}
Apply Paper~III-A,
Theorem~\ref{p3a:thm:threshold-event-descent}, to the parent coordinate
\(x=a\), threshold family \(b_k=a_k\), and threshold readout supplied by
Theorem~\ref{p3:thm:exact-spectral-resolution-staircase}.  The thresholds are
strictly increasing and locally finite.  Continuity and monotonicity give the
first-entry identities and prevent skipped levels.  Corollary~\ref{p3f:cor:sharp-frontier-discharges-generated-boundary}
identifies the frontier at level \(k\) with the full band
\(\mathcal P_x(k)=\bigoplus_{j=0}^{k}E_j^x\), so the unique retained step adds
exactly \(E_k^x\), and transitivity of inclusion gives the composition law.
The same Paper~III-A theorem gives the physical-time nonprojectability at a
strict crossing and the fibre-separating residual coordinate.  Since the
argument consumes the occurring scale history but supplies no equation or
initial datum for it, all final scope exclusions follow.
\end{proof}

\begin{corollary}[Exact transverse symbol at each spectral-frontier event]
\label{p3f:cor:frontier-transverse-shell-symbol}
Under the hypotheses of
Theorem~\ref{p3f:thm:spectral-frontier-event-process}, every reached first-entry
event \(\tau_k\) opens the quotient shell
\[
 \mathcal P_x(k)/\mathcal P_x(k-1)\simeq E_k^x.
\]
Paper~II-C,
Theorem~\ref{p2c:thm:born-transported-transverse-shell-reconstruction}, then
gives the exact Hopf-marked transverse symbol reconstruction
\[
 \boxed{
 \mathcal P_x(k)/\mathcal P_x(k-1)
 \xrightarrow{\sim} E_k^x
 \xrightarrow{\mathfrak H_k^x}
 \mathcal T_{x,k+1}^{\perp},
 \qquad
 N_{\partial x}(k)=(k+1)^2.}
\]
Thus the physical transverse-window label attached to the degree-\(k\)
position frontier is \(\mu_x=k+1\), not \(k\).  This is a shellwise spectral
reconstruction theorem; it is not a microscopic position-cell count, an
occupation number, an entropy, or a reconstruction of the entire accumulated
band from the single current shell.
\end{corollary}

\begin{corollary}[Exact Hopf-base spectral area density at first entry]
\label{p3f:cor:frontier-transverse-area-density}
At the first-entry event \(\tau_k\), let \(S^2_{H,x}(\tau_k)\) be the metric
Hopf base of the reciprocal round three-sphere of radius
\(a(\tau_k)\).  Its radius is \(a(\tau_k)/2\), hence
\[
 A_H(\tau_k)=\pi a(\tau_k)^2.
\]
Using
\[
 a(\tau_k)=\frac{\sqrt{k(k+2)}}{\Lambda_F},
 \qquad
 N_{\partial x}(k)=(k+1)^2,
\]
one obtains the exact density
\[
 \boxed{
 \frac{N_{\partial x}(k)}{A_H(\tau_k)}
 =\frac{\Lambda_F^2}{\pi}
  \frac{(k+1)^2}{k(k+2)}
 =\frac{\Lambda_F^2}{\pi}
  \left(1+\frac{1}{k(k+2)}\right).}
\]
Consequently
\[
 \boxed{
 \lim_{k\to\infty}
 \frac{N_{\partial x}(k)}{A_H(\tau_k)}
 =\frac{\Lambda_F^2}{\pi}.}
\]
This is an area law for transverse spectral-channel capacity.  It is not an
entropy-area law and uses no fundamental position lattice.
\end{corollary}

\begin{corollary}[Exact affine Hopf-area/channel identity]
\label{p3f:cor:v259-affine-hopf-area-channel}
At every reached first-entry event \(\tau_k\),
\[
 \boxed{
 N_{\partial x}(k)
 =1+\frac{\Lambda_F^2}{\pi}A_H(\tau_k),}
 \qquad
 \boxed{
 A_H(\tau_k)
 =\frac{\pi}{\Lambda_F^2}\bigl(N_{\partial x}(k)-1\bigr).}
\]
Thus the transverse spectral-channel count is exactly affine in the metric
Hopf-base area at first entry, not merely asymptotically proportional to it.
This is a spectral-capacity identity; it is not an entropy law, a spacetime
area quantization rule, or a microscopic position-cell count.
\end{corollary}

\begin{proof}
At first entry,
\(A_H(\tau_k)=\pi k(k+2)/\Lambda_F^2\) and
\(N_{\partial x}(k)=(k+1)^2\).  Since
\(k(k+2)=(k+1)^2-1\), the displayed identities follow immediately.
\end{proof}

\begin{remark}[Cumulative bulk versus current transverse frontier]
\label{p3f:rem:cumulative-bulk-versus-current-transverse-frontier}
The accumulated band and the current frontier have different exact ranks:
\[
 \dim\mathcal P_x(k)
 =\frac{(k+1)(k+2)(2k+3)}6=\Theta(k^3),
 \qquad
 \dim\mathcal T_{x,k+1}^{\perp}
 =(k+1)^2=\Theta(k^2).
\]
The shell rank is exactly the discrete derivative of cumulative scalar
capacity:
\[
 \boxed{
 \dim\mathcal P_x(k)-\dim\mathcal P_x(k-1)
 =(k+1)^2=N_{\partial x}(k).}
\]
Consequently the fraction carried by the current shell is
\[
 \boxed{
 \frac{N_{\partial x}(k)}{\dim\mathcal P_x(k)}
 =\frac{6(k+1)}{(k+2)(2k+3)}\longrightarrow0.}
\]
Accordingly, the current shell theorem cannot by itself reconstruct arbitrary
data in all lower shells from one current transverse window; indeed the
current shell is an asymptotically vanishing fraction of the accumulated
scalar capacity.  What is derived is the shellwise correspondence.  The
accumulated band can be represented, as a graded vector space, by the
history-tagged direct sum of the successive transverse shell records, but a
one-boundary bulk reconstruction or an entropy law requires an additional
theorem.
\end{remark}


\begin{corollary}[Exact spectral-event e-fold spacing]
\label{p3f:cor:v191-frontier-efold-spacing}
Under the hypotheses of
Theorem~\ref{p3f:thm:spectral-frontier-event-process}, define
$N(\widetilde t)=\log a(\widetilde t)+N_0$.  For any two reached first-entry
events $\tau_j,\tau_k$,
\[
 \boxed{
 N(\tau_k)-N(\tau_j)
 =\frac12\log\frac{k(k+2)}{j(j+2)}.}
\]
Using \(N_{\partial x}(r)-1=r(r+2)\), the same identity is exactly
\[
 \boxed{
 N(\tau_k)-N(\tau_j)
 =\frac12\log
 \frac{N_{\partial x}(k)-1}{N_{\partial x}(j)-1},}
\]
or equivalently
\[
 \boxed{
 N_{\partial x}(k)-1
 =\bigl(N_{\partial x}(j)-1\bigr)
 e^{2[N(\tau_k)-N(\tau_j)]}.}
\]
The frontier scale $\Lambda_F$ cancels.  These are exact background
frontier-event/channel-capacity relations; they are not by themselves
observable perturbation widths, sky-multipole identifications, occupation
numbers, or a replacement for a state/visibility history.
\end{corollary}

\begin{proof}
At first entry, Theorem~\ref{p3f:thm:spectral-frontier-event-process} gives
$a(\tau_k)=\sqrt{k(k+2)}/\Lambda_F$.  Taking the logarithm of the ratio of
two such radii gives the formula.
\end{proof}

\begin{corollary}[Comparison-band capture event]
\label{p3f:cor:v191-comparison-band-capture}
Using the Paper-II-C derived primitive resolution $m_x^{\rm res}=1$, assume that $\tau_2$
is reached.  Then
\[
 a(\tau_2)=\frac{\sqrt8}{\Lambda_F},\qquad
 \mathcal P_x(2)=E_0^x\oplus E_1^x\oplus E_2^x,
 \qquad\dim\mathcal P_x(2)=14.
\]
Paper~II-C, Corollary~\ref{p2c:cor:v191-quadratic-comparison-specialization},
gives
\[
 \operatorname{Sym}^2E_1^x\cong E_0^x\oplus E_2^x.
\]
Hence this is the first physical position frontier that contains both the
primitive linear sector and its complete quadratic scalar descendant, and it
is equal, as a scalar band, to the separately derived comparison band with
$m_{\rm cmp}=2$.  We call $\tau_2$ the \emph{comparison-band capture event}.
By cumulative retention, \(\mathcal P_x(2)\) remains a physical subband after
capture.  Therefore any static comparison theorem whose only missing x-side
premise is the occurrence of \(\mathcal P_x(2)\)---including the fixed chiral
clock comparison and exact \(2I\) quartic evaluation on that band---may use
the retained physical subband from \(\tau_2\) onward without a second
comparison-band occurrence premise.

The equality of band labels at this event does not identify the roles:
$m_{\rm cmp}$ remains a fixed comparison label while $m_x$ continues to grow
with the position frontier.  It does not select a state or action member,
identify \(\mathcal P_x(2)\) with the momentum carrier, or make the \(2I\)
quadrature exact on arbitrary later high-mode fields.
\end{corollary}

\begin{theorem}[History-tagged transverse shell stack is information-complete]
\label{p3f:thm:v260-history-tagged-shell-stack}
At any reached sharp-frontier level $k$, the shellwise Born--Hopf
isomorphisms of Paper~I-A and Paper~II-C assemble to a canonical graded
isomorphism
\[
 \boxed{
 \mathcal P_x(k)
 \xrightarrow{\ \sim\ }
 \bigoplus_{j=0}^{k}\mathcal T_{x,j+1}^{\perp}
 \cong
 \bigoplus_{\ell=0}^{k}(k-\ell+1)\,\mathscr Y_\ell(S_H^2).}
\]
Thus the \emph{ordered, level-tagged history} of transverse shell symbols is
information-complete for the accumulated scalar band.  A single current
boundary window is not: the history tags and their multiplicities are part of
the reconstruction data.
\end{theorem}

\begin{proof}
For every $j$, Theorem~\ref{p2c:thm:born-transported-transverse-shell-reconstruction}
gives
$E_j^x\cong\mathcal T_{x,j+1}^{\perp}$, while
$\mathcal P_x(k)=\bigoplus_{j=0}^kE_j^x$.  Taking the direct sum gives the
first isomorphism.  Since
$\mathcal T_{x,j+1}^{\perp}=\bigoplus_{\ell=0}^{j}\mathscr Y_\ell(S_H^2)$,
reordering the finite sum gives multiplicity $k-\ell+1$ for degree $\ell$.
The dimension identity
\[
 \sum_{\ell=0}^{k}(k-\ell+1)(2\ell+1)
 =\frac{(k+1)(k+2)(2k+3)}6
\]
confirms the rank.  No entropy or one-boundary holographic statement is used.
\end{proof}

\begin{corollary}[Nested boundary increment and accumulated-capacity law]
\label{p3f:cor:v260-boundary-increment-capacity}
At consecutive first-entry events,
\[
 \boxed{
 \mathcal T_{x,k+1}^{\perp}/\mathcal T_{x,k}^{\perp}
 \cong\mathscr Y_k(S_H^2),\qquad
 N_{\partial x}(k)-N_{\partial x}(k-1)=2k+1,}
\]
and therefore
\[
 \boxed{
 A_H(\tau_k)-A_H(\tau_{k-1})
 =\frac{\pi}{\Lambda_F^2}(2k+1).}
\]
Writing $N=N_{\partial x}(k)=(k+1)^2$ and $s=\sqrt N=k+1$, the accumulated
scalar capacity is
\[
 \boxed{
 D_k:=\dim\mathcal P_x(k)
 =\frac{s(s+1)(2s+1)}6,
 \qquad
 s=\sqrt{1+\frac{\Lambda_F^2}{\pi}A_H(\tau_k)}.}
\]
Hence current boundary capacity grows as area while accumulated scalar
capacity grows as area$^{3/2}$ asymptotically.  This strengthens, rather than
weakens, the firewall against identifying one current shell with the whole
accumulated bulk.
\end{corollary}

\begin{theorem}[Sharp-frontier operation-activation hierarchy]
\label{p3f:thm:v260-operation-activation-hierarchy}
Let the nondecreasing sharp frontier have reached $\mathcal P_x(K)$.
Then:
\begin{enumerate}[label=(\roman*),leftmargin=2.1em]
\item every fixed scalar spectral multiplier in the Paper-II-D domain is exact
      on $\mathcal P_x(K)$ immediately at $\tau_K$;
\item the compressed isometry-derivation identity
      $[D_X^{(K)},M_f^{(K)}]=M_{D_Xf}^{(K)}$ likewise requires no additional
      degree guard on the complete invariant band;
\item a multiplication word with symbol degrees $s_1,\ldots,s_r$ acting on
      $\mathcal P_x(K)$ is frontier-contained once
      \[
       \boxed{m_x\ge K+S,\qquad S:=\sum_{j=1}^r s_j,}
      \]
      so the occurrence of its guard band is generated at $\tau_{K+S}$.
\end{enumerate}
The exact e-fold guard delay is
\[
 \boxed{
 \Delta N_{K\to K+S}
 =\frac12\log\frac{(K+S)(K+S+2)}{K(K+2)}.}
\]
For a degree-two scalar profile this gives full pointwise retention of
$h^2,h^3,h^4$ at $\tau_4,\tau_6,\tau_8$, respectively.  Exact integrated
finite-design moments can occur earlier; in particular the degree-two $2I$
quartic comparison is already exact on the retained comparison subband after
$\tau_2$.  Thus linear/spectral/derivation activation and nonlinear product
activation must not be assigned the same frontier premise.
\end{theorem}

\begin{proof}
Paper~II-D, Theorem~\ref{p2d:thm:scalar-fixed-band-recovery} and
Corollary~\ref{p2d:cor:fixed-band-process-readout-recovery}, give the fixed-band
spectral statements.  Paper~II-C,
Corollary~\ref{p2c:cor:guarded-transvection-multiplication}, removes the guard
for the displayed single derivation commutator and gives the multiplication
word bound $m\ge K+S$.  Corollary~\ref{p3f:cor:v259-derived-plateau-finite-band-activation}
turns each finite guard-band requirement into the corresponding sharp-frontier
entry event.  The e-fold formula follows from
Corollary~\ref{p3f:cor:v191-frontier-efold-spacing}.
\end{proof}

\begin{corollary}[No finite e-fold accumulation and asymptotic radial spacing]
\label{p3f:cor:v260-frontier-event-density}
Adjacent first-entry events satisfy
\[
 \boxed{
 \delta N_k
 =\frac12\log\frac{(k+1)(k+3)}{k(k+2)}
 \sim\frac1k\longrightarrow0,}
\]
while for every fixed $j$,
$N(\tau_k)-N(\tau_j)\to\infty$.  Thus infinitely many frontier events do
not accumulate at finite e-fold.  Their physical-radius spacing is
\[
 \boxed{
 a_{k+1}-a_k
 =\frac{2k+3}{\Lambda_F\left[
 \sqrt{(k+1)(k+3)}+\sqrt{k(k+2)}\right]}
 \longrightarrow\frac1{\Lambda_F}.}
\]
Hence the staircase becomes logarithmically dense but asymptotically uniform
in physical radius, without thereby supplying a physical-time law.
\end{corollary}

\begin{corollary}[Protected coexact-shell availability after scalar capture]
\label{p3f:cor:v260-coexact-shell-availability}
On the protected reciprocal-odd Hodge--Dirac branch of Paper~I-C, for every
$k\ge1$,
\[
 \operatorname{span}\{*(df\wedge dg):f\in E_1^-,\ g\in E_k^-\}
 =\mathcal C_k^-.
\]
Therefore once the sharp frontier has retained the scalar ingredients through
$E_k^x$ and the protected reciprocal-odd identification is in force, no
additional \emph{mathematical vector-shell availability} premise is required
to span $\mathcal C_k^-$.  This is a statement about the allowed one-form
response space only: it does not actualize a gauge field, photon, QCD state,
vector frontier, or light-front representation.
\end{corollary}

\begin{theorem}[Casimir form, Hopf-charge activation, and isotropy of the sharp scalar frontier]
\label{p3f:thm:v261-casimir-hopf-frontier}
On a complete round scalar band Paper~II-D gives
\[
 \sum_{i=1}^{3}Q_i^2=a^2\Delta_a.
\]
Therefore the sharp scalar frontier of
Definition~\ref{p3:def:sharp-scalar-resolution-frontier} is exactly
\[
 \boxed{
 \Pi_{\Lambda_F}(a)
 =\mathbf1_{[0,(a\Lambda_F)^2]}
 \!\left(\sum_{i=1}^{3}Q_i^2\right).}
\]
At a reached frontier level \(k\),
\[
 \operatorname{spec}(Q_H|_{\mathcal P_x(k)})
 =\{-k,-k+1,\ldots,k\},
\]
so geometric Hopf charge \(\pm q\) first becomes available when
\(m_x=|q|\); for \(|q|\ge2\), its first-entry radius is
\[
 \boxed{a_{|q|}=\frac{\sqrt{|q|(|q|+2)}}{\Lambda_F}.}
\]
The complete-band multiplicities obey \(d_k(q)=d_k(-q)\), hence every odd
Hopf-charge moment vanishes in the canonical complete-band trace, and
\[
 \boxed{
 \tau_k(Q_iQ_j)=\delta_{ij}\frac{k(k+3)}5.}
\]
Thus the unmarked sharp frontier preserves the complete chosen chiral
\(SU(2)\) action and carries neither a preferred Hopf axis nor a charge-sign
bias.  These are geometric scalar-frontier statements; \(q\) is not
hypercharge, electric charge, color, family number, or a sky multipole.
\end{theorem}

\begin{proof}
The Casimir identity is Paper~II-D,
Corollary~\ref{p2d:cor:chiral-scalar-casimir-identity}; spectral functional
calculus gives the projector formula.  Paper~II-D,
Proposition~\ref{p2d:prop:complete-position-hopf-multiplicities}, supplies the
complete Hopf spectrum and sign-symmetric multiplicities.  Summing the shell
Casimirs through \(k\) gives
\(\tau_k(\sum_iQ_i^2)=3k(k+3)/5\); chiral isotropy divides this equally among
the three generators.  The event radius follows from
Corollary~\ref{p3:cor:spectral-resolution-thresholds}.
\end{proof}

\begin{corollary}[First possible protected transverse onset occurs at \(\tau_2\)]
\label{p3f:cor:v261-first-protected-transverse-onset}
On the protected reciprocal-odd Hodge--Dirac branch, if a scalar profile
\(h\) lies in one nonconstant Laplace eigenspace \(E_k^-\), then
\[
 \Delta_0^{-1}h=\lambda_k^{-1}h,
 \qquad
 dh\wedge d\Delta_0^{-1}h=0,
\]
and the protected second-order transverse response vanishes.  Hence on the
primitive sharp-frontier plateau
\[
 \boxed{m_x=1\quad\Longrightarrow\quad\alpha_2^{\perp}=0.}
\]
Nonzero protected transverse response requires mixing distinct scalar
eigenvalues.  Since the reciprocal-odd primitive sector contains only the
nonconstant shell \(E_1^-\), \(\tau_2\) is the first sharp-frontier event at
which such a response is mathematically possible.  This does not actualize a
mixed profile, gauge field, or physical excitation.
\end{corollary}

\begin{proof}
Paper~I-D, Theorem~\ref{p1b:thm:dirac-transverse-generation}, identifies the
protected response with the coexact projection generated by
\(dh\wedge d\Delta_0^{-1}h\).  A single Laplace eigenvalue makes the two
one-forms proportional, so their wedge vanishes.  The sharp-frontier
staircase supplies the first second eigenvalue at \(\tau_2\).
\end{proof}

\begin{corollary}[Scalar capture, coexact availability, and same-cutoff vector admission are distinct]
\label{p3f:cor:v261-scalar-vector-frontier-delay}
At the scalar first-entry event
\[
 a_k^{(0)}=\frac{\sqrt{k(k+2)}}{\Lambda_F},
\]
the same-index coexact one-form shell has
\[
 \boxed{
 \lambda_k^{(1),\mathrm{co}}(a_k^{(0)})
 =\Lambda_F^2\frac{(k+1)^2}{k(k+2)}
 =\Lambda_F^2\left(1+\frac1{k(k+2)}\right)
 >\Lambda_F^2.}
\]
Thus scalar retention, algebraic/Hodge generation of the associated coexact
shell, and admission of that shell by a hypothetical same-numerical-cutoff
one-form frontier are different finite-curvature statements.  If the latter
frontier were introduced, the exact delay would be
\[
 \Delta N_k^{0\to1}
 =\frac12\log\left(1+\frac1{k(k+2)}\right)>0.
\]
The physical frontier constituted here is scalar; this corollary forbids using
its staircase silently as a universal cutoff for every spin or bundle
Laplacian.
\end{corollary}


\begin{corollary}[Scalar-frontier type firewall]
\label{p3f:cor:v2180-scalar-frontier-type-firewall}
The constituted sharp frontier
\[
 \Pi_{\Lambda_F}=\mathbf1_{[0,\Lambda_F^2]}(\Delta_h)
\]
is a scalar-test-space constructor.  Its occurrence does not by itself
constitute a same-numerical-cutoff frontier
\[
 \mathbf1_{[0,\Lambda_F^2]}(\Delta_E)
\]
on an arbitrary vector, spinor, color, flavor, or associated bundle $E$.
Likewise, a presentation $\mathcal P_x(m_x)\otimes V_{\rm int}$ may be used
for a downstream physical sector only after that sector and its physical
operator/action word have been independently generated on the product.
A bundle frontier with the same numerical value $\Lambda_F$ therefore
requires its own licensed constructor theorem; it cannot be inferred from the
scalar staircase alone.
\end{corollary}

\begin{proof}
Corollary~\ref{p3f:cor:v261-scalar-vector-frontier-delay} already separates
scalar capture, Hodge generation and admission by a hypothetical same-cutoff
one-form frontier.  The same typing argument applies a fortiori to a distinct
physical bundle/operator sort.  GP2 transports retained generated structure;
it does not change the sort of the constituting spectral projector or create a
new field-sector occurrence.
\end{proof}

\begin{corollary}[Generated heat and sharp frontier are common functional-calculus readouts]
\label{p3f:cor:v261-heat-sharp-functional-relation}
Suppose the sharp scalar frontier and the Paper-III-E gravitational heat role
act on the same protected scalar operator \(L\), with
\[
 A_g(\Lambda_g)=e^{-L/\Lambda_g^2}.
\]
Then
\[
 \boxed{
 \Pi_{\Lambda_F}
 =\mathbf1_{[\exp(-\Lambda_F^2/\Lambda_g^2),1]}(A_g),
 \qquad
 [A_g,\Pi_{\Lambda_F}]=0.}
\]
On the optional scale-matched branch \(\Lambda_F=\Lambda_g\),
\[
 \boxed{\Pi_{\Lambda_F}=\mathbf1_{[e^{-1},1]}(A_g).}
\]
Hence no independent commutation or operator-compatibility premise remains
between these two readouts on their common-operator branch.  Their physical
occurrences and the scale match remain separately typed.
\end{corollary}

\begin{corollary}[One current transverse window cannot encode the accumulated scalar band]
\label{p3f:cor:v261-current-shell-rank-obstruction}
At sharp-frontier level \(k\),
\[
 D_k:=\dim\mathcal P_x(k)
 =\frac{(k+1)(k+2)(2k+3)}6,
 \qquad
 N_{\partial x}(k)=(k+1)^2,
\]
and
\[
 \boxed{D_k-N_{\partial x}(k)=D_{k-1}.}
\]
Therefore every linear map
\[
 T:\mathcal P_x(k)\longrightarrow\mathcal T_{x,k+1}^{\perp}
\]
obeys
\[
 \boxed{\dim\ker T\ge D_{k-1}.}
\]
The history tags in
Theorem~\ref{p3f:thm:v260-history-tagged-shell-stack} are therefore necessary
for lossless linear reconstruction; a single current boundary window is
insufficient by an exact rank defect equal to the entire previously
accumulated scalar band.
\end{corollary}


\begin{definition}[Exact-doubled K\"unneth/Born data]
\label{p3:def:local_kunneth_born}
On the exact doubled bundle \(E_U=TU\oplus T^*U\), define the anchor and the
standard Dorfman bracket by
\[
 \rho(X+\xi):=X,
 \qquad
 [X+\xi,Y+\eta]_D
 :=[X,Y]+\mathcal L_X\eta-\iota_Yd\xi.
\]
Define the fibrewise tensors
\begin{align}
 \eta_B(X+\xi,Y+\eta)&:=\eta(X)+\xi(Y),\\
 \omega_E(X+\xi,Y+\eta)&:=\eta(X)-\xi(Y),\\
 K(X+\xi)&:=X-\xi.
\end{align}
Then \(K^2=\mathrm{id}\), its \(\pm1\) eigenbundles are
\[
        F=TU\oplus0,
        \qquad
        G=0\oplus T^*U,
\]
and
\[
        \omega_E(A,B)=\eta_B(KA,B).
\]
Both \(F\) and \(G\) are maximally \(\eta_B\)-isotropic and involutive under
the Dorfman bracket.  Thus \((E_U,\eta_B,K)\) is the canonical split exact
K\"unneth/para-Hermitian datum over \(U\).

For any positive metric \(g\) on \(TU\), define
\[
 H_g(X+\xi,Y+\eta)
 :=g(X,Y)+g^{-1}(\xi,\eta).
\]
Then \((\eta_B,\omega_E,H_g)\) is the associated exact-doubled Born datum.
The fibrewise form \(\omega_E\) is not to be confused with the canonical
symplectic form \(\omega_{\rm can}=d\vartheta\) on the phase space \(M_U=T^*U\).
A connection-induced splitting \(TM_U\cong\pi^*E_U\) pulls these tensors to
\(TM_U\), but no such splitting is part of their definition.
\end{definition}

\begin{remark}[Claim strength and global gluing boundary]
\label{p3:rem:local_born_claim_strength}
Definition~\ref{p3:def:local_kunneth_born} makes no separate global
Born-manifold postulate.  The local phase spaces glue by the cotangent lifts of
Theorem~\ref{p3:thm:local_reciprocal_atlas}; the exact doubled bundles glue by
the functorial tangent--cotangent representation
\[
        \phi_*\oplus(\phi^{-1})^*.
\]
These transitions preserve \(\eta_B\), \(\omega_E\), \(K\), and the pulled-back
metric datum \(H_g\).  Hence the exact-doubled Born package is globally defined
on the reciprocal-position domain whenever the base metric data glue.
Independent \(B\)-field, \(\beta\)-field, or \(O(d,d)\) twists would change the
transition cocycle and therefore enlarge the constituted datum.  They are not
alternative realizations of the same declared untwisted doubled-role datum.

Any later global or homogeneous realization of this local package requires a
separate downstream theorem.  No such realization is used here.
\end{remark}

\begin{theorem}[Unique anchored no-insertion Born connection]
\label{p3:thm:local_kunneth_connection}
Let \(U\) be a reciprocal-position patch and let \(g\) be a positive metric on
\(TU\).  Let \(\nabla^g\) be its Levi--Civita connection and use the same symbol
for the induced dual connection on \(T^*U\).  Then
\begin{equation}
\label{p3:eq:anchored-born-connection}
 \nabla^B_{X+\xi}(Y+\eta)
 :=\nabla^g_XY+\nabla^g_X\eta
\end{equation}
defines a generalized connection on \(E_U=TU\oplus T^*U\) with the following
properties:
\begin{enumerate}[label=(\roman*),leftmargin=2.2em]
\item its dependence on the first argument is only through the canonical
      anchor \(\rho(X+\xi)=X\);
\item it preserves the splitting \(E_U=F\oplus G\) and the tensors
      \(\eta_B\), \(\omega_E\), \(K\), and \(H_g\);
\item its restriction to \(F\cong TU\) is exactly the Levi--Civita connection,
      so its ordinary projected torsion vanishes;
\item it is the unique generalized connection satisfying (i), preserving
      \(F\), \(G\), \(\eta_B\), and \(H_g\), and having torsion-free restriction
      to \(F\);
\item under reciprocal-position coordinate changes it transforms naturally
      with \(\phi_*\oplus(\phi^{-1})^*\).  Consequently the local connections
      glue whenever the base metrics glue.
\end{enumerate}
No horizontal distribution on \(T(T^*U)\), generalized-torsion convention, or
additional Born twist is required.
\end{theorem}

\begin{proof}
For \(f\in C^\infty(U)\), the Levi--Civita and dual-connection identities give
\begin{align*}
 \nabla^B_{f(X+\xi)}(Y+\eta)
 &=f\nabla^B_{X+\xi}(Y+\eta),\\
 \nabla^B_{X+\xi}\bigl(f(Y+\eta)\bigr)
 &=X(f)(Y+\eta)+f\nabla^B_{X+\xi}(Y+\eta).
\end{align*}
Thus \eqref{p3:eq:anchored-born-connection} is a generalized connection with
anchor \(\rho\).  It plainly preserves \(F\), \(G\), and \(K\).

Metric compatibility of \(\nabla^g\) and of its dual gives
\[
 X\bigl(\eta(Y)\bigr)
 = (\nabla^g_X\eta)(Y)+\eta(\nabla^g_XY).
\]
Applying this identity to both cross terms proves \(\nabla^B\eta_B=0\) and
\(\nabla^B\omega_E=0\).  The equalities \(\nabla^gg=0\) and
\(\nabla^gg^{-1}=0\) prove \(\nabla^BH_g=0\).  On \(F\), equation
\eqref{p3:eq:anchored-born-connection} reduces to
\[
        \nabla^B_XY=\nabla^g_XY,
\]
whose ordinary torsion is zero.

For uniqueness, let \(\widetilde\nabla\) satisfy the stated conditions.
Anchor-only dependence reduces its first argument to a vector field \(X\).
Preservation of \(F\) gives a connection \(D\) on \(TU\).  Since
\(H_g|_F=g\), preservation of \(H_g\) makes \(D\) metric compatible; the
assumed torsion-free restriction then forces \(D=\nabla^g\) by uniqueness of
the Levi--Civita connection.  Preservation of the evaluation pairing, which
is the \(F\)-\(G\) part of \(\eta_B\), forces the connection on \(G=T^*U\) to
be the dual of \(D\).  Hence \(\widetilde\nabla=\nabla^B\).

Finally, Levi--Civita formation is natural under diffeomorphisms, and the dual
connection transforms contragrediently.  Formula
\eqref{p3:eq:anchored-born-connection} is therefore invariant under
\(\phi_*\oplus(\phi^{-1})^*\), proving the overlap and gluing statement.
\end{proof}

\begin{definition}[Macroscopically local infrared sector]
\label{p3:def:macrolocal_sector}
A coarse-grained infrared sector is called \emph{macroscopically local} if, on an open cover by
patches \(U\), it admits local doubled data of the form above such that:
\begin{enumerate}[label=(\roman*),leftmargin=2.2em]
\item the compressed multiplication observables satisfy the quantitative
      separated-support estimate \eqref{p3:eq:quantified-qloc};
\item the position summand \(F=TU\oplus0\subset E_U\) is involutive under the standard Dorfman bracket;
\item the projected connection gives a Lorentzian metric effective dynamics whose leading
nontrivial term is two-derivative on each patch.
\end{enumerate}
\end{definition}

\begin{theorem}[Finite-propagation Fourier criterion for \((\mathrm{QLOC})\)]
\label{p3f:thm:qloc-fourier-criterion}
Let \(D\) be self-adjoint and suppose its wave group obeys the finite-
propagation relation with speed at most \(v\).  Write
\[
 Q_\Lambda=q(D^2/\Lambda^2)=r(D/\Lambda),
 \qquad r(s):=q(s^2).
\]
Assume \(r\) is bounded and the Fourier transform of \(r^2\) is integrable
with tails
\[
 \int_{|t|\ge R}|\widehat{r^2}(t)|\,dt
 \le C_N(1+R)^{-N}
 \quad\text{for every }N.
\]
Then
\[
 \|\mathbf1_EQ_\Lambda^2\mathbf1_F\|
 \le C_N'
 \bigl(1+\Lambda d_g(E,F)/v\bigr)^{-N},
\]
and the factorization in
Definition~\ref{p3:def:lat-pol-csa-gred-hypotheses} proves
\((\mathrm{QLOC})\).  A sufficient profile condition is
\(r^2\in\mathcal S(\mathbb R)\); compact Fourier support gives exact
finite-range vanishing beyond the corresponding propagation radius.  A
complete self-adjoint Dirac-type representative is a standard sufficient
geometric realization of the finite-propagation premise, but completeness and
Dirac type are not independent inputs to this abstract implication.
\end{theorem}

\begin{proof}
Fourier functional calculus gives
\[
 Q_\Lambda^2
 =\frac1{2\pi}\int_{\mathbb R}
 \widehat{r^2}(t)e^{itD/\Lambda}\,dt.
\]
Finite propagation makes
\(\mathbf1_Ee^{itD/\Lambda}\mathbf1_F=0\) whenever
\(|t|<\Lambda d_g(E,F)/v\).  Taking norms leaves only the stated Fourier
tail.  The commutator factorization and boundedness of \(Q_\Lambda\) complete
the estimate.
\end{proof}

\begin{proposition}[Davies--Gaffney Laplace-mixture criterion for
\((\mathrm{QLOC})\)]
\label{p3:prop:qloc-sufficient-criterion}
Let \(L=D^2\ge0\) extend to a complete representative satisfying the
Davies--Gaffney estimate.  Suppose the coarse profile has the representation
\[
 q(z)=\int_{(0,\infty)}e^{-tz}\,d\nu(t),
\]
where \(\nu\) is a finite real or complex Borel measure.  Put
\(\omega:=\nu*\nu\) and assume
For an all-orders conclusion assume
\begin{equation}
\label{p3f:eq:lap-mom}
 (\mathrm{LAP\!-!MOM}):\qquad
 \|\omega\|_{\rm TV}<\infty,
 \qquad
 M_m:=\int_0^\infty s^m\,d|\omega|(s)<\infty
 \quad(m=0,1,2,\ldots).
\end{equation}
Then the compressed observables
\(A_{a,\Lambda}=q(D^2/\Lambda^2)M_aq(D^2/\Lambda^2)\) satisfy
\eqref{p3:eq:quantified-qloc}.  Thus \((\mathrm{QLOC})\) is a theorem on this
analytic subdomain.  For a single prescribed decay order \(N\), it is enough
to assume total variation and one moment \(M_m<\infty\) with \(2m\ge N\);
all positive moments are needed only to obtain every polynomial order.
Outside this analytic subdomain, the displayed estimate remains an explicit
testable hypothesis rather than a consequence of doubled geometry.  Only
positive moments are required for this operator-norm statement; differentiated
pointwise kernel estimates may require additional small-time conditions or
negative moments.  In the pure heat branch \(q(z)=e^{-tz}\),
Davies--Gaffney gives the stronger
exponential form
\[
 \|\mathbf1_EQ_\Lambda^2\mathbf1_F\|
 \le
 \exp\!\left(-\frac{\Lambda^2d_g(E,F)^2}{8t}\right).
\]
\end{proposition}

\begin{proof}
Functional calculus and convolution give
\[
 Q_\Lambda^2=q(L/\Lambda^2)^2
 =\int_0^\infty e^{-sL/\Lambda^2}\,d\omega(s),
 \qquad \|Q_\Lambda\|\le\|\nu\|_{\rm TV}.
\]
For measurable sets \(E,F\) with \(d=d_g(E,F)>0\), Davies--Gaffney yields
\[
 \|\mathbf1_EQ_\Lambda^2\mathbf1_F\|
 \le\int_0^\infty
 \exp\!\left(-\frac{\Lambda^2d^2}{4s}\right)d|\omega|(s).
\]
Let \(R=\Lambda d\).  For \(R\ge1\) and every integer \(m\ge0\),
\(e^{-R^2/(4s)}\le C_ms^mR^{-2m}\); integration and
\eqref{p3f:eq:lap-mom} give arbitrary polynomial decay.  For \(R<1\), total
variation gives the uniform bound.  Choosing \(2m\ge N\) proves
\[
 \|\mathbf1_EQ_\Lambda^2\mathbf1_F\|
 \le C_N(1+\Lambda d)^{-N}.
\]
Finally the two products
\(Q_\Lambda M_aQ_\Lambda^2M_bQ_\Lambda\) and its reverse are bounded using
the separated-support estimate and
\(2\|Q_\Lambda\|^2\|a\|_\infty\|b\|_\infty\).  This proves
\eqref{p3:eq:quantified-qloc}.  For \(\nu=\delta_t\), one has
\(\omega=\delta_{2t}\), giving the displayed exponential bound.
\end{proof}


\section{One action, joint operators, and constrained variation}
\label{paper3f:sec:organized-action}
\subsection{Common-heat specialization of the joint kinetic trace}
\label{p3f:sec:common-heat-gauge-normalization}

The broad trace-closure theorem above classifies every admissible uninserted
parity-even quadratic curvature functional and therefore leaves one positive
coefficient.  The source-native heat branch is stronger: Paper~III-E already
generates the unit-preserving profile $F(u)=e^{-u}$, and on $\mathrm{JOINT}_x$
all retained metric and nonmetric terms are read from one fluctuated
position-sector Dirac datum.  The common position functional also fixes the
nonmetric role weight to one half.  The following theorem records the
consequence for the gauge coefficient without importing any later
representation multiplicity.

\begin{theorem}[Common-heat specialization removes the free kinetic scalar]
\label{p3f:thm:common-heat-specializes-kinetic-trace}
Use the premise-minimal heat role of Paper~III-E, the derived
$\mathrm{JOINT}_x^{\rm der}$ certificate, and
the unmarked closed spectral branch $\mathcal D_{\rm cl}$.  Let $D_A$ be the
resulting local self-adjoint Dirac-type datum on its least complex
four-dimensional Clifford module tensored with a finite unitary internal
carrier $E$.  Then the parity-even pure-gauge term in the common position
heat action is
\[
 \boxed{
 S_{\rm YM}^{\rm heat}[A]
 =\frac1{48\pi^2}
   \int_U\operatorname{Tr}_{E}
   \!\left(\mathcal F_{\mu\nu}^{\dagger}
            \mathcal F^{\mu\nu}\right)d\mu_g .}
\]
In particular, the positive coefficient $C$ of the broader invariant
classification is not an additional member on this common-heat branch: it is
the displayed heat-coefficient readout of the generated Dirac datum and
ordinary trace.
\end{theorem}

\begin{proof}
For a four-dimensional twisted Dirac operator the gauge part of the local
$a_4$ coefficient is
\[
 (4\pi)^{-2}\frac1{360}
 \operatorname{tr}\!\left(180E_F^2+30\Omega_F^2\right).
\]
With $E_F=\tfrac12\gamma^{\mu\nu}\mathcal F_{\mu\nu}$ and the rank-four
complex Clifford trace,
\[
 \operatorname{tr}_{\mathbb S}(E_F^2)
 =-2\operatorname{Tr}_{E}(\mathcal F_{\mu\nu}\mathcal F^{\mu\nu}),
 \qquad
 \operatorname{tr}_{\mathbb S}(\Omega_F^2)
 =4\operatorname{Tr}_{E}(\mathcal F_{\mu\nu}\mathcal F^{\mu\nu}).
\]
Hence the raw heat trace contributes
$\frac1{24\pi^2}\int\operatorname{Tr}_{E}(\mathcal F^\dagger\mathcal F)$.
The common position functional assigns every nonmetric raw term the already
derived factor $1/2$, giving the stated $1/(48\pi^2)$ coefficient.  An
independent rescaling of this term would no longer be the same generated heat
functional; a weighted trace would be an inserted endomorphism and is outside
the unmarked branch.
\end{proof}

\begin{remark}[Broad kinetic class versus source-native heat branch]
\label{p3f:rem:broad-versus-heat-kinetic-class}
The earlier theorem
$C\int\operatorname{Tr}(\Omega^\dagger\Omega)$ remains the correct
classification of the larger parity-even positive kinetic class.  It is not
retracted.  The present theorem says only that the stronger source-native
common-heat realization selects one member of that class once the local
Dirac/trace carrier is fixed.  Later physical action adoption may transport
this coefficient but may not recalibrate it without leaving the same branch.
\end{remark}

\begin{corollary}[Profile-invariant JOINT heat descendants]
\label{p3f:cor:profile-invariant-joint-heat-descendants}
On the premise-minimal branch of
Theorem~\ref{p3f:thm:common-heat-specializes-kinetic-trace}, the metric and
pure-gauge coefficients, their ordinary-trace variations, and every later
same-process descendant are independent of any unselected spectral-profile
representative.  They inherit the singleton fibre
\(mathfrak P^{\rm heat}_{\Lambda_{\rm EH}^2}\) of Paper~III-E.  A
descendant with an independently supplied profile, moment, or cutoff is an
enlarged process root and does not belong to the source-native prediction
chain.
\end{corollary}

\begin{proof}
Paper~III-E
Corollary~\ref{p3e:cor:singleton-heat-profile-fibre} fixes the whole heat
profile and its positive scale representative at each Einstein--Hilbert
coefficient.  The \(\mathrm{JOINT}_x\) theorem constructs the metric and
nonmetric terms from one fluctuated Dirac datum with that same functional,
while the theorem above obtains the gauge coefficient by ordinary trace and
the already generated nonmetric weight.  These are licensed same-process
maps, so Paper~III-B
Theorem~\ref{p3b:thm:gp-profile-invariance-factorization} applies.  Adding a
new profile changes the frozen root and is excluded by GP2 no-insertion.
\end{proof}


\subsection{Finite background-gauge quotient response}
\label{p3f:sec:v2140-background-gauge-response}

The local Yang--Mills action, Gauss quotient, and local gauge inverse above
already determine the infinitesimal gauge-orbit complex.  The result below
records the finite spectral-window representative needed by later one-loop
calculations.  It is a quotient-coordinate theorem, not a new particle or
interaction occurrence.

\begin{theorem}[Generated Yang--Mills background-gauge response complex]
\label{p3f:thm:v2140-background-gauge-response-complex}
Work on the Euclidean real-form sibling of
Theorem~\ref{p3f:thm:v174-generated-joint-gauge-phase-space-action}, and let
$\bar A$ be a smooth stationary background in a regular local quotient chart
where Proposition~\ref{p3f:prop:m8td-local-gauge-inverse} applies.  Let
$V_{0,N}\subset\Omega^0(\operatorname{ad}P)$ and
$V_{1,N}\subset\Omega^1(\operatorname{ad}P)$ be finite spectral windows
compatible with $d_{\bar A}$ and $d_{\bar A}^{\dagger}$.  Use the actual
finite integration measure and its transported density $\mu_N(\bar A)$,
including all inherited normalization and quotient factors.  Split off
\emph{all} exact zero modes in both operators below, not only stabilizer
modes; their collective-coordinate, stabilizer-volume or separately
retained treatment must be supplied by that same measure.  Assume constant
retained ranks on the comparison chart and a strictly positive physical
nonzero quadratic block.  A negative nonzero mode is outside this real
positive-Gaussian claim unless the actual quantization supplies its contour
and phase.  For the background-gauge coordinate
\[
 \chi_{\bar A}(a)=d_{\bar A}^{\dagger}a,
\]
the linearized orbit map and its Jacobian operator are
\[
 \epsilon\longmapsto d_{\bar A}\epsilon,
 \qquad
 D\chi_{\bar A}(d_{\bar A}\epsilon)
 =\Delta_0\epsilon,
 \qquad
 \boxed{\Delta_0=d_{\bar A}^{\dagger}d_{\bar A}}.
\]
Adding the unit background-gauge coordinate term
$\frac12\|d_{\bar A}^{\dagger}a\|^2$ to the quadratic representative of the
already generated Euclidean Yang--Mills action gives the minimal one-form
operator
\begin{equation}
\label{p3f:eq:v2140-background-vector-operator}
 \boxed{
 (\Delta_1 a)_\mu
 =-D_{\bar A}^{\nu}D_{\bar A\nu}a_\mu
   -2[\bar F_{\mu\nu},a^\nu]
   +R_{\mu}{}^{\nu}a_\nu .}
\end{equation}
Consequently the regular finite background-gauge representative of the
quotient has the relative Gaussian response
\begin{equation}
\label{p3f:eq:v2140-background-gauge-relative-response}
 \boxed{
 \Gamma_{N}^{(1)}(\bar A)-\Gamma_{N}^{(1)}(0)
 =\frac12\log\frac{\det'\Delta_1(\bar A)}{\det'\Delta_1(0)}
  -\log\frac{\det'\Delta_0(\bar A)}{\det'\Delta_0(0)}
  -\log\frac{\mu_N(\bar A)}{\mu_N(0)} .}
\end{equation}
Here each prime omits every exact zero eigenvalue of that operator, with
its separately declared treatment retained in $\mu_N$ or as an explicit
collective coordinate.  The second determinant is the Jacobian of the
gauge-orbit quotient.  It is not a physical scalar particle determinant and
creates no additional particle owner.  The measure ratio cancels only when
its equality is actually proved.  On a licensed contour continuation,
negative-mode phases and contour Jacobians are retained as additional
terms; replacing a determinant by its absolute value without them does not
extend this theorem.
\end{theorem}

\begin{proof}
Gauge covariance of the generated action gives the infinitesimal orbit
$a\mapsto a+d_{\bar A}\epsilon$.  Differentiating the displayed gauge
coordinate along that orbit gives
$d_{\bar A}^{\dagger}d_{\bar A}\epsilon$, so ordinary finite-dimensional
change of variables on $V_{0,N}$ gives the determinant Jacobian $\det'\Delta_0$.
The second variation of the generated Euclidean Yang--Mills action plus the
unit gauge-coordinate quadratic term is the standard Weitzenbock form on
adjoint one-forms, which is
\eqref{p3f:eq:v2140-background-vector-operator}; the Ricci term acts only on
the one-form index.  On the stipulated positive physical nonzero block,
finite Gaussian integration in the actual measure gives
$\mu_N(\bar A)(\det'\Delta_1)^{-1/2}$, while the quotient Jacobian
contributes $\det'\Delta_0$.  Every omitted zero direction is handled by
the separately supplied measure/collective-coordinate rule.  Taking minus
the logarithm and dividing by the corresponding zero-background
representative gives
\eqref{p3f:eq:v2140-background-gauge-relative-response}.  The local inverse
hypothesis removes non-stabilizer gauge-copy degeneracy on the stated chart.
\end{proof}

\begin{corollary}[Gauge-orbit determinant is auxiliary response data]
\label{p3f:cor:v2140-ghost-not-particle}
On the domain of
Theorem~\ref{p3f:thm:v2140-background-gauge-response-complex}, any Grassmann
presentation of $\det'\Delta_0$ is an algebraic representation of the
already generated quotient Jacobian.  It does not constitute an additional
scalar field occurrence, particle species, or independent action owner.
\end{corollary}

\begin{remark}[Finite quotient export versus quantum completion]
The theorem supplies the action-owned quadratic operators and finite regular
quotient Jacobian.  It does not construct an infinite-dimensional gauge-field
measure, prove regulator removal, or select a renormalization prescription.
Those remain later quantum/continuum interfaces.
\end{remark}

\subsection{Degree-two scalar comparison action and exceptional finite sources}
\label{p3f:sec:v180-degree-two-scalar-comparison-action}

The statements in this section use the position-side comparison band
$m_{\rm cmp}=2$.  The physical compact momentum truncation remains $m_p=1$,
and the actual frontier occurrence/history remains contingent, while $m_x$ is
the derived least-resolution statistic of that history and may grow beyond two.


\begin{theorem}[Natural Born-even second-order relative scalar selector]
\label{p3f:thm:v181-natural-relative-scalar-selector}
On the generated relative base
$K_{\rm rel}\cong SU(2)_{\rm rel}\times U(1)_{\rm rel}$, consider scalar
local differential operators of order at most two that are bi-invariant on
the $SU(2)$ factor, translation invariant on the relative circle, and even
under Born inversion.  Every such operator has the form
\[
 \boxed{L=A\,\mathsf C_{SU(2)}+B\,N_{\rm rel}^2+C\,I,}
\]
where $N_{\rm rel}$ is the generator of
Corollary~\ref{p3b:cor:v181-relative-scalar-generator}.  No Born-even term
linear in $N_{\rm rel}$ and no mixed first-order $SU(2)$--circle term is
available in this symmetry class.
\end{theorem}
\begin{proof}
Use the relative-group classification in Paper~I-C,
Corollary~\ref{p1b:cor:second-order} \cite{Kugelmann2026Paper1C},
on \(K_{\rm rel}\) with its retained relative-circle generator.
In differential order at most two, its inversion-even algebra
\(\mathbb C[\mathsf C_{SU(2)},N_{\rm rel}^2]\) has precisely the
three displayed terms, with the same symmetry class and coefficients.
\end{proof}

\begin{corollary}[Clock-linked degree-two selector spectrum]
\label{p3f:cor:v181-degree-two-relative-selector-spectrum}
On the $m_{\rm cmp}=2$ scalar comparison, the eigenvalues of $L$ are
\[
 \boxed{C,\quad 3A+B+C,\quad 8A+C,\quad 8A+4B+C}
\]
with multiplicities $1,4,3,6$, respectively, by Paper~I-C,
Corollary~\ref{p1b:result:charge-linked-degree-two-spectra-and-polynomial-integrals}
\cite{Kugelmann2026Paper1C}, using the retained clock--Hopf charge link.
The framework supplies no
independent scalar-circle metric modulus on this native branch, so this
classification fixes the form of the selector but not the ratio $B/A$ or an
absolute numerical cutoff.
\end{corollary}

\begin{theorem}[Exact finite representation of the local quartic scalar action]
\label{p3f:thm:v180-exact-finite-quartic-scalar-action}
Freeze the same stationary round spatial metric, compatible background, and
nonzero local quartic scalar block admitted by the native local-action theorem.
Let $\{f_A\}_{A=1}^{14}$ be a real orthonormal basis of
\[
 \mathcal P_x(2)=E_0^x\oplus E_1^x\oplus E_2^x,
 \qquad
 \mathsf C:=a^2\Delta_a=\operatorname{diag}(0,3I_4,8I_9),
\]
and write
\[
 h(t,x)=V_a^{-1/2}\sum_Az_A(t)f_A(x),\qquad V_a=2\pi^2a^3.
\]
For a canonically normalized real scalar direction with local potential
$\frac12m_{\rm loc}^2h^2+\lambda vh^3+\frac\lambda4h^4$, the restricted mode
Lagrangian is
\[
\begin{aligned}
 L_{\rm mode}={}&\frac12\dot z^T\dot z
 -\frac12z^T\bigl(a^{-2}\mathsf C+m_{\rm loc}^2I\bigr)z\\
 &-\frac{\lambda v}{\sqrt{V_a}}T_{ABC}z_Az_Bz_C
 -\frac{\lambda}{4V_a}T_{ABCD}z_Az_Bz_Cz_D,
\end{aligned}
\]
where, for $r=2,3,4$,
\[
 \boxed{
 T_{A_1\ldots A_r}
 =\int_{S^3}f_{A_1}\cdots f_{A_r}\,d\sigma
 =\frac1{120}\sum_{x\in2I}f_{A_1}(x)\cdots f_{A_r}(x).}
\]
Thus the finite $2I$ expression reproduces the kinetic pairing, Hessian,
cubic vertex, quartic vertex, and every restricted derivative of this local
polynomial action exactly.
\end{theorem}
\begin{proof}
The physical kinetic normalization and \(V_a^{-1/2}\) fix the
quadratic mode pairing and the displayed powers of \(V_a\).
Paper~I-C, Theorem~\ref{p1b:thm:design} and
Corollary~\ref{p1b:result:charge-linked-degree-two-spectra-and-polynomial-integrals}
\cite{Kugelmann2026Paper1C}, give exact kinetic pairings and polynomial
moments on this \(14\)-dimensional comparison band. Applying those
identities at each supplied time to the already normalized local density
gives \(L_{\rm mode}\); polynomial equality gives all restricted vertices.
\end{proof}

\subsection{Untruncated scalar harmonic dynamics and occurrence boundaries}
\label{p3f:subsec:v3067-scalar-cascade-and-guard}
\begin{theorem}[Untruncated quartic evolution has no nonconstant finite harmonic invariant band]
\label{p3f:thm:v15-quartic-finite-band-no-go}
On the frozen-round one-real-scalar comparison branch of
Theorem~\ref{p3f:thm:v180-exact-finite-quartic-scalar-action},
write $F_K=\bigoplus_{j=0}^{K}E_j^x$ and suppose the full
Euler--Lagrange equation is
\[
 \ddot h-\Delta_a h+m_{\rm loc}^2h+3\lambda v h^2+\lambda h^3=0,
 \qquad \lambda\ne0,
\]
with no external position-dependent source or other field contributing
opposite grade-$3K$ force.  Here $\Delta_a$ has eigenvalues
$-j(j+2)/a^2$ on $E_j^x$.  No spatially nonconstant solution can
remain in a fixed finite band $F_K$ on a nonempty open time interval.
In fact, for every $h\in F_K$ with nonzero highest harmonic
$h_K\in E_K^x$ and $K\ge1$, one has
\[
 \boxed{P_{E_{3K}^x}(h^3)\ne0.}
\]
\end{theorem}
\begin{proof}
Let $H_K\ne0$ be the homogeneous harmonic polynomial of degree $K$
representing $h_K$ in $\mathbb R^4$.  Fischer decomposition of
$H_K^3$ shows that vanishing of its degree-$3K$ harmonic part is
equivalent to divisibility by $Q=\sum_{a=1}^{4}x_a^2$.
This rank-four quadratic is irreducible over $\mathbb C$ and therefore
prime in $\mathbb C[x_1,\ldots,x_4]$.  If $Q\mid H_K^3$,
then $Q\mid H_K$, which contradicts the direct Fischer decomposition
of a nonzero harmonic polynomial. Thus $P_{E_{3K}}h^3\ne0$.
The time and Laplace derivatives remain in $F_K$, and $h^2\in F_{2K}$;
projecting the full field equation onto $E_{3K}$ gives a contradiction.
A Galerkin-restricted equation $P_{F_K}\mathcal E[h]=0$ does not
include this omitted projection and is not contradicted.
\end{proof}

\begin{proposition}[Exact first off-band quartic source and exceptional cubature alias]
\label{p3f:prop:v15-quartic-offband-cubature}
For $u_1$ a unit-round coordinate and $U_j$ the zonal
$SU(2)$ character harmonic of degree $j$, the exact identity
\[
 U_2^3=U_0+3U_2+2U_4+U_6
\]
shows that the equation with $h(0,u)=B U_2(u_1)$ and
$\dot h(0,u)=0$ has $P_{E_6}\ddot h(0)=-\lambda B^3U_6$.
The exact overlap is
\[
 \boxed{\int_{S^3}U_6U_2^3\,d\sigma=1,
 \quad \frac1{120}\sum_{q\in2I}U_6(q_0)U_2(q_0)^3=2,
 \quad \sum_{z\in V\cup F}w_zU_6(z_0)U_2(z_0)^3=1.}
\]
In the last term the existing positive 720-node weights are
$w_v=1/504$ on the 120 vertices and $w_f=2/1575$ on the
600 tetrahedral facet centers.
\end{proposition}
\begin{proof}
Clebsch--Gordan multiplication of the $SU(2)$ characters gives the
first identity, and their orthonormality gives the continuum value.
The restricted $2I$ character has norm squared
$\langle U_6,U_6\rangle_{2I}=2$ and is orthogonal to
$U_0,U_2,U_4$, giving the 120-node value.  The integrand has
coordinate degree twelve; Paper I-B's positive degree-nineteen
vertex--facet cubature integrates it exactly, giving the third value.
\end{proof}


\begin{corollary}[Exact quartic comparison and off-band-force cubature domains]
\label{p3f:cor:v15-cubature-comparison-bandwidth}
On the frozen-round real-scalar polynomial comparison domain of
Theorem~\ref{p3f:thm:v180-exact-finite-quartic-scalar-action},
let a positive scalar cubature be exact through ambient polynomial degree
$d$. The restricted quartic action on the complete scalar harmonic band
$F_K$ is integrated exactly whenever $4K\le d$. Every weak cubic-force
pairing against test harmonics through degree $3K$ is integrated exactly
whenever $6K\le d$. For the existing $120$-vertex, $720$-node
vertex--facet, and $12{,}720$-node scalar rules, with degrees $d=11,19,41$
respectively, the respective sufficient whole-band bounds are
\[
 \begin{array}{c|ccc}
 d & 11 & 19 & 41 \\\hline
 K\text{ for full quartic action} &2&4&10\\
 K\text{ for full cubic off-band force} &1&3&6
 \end{array}
\]
The extended rules are the mathematical source geometries of
Paper~I-B, not a new physical set of spatial detectors or an
additional physical position test frontier. These are sufficient
uniform polynomial exactness bounds, not no-go statements for every
individual higher-degree moment.
\end{corollary}
\begin{proof}
Every product of four harmonics in $F_K$ has ambient degree at most
$4K$. The full weak cubic force pairs a degree-at-most-$3K$ test
harmonic against a cubic field product of degree at most $3K$, giving
ambient degree at most $6K$. The claimed integers follow by taking
integer floors of $d/4$ and $d/6$ and using the established positive
finite integration rules.
\end{proof}

\begin{theorem}[Strict instability of nonconstant stationary real scalars on round $S^3$]
\label{p3f:thm:v15-round-scalar-static-instability}
Suppose $a>0$, $V\in C^2$, and $h\in C^3(S^3(a),\mathbb R)$
solves $-\Delta_a h+V'(h)=0$ for a frozen-round scalar-only action
with positive quadratic gradient energy. If $h$ is spatially nonconstant,
its full scalar Jacobi operator $\mathcal J=-\Delta_a+V''(h)$ obeys
\[
 \boxed{\lambda_{\min}(\mathcal J)\le-2/a^2<0.}
\]
\end{theorem}
\begin{proof}
The Bochner identity and $\mathrm{Ric}=2a^{-2}g$ give
\[
 \int\bigl(|\nabla^2h|^2+V''(h)|\nabla h|^2
            +2a^{-2}|\nabla h|^2\bigr)=0.
\]
The Sobolev test function $\eta=|\nabla h|$ satisfies
$|\nabla\eta|^2\le|\nabla^2h|^2$ almost everywhere by Kato.
Hence its Jacobi Rayleigh quotient is at most $-2/a^2$.
The result concerns the complete scalar variation space, not a
Galerkin subspace or a coupled metric/gauge Hessian.
\end{proof}

\begin{proposition}[Action-generated leading odd harmonic cascade]
\label{p3f:prop:v15-local-odd-harmonic-cascade}
On the untruncated frozen-round scalar equation of
Theorem~\ref{p3f:thm:v15-quartic-finite-band-no-go}, suppose an
analytic local solution has initial conditions
$h(0,u)=A u_1$, $\dot h(0,u)=0$, and $A\lambda\ne0$.
Define $b_0=1$ and
\[
 b_{n+1}=\frac{\sum_{i+j+k=n}b_ib_jb_k}
 {(2n+2)(2n+1)} >0.
\]
Then the first nonzero Taylor coefficient in every odd shell is
\[
 \boxed{P_{E_{2n+1}}h(t)
 =\frac{(-1)^n b_n\lambda^n A^{2n+1}}{2^{2n+1}}
 t^{2n}U_{2n+1}(u_1)+O(t^{2n+2}),\quad n\ge1.}
\]
In particular the $E_3,E_5,E_7$ coefficients are
$-\lambda A^3t^2/16$, $\lambda^2A^5t^4/256$, and
$-3\lambda^3A^7t^6/10240$, respectively.
\end{proposition}
\begin{proof}
Expand $h=\sum_{n\ge0}h_nt^{2n}$ and use
$\Delta_a f(x)=a^{-2}[(1-x^2)f''-3xf']$ on zonal polynomials.
Induction gives $\deg h_n\le2n+1$. The coefficient of
$x^{2n+3}$ in $h_{n+1}$ can only come from the cubic product;
its recurrence is the displayed positive $b_n$ equation with alternating
signs. Since $U_{2n+1}$ has leading coefficient $2^{2n+1}$ and
lower time orders cannot contain harmonic degree $2n+1$, the result
follows by harmonic projection.
\end{proof}

\begin{remark}[Dynamical instability band is not a GP2 test-frontier occurrence]
\label{p3f:rem:v15-instability-not-test-frontier}
At a supplied homogeneous stationary scalar background $\bar h$
with $V_{\rm eff}''(\bar h)=-\mu^2<0$ on a frozen round branch,
linear scalar modes have squared frequencies
$k(k+2)/a^2-\mu^2$.  The exact linearly unstable degrees are
$k(k+2)<a^2\mu^2$, with crossing radii
$a_k=\sqrt{k(k+2)}/\mu$.  This is an existing-action Hessian
criterion for the growth of field fluctuations, not the separate
occurrence of the complete sharp physical test frontier
$\mathsf{SF}_{\Lambda_F}$ or a licence to identify
$\mu=\Lambda_F$.  Exact symmetric zero initial data remains symmetric
under deterministic evolution.  No scalar seed, source-owned detector,
physical state, or $\mathsf{DYN}$ record is generated by the
present mathematical identities alone.
\end{remark}


\begin{proposition}[Exact native quartic under an explicitly licensed rooted pullback]
\label{p3f:prop:v243-rooted-native-quartic-test}
On a field owner where the \emph{actual} native scalar quartic action
is shown, by a separate source-incidence certificate, to use the
rooted finite comparison of
Theorem~\ref{p1b:thm:v243-rooted-register-functional}, suppose its
scalar profile is $\Phi(q)=v+f_a(q)$ with $v\in\mathbb R$,
$a\in G_\perp^{\mathbb R}$, and that both the quadratic and quartic
terms use the same unrenormalized rooted functional.  Write
$r^2=\sum_ja_j^2$, $c=a_0a_1a_2$, and retain the inherited
coefficient $\alpha_\Phi=(16\pi^2)^{-1}$.  Then the \emph{mathematical}
finite-action pullback of the native potential is exactly
\begin{equation}\label{p3f:eq:v243-rooted-quartic-pullback}
 \boxed{\begin{aligned}
 \alpha_\Phi^{-1}V_{\rm root}(v,a)={}&
 \frac{24}{25}(v^4-2\Lambda_g^2v^2)\\
 &+\left(\frac{144}{25}v^2-\frac{48}{25}\Lambda_g^2\right)r^2\\
 &-\frac{24}{25}vc+\frac{131}{50}r^4.
 \end{aligned}}
\end{equation}
The cubic is fixed by the same quartic parent and the same deleted
register, rather than by an independently introduced cubic coupling.
If the corresponding kinetic quadratic metric is the geometric
$24I/25$, the canonically rescaled cubic is
$-\alpha_\Phi v\sqrt{25/24}\,
 a_{{\rm can},0}a_{{\rm can},1}a_{{\rm can},2}$.
This proposition neither proves that the physical $\Phi$ has the
specified source/target type nor that the actual parent contains this
marked pullback; it does not determine $v$ from the complete
stationarity equations, furnish a chiral Yukawa vertex, or establish
an interacting pole.
\end{proposition}
\begin{proof}
Expand $(v+f_a)^4-2\Lambda_g^2(v+f_a)^2$ before taking
$\tau_{\rm root}$, and use the first four moments from
Theorem~\ref{p1b:thm:v243-four-atom-family-law}:
\[
\begin{aligned}
 \tau_{\rm root}(1)&=24/25,& \tau_{\rm root}(f_a)&=0,\\
 \tau_{\rm root}(f_a^2)&=24r^2/25,& \tau_{\rm root}(f_a^3)&=-6c/25,\\
 \tau_{\rm root}(f_a^4)&=131r^4/50.
\end{aligned}
\]
The binomial coefficients $6,4$ generate respectively the
quadratic $v^2$ and cubic $v$ terms in the displayed equation.
Substitute $a=\sqrt{25/24}\,a_{\rm can}$ for the final coefficient.
\end{proof}

\begin{theorem}[Global vacuum set and raw Hessian of the rooted finite quartic]
\label{p3f:thm:v245-rooted-quartic-global-minima}
On precisely the real rooted-comparison domain and positive
$\alpha_\Phi$ of Proposition~\ref{p3f:prop:v243-rooted-native-quartic-test},
with $\tau_{\rm root}=\tau_{2I}-\frac1{25}\tau_{2T}$, the potential
$V_{\rm root}(v,a)$ has the exact lower bound and all global minimizers
\[
 \boxed{V_{\rm root}(v,a)\ge
 -\frac{24}{25}\alpha_\Phi\Lambda_g^4,\qquad
 \operatorname*{argmin}_{v\in\mathbb R,\,a\in G_\perp^{\mathbb R}}
 V_{\rm root}=\{(\Lambda_g,0),(-\Lambda_g,0)\}.}
\]
At each minimum its \emph{raw coefficient-space} Hessian is
\[
 \boxed{D^2V_{\rm root}(\pm\Lambda_g,0)
  =\frac{192}{25}\alpha_\Phi\Lambda_g^2 I_3}
\]
on $\mathbb Rv\oplus G_\perp^{\mathbb R}$ in the displayed coefficient
metric. This does not claim the stationary vacuum, physical scalar
mass matrix or family hierarchy of the \emph{complete} action.
\end{theorem}
\begin{proof}
The normalized averages over $2I$ and $2T\subset2I$ assign
$\tau_{\rm root}$ the strictly positive point weights
$w(q)=1/150$ for $q\in2T$ and $w(q)=1/120$ for
$q\in2I\setminus2T$, with total mass $24/25$. Thus
\[
 V_{\rm root}=\alpha_\Phi\tau_{\rm root}\bigl(
 ((v+f_a)^2-\Lambda_g^2)^2\bigr)
 -\frac{24}{25}\alpha_\Phi\Lambda_g^4.
\]
Equality forces $(v+f_a(q))^2=\Lambda_g^2$ at every icosian point.
The squared expression is a nonnegative degree-at-most-eight
polynomial on the round $S^3$ because $f_a$ has degree two.
The $2I$ eleven-design identity makes its round integral zero;
continuity and positivity force it to vanish on the entire sphere.
Connectedness makes the continuous map $v+f_a$ constant, with value
$+\Lambda_g$ or $-\Lambda_g$. The zero mean and injectivity of the
quadratic family embedding imply $a=0$ and $v=\pm\Lambda_g$.
Differentiating Eq.~\eqref{p3f:eq:v243-rooted-quartic-pullback}
there gives the displayed Hessian on the one scalar plus two
transverse-family real coordinates.
\end{proof}
\begin{remark}[Vacuum scope and kinetic normalization]
\label{p3f:rem:v245-rooted-vacuum-action-scope}
The isolated finite scalar potential alone leaves the two family-plane
quadratic curvatures degenerate; its nonzero cubic $-24vc/25$ does not
select a nonzero family-plane vacuum.  The physical scalar masses require
the \emph{actual} kinetic Gram, source incidence, metric curvature coupling
and full stationary action. Neither the physical Yukawa hierarchy nor a
physical $\mathbb Z_2$ breaking claim follows merely from this potential.
\end{remark}

\begin{remark}[Precise residual same-action tests; no asserted one-condition closure]
\label{p3f:rem:v243-one-root-action-gap}
The geometric comparisons, the quartic pullback above and the
neutral map of Theorem~\ref{p1b:thm:v243-neutral-five-isometry}
share one finite geometric origin but do not prove one universal
physical action.  A physical closure requires the \emph{actual}
preterminal source incidence and mixed action derivatives, including
the full kinetic, gauge, metric, normal-source and boundary jets,
plus the separate physical state/pole conditions in their native
sorts.  The conormal cap amplitude of
Theorem~\ref{p3e:thm:v229-conormal-cap}, the whole-action Newton
normalization of Proposition~\ref{p3f:prop:v224-normalization-transport}
and the actual K56 mass-incidence word are not outputs of an
unlicensed finite register trace.  No conjunction of these
requirements is presented as a new single geometric axiom.
\end{remark}

\begin{proposition}[Exact singlet--plane Schur response of the rooted phase Gram]
\label{p3f:prop:v244-rooted-phase-schur-order}
On the fixed finite source comparison of
Proposition~\ref{p1:prop:v242-complete-transverse-return} and
Corollary~\ref{p1:cor:v244-quarter-turn-full-gram}, restrict the
quarter-turned Gram to its three labelled family columns:
\[
 D_f=\operatorname{diag}(0,1,-1),\quad
 G_f(t)=3(1+t^2)I_3+4tD_f,\quad
 c=\tfrac1{\sqrt3}(1,1,1)^T,\quad Q_\perp=I_3-cc^\dagger.
\]
Suppose, \emph{for this finite Hessian diagnostic}, that \(G_f(t)\)
is the positive common-action quadratic kernel and the \(Q_\perp\)
plane is an actually shared stationary interior.  Projecting onto
\(c\) before elimination gives \(K_{\rm proj}=3(1+t^2)\), whereas
eliminating the plane first gives the exact scalar Schur coefficient
\[
 \boxed{K_c(t)=\frac{9(1+t^2)(9t^4+2t^2+9)}
 {27t^4+38t^2+27},\qquad
 \Delta_c(t)=K_{\rm proj}-K_c
 =\frac{96t^2(1+t^2)}{27t^4+38t^2+27}\ge0.}
\]
For the \emph{comparison parameter} \(t\),
\[
 \boxed{K_{\rm proj}''(0)=6,\quad K_c''(0)=-\frac{10}{9},
 \quad\|R_{\rm load}(0)\|^2=\frac{32}{3},
 \quad C(0)=3I_2.}
\]
These are exact finite-source identities, not an identification of
\(t\) with physical curvature or of \(G_f\) with the physical
joint-Dirac Hessian.  The required occurrence, kinetic/metric form,
source frame, action and common-interior incidence must be certified
separately before the term can enter a physical \(f_{RR}\), pole or
source response.
\end{proposition}
\begin{proof}
The matrix is positive for all real \(t\).  Its Schur reduction is
\(K_c=(c^\dagger G_f^{-1}c)^{-1}\); the rational simplification
and second-jet load evaluation are given in
Proposition~\ref{p3g:prop:v244-rooted-phase-schur-proof}.  They
also follow from Theorem~\ref{p3f:thm:v199-second-schur-jet}.
\end{proof}

\begin{remark}[Rooted source averaging, nonlinear vertices and quotient order]
\label{p3f:rem:v244-rooted-action-order-firewall}
The \(C_3\) orbit average kills \(D_f\), hence averaging this
\emph{linear kernel} before elimination gives \(3(1+t^2)\).
It does not equal the scalar response obtained by eliminating the
oriented common interior first.  It does \emph{not} erase the distinct
\(C_3\)-invariant rooted cubic
\(\tau_{\rm root}(f_a^3)=-6a_0a_1a_2/25\) of
Theorem~\ref{p1b:thm:v243-four-atom-family-law} or change the
quartic parent of Proposition~\ref{p3f:prop:v243-rooted-native-quartic-test}.
For that marked functional,
\(\tau_{\rm root}(f_a^4)-(24/25)\tau_{2I}(f_a^4)
=7\|a\|^4/250\); a single \(24/25\) multiplier is not a
valid replacement for its nonlinear vertex word.
The source/orbit average, nonlinear function, source projection,
and stationary shared-interior quotient have different owner types.
Their ordering is physical only after the occurring action prescribes
it.  Retain the original unnormalized \(25\)-register trace unless
the actual kinetic/action normalization licenses its rescaling;
never add the same Schur response again as a separate cap term.
In particular, an operator-valued phase load is not by itself the
rank-seven gravitational normal defect or the constant \(h^4\)
boundary source of a proposed complete cap.
\end{remark}


\begin{remark}[Product and truncation scope]
The pointwise products are taken in the full $120$-point function algebra
before the normalized spatial trace.  They are not products of operators
first projected back into the $14$-dimensional harmonic image.  Interacting
solutions of the restricted system need not solve the unrestricted field
equations, because products can source higher harmonics.
\end{remark}

\begin{corollary}[Same-action chiral Noether charge and free Casimir defect]
\label{p3f:cor:v180-scalar-clock-noether-mass-defect}
The continuous round chiral action preserves the finite polynomial action even
when a generic chiral element does not permute the $120$ vertices.  For an
infinitesimal real antisymmetric mode generator $T_i$,
\[
 \mathcal J_i=p^TT_i z,\qquad p=\dot z,
\]
is conserved on the stated stationary same-action domain.  Its quadratic
positive-frequency branch is
\[
 H_{\rm sc,free}=\sqrt{E_\star^2\mathsf C+m_{\rm loc}^2I},
 \qquad
 \boxed{H_{\rm sc,free}^2-\sum_iG_i^2=m_{\rm loc}^2I}.
\]
The cubic and quartic terms remain in the interacting Hamiltonian.
\end{corollary}

\begin{theorem}[Cell incidence is exactly the constant-plus-quadratic band]
\label{p3f:thm:v180-incidence-harmonic-resolution}
Let $U_B$ be the kinetically normalized rank-$17$ $2I$ incidence map already
constructed in Paper~V-B, let $j_H$ insert a marked Hurwitz $2T$ cell, and set
$R_H=j_H^\dagger U_B$.  On that cell let $P_0,P_1,P_2$ denote the sampled
round harmonic projectors of ranks $1,4,9$.  Then
\[
 \boxed{R_HR_H^\dagger=P_0+\frac4{15}P_2,\qquad
 \operatorname{im}R_H=\operatorname{im}(P_0+P_2).}
\]
Thus the old rank-ten restriction is precisely $\mathcal H_0\oplus\mathcal
H_2$ and contains no degree-one harmonic.  In orthonormal scalar mode
coordinates its kinetic Gram is
$\operatorname{diag}(1,(4/15)I_9)$, so canonical normalization multiplies
the nine quadratic coordinates by $\sqrt{15}/2$.  The raw $4/15$ therefore
multiplies both kinetic and potential pullbacks before normalization and is not
an additional mass factor.
\end{theorem}

\begin{corollary}[Hurwitz return and canonical polar invariance]
\label{p3f:cor:v180-hurwitz-return-polar-invariance}
Compress the existing unital Hurwitz projector $P_{17}$ to the complete
$14$-dimensional sampled band $P=P_0+P_1+P_2$.  The compressed response has
spectrum
\[
 1^{(8)}\oplus(1/2)^{(3)}\oplus(1/6)^{(3)}.
\]
On the even ten-dimensional band $P_0+P_2$, the combined incidence--Hurwitz
return has kinetic Gram spectrum
\[
 1^{(1)}\oplus(4/15)^{(3)}\oplus(1/15)^{(6)},
\]
and its canonical polar map is exactly the same as for the incidence source
alone.  Hence the two specified source words give the same canonically
parametrized isolated scalar action on their common ten-dimensional image,
while their parent label/source roles remain distinct.
\end{corollary}

\begin{corollary}[Complete-frontier mode count]
\label{p3f:cor:complete-frontier-mode-density}
For the existing complete sharp position frontier $m_x(a)$,
\[
 D_{m_x}=\frac{(m_x+1)(m_x+2)(2m_x+3)}6,
 \qquad
 \frac{D_{m_x(a)}}{2\pi^2a^3}\longrightarrow
 \frac{\Lambda_F^3}{6\pi^2}.
\]
The currently exposed frontier shell has exact transverse rank
\[
 \boxed{N_{\partial x}(m_x)=(m_x+1)^2,}
\]
and the successor event $m_x\to m_x+1$ adds the next shell, of rank
$(m_x+2)^2$.  Thus the cumulative scalar test capacity is cubic while the
instantaneous shell capacity is quadratic.  These are mode/test capacities,
not occupation numbers, entropies, or counts of ontic spatial cells, and they
are unrelated to the fixed comparison value $m_{\rm cmp}=2$ except where
that finite test band is explicitly invoked.
\end{corollary}


\subsection{Quadratic sibling same-action response}
\label{p3f:sec:v189-quadratic-sibling-response}
The complete quadratic $M_8$ parent descends through Paper~III-A into three
typed local variations: trace/homothety, symmetric-traceless metric shear, and
antisymmetric frame/connection directions.  Once these belong to one
occurring action owner $S_\star$, their first and second variations are
same-owner descendants rather than independently supplied cross-sector maps.

\begin{theorem}[Ward--Schur inheritance of the TT K16 split]
\label{p3f:thm:v189-ward-schur-k16-tt}
Assume the physical boundary trace/splitting preserves a TT subspace and the
interior Hessian block is invertible on the stated reduced domain.  Then the
same-action Schur complement
\[
 K_\partial=H_{\partial\partial}
 -H_{\partial I}H_{II}^{-1}H_{I\partial}
\]
preserves every projectable Ward null relation.  At the $K=16$ junction the
scalar distributional block found in Paper~III-E remains absent from the
reduced TT Hessian, whereas the finite one-sided $F$ normalization is inherited
by the TT boundary block.
\end{theorem}

\begin{remark}[Block-preservation domain]
The theorem does not assert scalar/vector/tensor decoupling on a generic
anisotropic or mode-mixing background.  When the background or boundary
condition fails to preserve the TT subspace, the full coupled Hessian must be
reduced before any sector statement is made.
\end{remark}


\subsection{Finite metric-dependent realization of the complete action}
\label{p3f:sec:v1148-finite-metric-construction}

The constructor here consumes the complete local action, including the
clocked fermion member of Theorem~\ref{p3f:thm:v177-local-clocked-fermion-jet}
and Corollary~\ref{p3f:cor:v1122-root-generated-fermion-member-handoff},
and every gauge, scalar, curvature and constraint sector with consumed jets.
The native potential and curvature coefficient are
\[
 \alpha_{\rm loc}=\frac1{16\pi^2},\qquad
 \mathscr F(\Phi)=\kappa_J-\frac{\alpha_{\rm loc}}6
       \operatorname{Tr}_F\Phi^2,\qquad
 V(\Phi)=\alpha_{\rm loc}\operatorname{Tr}_F
                      (\Phi^4-2\Lambda_g^2\Phi^2).
\]
These are the inherited action coefficients, not new interaction weights.
The claim domain is an already licensed finite action/quantization word,
with its actual clock, field and dual-field embeddings, ordering, constraint
and occupation projectors, and vacuum allocation. Kinetic forms are positive
on an open fixed-rank chart; all consumed first, mixed second and functional
source jets exist. A finite CAR source or an actually finite joint bosonic
occupation space is required. Finite bosonic mode rank alone is insufficient.
No source member or physical cutoff is selected by this construction.

\begin{theorem}[Complete kinetic normalization and interaction-slot transport]
\label{p3f:thm:v1148-finite-metric-constructor}
Let $J(x,t)$ be the actual finite field embedding. Write the parent
symmetrized first-order fermion action with $G_p\succ0$,
$B_p^\dagger=-B_p$ and $K_p^\dagger=K_p$. Its pullback has
\begin{align*}
 G&=J^\dagger G_pJ,&K&=J^\dagger K_pJ,\\
 B&=J^\dagger B_pJ+\tfrac12
        (J^\dagger G_p\dot J-\dot J^\dagger G_pJ),\\
 L_F&=\tfrac{i\hbar}2(c^\dagger G\dot c-\dot c^\dagger Gc)
        +i\hbar c^\dagger Bc-c^\dagger Kc-V_{\rm inherited}.
\end{align*}
Put $S=G^{1/2}$, $W=S^{-1}$ and $Z=JW$. In the canonical coefficient
$q=Sc$, the energy and connection are
\[
 e=WKW,\qquad
 \gamma=WBW+\tfrac12(W\dot S-\dot SW),\qquad
 \widehat H=d\Gamma(e)-i\hbar d\Gamma(\gamma)
                 +\widehat V_{\rm inherited,can}+\varepsilon_0I.
\]
Every creation slot of an inherited vertex receives $Z^*$ and every
annihilation slot receives $Z$; the other quantum sectors use their own
inherited pullbacks. All consumed coefficient, density, field, occupation,
constraint and reduction jets are differentiated. The central energy
$\varepsilon_0$ remains part of the complete action family.

The noncommuting root and inverse jets are uniquely determined by
\begin{align*}
 SS_a+S_aS&=G_a,\\
 SS_{ab}+S_{ab}S&=G_{ab}-S_aS_b-S_bS_a,\\
 W_a&=-WS_aW,\\
 W_{ab}&=WS_bWS_aW+WS_aWS_bW-WS_{ab}W,\\
 Z_a&=J_aW+JW_a,\\
 Z_{ab}&=J_{ab}W+J_aW_b+J_bW_a+JW_{ab}.
\end{align*}
Consequently the complete finite metric/source jets are computable
descendants of the licensed parent without a Wick closure.
\end{theorem}
\begin{proof}
Substitute $Jc$ and its time derivative into the parent action, retaining
the two embedding-derivative terms in its symmetrized kinetic part.
Substitution of $c=Wq$ gives the displayed $e$ and $\gamma$;
$\gamma^\dagger=-\gamma$ makes the lifted generator self-adjoint.
The positive-root Sylvester operator has eigenvalues $s_i+s_j>0$;
twice differentiating $S^2=G$ therefore gives its unique jets.
Differentiating $SW=I$ and $Z=JW$ proves the remaining formulas.
Substitution into each inherited monomial gives the slot rule; the product
rule differentiates every slot and its coefficient. This is finite matrix
quantization of the pulled-back word, not a claim that its trajectories
equal unrestricted parent trajectories.
\end{proof}

For example, where a quartic fermion vertex is actually inherited,
\[
 \widehat V_4=\sum_{i<j,k<l}\mathsf V_{ij;kl}
                   a_i^\dagger a_j^\dagger a_l a_k,
 \qquad \mathsf V_F=(\Lambda^2Z)^\dagger\mathsf V_p(\Lambda^2Z).
\]
Its mixed derivative is the sum of all nine placements of the two
derivatives among the three factors, including their mixed jets.
This tensor formula introduces no additional native four-fermion term.
The native scalar quartic and quantum-scalar Yukawa interactions already
belong to the joint action.

For the native restriction-before-Legendre scalar word, its kinetic part is
\[
 L_\Phi=\tfrac12\dot z^TM\dot z-\dot z^TCz
       +\tfrac12z^TDz-\tfrac12z^TKz-U(z),\qquad M\succ0.
\]
Direct Legendre transformation gives
\[
 p=M\dot z-Cz,\qquad
 H_\Phi=\tfrac12(p+Cz)^TM^{-1}(p+Cz)
                         +\tfrac12z^T(K-D)z+U(z).
\]
All actual affine, temporal-gauge, curvature and connection terms are
retained in this step. In particular $D$ is formed from the parent product
of both shift derivatives before projection; $D=C^TM^{-1}C$ requires a
closed derivative image and is not imposed. Ordering and occupation
compression are taken in the licensed word. Finite CCR boundary defects
and source-dependent occupation projectors remain in its jets.
Lemma~\ref{p4:lem:v2174-weighted-ordering-boundary} evaluates the exact
multimode top-edge difference between products of compressed quadratures
and compression of the parent product. The licensed word retains its
actual ordering and central allocation; neither comparison is chosen
for a missing member, and its top projector is not assumed to kill the
full Gibbs support.


\subsubsection{Native metric-density vertices and exact geometric evaluation}
\label{p3f:sec:v1148-metric-density-coefficients}
For a real unit harmonic $Q$ of degree $n>2$, put
$\lambda=n(n+2)$, $H_{ij}=D_iD_jQ$ and $u=-2\psi QI+2EH$.
Retain the full scalar metric chart before restriction:
\[
 g_{00}=-a^2(1+2AQ),\quad g_{0i}=a^2BD_iQ,\quad
 h_{ij}=a^2(\gamma_{ij}+u_{ij}).
\]
Then
\[
\begin{aligned}
 \ell&=\sqrt{1+2AQ+B^2DQ^T(I+u)^{-1}DQ},\\
 J_h&=\sqrt{\det(I+u)},& N&=a\ell,\\
 \sqrt{-g}&=a^4\sqrt\gamma\,\ell J_h.
\end{aligned}
\]
With $f_a=a^4/c$, the fixed reference scalar-coefficient potential is
\[
 \widehat H_V=f_a\int_{S^3}\ell J_h\widehat{V(\Phi)}\,d\sigma_\gamma.
\]
It is generally noncentral. In port order $(A,B,\psi,E)$, the density
first jets are $(Q,0,-3Q,-\lambda Q)$ and the second jets are
\[
 (c_{ab})=\begin{pmatrix}
 -Q^2&0&-3Q^2&-\lambda Q^2\\
 0&|DQ|^2&0&0\\
 -3Q^2&0&3Q^2&\lambda Q^2\\
 -\lambda Q^2&0&\lambda Q^2&\lambda^2Q^2-2|D^2Q|^2
 \end{pmatrix}.
\]
Thus its density Hessian is $f_a\int c_{ab}\widehat V\,d\sigma_\gamma$.
Expansion of the square root and determinant proves the formula, including
the shift-squared lapse contact. Moving scalar fields add the actual slot
jets. This is not the complete Hessian with kinetic and curvature terms
omitted. A homogeneous potential density has $\bar p_V=-\bar e_V$;
its density contacts cancel in the normalized Newtonian stress conversion
of Paper~VI-A, while its operator still affects the state and memory.

Clock-normal fermion half-density normalization cancels $J_h$ between the
two Yukawa field legs. With the actual normalized Yukawa block $R_Y$,
\[
 \widehat H_Y=\int\chi^\dagger\ell R_Y(a,\widehat\Phi)\chi\,d\sigma_\gamma.
\]
At fixed scalar coefficients its density jets are
$D_A^Y=\int Q\chi^\dagger R_Y\chi$, $D_\psi^Y=0$,
$Q_{AA}^Y=-\int Q^2\chi^\dagger R_Y\chi$,
$Q_{A\psi}^Y=Q_{\psi\psi}^Y=0$ and
$Q_{BB}^Y=\int|DQ|^2\chi^\dagger R_Y\chi$.
The other field, frame and connection jets are retained separately.

\begin{corollary}[Exact native metric-weighted polynomial coefficients]
\label{p3f:cor:v1148-metric-haar-coefficients}
On the local-action integral word, Paper~I-C,
Theorem~\ref{p1b:thm:design}, supplies $V_\gamma=2\pi^2$ and
\[
 \int\prod_{i=0}^3y_i^{2k_i},\frac{d\sigma_\gamma}{V_\gamma}
 =\frac{\prod_i(2k_i-1)!!}{2^K(K+1)!},\qquad K=\sum_i k_i;
\]
an odd exponent gives zero. Each polynomial vertex with its retained
metric-jet weights is consequently an exact finite moment sum.
For a homogeneous harmonic polynomial $F_n$,
\begin{align*}
 |DF_n|^2&=|\nabla F_n|^2-n^2F_n^2,\\
 |D^2F_n|^2&=|\nabla^2F_n|^2-2(n-1)^2|\nabla F_n|^2
                              +n^2(n^2+2)F_n^2.
\end{align*}
These weights are polynomial and use the same moment evaluator.
\end{corollary}
\begin{proof}
Expand the finite vertex polynomial into monomials and apply the inherited
Haar formula. With $\Pi=I-yy^T$, the sphere Hessian is
$\Pi\nabla^2F_n\Pi-nF_n\Pi$. Euler homogeneity and
$\operatorname{Tr}\nabla^2F_n=0$ give the two norm identities.
Their products with the vertex polynomial are again finite moment sums.
\end{proof}
The degree-11 $2I$ design evaluates the native round quartic on its stated
degree domain; it does not establish all metric-weighted second jets.
The evaluation degree bound is not a physical cutoff. A literal vertex-sum
owner must be differentiated as that sum unless its full metric family is
proved equivalent to the integral word. Paper~III-G, Appendix~\ref{p3f:app:v1148-metric-quartic-entry}
computes an exact quartic tensor entry and the $1/32$ discrepancy from
extending fixed $2I$ cubature beyond its certified degree.

\subsubsection{Free scalar-metric Dirac specialization}
\label{p3f:sec:v1148-free-metric-dirac}
On the constant homogeneous generated mass branch, normalize the local
spinor by its clock-normal form and spatial half-density, using the
future-normal positive square-root coframe. Expand in the actual reference
component band $P_m$. This gives kinetic metric $I$. The full action retains
lapse, shift, spatial Dirac and spin-connection terms before $B=E=0$.
For one scalar harmonic, the spatial matrices $I$ and $D^2Q$ commute with
their time derivatives; their rotational temporal coframe connection
vanishes and half-density normalization cancels the symmetric time term.
This supplies instantaneous scalar jets on this specific frame branch.
General moving carriers and multiple noncommuting tensor directions use
the functional source theorem of Paper~IV.

On the Newtonian leaf put $T=\hbar D_\gamma^{\rm spatial}$,
$R=ac\beta_D\mathsf m$, $\ell=\sqrt{1+2AQ}$ and
$v=(1-2\psi Q)^{-1/4}$. The energy-complete one-body member is
\[
 h_\eta=P_m[\tfrac12\{\ell,vTv\}+\ell R]P_m.
\]
The conformal spinor law
$D_{e^{2\omega}\gamma}=e^{-2\omega}D_\gamma e^\omega$
becomes $e^{-\omega/2}D_\gamma e^{-\omega/2}$ after half-density
normalization; lapse symmetrization supplies its Hermitian gradient term.
Differentiation therefore gives
\begin{align*}
 D_A&=P_m[\tfrac12\{Q,T\}+QR]P_m,&
 D_\psi&=P_m\tfrac12\{Q,T\}P_m,\\
 Q_{AA}&=-P_m[\tfrac12\{Q^2,T\}+Q^2R]P_m,\\
 Q_{A\psi}&=P_m\tfrac14\{Q,\{Q,T\}\}P_m,\\
 Q_{\psi\psi}&=P_m[\tfrac54\{Q^2,T\}+\tfrac12QTQ]P_m.
\end{align*}
These are parent multiplication words with outer compression. Lift by
$d\Gamma$ and add the actual central derivatives. Products of individually
compressed factors require the consumed guard/invariance certificate.
Paper~VI-A proves the guarded first-response Ward and contact specialization;
no generic interacting or finite-amplitude Ward closure is inferred here.


\subsection{Complete finite quantum source realization of the causal Schur block}
\label{p3f:sec:v1145-finite-quantum-source-schur}

The general causal reduction consumes both response and source columns.
For an admitted system
\[
 \begin{pmatrix}A&B\\D&C\end{pmatrix}\binom{w}{z}
 =\binom{j_w}{j_z},
\]
with an actual retarded source inverse, elimination gives
$K_R=A-BC_R^{-1}D$ and $K_Rw=j_w-BC_R^{-1}j_z$.
This is the same-action result of
Theorem~\ref{p3f:thm:v1144-causal-dynamic-schur}, restated here as the
load-bearing antecedent.  Its finite quantum realization uses Paper~IV's
source calculus; that calculus does not depend on the application below.

\begin{corollary}[Exact finite interacting quantum source block]
\label{p3f:cor:v1145-finite-quantum-source-schur}
Use an already constituted active-stress scalar port $v=\delta\phi=-w/2$
and the complete finite source owner of Paper~IV,
Theorem~\ref{p4:thm:v4104-finite-interacting-source-response}.
Let $\mathcal R_\star$ be the inherited Hilbert/Riesz conversion of the
force covector into the normalized scalar source, including its actual Gram
and gravitational pairing.  Form covectors and full action jets before this
conversion.  Put
\[
 (\mathcal C_Qw)(t)=\frac{i}{\hbar}[D_bw^b,\bar\rho],\qquad
 \mathcal B_QX=\tfrac12\mathcal R_\star
                    \{\operatorname{Tr}(D_bX)\}_b,
\]
and retain the complete local scalar block $A$, including
\[
 \mathcal V_Q^{\rm loc}=\tfrac12\mathcal R_\star\bar Q
\]
and every additional consumed measure, connection and mixed-metric contact.
The exact coupled equations and their causal elimination are
\[
 \mathscr L\sigma=\tfrac12\mathcal C_Qw+q,
 \qquad Aw-2\mathcal B_Q\sigma=j_{\rm grav},
\]
\[
 \boxed{K_Q^R=A-\mathcal B_Q\mathscr L_R^{-1}\mathcal C_Q,
 \qquad K_Q^Rw=j_{\rm grav}+2\mathcal B_Q\sigma_{\rm ind}.}
\]
Equivalently, the source contribution is
\[
 \mathcal V_Q^R=\mathcal V_Q^{\rm loc}
                  +\tfrac12\mathcal R_\star\Pi^R.
\]
No positivity of the quantum contact is asserted.  Interactions and initial
source data remain part of the same owner.  Source-dependent connection or
metric-derivative terms use the complete functional variation and associated
contact blocks, not an instantaneous derivative with those terms omitted.
\end{corollary}
\begin{proof}
The inherited $v=-w/2$ converts the finite density tangent into the first
coupled equation.  Insert its unitary retarded solution in the second.
The factors $1/2$ and $2$ cancel in the memory Schur term; the independent
source contribution retains the factor $2$.  The commutator formula in
Paper~IV identifies this term with $\mathcal R_\star\Pi^R/2$.
\end{proof}

The finite interacting part of general source dynamics is thus closed on
constituted finite ports.  A full moving metric application additionally
requires the complete metric-dependent source family and its consumed stress
Ward identities.  Finite unitarity alone does not imply those identities or
license a kinetic environmental projector on every quantum state.
For an already constituted complete finite joint source owner and licensed
finite reduction, Paper~IV,
Corollary~\ref{p4:cor:finite-reduced-source-response}, supplies the exact
reduced contact/memory response.  Source dependence of the reduction retains
$D\mathcal R_\phi\bar\rho_J$ and the consumed second/mixed reduction jets;
these terms are contacts of the existing owner, not independent coefficients.
The resulting outputs enter the same causal Schur block with the actual
source columns and inherited Hilbert/Riesz pairing.  No closed autonomous
Markov equation for the reduced density is inferred.  A physical subsystem
or bath, correlation/assignment domain, full metric/Ward owner, collisional
limit, unbounded-field domain and continuum bounds remain antecedents where
consumed.  The finite reduced response computation is thus closed-dependent
on the complete joint owner and licensed reduction.

\subparagraph{Parent origin of the metric source.}
The finite response calculus consumes the existing complete action/source
packet; its metric family must be formed by the licensed parent
action/quantization word before lapse, shift, curvature or shear directions
are restricted.  Once that word and its consumed source/frame data exist,
Hilbert derivatives, contacts and the retarded finite response are descendants
as in Corollary~\ref{p3a:cor:complete-parent-descendant-reuse}, not new stress
owners at the RAR consumer.  A generated complete family may be consumed as
a family where the target permits it.  The actual full metric word,
projectability and quantum symmetry/state/measure Ward antecedents remain
explicit certificates; the one-port response does not determine discarded
metric jets.  Regular coupled invertibility and further source/critical
domains are later claim-domain checks.

\subparagraph{Constructed metric and functional-source handoff.}
Theorem~\ref{p3f:thm:v1148-finite-metric-constructor} supplies the complete
licensed finite metric family and slotwise interaction jets.
Paper~IV, Theorem~\ref{p4:thm:v4106-functional-hilbert-response}, computes
its functional insertion and connection/derivative contacts with the same
retarded inverse. The canonical Ward defect and its first variation in
Theorem~\ref{p4:thm:v4106-canonical-stress-ward} test the actual transport.
Construction and response close on their declared finite word; general
source covariance/projectability and consumed regularity remain explicit.
The guarded free specialization is proved in Paper~VI-A.


\subsubsection{Same-parent successive action reduction}
\label{p3f:sec:same-parent-reduction-reuse}
\begin{corollary}[Complete action reuse under regular reduction composition]
\label{p3f:cor:complete-reduction-reuse}
Two licensed successive reductions of one complete parent action have the
same descendant on the common uniquely generated regular stationary branch
of Corollary~\ref{p1b:cor:complete-successive-reductions}, provided both
routes consume the same retained output and normalization packet. All
existing source and metric jets agree. An actual Gaussian quantum route
also consumes its determinant and measure jets. A classical stationary
route contributes only its classically reduced action.
\end{corollary}
\begin{proof}
The common-parent reduction equality holds as an equality of functions,
including the inherited normalized Gaussian factors when that integration
is licensed. Differentiating the complete equality proves the jet claim.
The existing complete-parent reuse and II-A co-consumption rules retain
those generated outputs. They add no determinant to a classical word.
\end{proof}
Route equality requires the same branch and endpoints; it does not compare
different vacuum allocations or omit a metric-dependent central term.


\subsection{Tensor columns of the generated Dirac source}
\label{p3f:sec:generated-tt-dirac-source}
Consume the clocked fermion member of
Theorem~\ref{p3f:thm:v177-local-clocked-fermion-jet} and
Corollary~\ref{p3f:cor:v1122-root-generated-fermion-member-handoff},
canonically normalized by
Theorem~\ref{p3f:thm:v1148-finite-metric-constructor}. The varied metric
is reciprocal-position spacetime after its causal/coframe realization;
the compact momentum geometry and $m_p=1$ remain fixed.
The ordinary real mass coefficient is the actual generated member. If it
varies along the history, its mass owner and work terms are retained.
No occupied source is inserted into the native fresh neutral root.

Fix real transverse traceless harmonics on the unit round slice,
\[
 \nabla^iY^A_{ij}=0,\quad \gamma^{ij}Y^A_{ij}=0,\quad
 \int_{S^3}Y^A_{ij}Y^{B,ij}\,d\sigma=\delta_{AB},
 \qquad V_\gamma=2\pi^2.
\]
Use the metric-component chart
$g=-dt^2+a^2(\gamma+u)$, $u=\sum_Ah_AY^A$, and the canonical half-density
frame. All parent metric and embedding jets are formed before the finite
restriction $P_F$; all consumed transition images and contact words are
retained. Equality to a larger smooth parent consumes its own actual
comparison guard, rather than changing the fundamental momentum cutoff.

\begin{theorem}[TT vertex, Killing decoupling, and complete contact]
\label{p3f:thm:tt-dirac-vertex-contact}
For $\bar H_D=T+m\beta_D$, $T=(\hbar/a)\mathcal D_\gamma$, the first
one-body TT vertex, in the convention $c_i^2=-I$, is
\[
 D_A^{(h)}=P_F\left[-\frac{\hbar}{2a}
                 Y^A_{ij}c_i\nabla_j\right]P_F.
\]
Its quantum insertion is the native CAR lift $d\Gamma(D_A^{(h)})$.
It annihilates the lowest spatial Killing-spinor band on both sides:
\[
 D_A^{(h)}P_{\mathcal K_{x,0}}
       =P_{\mathcal K_{x,0}}D_A^{(h)}=0.
\]
Consequently a free source covariance supported on that complete band has
zero induced TT expectation kernel at every $L$, including on an evolving
round background that preserves the band.

On a homogeneous, isotropic, parity-balanced unpolarized free-source
domain, retain the vanishing background vector, spin and axial bilinears
that enter connection contacts. Set
\[
 V=a^3V_\gamma,\qquad p_D=\frac{\operatorname{tr}(CT)}{3V},\qquad
 Q_{AB}^{(h),D}=\left.\partial_A\partial_B H_D\right|_0.
\]
The same-action contact expectation and Hilbert conversion are
\begin{align*}
 \operatorname{tr}(CQ_{AB}^{(h),D})&=\tfrac34a^3p_D\delta_{AB},\\
 \delta\Pi_A&=-\frac2{a^3}\delta F_A^{(h)}+\bar p\,h_A.
\end{align*}
An inherited field-independent local constant density $V_0$, with
$p_0=-V_0$, contributes $Q_{AB}^{(h),0}=\tfrac12a^3p_0\delta_{AB}$.
For this Dirac-plus-constant domain, $\bar p=p_D+p_0$ and therefore
\[
 \boxed{\delta\Pi_A=-\tfrac12p_Dh_A
                    -\frac2{a^3}\delta F_{{\rm ind},A}.}
\]
The induced force includes its actual transported initial tangent when
one is owned. Other consumed source sectors and metric-dependent central
functionals contribute their complete jets separately.
\end{theorem}
\begin{proof}
Differentiating the inverse orthonormal frame gives $-u/2$. The fixed-spinor
variation of its spatial connection, in the convention of
\cite[Theorem~5.1]{BaerGauduchonMoroianu2005}, completes this to
\[
 \delta\mathcal D=-\tfrac12u_{ij}c_i\nabla_j
                 +\tfrac14c(d\operatorname{tr}u-\operatorname{div}u).
\]
Both latter terms and the first volume half-density variation vanish for
TT $u$. Differentiation of the normalized self-adjoint Hamiltonian gives
the displayed self-adjoint insertion. For a Killing spinor,
$\nabla_j\psi=\kappa c_j\psi$, $\kappa=\pm\tfrac12$.
Symmetry and tracelessness give
$Y_{ij}c_ic_j=-\operatorname{tr}(Y)I=0$.
Self-adjointness proves the left identity as well. The free commutator
response then vanishes because the occupied covariance is supported on
this annihilated invariant band.

For the contact, expand the inverse frame and volume in the fixed chart:
\[
 (I+u)^{-1/2}=I-\tfrac12u+\tfrac38u^2+O(u^3),\qquad
 J=\sqrt{\det(I+u)}.
\]
Thus $\partial_{AB}(I+u)^{-1/2}=\tfrac38(Y^AY^B+Y^BY^A)$.
Isotropy fixes the kinetic contraction to $a^3p_D$ times the matrix
trace; harmonic normalization yields the factor $3/4$.
The real mass term is metric independent after half-density normalization.
The remaining connection terms pair with the excluded background bilinears;
their expectations vanish on the stated domain, not as operator identities.
Also $J_{AB}=-\tfrac12\operatorname{tr}(Y^AY^B)$, giving the constant
density contact. Finally the complete force obeys
\[
 F_A^{(h)}=-\frac{a^3}{2}\int J\,
     \operatorname{tr}\bigl(T^i{}_j(I+u)^{-1}Y^A\bigr)d\sigma.
\]
Differentiating its inverse metric at $T^i{}_j=\bar p\delta^i{}_j$
gives the $\bar p h_A$ conversion. Substitution gives $-p_Dh_A/2$;
the local constant pressure cancels. This proves the complete contact,
even when the linear vertex annihilates the occupied band.
\end{proof}
A nonlinear change of metric coordinates retains its second chain-rule
contact. A more general source uses the functional Hilbert insertions of
Paper~IV, Theorem~\ref{p4:thm:v4106-functional-hilbert-response}; the
evaluated unpolarized contact is not substituted for that complete operator.

\begin{corollary}[Action incidence of the level-one quadrupole]
\label{p3f:cor:level-one-quadrupole-incidence}
In a left-invariant unit-sphere frame, the component-orbit block is
$\mathcal D_j=\tfrac32I+2\sum_i\sigma_i\otimes J_i^{(j)}$.
For $j=\tfrac12$, its coupled triplet has eigenvalue $+5/2$ and spectator
dimension two. A homogeneous normalized $L=2$ direction
$Y_A=S_A/\sqrt{V_\gamma}$, with
$\operatorname{Tr}(S_AS_B)=\delta_{AB}$, has insertion
\[
 d_A=-\frac{\hbar}{a\sqrt{V_\gamma}}
               \sum_{ij}S_{A,ij}\sigma_i\otimes J_j^{(j)}.
\]
On the triplet this is $(\hbar/(a\sqrt{V_\gamma}))S_A$ in its Cartesian
spin-one frame. The ordinary energy-sign double supplies the factor $t_3$.
Thus $X\mapsto I_2\otimes X$ embeds the right $M_3$ quadrupole of
Corollary~\ref{p1:cor:level-one-shear-tt-operators} into
$\operatorname{End}(R_A)$, $R_A=(\tfrac12,1)$, with this action-fixed
coefficient. Parity gives the left counterpart. This is a nonzero source
incidence between distinct carriers.
\end{corollary}
\begin{proof}
The displayed Dirac block is the spin connection $3/2$ plus the component
derivative in the invariant frame. Differentiate its inverse frame using
the theorem. Symmetric spin coupling on the triplet identifies the insertion
with $-S_{ij}Q_{ij}$. In the Cartesian spin-one representation this equals
$S$, proving the coefficient and rank five. The spectator factor supplies
the stated equivariant embedding. Paper~III-G, Appendix~\ref{p3f:app:tt-angular-contacts}
evaluates the full angular Gram and an independent metric-spectrum contact.
\end{proof}
The incidence is available whenever that spatial source block is actually
retained; it selects no occupation. Paper~VI-A evaluates an occupied
first-higher-band response conditionally and separately proves the native
free-root tensor covariance.

\begin{corollary}[First higher-band icosahedral mirror transfer of the generated metric vertex]
\label{p3f:cor:kneser-spin-five-mirror-metric-vertex}
Continue with the action-owned homogeneous trace-free rank-two metric
vertex of Corollary~\ref{p3f:cor:level-one-quadrupole-incidence},
without changing its action coefficient. The upper coupled branch of the
$j=9/2$ orbital Dirac block has total angular momentum $\ell=5$ and
restricts, under an occurring marked $A_5$, as
\[
 R_5\big|_{A_5}\cong\mathbf3\oplus\mathbf3'\oplus\mathbf5.
\]
This is the first integer angular band containing both inequivalent
icosahedral triplets. Write $K_i$ for its spin-five generators,
$P_3,P_{3'}$ for the corresponding character projectors and
$\kappa=\hbar/(a\sqrt{V_\gamma})$. On the upper coupled band the
already generated metric insertion is
\[
 d^{(5)}(S)=-\frac\kappa5\sum_{ab}S_{ab}
 \frac{K_aK_b+K_bK_a}{2}.
\]
For a unit direction $x\in\mathbb R^3$, put
$Q_x=xx^{\mathsf T}-I_3/3$, $S_x=\sqrt{3/2}\,Q_x$ and
$C_x=P_{3'}d^{(5)}(S_x)P_3$. Then
\[
 \boxed{\operatorname{Tr}(C_x^\dagger C_x)=\frac{189}{25}\kappa^2}
\]
for every unit $x$. For the three symmetry-axis types of the marked
packet, the exact singular-value multisets are
\[
\begin{array}{c|c}
 \text{threefold}& (\sqrt{21}/5)|\kappa|\,(1,2,2)\\
 \text{twofold}& (3\sqrt{21}/10)|\kappa|\,(\varphi,1,\varphi^{-1})\\
 \text{fivefold}& (3\sqrt{21}/5)|\kappa|\,(1,0,0),
\end{array}
\qquad \varphi=\frac{1+\sqrt5}{2}.
\]
These are source-vertex matrix-block strengths, not physical energy
poles. The full axial spin-five metric insertion remains unitarily
conjugate under spatial rotations and has axis-independent spectrum;
the listed differences exist only after resolving the $\mathbf3\to
\mathbf3'$ block of an occurring icosahedral marking.
\end{corollary}
\begin{proof}
The stated restriction of the eleven-dimensional spin-five space
follows by evaluating its character on the five conjugacy classes of
$A_5$. The $\ell=j+1/2$ highest-weight state gives the coefficient
$-\kappa/(j+1/2)=-\kappa/5$ for the projected rank-two tensor; by
rotational equivariance this determines the complete rank-two vertex.
For $T_z=K_z^2-10I$ the finite-group character projectors satisfy the
exact identity
\[
 \operatorname{Tr}(P_3T_zP_{3'}T_z)=126.
\]
This identity is independently certified by exact arithmetic in
$\mathbb Q(\sqrt5,i)$ using the $120$ binary-icosahedral matrices in
$\operatorname{Sym}^{10}(\mathbb C^2)$. The quintet occurs exactly once
in $\operatorname{Hom}(\mathbf3,\mathbf3')$, so the projected
rank-two map is proportional to the already established Petersen
mirror map of Theorem~\ref{p1:thm:petersen-mirror-envelope}. Combining
the norm $126$ with the unit shear factor $\sqrt{3/2}$ gives
$(3/2)(126/25)\kappa^2=189\kappa^2/25$.
The symmetry-axis singular ratios follow from the exact Petersen
quintet operator and its fivefold, threefold and twofold stabilizers;
all three displayed multisets have the stated common squared norm.
No such compression is a reducing, physically occupied spectral band
without an independently established source and state selection.
\end{proof}


\subsubsection{Complete chiral-source handoff to the QCD response}
\label{p3f:sec:chiral-source-handoff}
The full kinetic pairing and every inherited interaction slot in
Theorem~\ref{p3f:thm:v1148-finite-metric-constructor} precede scalar
response evaluation. Paper~IV supplies the finite stationary bound
of Theorem~\ref{p4:thm:finite-energy-source-certificate};
Paper~VII-A specializes it in
Theorem~\ref{p10:thm:chiral-energy-source-certificate}.
Both consume the existing complete action and source. Nonlinear or
moving charts retain their contacts rather than being declared linear.

The intrinsic momentum Dirac, the Cartan-twisted defining-doublet
operator, and a momentum multiplication symbol are distinct from the
full QCD Dirac and singlet many-body Hamiltonian. The round
spatial representative in Section~\ref{p3f:sec:m8td-spatial-operators}
does not turn a reciprocal position readout into a new primitive
spatial carrier, or identify its radius with the momentum radius.
Use the metric-compatible transport of
Corollary~\ref{p2c:cor:v244-reciprocal-spin-dirac-transport}.
The protected geometric zero and complementary gap in
Theorem~\ref{p3f:thm:m8td-protected-zero} are retained on their
proven domain; they do not determine the action-weighted QCD mode density.
The complete Schur source and norm are explicitly consumed in
Theorem~\ref{p10:thm:complete-schur-scalar-source}, as a descendant
of the existing complete regular reduction, without a new action owner.


\subsection{Even stationary Schur reduction preserves the fermion quadratic jet}
Work in a finite Grassmann coordinate presentation of the retained fermion
fields. Let $\mathfrak m$ be the ideal generated by their odd coordinates.
Its nilpotence is algebraic. Fermion degree below means degree in
$\mathfrak m$, not a perturbation in the physical pair strength and not an
occupation assumption. Pairing and normal bilinears both have degree two.

Let $y$ collect the retained bosonic boundary variables and the actual
source/metric/history arguments. On the existing even action patch write
\begin{equation}
 S(u;y,\psi)=S_0(u;y)+S_2(u;y,\psi)+S_4(u;y,\psi)+\cdots,
 \label{p3f:v2171:eq:degreeexpansion}
\end{equation}
where $S_{2j}$ has fermion degree $2j$. The coefficients include every
action term and inherited source contact consumed by the claim. This expansion
does not discard any higher degree. Let $u_0(y)$ solve
$D_uS_0(u_0;y)=0$ on the actual selected regular stationary branch, and let
\[
 C(y)=D_u^2S_0(u_0;y)
\]
be invertible on the already declared quotient. $C$ is the actual interior
Hessian/inverse domain, not a fitted stiffness.

\begin{theorem}[Exact degree-two retention and terminating interior lift]
\label{p3f:v2171:thm:retention}
On this patch the stationary interior solution has an even expansion
\[
 u_*(y,\psi)=u_0(y)+u_2(y,\psi)+u_4(y,\psi)+\cdots
\]
with $u_*-u_0\in\mathfrak m^2$. Its first nonzero term is
\begin{equation}
 u_2=-C^{-1}j_2,\qquad j_2=D_uS_2(u_0;y,\psi).
 \label{p3f:v2171:eq:lift}
\end{equation}
The complete reduced action $S^\partial(y,\psi)=S(u_*;y,\psi)$ obeys
\begin{equation}
 \boxed{S^\partial_2(y,\psi)=S_2(u_0(y);y,\psi).}
 \label{p3f:v2171:eq:quadratic}
\end{equation}
The higher even terms are determined recursively by the same inverse $C$.
For finitely many odd coordinates this recursion terminates at the maximal
exterior degree.
\end{theorem}
\begin{proof}
Evenness excludes a term linear in $\psi$. At degree zero the stationary
equation is the body equation. At degree two it is $Cu_2+j_2=0$.
At every subsequent even degree the unknown $u_{2j}$ occurs linearly as
$Cu_{2j}$; the remainder depends only on previously determined degrees and
the complete action coefficients. Invertibility fixes the lift recursively.
No infinite series in odd variables remains on a finite coordinate patch.

In $S_0(u_0+\delta u)$ the linear term vanishes by body stationarity and
the quadratic displacement starts at fermion degree four. In
$S_2(u_0+\delta u)$ the first displacement term also starts at degree four.
Every other action term starts at degree four or higher. This proves
\eqref{p3f:v2171:eq:quadratic} without assuming the higher terms vanish.
\end{proof}

Equivalently, at zero odd source coordinates the fermion--even-interior
Hessian blocks vanish. The even-interior Schur correction to the fermion
quadratic Hessian is therefore zero there. This is evaluation of an action
coefficient. It is not a claim that the physical cyclic state has zero
fermion occupation, zero covariance, or zero anomalous expectation.
An additional odd-mode elimination can have a nonzero Schur correction;
the theorem does not remove it.

\begin{corollary}[Consumed source jets retain stationary-background terms]
\label{p3f:v2171:cor:sourcejets}
For an actual retained source direction $a$,
\begin{equation}
 D_aS^\partial_2
 =\left(D_aS_2\right)_{u_0}
  +\left(D_uS_2\right)_{u_0}D_au_0,
 \qquad
 D_au_0=-C^{-1}D_aD_uS_0.
 \label{p3f:v2171:eq:sourcejets}
\end{equation}
Higher consumed jets are obtained by differentiating this same stationary
identity, including the inverse, frame and embedding derivatives where used.
If a provenance mark only multiplies an already existing fermion constituent
and does not change $S_0$, its body displacement is zero; its quadratic
marked coefficient passes through the even reduction by evaluation at $u_0$.
\end{corollary}
\begin{proof}
Differentiate \eqref{p3f:v2171:eq:quadratic} and the body stationary equation.
The last clause uses $D_aS_0=0$, not a fixed-background assumption for arbitrary
physical gauge, metric or history sources. A provenance mark is bookkeeping
for an existing constructor constituent, not a newly introduced potential.
\end{proof}

\subsection{The induced quartic action is a required constituent}
\begin{theorem}[Complete first stationary contact]
\label{p3f:v2171:thm:quartic}
Using the actual even-interior Hessian pairing, the degree-four reduced
action is
\begin{equation}
 \boxed{S^\partial_4
 =S_4(u_0;y,\psi)-\tfrac12\langle j_2,C^{-1}j_2\rangle.}
 \label{p3f:v2171:eq:quartic}
\end{equation}
In a functional domain this is the corresponding full Green-operator
pairing, which need not be local. If $j_2=\sum_\alpha j_{2,\alpha}$,
the required contact includes every term
\begin{equation}
 -\tfrac12\sum_{\alpha,\beta}
 \langle j_{2,\alpha},C^{-1}j_{2,\beta}\rangle.
 \label{p3f:v2171:eq:crosscontacts}
\end{equation}
\end{theorem}
\begin{proof}
At degree four the body contributes $\tfrac12\langle u_2,Cu_2\rangle$,
the quadratic fermion action contributes $\langle j_2,u_2\rangle$,
and the direct quartic term contributes $S_4(u_0)$. Substitute
\eqref{p3f:v2171:eq:lift}. Even-degree Grassmann expressions commute, so the
ordinary symmetric Hessian contraction supplies the stated factor.
\end{proof}
This term comes from the already eliminated interior and must remain in the
particle-boundary response. Reintroducing an independently evolving scalaron
to account for the same contact would count the eliminated parent again.
Conversely, retaining only the quadratic particle action would discard a
nonzero generated constituent unless the actual contact is proved zero on
the consumed domain. The actual action-to-Hamiltonian map, normal ordering,
metric and any induced connection remain inherited; the action coefficient
is not identified with a selected numerical Fock Hamiltonian by fiat.

\subsection{Reuse of the already owned boundary pair}
V-A, \texttt{p5:thm:v3105-k16-bl2-pair-corealization}, supplies the physical
neutral-right pairing coordinate on the occurring boundary owner, and IV
fixes its generated normalization. It is not a new bulk scalaron-decay
coupling. The following identity applies precisely to a constituent already
defined on the retained boundary variables.

\begin{theorem}[A boundary-only constituent commutes with interior reduction]
\label{p3f:v2171:thm:boundarypair}
Let $B(y,\psi)$ be an already generated constituent on the retained boundary,
independent of the eliminated interior coordinate $u$. On the same regular
branch,
\begin{equation}
 \boxed{\operatorname{Red}_u(S_N+B)=\operatorname{Red}_u(S_N)+B.}
 \label{p3f:v2171:eq:pairreuse}
\end{equation}
This holds for the complete action, not only its quadratic part.
\end{theorem}
\begin{proof}
$D_uB=0$ makes the stationary interior equations and their actual selected
branch identical. Substitution therefore leaves $B$ unchanged and reduces
$S_N$ by the existing operation. Boundary/source derivatives still include
the actual dependence of $B$ on $y$, its field trace and its kinetic frame.
\end{proof}
This proves reuse of a boundary-only pair in the already reduced normal
owner. The displayed constant-interior bookkeeping identity does not assert
an additional physical bulk pair occurrence. It does not invent an interior
dependence of that pair, and it does
not derive a physical normal-current trace from pair faithfulness. If a
different claimed constituent really depends on $u$, it belongs in the
complete $S_2$ and the contact formula \eqref{p3f:v2171:eq:quartic}; its interior
dependence must actually be generated. No such additional dependence is
postulated here.

For a normal action with the existing neutral-interior charge symmetry,
the selected stationary branch is equivariant: rotating the retained
fermion variables maps the interior solution to the solution of the same
stationary equation, and local uniqueness makes the two agree. Thus its
complete reduced action retains that symmetry, including all the contacts.
This is a specialization of the existing charge-graded stationary descent
in III-F, not a new conservation assumption derived from a generic $H_0$.

\subsection{Baryon conservation survives the native stationary contacts}
\label{p3f:v2171:sec:baryonexport}
V-A, \texttt{p5:thm:split-phase-proton-stability}, proves that the declared
split-phase gauge and charged/neutral Dirac vertices preserve $B$.
The native BL2 boundary constituent acts on the neutral lepton support and
also has $\Delta B=0$. The following result propagates this representation
statement through the stationary operation already used in this continuation.

\begin{theorem}[Baryon-equivariant all-order stationary descent]
\label{p3f:v2171:thm:baryonstationary}
Let the complete consumed classical parent and its boundary BL2 constituent
be invariant under the baryon action. Let its actual regular stationary
solution be unique on an invariant patch, with the inherited metric,
quantization and constraint maps baryon-equivariant. Then the reduced action
and every induced contact are baryon invariant. In a compatible canonical
charge frame their full finite Hamiltonian satisfies
\begin{equation}
 [Q_B,H_{\rm split+BL2,red}(t)]=0.
 \label{p3f:v2171:eq:baryoncommutator}
\end{equation}
Consequently on this finite Hamiltonian domain, from the neutral inherited
birth state,
\begin{equation}
 q_B(t)=0
 \quad\hbox{for all times for which this owner remains complete},
 \label{p3f:v2171:eq:baryonzero}
\end{equation}
even if its actual lepton charge becomes nonzero.
\end{theorem}
\begin{proof}
Write the baryon transformation as $g_\theta$ and the eliminated stationary
solution as $u_*(y,\psi)$. Parent invariance carries a solution to a solution
at the transformed retained arguments. Uniqueness gives
$u_*(g_\theta(y,\psi))=g_\theta u_*(y,\psi)$.
Substitution therefore makes $S_{\rm red}$ invariant. Differentiation in
the baryon action proves that each retained Grassmann coefficient is
invariant, including the full $-j_2C^{-1}j_2/2$ crosscontact and all higher
terms. The canonical pullback, connection and inherited ordering are
equivariant by hypothesis, so their CAR lift commutes with $Q_B$.
The lepton-only pair commutes with $Q_B$ as well. Apply
Theorem~\ref{p4:v2171:thm:fullcharge}, with the covariantly constant charge, and
the inherited neutral initial value.
\end{proof}
This is the exact scope of the specified finite descendant. It is not a
claim that the framework forbids an anomalous physical baryon current.
Such a current consumes a regulated quantum response or a different owned
charge-violating constituent, which is not manufactured by stationary
substitution or by the nonzero CP family invariant.


Paper~III-G, Appendix~\ref{p3f:v2171:app:stationarychecks} verifies the higher
contact recursion in finite exterior-algebra examples; these are formula
checks rather than selected native action coefficients.


\subsubsection{Claim-complete source-cyclic determination}
\label{p3f:sec:v2186-source-cyclic-handoff}
For a fixed stationary tagged pole claim, the consumed object is the actual
complete action on its source-cyclic support quotient.  Paper~IV,
Lemma~\ref{p4:lem:v2186-exact-source-cyclic-action} and
Theorem~\ref{p4:thm:v2186-common-source-cyclic-polynomial}, prove exact
absence of complementary re-entry and the common-family polynomial
criterion.  Thus either determine that actual quotient or prove its
polynomial common across the independently licensed source/action family.
A complete finite member consumed by
Theorem~\ref{p3f:thm:v1148-finite-metric-constructor} is one sufficient
route; uniqueness of the entire quantum representation is not the
pole-readout criterion.  The complete parent must still justify the
source restriction and actual support, retaining every contributing
coupling, boundary and mixed return.

For this fixed pole, a scalar shift $H\mapsto H+cI$ leaves the source
commutator, ground projector and polynomial unchanged, by Paper~IV,
Corollary~\ref{p4:cor:v2186-scalar-center-pole-invisibility}.  The scalar
value is therefore not an independent pole-location selector.  The
complete metric/source jets, central derivatives and vacuum allocation
remain consumed by the metric, stress, absolute-energy and action-value
claims of the finite-metric constructor.  Neither this scope reduction
nor common-family availability evaluates the actual electron external,
dual-field and complete remainder returns.  Paper~III-G, Appendix~
\ref{p3f:app:v2174-native-capacity-domain} records the still uninstantiated
complete-member route and its exact first source equation; C01/A05 physical
identification and consumed limits remain open.


\subsection{common-profile action-response descent}
\label{p3f:sec:v164-common-profile-response}

\begin{theorem}[Same-action profile invariance]
\label{p3f:thm:v164-same-action-profile-invariance}
On the generated common-heat branch, the joint action is represented by
\[
 \operatorname{Tr}e^{-D^2/\Lambda_{\rm eff}^2},
 \qquad
 \Lambda_{\rm eff}^2=\Lambda^2/a.
\]
Every gauge, scalar, gravitational, and mixed response operator obtained from
that action by field variation is invariant under
\[
 (a,\Lambda)\longmapsto(\rho a,\sqrt\rho\Lambda),
 \qquad \rho>0.
\]
In particular, response is defined by
\[
 \mathcal R=\delta S_\star,
 \qquad K=\delta^2S_\star,
 \qquad V=\delta^3S_\star,
\]
not by a separate response profile or normalization.  Numerical response
values still consume the actual state/history and the declared unit map.
\end{theorem}

\subsection{General same-action response descent}
\label{p3f:sec:v236-same-action-response}
\begin{theorem}[Same-action force, Hessian and vertex descent]
\label{p3f:thm:v236-same-action-response-descent}
Let \(S\) be a claim-complete twice/three-times differentiable physical action
on a declared configuration owner \(\mathcal C\), and let \(x\in\mathcal C\)
be the actual configuration member.  Then the force, Hessian and cubic vertex
\[
 E_x=\delta S_x,\qquad K_x=\delta^2S_x,\qquad V_x=\delta^3S_x
\]
are generated jets of the same action rather than independent physical laws.
For a linear physical inclusion \(i:\mathcal H\hookrightarrow T_x\mathcal C\),
\[
 K_{\mathcal H}=i^*K_xi.
\]
If \(T_x\mathcal C=\mathcal H\oplus\mathcal I\), the internal variables are
stationary, and \(K_{\mathcal I\mathcal I}\) is invertible on the declared
gauge/constraint quotient, then the reduced response is the Schur complement
\[
 \boxed{K_{\rm eff}=K_{\mathcal H\mathcal H}
 -K_{\mathcal H\mathcal I}K_{\mathcal I\mathcal I}^{-1}
 K_{\mathcal I\mathcal H}.}
\]
Consequently any downstream response role proved to equal one of these same-action
descendants is generated on that owner.  This theorem does not create the
action, actual background, gauge quotient, inverse domain or calibration.
\end{theorem}

\section{Scoped minimal complete Lorentzian gravity and source action}
\label{p3f:sec:v169-minimal-complete-action}

This section assembles and completes the variational presentation of
already inherited results; it does not select a new gravity profile.
\emph{Complete} means that every retained variable has its Euler equation,
the lapse constraint is retained, and the specified endpoint variations
and source backreaction belong to the same action. The domain is the
certified nonlinear metric--\(f(R)\) sector, with an optional actual native
free Dirac mass member. This free subsystem is not identified with the full
direct K16 particle packet. The claim-scoped residual action requirements,
and the exact global nonlocal microscopic frontier of
Remark~\ref{p3f:rem:v118-parent-first-full-action-frontier}, remain unchanged.

\subsection{Inherited coefficient and canonical boundary completion}
Use \(\hbar=c=1\), signature \((-+++)\), and \(C:=C_7=C_{56}\).
In this section \(M^2:=U/4\) denotes the normalized gravitational
coefficient of Paper~VI-A, Eq.~\eqref{p8:eq:Lagrangian}, not a second raw
represented joint-Dirac coefficient or independently chosen mass scale.
It is the scalar-response coefficient coordinate
\(M^2/2=1/(16\pi G_{\rm scal})\) in these units.  Claiming that the
following action is the physical descendant of the raw joint parent
requires the actual projection and source normalization of
Proposition~\ref{p3f:prop:v2211-newton-owner-matching}; the trace ratio
alone does not establish this action identity.
Theorem~\ref{p3f:thm:v263-dressed-born-legendre-descent} supplies exactly
\begin{align}
\label{p3f:eq:v169-inherited-hessian}
 D_2(R)&=(8\sqrt\alpha)^{n_2^+(R)},&
 \chi(R)&=D_2(R)
 \frac{(1+n_7^+(R))(R^2-C^2)}{(U+C)R^2},\\
\label{p3f:eq:v169-inherited-gravity-coefficient}
 f(R)&=R-C+\int_C^R(R-x)\chi(x)\,dx,&
 F(R)&=1+\int_C^R\chi(x)\,dx .
\end{align}
The coefficient data are inherited. Neither the secant readout nor an
\(R^2\) approximation replaces the local Hessian \(\chi=f_{RR}\).
On the active interval \(R>C\), \(\chi>0\);
\(f(C)=0\), \(F(C)=1\), and \(\chi(C)=0\).

\begin{theorem}[Canonical endpoint completion of the inherited gravity action]
\label{p3f:thm:v169-complete-gravity-action}
Let \(\mathcal M\) have smooth canonical non-null boundary, outward unit
normal \(n\), \(\epsilon=n_\mu n^\mu=\pm1\), induced metric \(h\), and
\(K=h^{\mu\nu}\nabla_\mu n_\nu\). Fix \(h_{ij}\) and \(F=f'(R)\)
there. On a regular active interval the complete gravitational
Dirichlet action is
\begin{equation}
\label{p3f:eq:v169-gravity-action}
\boxed{
 S_g=\frac{M^2}{2}\int_{\mathcal M}\sqrt{-g}\,f(R)\,d^4x
 +M^2\sum_{\partial\mathcal M}\epsilon
       \int\sqrt{|h|}\,F(R)K\,d^3x .}
\end{equation}
Interior compactly supported variations do not need the endpoint term.
On the closed \(S^3\) cosmological region only spacelike temporal
endpoints occur. No corner, null-boundary, asymptotic or noncanonical
wall completion is asserted.
\end{theorem}
\begin{proof}
On \(\chi>0\), invert \(F=f'(r)\) on the chosen interval and define
\begin{equation}
\label{p3f:eq:v169-legendre-potential}
 P(F):=Fr(F)-f(r(F)),\qquad P_F=r(F).
\end{equation}
The Jordan presentation is
\begin{equation}
\label{p3f:eq:v169-jordan-action}
 S_{g,J}=\frac{M^2}{2}\int\sqrt{-g}\,[FR-P(F)]\,d^4x
       +M^2\sum_\partial\epsilon\int\sqrt{|h|}\,FK\,d^3x .
\end{equation}
Interior variation of \(F\) gives \(R=P_F=r(F)\), hence
\(F=f'(R)\) and \(FR-P=f(R)\). Conversely every metric-action
solution on this interval supplies that \(F\) and satisfies the auxiliary
equation. Thus these presentations are exactly equivalent on this
domain; \(F\) is the scalaron already present, not another physical scalar.

The curvature variation contains normal derivatives of \(\delta g\).
Variation of the \(2FK\) integrand cancels those derivatives with the
same coefficient. The residual boundary terms multiply \(\delta h_{ij}\)
or \(K\delta F\), and therefore vanish under the specified data.
Since \(\delta F=\chi\,\delta R\), fixing \(F\) is equivalent to
fixing \(R\) on the active interval. Fixing only \(h_{ij}\) is not
the higher-derivative Dirichlet problem asserted here
\cite{DyerHinterbichler2009BoundaryTerms}.
\end{proof}

For a retained covariant source action \(S_s\), put
\(T_{\mu\nu}:=-2(\sqrt{-g})^{-1}\delta S_s/\delta g^{\mu\nu}\).
The same interior variation gives
\begin{equation}
\label{p3f:eq:v169-metric-euler}
\boxed{
 \begin{aligned}
 FR_{\mu\nu}-\tfrac12f g_{\mu\nu}
 +(g_{\mu\nu}\Box-\nabla_\mu\nabla_\nu)F&=M^{-2}T_{\mu\nu},\\
 3\Box F+FR-2f&=M^{-2}T .
 \end{aligned}}
\end{equation}
Indeed \(\delta R\) is \(R_{\mu\nu}\delta g^{\mu\nu}\) plus the
contracted second derivatives of \(\delta g^{\mu\nu}\); two
integrations by parts give the displayed derivatives of \(F\).
This specializes Theorem~\ref{p3f:thm:general-local-fR-variation},
without an additional native-cyclic \(\mathrm{REN}_0\) subtraction.
For \(F=f'(R)\) the metric tensor on the left has identically zero
divergence. The source Euler equations supply its corresponding Ward
identity and propagation of the metric constraint.

There is one gravitational functional. Its linear curvature term,
nonlinear response and terminal intercept already belong to \(f\).
Another Einstein term, another \(-C\) subtraction, a second momentum
scalaron or an extra Born half-factor changes it. A source genuinely
depending on \(R\) must retain its full curvature variation: replacing
that \(R\) by the bare inverse \(r(F)\) can change the theory.
If required, the exact constrained presentation is
\(\int\sqrt{-g}\,[M^2f(r)/2+\Pi_R(R-r)]+S_s[g,r]\):
\(\delta\Pi_R\) enforces \(r=R\), and \(\delta r\) retains both
the gravitational and source curvature derivatives. The free source
below has no such dependence.

\subsection{Explicit free source and derived scalaron interaction}
\begin{theorem}[Closed gravity--Dirac subsystem on the native free branch]
\label{p3f:thm:v169-explicit-free-dirac-source}
Retain an actually generated constant invariant Dirac mass member \(m\),
with no omitted gauge or scalar force, on the admitted time-oriented
Lorentzian spin region. Use its inherited spin structure and
clock-hermitizing pairing, whose future-normal contraction is the
positive kinetic form, and the Clifford convention
\(\{\gamma^\mu,\gamma^\nu\}=-2g^{\mu\nu}I\).
The canonically normalized source action is
\begin{equation}
\label{p3f:eq:v169-dirac-action}
 S_D=\int\sqrt{-g}\left[
 \frac{i}{2}\{\bar\psi\gamma^\mu\nabla_\mu\psi
       -(\nabla_\mu\bar\psi)\gamma^\mu\psi\}
 -m\bar\psi\psi\right]d^4x .
\end{equation}
Source variations have compact interior support or fixed compatible
temporal endpoint data. Then \(S_{\min}=S_g+S_D\) is a closed
classical variational system on this sector: the Dirac equations and
the same-action Dirac stress supply the source in
Eq.~\eqref{p3f:eq:v169-metric-euler}. No independent scalaron--fermion
coupling or decay efficiency is introduced.
\end{theorem}
\begin{proof}
Independent field and dual variation gives their first-order Dirac
equations; the symmetrized expression differs from the canonical clock
form by its compatible endpoint term. The future-normal positive pairing
is normalized by Theorem~\ref{p3f:thm:v1148-finite-metric-constructor}.
The common raw fermion weight has been absorbed into that normalization
and is not reapplied to its canonical Hamiltonian. Variation of the same
metric before restriction defines its stress and on-equation Ward law.

On \(F>0\), set
\begin{equation}
\label{p3f:eq:v169-einstein-variables}
 g_E=Fg_J,\qquad
 \phi=\sqrt{\tfrac32}\,M\log F,\qquad
 \psi_E=F^{-3/4}\psi_J .
\end{equation}
The conformal curvature identity gives the Einstein term and
\(-(\nabla_E\phi)^2/2\). Its boundary derivative of \(\log F\)
cancels that from the transported \(FK\) term, leaving ordinary
Einstein GHY. The exact potential and mass are
\begin{equation}
\label{p3f:eq:v169-einstein-potential-and-mass}
\boxed{
 V_E(\phi)=\frac{M^2P(F)}{2F^2},\qquad
 m_E(\phi)=\frac{m}{\sqrt F}
          =m\exp[-\phi/(\sqrt6M)] .}
\end{equation}
Here \(\sqrt{-g_J}=F^{-2}\sqrt{-g_E}\) and
\(\bar\psi_J\psi_J=F^{3/2}\bar\psi_E\psi_E\), giving exactly
the remaining factor \(F^{-1/2}\) in the mass term. The kinetic
term is conformally transported, not independently reweighted. Thus
\[
 S_{\min,E}=\int\sqrt{-g_E}\left[
 \frac{M^2}{2}R_E-\frac12(\nabla_E\phi)^2-V_E
 +\mathcal L_{D,E}(m_E)\right]d^4x+S_{\rm GHY,E}.
 \]
\end{proof}

At an active constant \(F_*>1\), expansion gives the derived vertex
\begin{equation}
\label{p3f:eq:v169-derived-mass-vertex}
 \mathcal L_{\rm int}
 =\frac{m}{\sqrt6M\sqrt{F_*}}\,
       \delta\phi\,\bar\psi_E\psi_E+O(\delta\phi^2).
\end{equation}
The exact action keeps the full exponential. Its \(F_*=1\)
coefficient is an active-side limiting value, not a terminal decay
vertex after removal of the scalaron. At \(m=0\) the classical
conformal mass coupling vanishes. A trace anomaly, collision term or
other interaction requires its actual sector. Any consumed central
density keeps its inherited allocation between source and gravity.

\subsection{Closed-FLRW lapse, momenta and phase action}
\begin{theorem}[Constraint-preserving homogeneous reduction]
\label{p3f:thm:v169-flrw-lapse-and-phase-action}
For \(ds^2=-N(t)^2dt^2+a(t)^2d\Omega_{3,k}^2\), put
\(k=1\), \(V_\gamma=2\pi^2\), \(D_t=N^{-1}d/dt\),
and \(H=D_ta/a\). The Jordan bulk and temporal endpoint reduce to
\begin{equation}
\label{p3f:eq:v169-flrw-lagrangian}
\boxed{
 L_g=\frac{M^2V_\gamma}{2}
 \left[-\frac{6Fa\dot a^2}{N}
       -\frac{6a^2\dot a\dot F}{N}
       +6kNaF-Na^3P(F)\right].}
\end{equation}
For an instantaneous homogeneous source with complete proper energy
\(\widehat H_0(a)\), lapse variation before \(N=1\) gives
\begin{equation}
\label{p3f:eq:v169-flrw-constraint}
\boxed{
 3F(H^2+k/a^2)=\frac{P(F)}2-3HD_tF+\frac{\rho_s}{M^2},
 \qquad \rho_s=\frac{\langle\widehat H_0\rangle}{V_\gamma a^3}.}
\end{equation}
\end{theorem}
\begin{proof}
The curvature equals
\[
 R=6\left[\frac{\ddot a}{aN^2}+\frac{\dot a^2}{a^2N^2}
 -\frac{\dot a\dot N}{aN^3}+\frac{k}{a^2}\right].
 \]
Integration by parts and the temporal \(FK\) term give \(L_g\).
Varying \(L_g-N\langle\widehat H_0\rangle\) in \(N\) gives
the constraint. Variation in \(F\) gives \(R=P_F\), while the trace gives
\[
 D_t^2F+3HD_tF+\frac{2f-RF}{3}
       =\frac{\rho_s-3p_s}{3M^2}.\qedhere
 \]
\end{proof}

With \(B_g:=M^2V_\gamma\), direct Legendre transformation gives
\begin{align}
\label{p3f:eq:v169-flrw-momenta}
 p_a&=-6B_gFa\dot a/N-3B_ga^2\dot F/N,&
 p_F&=-3B_ga^2\dot a/N,\\
\label{p3f:eq:v169-flrw-hamiltonian}
 \mathcal H_g&=-\frac{p_ap_F}{3B_ga^2}
 +\frac{Fp_F^2}{3B_ga^3}-3B_gkaF+\frac{B_g}{2}a^3P(F).
\end{align}
Solving successively the two momentum relations for \(\dot a/N\)
and \(\dot F/N\) proves the displayed Hamiltonian. The phase action is
\[
 \int[p_a\dot a+p_F\dot F
       -N(\mathcal H_g+\langle\widehat H_0\rangle)]dt
       +S_{s,\rm kin}.
 \]
The primary lapse momentum vanishes, and the secondary constraint is
\(\mathcal H_g+\langle\widehat H_0\rangle=0\). Fixing the lapse
before variation deletes this equation. Local Cauchy closure consumes
a regular active chart and constrained initial data. Moving sources
with velocity insertions instead retain the complete momenta and
principal matrix of
Proposition~\ref{p3f:prop:complete-source-cauchy-matrix}.

\subsection{Concrete finite source, stress and pair block}
On the actual fixed-rank homogeneous free Dirac band of
Section~\ref{p3f:sec:v1148-free-metric-dirac}, half-density
normalization removes the \(a^3\) kinetic measure. In its fixed
coherent frame the complete energy is
\begin{equation}
\label{p3f:eq:v169-finite-free-source}
\boxed{
 h_0(a)=P_m[a^{-1}\mathcal D_\gamma+m\beta_D]P_m,\qquad
 \widehat H_0=d\Gamma(h_0)+\varepsilon_0(a)I .}
\end{equation}
The band and mass are inherited data, not new selectors. Retain the
invariant complete conjugate-sign band, with \(\beta_D^2=I\)
and \(\{\mathcal D_\gamma,\beta_D\}=0\). The homogeneous coframe
has no rotational clock connection on this branch. General moving
carriers retain the normalization connection and every interaction
slot of Theorem~\ref{p3f:thm:v1148-finite-metric-constructor}.
Discarding the negative component or \(\varepsilon_0\) changes
the owner and its metric source.

\begin{theorem}[Complete finite source and homogeneous work]
\label{p3f:thm:v169-finite-source-work}
For supplied initial density \(\sigma_i\), the mathematical amplitude
presentation \(\Psi(t_i)=\sqrt{\sigma_i}\) has real action
\begin{equation}
\label{p3f:eq:v169-homogeneous-joint-action}
\boxed{
 S_{\rm hom}=\int\left[
 L_g+\frac{i}{2}\operatorname{Tr}
       (\Psi^\dagger\dot\Psi-\dot\Psi^\dagger\Psi)
 -N\operatorname{Tr}(\Psi^\dagger\widehat H_0\Psi)
 \right]dt .}
\end{equation}
It gives \(i\dot\Psi=N\widehat H_0\Psi\) and
\(\sigma=\Psi\Psi^\dagger\), with averaged homogeneous ports
\begin{equation}
\label{p3f:eq:v169-finite-source-stress}
 i\dot\sigma=N[\widehat H_0,\sigma],\qquad
 \rho_s=\frac{\operatorname{Tr}(\sigma\widehat H_0)}{V_\gamma a^3},
 \qquad
 p_s=-\frac{\operatorname{Tr}(\sigma\partial_a\widehat H_0)}
                         {3V_\gamma a^2}.
\end{equation}
They satisfy the exact work identity
\begin{equation}
\label{p3f:eq:v169-finite-source-conservation}
\boxed{D_t\rho_s+3H(\rho_s+p_s)=0 .}
\end{equation}
Their embedding as an exact FLRW source additionally needs homogeneous
energy, isotropic stress and vanishing spatial flux. An actual
spatial-isometry-invariant full multiplet/state is a sufficient
certificate of these properties, not a universally necessary
state characterization. The averaged minisuperspace identity alone
does not establish them for every finite state.
\end{theorem}
\begin{proof}
Independent amplitude and adjoint variation gives their Schr\"odinger
equations. They preserve positivity, trace and spectrum of \(\sigma\);
the amplitude presents a supplied state rather than prepares it.
The lapse and scale variations give the two ports, including the
actual derivatives of \(\varepsilon_0\). Trace cyclicity gives
\(\operatorname{Tr}(\dot\sigma\widehat H_0)=0\). Differentiate
\(V_\gamma a^3\rho_s=\operatorname{Tr}(\sigma\widehat H_0)\)
to obtain the work identity and hence constraint propagation when
combined with the gravitational Euler equations. No Wick factorization
or one-body-state assumption is used.
\end{proof}

This is expectation-valued finite matter coupled to classical gravity,
not quantization of gravity. The state is Cauchy data, not a condition
for defining the bare gravitational action. The native guard, Ward
and projectability hypotheses remain necessary for an inhomogeneous
finite metric extension.

On a paired \(\pm\lambda\) block, put \(p=|\lambda|/a\),
\(p_i=|\lambda|/a_i\), and \(E_i=\sqrt{p_i^2+m^2}\). Fixed
initial diagonalization and the particle-hole CAR presentation give
\begin{align}
\label{p3f:eq:v169-explicit-pair-block}
 A_\lambda(a)&=\frac{pp_i+m^2}{E_i},&
 B_\lambda(a)&=\frac{m(p-p_i)}{E_i},\\
 \widehat H_{\lambda,0}
 &=A_\lambda(n_a+n_b-1)
   +B_\lambda(a^\dagger b^\dagger+ba)+\varepsilon_0 I .
\end{align}
Rotating \(p\sigma_1+m\sigma_3\) by its fixed initial angle gives
\(A_\lambda\sigma_3+B_\lambda\sigma_1\); replacing its negative
component by the hole creation operator proves the displayed ordered
word. Its even block in order \((|ab\rangle,|0\rangle)\) is
the same matrix, plus its central allocation. This is an explicit pair
coupling, not a fitted decay rate. The formula consumes \(E_i>0\).
At \(m=0\) and \(\lambda\ne0\), \(B_\lambda=0\), so the
conformal vacuum produces no pairs. A zero \(m=\lambda=0\) block
has zero bilinear generator and is treated directly, not divided by
\(E_i\). A moving instantaneous frame instead retains the equivalent
\(i(D_t\theta)(a^\dagger b^\dagger-ba)/2\) in the proper
generator \(\widehat H_0\), or
\(i\dot\theta(a^\dagger b^\dagger-ba)/2\) in the coordinate
generator \(N\widehat H_0\). Here
\(\theta=\operatorname{atan2}(p,m)\) is chosen on a continuous
\(E>0\) branch, including the massless sector; deleting the
connection is not equivalent. This distinction is retained before
fixing the lapse.

Frame consistency follows from
\[
 \begin{gathered}
 a_E=\sqrt F\,a_J,\quad N_E=\sqrt F\,N_J,\quad
 a_E^{3/2}\psi_E=a_J^{3/2}\psi_J,\\
 N_E[\mathcal D_\gamma/a_E+(m/\sqrt F)\beta_D]
 =N_J[\mathcal D_\gamma/a_J+m\beta_D].
 \end{gathered}
 \]
The central allocation is transported likewise,
\(N_E\varepsilon_{0,E}=N_J\varepsilon_{0,J}\);
no new normal ordering is chosen in the other frame.
In Einstein variables put \(D_E=N_E^{-1}d/dt\) and
\(H_E=D_Ea_E/a_E\). The actual source derivative
\[
 J_\phi=(V_\gamma a_E^3)^{-1}
             \operatorname{Tr}(\sigma\partial_\phi\widehat H_{0,E})
 \]
gives \(D_E\rho_D+3H_E(\rho_D+p_D)=J_\phi D_E\phi\) and
\(D_E^2\phi+3H_ED_E\phi+V_{E,\phi}=-J_\phi\).
Its mass contribution is
\(-m_E\operatorname{tr}(C_D\beta_D)/(\sqrt6M V_\gamma a_E^3)\),
with \(C_D\) the exact one-body expectation matrix; every actual
central derivative remains. Source and scalaron work cancel exactly,
with no independently assigned exchange current.

\subsection{Reheating, terminal strata and preserved frontiers}
The free source is not shown to equal the direct K16 packet; its
production and stress require the native boundary-owner matching
before that identification. Section~\ref{p3f:sec:v1134-successor-cell-schur}
and Remark~\ref{p3f:rem:v1134-stationary-descendant-firewall}
already supply the relevant Schur descendant and no-double-counting
rule. The reduced response is not appended again to its retained bulk.
The K8--K16 readout \(\eta(K)=(K-8)/(24-K)\) is not a new
scalar field. Its cumulative release \(\mathcal I_{16}\eta\)
and current \(Q_\partial=\mathcal I_{16}\dot\eta\) do not
determine all off-shell source jets or an independent action.
For radiation-like output, its physical density instead is
\begin{equation}
\label{p3f:eq:v169-radiation-dilution}
 \rho_r(t)=\frac{a_i^4\rho_{r,i}}{a(t)^4}
 +\frac{\mathcal I_{16}}{a(t)^4}
       \int_{t_i}^{t}a(s)^4\dot\eta(s)\,ds ,
\end{equation}
obtained by integrating \(\dot\rho_r+4H\rho_r=Q_\partial\)
in proper time. This is the existing dilution accounting, not an
extra efficiency or a replacement for produced-state evolution.
Particle branching or collisions consume their actual interactions.
Generated Yang--Mills, odd-scalar and Yukawa slots are retained when
consumed, not added by default to vacuum gravity. A retained nonminimal
scalar, for example, changes the curvature coefficient to
\(\mathcal A=M^2F/2-\operatorname{Tr}_F\Phi^2/(96\pi^2)\);
its matching canonical endpoint is
\(2\sum_\partial\epsilon\int\sqrt{|h|}\,\mathcal A K\),
with \(\Phi\) also fixed there.

At K16, \(\chi>0\), and the scalaron persists with its canonical
momenta and complete source Cauchy data. With no surface action,
matching is of those canonical data, not only \(a,F\); the flux law of
Theorem~\ref{p3f:thm:v266-gp-exact-k16-junction-condition}
remains unchanged. No instantaneous scalaron reset is inserted.
At the terminal endpoint the inverse \(r(F)\) is not continued
through \(\chi=0\). On the separately licensed terminal stratum
\(F=1\) is fixed and has no scalaron tangent. Taking the one-sided
value \(\lim_{F\downarrow1}P(F)=C\) restricts the Jordan action to
\begin{equation}
\label{p3f:eq:v169-terminal-affine-action}
 S_{g,\rm term}=\frac{M^2}{2}\int\sqrt{-g}\,(R_g-C)\,d^4x
       +M^2\sum_\partial\epsilon\int\sqrt{|h|}\,K\,d^3x .
\end{equation}
It gives \(\Lambda=C/2\) and a formal Einstein vacuum at \(R_g=2C\), not \(R_g=C\); if this terminal representative applies only for \(R_g\le C\), the formal root is outside its curvature domain. See Remark~\ref{p3f:rem:v346-terminal-affine-extrapolation}.
The latter would come from varying a scalaron absent on that quotient.
The reduced auxiliary \(f(r)+f'(r)(R-r)\) has equation
\(f''(r)(R-r)=0\), loses active equivalence at the endpoint, and
is not used to enforce a terminal scalar equation. The intercept
is already fixed by
Corollary~\ref{p3f:cor:v2118-native-cyclic-ren0-firewall}.

An occurring rank-changing handoff still consumes native collar,
source and constraint matching. Writing the affine restriction proves
neither automatic physical arrival nor a smooth active pole-mass limit.
Noncanonical boundaries, a general finite local Ward extension, the
actual produced-sector trace and thermal history, and the exact global
nonlocal microscopic parent retain their previous claim-specific
requirements. This section completes the declared variational
presentation and explicit free subsystem, not those different targets.


\setlength{\emergencystretch}{2em}
\subsection{Variational dynamical asymmetry}
\label{p3:subsec:scalars-vs-vectors}
\label{p3:sec:variational_asymmetry}

\begin{theorem}[Variational asymmetry under \((\mathrm{VDOM}_x^{\min})\)]
\label{p3:thm:completed-carrier-dynamical-asymmetry}
\label{p3:prop:variational-asymmetry}
\label{p3:thm:asymmetry-rigorous}
Assume the premise-minimal Paper-II-C/GP2 reciprocal-role domain completed in
Paper~III-A, including the generated reciprocal generator-germ localization,
and the fixed momentum seed
\[
 \mathcal H_p\cong\mathcal H_1,\qquad m_p=1,
\]
and \((\mathrm{VDOM}_x^{\min})\).  Let the reciprocal-position metric occur in the
fluctuated position operator while the compact momentum representation and
its constituted geometric data are held fixed; the Levi--Civita and dual
connections co-vary with the position metric.  Then the momentum data have
no Euler--Lagrange equation in this branch, whereas the position metric is a
variational field.  Consequently all local macroscopic metric dynamics is
carried by the reciprocal-position polarization.  On the inherited
scalar-succession orientation, a homogeneous solution is expanding exactly
where \(\partial_{\widetilde t}a>0\) and contracting where it is negative.
The polarization theorem does not determine that derivative.
The Euler--Lagrange asymmetry itself uses only the declared variation space
\((\delta q_x,0)\) and the occurrence of the metric in \(q_x\) but not in the
fixed datum \(q_p\).  The preceding spectral, cotangent, Born, and momentum-
level statements certify the provenance of these roles; they are not further
mathematical premises of the variation-space implication.
\end{theorem}

\begin{proof}
Paper~II-B fixes the momentum carrier and its derived label before Paper~III.
Paper~III-A completes the smooth reciprocal role supplied by the
premise-minimal Paper-II-C/GP2 closure, without imposing a finite position
level.  The optional metric-curry certificate realizes the full profile
family but is not used here.  Paper~III-A's generated generator-germ theorem
identifies the local reciprocal readout role with the cotangent target without
an additional bridge occurrence.  By \((\mathrm{VDOM}_x^{\min})\),
the compact momentum representation is held fixed, while the
reciprocal-position metric and connection enter the fluctuated position
operator and are varied.  The Euler--Lagrange variation therefore acts on the
position metric but not on the momentum radius or finite momentum carrier.
The local cotangent/Born gluing theorem places this metric data on the position
summand \(F=TU\oplus0\subset E_U\).  A homogeneous specialization determines
the evolution equation only after the action branch is fixed; it does not
prove the sign of \(\partial_{\widetilde t}a\) on a solution.
\end{proof}

\begin{remark}[Coarse homogeneous entropy diagnostic deferred]
The coarse macrocell entropy identity is a consistency diagnostic of an
already chosen expanding homogeneous orientation and is not used in the
positive locality/backreaction proof.  It is deferred to
Paper~III-G, Appendix~\ref{p3f:app:coarse-entropy}.
\end{remark}

\begin{theorem}[Reciprocal-position dynamical sector under the minimal action parent]
\label{p3:thm:locality_selection}
On Paper~III-E's generated gravitational heat process, whose one-profile
additive form belongs to the loop-derived class \(\Dcl\), together with
\((\mathrm{VDOM}_x^{\min})\),
Theorem~\ref{p3:thm:sectorwise_additivity},
Theorem~\ref{p3:thm:local_kunneth_connection},
and Proposition~\ref{p3:prop:variational-asymmetry}, any coarse-grained sector admitting a local two-derivative Lorentzian metric description on
reciprocal-position patches must be realized on the position summand of the exact doubled bundle
\[
F=TU\oplus0\subset E_U,
\]
not on the compact momentum sector.
Equivalently, the compact momentum geometry supplies fixed spectral/UV data, while any candidate
macroscopic local metric sector is realized only on the reciprocal-position side.
\end{theorem}

\begin{proof}
By Theorem~\ref{p3:thm:local_kunneth_connection}, the unique anchored no-insertion connection on
the exact doubled bundle restricts on the position summand \(F=TU\oplus0\) to the ordinary
Levi--Civita connection of the effective metric on \(U\).  Thus any candidate local
two-derivative metric dynamics must be realized through the base position datum whose metric is
actually varied in the action principle.

Proposition~\ref{p3:prop:variational-asymmetry} distinguishes the two reciprocal sides:
the momentum radius \(\Rmom\) is fixed representation data inherited from the compact momentum
kinematics and is not varied in the spectral action, whereas the position-sector metric data enter
through the fluctuated operator \(D_x\mapsto D_{x,A}\) and are varied. Therefore the compact
momentum sector cannot carry the local dynamical metric degrees of freedom required for a local
two-derivative Lorentzian description.

Hence any coarse-grained sector admitting such a local metric description must be realized on the
reciprocal-position summand \(F=TU\oplus0\subset E_U\) over \(U\), not on the compact momentum sector.
\end{proof}

\begin{definition}[Analytic localization multiplier]
\label{p3:def:analytic-localizer}
Let \(U_0\Subset U\) be reciprocal-position charts and choose
\(\varphi_U\in C_c^\infty(U)\) with \(0\le\varphi_U\le1\), equal to one on a
neighborhood of \(\overline{U_0}\).  In expressions
\[
 \Tr\!\left(\varphi_U e^{-D^2/\Lambda_g^2}\right),
\]
\(\varphi_U\) is an analytic localization multiplier for the generated smooth-parent
heat-kernel asymptotic.  It is not a physical channel projector, sector
selector, or additional \(\Dcl\) profile, and it does not change finite-role
content.  The smeared trace extracts coefficients on \(U\).  Every variation
used below to infer a local field equation is restricted to compact support in
\(U_0\), where \(\varphi_U\equiv1\); hence no derivative-of-localizer term is
discarded.
\end{definition}

\begin{lemma}[Diagonal natural scalar-density handoff preserves the full variation]
\label{p3f:lem:diagonal-eh-full-variation}
For a natural pure-metric scalar density \(\ell\), let
\[
 S_{{\rm B},\ell}(g_x,g_p)
 =\tfrac12 I_{x,\ell}(g_x)+\tfrac12 I_{p,\ell}(g_p),
 \qquad j_{\rm B}(g)=(g,\beta_*g),
\]
with \(I_{p,\ell}(\beta_*g)=I_{x,\ell}(g)\).  Then
\[
 S^{\rm eff}_{x,\ell}:=j_{\rm B}^*S_{{\rm B},\ell}=I_{x,\ell},
 \qquad
 \delta S^{\rm eff}_{x,\ell}=\delta I_{x,\ell}.
\]
Varying a single off-diagonal sector while freezing the other would instead
give \(\tfrac12\delta I_{x,\ell}\); that is a different matching choice and
is not the Born-diagonal pullback.  The legacy \((\mathrm{EH}_x)\) statement
is the specialization \(\ell=\scal\).
More generally, coefficients \(\alpha\) and \(\beta\) give
\[
 j_{\rm B}^*\bigl(\alpha I_{x,\ell}+\beta I_{p,\ell}\bigr)
 =(\alpha+\beta)I_{x,\ell}.
\]
Thus equality of the individual weights is unnecessary; preservation of one
full copy is exactly the condition \(\alpha+\beta=1\).
\end{lemma}

\begin{proof}
Pullback along \(j_{\rm B}\) makes both half-weighted terms equal to
\(\tfrac12I_{x,\ell}(g)\).  Their sum and its differential are therefore the
full functional and full differential.
\end{proof}

\begin{theorem}[Local spectral asymptotics and protected pure-metric branch]
\label{p3:thm:local_spectral_asymptotics}
Assume
\[
 (\mathrm{VDOM}_{x}^{\min})
 +(\mathrm{GRAV\!-\!ASYM})
 +\mathsf C_{\rm IR}^{\rm der}.
\]
Let \(U_0\Subset U\) and \(\varphi_U\) be as in
Definition~\ref{p3:def:analytic-localizer}.  In the semiclassical derivative
expansion regime, define the raw position trace
\[
 S^{\rm raw}_{x,U}
 :=\Tr\!\left(\varphi_U
 e^{-D_{x,A}^2/\Lambda_g^2}\right).
\]
It has the local smooth-parent asymptotic expansion
\[
 S^{\rm raw}_{x,U}
 \sim\sum_{n\ge0}\Lambda_g^{4-n}
 a_n(D_{x,A,U}^2;U),
\]
whose coefficients are local invariants of the projected metric and bundle
connection \cite{Vassilevich2003HeatKernel}.  Let
\(S^{\rm raw}_{x,U,{\rm met}}\) denote the sum of its constituents built only
from the metric and its Levi--Civita curvature.  With
\(S_{{\rm B},U,{\rm met}}:=\tfrac12S^{\rm raw}_{x,U,{\rm met}}
+\tfrac12S^{\rm raw}_{p,U,{\rm met}}\), set
\[
 S^{\rm full}_{U,{\rm met}}
 :=j_{\rm B}^{*}S_{{\rm B},U,{\rm met}}.
\]
Paper~III-B assigns the raw position output the role weight \(1/2\), whereas
the derived natural scalar-density pullback gives one full Born-diagonal copy
of every term in \(S^{\rm raw}_{x,U,{\rm met}}\).  Consequently the effective
gravitational
derivative expansion has the form
\[
 S_U^{\rm grav}\sim\int_U\!\left(
 \alpha_0\Lambda_g^4
 +\alpha_2\Lambda_g^2\scal[g]
 +\alpha_4\mathcal L_{4\partial}+\cdots\right)\sqrt{|g|}\,d^4x,
\]
with \(\alpha_0,\alpha_2\), and every pure-metric higher coefficient carrying
the full Born-diagonal normalization inherited from Paper~III-E.  In
particular the protected internal support coefficient of both the vacuum and
Einstein terms is \(1/8\).  Nonmetric bundle and matter constituents remain
half-weighted.  This is a Euclidean local derivative expansion and uses no
Wick hypothesis; an exact spectral trace may additionally retain controlled
nonlocal remainders.
\end{theorem}

\begin{proof}
The local reciprocal atlas and anchored doubled connection supply
\(T^*U\), \(E_U=TU\oplus T^*U\), and the Levi--Civita restriction on
\(F=TU\oplus0\).  Paper~III-E's generated heat theorem and Paper~III-B's loop
theorem supply one \(\Dcl\) profile and the sectorwise additive functional
\(\lambda_{\rm B}\circ\mathfrak T_{\rm sec}\); hence its position output is
\(\tfrac12S^{\rm raw}_{x,U}\), not the unweighted raw trace.  In the smooth
semiclassical regime, the standard heat-kernel expansion gives the displayed
local invariants.  Under
\(\mathsf C_{\rm IR}^{\rm der}+(\mathrm{GRAV\!-\!ASYM})\),
Paper~III-B's co-realization theorem derives the rank-one real-channel
constituent certified by \((\mathrm{CORN})\).  Paper~III-E's Born-orbit
theorem derives the diagonal full-copy statement certified by
\((\mathrm{BD}_{\rm EH})\), fixing the full Einstein coefficient without
adding either tag as a hypothesis.
The same naturality proof applies to every pure-metric scalar density, so the
Born-diagonal pullback freezes the entire pure-metric coefficient family
before the derived \((\mathrm{VDOM}_x^{\min})\) varies only the local action
arguments.  No new factor of
one half is introduced in those channels, and no nonmetric channel is doubled.
\end{proof}

\begin{corollary}[Conditional Lorentzian action interpretation]
\label{p3f:cor:local-asymptotics-lorentzian-interpretation}
The Euclidean expansion of
Theorem~\ref{p3:thm:local_spectral_asymptotics} admits a Lorentzian action
interpretation only under the claim-specific continuation and
same-variable tags required by Paper~III-E, for example
\((\mathrm{WICK}_{2{\rm cyc}})+(\mathrm{WICK}_{\rm same})\) for the
corresponding two-point branch.  At the III-F stage a finite reconstruction
claim additionally requires \((\mathrm{WICK}_{\rm OS,fin})\).  On the
separately constituted positive-kernel/Hamiltonian-generated domains of
Paper~IV, the corresponding finite same-carrier/cyclic/reflection-positive
tags are \textsc{derived-later}.  Paper~IV also derives the Lorentzian action child for stationary stable
quadratic bosonic blocks from the same complex-time free-law parent.  The
later v1.74--v1.75 results derive the local parity-even Yang--Mills,
Einstein--Hilbert interior, and scalar-dependent through-\(\Phi^4\) action constituents on
their printed domains.  They do not promote odd fermionic terms, the exact
nonlocal heat trace, uncontrolled higher expansions, or an interacting
infinite-volume functional to a Lorentzian action.  Those claims retain the
residual \((\mathrm{ACT}_{C}^{\rm res})\) and, where applicable,
\((\mathrm{WICK}_{\rm int,\infty})\) interfaces.
\end{corollary}

\begin{corollary}[Derived-later stationary quadratic action branch]
\label{p3f:cor:stationary-quadratic-action-derived-later}
On the stationary stable finite free-law domain of
Paper~IV, Theorem~\ref{p4:thm:stationary-quadratic-real-form-pair}, the
Euclidean and Lorentzian quadratic action forms, their Euler operators, and
their positive-time dynamics are co-realized restrictions of one complex-time
parent.  Denote this theorem-status certificate by
\[
 (\mathrm{ACT}_{2,\mathrm{stat}}^{\rm der}).
\]
Accordingly, a III-F claim restricted to such a quadratic block does not need
an independent \((\mathrm{CONT}_{\rm cl})\) or same-carrier Wick premise for
that block.  The certificate is not a replacement for the later local
joint-Dirac action theorems: it does not determine the nonlinear Yang--Mills,
Einstein--Hilbert, or scalar coefficients.  Conversely, the v1.74--v1.75 local
bosonic closure does not determine the fermionic real structure, determinant
or Pfaffian phase, exact nonlocal/higher-order functional, or noncanonical
boundary/global completion.
\end{corollary}

\begin{theorem}[Parity-even uninserted kinetic trace closure]
\label{p3f:thm:joint-connection-uninserted-trace-closure}
On \((\mathrm{JOINT}_x^{\rm der})\), let \(E\to U\) be its generated finite
joint carrier with connection \(\nabla^E\) and curvature
\(\Omega^E\in\Omega^2(U,\End E)\).  Work in the unmarked \(\Dcl\) branch:
the action may use the generated Hilbert metric, adjoint, volume form, and
ordinary finite trace, but no new central projector, factor label, profile,
or positive endomorphism is inserted.  Restrict to real, derivative-free,
parity-even (equivalently full-frame \(O(4)\)-invariant on the Euclidean
representative), positive quadratic kinetic functionals.  Then every such
nonzero single-trace curvature functional is, up to one overall constant
\(C>0\),
\begin{equation}
\label{p3f:eq:joint-connection-canonical-trace-action}
 S_{\rm conn}
 =C\int_U\operatorname{Tr}_{E}
   \bigl((\Omega^E)^{\dagger}\Omega^E\bigr)\,d\mu_g .
\end{equation}
If the connection decomposes functorially as
\[
 \nabla^E
 =\nabla^{E_1}\otimes1\otimes\cdots
 +\cdots+1\otimes\cdots\otimes\nabla^{E_r},
\]
all component kinetic trace weights are therefore the representation indices
obtained by expanding the one trace in
\eqref{p3f:eq:joint-connection-canonical-trace-action}.  Independent
factor coefficients require an inserted central weight
\(Z=\sum_a z_aP_a\) and hence lie outside the unmarked generated branch.
Thus a common kinetic normalization is a theorem on the generated joint
family; no separate connection-occurrence residual remains.  The independent
orientation-odd density
\[
 \int_U\operatorname{Tr}_E(\Omega^E\wedge\Omega^E)
\]
is also natural, local, quadratic, single-trace, and uninserted, but it is
parity odd and therefore lies outside this kinetic-closure theorem.  Positivity
alone would not exclude its presence.
\end{theorem}

\begin{proof}
At a point of an oriented Euclidean four-dimensional representative, the
Hodge decomposition is
\(\Lambda^2=\Lambda^2_+\oplus\Lambda^2_-\).  The two summands are inequivalent
irreducible \(SO(4)\)-modules.  Hence every real invariant symmetric quadratic
form on the two-form factor is
\[
 c_+\|\Omega_+^E\|^2+c_-\|\Omega_-^E\|^2,
\]
with the literal ordinary carrier trace furnishing the endomorphism pairing.
Equivalently this is a linear combination of
\(\operatorname{Tr}(\Omega^E\wedge *\Omega^E)\) and
\(\operatorname{Tr}(\Omega^E\wedge\Omega^E)\).  Orientation reversal
exchanges the two Hodge summands, so parity invariance forces
\(c_+=c_-\).  The theta coefficient therefore vanishes; positivity makes the
remaining common coefficient strictly positive.  This proves
\eqref{p3f:eq:joint-connection-canonical-trace-action}.  Expanding the tensor
sum gives the usual representation-index weights from the same trace.
Changing one factor weight while leaving another fixed is equivalent to
replacing \(\operatorname{Tr}_E\) by
\(A\mapsto\operatorname{Tr}_E(ZA)\) for a non-scalar central
endomorphism \(Z\).  Such \(Z\) is precisely an extra insertion excluded by
the unmarked \(\Dcl\) branch.  A possible double-trace central term is likewise
not a single-trace closure and vanishes on the projective/traceless curvature
sector.  Thus the stated parity-even kinetic sector has no independent
common-normalization clause.  The displayed theta density demonstrates why
the parity qualifier is indispensable.
\end{proof}


\begin{definition}[Native local Yang--Mills phase-space domain]
\label{p3f:def:v174-native-yang-mills-phase-space-domain}
Use the derived \((\mathrm{JOINT}_x^{\rm der})\) certificate, the complete
\((\mathrm{CAUS})\) domain, and the parity-even positive uninserted kinetic
trace branch of
Theorem~\ref{p3f:thm:joint-connection-uninserted-trace-closure}.  Work on the
smooth stationary product representative
\[
 U=I_{\widetilde t}\times\Sigma,
 \qquad
 g_{\rm op}=-c_{\rm op}^{2}d\widetilde t^{\,2}+h,
 \qquad c_{\rm op}>0,
\]
with either closed \(\Sigma\) or the compact-support boundary rule of the
minimal local variation domain.  Remain on the native undeformed
same-principal-owner branch: the frozen grammar contains no independent
clock, spatial cometric, gauge constitutive tensor, medium response,
anisotropy, or other principal-order deformation.  Let \(P\to U\) be the
generated joint principal bundle.  Let
\[
 Q_\rho^\flat:\operatorname{ad}P\longrightarrow\operatorname{ad}^{*}P
\]
be the positive invariant fibre pairing induced by the one ordinary
joint-carrier trace, including its one common coefficient and the generated
factor representation indices, and put
\(Q_\rho^\sharp=(Q_\rho^\flat)^{-1}\).

The target is the local, real, parity-even, gauge-invariant, metric-natural,
two-derivative pure-gauge action and its regular constrained cotangent
presentation, modulo exact boundary terms.  No gauge-field member, temporal gauge, matter current, theta term,
boundary edge mode, finite regulator action, state, or history is selected by
this domain.
\end{definition}

\begin{lemma}[Clocked complex real form of the complete Yang--Mills curvature]
\label{p3f:lem:v174-yang-mills-complex-connection-real-form}
On the analytic first-jet subdomain of the domain above, additionally
complexify the oriented clock line and adjoint coefficient bundle, extend
\(Q_\rho\) and \(h\) complex bilinearly, and require the connection first-jet
germ and gauge parameter to be holomorphic on a contour sector joining the
Euclidean and oriented Lorentzian real slices, with no additional contour-end
term under the stated boundary rule.  Let
\[
 U_E=I_\tau\times\Sigma,
 \qquad
 g_E=c_{\rm op}^{2}d\tau^2+h
\]
be the Euclidean real-form sibling of the generated clock-product patch.
Write the same complexified connection one-form in its Euclidean and
Lorentzian real-form presentations as
\begin{equation}
\label{p3f:eq:v174-yang-mills-clocked-connection-real-form}
 A_E=A_\tau\,d\tau+a(\tau),
 \qquad
 A_L=A_0\,d\widetilde t+a(\widetilde t).
\end{equation}
Under the already oriented continuation
\begin{equation}
\label{p3f:eq:v174-yang-mills-clocked-time-continuation}
 \tau=i\widetilde t,
 \qquad d\tau=i\,d\widetilde t,
 \qquad \partial_\tau=-i\partial_{\widetilde t},
 \qquad A_\tau=-iA_0,
\end{equation}
the complete nonlinear curvature obeys
\begin{equation}
\label{p3f:eq:v174-yang-mills-curvature-real-form}
 F_{\tau i}
 =\partial_\tau A_i-\partial_iA_\tau+[A_\tau,A_i]
 \longmapsto -iF_{0i},
 \qquad
 F_{ij}\longmapsto F_{ij}.
\end{equation}
Thus, with
\[
 \mathcal E_E:=c_{\rm op}^{-1}
   (\partial_\tau a-d_aA_\tau),
 \qquad
 \mathcal E_L:=c_{\rm op}^{-1}
   (\partial_{\widetilde t}a-d_aA_0),
\]
the already generated parity-even ordinary-trace Euclidean curvature action
and its Lorentzian real-form child are
\begin{align}
\label{p3f:eq:v174-yang-mills-action-real-form-pair}
 S_{{\rm YM},E}[A_E]
 &=\frac{c_{\rm op}}2\int d\tau\int_\Sigma
   \left(\|\mathcal E_E\|_{Q_\rho,h}^{2}
        +\|B_a\|_{Q_\rho,h}^{2}\right)d\mu_h,\\
 S_{{\rm YM},L}[A_L]
 &=\frac{c_{\rm op}}2\int d\widetilde t\int_\Sigma
   \left(\|\mathcal E_L\|_{Q_\rho,h}^{2}
        -\|B_a\|_{Q_\rho,h}^{2}\right)d\mu_h.
\end{align}
On every stated boundary domain whose endpoint terms are transported or
vanish,
\begin{equation}
\label{p3f:eq:v174-yang-mills-action-weight-real-form}
 S_{{\rm YM},E}\longmapsto-iS_{{\rm YM},L},
 \qquad
 e^{-S_{{\rm YM},E}/\hbar}
 \longmapsto e^{iS_{{\rm YM},L}/\hbar}.
\end{equation}
The map includes the commutator term in
\eqref{p3f:eq:v174-yang-mills-curvature-real-form}; hence it continues the
cubic and quartic Yang--Mills interactions and is not merely a free-field or
physical-polarization statement.  This is a local holomorphic first-jet
real-form theorem.  It does not assert an OS measure, analytic continuation of
arbitrary global field histories, or an interacting infinite-volume limit.
\end{lemma}

\begin{proof}
Equality of the complexified one-form in
\eqref{p3f:eq:v174-yang-mills-clocked-connection-real-form} gives
\(A_\tau d\tau=A_0d\widetilde t\), hence \(A_\tau=-iA_0\).
Substitution into the full polynomial curvature gives
\eqref{p3f:eq:v174-yang-mills-curvature-real-form}, including
\([A_\tau,A_i]\mapsto-i[A_0,A_i]\).  Therefore
\(\mathcal E_E\mapsto-i\mathcal E_L\), whereas \(B_a\) is unchanged.
Substitution together with \(d\tau=i\,d\widetilde t\) proves
\eqref{p3f:eq:v174-yang-mills-action-weight-real-form}.  No gauge-field
member or contour-dependent coefficient is selected.
\end{proof}

\begin{theorem}[Generated joint gauge phase space and local Lorentzian action]
\label{p3f:thm:v174-generated-joint-gauge-phase-space-action}
On the domain of
Definition~\ref{p3f:def:v174-native-yang-mills-phase-space-domain}, write the
generated connection and curvature uniquely as
\[
 A=A_0\,d\widetilde t+a(\widetilde t),
 \qquad
 F_A=d\widetilde t\wedge
       \bigl(\dot a-d_aA_0\bigr)+B_a,
 \qquad
 B_a=d_\Sigma a+a\wedge a,
\]
and put
\[
 \mathcal E_A
 :=c_{\rm op}^{-1}\bigl(\dot a-d_aA_0\bigr).
\]
Put
\(\mathcal A_\Sigma:=\operatorname{Conn}(P_\Sigma)\).  The regular
connection-field cotangent role
\begin{equation}
\label{p3f:eq:v174-regular-yang-mills-phase-space}
 \mathcal P_{\rm YM,\Sigma}^{\rm reg}
 :=\mathcal A_\Sigma
   \times\Omega_c^2(\Sigma,\operatorname{ad}^{*}P_\Sigma)
 \cong T_{\rm reg}^{*}\mathcal A_\Sigma
\end{equation}
is not the particle cotangent space \(T^*U\).  It has the generated
tautological and symplectic forms
\begin{equation}
\label{p3f:eq:v174-yang-mills-tautological-symplectic-forms}
 \Theta_{\rm YM}|_{(a,\Pi)}(\delta a,\delta\Pi)
 :=\int_\Sigma\langle\Pi\wedge\delta a\rangle,
 \qquad
 \Omega_{\rm YM}:=-d\Theta_{\rm YM}.
\end{equation}
The cotangent-lifted spatial gauge action has equivariant moment map, up to a
simultaneous conventional overall sign,
\begin{equation}
\label{p3f:eq:v174-yang-mills-gauss-moment-map}
 \mu_{\rm G}(a,\Pi)=d_a\Pi,
\end{equation}
so \(\mu_{\rm G}^{-1}(0)\) is the source-free Gauss locus.

For an adjoint-valued spatial form \(\eta\) and an adjoint-dual-valued
spatial two-form \(\Pi\), use the pointwise positive norms determined by
\((Q_\rho,h)\), equivalently
\[
 \|\eta\|_{Q_\rho,h}^{2}d\mu_h
 :=\langle Q_\rho^\flat\eta\wedge *_h\eta\rangle,
 \qquad
 \|\Pi\|_{Q_\rho^{-1},h}^{2}d\mu_h
 :=\langle\Pi\wedge *_hQ_\rho^\sharp\Pi\rangle.
\]
Then the native branch has the unique positive local Hamiltonian
\begin{equation}
\label{p3f:eq:v174-generated-joint-yang-mills-hamiltonian}
 \boxed{
 H_{\rm YM}(a,\Pi)
 =\frac{c_{\rm op}}2\int_\Sigma\left[
   \|\Pi\|_{Q_\rho^{-1},h}^{2}
  +\|B_a\|_{Q_\rho,h}^{2}
 \right]d\mu_h .}
\end{equation}
Its constrained phase-space action is
\begin{equation}
\label{p3f:eq:v174-generated-joint-yang-mills-phase-action}
 S_{\rm ph}[a,\Pi,A_0]
 =\int d\widetilde t\left[
   \int_\Sigma\langle\Pi\wedge\dot a\rangle
  -H_{\rm YM}(a,\Pi)
  +\int_\Sigma\langle A_0,d_a\Pi\rangle
 \right].
\end{equation}
Its algebraic momentum equation is
\begin{equation}
\label{p3f:eq:v174-yang-mills-legendre-relation}
 \Pi=*_hQ_\rho^\flat(\mathcal E_A).
\end{equation}
Eliminating \(\Pi\) gives the full nonlinear local Lorentzian Yang--Mills
action
\begin{equation}
\label{p3f:eq:v174-generated-local-lorentzian-yang-mills-action}
 \boxed{
 S_{{\rm YM},L}[A]
 =-\frac14\int_U\sqrt{-g_{\rm op}}\,
 Q_\rho(F_{\mu\nu},F^{\mu\nu})\,d^4x .}
\end{equation}
The cubic and quartic nonabelian terms and the temporal Gauss coupling are
therefore restrictions of the same generated curvature
\(F_A=dA+A\wedge A\); they are not separate interaction selectors.
Variation gives
\[
 d_A\bigl(*_{g_{\rm op}}Q_\rho^\flat F_A\bigr)=0,
\]
whose normal component is \(d_a\Pi=0\).  The generated reduced phase-space
role is
\begin{equation}
\label{p3f:eq:v174-yang-mills-reduced-phase-space}
 \mathcal P_{\rm YM,\Sigma}^{\rm red}
 :=\mu_{\rm G}^{-1}(0)\mathbin{/\mkern-6mu/}\mathcal G_\Sigma,
\end{equation}
understood as the typed stratified quotient/groupoid at reducible
connections, and \(H_{\rm YM}\ge0\) descends to it.

The complete incidence package is retained familywise:
\begin{equation}
\label{p3f:eq:v174-yang-mills-package-corealization}
 \operatorname{CoReal}\!\left(
 S_{{\rm YM},L},\Theta_{\rm YM},\Omega_{\rm YM},\mu_{\rm G},
 H_{\rm YM},\mathcal P_{\rm YM,\Sigma}^{\rm red}
 \ \middle|\ \mathfrak D_J,\mathrm{CAUS},Q_\rho
 \right).
\end{equation}
Denote the theorem-status certificates by
\[
 \boxed{(\mathrm{JPHASE}_{\rm YM}^{\rm der})},
 \qquad
 \boxed{(\mathrm{GAUSS}_{\rm YM}^{\rm der})},
 \qquad
 \boxed{(\mathrm{ACT}_{{\rm YM},U}^{\rm der})}
 \equiv
 \boxed{(\mathrm{ACT}_{\rm YM,prod}^{\rm der})}.
\]
They discharge the local parity-even pure-gauge restriction of
\((\mathrm{ACT}_{\rm cl})\) on this theorem domain through the direct
Lorentzian-action branch \((\mathrm{LOR}_{\rm dir})\).  On the analytic
first-jet subdomain of
Lemma~\ref{p3f:lem:v174-yang-mills-complex-connection-real-form}, the same
action is also the oriented continuation child of the Euclidean curvature
action, so the gauge-sector restriction of \((\mathrm{CONT}_{\rm cl})\) is
derived there as well.
\end{theorem}

\begin{proof}
The clock product gives the displayed connection and curvature decomposition.
The regular evaluation dual of the generated affine connection role supplies
\eqref{p3f:eq:v174-regular-yang-mills-phase-space} and its tautological form.
For an infinitesimal gauge parameter \(\epsilon\), invariance of \(Q_\rho\)
and the stated boundary rule give
\[
 \iota_{\xi_\epsilon}\Omega_{\rm YM}
 =d\left[-\int_\Sigma
          \langle\epsilon,d_a\Pi\rangle\right],
\]
which proves the momentum-map statement, up to the declared overall sign
convention.

Pointwise metric naturality and gauge invariance leave, at differential order
two, only the parity-even Lorentzian curvature contraction
\(Q_\rho(F_A,*_{g_{\rm op}}F_A)\), up to one overall scalar; the other
independent four-dimensional quadratic density is the orientation-odd
\(Q_\rho(F_A,F_A)\), which is excluded by the theorem domain.  Restriction to
purely spatial curvature must reproduce the already generated ordinary-trace
magnetic term, fixing that one scalar.  The undeformed same-principal-owner
condition forbids a second electric constitutive coefficient, and positivity
of \(Q_\rho\) fixes the healthy simultaneous Hamiltonian sign.  Hence
\eqref{p3f:eq:v174-generated-joint-yang-mills-hamiltonian} contains no
independent electric coefficient, Lorentzian action member, or Hamiltonian
member.  On the analytic first-jet subdomain,
Lemma~\ref{p3f:lem:v174-yang-mills-complex-connection-real-form}
independently gives the same electric/magnetic coefficient and sign by
oriented real-form restriction.  The following principal-symbol theorem then
verifies that the resulting Euler operator has the operational cone; that
consequence is not used circularly in this construction.

Variation of
\eqref{p3f:eq:v174-generated-joint-yang-mills-phase-action} in \(A_0\) gives
\(d_a\Pi=0\).  Variation in \(\Pi\) gives
\[
 \dot a-d_aA_0=c_{\rm op}*_hQ_\rho^\sharp\Pi,
\]
equivalent to \eqref{p3f:eq:v174-yang-mills-legendre-relation}.  Substitution
yields
\[
 \frac{c_{\rm op}}2\int d\widetilde t\int_\Sigma
 \left[\|\mathcal E_A\|_{Q_\rho,h}^{2}
       -\|B_a\|_{Q_\rho,h}^{2}\right]d\mu_h.
\]
Using
\(d\mu_{g_{\rm op}}=c_{\rm op}d\widetilde t\,d\mu_h\) and the curvature
split proves
\eqref{p3f:eq:v174-generated-local-lorentzian-yang-mills-action}.  Since both
\(B_a\) and \(d_aA_0\) use the generated covariant derivative, the identity
contains the complete local nonabelian interaction.  Gauge invariance lets the
Hamiltonian descend through
\eqref{p3f:eq:v174-yang-mills-reduced-phase-space}, and positivity follows
from \(c_{\rm op}>0\) and the positive trace pairing.  Finally, Paper-II-A
rooted incidence retains every displayed object and map as one family-valued
constructor output, proving
\eqref{p3f:eq:v174-yang-mills-package-corealization} without selecting a field
member or history.
\end{proof}

\begin{corollary}[Gauge reduction and quadratic reciprocal-dynamics discharge]
\label{p3f:cor:v174-yang-mills-gauge-reduction-wick-discharge}
On the theorem domain above, reduction at \(\mu_{\rm G}=0\) supplies the
separately valid gauge reduction required by Paper~IV,
Theorem~\ref{p4:thm:stationary-quadratic-real-form-pair}.  Hence every
stationary stable finite physical-polarization block of the generated
Yang--Mills action inherits the already derived complex-time Euclidean/
Lorentzian quadratic action pair and reciprocal state/test dynamics without a
second gauge-reduction or Wick premise.  This finite quadratic consistency
corollary does not construct an actual state, path measure, interacting
continuum limit, or global reduced phase space.
\end{corollary}

\begin{remark}[Exact boundary of the Yang--Mills phase-space closure]
\label{p3f:rem:v174-yang-mills-phase-space-boundary}
Dropping parity restores the independent theta/\(E\!\cdot\!B\) branch and
shifts the canonical momentum.  Coupled matter replaces the source-free Gauss
law by a charge-current equation and requires the Lorentzian matter action.
The v1.75 native odd-scalar action now supplies that matter action on the
common local \(A+\Phi\) branch; Corollary~\ref{p3f:cor:v176-coupled-gauge-scalar-moment-map}
therefore derives its scalar matter current and coupled Gauss law.  Fermionic
matter current remains outside that corollary.
A boundary or noncompact slice outside the stated support rule requires
surface charges and edge data.  Time-dependent lapse/shift data require the
corresponding ADM phase-space form.  The theorem neither selects Paper~VII-A's
finite regulator action representative or finite Haar/transfer measure nor
proves its universality or continuum limit.  It also does not close the fermionic-reality and fermionic-current, exact
nonlocal or dimension-six-or-higher, constant odd/Majorana-offset,
noncanonical-boundary/global, actual solution, state, or history portions of
the complete action interface.  The local Einstein and scalar constituents
are derived separately below.
\end{remark}

\begin{theorem}[Joint gauge principal-symbol inheritance on the Lorentzian branch]
\label{p3f:thm:joint-gauge-principal-symbol-inheritance}
Use the derived \((\mathrm{JOINT}_x^{\rm der})\) and the complete upstream
\((\mathrm{CAUS})\) domain.  On the native stationary-product,
parity-even pure-gauge branch, use the derived
\((\mathrm{ACT}_{{\rm YM},U}^{\rm der})\) certificate of
Theorem~\ref{p3f:thm:v174-generated-joint-gauge-phase-space-action}.  On a
broader branch outside that theorem domain, retain a claim-specific Lorentzian
realization that preserves field identity and the two-derivative coefficient
and uses the same generated metric representative as the operational conformal
class, up to a positive conformal factor.  Use the smooth globally hyperbolic
product representative and standard Lorenz principal gauge on the evolution
domain.  Let \(E_a\) be any functorial gauge-factor component of the occurring
joint carrier that is nontrivial and whose internal two-derivative kinetic
form is nondegenerate.  On the parity-even positive trace branch the form is
the scalar representation index \(Z_a>0\) times the one common overall
constant \(C>0\), and the quadratic gauge action has the local Lorentzian
principal form
\[
 S^{(2)}_a
 =-\frac{CZ_a}{4}\int_U\!\sqrt{-g_{\rm op}}\,
 g_{\rm op}^{\mu\rho}g_{\rm op}^{\nu\sigma}
 F^a_{\mu\nu}F^a_{\rho\sigma}\,d^4x
\]
(up to the stated conformal representative).  Modulo the gauge direction, or
after a standard principal gauge fixing, the Euler--Lagrange principal symbol
is a nonzero internal factor times
\[
 (g_{\rm op}^{\alpha\beta}k_\alpha k_\beta)\delta^\nu{}_{\rho}
 -k^\nu k_\rho .
\]
Consequently every physical gauge polarization has characteristic set
\[
 \boxed{g_{\rm op}^{\mu\nu}k_\mu k_\nu=0}
\]
and therefore front speed
\[
 \boxed{c_{{\rm front},a}=c_{\rm op}.}
\]
An invertible internal change of gauge basis and any zeroth-order mass matrix
leave this characteristic set unchanged.  Only invertibility of the internal
kinetic form is used to identify the characteristic cone; positivity is the
separate healthy-energy condition supplied by the positive trace branch.
\end{theorem}

\begin{proof}
Theorem~\ref{p3f:thm:joint-connection-uninserted-trace-closure} supplies one
quadratic curvature trace and contracts the two-form indices with the generated
metric; expanding the joint trace supplies positive representation-index
factors on that branch.  More generally, an invertible internal kinetic form
multiplies the same spacetime symbol.  On the native stationary-product branch,
Theorem~\ref{p3f:thm:v174-generated-joint-gauge-phase-space-action} constructs
the Lorentzian action from the generated constrained Hamiltonian and transports
this two-derivative metric contraction to the operational metric owner without
a new coefficient.  On the broader branch the same sentence remains the
explicitly supplied Lorentzian-realization premise.  Either route discharges
the hypotheses of Paper~II-A,
Theorem~\ref{p2a:thm:undeformed-host-principal-inheritance}, for each physical
gauge polarization.  Variation of the Yang--Mills quadratic term gives the
displayed one-form principal symbol.  Its gauge-null direction is removed by
quotienting physical polarizations or by principal gauge fixing; the remaining
characteristic condition is precisely the metric null equation.  Multiplying
by \(C Z_a\), making an invertible transformation in internal factor space, or
adding a zeroth-order mass matrix cannot change the principal zero set.
In Lorenz gauge the remaining evolution operator is normally hyperbolic on the
stated globally hyperbolic product domain, so the standard finite-propagation
theorem identifies this characteristic cone with the wave-front boundary
\cite[Theorems~3.2.11--3.2.12]{BarGinouxPfaffle2007}.  Thus the analytic front-speed hypothesis of the
Paper~II-A theorem is discharged rather than assumed from common support.
\end{proof}

\begin{theorem}[Generated interior Einstein action from the joint Dirac owner]
\label{p3f:thm:v175-generated-joint-einstein-action}
Use the native common-heat branch of
Theorem~\ref{p3f:thm:common-heat-specializes-kinetic-trace}, the generated
joint Dirac family of
Theorem~\ref{p3f:thm:generated-fibre-spin-joint-operator}, and the complete
clocked \((\mathrm{CAUS})\) domain.  Work on a smooth four-dimensional
clock-adapted spin chart, retain the parity-even local heat-coefficient branch,
and use compactly supported interior variations in
\(U_0\Subset U\).  For the metric and connection part
\[
 D_A=c_g\circ\bigl(\nabla^{\mathbb S_g}\otimes1_F
                 +1\otimes\nabla_A^F\bigr),
\]
write
\[
 D_A^2=-\bigl(g_E^{\mu\nu}\nabla_\mu\nabla_\nu+E_A\bigr).
\]
In the Laplace-type convention of
Theorem~\ref{p3:thm:local_spectral_asymptotics}, the Lichnerowicz term is
\[
 E_A\big|_{R_E}=-\frac14R_E\,1.
\]
Hence the linear-curvature part of the joint-Dirac coefficient is negative,
\begin{equation}
\label{p3f:eq:v175-joint-dirac-einstein-a2}
 a_2(D_A^2)\big|_{R_E}
 =(4\pi)^{-2}\int_{U_0}
 \operatorname{tr}_{\mathbb S_g\otimes F}
 \left(E_A+\frac16R_E\right)d\mu_E
 =-(4\pi)^{-2}\frac{\operatorname{tr}(1)}{12}
   \int_{U_0}R_E\,d\mu_E .
\end{equation}
After the already derived positive support and Born-diagonal weights are
included, write its \emph{raw represented joint coefficient} as
\(\kappa_J>0\).  The generated Euclidean
interior action germ is therefore
\begin{equation}
\label{p3f:eq:v175-einstein-euclidean-interior-action}
 I_{E,U_0}^{\rm EH}=-\kappa_J\int_{U_0}R_E\,d\mu_E
 \qquad\text{modulo compact-support exact divergences.}
\end{equation}

On the additional analytic clock--ADM first-jet subdomain, require the lapse,
spatial metric, and their first clock derivatives to extend holomorphically to
a contour sector joining the Euclidean and oriented Lorentzian real slices,
with no extra contour-end term.  In a zero-shift clock chart use
\[
 \tau=i\widetilde t,
 \qquad d\tau=i\,d\widetilde t,
 \qquad K_{E,ij}=-iK_{L,ij}.
\]
Then the interior ADM representative obeys
\begin{equation}
\label{p3f:eq:v175-einstein-clock-real-form}
\begin{aligned}
 I_{E,U_0}^{\rm EH}&\longmapsto-iS_{L,U_0}^{\rm EH},\\
 S_{L,U_0}^{\rm EH}&=\kappa_J\int_{U_0}R_L\,d\mu_L\\
 &\quad\text{modulo the same transported exact divergences.}
\end{aligned}
\end{equation}
Thus the conventional positive Einstein--Hilbert action orientation is a
generated local interior real-form descendant of the joint Dirac owner.  Denote
this certificate by
\[
 \boxed{(\mathrm{ACT}_{{\rm J,EH},U_0}^{\rm int,der})}.
\]
Together with the independently derived v1.74 Yang--Mills certificate, the
intersection of their printed domains carries the co-realized package
\[
 \boxed{(\mathrm{ACT}_{{\rm J,EH+YM},U_0}^{\rm loc,der})}
 :=(\mathrm{ACT}_{{\rm J,EH},U_0}^{\rm int,der})
  +(\mathrm{ACT}_{{\rm YM},U}^{\rm der}).
\]
This is a common-root action package over the same clocked metric, joint
connection, and ordinary-trace ancestry.  It does not by itself calibrate
Newton's constant, assert a relative empirical matching scale, or supply a
global boundary action.  The positive \(+R_E/6\) coefficient of
Paper~III-E's unit-preserving scalar Laplacian remains valid in its
scalar-resolution/readout sort; it is not the sign-bearing coefficient of this
joint-Dirac action owner.
\end{theorem}

\begin{proof}
The four-dimensional coefficient formula is
\(
 a_2=(4\pi)^{-2}\int\operatorname{tr}(E+R/6)
\)
on the stated smooth local domain.  Substituting the Lichnerowicz value
\(E_A|_R=-R/4\) gives
\eqref{p3f:eq:v175-joint-dirac-einstein-a2}.  The internal trace, protected
support weight, and Born-diagonal pure-metric pullback are positive and
therefore cannot change its sign; they are absorbed into the positive generated
coefficient \(\kappa_J\).

Modulo the exact divergence removed by the compact-support interior rule, the
Euclidean ADM integrand is
\[
 -\kappa_JN_E\sqrt h
 \left({}^{(3)}R+K_E^2-K_{E,ij}K_E^{ij}\right).
\]
The stated holomorphic first-jet continuation sends it to \(-i\) times
\[
 \kappa_JN_L\sqrt h
 \left({}^{(3)}R+K_{L,ij}K_L^{ij}-K_L^2\right),
\]
which is the Lorentzian Einstein--Hilbert interior action germ in this paper's
conventions.  The v1.74 constrained phase-space theorem independently derives
the healthy Yang--Mills action from the same clocked metric, joint connection,
and positive trace pairing; on the analytic first-jet overlap its connection
real form agrees with this oriented clock continuation.  Paper~II-A's rooted
image theorem therefore retains both restrictions as one generated co-realized
package.  No exact spectral-trace continuation, Newton calibration, physical
boundary class, or actual metric/connection member is used.
\end{proof}

\begin{corollary}[Born--causal Newton coordinate of the generated joint-EH owner]
\label{p3f:cor:v267-joint-eh-newton-readout}
For the generated raw represented coefficient \(\kappa_J>0\) of the
joint-Dirac Einstein real-form theorem, define
\[
 \UEHjoint^{\rm raw}:=8\kappa_J,\qquad \hBorn:=2\pi\hbar.
\]
For backward notation compatibility within this paper, an unqualified
\(\UEHjoint\) denotes this raw represented coefficient.  It is not the
primitive momentum-curvature owner \(\mathcal K_p=6/\Rmom^2,\qquad U=(3/\pi)\mathcal K_p\quad\text{only under }(\mathrm{CELL}_{K\to U})\).
Paper~III-A's three-dimensional causal determinant is
\(J_c=c_{\rm op}^3\).  The associated Newton coupling coordinate is
\begin{equation}
\label{p3f:eq:v267-joint-eh-newton-readout}
 \boxed{G_{J,\mathrm{raw}}:=\frac{J_c}{\hBorn\UEHjoint^{\rm raw}}
 =\frac{c_{\rm op}^3}{\hBorn\UEHjoint^{\rm raw}}.}
\end{equation}
Equivalently,
\[
 \boxed{\hbar\kappa_J=\frac{c_{\rm op}^3}{16\pi G_{J,\mathrm{raw}}}.}
\]
No observed Newton constant enters this definition.  Identifying this raw
action coordinate with a measured coupling first requires the actual
physical action/source projection of
Proposition~\ref{p3f:prop:v2211-newton-owner-matching}.  Numerical
comparison is metrology only after that physical normalization is fixed.
Paper~V-A later closes the source-native common-heat carrier and proves the
exact coefficient transport
\[
 \UEHjoint^{\rm raw}=96\,U.
\]
Thus the former equality bridge
\((\mathrm{MATCH}_{J/\mathrm{scal}}^{\rm EH}):\UEHjoint=U\) is retired on
that branch.  A claim that identifies a physical Newton/Planck owner with the
raw represented coefficient, or with a primitive descendant of it, remains a
separately typed owner statement.
\end{corollary}
\begin{proof}
Substitute \(\UEHjoint^{\rm raw}=8\kappa_J\), \(\hBorn=2\pi\hbar\), and
\(J_c=c_{\rm op}^3\) into the definition and rearrange.  No cross-owner
coefficient identification is used.
\end{proof}

\begin{proposition}[Distinct Newton coordinates and the actual matching test]
\label{p3f:prop:v2211-newton-owner-matching}
Retain the exact common-heat relation \(U_J^{\rm raw}=96U\), common
clock units and \(\hBorn=2\pi\hbar\).  Define the scalar-response
coordinate separately by
\[
 G_{\rm scal}:=\frac{c_{\rm op}^3}{\hBorn U}.
\]
Then
\[
 \boxed{G_{J,\rm raw}=\frac{G_{\rm scal}}{96},\qquad
 \hbar\kappa_J=12\hbar U,\qquad
 \frac{c_{\rm op}^3}{16\pi G_{\rm scal}}=\frac{\hbar U}{8}.}
\]
These are unequal coefficient coordinates, not two measured values of
one already matched Newton coupling.

For an actual claim-complete projection/reduction \(\mathcal P\) of the
joint action, let \(c_{\rm EH}^{\rm phys}\) be the coefficient of
\(\int\sqrt{-g}R\) in its physical low-derivative representative, with
the actual metric, matter/source, clock and boundary maps retained.
Matching the scalar-response action requires, in particular,
\[
 c_{\rm EH}^{\rm phys}=\hbar U/8
                    =(\hbar\kappa_J)/96.
\]
This coefficient test is necessary, not sufficient: the normalized
source/stress and consumed mixed metric, field and boundary jets must
also agree.  No division of the joint action or of one constituent by
\(96\) is licensed merely by knowing its raw trace coefficient.
\end{proposition}
\begin{proof}
Substitute \(U_J^{\rm raw}=96U\) in the raw-coordinate definition;
the displayed coefficient relations follow by rearrangement.
The physical Einstein normalization
\(c_{\rm EH}^{\rm phys}=c_{\rm op}^3/(16\pi G_{\rm phys})\) identifies
\(G_{\rm phys}=G_{\rm scal}\) only when its actual projected coefficient
equals \(\hbar U/8\).  The Hilbert source and its mixed responses are
derivatives of the same projected action.  Equality of one coefficient
does not identify those other derivatives, while a complete action
identity with the retained maps identifies them by differentiation.
Hence raw trace counting supplies the ratio, not the missing projection.
\end{proof}

\begin{definition}[Physical action at a specified reference]
\label{p3f:def:v3041-adopted-metric-assembly}
On a common smooth metric, matter, source and boundary chart, write
\(S_J(g,X;b)\) for the ordinary joint action with its original Born
allocations and clock. Here \(b\) denotes declared boundary and source
data held fixed when forming an interior variation. Specify a reference
section \(i_v(g,b)=(g,X_v(g,b);b)\) independently of the force being
predicted. The label \(v\) and the prescription for the section are fixed.
Let \(M_{\rm rp}(g;b)\) be the generated marked geometric action on its
licensed product and clock domain, with the endpoint appropriate to the
same variational problem. For this reference the adopted physical action is
\begin{equation}
 \label{p3f:eq:v3041-physical-reference-action}
 M_{J,v}=i_v^*S_J,\qquad
 \boxed{S_{\rm phys}^{(v)}(g,X;b)
    =M_{\rm rp}(g;b)+S_J(g,X;b)-M_{J,v}(g;b).}
\end{equation}
This is a choice of physical action. The ordinary Einstein action and
source theorem, Theorems~\ref{p3f:thm:v175-generated-joint-einstein-action}
and~\ref{p3:thm:common-position-spectral-origin}, and the two coordinates
in Proposition~\ref{p3f:prop:v2211-newton-owner-matching} continue to
describe \(S_J\) and its comparison. The marked trace is the one in
Paper~I-C, Proposition~\ref{p1b:prop:v224-retained-primitive}.
On a pure geometric product \(B_g\otimes I_F\), tracing
\(P_+/48\) gives \(1/96\) of the ordinary trace since the two internal
grading spaces have equal rank and \(\dim F=48\). This fixes the marked
Einstein coefficient \(\nu=U/8\) in natural units. A reference carrying
scalar energy, gauge curvature or boundary work contributes its full
\(M_{J,v}\) in the subtraction. Setting the added gauge potential to zero
can leave the curved Born connection at the reference.
\end{definition}

\begin{proposition}[Reference centering and calibrated tensor exchange]
\label{p3f:prop:v3041-reference-exchange}
On a chart with connected nonmetric fibres,
Eq.~\eqref{p3f:eq:v3041-physical-reference-action} is the unique action
whose restriction to \(i_v\) is \(M_{\rm rp}\) and whose fixed-\((g,b)\)
differential in every nonmetric direction equals \(d_XS_J\).
Every existing mixed derivative with at least one independently varied
\(X\) leg agrees as well. For a canonical scalar \(h\) with
\(F_J(h)=\kappa_J-h^2/12\) and potential \(V(h)\), choose the constant
reference \(h=v\) with zero field strengths and all other material work
absent there. If the marked local metric density is
\(\nu R-V_{\rm rp}\) at two derivatives, the adopted density in this sector is
\[
 \left[\nu-\frac{h^2-v^2}{12}\right]R
 -\tfrac12(\nabla h)^2-\bigl[V_{\rm rp}+V(h)-V(v)\bigr].
\]
Here \(V_{\rm rp}\) is the zero-order contribution of the declared
geometric action on this domain, including any established vacuum
completion. The full coefficient \(F_v(h)=\nu-(h^2-v^2)/12\) multiplies
the canonical non-null Einstein boundary term with relative factor two
when that boundary problem is admitted. A constant reference which is
also a solution requires \(R_0=2V_{\rm rp}/\nu\) and
\(V'(v)=-vR_0/6\) on this source-free scalar branch.

In particular consider a stationary flat reference in the admitted
interior (or with compatible endpoints), with \(F_v(v)=\nu>0\).
Assume a single Einstein tensor pole with its regular normalized Hilbert
vertex, in a common metric and clock. After all physical source,
boundary and constraint reductions, require that additional fields,
state and measure supply no further leading tensor stiffness or mixed
pole and no independent singular source contact. For canonical conserved
probes and spacelike transfer \(q=k_E^2\downarrow0\), the tensor part of
the connected exchange has, up to the common propagator sign,
\begin{equation}
 \label{p3f:eq:v3041-calibrated-pole}
 \lim_{q\downarrow0}q\,\mathcal A^{\rm tensor}_{12}(q)=
 \frac{T_1^{\mu\nu}T_{2,\mu\nu}-\tfrac12
       T_{1\,\mu}^{\ \mu}T_{2\,\nu}^{\ \nu}}{2\nu},\qquad
 G_{\rm tensor}=\frac{c_{\rm op}^3}{16\pi\hbar\nu}=G_{\rm scal}.
\end{equation}
Other long-range exchanges must be evaluated or separated for a
total-force claim. Existence and stability of the required history and
an identification with an observed Newton experiment have their own
matter, source and boundary requirements.
\end{proposition}
\begin{proof}
At fixed \((g,b)\) the difference of two actions with the specified
nonmetric differential is constant on the connected fibre. Evaluation
on the section fixes this function to \(M_{\rm rp}-M_{J,v}\) for the
difference from \(S_J\). Its nonmetric derivatives vanish. Subtracting
the scalar density at constant \(v\) gives the displayed \(F_v\) and
potential. Variation before imposing \(h=v\) gives the two stationary
conditions. Equivalently, the ordinary scalar and gauge Euler--Noether
identities of Theorem~\ref{p3f:thm:v176-local-bosonic-euler-noether}
remain the identities of the relative action, with the change to the
pure metric Euler derivative obtained by varying \(M_{\rm rp}-M_{J,v}\).
In the canonical strain convention the source vertex is
\(h_{\mu\nu}T^{\mu\nu}/2\). Inverting the linearized Einstein equation
\(2\nu G_{\mu\nu}^{(1)}=T_{\mu\nu}\) on conserved sources gives
Eq.~\eqref{p3f:eq:v3041-calibrated-pole}. In a scalar sector with
\(F'_v(v)\ne0\) the trace couples to the scalar as well. A regular
massive scalar adds a finite-range exchange, while a massless scalar
requires its own residue. Neither is included in the displayed tensor part.
\end{proof}

\begin{remark}[Reference work and scope of the action selection]
\label{p3f:rem:v3041-reference-work}
Metric and boundary derivatives of \(M_{J,v}\) are derivatives of the
pullback, including those of \(X_v(g,b)\), the domain and density.
If \(v\) itself varies with an applied source, a selected state or a
dynamical field, differentiating the resulting composition adds the
corresponding first and second contacts. Fixed-\(v\) identities alone
do not evaluate that different response. Likewise, marking the whole
pure metric word on a product preserves its relative higher-order and
vacuum coefficients on the channels where that word and its marked
continuation are established. An isolated substitution in the Einstein
coefficient is a different action. The existing conditional scalar
reference comparison in Paper~III-G, Section~\ref{p3f:app:v2216-reference-centered-matching},
gives the same two-derivative centering in its domain. At a fixed external
metric and boundary, suppose the material source \(J\) leaves the
reference subtraction and marked action independent of \(J\), and the
same material integration domain, measure and regulator define both
functionals. In dimensional action units the conditional identity is
\[
 Z_{\rm phys}^{(v)}[g,J]
 =\exp\!\left(\frac{i}{\hbar}[M_{\rm rp}(g;b)-M_{J,v}(g;b)]\right)
     Z_J[g,J],
\]
where the subscript \(J\) on the right designates the ordinary joint
action. The factor is independent of the integrated material fields and
of the probe. When \(Z_J[g,0]\ne0\), it cancels from the normalized
generating functional. Every connected matter correlator obtained by
differentiating its logarithm therefore agrees on the fixed background,
at every order where the functional or perturbation series is defined.
Integrating over \(g\) retains this metric-dependent factor inside the
integral, and probe-dependent reference or measure factors require their
own derivatives. This conditional identity supplies neither a quantum
measure nor protection of the physical graviton pole. Quantizing first can
introduce determinants, state work and measure derivatives. The selection
in Definition~\ref{p3f:def:v3041-adopted-metric-assembly} supplies no
global quantum measure or stationary vacuum, and by itself establishes
no equality between reductions performed in different orders.
\end{remark}
% BEGIN proposed v3.071: reference-independence of fixed-background off-diagonal matter currents.
\begin{corollary}[Exact reference-independence of nonmetric transition vertices]
\label{p3f:cor:v371-reference-independent-transition-current}
On the adopted action family of
Definition~\ref{p3f:def:v3041-adopted-metric-assembly}, fix the metric
$g$ and a common matter Hilbert/source chart. Let $b$ be any declared
external source/boundary datum and let $i_v(g,b)$ be an admissible
reference section that is prescribed independently of the fluctuating
matter variable $X$. Write
\[
 C_v(g;b):=M_{\rm rp}(g;b)-S_J(g,X_v(g,b);b),\qquad
 S_{\rm phys}^{(v)}(g,X;b)=S_J(g,X;b)+C_v(g;b).
\]
Whenever the indicated derivatives exist, every vertex carrying at
least one independently varied nonmetric matter leg is exactly
reference-independent:
\[
 \boxed{D_X^mD_b^n S_{\rm phys}^{(v)}
        =D_X^mD_b^nS_J\quad(m\ge1,\ n\ge0),}
\]
with the metric and the prescribed reference section held to their
declared meanings. This holds for every reference $v$, with no
stationarity, Hessian-invertibility or uniqueness requirement.
In particular the $b$-current insertion on the matter space can
change only by a scalar multiple of the identity,
\[
 J_b^{(v)}=J_b^{J}+(D_b C_v)\,I_{\rm matter},
 \qquad
 \boxed{\langle B|J_b^{(v)}|D\rangle
        =\langle B|J_b^J|D\rangle\quad(B\perp D).}
\]
For a common, $v$-independent fixed-background matter measure and
integration domain, the normalized generating functional of arbitrary
independent matter probes $j$ obeys
\[
 \frac{Z_{\rm phys}^{(v)}[g,b;j]}{Z_{\rm phys}^{(v)}[g,b;0]}
 =\frac{Z_J[g,b;j]}{Z_J[g,b;0]},
\]
provided the denominators are nonzero and $j$ is not included in the
reference prescription. Thus normalized fixed-background matter
correlators and their off-diagonal transition matrix elements do
\emph{not} inherit the freedom in $v$. They still require the actual
$S_J$ current, physical pole states, normalization and open channels.
This conclusion does not extend automatically to metric integration,
source-dependent measures, independent probe dependence of the chosen
reference, or metric-mediated contributions to a full observable width.
\end{corollary}
\begin{proof}
The reference term $C_v(g;b)$ is independent of $X$ by definition.
One $X$ derivative annihilates it; all further derivatives commute
where defined. Variation with respect to an external probe $b$ yields
$D_bS_J+D_bC_v$, the second term acting as a c-number on the common
matter Hilbert space. Its matrix element between orthogonal states
vanishes exactly. For fixed $g,b$ the path-integral weight is multiplied
by the field-independent phase $\exp(iC_v/\hbar)$, which cancels from the
ratio of generating functionals. No assertion about the full metric
path integral is used.
\end{proof}
% END proposed v3.071 current-source reference closure.

\begin{theorem}[Local on-shell reference section and exact classical Hessian cancellation]
\label{p3f:thm:v3042-onshell-reference-schur}
Work on a fixed regular gauge/constraint/boundary quotient in a Banach chart
of metrics $g$ and material variables $X$. Let the ordinary same-owner action
$S_J(g,X;b)$ be $C^3$ near $(g_0,X_0)$ at fixed external data $b$.
Suppose $D_XS_J(g_0,X_0;b)=0$ and
$H=D_X^2S_J(g_0,X_0;b)$ is a bounded isomorphism on the retained
physical material tangent. Then the implicit-function theorem produces a
unique local $C^2$ stationary section $v(g)$ with $v(g_0)=X_0$ and
$D_XS_J(g,v(g);b)=0$. In a covariant chart, uniqueness makes the section
equivariant under symmetries that preserve the seed and the quotient.

For the \emph{adopted} action
\[
 S_{\rm phys}^{(v)}(g,X;b)=M_{\rm rp}(g;b)
 +S_J(g,X;b)-S_J(g,v(g);b),
\]
write $Y=X-v(g)$. Then, for all $g$ in the stationary neighbourhood,
\[
 S_{\rm phys}^{(v)}(g,Y=0;b)=M_{\rm rp}(g;b),\qquad
 D_Y S_{\rm phys}^{(v)}(g,0;b)=0,\qquad
 D_gD_Y S_{\rm phys}^{(v)}(g,0;b)=0.
\]
At $(g_0,0)$ the classical Hessian is block diagonal with entries
$D_g^2M_{\rm rp}$ and $H$. Equivalently the stationary metric Schur
complement calculated at fixed original $X$ is exactly
\[
 \boxed{\quad H_{gg}^{\rm phys}
  -H_{gX}^{\rm phys}H^{-1}H_{Xg}^{\rm phys}
  =D_g^2M_{\rm rp}.\quad}
\]
For any two independent linear probe insertion vectors $(j_g,j_X)$ held
external to the choice of reference section, the exact Gaussian tree
cross response is
\[
 \boxed{\mathcal A_{12}=
 (j_{g,1}+V^*j_{X,1})^*(D_g^2M_{\rm rp})^{-1}
 (j_{g,2}+V^*j_{X,2})+j_{X,1}^*H^{-1}j_{X,2}.}
\]
The induced metric source $j_g+V^*j_X$ is fixed by the same action;
when $H^{-1}$ and this induced source are infrared regular, they do not
supply an additional massless tensor pole. Other independent massless
exchanges are not excluded by this formula.
Thus a separate \emph{classical} mixed-Hessian cancellation assumption is
unnecessary on this regular on-shell branch. This does not prove existence
of $(g_0,X_0)$, positivity of $H$, a mass gap, source-free Lorentzian
propagation, quantum determinant cancellation or physical action occurrence.
\end{theorem}
\begin{proof}
The implicit-function theorem applies to $D_XS_J=0$, since its $X$
derivative at the seed is the isomorphism $H$. Define
$\widehat S_J(g,Y)=S_J(g,v(g)+Y)$. Then
$S_{\rm phys}^{(v)}=M_{\rm rp}+\widehat S_J(g,Y)-\widehat S_J(g,0)$,
and $D_Y\widehat S_J(g,0)=0$ for each $g$. Differentiation gives the
block diagonal Hessian in $(g,Y)$ coordinates. At the seed, write
$V=D_gv$ and differentiate $D_XS_J(g,v(g))=0$ to find
$S_{J,Xg}=-HV$. The Hessian in the original $(g,X)$ coordinates has
$H_{gg}^{\rm phys}=M_{\rm rp,gg}+V^*HV$,
$H_{gX}^{\rm phys}=-V^*H$, $H_{XX}^{\rm phys}=H$.
Its Schur complement is $M_{\rm rp,gg}$ exactly.
\end{proof}

\begin{proposition}[The independent canonical fermion source has fixed gravitational charge]
\label{p3f:prop:v3042-canonical-dirac-hilbert-source}
On the local Lorentzian Dirac branch where the same action has the canonical
vierbein/spin connection coupling and a physical isolated one-particle
fermion pole, let the normalized one-particle state obey
$P_\mu|p\rangle=p_\mu|p\rangle$ and the standard relativistic
state normalization. Its on-shell Hilbert tensor is obtained by
varying $S_J$ with respect to $g$; the reference subtraction contains no
independent fermionic fluctuation. Conservation/translation charge and
canonical normalization give the forward matrix element
\[
 \boxed{\langle p|T_{\mu\nu}(0)|p\rangle=2p_\mu p_\nu}
\]
with the conventional stripping of momentum delta functions and on-shell
spin labels. In particular, for a canonical free Dirac wavefunction
$\bar u(p)\gamma_\mu u(p)=2p_\mu$, the Belinfante--Hilbert tensor
$T_{\mu\nu}=\frac i4\bar\psi(\gamma_\mu
\overleftrightarrow D_\nu+\gamma_\nu\overleftrightarrow D_\mu)\psi$
has that forward matrix element. Accordingly no extra freely chosen
\emph{relative} gravitational vertex weight exists for such probes.
This proposition does not construct the particle poles, their masses,
or their occupied state on the Einstein reference.
\end{proposition}

\begin{proposition}[Closed three-sphere mean-zero inverse and local Newton singularity]
\label{p3f:prop:v3042-S3-Newton-local-parametrix}
On the round position $S^3(R)$, for geodesic angle $0<\chi<\pi$,
\[
 G_R(\chi)=\frac{(\pi-\chi)\cot\chi-\tfrac12}{4\pi^2R}
\]
is the unique mean-zero scalar Green distribution solving
$-\Delta_{S^3(R)}G_R=\delta-1/(2\pi^2R^3)$.
For $r=R\chi\downarrow0$ it obeys
\[
 G_R=\frac1{4\pi r}-\frac3{8\pi^2R}+O(r/R^2),\qquad
 \frac{-4\pi R^2\chi^2}{R}\,\partial_\chi G_R
 =1+\frac{\chi^2}3+O(\chi^3).
\]
Hence any licensed static Poisson reduction with Einstein coefficient
$\nu=U/8$ has local two-probe potential coefficient
$1/(16\pi\nu)=1/(2\pi U)$ without requiring a globally noncompact
flat slice. The equality with the complete time-dependent gravitational
constraint, including its lower-order curvature and matter terms,
is \emph{not} inferred from the scalar Green function alone.
\end{proposition}

\begin{remark}[Unresolved physical gates]
\label{p3f:rem:v3042-onshell-reference-gates}
The on-shell selection theorem does not manufacture its initial stationary
member. On the native exact positive parallel scalar heat branch of
Theorem~\ref{p3f:thm:v250-exact-parallel-scalar-heat}, a nonzero parallel
scalar has strictly nonzero radial derivative, which cannot be cancelled
by a purely metric reference subtraction because its $X$ derivative
vanishes. The v3.041 adoption of the marked action is a model decision,
not an output of G0--GP3. Even on a stationary classical section,
integrating Gaussian material fluctuations produces a generally
metric-dependent $\frac\hbar2\operatorname{Tr}\log H_{XX}(g)$; the
classical Schur identity gives no all-loop protection of the Einstein pole.
\end{remark}


\begin{corollary}[All-order stationary-branch metric descent]
\label{p3f:cor:v3043-stationary-allorder-descent}
On the regular local stationary branch of
Theorem~\ref{p3f:thm:v3042-onshell-reference-schur}, consider the
\emph{adopted} action of
Definition~\ref{p3f:def:v3041-adopted-metric-assembly}.
At any finite derivative order $r$ for which the stationary section,
ordinary action and marked metric action have the requisite smoothness,
let $Y=X-v(g)$ and
$\widehat S_J(g,Y;b)=S_J(g,v(g)+Y;b)$ at fixed external data $b$.
Then
\[
 S_{\rm phys}^{(v)}(g,Y;b)=M_{\rm rp}(g;b)
    +\widehat S_J(g,Y;b)-\widehat S_J(g,0;b),
 \qquad D_Y\widehat S_J(g,0;b)=0.
\]
Consequently, on this branch the exact classical metric-only reduced action
and each of its defined derivatives are
\[
 \boxed{W_{\rm cl}^{(v)}(g;b)
    :=S_{\rm phys}^{(v)}(g,v(g);b)=M_{\rm rp}(g;b),
 \qquad D_g^nW_{\rm cl}^{(v)}=D_g^nM_{\rm rp}\quad(0\le n\le r).}
\]
In the translated $(g,Y)$ chart one also has
$D_g^pD_YS_{\rm phys}^{(v)}(g,0;b)=0$ whenever defined.
This does \emph{not} assert vanishing of
$D_g^pD_XS_{\rm phys}^{(v)}$ with the \emph{original} $X$ held
fixed; its mixed Hessian is generally nonzero and is exactly the one
accounted for by the Schur identity of the cited theorem.
\end{corollary}
\begin{proof}
The first equation is a reparametrization of
Eq.~\eqref{p3f:eq:v3041-physical-reference-action}.
The stationary equation $D_XS_J(g,v(g);b)=0$ is
$D_Y\widehat S_J(g,0;b)=0$.  Setting $Y=0$ cancels both ordinary
terms \emph{as functions of $g$}, so differentiation proves every
claimed metric-only derivative identity.  Differentiating the
identically zero first material variation at fixed $Y$ proves the
translated-chart mixed identities.  In the original chart the
stationary-section derivative $D_gv$ produces the nonzero
mixed Hessian already stated in the source theorem.
\end{proof}

\begin{corollary}[Material-jet invariance under pure metric reference work]
\label{p3f:cor:v3043-fixed-material-jet-invariance}
On the same fixed $(g,X;b)$ chart, at all existing derivative orders,
\[
 \boxed{D_g^pD_X^qS_{\rm phys}^{(v)}(g,X;b)
       =D_g^pD_X^qS_J(g,X;b),\qquad q\ge1,}
\]
where the metric derivatives hold the \emph{original} $X$ fixed.
Thus the same fixed-background, same-state and same-boundary normalized
material kinetic and chiral mixed source Hessians are unchanged by
reference centering.  This includes the inputs $Z_L,Z_R,D_{\rm ch}$ to
Proposition~\ref{p3f:prop:v239-same-action-electron-pencil}; the identity
\emph{does not} evaluate those inputs or establish
$\mathsf{Cert}_{\rm hier}$ of
Definition~\ref{p5:def:v228-hierarchy-origin}.  It does not assert
that subsequent metric integration or changes of state preserve poles.
\end{corollary}
\begin{proof}
The difference from $S_J$ is the function
$M_{\rm rp}(g;b)-S_J(g,v(g);b)$, independent of the original $X$.
Every derivative containing at least one $X$ leg annihilates it.
\end{proof}

\begin{proposition}[Stationary-reference component invariance]
\label{p3f:prop:v3043-stationary-component}
Fix $g,b$. Along each piecewise $C^1$ path $t\mapsto v_t$ of
material stationary points of $S_J(g,\cdot;b)$, the stationary action
value is constant.  Therefore two stationary sections whose values
lie in the same such path component for each $g$ define the same
reference subtraction and the same centered action.  Different
regular components can, in general, give distinct metric-dependent
stationary pullbacks; local regularity does not choose between them.
\end{proposition}
\begin{proof}
On every smooth part of the path,
$\frac{d}{dt}S_J(g,v_t;b)=
\langle D_XS_J(g,v_t;b),\dot v_t\rangle=0$.
Continuity extends equality across finitely many corners.
For the last warning, the elementary smooth action
$S_J(g,x)=(x^2-1)^2/4+\epsilon gx$
has two regular stationary seeds $x=\pm1$ at $g=0$; their
stationary pullbacks have first metric derivatives $\pm\epsilon$.
This is a mathematical branch-selection example, not a new Born--Hopf
physical endpoint.
\end{proof}

\begin{proposition}[Fixed-geometry matter probes do not select a metric action]
\label{p3f:prop:v3043-matter-metric-action-discriminator}
On a common regular chart suppose $S_2(g,X;b)=S_1(g,X;b)+C(g;b)$
with the same matter integration domain, measure and regulator,
where $C$ depends on neither the integrated matter fields nor the
independent matter probe.  The normalized fixed-geometry matter
correlators of these actions coincide whenever defined (the
existing identity is recorded in
Remark~\ref{p3f:rem:v3041-reference-work}).
For any common stationary section $v(g)$, however,
\[
 W_2(g;b)-W_1(g;b)=C(g;b),\qquad
 D_g^nW_2-D_g^nW_1=D_g^nC
\]
for every defined order.  Thus fixed-geometry matter data alone
cannot determine the metric-only action modulo the addition of $C$.
Under the \emph{separate} common-Hilbert-vertex and single regular
tensor-pole hypotheses of Propositions
\ref{p3f:prop:v3041-reference-exchange} and
\ref{p3f:prop:v3036-gapped-schur-pole}, the printed coefficient
coordinates $12U$ (ordinary regular-gapped tensor packet) and
$U/8$ (marked stationary tensor packet) give the conditional
ratio of their leading $1/k^2$ tensor residues
\[
 \boxed{\mathcal R_{\rm marked}/\mathcal R_{\rm ordinary}=96.}
\]
This ratio compares the specified \emph{mathematical} tensor packets;
it is neither two measured gravitational couplings nor proof that
both packets are actual admissible GP3 physical endpoints.
\end{proposition}
\begin{proof}
At fixed geometry, the factor $\exp(iC/\hbar)$ multiplies the
un-normalized material generating functional but cancels from its
normalized ratio, exactly as in the cited existing remark.  On the
same stationary section $S_2-S_1=C$, and metric differentiation
proves the second identity.  In the stated tensor comparison domain
the residue scales inversely with the kinetic coefficient at a fixed
canonical source vertex, so $(12U)/(U/8)=96$.
Without the stated common-vertex and pole conditions no physical
residue ratio follows from the coefficient comparison alone.
\end{proof}

\begin{remark}[Physical selection remains a separate occurrence problem]
\label{p3f:rem:v3043-reference-occurrence-scope}
The all-order cancellation is an immediate formal consequence of
the \emph{adopted} action and regular stationary section; it adds no
independent action-selection theorem to G0--GP3.  Pure metric reference
work cannot supply the source-owned chiral hierarchy arrows, while
fixed-background matter correlators cannot distinguish metric action
representatives.  An unconditional observed Newton response requires
physical action incidence or a proved residue-invariance theorem on
the actual admissible owner family, an occurring regular seed, and
its normalized source/measure/measurement bridges.  The charged
preterminal-source and interacting-pole obligations remain separate.
\end{remark}


\begin{remark}[Conditional canonical Dirac boundary completion]
\label{p3f:rem:v175-canonical-dirac-boundary-completion}
If a smooth non-null boundary is separately constituted, let
\[
 \chi:=\gamma_n\gamma_g,
 \qquad \Pi_\pm:=\frac12(1\pm\chi),
\]
and impose the local self-adjoint Dirac domain
\[
 \Pi_-\psi|_{\partial U}=0,
 \qquad
 \Pi_+\bigl(\nabla_n-\tfrac12K_E\bigr)\Pi_+\psi|_{\partial U}=0.
\]
For the rank-four complex spin module the boundary-inclusive coefficient is
\[
 a_2(D_A^2,\chi)
 =(4\pi)^{-2}\frac16\left[
 \int_U\operatorname{tr}(6E_A+R_E)
 +\int_{\partial U}\operatorname{tr}(2K_E+12S)
 \right],
 \qquad S=-\frac12K_E\Pi_+.
\]
Its bulk and boundary curvature coefficients are respectively
\(-1/3\) and \(-2/3\), before the common positive internal/support factor.
Thus the same owner gives
\[
\begin{aligned}
 I_{E,U}^{\rm EH+GHY}
 &=-\kappa_J\left(\int_U R_E\,d\mu_E+2\int_{\partial U}K_E\,d\mu_{\partial,E}\right)\\
 &\longmapsto-i\kappa_J\left(\int_U R_L\,d\mu_L+2\int_{\partial U}K_L\,d\mu_{\partial,L}\right).
\end{aligned}
\]
Accordingly, on the union of the compact-support interior domain and this
canonical local boundary domain, write
\[
 \boxed{(\mathrm{ACT}_{{\rm J,EH},U}^{\rm der})}.
\]
The boundary extension is conditional on the printed boundary constitution;
it is not part of the unrestricted interior theorem.  Null, corner,
asymptotic, mixed, and otherwise noncanonical boundary classes remain open.
\end{remark}

\begin{theorem}[Native local scalar action from the generated odd-field family]
\label{p3f:thm:v175-native-local-scalar-action}
On the same native common-heat local domain, let
\[
 D_{A,\Phi}=D_A+\gamma_g\otimes\Phi,
 \qquad \Phi=\Phi^\dagger,
\]
be the generated joint operator with
\(
 D_\mu\Phi=\partial_\mu\Phi+[A_\mu,\Phi]
\).
After applying the already derived coefficient \(1/2\) to the nonmetric
position-role output, the scalar-dependent local heat coefficients through
canonical dimension four give
\begin{equation}
\label{p3f:eq:v175-native-euclidean-scalar-action}
 I_{E,U_0}^{\Phi,\le4}
 =\frac1{16\pi^2}\int_{U_0}\operatorname{Tr}_F\left[
   (D_\mu\Phi)(D^\mu\Phi)+\Phi^4
   +\left(\frac{R_E}{6}-2\Lambda_g^2\right)\Phi^2
   -\frac13\Box_E\Phi^2
 \right]d\mu_E .
\end{equation}
On the interior unit-localizer domain with compactly supported fields and
variations, the displayed divergence integrates to zero.

For the Lorentzian real-form statement, additionally require the scalar and
connection first-jet germs to be holomorphic on the same oriented contour
sector as in
Lemma~\ref{p3f:lem:v174-yang-mills-complex-connection-real-form}, with no extra
contour-end term.  Then
\[
 \tau=i\widetilde t,
 \qquad D_\tau\Phi=-iD_{\widetilde t}\Phi
\]
gives
\begin{equation}
\label{p3f:eq:v175-native-lorentzian-scalar-action}
 I_{E,U_0}^{\Phi,\le4}\longmapsto-iS_{L,U_0}^{\Phi,\le4},
\end{equation}
where
\begin{equation}
\label{p3f:eq:v175-native-lorentzian-scalar-action-explicit}
 S_{L,U_0}^{\Phi,\le4}
 =\frac1{16\pi^2}\int_{U_0}\operatorname{Tr}_F\left[
  -(D_\mu\Phi)(D^\mu\Phi)
  -\frac{R_L}{6}\Phi^2
  +2\Lambda_g^2\Phi^2
  -\Phi^4
 \right]d\mu_L .
\end{equation}
Consequently the scalar kinetic term, curvature coupling, native quadratic and
quartic local coefficients, and the gauge--scalar cubic and quartic terms
contained in \(D_\mu\Phi\) are same-owner descendants.  Denote the certificate
by
\[
 \boxed{(\mathrm{ACT}_{\Phi,U_0,\le4}^{\rm loc,der})}.
\]
On the common overlap with the Einstein and v1.74 Yang--Mills theorems this
yields the restricted local-bosonic certificate
\[
 \boxed{(\mathrm{ACT}_{{\rm J,EH+YM+\Phi},U_0}^{\rm loc,der})}
 :=(\mathrm{ACT}_{{\rm J,EH},U_0}^{\rm int,der})
  +(\mathrm{ACT}_{{\rm YM},U}^{\rm der})
  +(\mathrm{ACT}_{\Phi,U_0,\le4}^{\rm loc,der}).
\]
This symbol means only the local parity-even Einstein, Yang--Mills, and
scalar-dependent canonical-dimension-at-most-four subfunctional.  It does not
include the pure-metric vacuum term, curvature-squared channels, the exact
nonlocal heat functional, or a selected odd-field/Yukawa member.
\end{theorem}

\begin{proof}
In the paper's Laplace-type convention the odd-field part of the square has
\[
 E_{A,\Phi}=E_A-\Phi^2-i\gamma^\mu\gamma_gD_\mu\Phi.
\]
The Clifford trace of the term linear in
\(\gamma^\mu\gamma_gD_\mu\Phi\) vanishes, while its square gives
\[
 \operatorname{tr}_{\mathbb S_g}
 \left(-i\gamma^\mu\gamma_gD_\mu\Phi\right)^2
 =4(D_\mu\Phi)(D^\mu\Phi).
\]
Substitution into the four-dimensional \(a_2\) and \(a_4\) coefficients gives,
before the nonmetric position-role weight,
\[
 a_2\big|_{\Phi^2}
 =-\frac4{(4\pi)^2}\int\operatorname{Tr}_F(\Phi^2)d\mu_E
\]
and
\[
 a_4\big|_\Phi
 =\frac1{(4\pi)^2}\int\operatorname{Tr}_F\left[
  2(D\Phi)^2+2\Phi^4+\frac13R_E\Phi^2
  -\frac23\Box_E\Phi^2
 \right]d\mu_E .
\]
Multiplication by the generated factor \(1/2\), with heat time
\(\Lambda_g^{-2}\), proves
\eqref{p3f:eq:v175-native-euclidean-scalar-action}.  The complete covariant
derivative contains the generated connection, so its expansion already owns
the gauge--scalar interactions.  On the additional holomorphic first-jet
domain the clock substitution changes the temporal kinetic term and curvature
real form together, yielding
\eqref{p3f:eq:v175-native-lorentzian-scalar-action-explicit}.  Independent
changes of the kinetic, curvature, quadratic, quartic, or gauge--scalar
coefficients would insert a new central weight into one joint-operator trace.
The rooted generated-image theorem therefore co-realizes the complete local
scalar-action family without selecting a field member.
\end{proof}

\begin{corollary}[Generated action-owner dependency]
\label{p3f:cor:generated-action-owner-dependency}
On their stated local theorem domains, the native Yang--Mills phase-space
action, joint-Dirac Einstein action, and odd-scalar action of this paper are
generated structural descendants of the already generated reciprocal,
connection, Dirac, clock, and metric-role data.  By the Paper-II-A typed
structural-realization equivalence, no additional premise asserting the
``physicality'' of these action owners is required once their licensed
constructor words have been discharged.

This does not choose numerical couplings, a boundary completion, a field or
solution member, a state/history, or any separately stated analytic
continuation or convergence domain.
\end{corollary}

\begin{corollary}[Joint-owner coefficient non-insertion]
\label{p3f:cor:joint-owner-coefficient-noninsertion}
On the native common-heat local domain, the Einstein, Yang--Mills and odd-scalar
terms generated from the same joint Dirac/trace owner are one action-constructor
packet.  In particular the scalar kinetic, curvature, quadratic, quartic and
gauge--scalar coefficients displayed in
Theorem~\ref{p3f:thm:v175-native-local-scalar-action} cannot be rescaled
independently while retaining the same frozen constructor.  A different
central weight, trace weight, profile or sector coefficient is a new datum and
requires a separately generated owner.
\end{corollary}

\begin{proof}
The heat-coefficient computation and common ordinary trace are part of one
licensed constructor word.  Paper~III-A
Theorem~\ref{p3a:thm:generated-action-constructor-closure} therefore retains
the whole generated packet and forbids independent deletion or rescaling of
one descendant term.
\end{proof}

\begin{remark}[Exact boundary of the v1.75 local action closure]
\label{p3f:rem:v175-local-bosonic-action-boundary}
The native exact heat potential on a constant scalar ray is convex and strictly
decreasing whenever its positive odd-field square is nonzero; the formal
minimum of the quartic truncation lies at the edge of its own uniform
small-field domain.  Thus
\((\mathrm{ACT}_{\Phi,U_0,\le4}^{\rm loc,der})\) does not discharge
\((\mathrm{VAC\!-\!EFF})\).  The pure-metric local finite-jet real forms are derived in Theorem~\ref{p3f:thm:finite-local-bosonic-clock-real-form}. Fermionic Lorentzian reality/phase, exact nonlocal or uncontrolled higher expansions, constant odd/Majorana offsets, noncanonical/global boundaries, continuum reconstruction and actual solutions, states and histories retain their claim-specific domains.
\end{remark}


\subsection{Local Noether and constrained-Hamiltonian closure of the joint \texorpdfstring{\(A+\Phi+g\)}{A+Phi+g} action}
\label{p3f:sec:v176-local-noether-constraint-closure}

\begin{definition}[Restricted local bosonic Noether domain]
\label{p3f:def:v176-local-bosonic-noether-domain}
On the overlap of the v1.74 Yang--Mills theorem and the v1.75 joint-Dirac
Einstein and native odd-scalar theorems, put
\[
 \alpha:=\frac1{16\pi^2},\qquad
 X:=\operatorname{Tr}_F(\Phi^2),\qquad
 \mathscr F(\Phi):=\kappa_J-\frac{\alpha}{6}X,
\]
and
\[
 V(\Phi):=\alpha\operatorname{Tr}_F
 \bigl(\Phi^4-2\Lambda_g^2\Phi^2\bigr).
\]
Then the already-derived restricted common action is merely rewritten as
\begin{equation}
\label{p3f:eq:v176-local-bosonic-action-repackaged}
 S_{\rm loc}[g,A,\Phi]
 =\int_{U_0}\!\sqrt{-g}\,
 \left[
 \mathscr F(\Phi)R
 -\frac14Q_\rho(F_{\mu\nu},F^{\mu\nu})
 -\alpha\operatorname{Tr}_F(D_\mu\Phi D^\mu\Phi)
 -V(\Phi)
 \right]d^4x .
\end{equation}
This definition adds no action term.  Every result below is restricted to the
common local, parity-even, compact-support/interior and regular first-jet
domains already printed for the three constituent theorems.
\end{definition}

\begin{theorem}[Native same-owner gauge and diffeomorphism invariance]
\label{p3f:thm:v166-native-local-action-invariance}
On the domain of
Definition~\ref{p3f:def:v176-local-bosonic-noether-domain}, let
\(u\) be a constituted smooth internal-gauge transformation with the same
compact-support/interior behavior already required by the local action, and let
\(\varphi\) be a constituted diffeomorphism supported in the same interior
domain.  Assume in each case that the transformation preserves every
retained source, coupling, orientation, grading, localizer, boundary/interface
datum, and other mark entering the complete action-owning presentation.  Then
these transformations are unmarked automorphisms of that complete presentation,
and
\begin{equation}
\label{p3f:eq:v166-native-gauge-action-invariance}
 \boxed{S_{\rm loc}[g,A^u,\Phi^u]=S_{\rm loc}[g,A,\Phi],}
\end{equation}
\begin{equation}
\label{p3f:eq:v166-native-diffeomorphism-action-invariance}
 \boxed{S_{\rm loc}[\varphi^*g,\varphi^*A,\varphi^*\Phi]
 =S_{\rm loc}[g,A,\Phi].}
\end{equation}
Thus the action-invariance antecedent used below is a native same-owner GP2
consequence on this branch, not an independently supplied physical Noether
premise.

The conclusion does not apply to a transformation that changes any retained
source, coupling, boundary datum, orientation, grading, state-dependent mark,
or other part of the complete action owner.  Such a transformation is not an
unmarked automorphism of the antecedent presentation and may represent explicit
symmetry breaking or a different physical input.  An occurring state may also
fail to be invariant even when the action is invariant; this does not alter the
action-level statement.
\end{theorem}

\begin{proof}
For the internal gauge action,
\[
 A_\mu\mapsto A_\mu^u,
 \qquad
 \Phi\mapsto\Phi^u:=\rho(u)\Phi\rho(u)^{-1},
\]
so that
\[
 F_{\mu\nu}\mapsto\operatorname{Ad}_uF_{\mu\nu},
 \qquad
 D_\mu\Phi\mapsto\rho(u)(D_\mu\Phi)\rho(u)^{-1}.
\]
The already-generated pairing \(Q_\rho\) is adjoint invariant and the finite
trace is cyclic.  Hence the Yang--Mills term, scalar kinetic term,
\(\operatorname{Tr}_F(\Phi^2)R\), and \(V(\Phi)\) are individually invariant,
while the metric is unchanged.  This proves
\eqref{p3f:eq:v166-native-gauge-action-invariance}.

For the diffeomorphism action, every factor in
\eqref{p3f:eq:v176-local-bosonic-action-repackaged} is built tensorially from
the common fields and their natural covariant derivatives, and the complete
integrand is a scalar density.  Pullback therefore transports the integrand
naturally; compact support in the interior unit-localizer domain removes any
new boundary contribution, proving
\eqref{p3f:eq:v166-native-diffeomorphism-action-invariance}.  The
gauge-covariant infinitesimal diffeomorphism tangent used below differs from
the ordinary Lie derivative of the connection by an infinitesimal gauge
transformation; because both constituent transformations leave the same action
invariant, that gauge-covariant combination is an invariant tangent as well.

These direct calculations realize, on this concrete common-action branch, the
trivial-target specialization of Paper~II-A v3.68's GP2 same-owner
equivariance theorem: a scalar descendant of an unmarked automorphism of its
complete frozen input cannot acquire a new presentation-dependent mark.  The
scope firewall is exactly the hypothesis that the transformation be an
unmarked automorphism of the \emph{complete} action-owning presentation.
\end{proof}

\begin{theorem}[Covariant Euler--Noether packet for the local bosonic action]
\label{p3f:thm:v176-local-bosonic-euler-noether}
On the domain of
Definition~\ref{p3f:def:v176-local-bosonic-noether-domain}, the scalar equation
is
\begin{equation}
\label{p3f:eq:v176-scalar-euler-equation}
 \boxed{D_\mu D^\mu\Phi-\frac R6\Phi
 +2\Lambda_g^2\Phi-2\Phi^3=0.}
\end{equation}
The variation with respect to the gauge connection defines the type-safe
scalar matter current \(J_\Phi^\mu\in\mathfrak g^*\) by
\begin{equation}
\label{p3f:eq:v176-scalar-gauge-current}
 \boxed{
 \langle J_\Phi^\mu,\epsilon\rangle
 :=2\alpha\operatorname{Tr}_F
 \bigl([\Phi,D^\mu\Phi]\,\rho(\epsilon)\bigr),
 \qquad\epsilon\in\mathfrak g,}
\end{equation}
with the overall sign fixed so that the gauge Euler equation is
\begin{equation}
\label{p3f:eq:v176-coupled-yang-mills-equation}
 \boxed{D_\mu(Q_\rho^\flat F^{\mu\nu})=J_\Phi^\nu.}
\end{equation}
No independent matter-charge coefficient is introduced.

Writing
\[
 T_{\mu\nu}^{\Phi,\min}
 :=2\alpha\operatorname{Tr}_F(D_\mu\Phi D_\nu\Phi)
 -g_{\mu\nu}\left[
 \alpha\operatorname{Tr}_F(D_\rho\Phi D^\rho\Phi)+V(\Phi)\right],
\]
the metric Euler equation is
\begin{equation}
\label{p3f:eq:v176-metric-euler-equation}
 \boxed{
 2\mathscr F G_{\mu\nu}
 +2(g_{\mu\nu}\Box-\nabla_\mu\nabla_\nu)\mathscr F
 =T_{\mu\nu}^{\rm YM}+T_{\mu\nu}^{\Phi,\min}.}
\end{equation}

Let \(\mathcal E_A^\mu,\mathcal E_\Phi,\mathcal E_g^{\mu\nu}\) denote the
Euler derivatives in the convention
\[
 \delta S=\int\sqrt{-g}\,
 [\mathcal E_g^{\mu\nu}\delta g_{\mu\nu}
 +\langle\mathcal E_A^\mu,\delta A_\mu\rangle
 +\operatorname{Tr}_F(\mathcal E_\Phi\delta\Phi)] .
\]
The native same-owner invariances
\eqref{p3f:eq:v166-native-gauge-action-invariance} and
\eqref{p3f:eq:v166-native-diffeomorphism-action-invariance}, differentiated
with the already licensed local variational calculus, give the off-shell
Noether-II identities
\begin{align}
\label{p3f:eq:v176-gauge-noether-identity}
 \boxed{D_\mu\mathcal E_A^\mu
 +\operatorname{pr}_{\mathfrak g^*}[\Phi,\mathcal E_\Phi]=0,}
\\
\label{p3f:eq:v176-diffeomorphism-noether-identity}
 \boxed{-2\nabla_\mu\mathcal E_g^\mu{}_{\nu}
 +\langle\mathcal E_A^\mu,F_{\nu\mu}\rangle
 +\operatorname{Tr}_F(\mathcal E_\Phi D_\nu\Phi)=0.}
\end{align}
Consequently, on the scalar equation,
\begin{equation}
\label{p3f:eq:v176-current-conservation}
 \boxed{D_\mu J_\Phi^\mu=0,}
\end{equation}
and on the joint matter equations the total stress is covariantly conserved.
For any fixed covariant split of the gauge and scalar stresses, their
individual divergences are equal and opposite and are fixed by the Lorentz
force density \(\langle J_\Phi^\mu,F_{\nu\mu}\rangle\), with its sign fixed by
the same global Euler-equation convention.  Hence the gauge--scalar exchange
is a solution functional of the same action, not an independent source law.
\end{theorem}

\begin{proof}
Variation of the scalar action gives
\eqref{p3f:eq:v176-scalar-euler-equation}.  Since
\(\delta_A(D_\mu\Phi)=[\delta A_\mu,\Phi]\), cyclicity of the finite trace
gives \eqref{p3f:eq:v176-scalar-gauge-current} and
\eqref{p3f:eq:v176-coupled-yang-mills-equation}.  The standard variation of
\(\sqrt{-g}\,\mathscr F R\) gives
\(\mathscr F G_{\mu\nu}+(g_{\mu\nu}\Box-\nabla_\mu\nabla_\nu)\mathscr F\),
which proves \eqref{p3f:eq:v176-metric-euler-equation}.
Theorem~\ref{p3f:thm:v166-native-local-action-invariance} supplies the exact
same-owner gauge and diffeomorphism invariances before any conservation
statement is used.  Differentiating those identities and inserting the
infinitesimal gauge variations
\(\delta_\epsilon A_\mu=-D_\mu\epsilon\),
\(\delta_\epsilon\Phi=[\epsilon,\Phi]\), and the gauge-covariant
diffeomorphism variations
\(\delta_\xi A_\mu=\xi^\rho F_{\rho\mu}\),
\(\delta_\xi\Phi=\xi^\rho D_\rho\Phi\),
\(\delta_\xi g_{\mu\nu}=2\nabla_{(\mu}\xi_{\nu)}\), then integrating by
parts against arbitrary compactly supported parameters, proves the two
off-shell identities.  Finally,
\(D_\mu[\Phi,D^\mu\Phi]=[\Phi,D^2\Phi]\), and the nondifferential right-hand
side of \eqref{p3f:eq:v176-scalar-euler-equation} commutes with \(\Phi\),
which proves \eqref{p3f:eq:v176-current-conservation}.  The stress statement
is the on-shell specialization of the diffeomorphism identity.  Thus the
Noether-II identities used here are ordinary variational consequences of the
already-derived same-owner invariances; they introduce no additional physical
symmetry or conservation postulate.
\end{proof}

\subsection{Kinetic-normalized scalar and gauge response operators}
\label{p3f:sec:v179-normalized-action-responses}

This construction differentiates the native joint action of
Theorem~\ref{p3f:thm:v175-native-local-scalar-action} on its stated local,
interior and analytic-continuation domain.  It uses the same rooted field
owner for the kinetic and potential derivatives.  In particular, the
normalization inheritance of Paper II-A,
Theorem~\ref{p2:thm:gp2-generated-owner-normalization-inheritance}, is
applied after that common ancestry is established.

\begin{theorem}[Kinetic-normalized local scalar Hessian]
\label{p3f:thm:v179-normalized-scalar-hessian}
Put \(C=1/(16\pi^2)\) and \(\mu_R^2=2\Lambda_g^2-R/6\);
the notation \(\mu_R^2\) does not assert positivity.
Hold the metric and connection fixed while taking the local algebraic
scalar block at a background \(\Phi_0=\Phi_0^\dagger\).
Let \(\mathcal R:V\to\operatorname{Herm}_{\rm odd}(F)\) be an admitted
linear source, with trace pairing
\(\langle X,Y\rangle=\operatorname{Re}\operatorname{Tr}_F(XY)\).
On the positive support of \(G=\mathcal R^\dagger\mathcal R\), set
\[
 Z=2CG,\qquad W=\mathcal R G^{-1/2}.
\]
Then the quadratic potential matrix and normalized mass block are
\begin{equation}
\label{p3f:eq:v179-scalar-hessian-compiler}
\begin{aligned}
 \mathcal M_{\Phi_0}^2(X)
 &=2(\Phi_0^2X+\Phi_0X\Phi_0+X\Phi_0^2)-\mu_R^2X,\\
 K&=2C\mathcal R^\dagger\mathcal M_{\Phi_0}^2\mathcal R,\\
 Z^{-1/2}KZ^{-1/2}&=W^\dagger\mathcal M_{\Phi_0}^2W,\qquad
 W^\dagger W=I .
\end{aligned}
\end{equation}
A source kernel may be quotiented only if no retained sector of the
claim-complete action uses it.  For a moving source, its induced connection
and derivatives remain in the kinetic action.
\end{theorem}
\begin{proof}
For \(V(\Phi)=C\operatorname{Tr}_F(\Phi^4-\mu_R^2\Phi^2)\), the first
gradient is \(C(4\Phi^3-2\mu_R^2\Phi)\). Its derivative is
\(2C\mathcal M_{\Phi_0}^2\), and pullback through \(\mathcal R\) gives \(K\).
The kinetic term \(-C\|\mathcal R\,d\varphi\|^2\) gives \(Z=2CG\).
Positive normalization and the polar isometry are those of Paper~I-D,
Proposition~\ref{p1b:prop:kinetic} \cite{Kugelmann2026Paper1D}, applied
on this positive source support. The restriction to a claim-complete
quotient follows from Theorem~\ref{p3a:thm:claim-complete-action-descent}.
\end{proof}

\begin{corollary}[Source-coordinate invariance of the normalized scalar Hessian]
\label{p3f:cor:v177-source-coordinate-polar-invariance}
Under the hypotheses of
Theorem~\ref{p3f:thm:v179-normalized-scalar-hessian}, let
$A_V:V\to V$ be invertible on the positive source support and replace the
source presentation by
\[
 \mathcal R' = \mathcal R A_V.
\]
Let $M_B^2=Z^{-1/2}KZ^{-1/2}$ and $(M_B^2)'$ denote the corresponding
canonically normalized Hessian blocks.  Then
\[
 \boxed{(M_B^2)'=U_{A_V}^\dagger M_B^2U_{A_V}},
 \qquad
 U_{A_V}:=G^{1/2}A_V(A_V^\dagger G A_V)^{-1/2},
\]
with $U_{A_V}^\dagger U_{A_V}=I$.  Hence invertible changes of source
coordinates, including a common nonzero source amplitude, cannot change the
physical bosonic spectrum on a fixed normalized source image and background
Hessian.
\end{corollary}
\begin{proof}
Apply the coordinate-invariance statement of Paper~I-D,
Proposition~\ref{p1b:prop:kinetic} \cite{Kugelmann2026Paper1D}, to
\(\mathcal R\) and \(\mathcal R A_V\) on their common positive image.
Its unitary polar factor is precisely the displayed \(U_{A_V}\), so the
compression in Theorem~\ref{p3f:thm:v179-normalized-scalar-hessian}
transforms by the asserted unitary conjugation.
\end{proof}


\subsection{Stationary physical Hessians and generated propagation operators}
\label{p3f:sec:v244-physical-hessian-propagation}

The local action is the owner in this subsection.  A ``wave operator'' is
therefore not supplied independently when the same action, a stationary
solution, a valid physical reduction, and a stable quadratic Hessian are
already available.

\begin{theorem}[Stationary physical-Hessian propagation theorem]
\label{p3f:thm:v244-stationary-physical-hessian-propagation}
Let \(S[\Phi]\) be an already constituted local action on one of its stated
real Lorentzian domains and let \(\bar\Phi\) be a stationary solution.  After
all gauge and constraint reduction required by the claim, suppose the real
physical perturbation \(\eta\) has quadratic action
\begin{equation}
\label{p3f:eq:v244-physical-quadratic-action}
 S^{(2)}[\eta]
 =\frac12\int dt\,
 \left[
 \langle\dot\eta,M\dot\eta\rangle
 -\langle\eta,L\eta\rangle
 \right],
\end{equation}
with time-independent \(M=M^*>0\) and \(L=L^*\ge0\) on the physical block.
Then the Euler equation is
\[
 M\ddot\eta+L\eta=0.
\]
With the canonical kinetic coordinate \(q=M^{1/2}\eta\), define
\begin{equation}
\label{p3f:eq:v244-generated-frequency-operator}
 \boxed{\Omega^2:=M^{-1/2}LM^{-1/2}\ge0}.
\end{equation}
Then
\begin{equation}
\label{p3f:eq:v244-generated-lorentzian-operator}
 \boxed{K_L:=\partial_t^2+\Omega^2},
 \qquad K_Lq=0,
\end{equation}
and the canonically normalized quadratic action is
\begin{equation}
\label{p3f:eq:v244-canonical-quadratic-action}
 \boxed{
 S_L^{(2)}[q]
 =\frac12\int dt\,
 \left(\|\dot q\|^2-\|\Omega q\|^2\right).}
\end{equation}
Thus, on this stationary stable physical domain, the quadratic action,
kinetically normalized Hessian and Lorentzian propagation operator are the
same same-owner datum in three presentations; no independent second-order
wave-law premise is needed.

The displayed normal form is not asserted for a parity-odd or gyroscopic
quadratic sector containing a nontrivial mixed term such as
\(\langle\dot\eta,B\eta\rangle\), nor for a time-dependent kinetic metric.
Those cases require their own physical Hessian reduction.
\end{theorem}

\begin{proof}
Varying \eqref{p3f:eq:v244-physical-quadratic-action} and integrating the
kinetic term by parts gives
\[
 \delta S^{(2)}
 =-\int dt\,\langle\delta\eta,M\ddot\eta+L\eta\rangle,
\]
up to the stated endpoint rule.  Since \(M\) is positive and stationary,
\(q=M^{1/2}\eta\) is an invertible physical coordinate change and
\(\ddot q=M^{1/2}\ddot\eta\).  Multiplication by \(M^{-1/2}\) yields
\eqref{p3f:eq:v244-generated-lorentzian-operator}; substitution into the
quadratic action yields \eqref{p3f:eq:v244-canonical-quadratic-action}.
Positivity of \(\Omega^2\) follows by congruence from \(L\ge0\).
\end{proof}

\begin{corollary}[Scalar harmonic propagation from the normalized action Hessian]
\label{p3f:cor:v244-scalar-harmonic-propagation}
On a stationary round spatial branch of the scalar action, let
\(\widehat M_{\Phi_0}^2=Z^{-1/2}KZ^{-1/2}\) be the normalized internal
Hessian of Theorem~\ref{p3f:thm:v179-normalized-scalar-hessian}.  If an
internal normalized eigenvector has nonnegative Hessian contribution
\(\nu_A^2\), then the corresponding scalar perturbation obeys
\[
 \boxed{
 K_A=\partial_t^2+c_{\rm op}^2\Delta_0+\nu_A^2.}
\]
For a scalar harmonic \(Y_{k\alpha}\) with
\(\Delta_0Y_{k\alpha}=\lambda_k^{(0)}Y_{k\alpha}\),
\[
 \boxed{\omega_{A,k}^2=c_{\rm op}^2\lambda_k^{(0)}+\nu_A^2.}
\]
If a downstream physical mass calibration identifies
\(\nu_A^2=m_A^2c_{\rm op}^4/\hbar^2\), this becomes the usual massive
relativistic finite-mode formula.  The calibration is not created by the
Hessian theorem itself.
\end{corollary}

\begin{corollary}[Source-free transverse gauge propagation]
\label{p3f:cor:v244-transverse-gauge-propagation}
On the stationary product domain of
Theorem~\ref{p3f:thm:v174-generated-joint-gauge-phase-space-action}, linearize
an unbroken Abelian gauge direction about the trivial spatial connection and
perform the valid source-free Gauss/gauge reduction to a physical spatial
one-form \(a_T\) with \(\delta a_T=0\).  After whitening the common positive
internal trace pairing, the quadratic action is, up to its inherited common
positive scalar,
\begin{equation}
\label{p3f:eq:v244-transverse-gauge-action}
 S_T^{(2)}[a_T]
 =\frac1{2c_{\rm op}}\int dt\,
 \left(\|\dot a_T\|^2-c_{\rm op}^2\|da_T\|^2\right).
\end{equation}
Hence
\begin{equation}
\label{p3f:eq:v244-transverse-gauge-operator}
 \boxed{
 K_{0,T}=\partial_t^2+c_{\rm op}^2\Delta_1^{\rm co}.}
\end{equation}
On an oriented three-dimensional spatial slice, writing \(C=*d\) on coexact
one-forms gives
\[
 \boxed{K_{0,T}=\partial_t^2+c_{\rm op}^2C^2}.
\]
This is the physical transverse operator generated by the Yang--Mills action
and Gauss reduction.  It does not select a gauge-field state or history, and
it does not include a Higgs mass Hessian or a longitudinal massive-vector
mode.
\end{corollary}

\begin{proof}
At the Abelian trivial background the quadratic curvature is \(F=dt\wedge
\dot a_T+da_T\).  Substitution into the already generated Lorentzian
Yang--Mills action gives \eqref{p3f:eq:v244-transverse-gauge-action}.  Its
variation gives
\(\ddot a_T+c_{\rm op}^2\delta d\,a_T=0\).  Because
\(\delta a_T=0\), one has \(\delta d\,a_T=\Delta_1^{\rm co}a_T\).  On
coexact one-forms in dimension three, \((*d)^2=\Delta_1^{\rm co}\), which
gives the curl form.
\end{proof}

\begin{theorem}[Connection of a moving positive kinetic normalization]
\label{p3f:thm:v199-rgas-kinetic-connection}
On a fixed finite-dimensional positive kinetic support, let
\(Z_L(\ell),Z_R(\ell)\) be differentiable positive Hermitian matrices,
\(S_a=Z_a^{1/2}\), and \(A_a=\dot S_aS_a^{-1}\), where a dot denotes
\(\partial_\ell\).  For a differentiable mass block \(D\), set
\(\widehat D=S_L^{-1}DS_R^{-1}\).  Without assuming commutativity,
\begin{equation}
\label{p3f:eq:v199-rgas-kinetic-connection}
 \dot{\widehat D}+A_L^\dagger\widehat D+\widehat D A_R
       =S_L^{-1}\dot D S_R^{-1},\qquad
 A_a+A_a^\dagger=S_a^{-1}\dot Z_aS_a^{-1}.
\end{equation}
The same congruence formula holds for a bosonic Hessian \(K\), with the
same kinetic matrix on its two sides.  For a scalar \(Z\),
\(A=\tfrac12\dot{\log Z}=-\eta/2\), where \(\eta=-\dot{\log Z}\).
\end{theorem}
\begin{proof}
Use the moving positive normalization of Paper~I-D,
Proposition~\ref{p1b:prop:kinetic} \cite{Kugelmann2026Paper1D}.
Here \(S_a^{-1}\dot S_a=A_a^\dagger\) and
\(\dot S_a^{-1}=-S_a^{-1}\dot S_aS_a^{-1}\); differentiating
\(\widehat D\) gives the first identity. Differentiating \(Z_a=S_a^2\)
and multiplying by \(S_a^{-1}\) on both sides gives the second.
These formulas hold on the fixed positive support and do not cross a rank change.
\end{proof}

\begin{corollary}[Moving spectral frame]
\label{p3f:cor:v199-rgas-spectral-frame}
For a differentiable unitary frame \(U(\ell)\), let
\(\Omega=U^\dagger\dot U\) and
\(\widetilde K=U^\dagger\widehat K U\).  Then
\[
 \dot{\widetilde K}+[\Omega,\widetilde K]
                  =U^\dagger\dot{\widehat K}U .
\]
This holds on a differentiable frame patch, including a degenerate
cluster when its frame is chosen differentiably.
\end{corollary}
\begin{proof}
Differentiate \(U^\dagger U=I\), so
\(\dot U^\dagger U=-\Omega\), and apply the product rule.
\end{proof}
These connections transport the normalized action and physical sources
together.  An admitted invertible field change preserves the scalaron
pole; it does not delete it as a normalization.
Theorem~\ref{p12:thm:v199-rgas-moving-eaa} gives the finite quantum flow
with its regulator and measure retained.


\begin{remark}[Canonical physical-visibility criterion]
\label{p3f:rem:v177-canonical-physical-visibility}
A raw finite factor carried only by a source coordinate is not by itself a
physical mass multiplier.  On the local scalar branch it can survive
canonical normalization only by changing the normalized source image,
changing the background Hessian on that image, or entering kinetic and
potential jets asymmetrically.  This criterion applies in particular to raw
incidence ranks, projector norms and common scalar amplitudes; it does not
forbid a separately generated asymmetric same-action insertion.
\end{remark}

\begin{corollary}[Nonlinear source chain rule and the Noether mass map]
\label{p3f:cor:v179-noether-normalized-response}
For a nonlinear admitted field source \(\Phi=\mathcal F(\varphi)\),
\[
 D^2(V\circ\mathcal F)(v,w)
 =D^2V(D\mathcal Fv,D\mathcal Fw)
   +DV(D^2\mathcal F(v,w)).
\]
At a full stationary matrix point, a permitted conjugation symmetry has
\(\mathcal M_{\Phi_0}^2([\xi,\Phi_0])=0\).  Away from that stationarity
the orbit-acceleration term in the displayed chain rule must be retained.
For the gauge kinetic pairing \(Q_\rho\) of the same joint action, put
\(\mathcal J_{\Phi_0}(\xi)=[\rho(\xi),\Phi_0]\).  On its positive physical
gauge support the scalar contribution to the gauge-boson mass operator is
\begin{equation}
\label{p3f:eq:v179-gauge-mass-commutator}
 M_{\rm gauge}^2
 =2C Q_\rho^{-1/2}\mathcal J_{\Phi_0}^\dagger
                      \mathcal J_{\Phi_0}Q_\rho^{-1/2}.
\end{equation}
The scalar Ward null and the gauge-boson mass form are two derivatives of
one coupled action.  The stationary premise in this corollary is used only
for the scalar Ward-null statement.  The gauge Gram
\eqref{p3f:eq:v179-gauge-mass-commutator} is valid for any fixed occurring
odd zero-order source/background on the stated positive kinetic domain.
\end{corollary}
\begin{proof}
Differentiate the composite potential twice.  For a conjugation orbit,
the acceleration is \([\xi,[\xi,\Phi_0]]\).
The gradient of a conjugation-invariant trace polynomial is covariant;
at a zero of the full gradient its derivative along the symmetry orbit
vanishes.  Finally,
\(D_\mu\Phi_0=[\rho(A_\mu),\Phi_0]\); the scalar kinetic term supplies
\(2C\mathcal J^\dagger\mathcal J\) in the quadratic mass convention.
Normalizing by the existing gauge pairing gives
\eqref{p3f:eq:v179-gauge-mass-commutator}, consistently with the current and
Noether identities of Theorem~\ref{p3f:thm:v176-local-bosonic-euler-noether}.
\end{proof}

\begin{corollary}[Gauge Gram and stabilizer without a scalar-stationarity premise]
\label{p3f:cor:v111-gauge-gram-no-stationarity}
Let \(\Phi_0\) be any fixed occurring odd zero-order background on the local
scalar-action domain and define
\[
 \mathcal J_{\Phi_0}(\xi)=[\rho(\xi),\Phi_0].
\]
Before kinetic whitening the algebraic gauge quadratic form is
\[
 2C\,\mathcal J_{\Phi_0}^\dagger\mathcal J_{\Phi_0}.
\]
Hence
\[
 \boxed{\ker(\mathcal J_{\Phi_0}^\dagger\mathcal J_{\Phi_0})
 =\ker\mathcal J_{\Phi_0}
 =\{\xi:[\rho(\xi),\Phi_0]=0\}.}
\]
Positive invertible gauge-kinetic normalization transports this null space
isomorphically.  No stationarity condition on the scalar potential is used.
Stationarity is required only for the separate scalar Ward-null statement
\(\mathcal M_{\Phi_0}^2([\xi,\Phi_0])=0\).
\end{corollary}

\begin{proof}
For every gauge generator \(\xi\),
\[
 \langle\xi,\mathcal J_{\Phi_0}^\dagger
 \mathcal J_{\Phi_0}\xi\rangle
 =\|\mathcal J_{\Phi_0}\xi\|^2,
\]
so the Gram kernel equals the kernel of \(\mathcal J_{\Phi_0}\).
The identity \(D_\mu\Phi_0=[\rho(A_\mu),\Phi_0]\) shows that this Gram is
supplied by the scalar covariant kinetic term itself.  Conjugation by the
positive invertible kinetic square root preserves nullity.  The stationary
hypothesis in Corollary~\ref{p3f:cor:v179-noether-normalized-response}
was used for the scalar Hessian Ward null, not for this kinetic Gram identity.
\end{proof}

For background eigenvalues \(a,b\), the algebraic off-diagonal coefficient
of \(\mathcal M_{\Phi_0}^2\) is
\(2(a^2+ab+b^2)-\mu_R^2\).  These are eigenvalues of the actual field
background, not automatically hierarchy energies.
The physical Higgs source is supplied explicitly in
Theorem~\ref{p5:thm:v179-canonical-higgs-direction}.
The Hurwitz and Clifford source metrics remain those of their own carriers
until an admitted field map transports them into this action.

\subsection{Radial normalization with the retained scalar--metric coupling}
\label{p3f:sec:v179-radial-metric-response}

\begin{proposition}[Local scalar--tensor radial response]
\label{p3f:prop:v179-radial-einstein-response}
On a constant linear odd-field ray with canonical real coordinate \(h\)
and no constant offset, \(C\operatorname{Tr}_F\Phi(h)^2=h^2/2\).
The curvature coefficient in the joint local action is therefore
\[
 F(h)=\kappa_J-h^2/12 .
\]
Restrict to \(F>0\), keep a fixed positive reference \(F_{\rm ref}\), and
write the radial Jordan action as
\(F R-\tfrac12(\partial h)^2-V_J(h)\).
Under \(g_E=(F/F_{\rm ref})g\), retaining the appropriate boundary term,
\begin{equation}
\label{p3f:eq:v179-einstein-radial-metric}
 Z_E=F_{\rm ref}\left(\frac1F+\frac{3(F')^2}{F^2}\right)
     =\frac{F_{\rm ref}\kappa_J}{F^2},\qquad
 V_E=\left(\frac{F_{\rm ref}}F\right)^2 V_J .
\end{equation}
At a stationary point \(F V_J'-2F'V_J=0\), its normalized radial block is
\begin{equation}
\label{p3f:eq:v179-einstein-radial-mass}
 m_{E,\rm rad}^2
 =\frac{F_{\rm ref}}{\kappa_J F}
      (F V_J''-F'V_J'-2F''V_J).
\end{equation}
\end{proposition}
\begin{proof}
The curvature term in
Definition~\ref{p3f:def:v176-local-bosonic-noether-domain} gives \(F\)
after the same kinetic normalization.  The four-dimensional conformal
curvature transformation, integrated with its retained boundary term,
adds \(3F_{\rm ref}(F')^2/F^2\) to the radial kinetic coefficient.
Since \(F'=-h/6\), \(F+3(F')^2=\kappa_J\), proving the simplification.
Differentiate \(V_E=F_{\rm ref}^2V_J/F^2\) and impose its stationary
condition.  Division by \(Z_E\) gives the mass formula.
\end{proof}

A contributing vacuum-energy term must remain in \(V_J\).
For a Minkowski stationary point with \(V_J=V_J'=0\) and
\(F_{\rm ref}=F(v)\), the formula becomes
\(m_{E,\rm rad}^2=F(v)V_J''(v)/\kappa_J\).
This is a local scalar--tensor restriction of the joint action, not a
replacement of its full nonlinear gravitational or effective response.
For other retained fluctuations \(\chi\), the complete quadratic
elimination on an invertible block is
\[
 \Gamma_{\rm red}(p)=\Gamma_{hh}(p)
 -\Gamma_{h\chi}(p)\Gamma_{\chi\chi}(p)^{-1}\Gamma_{\chi h}(p).
\]
Its kinetic residue and pole must be extracted from that same reduced
operator, with the same Ward and boundary data.  The source norm alone
does not permit discarding a coupled sector.


\begin{remark}[Internal same-owner status of the local Noether packet]
\label{p3f:rem:v166-internal-noether-status}
Theorem~\ref{p3f:thm:v166-native-local-action-invariance} supplies the
physical action-invariance antecedent from the generated complete action owner.
Theorem~\ref{p3f:thm:v176-local-bosonic-euler-noether} then uses only the
already licensed local first-variation and integration-by-parts calculus.
Accordingly, on this branch the current, Noether-II identities, coupled
moment-map relation, constraint zero locus, and differentiated Ward relations
are internal theorem closure in the sense of Paper~II-A v3.68 and Paper~III-A
v1.49.  This retyping does not weaken the analytic and owner firewalls below:
global charges, boundary Hamiltonians, quantum Ward identities, well-posed
flows, and actual histories retain their separately stated hypotheses.
\end{remark}

\begin{corollary}[Coupled gauge--scalar moment map and fibre-product constraint]
\label{p3f:cor:v176-coupled-gauge-scalar-moment-map}
Let \(p_\Phi\) be the scalar canonical momentum on a spacelike slice
\(\Sigma\).  The scalar gauge action has moment-map contribution
\begin{equation}
\label{p3f:eq:v176-scalar-moment-map}
 \langle\mu_\Phi,\epsilon\rangle
 =\int_\Sigma\operatorname{Tr}_F
 \bigl(p_\Phi[\rho(\epsilon),\Phi]\bigr),
\end{equation}
and the v1.74 pure-gauge moment map is replaced on the coupled branch by
\begin{equation}
\label{p3f:eq:v176-coupled-gauss-constraint}
 \boxed{
 \mathcal G
 =d_a\Pi+\operatorname{pr}_{\mathfrak g^*}[\Phi,p_\Phi]\approx0,}
\end{equation}
up to the single overall Gauss-sign convention already present in the pure
Yang--Mills theorem.  Thus
\[
 \mu_{\rm tot}=\mu_{\rm YM}+\mu_\Phi
\]
and the constrained descendant is the generated fibre product
\begin{equation}
\label{p3f:eq:v176-gauge-scalar-fibre-product}
 \boxed{
 \mu_{\rm tot}^{-1}(0)
 \cong T^*\mathcal A_\Sigma
 \times_{\mathfrak g_\Sigma^*}
 T^*\mathcal V_\Phi,}
\end{equation}
with the sign on one leg chosen according to that convention.  The relation is
generated internally by the common action, its native same-owner invariance,
and the resulting Noether-II identity; it is not inserted retroactively into
the earlier unconstrained role product and requires no second relation owner.
\end{corollary}

\begin{theorem}[Regular ADM Legendre map and complete local bulk constraint packet]
\label{p3f:thm:v176-adm-constraint-packet}
Choose the standard clock-adapted ADM split
\[
 ds^2=-N^2dt^2+h_{ij}(dx^i+N^i dt)(dx^j+N^j dt),
 \qquad
 K_{ij}=\frac1{2N}(\dot h_{ij}-D_iN_j-D_jN_i),
\]
and put \(v_\Phi:=D_\perp\Phi\), \(E_i:=F_{i\perp}\).  Modulo the local
compact-support/divergence terms allowed by the present interior theorem, the
canonical momenta are
\begin{align}
\label{p3f:eq:v176-adm-metric-momentum}
 \pi^{ij}&=\sqrt h\left[
 \mathscr F(K^{ij}-Kh^{ij})-h^{ij}D_\perp\mathscr F\right],\\
\label{p3f:eq:v176-adm-scalar-momentum}
 p_\Phi&=2\alpha\sqrt h\left(v_\Phi+\frac K3\Phi\right),\\
\label{p3f:eq:v176-adm-gauge-momentum}
 \Pi^i&=\sqrt h\,Q_\rho^\flat(E^i).
\end{align}
Writing \(\pi=h_{ij}\pi^{ij}\), the scalar--metric trace block satisfies the
exact identity
\begin{equation}
\label{p3f:eq:v176-dilatational-momentum-combination}
 \boxed{
 \mathcal C_{\rm dil}
 :=\pi-\frac12\operatorname{Tr}_F(\Phi p_\Phi)
 =-2\kappa_J\sqrt h\,K.}
\end{equation}
Moreover its Schur complement is
\begin{equation}
\label{p3f:eq:v176-gravity-scalar-legendre-determinant}
 \boxed{
 \Delta_{g\Phi}
 :=\mathscr F+\frac{3}{2\alpha}
 \|\nabla_\Phi\mathscr F\|^2
 =\kappa_J>0.}
\end{equation}
Hence the conformally coupled scalar--metric trace block carries no additional
primary conformal constraint on this branch.  The traceless metric block
still requires the separate regularity condition \(\mathscr F(\Phi)\ne0\)
(and a positive effective graviton branch requires \(\mathscr F>0\)).
On that regular branch the inverse Legendre map is
\begin{align}
 K_T^{ij}&=\frac{\pi_T^{ij}}{\sqrt h\,\mathscr F},\\
 K&=-\frac{\mathcal C_{\rm dil}}{2\kappa_J\sqrt h},\\
 v_\Phi&=\frac{p_\Phi}{2\alpha\sqrt h}
 +\frac{\mathcal C_{\rm dil}}{6\kappa_J\sqrt h}\Phi.
\end{align}
The bulk Hamiltonian is a sum of multipliers and constraints,
\begin{equation}
\label{p3f:eq:v176-bulk-constraint-hamiltonian}
 \boxed{
 H_{\rm bulk}=\int_\Sigma
 \left[N\mathcal H_\perp+N^i\mathcal H_i
 +\langle A_\perp,\mathcal G\rangle\right]
 +H_{\rm primary},}
\end{equation}
where
\begin{align}
\mathcal H_{g\Phi}^{\rm kin}
={}&\frac{1}{\sqrt h\,\mathscr F}
 \left(\pi_{ij}\pi^{ij}-\frac13\pi^2\right)
 +\frac{1}{4\alpha\sqrt h}\operatorname{Tr}_F(p_\Phi^2)
 -\frac{\mathcal C_{\rm dil}^2}{6\kappa_J\sqrt h},\\
\mathcal H_{g\Phi}^{\rm pot}
={}&\sqrt h\left[-\mathscr F\,{}^{(3)}R
 +2D_iD^i\mathscr F
 +\alpha\operatorname{Tr}_F(D_i\Phi D^i\Phi)+V(\Phi)\right],\\
\mathcal H_{\rm YM}
={}&\frac1{2\sqrt h}\langle\Pi_i,Q_\rho^\sharp\Pi^i\rangle
 +\frac{\sqrt h}{4}Q_\rho(F_{ij},F^{ij}),\\
\boxed{\mathcal H_\perp}
={}&\mathcal H_{g\Phi}^{\rm kin}
 +\mathcal H_{g\Phi}^{\rm pot}+\mathcal H_{\rm YM}\approx0,\\
\boxed{\mathcal H_i}
={}&-2D_j\pi^j{}_i
 +\langle\Pi^j,F_{ij}\rangle
 +\operatorname{Tr}_F(p_\Phi D_i\Phi)\approx0,
\end{align}
with \(\mathcal G\) given by
\eqref{p3f:eq:v176-coupled-gauss-constraint}.  The primary multiplier
constraints are
\[
 \boxed{p_N\approx0,\qquad p_{N^i}\approx0,\qquad\Pi_\perp\approx0.}
\]
Their preservation gives respectively
\(\mathcal H_\perp\), \(\mathcal H_i\), and \(\mathcal G\).
\end{theorem}

\begin{proof}
The ADM decomposition of \(\sqrt{-g}\,\mathscr F R\) contains the kinetic
cross term \(-2N\sqrt h\,K D_\perp\mathscr F\).  Since
\(\nabla_\Phi\mathscr F=-\alpha\Phi/3\), differentiation with respect to
\(\dot h_{ij}\), \(v_\Phi\), and the gauge velocity gives
\eqref{p3f:eq:v176-adm-metric-momentum}--
\eqref{p3f:eq:v176-adm-gauge-momentum}.  Taking the metric trace and
subtracting one half of \(\operatorname{Tr}_F(\Phi p_\Phi)\) cancels the
scalar normal velocity and yields
\eqref{p3f:eq:v176-dilatational-momentum-combination}, because
\(\mathscr F+\alpha X/6=\kappa_J\).  Likewise
\[
 \|\nabla_\Phi\mathscr F\|^2=\alpha^2X/9,
\]
which proves
\eqref{p3f:eq:v176-gravity-scalar-legendre-determinant}.  Inverting the
regular trace/traceless blocks and performing the Legendre transform gives the
displayed Hamiltonian densities.  Lapse, shift, and \(A_\perp\) carry no
independent velocities and therefore supply the three multiplier constraints.
\end{proof}

\begin{corollary}[First-class propagation and reduced local phase-space role]
\label{p3f:cor:v176-first-class-constraint-propagation}
On the regular interior branch of
Theorem~\ref{p3f:thm:v176-adm-constraint-packet}, let
\(G[\epsilon]\), \(D[\xi]\), and \(H[N]\) be the smeared Gauss,
gauge-covariant spatial-diffeomorphism, and normal constraints.  Their bulk
Poisson algebra closes with the usual field-dependent structure functions:
\[
\begin{aligned}
 \{G,G\}&\sim G,&\{D,G\}&\sim G,&\{D,D\}&\sim D,\\
 \{D,H\}&\sim H,&\{G,H\}&=0,&\{H,H\}&\sim D\pmod G.
\end{aligned}
\]
Consequently the same-action Hamiltonian evolution, constrained by the
native Noether-II identities, preserves the complete constraint locus,
\begin{equation}
\label{p3f:eq:v176-constraint-propagation}
 \boxed{\mathcal C(t_0)=0\Longrightarrow\mathcal C(t)=0,
 \qquad
 \mathcal C=(\mathcal G,\mathcal H_i,\mathcal H_\perp).}
\end{equation}
The corresponding physical bulk role is therefore the stratified/groupoid
reduction of the joint zero locus by the generated internal-gauge and local
diffeomorphism actions.  No globally smooth quotient, ADM boundary energy, or
physical clock is asserted.
\end{corollary}

\begin{corollary}[Noether--Ward derivative tower]
\label{p3f:cor:v176-noether-ward-derivative-tower}
Let \(\varphi^I=(g,A,\Phi)\) and let
\(R_\alpha^I(\varphi)\) be a gauge or diffeomorphism tangent.  From
\(S_{,I}R_\alpha^I=0\), differentiation and restriction to an on-shell
background give
\begin{equation}
\label{p3f:eq:v176-hessian-ward-null}
 \boxed{K_{IJ}R_\alpha^J=0,
 \qquad K:=\delta^2S,}
\end{equation}
and one further derivative gives
\begin{equation}
\label{p3f:eq:v176-cubic-ward-identity}
 \boxed{
 V_{IJK}R_\alpha^K
 +K_{KJ}R_{\alpha,I}^K
 +K_{KI}R_{\alpha,J}^K=0,
 \qquad V:=\delta^3S.}
\end{equation}
Higher action variations obey the corresponding differentiated Ward
relations.  These identities constrain the physical response tensors after
constraint/gauge reduction; they do not imply that an allowed transverse
vertex coefficient is nonzero.
\end{corollary}

\begin{corollary}[Restricted broken-Weyl trace identity]
\label{p3f:cor:v176-restricted-broken-weyl-trace}
For the local truncation
\eqref{p3f:eq:v176-local-bosonic-action-repackaged} only, the infinitesimal
four-dimensional Weyl variation
\(\delta_\sigma g_{\mu\nu}=2\sigma g_{\mu\nu}\),
\(\delta_\sigma\Phi=-\sigma\Phi\), \(\delta_\sigma A_\mu=0\)
leaves the gauge, quartic, kinetic, and conformal \(R\Phi^2/6\) package
invariant up to the local divergence.  On the joint equations the explicit
breaking terms therefore obey
\begin{equation}
\label{p3f:eq:v176-restricted-weyl-trace}
 \boxed{2\kappa_JR+4\alpha\Lambda_g^2
 \operatorname{Tr}_F(\Phi^2)=0.}
\end{equation}
\end{corollary}

\begin{remark}[Exact v1.76 Noether-domain firewall]
\label{p3f:rem:v176-noether-domain-firewall}
The preceding packet belongs only to the restricted common local
Einstein--Yang--Mills--odd-scalar action on its proved interior/regular domain.
It does not derive the Lorentzian fermionic current before
\((\mathrm{REAL}_{F,L})\), a global or edge charge before a boundary owner,
an ADM/Brown--York/Iyer--Wald boundary Hamiltonian, an exact-nonlocal quantum
measure, anomaly-free quantum Ward identities, a physical clock or actual
history, the nonzero K16 response coefficients, the light-front species
projectors, \(YFIX_5\), \(\mathrm{VAC\!-\!EFF}\), the RAR
\(B_{48}\to J_\partial\) map, or the critical \(\sqrt y\) lift.  In
particular, neither the trace identity
\eqref{p3f:eq:v176-restricted-weyl-trace} nor the ADM packet may be inserted
into the nonlinear \(f(R)\)/scalaron branch without a separately proved
common-action ancestry.  The present internalization removes an independent
physical Noether premise only; it does not turn boundary regularity, anomaly
cancellation, Cauchy well-posedness, or any other analytic-certification
hypothesis into a GP consequence.
\end{remark}

\subsection{A common finite-jet bosonic real form}

\begin{definition}[Common holomorphic finite-jet owner]
\label{p3f:def:common-holomorphic-finite-jet-owner}
A common owner for a local heat truncation consists of one
nondegenerate holomorphic metric/connection/scalar configuration,
the field and clock jets needed by that truncation, its analytic
coefficient functions, and a common oriented interior or explicitly
transported boundary problem.  Its real slices use \(\tau=it\)
and the same invariant tensor contractions.  Write
\[
 I_{E,K}=\sum_{a\leq K}c_a\int P_{a,E}\,d\mu_E
\]
for the finite local polynomial expression.  The continuation
requires no summation of the full heat asymptotic series.
\end{definition}

\begin{theorem}[Finite local bosonic clock real form]
\label{p3f:thm:finite-local-bosonic-clock-real-form}
On this common owner, and only on its stated domain,
\[
 I_{E,K}\longmapsto-iS_{L,K},\qquad
 S_{L,K}=-\sum_{a\leq K}c_a\int P_{a,L}\,d\mu_L.
\]
This includes the local \(a_0,a_2\) and parity-even pure-metric
\(a_4\) terms, and any specified finite local truncation for which
the required common jets and contractions are supplied.  Actual
covariant symmetries of the claim-complete owner give its local
Euler and Noether identities.
\end{theorem}
\begin{proof}
Each local scalar \(P_a\) is the restriction of the same finite
holomorphic contraction on both slices.  Nondegeneracy makes the
inverse metric and contractions holomorphic on the chosen domain.
The clock substitution transports the oriented volume to
\(d\mu_E=i\,d\mu_L\).  Therefore the continued integral is
\(i\sum c_a\int P_{a,L}d\mu_L=-iS_{L,K}\).
The finite sum permits termwise differentiation through the
declared jet order.  Applying
Theorem~\ref{p3a:thm:claim-complete-action-descent} to its actual
symmetries proves the last assertion, including transported
boundary variations when those are part of the owner.
\end{proof}

This theorem extends the covariant local real-form and Noether
domain.  It does not extend the two-derivative ADM or constraint
calculation to curvature-squared actions: their additional
derivative data require a new Legendre and constraint analysis.
Exact nonlocal spectral actions, global continuation, fermionic
reality, noncanonical boundary terms and renormalized vacuum
selection retain their own claim-specific requirements.


\begin{theorem}[Joint spectral action and conditional stationarity]
\label{p3:thm:common-position-spectral-origin}
Use the derived \((\mathrm{JOINT}_x^{\rm der})\) together with all hypotheses of
Theorem~\ref{p3:thm:local_spectral_asymptotics} and remain in its local
semiclassical domain.  Let
\(U_0\Subset U\) be the interior unit-localizer domain of
Definition~\ref{p3:def:analytic-localizer}.  Let \(\Phi\) denote the retained
bosonic local fields and \(\Psi\) any retained fermionic local fields in one
fluctuated position-sector Dirac datum \(D_{x,A}(g,\Phi)\).  Write
\(S^{\rm raw}_{F,U}\) for its unweighted local fermionic quadratic form and set
\[
 S^{\rm role}_{F,U}:=\frac12S^{\rm raw}_{F,U}.
\]
Canonical normalization of \(\Psi\), when performed, is applied afterward and
uniformly to this form and its interactions.  Define the effective position
functional
\begin{equation}
\label{p3:eq:common-position-spectral-action}
 S^{\rm common}_{x,U}[g,\Phi,\Psi]
 :=S^{\rm full}_{U,{\rm met}}
   +\frac12\!\left(S^{\rm raw}_{x,U}
                  -S^{\rm raw}_{x,U,{\rm met}}\right)
   +S^{\rm role}_{F,U}.
\end{equation}
Thus the derived Born pullback supplies one full copy of every pure-metric
scalar-density channel, while every nonmetric raw position term and the
pre-rescaling fermionic form retain the position-role coefficient inherited
from \(\Dcl\).  Then the
metric and retained fields are entries of one operator datum, its derivative
expansion can be decomposed by invariant type as
\[
 S^{\rm common}_{x,U}
 =S_U^{\rm grav}[g]+S^{\rm phys}_{{\rm RP},U}[g,\Phi,\Psi]
  +S_U^{\rm higher}[g,\Phi,\Psi],
\]
and this bookkeeping decomposition does not constitute independently weighted
geometry and matter actions.  Independent species-dependent gravitational
charges, source projectors, or spectral profiles would enlarge the declared
\(\Dcl+(\mathrm{JOINT}_x)\) datum.  Theorem~\ref{p3f:thm:v174-generated-joint-gauge-phase-space-action} supplies
the Lorentzian variational role of the full local parity-even Yang--Mills
constituent.  Theorems~\ref{p3f:thm:v175-generated-joint-einstein-action}
and~\ref{p3f:thm:v175-native-local-scalar-action} supply, on their additional
analytic first-jet and interior domains, the Einstein--Hilbert action germ and
the scalar-dependent canonical-dimension-at-most-four action.  Their common
positive result is the restricted action-bearing joint-Dirac package
\((\mathrm{ACT}_{{\rm J,EH+YM+\Phi},U_0}^{\rm loc,der})\); it is not a claim that
every bosonic heat coefficient through canonical dimension four has acquired a
Lorentzian variational role.  In particular, the positive scalar-Laplacian
Einstein coefficient inside \(S^{\rm full}_{U,{\rm met}}\) remains a
resolution/readout coefficient.  The physical local Einstein action uses the
negative Euclidean coefficient of the joint-Dirac owner.  Their raw
coefficient comparison is evaluated downstream once the common-heat carrier is
closed; it is not an equality premise.

For stationarity and Hilbert-tensor conclusions fix a claim \(C\) and its complete owner (Definition~\ref{p3a:def:claim-complete-action-owner}). The residual notation denotes only the consumed, still-unrealized sectors of that claim:
\[
 (\mathrm{ACT}_{C}^{\rm res}) :=\{\text{required unresolved action sectors of }C\}.
\]
Candidate entries are the required fermionic reality/phase, exact nonlocal action, offset/Majorana action and noncanonical/global boundary sectors. The local pure-metric finite-jet real forms are derived by Theorem~\ref{p3f:thm:finite-local-bosonic-clock-real-form} on its common-owner domain; outside that domain the specific continuation requirement remains. No universal missing-action premise is imposed on claims that consume none of these sectors.
Section~\ref{p3f:sec:v169-minimal-complete-action} assembles the inherited
nonlinear gravity coefficient, its canonical non-null endpoint completion,
and an explicit native free Dirac source on their declared domains. This
scoped variational completion does not discharge
\((\mathrm{ACT}_{C}^{\rm res})\) for sectors absent from that calculation.
The later source-native common-heat carrier evaluation fixes the raw
cross-owner transport rather than an equality premise:
\[
 \boxed{
 \UEHjoint^{\rm raw}=96\Lambda_{\rm EH}^2,
 \qquad
 \kappa_J=12\Lambda_{\rm EH}^2.}
\]
Common rooted provenance alone would not have supplied this relation; it is
fixed only after the generated carrier rank and ordinary trace are known.
This raw transport is not a constituent of
\((\mathrm{ACT}_{C}^{\rm res})\).  Framework v2.88 fixes the primitive
Newton/Planck curvature at the pre-amplification momentum-curvature owner
\(\mathcal K_p=6/\Rmom^2,\qquad U=(3/\pi)\mathcal K_p\quad\text{only under }(\mathrm{CELL}_{K\to U})\); the represented coefficient \(U_J^{\rm raw}=96U\) is a
distinct action readout, not an alternative curvature owner.  Write
\[
 (\mathrm{ACT}_{\rm cl})
 :=(\mathrm{ACT}_{{\rm J,EH+YM+\Phi},U_0}^{\rm loc,der})
   +(\mathrm{ACT}_{C}^{\rm res}).
\]
Each residual constituent may be supplied by a direct Lorentzian realization
or a specified continuation preserving field identity, overall sign,
coefficient, and variational role.  If \((\mathrm{SOLHIST}_{U}^{\rm phys})\) is also supplied,
then for variations compactly supported in \(U_0\), metric variation of this
Lorentzian common action gives
\begin{equation}
\label{p3:eq:common-action-stationarity}
 \frac{\delta S_U^{\rm grav}}{\delta g^{\mu\nu}}
 +\frac{\delta S^{\rm phys}_{{\rm RP},U}}{\delta g^{\mu\nu}}
 +\frac{\delta S_U^{\rm higher}}{\delta g^{\mu\nu}}=0,
\end{equation}
with
\[
 T^{\rm RP}_{\mu\nu}
 :=-\frac{2c_{\rm op}}{\sqrt{-g}}\frac{\delta S^{\rm phys}_{{\rm RP},U}}
 {\delta g^{\mu\nu}}.
\]
Equivalently,
\(\delta S^{\rm phys}_{{\rm RP},U}
=-(2c_{\rm op})^{-1}\int_{U_0}\sqrt{-g}\,
T^{\rm RP}_{\mu\nu}\delta g^{\mu\nu}\,d^4x\).  This convention is used
throughout and is the one compatible with the displayed \(c_{\rm op}^{-4}\)
source coefficient.
These equalities are local derivative-expansion equations on \(U_0\); the exact
spectral trace may carry the controlled nonlocal remainders already separated
in Theorem~\ref{p3:thm:local_spectral_asymptotics}.
\end{theorem}

\begin{proof}
The nonzero represented connection and single-datum ancestry are the
generated content of \((\mathrm{JOINT}_x^{\rm der})\) by
Theorem~\ref{p3f:thm:generated-fibre-spin-joint-operator}.
Theorem~\ref{p3f:thm:joint-same-datum-relation-closure} then derives the
common operator restrictions, curvature, and profile on the generated
\(\Dcl\)-typed gravitational process.  Current
Paper~III-B first forms the role-valued sector trace and then applies
\(\lambda_{\rm B}\), giving coefficient \(1/2\) to the position output.
Paper~III-E's natural scalar-density theorem replaces the two-sector-combined
pure-metric family by \(S^{\rm full}_{U,{\rm met}}\), which proves
\eqref{p3:eq:common-position-spectral-action} without doubling retained matter,
bundle, or fermionic invariants.  The local
heat-kernel coefficients are joint invariants of the metric and retained
bundle/connection fields, and any fermionic quadratic form uses the same
covariant operator.  Sorting the single Euclidean functional by local invariant type gives the
bookkeeping decomposition; independent weights or profiles would constitute a
different datum.  None of this occurrence bookkeeping implies stationarity.
Under \((\mathrm{ACT}_{\rm cl})\) and the claim-local \((\mathrm{SOLHIST}_{U}^{\rm phys})\) token,
\(\varphi_U=1\) near \(\overline{U_0}\), so compactly supported variations in
\(U_0\) produce no derivative-of-localizer terms.  Ordinary metric variation
then gives the stationarity identity and Hilbert tensor there.
\end{proof}

\begin{theorem}[Clock-hermitizing line and local uninserted fermion jet]
\label{p3f:thm:v177-local-clocked-fermion-jet}
Restrict Theorem~\ref{p3:thm:common-position-spectral-origin} to a local
clocked Lorentzian branch on which the complex four-dimensional spinor factor
carries the irreducible Clifford representation $c$ of the generated
Lorentzian coframe.  A Clifford-hermitizing matrix $A_{\rm h}$ satisfies
\[
 A_{\rm h}^\dagger=A_{\rm h},\qquad
 (A_{\rm h}c(v))^\dagger=A_{\rm h}c(v)
 \quad\text{for every real tangent vector }v.
\]
The hermitizing matrices form one nonzero real line.  If $n_{\rm clk}$ is the
admissible future unit clock normal, its sign is fixed by demanding that the
clock-normal Hilbert form
\[
 h_{n_{\rm clk}}(\psi,\phi)
 :=\psi^\dagger A_{\rm h}c(n_{\rm clk})\phi
\]
be positive definite.

On the \emph{uninserted} local classical branch---meaning that no additional
positive left/right internal wavefunction operator is inserted beyond the
existing finite Hilbert pairing---the complete fermionic quadratic form of
\eqref{p3:eq:common-position-spectral-action} has one common positive
coefficient $c_F$ on its kinetic and Yukawa blocks.  For every generated
Dirac/Yukawa member $D_Y$ its same-action local jet is therefore
\[
 \boxed{
 \mathfrak J_F(D_Y)
 =(\operatorname{Owner},\operatorname{Occ},Z_L,D,Z_R)
 =(\operatorname{Owner},\operatorname{Occ},
   c_FI_L,c_FD_Y,c_FI_R).}
\]
Consequently canonical fermion normalization gives
\[
 \boxed{\widehat D_Y
 =Z_L^{-1/2}DZ_R^{-1/2}=D_Y.}
\]
The theorem fixes the local classical kinetic pairing and the relative
placement of an already occurring Yukawa member.  It does not select that
member and does not establish a global determinant/Pfaffian, Majorana real
structure, PCT theorem, anomaly cancellation beyond the already proved
finite statements, or a regulated quantum measure.
\end{theorem}
\begin{proof}
If $A_{\rm h}$ and $B_{\rm h}$ hermitize the same irreducible spinor
Clifford representation, their defining identities imply that
$A_{\rm h}^{-1}B_{\rm h}$ commutes with every $c(v)$.  The complex Clifford
image is irreducible, so Schur's lemma gives
$B_{\rm h}=\lambda A_{\rm h}$.  Hermiticity forces $\lambda\in\mathbb R^*$.
The future clock normal selects the sign for which
$A_{\rm h}c(n_{\rm clk})$ is positive.  On the tensor product with the
finite internal carrier, the commutant contains internal endomorphisms; the
uninserted hypothesis retains the already generated identity pairing there
rather than promoting an arbitrary commutant element to new data.

Equation~\eqref{p3:eq:common-position-spectral-action} multiplies the entire
raw local fermionic quadratic form by one position-role coefficient before
canonical normalization.  Therefore the left kinetic block, the Yukawa block
and the right kinetic block inherit the same positive scalar $c_F$, giving
$(Z_L,D,Z_R)=c_F(I_L,D_Y,I_R)$.  The two inverse square roots cancel that
common scalar exactly.
\end{proof}

\begin{corollary}[Independently generated finite member needs no duplicate source occurrence]
\label{p3f:cor:v1122-root-generated-fermion-member-handoff}
Let $D_Y$ be a specific finite Dirac/Yukawa member which is itself the required
output of a licensed finite geometric constructor on an occurring root, and
assume its owner is the generated finite Dirac/Yukawa slot consumed by
Theorem~\ref{p3f:thm:v177-local-clocked-fermion-jet}.  Then on the local
clocked uninserted classical branch no second occurrence token is required for
that same member.  Paper~II-A generated-output/no-idle co-realization and the
fermion-jet constructor give
\[
 \boxed{\mathfrak J_F(D_Y)
 = (\operatorname{Owner},\operatorname{Occ},
 c_F I_L,c_FD_Y,c_F I_R),\qquad \widehat D_Y=D_Y.}
\]
This corollary does not select a member of an unresolved owner family.  It
applies only after the specific member has independent licensed root
provenance.
\end{corollary}

\begin{remark}[Inserted wavefunction and quantum firewall]
\label{p3f:rem:v177-fermion-insertion-quantum-firewall}
A separately generated positive internal insertion $K_L\oplus K_R$ would
change the jet to nontrivial $Z_L,Z_R$ and must be normalized explicitly; it
is not licensed by the spinor hermitizing-line theorem.  Likewise the local
bilinear real form does not determine the phase or global definition of a
fermion determinant or Pfaffian.  Those are distinct owner/domain questions.
\end{remark}


\section{The complete local bilinear parent and its missing coefficient}
\label{p3f:sec:v246-bilinear-parent}
\begin{proposition}[Coefficientwise local common-parent differentiation]
\label{p3f:prop:v246-bilinear-parent-jets}
On a fixed finite local window, let $q$ contain the retained commuting
configuration/source coefficients and expand in independent ordered odd
coordinates $\bar\psi_a,\psi_b$. Suppose the already constituted classical
owner is quadratic in these coordinates and has the form
\[
 S(q,\bar\psi,\psi)=A(q)+c_F\sum_{a,b}\bar\psi_a\mathcal D_{ab}(q)\psi_b,
 \qquad c_F>0.
\]
Here $A$ includes all bosonic, source and boundary constituents consumed
on this chart, and $\mathcal D$ is the actual finite coefficient of the
same local Dirac form. With $[\bar\psi_a\psi_b]$ denoting extraction of
the coefficient in this declared ordering, its complete commuting two-jet
and mixed odd coefficients obey
\begin{align*}
 \partial_i S&=\partial_i A+c_F\bar\psi(\partial_i\mathcal D)\psi,\\
 \partial_i\partial_j S&=\partial_i\partial_j A
           +c_F\bar\psi(\partial_i\partial_j\mathcal D)\psi,\\
 [\bar\psi_a\psi_b]S&=c_F\mathcal D_{ab},\\
 [\bar\psi_a\psi_b]\partial_i S&=c_F\partial_i\mathcal D_{ab},\qquad
 [\bar\psi_a\psi_b]\partial_i\partial_j S
      =c_F\partial_i\partial_j\mathcal D_{ab}.
\end{align*}
Thus the scalar--fermion and connection--fermion vertices of this parent
are derivatives of the same $\mathcal D(q)$; they are not independently
chosen sector terms. In particular, at zero odd background the commuting
Hessian is $D^2A$, with no direct bilinear contribution, and does not
by itself determine the marked coefficient $\mathcal D(q)$.
\end{proposition}
\begin{proof}
Expand the action in the fixed ordered Grassmann basis and differentiate
its ordinary coefficients. Coefficient extraction commutes with every
commuting partial derivative, giving all formulas. Setting the odd
coordinates to zero eliminates every bilinear from the commuting Hessian.
This avoids an implicit change of left/right Grassmann derivative convention.
\end{proof}
\begin{corollary}[Exact conditional discharge of kinetic normalization]
\label{p3f:cor:v246-canonical-member-handoff}
On the uninserted local clock branch of
Theorem~\ref{p3f:thm:v177-local-clocked-fermion-jet}, once the actual
charged chiral coefficient $\mathcal D_{LR,s}(q_*)$ has independent
source incidence, its mass coefficient is exactly
\[
 M_s=(c_FI_L)^{-1/2}(c_F\mathcal D_{LR,s})(c_FI_R)^{-1/2}
     =\mathcal D_{LR,s}(q_*).
\]
Consequently the remaining charged contractions consume this one actual
coefficient. They require no second kinetic-scale calculation on this
branch. Specifying only $A$ and its exact heat derivatives does not supply
the coefficient: all actions
$A+c_F\bar\psi(\mathcal D+T)\psi$, with $T$ a smooth typed coefficient,
have the same zero-odd-background commuting restriction and Hessian.
This last family proves information loss of that restriction; it does not
assert that any nonzero $T$ has framework provenance. If a joint spectral
constructor ties $A$ to $\mathcal D$, changing $\mathcal D$ can also
change $A$; this family is not a claim of equality of complete
source-dependent joint spectral actions.
\end{corollary}
\begin{proof}
The inherited positive scalar cancels on the two kinetic legs. The family
statement follows from the proposition. Selecting its physical member still
requires the source constructor, including $\mathsf{Cert}_{\rm hier}$ on
Paper~V-A's certified charged packet domain.
\end{proof}


\begin{proposition}[Conditional same-action two-frame chiral normalization]
\label{p3f:prop:v221-relative-chiral-comparison-frame}
Let $j_0,j_1:G\to\mathscr E$ be any two isometric realizations of
one Hermitian family source with $j_0^*j_1=tI_G$, where
$t=\sqrt{24/25}$ is supplied by the geometric Hurwitz frame of
Paper~I-C, Corollary~\ref{p1b:cor:v221-hurwitz-principal-angle}.
Consider a \emph{specified} port-augmented quadratic chiral action
in which the left and right frames are $j_{a_L},j_{a_R}$, with
$a_L,a_R\in\{0,1\}$, and an underlying family source $D_0$ acts
trivially on the comparison port.  Its same-action blocks obey
\[
 Z_L=c_Fj_{a_L}^*j_{a_L},\qquad
 Z_R=c_Fj_{a_R}^*j_{a_R},\qquad
 D=c_F(j_{a_L}^*j_{a_R})D_0.
\]
After canonical normalization,
\[
 \boxed{Z_L^{-1/2}DZ_R^{-1/2}
   =\begin{cases}D_0,&a_L=a_R,\\
                   tD_0,&a_L\ne a_R.
     \end{cases}}
\]
The same conclusion follows using the unnormalized rooted
$Q_{r_0}F$ and its inherited kinetic Gram $\eta_H I$.
\end{proposition}
\begin{proof}
Each frame is an isometry.  Therefore both kinetic Grams equal
$c_FI$.  Same-frame cross Grams are $I$ and opposite-frame
cross Grams are $tI$, giving the result.  In the unnormalized
comparison the common rooted Gram factors cancel on two rooted
chiral legs, while a single rooted leg contributes $\sqrt{\eta_H}$.
\end{proof}
\begin{remark}[No free source-selection theorem from a common kinetic trace]
\label{p3f:rem:v221-root-port-not-automatic-mass}
This proposition classifies \emph{specified port-augmented actions}.
The uninserted jet of
Theorem~\ref{p3f:thm:v177-local-clocked-fermion-jet} contains an
already generated $D_Y$ and does not itself assign $j_0$ or $j_1$
to any physical chiral leg.  In particular the same rooted geometry
permits both $j_L=j_R=j_1$, which yields no mass correction, and
$j_L=j_0,j_R=j_1$, which yields one factor $t$.  Thus no
physical charged comparison exponent follows from the uninserted jet
alone, and universal rescaling of kinetic and mass terms cancels.
\end{remark}


\begin{theorem}[Generated stationary-solution family under the classical-action bridge]
\label{p3f:thm:v173-stationary-solution-family}
On a complete \((\mathrm{ACT}_{\rm cl})\) branch of
Theorem~\ref{p3:thm:common-position-spectral-origin}, let \(S_L\) denote the
resulting Lorentzian common action on the generated local configuration role
\(\mathcal C_{\rm loc}(U)\).  Declare before evaluating a configuration the
certified variational word
\[
\begin{gathered}
 [S_L]\longmapsto dS_L\longmapsto\operatorname{Eq}(dS_L,0)\\
 \longmapsto\mathfrak M_{\rm stat,U}([S_L])
 :=\frac{\substack{\{q\in\mathcal C_{\rm loc}(U):\ dS_L|_q=0,\\
 \mathfrak D^{\rm GP2}_{\rm joint}(S_L,q)=0\}}}
 {(\cong_{\rm typed},\sim_{\rm gauge},\sim_{\rm cal})}.
\end{gathered}
\]
where the quotient relations and target family sort are frozen parts of the
word and the presentation lies in the finite/exact domain required by
Theorem~\ref{p2a:thm:generated-exact-solution-moduli}.  Then:
\begin{enumerate}[label=(\roman*),leftmargin=2em]
\item if \(\mathfrak M_{\rm stat,U}([S_L])\neq\varnothing\), it is generated
      role-level family content and is co-realized with the complete action
      root; if it is empty, the requested stationary role is FRC-obstructed;
\item a unique GP3-least class is retained when it exists, and otherwise the
      complete unresolved stationary family is retained without selecting a
      member;
\item every licensed solution-natural response, Hessian, vertex, restriction,
      or readout constructor acts on the whole family and retains its output as
      a family or graph; a value constant on every surviving class is common
      generated content;
\item a particular physical solution is an additional instantiational token
      \((\mathrm{SOLHIST}_{U}^{\rm phys})\), not an antecedent for constructing
      the Euler--Lagrange map or its complete zero locus.
\end{enumerate}
Denote the family-level conclusion by
\[
 \boxed{(\mathrm{STATFAM}_x^{\rm der}).}
\]
The theorem does not prove nonemptiness, stability, a broken vacuum, an actual
solution, a probability/amplitude measure on the family, or a global/continuum
extension.
\end{theorem}

\begin{proof}
Differentiation is a generated operation on the already adopted variational
action, and the zero equalizer and licensed quotient are exactly the
predeclared exact-solution constructor word.  Theorem~\ref{p2a:thm:generated-exact-solution-moduli}
therefore gives the empty/nonempty alternatives, GP3 retention, common-readout
inheritance, and rooted co-realization.  Corollary~\ref{p2a:cor:v173-whole-family-parent-reuse}
permits any later certified constructor to consume the complete family node
without evaluating a fibre.  Evaluating at \(q_{\rm world}\) instead inserts a
member and is consequently the separate instantiational token stated in~(iv).
\end{proof}

\begin{corollary}[Classical stationary-family provider for the cascade]
\label{p3f:cor:v257-classical-family-provider}
At a cascade node which is a complete finite/exact
$(\mathrm{ACT}_{\rm cl})$ action output, assume the target, equalizer,
quotient and retained joint admissibility are those of
Theorem~\ref{p3f:thm:v173-stationary-solution-family}.
Its certified word
\[
 [S_L]\longmapsto dS_L\longmapsto\operatorname{Eq}(dS_L,0)
 \longmapsto\mathfrak M_{\rm stat,U}([S_L])
\]
supplies that node's stationary-family constructor in
Theorem~\ref{p2a:thm:rooted-generated-parent-cascade}.
On the nonempty admissible branch this family is a generated descendant
without a second independent occurrence token. A physical history
remains a separate instantiation, and unrelated parent roles require
their own family word.
\end{corollary}
\begin{proof}
Compose the licensed action-output word with III-F's licensed family
word, including its frozen quotient and target. Use licensed output
retention and reuse. The exact-solution-moduli theorem applies to that
composite. No generic solution constructor is inferred for another role.
\end{proof}


\begin{corollary}[Local stationary-family host and covariant history identity]
\label{p3f:cor:v173-local-stationary-family-host}
On the nonempty branch of
Theorem~\ref{p3f:thm:v173-stationary-solution-family}, the minimal closure
\[
 \mathcal P_{U}^{\rm sol,fam}
 :=\operatorname{Cl}^{\min}_{\rm GP2}
 \bigl(\mathfrak D_J,[S_L],dS_L,
       \mathfrak M_{\rm stat,U}([S_L])\bigr)
\]
is a generated and co-realized local structural host.  A member
\(q=(g,\Phi,\Psi)\) is a tuple of sections over the whole spacetime domain
\(U\); hence an actual member is already the covariant on-shell history on
\(U\).  No second history occurrence is required for that same tuple.
This identity does not supply Cauchy data, prove nonlinear well-posedness,
select a quantum state or preparation, or extend the local solution globally.
\end{corollary}

\begin{corollary}[Local discharge of the historical Einstein action-role bridge]
\label{p3f:cor:act-cl-discharges-act-eh}
On the v1.75 analytic clock--ADM first-jet and compact-support interior domain,
Theorem~\ref{p3f:thm:v175-generated-joint-einstein-action} gives
\begin{equation}
\label{p3f:eq:act-cl-eh-restriction}
 \operatorname{Res}_{\rm EH}
 \bigl(\mathrm{ACT}_{{\rm J,EH+YM+\Phi},U_0}^{\rm loc,der}\bigr)
 =\mathrm{ACT}_{{\rm J,EH},U_0}^{\rm int,der}.
\end{equation}
Thus the action-role component of the claim-scoped historical condition
\(\mathrm{ACT}_{\rm EH}^{\rm cl}\) used in Paper~III-E is
\textsc{derived-later} on this local interior domain without waiting for the
fermionic, exact-nonlocal, offset, or noncanonical-boundary completion of the
full action.  The later common-heat carrier evaluation fixes the raw relation to
Paper~III-E's scalar-resolution scale as \(U_J^{\rm raw}=96U\).  The pre-amplification momentum curvature is independently fixed; its identification with the primitive Newton/Planck curvature is conditional on the evaluated cell port at \(\mathcal K_p=6/\Rmom^2,\qquad U=(3/\pi)\mathcal K_p\quad\text{only under }(\mathrm{CELL}_{K\to U})\); no \((\mathrm{GRAV\!-\!DESC}_{J\to\rm prim})\) theorem is needed for that identification.  Outside the
printed domain the historical conditional statement remains properly typed.
\end{corollary}

\begin{proof}
Theorem~\ref{p3f:thm:v175-generated-joint-einstein-action} constructs the
Lorentzian Einstein--Hilbert interior action germ itself, with its sign,
coefficient owner, and variational role.  Restriction of the generated local
package to that invariant summand is therefore exactly the displayed
certificate.  No converse from a later response, no Newton calibration, no
physical boundary completion, and no unproved full-action completion is used.
\end{proof}

\begin{definition}[Certified action/contact and local-excess response domain]
\label{p3f:def:v2211-certified-response-domain}
The \emph{certified action/contact branch} means that the actual
action-density coefficient has passed the normalized root and zero
connection-defect test of
Theorem~\ref{p3b:thm:v2211-action-contact-certificate}.  Only on this
branch is its coefficient dual identified with the contact determinant
dual used by CRP4.  The \emph{certified local-excess branch} additionally
passes the actual fixed-depth Euler-excess character test of
Theorem~\ref{p3d:thm:v2211-excess-character-certificate}.
Neither certificate is inferred from the action's occurrence, scalar
secant values, an abstract one-dimensional character or an already
reconstructed \(f(R)\).

Every same-action use below of the CRP4 coefficient connection imports the
first certificate.  Every physical use of the printed fixed-depth
\(\chi_7\) and its dressed descendant additionally imports the second.
Conditional reconstruction from those supplied coefficients remains exact
whether or not their actual physical ports have yet been evaluated.
The local joint-Dirac action, finite source variations and
Noether/Cauchy/Schur results that consume neither identification retain
their own independent domains.
\end{definition}

\begin{corollary}[Generated action response under configuration and coefficient-family variation]
\label{p3f:cor:generated-action-response-variation}
Let \(S_\star\) be the occurring scalar action of Paper~III-B.  The local
field variables retain the minimal configuration tangent of
Theorem~\ref{p3f:thm:minimal-configuration-tangent-vdom}.  Independently,
let \(B_{\rm sd}\) be the already generated scheme--depth base of
Paper~III-C and regard the action as the parameterized family on
\[
 \widetilde{\mathcal C}:=B_{\rm sd}\times\mathcal C_{\rm loc}(U),
 \qquad
 d_{\rm tot}=d_{\rm sd}+d_{\rm cfg}.
\]
Here \(d_{\rm cfg}\) varies only the local metric/field arguments, whereas
\(d_{\rm sd}\) differentiates only the generated coefficient family.
Paper~III-B Theorem~\ref{p3b:thm:generated-action-response-owner}
constructs the action-density dual and its initial section.
On the additional certified action/contact branch of
Definition~\ref{p3f:def:v2211-certified-response-domain},
\[
 \mathcal L_U=\mathcal L_{\det}^{-1},\qquad
 a_U=\operatorname{id},\qquad
 U=\Lambda_{\rm EH}^{2},\qquad\kappa_U=1.
\]
The coefficient-family differential \(d_{\rm sd}S_\star\), evaluated through
the frozen determinant pairing of
Definition~\ref{p3b:def:scalar-action-coefficient-compiler}, is the generated
physical response covector.  The local Euler--Lagrange equations continue to
use \(d_{\rm cfg}S_\star\); the two tangent sorts are not identified.  The
second variation along two generated rank-four coefficient directions has
owner
\[
 \operatorname{Sym}^2(\mathcal L_{1/4})
 \cong\mathcal L_{1/4}^{\otimes2}.
\]
Consequently \(\mathrm{ACT}_{\rm MRESP}\) and its bilinear response are
generated on the action line; their contact/PQTI identification uses the
displayed certificate and the actual rank-four source directions.
This corollary does not select the spectral profile,
a dressed Laurent member, a Lorentzian/OS continuation, stationarity, an
actual solution/history, or an empirical unit map.
\end{corollary}

\begin{proof}
The action-density owner, evaluation pairing and initial section are the
cited Paper~III-B constructor.  Its identification with the contact
determinant line follows only on the additional certificate domain: the
normalized root and vanishing relative connection defect transport the
same section throughout the connected parameter interval.  The product family
\(B_{\rm sd}\times\mathcal C_{\rm loc}(U)\) carries the ordinary product
differential, so varying a scheme/depth parameter does not add a local field
argument and varying a local field does not move the coefficient-family base.
Differentiating the already frozen coefficient slot therefore gives
\(d_{\rm sd}S_\star\) as a generated descendant of the same action parent,
while the usual variational calculus gives \(d_{\rm cfg}S_\star\).
Bilinearity of the generated second response gives the tensor-square owner.
No independent response or extra configuration variable is inserted.
Without the certificate the same product differentiation remains valid
on the action-density line, while its relative contact connection defect
must be retained; it is not set to zero by the existence of a dual pairing.
\end{proof}

\begin{theorem}[Coefficient-family evaluation response on the bare Born leaf]
\label{p3f:thm:born-coefficient-family-first-response}
Fix one scheme \(\sigma\) on the certified action/contact branch of
Definition~\ref{p3f:def:v2211-certified-response-domain}, and
write
\[
 L:=L_1^{(\sigma)},\qquad
 R(q):=Ue^{-Lq},\qquad
 C:=Ue^{-7L},\qquad
 B:=CU,
\]
with \(0<C<U\).  On the calibrated bare two-root leaf of Paper~III-D set
\[
 u_B:=\log\frac{R}{\sqrt B}
      =L\left(\frac72-q\right),
 \qquad
 \mathscr H_{C,U}(R):=R g_{\rm eff}(R)
 =\frac{R^2+B}{U+C}.
\]
Then the fixed-scheme specialization of Paper~III-C's fundamental
scheme--depth connection is
\begin{equation}
\label{p3f:eq:born-coefficient-connection}
 \boxed{\nabla_U^{\rm sd}=d_{\rm sd}+L\,dq=d_{\rm sd}-du_B,\qquad
        \nabla_U^{\rm sd}R=0.}
\end{equation}
Let \(\mu_R\in\mathcal L_{\det}\) be the horizontal dual section normalized
by
\(\operatorname{ev}_{\det}(R,\mu_R)=1\).  Then
\begin{equation}
\label{p3f:eq:born-response-scalarization}
 \boxed{
 \operatorname{ev}_{\det}(\mathscr H_{C,U},\mu_R)=g_{\rm eff},\qquad
 d_{\rm sd}g_{\rm eff}
 =\operatorname{ev}_{\det}(\nabla_U^{\rm sd}\mathscr H_{C,U},\mu_R).}
\end{equation}
Moreover
\begin{equation}
\label{p3f:eq:born-first-response-section}
 \boxed{
 \nabla_U^{\rm sd}\mathscr H_{C,U}
 =\frac{R^2-B}{U+C}\,du_B,\qquad
 \mathfrak f_{\rm nl}^{(1)}(R)
 :=\mathcal E\mathscr H_{C,U}(R)
 =\frac{R^2-B}{U+C},}
\end{equation}
where \(\mathcal E:=R\partial_R-1\).  Thus the generated nonlinear
first-response section has the same bare Hessian as the finite interpolation,
\[
 (\mathfrak f_{\rm nl}^{(1)})''=\mathscr H_{C,U}''=\frac{2}{U+C},
\]
but is fixed modulo only the parallel coefficient-line kernel
\begin{equation}
\label{p3f:eq:born-parallel-kernel}
 \boxed{\ker\nabla_U^{\rm sd}=\operatorname{span}_{\mathbb R}\{R\}.}
\end{equation}
\end{theorem}

\begin{proof}
At fixed \(\sigma\), Paper~III-C
Proposition~\ref{p3c:prop:full-scheme-depth-determinant-connection} gives
\(\nabla_U^{\rm sd}=d_{\rm sd}+L\,dq\), while the generated continuation
section is \(R(q)=Ue^{-Lq}\).  Since
\(du_B=-L\,dq\), equation~\eqref{p3f:eq:born-coefficient-connection}
follows and \(R\) is horizontal.  Preservation of the frozen evaluation
pairing gives the dual connection and a horizontal \(\mu_R\) with local
component proportional to \(R^{-1}\); the chosen normalization makes that
component exactly \(R^{-1}\) in the corresponding trivialization.

Paper~III-D Corollary~\ref{p3d:cor:two-scale-response-normalizations} gives
\(\mathscr H_{C,U}=Rg_{\rm eff}\).  Since multiplication of the horizontal
coefficient section \(R\) by the scalar readout \(g_{\rm eff}\) remains in
\(\mathcal L_U\), the frozen pairing gives
\(\operatorname{ev}_{\det}(\mathscr H_{C,U},\mu_R)=\mathscr H_{C,U}/R
=g_{\rm eff}\).  The covariant Leibniz rule and horizontality of \(\mu_R\)
then give the second identity in
\eqref{p3f:eq:born-response-scalarization}.

For every weight-one coefficient section \(F(R)\), one has
\(d_{\rm sd}R=R\,du_B\), hence
\[
 \nabla_U^{\rm sd}F
 =d_{\rm sd}F-F\,du_B
 =(RF_R-F)\,du_B
 =\mathcal EF\,du_B.
\]
Applying this to \(\mathscr H_{C,U}=(R^2+B)/(U+C)\) gives
\eqref{p3f:eq:born-first-response-section}.  The certified action/contact
identification \(a_U=\operatorname{id}\) lifts the scalar coefficient of this
same-action response back to the same determinant-dual coefficient line, so
no independent response owner is introduced.  Finally
\(\nabla_U^{\rm sd}h=0\) is \(Rh_R-h=0\), whose solutions on the connected
positive interval are exactly \(h=\lambda R\), proving
\eqref{p3f:eq:born-parallel-kernel}.
\end{proof}

\begin{remark}[Common-action/CRP4 dependency firewall]
\label{p3f:rem:common-action-crp4-cycle-firewall}
The two closure chains have different directions:
\[
 \mathrm{CRP4}/a_U
 \longrightarrow
 \left(C/U=\alpha_\sigma^{-7},\mathfrak f_{\rm sp}\right)
 \longrightarrow
 \text{the imported smooth }f(R)\text{ action},
\]
whereas the noncircular v1.75 action route is
\[
 (\mathrm{JOINT}_x^{\rm der})
 +\mathsf{Heat}_{J}^{\rm der}
 +(\mathrm{CAUS})
 \longrightarrow
 \mathrm{ACT}_{{\rm J,EH},U_0}^{\rm int,der}.
\]
This second chain constructs the local Einstein action-role bridge directly
from the joint Dirac owner; it does not wait for the fermionic, exact-nonlocal,
or noncanonical-boundary completion of \(\mathrm{ACT}_{\rm cl}\).  Reversing
from the already imported smooth \(f(R)\) response to infer its ancestry
remains circular, because that response was formed using the owner relation.
The noncircular candidate route occurs earlier:
\[
 S_\star+\operatorname{ev}_{\det}
 \longrightarrow
 \mathcal L_U=\mathcal L_{\det}^{-1}
 \longrightarrow dS_\star\longrightarrow d^2S_\star .
\]
Paper~III-B freezes the scalar-action coefficient compiler before response
evaluation.  Its actual coefficient must additionally pass
Theorem~\ref{p3b:thm:v2211-action-contact-certificate} before the displayed
identification with the contact determinant dual is available.
Then
Corollary~\ref{p3f:cor:generated-action-response-variation} transports its
outputs under ordinary variation.  The compiler alone constructs a dual
line but does not prove the missing representation equality.  Thus this
is an explicit test of CRP4 without using the already assembled response
action; the forbidden reverse inference remains forbidden.
\end{remark}

\begin{corollary}[Metric backreaction in the generated joint-EH coupling coordinate]
\label{p3:cor:common-action-metric-backreaction}
Assume the hypotheses of Theorem~\ref{p3:thm:common-position-spectral-origin}
and one claim-local \((\mathrm{SOLHIST}_{U}^{\rm phys})\) member.  For the local
bosonic truncation use the already derived
\((\mathrm{ACT}_{{\rm J,EH+YM+\Phi},U_0}^{\rm loc,der})\); if the sourced stress tensor
contains fermionic, exact-nonlocal, offset/Majorana, or noncanonical-boundary
terms, additionally supply the corresponding constituents of
\((\mathrm{ACT}_{C}^{\rm res})\).  Assume
\((\mathrm{REN}_0(\Lambda_{\rm ren}))\).

The generated joint Einstein coefficient \(\kappa_J>0\) defines
\[
 \UEHjoint:=8\kappa_J,
 \qquad
 G_{J,\mathrm{raw}}:=\frac{c_{\rm op}^{3}}{\hBorn\UEHjoint^{\rm raw}},
 \qquad \hBorn:=2\pi\hbar,
\]
as in Corollary~\ref{p3f:cor:v267-joint-eh-newton-readout}.  Thus no
observed Newton constant is an additional premise of this corollary.  The
dimensional physical action is \(\hbar\) times the generated dimensionless
Lorentzian action, and the Hilbert tensor is formed from the correspondingly
dimensional matter functional.

Under the exact alternative \((\mathcal E_{\ge4}=0)\), the interior domain
\(U_0\) obeys
\[
 \boxed{G_{\mu\nu}+\Lambda_{\rm ren}g_{\mu\nu}
 =\frac{8\pi G_{J,\mathrm{raw}}}{c_{\rm op}^4}T^{\rm RP}_{\mu\nu}.}
\]
The vacuum-term-free equation is the specialization \((\mathrm{REN}_0^0)\),
not a consequence of merely specifying the zero-order variation.  Under the
effective alternative \((\mathrm{IR}_{2\partial}(\varepsilon_*))\), one instead has
\[
 G_{\mu\nu}+\Lambda_{\rm ren}g_{\mu\nu}
 =\frac{8\pi G_{J,\mathrm{raw}}}{c_{\rm op}^4}T^{\rm RP}_{\mu\nu}
 +\mathcal E_{\mu\nu}^{\ge4},
 \qquad
 \|\mathcal E^{\ge4}\|_{\mathcal X}
 \le\varepsilon_*\,\|\mathcal E^{\rm EH}\|_{\mathcal X},
\]
in the declared normed field/domain class \(\mathcal X\).  If neither
higher-order alternative is imposed, the higher-curvature variation must be
retained without an error claim.  A comparison
\(G_{J,\mathrm{raw}}=G_N^{\rm obs}\) additionally consumes the physical
metric/source normalization of
Proposition~\ref{p3f:prop:v2211-newton-owner-matching}; once that
identification is proved, numerical comparison is a
metrological validation statement and is not used to construct
\(\kappa_J\), \(\UEHjoint\), or the action sign.
\end{corollary}

\begin{proof}
Theorem~\ref{p3f:thm:v175-generated-joint-einstein-action} supplies the
Lorentzian Einstein--Hilbert interior action role, its positive coefficient
ray, and its joint-Dirac ownership.  The algebraic coordinate identity
\(\hbar\kappa_J=c_{\rm op}^3/(16\pi G_{J,\mathrm{raw}})\) follows from
\(\UEHjoint^{\rm raw}=8\kappa_J\) and \(\hBorn=2\pi\hbar\); it is not an empirical
normalization premise.  A GHY term is available only under
Remark~\ref{p3f:rem:v175-canonical-dirac-boundary-completion}.  The separate
\((\mathrm{REN}_0(\Lambda_{\rm ren}))\) condition retains the complete
zero-derivative variation as \(\Lambda_{\rm ren}g_{\mu\nu}\).  Substitution
of the full diagonal EH variation from
Lemma~\ref{p3f:lem:diagonal-eh-full-variation} and the Hilbert definition into
\eqref{p3:eq:common-action-stationarity} gives the equation with the
higher-curvature residual.  The exact alternative sets it to zero; the
effective alternative applies its declared norm bound.
\end{proof}


\begin{definition}[Renormalized vacuum baseline and quantum excitation source]
\label{p3f:def:v2105-vacuum-excitation-source-split}
On any claim-local quantized matter/perturbation sector for which the
renormalized zero-order gravitational coordinate and the quantum Hamiltonian
are both defined, write the state-independent vacuum offset and the positive
excitation operator as
\begin{equation}
\label{p3f:eq:v2105-vacuum-excitation-hamiltonian-split}
 \boxed{\widehat H=E_{\rm vac}^{\rm ren}\mathbf1+\widehat H_{\rm exc},
 \qquad \widehat H_{\rm exc}\succeq0.}
\end{equation}
The metric-proportional state-independent part is owned by the already
separate zero-order coordinate \((\mathrm{REN}_0)\).  The state-dependent
remainder is the excitation Hilbert source, denoted
\(T_{\mu\nu}^{\rm exc}[\omega]\), obtained from the same claim-complete
action after that vacuum offset has been assigned to the zero-order
coordinate.  Thus a finite renormalization
\(E_{\rm vac}^{\rm ren}\mapsto E_{\rm vac}^{\rm ren}+\Delta E\) changes the
zero-order gravitational representative but does not change
\(\ker\widehat H_{\rm exc}\), the excitation occupation content, or any
positive-energy complex structure constructed from the quadratic Hessian.
This is a typing split, not a numerical solution of the absolute vacuum
renormalization problem.
\end{definition}

\begin{definition}[QROOT quantum-source-completeness bridge]
\label{p3f:def:v2105-qroot}
Let Paper~VI-A supply an occurring complete-cap fresh-root branch with its
same-action homogeneous datum and exact source-free root equations
\(\rho_0=p_0=0\).  The statement
\begin{equation}
\label{p3f:eq:v2105-qroot}
 \boxed{(\mathrm{QROOT}):\qquad
 \rho_0=p_0=0
 \ \Longrightarrow\ 
 \langle T_{\mu\nu}^{\rm exc}\rangle_0=0}
\end{equation}
is the claim that the fresh-root source-free condition is complete also with
respect to the renormalized quantum excitation Hilbert source of
Definition~\ref{p3f:def:v2105-vacuum-excitation-source-split}.  In particular,
on a homogeneous positive-energy block it implies
\(\operatorname{Tr}(\rho_0\widehat H_{\rm exc})=0\).

The bridge is distinct from \((\mathrm{REN}_0)\): QROOT concerns vanishing
excitation above the vacuum baseline, whereas \((\mathrm{REN}_0)\) owns the
absolute state-independent zero-order term.  The classical datum alone does
not prove QROOT; on the native complete-cap cyclic perturbation branch the
additional licensed neutral Fock constructor and source-response theorem close
it in Theorem~\ref{p3f:thm:v1128-qroot-neutral-constructor-closure}.
\end{definition}

\begin{remark}[QROOT is not a repeated mode-admission vacuum rule]
\label{p3f:rem:v2105-qroot-no-readmission}
The sharp spectral frontier resolves additional operational harmonics but is
not an ontological creation event for each mode.  If QROOT is proved, the
quantum state is fixed once on the fresh daughter branch and later finite
frontiers inherit restrictions of its same-action evolution.  No rule of the
form \(\rho_{m}\mapsto\rho_m\otimes|0\rangle\langle0|\) is introduced at each
resolution threshold.
\end{remark}

\subsection{Metric field equation of the imported smooth response}
\label{p3:sec:master-field-equation}

\begin{theorem}[General local metric-\(f(\scal)\) variation]
\label{p3f:thm:general-local-fR-variation}
Let \(g\in C^2_{\rm loc}\) be a nondegenerate Lorentzian metric and define
\[
 S[g,\psi]
 =\frac{c_{\rm op}^3}{16\pi G}
 \int_U\sqrt{-g}\,[f(\scal)-2\Lambda_{\rm ren}]\,d^4x
 +S_{\rm m}[g,\psi]
\]
whenever the displayed action, or equivalently all compactly supported
relative action differences used below, is meaningful.  For every admitted
inverse-metric variation
\(g_\varepsilon^{\mu\nu}=g^{\mu\nu}+\varepsilon h^{\mu\nu}\) with
\(h\in C_c^\infty(U_0)\) and \(K=\operatorname{supp}h\), assume that for
sufficiently small \(\varepsilon\):
\begin{enumerate}[label=(\roman*),leftmargin=2.2em]
\item \(g_\varepsilon\) remains a \(C^2\) nondegenerate metric and the action
      difference is defined;
\item some neighborhood \(V\supset K\) and open interval \(I\) satisfy
      \(\scal[g_\varepsilon](V)\subset I\) and \(f\in C^1(I)\);
\item \(\phi:=f'(\scal[g])\in C^2(V)\).
\end{enumerate}
Use the single matter convention
\begin{equation}
\label{p3f:eq:matter-tensor-convention}
 (\mathrm{T\!-!CONV}):\qquad
 T_{\mu\nu}:=-\frac{2c_{\rm op}}{\sqrt{-g}}
 \frac{\delta S_{\rm m}}{\delta g^{\mu\nu}},
 \qquad
 \delta S_{\rm m}
 =-\frac1{2c_{\rm op}}\int_U\sqrt{-g}\,
 T_{\mu\nu}\delta g^{\mu\nu}\,d^4x.
\end{equation}
Assume \((\mathrm{VDOM}_x^{\min})+(\mathrm{SOLHIST}_{U}^{\rm phys})
+(\mathrm{REN}_0(\Lambda_{\rm ren}))\), with compactly supported
metric variations in \(U_0\Subset U\), and assume either that this is a
directly occurring Lorentzian action
\((\mathrm{LOR}_{\rm dir})\), or that the displayed classical local action
has been continued with a fixed sign/field convention
\((\mathrm{CONT}_{\rm cl})\).  Then on \(U_0\)
\[
 \boxed{
 f'(\scal)R_{\mu\nu}-\frac12f(\scal)g_{\mu\nu}
 +(g_{\mu\nu}\Box-\nabla_\mu\nabla_\nu)f'(\scal)
 +\Lambda_{\rm ren}g_{\mu\nu}
 =\frac{8\pi G}{c_{\rm op}^4}T_{\mu\nu}.}
\]
The theorem is valid for every response in this local classical domain and
does not require \((\mathrm{MRESP})\).  The stronger global conditions
\(g\in C^4_{\rm loc}\), \(\scal[g]\in C^2_{\rm loc}(U;I)\), and
\(f\in C^3(I)\) are convenient sufficient conditions, not necessary ones.
Only \(f\in C^1\) and \(f'(\scal[g])\in C^2\) near each variation support
enter the proof.
The classical continuation tag is weaker than
an Osterwalder--Schrader or interacting quantum reconstruction claim.
\end{theorem}

\begin{proof}
Outside \(K\), \(g_\varepsilon=g\), so the relative action difference cancels
identically; no regularity of \(f\) beyond the neighborhood \(V\) is used.
The standard metric variation gives
\(\delta\scal=R_{\mu\nu}\delta g^{\mu\nu}
+g_{\mu\nu}\Box\delta g^{\mu\nu}
-\nabla_\mu\nabla_\nu\delta g^{\mu\nu}\).  Two integrations by parts are
licensed by compact support in \(U_0\); the matter variation defines
\(T_{\mu\nu}\) through \eqref{p3f:eq:matter-tensor-convention}.  The
variation of \(-2\Lambda_{\rm ren}\) supplies
\(\Lambda_{\rm ren}g_{\mu\nu}\).  Comparing the prefactor
\(c_{\rm op}^3/(16\pi G)\) with the factor
\((2c_{\rm op})^{-1}\) in the matter variation gives precisely
\(8\pi G/c_{\rm op}^4\).  Stationarity gives the displayed equation.
\end{proof}

\begin{remark}[Weak metric-response branch]
If the classical regularity above is unavailable, the same equation may be
read distributionally only after separately assuming
\(\phi:=f'(\scal)\in H^2_{\rm loc}\), requiring every remaining product to
lie in \(L^1_{\rm loc}\), and declaring the compactly supported test-tensor
space and boundary conditions.  A bare \(f\in C^2\) hypothesis licenses
neither this weak branch nor the pointwise classical equation.
\end{remark}


\begin{theorem}[Joint TERM--REN zero-order redundancy]
\label{p3f:thm:term-ren-zero-order-redundancy}
In the action domain of Theorem~\ref{p3f:thm:general-local-fR-variation}, let
\(F_C=f(C)\), \(Z_C=f'(C)\), and let \(\Lambda_{\rm ren}\) be the
zero-derivative member.  For every curvature-valued
\(\delta\in\mathcal L_{\scal}\), the simultaneous
replacement
\[
 F_C\longmapsto F_C+\delta,
 \qquad
 \Lambda_{\rm ren}\longmapsto\Lambda_{\rm ren}+\frac{\delta}{2},
\]
with the Hessian and slope data fixed, leaves the complete gravitational action
unchanged.  Hence the action-level zero-order endpoint coordinate is
\[
 \boxed{\overline F_C:=F_C-2\Lambda_{\rm ren}},
\]
and \(F_C\) and \(\Lambda_{\rm ren}\) are not independent physical
coordinates once they are realized in the same action.  This statement does
not alter Paper~III-E's raw response-function endpoint atlas before REN is
composed with it.
After a curvature unit is fixed, \(\delta\) may of course be represented by a
real number; the typed shift itself has the same physical dimension as
\(f(\scal)\) and \(\Lambda_{\rm ren}\).
\end{theorem}

\begin{proof}
Changing the endpoint value by \(\delta\) changes the reconstructed response
by the additive constant \(f(R)\mapsto f(R)+\delta\) while leaving
\(f'\) and \(f''\) unchanged.  Therefore
\[
 [f(R)+\delta]-2\left(\Lambda_{\rm ren}+\frac\delta2\right)
 =f(R)-2\Lambda_{\rm ren}
\]
identically.  The matter action is untouched, so the full action is identical.
\end{proof}

\begin{proposition}[Legendre potential of the complete same action]
\label{p3f:prop:complete-action-legendre-potential}
In the local metric-action domain above, keep \(\Lambda_{\rm ren}\)
constant on the branch and fix \(U>0\). Define
\[
 \overline f=f-2\Lambda_{\rm ren},\qquad
 \overline{\mathcal F}=U\overline f,\qquad
 \mathcal V(R)=R\overline{\mathcal F}'(R)-\overline{\mathcal F}(R).
\]
The complete coefficient and \(\mathcal V\) are invariant under the joint
TERM--REN redistribution. On a local interval where
\(\overline{\mathcal F}''\ne0\), the same action has the equivalent
auxiliary-field representation
\[
 \overline{\mathcal F}(r)+\overline{\mathcal F}'(r)(R-r)
   =XR-V(X),\quad
 X=\overline{\mathcal F}'(r),\quad
 V(X)=rX-\overline{\mathcal F}(r).
\]
On the auxiliary equation, \(r=R\) and \(V(X(R))=\mathcal V(R)\).
In particular the potential is generated from a supplied action, rather than
introduced as an independent matter field.
\end{proposition}
\begin{proof}
The complete coefficient is unchanged by
Theorem~\ref{p3f:thm:term-ren-zero-order-redundancy}; its Euler remainder is
therefore unchanged as well. Variation with respect to \(r\) gives
\(\overline{\mathcal F}''(r)(R-r)=0\). Nondegeneracy and the inverse
function theorem justify the local variable \(X\), and substitution proves
the displayed potential. Equivalently \(V_X=r\), so the auxiliary equation
in the \(X\) presentation is again \(R=r\).
\end{proof}

The local evaluation lift of Paper~I-C,
Proposition~\ref{p1b:prop:evaluation-legendre-contact-lift}, applies here
to the actual pair \(r,X=\overline{\mathcal F}'(r)\), with
\(z=\overline{\mathcal F}(r)\) and \(c=0\). It generates \(V\) from
that supplied graph. Neither the reciprocal response coordinate \(B/R\)
nor a time-evolution momentum is this conjugate coefficient.
Proposition~\ref{p3f:prop:vacuum-shift-symplectic-scope} shows why
equality of off-shell symplectic currents alone cannot fix the complete
zero-order coefficient.


\begin{corollary}[Generated bare normalized potential and the dressed target]
\label{p3f:cor:v2187-normalized-born-potential-scope}
On the generated bare two-root branch of
Theorems~\ref{p3f:thm:born-coefficient-family-first-response} and
\ref{p3f:thm:bare-born-potential-criterion}, put
\(\overline{\mathcal F}_{\rm reg}=U\overline f_{\rm reg}\).  Its normalized
same-action potential is already
\[
 \boxed{\mathcal N_{V,\rm bare}(R)
 :=\frac{R\overline{\mathcal F}_{\rm reg}'-
                \overline{\mathcal F}_{\rm reg}}{UR}
 =\frac{R+CU/R}{U+C}=g_{\rm eff}(R).}
\]
Consequently it is Born-even under \(R\mapsto CU/R\), with
\(\mathcal N_{V,\rm bare}(C)=\mathcal N_{V,\rm bare}(U)=1\).
The terminal slope and bare affine selector are derived by that upstream
coefficient-family constructor, not added potential or junction premises.

The dressed action and Euler response \(P_{2\mid7}=Rf_R-f\) are constructed
by Theorem~\ref{p3f:thm:v263-dressed-born-legendre-descent}.  For its fixed
normalization \(\overline{\mathcal F}=Uf\), the corresponding coefficient
target is \(\mathcal N_{V,2\mid7}=P_{2\mid7}/R\).  Construction of this
target does not itself assert its pointwise Born invariance.  A claim of
normalized potential equivariance or reciprocal endpoint equality must name
the actual complete coefficient and target action, retaining its density,
TERM--REN allocation, depth/level, sources, boundaries and calibration marks.
An internal pointwise-equality conclusion from Born inheritance uses a
trivial target action and a Born-preserved complete input word.  If a retained
mark changes, the target
is instead the independently licensed reciprocal-branch transport of
Paper~II-A, Proposition~\ref{p2a:prop:typed-born-branch-transport}; unequal
marked branches are not required to have identical scalar representatives.
\end{corollary}
\begin{proof}
The bare constructor theorem gives
\(R\overline f_{{\rm reg},R}-\overline f_{\rm reg}
=(R^2+CU)/(U+C)\).  Multiplying by \(U\), dividing by \(UR\), and exchanging
\(R\) with \(CU/R\) proves the first statements, including the endpoints.
The dressed Euler identity gives the displayed normalized coefficient target.
The remaining distinctions are precisely the unmarked-input and target-sort
hypotheses of Paper~II-A, Theorem~\ref{p2a:thm:generated-born-inheritance},
and its typed reciprocal-branch transport.  They neither reopen the bare
constructor nor certify an unevaluated dressed/marked potential comparison.
\end{proof}


\section{Scalaron, finite curvature response, and cap conditions}
\label{paper3f:sec:organized-scalar}
\subsection{Fixed-momentum candidate vacua and canonical response readouts}
\label{p3f:sec:v2218-fixed-momentum-candidate-readouts}
The physical momentum seed remains $m_p=1$ with its constituted
$\Rmom$ fixed.  Paper~III-D,
Definition~\ref{p3d:def:v2218-fixed-momentum-candidate}, constructs one
finite native-data composition
\[
 f_Q(R)=R-C+\frac{4QR^2}{D}\Phi(s,y),\qquad
 D=U+C,\quad Q=8\sqrt\alpha,\quad H=\log(U/C),
\]
where $a=H/2$, $s=\log(R/\sqrt{UC})$ and
$y=(1-C^2/R^2)[\log(R/C)/H]^2$.
The fixed extension of $Q$ and this source incidence are declared
composition choices.  No theorem identifies them with the actual physical
metric/heat action merely because their ingredients are generated.  The
following results test this candidate as a metric coefficient; they do not
replace the certified active Laurent--Legendre branch of
Theorem~\ref{p3f:thm:v263-dressed-born-legendre-descent}, its action/contact
and local-excess certificates, or the existing vacuum and source obligations.
All derivatives below retain the complete finite-value contacts of
Paper~III-D, Proposition~\ref{p3d:prop:v2218-finite-orbit-contacts}.

\begin{proposition}[Regular late-vacuum branch of the fixed-$Q$ candidate]
\label{p3f:prop:v2218-candidate-ir-vacuum}
For the candidate above, put $R=Cz$ and define the regular function
\[
 w_0(z):=4z^2\Phi\left(-a+\log z,
       (1-z^{-2})\frac{\log^2z}{H^2}\right),\qquad
 \mu:=\frac{QC}{D}.
\]
Then $f_Q=C(z-1)+(QC^2/D)w_0(z)$ and its source-free constant-curvature
vacuum equation is exactly
\begin{equation}
 \boxed{\frac{2f_Q-Rf_{Q,R}}C
 =z-2+\mu\{2w_0(z)-zw_0'(z)\}=0.}
 \label{p3f:eq:v2218-candidate-ir-vacuum}
\end{equation}
With $H$ fixed and $\mu$ regarded as a bookkeeping parameter, there is a
unique local vacuum branch through $(z,\mu)=(2,0)$.  Its expansion is
\begin{equation}
 \boxed{R_{\rm dS}=2C+\frac{QC^2}{D}
       \{2w_0'(2)-2w_0(2)\}
       +O(Q^2C^3/D^2).}
 \label{p3f:eq:v2218-candidate-vacuum-shift}
\end{equation}
Along this branch,
\[
 F_Q=1+\mu w_0',\qquad K_Q:=f_{Q,RR}=\frac QD w_0'',\qquad
 F_Q-RK_Q=1+\mu(w_0'-zw_0'').
\]
If $w_0''(2)>0$, all three quantities are positive for sufficiently small
positive $\mu$.  On the unextended metric-$f_Q$ constant-vacuum domain the
trace-linearized scalaron mass is consequently positive and satisfies
\[
 \frac{m_{\rm sc}^2}{U}
 =\frac{D}{3QUw_0''(2)}+O(C/U).
\]
The error statements are at fixed finite $H,Q$ on this local parameter
branch; they are not uniform claims over changing hierarchy data.
\end{proposition}
\begin{proof}
The exact finite orbit word is real analytic in a neighbourhood of $z=2$.
Differentiating its displayed rescaling gives the vacuum equation and the
three coefficient formulas.  The derivative of the vacuum equation with
respect to $z$ is one at $(2,0)$.  The implicit-function theorem therefore
gives the local branch, and its first $\mu$ derivative is
$2w_0'(2)-2w_0(2)$.  Continuity proves the stated positivity under
$w_0''(2)>0$.  Inserting these coefficients into the trace-linearized
metric-$f$ mass $(F_Q-RK_Q)/(3K_Q)$ gives the expansion.  A full coupled
physical pole would additionally consume its actual constraints, mixed
blocks, sources and nonlocal response.
\end{proof}

\begin{remark}[Executed native infrared evaluation; numerical scope]
\label{p3f:rem:v2218-candidate-ir-numerics}
On the scheme-zero member $\log\alpha=4\pi^2+1/(6\sqrt2)$, direct
high-precision differentiation of the finite $\Psi$ word gives
\[
 \begin{aligned}
 w_0(2)&\simeq9.40426816\times10^{-6},&
 w_0'(2)&\simeq2.61235223\times10^{-5},\\
 w_0''(2)&\simeq4.55386010\times10^{-5}>0.&
 \end{aligned}
\]
A direct 190-digit working-precision solution of the complete vacuum
equation verifies the local positive branch at the native parameter value:
\[
 \begin{aligned}
 \frac{R_{\rm dS}-2C}{2C}
   &\simeq2.23583777\times10^{-116},\\
 F_Q(R_{\rm dS})-1&\simeq3.49345474\times10^{-116},\\
 (F_Q-RK_Q)(R_{\rm dS})-1
   &\simeq-8.68614643\times10^{-116},\\
 UK_Q(R_{\rm dS})&\simeq1.44440898\times10^5,\\
 m_{\rm sc}^2/U&\simeq2.30774897\times10^{-6}>0.
 \end{aligned}
\]
Thus this candidate has a numerically verified stable late metric-$f_Q$
vacuum extremely close to the Einstein value $\Lambda=C/2$, where
$\Lambda_{\rm eff}=R_{\rm dS}/4$ on that vacuum.  These computations
evaluate the candidate; they do not turn numerical signs into a global
positivity theorem or prove its physical action incidence.  They require
neither an imposed scalaron freeze at $R=C$ nor a variation of $\Rmom$.
\end{remark}

\begin{remark}[Finite-band inflation test of the same candidate]
\label{p3f:rem:v2218-candidate-inflation-numerics}
Set $F= f_{Q,R}$, $K=f_{Q,RR}$, $P=RF-f_Q$, $B=2f_Q-RF$.
On its regular Einstein scalar chart the potential slow-roll diagnostics
and potential e-fold integral are
\[
 \epsilon_V=\frac{B^2}{3P^2},\qquad
 \eta_V=\frac23\left\{\frac{F(F-RK)}{KP}-\frac{2B}{P}\right\},
 \qquad
 N_V=\frac32\int_{R_{\rm end}}^R\frac{PK}{FB}\,dR.
\]
Paper~III-D, Theorem~\ref{p3d:thm:v2218-fixed-candidate-endpoints}, bounds
the upper-band finite word by its almost quadratic expression
\[
 f_Q(R)\simeq R-C+
       \frac{Q(R^2-C^2)}D\left[\frac{\log(R/C)}H\right]^2.
\]
This slowly varying coefficient produces a hilltop, rather than a globally
constant plateau.  Full finite-word evaluation on the native scheme-zero
member gives a potential maximum at
$R/U\simeq4.65270534\times10^{-8}$ with
$\eta_V\simeq-0.00508349$, and a left-branch end
$\epsilon_V=1$ at $R_{\rm end}/U\simeq2.12590731\times10^{-10}$.
The potential integral on the lower-$R$ side gives the following values.
\begin{table}[htbp]
\centering\small
\begin{tabular}{crrr}
\toprule
$N_V$&$R/U$&$\epsilon_V$&$\eta_V$\\
\midrule
50&$1.12091611\times10^{-8}$&$1.97338272\times10^{-4}$&$-0.02125991$\\
55&$1.21132352\times10^{-8}$&$1.60275850\times10^{-4}$&$-0.01968061$\\
60&$1.29927995\times10^{-8}$&$1.32158344\times10^{-4}$&$-0.01835286$\\
\bottomrule
\end{tabular}
\caption{Numerical potential slow-roll test of the fixed-$Q$ candidate,
not a selected physical cosmological history.}
\label{p3f:tab:v2218-candidate-slow-roll}
\end{table}
The full finite value and its first two logarithmic curvature derivatives
were checked at these points against the upper-band expression.  The moving
lower eigenvalue exceeds 120 and its intrinsic $\Psi$ value-error bound is
below $4\times10^{-105}$.  Independent 48- and 64-node high-precision Gauss
quadratures of the full finite-word integrand reproduce the three $N_V$
values, with quadrature differences below $3\times10^{-42}$.
A 205-point scan, refined across the $\Psi$ crossover, finds $F\ge1$ and
$K>0$ at every sampled point of the open interval.  It supplies sampled
numerical evidence only; the terminal value is $K(C)=0$ and no global
analytic sign theorem follows from this scan.

The same candidate thus supplies an inflationary coefficient example and a
late vacuum without fitting an intermediate switch.  $N_V$ is a potential
slow-roll integral, not the e-fold count of an exact coupled Lorentzian
solution.  No primordial amplitude, perturbation state, reheating, complete
UV-to-IR history, observational prediction, or physical owner selection is
established by these computations.  The candidate's source incidence and
constant extension of $Q$ remain unselected physical-action maps.
\end{remark}

\begin{corollary}[Canonical response rate on a fixed visibility parent]
\label{p3f:cor:v2218-fixed-visibility-canonical-rate}
Independently of the candidate test, retain the existing certified active
Laurent--Legendre action of
Theorem~\ref{p3f:thm:v263-dressed-born-legendre-descent}.  On an admitted
branch identifying its active curvature with the hierarchy readout, put
\[
 R_\tau=C\exp(He^{-\tau}),\qquad
 A=\sqrt{3/2}M,\quad M^2=U/4,\qquad \chi=f_{RR}.
\]
The exact canonical scalar and its response-parameter derivative are
\begin{equation}
 \boxed{\phi(\tau)=A\log F(R_\tau),\qquad
 \frac{d\phi}{d\tau}
 =-AH e^{-\tau}\frac{R_\tau\chi(R_\tau)}{F(R_\tau)}.}
 \label{p3f:eq:v2218-fixed-visibility-canonical-rate}
\end{equation}
For this existing action the terminal jet gives
\[
 F(R_\tau)-1=\frac{CH^2}{U+C}e^{-2\tau}+O(e^{-3\tau}),\qquad
 \phi(\tau)=\frac{ACH^2}{U+C}e^{-2\tau}+O(e^{-3\tau}).
\]
For the distinct composed candidate $f_Q$, the same canonical construction
instead has the endpoint coefficient
$\phi_Q(\tau)=3AQC\,e^{-2\tau}/(U+C)+O(e^{-3\tau})$.
Both statements hold with the compact momentum carrier and radius fixed.
The exponent is a rate with respect to the visibility/depth parameter;
neither formula identifies $\tau$ with physical proper time or asserts
that an actual history reaches the endpoint.
\end{corollary}
\begin{proof}
The canonical metric-$f$ scalar is $A\log F$ and
$dR_\tau/d\tau=-He^{-\tau}R_\tau$, proving the exact derivative.
For the certified Laurent action,
$\chi(C+x)=2x/[C(U+C)]+O(x^2)$ and hence
$F=1+x^2/[C(U+C)]+O(x^3)$.
Substituting $x=CH e^{-\tau}+O(e^{-2\tau})$ proves its asymptotics.
For $f_Q$, Paper~III-D's exact candidate terminal jet gives instead
$F_Q=1+3Qx^2/[C(U+C)H^2]+O(x^3)$, proving the second coefficient.
These are different coefficient graphs with a common fixed-parent
parameterization; their canonical normalizations have not been equated.
\end{proof}


\subsection{Carrier-four Legendre response and distinct UV/IR dynamical reductions}
\label{p3f:sec:v2219-carrier-four-uv-ir}
The previous Theorem~\ref{p3f:thm:v263-dressed-born-legendre-descent}
remains an algebraically valid depth-two \emph{comparison}.  The active
carrier-source selection is Theorem~\ref{p3d:thm:v2219-carrier-four-dressing};
it replaces $D_2$ by $D_4$, without inserting another microscopic carrier.
\begin{proposition}[Rank-four generated positive-response representative]
\label{p3f:prop:v2219-carrier4-legendre}
On the certified fixed-depth excess/action port and the positive response
chart $R>C_7$, the selected susceptibility and its integral are
\begin{align}
\chi_{4\mid7}(R)&=Q^{n_4^+(R)}
\frac{(1+n_7^+(R))(R^2-C_7^2)}{(U+C_7)R^2},\\
 f_{4\mid7}^{+}(R)&=R-C_7+\int_{C_7}^{R}(R-s)\chi_{4\mid7}(s)\,ds.\label{p3f:eq:v2219-carrier4-legendre}
\end{align}
It is the unique $C^2$ representative with this Hessian and the
terminal data $(f,f_R)(C_7)=(0,1)$. The coefficient is smooth at the
Born midpoint $C_2$; the positive-part third-jet jump first appears
at $C_4$, with value Eq.~\eqref{p3e:eq:v2219-c4-third-jet}.
\end{proposition}
\begin{proof}
The coefficient character $\rho_Q(s)=Q^s$ of
Theorem~\ref{p3f:thm:v263-laurent-coefficient-character} is independent
of which generated scalar exponent is evaluated. Apply it to $n_4^+$,
then integrate the positive continuous rank-seven local second jet twice.
The change of one-sided derivative at $C_4$ follows from the Laurent
positive-part zero and the product rule.
\end{proof}
\begin{theorem}[Coherent UV scalaron on a closed universe: conditional two-sided branch]
\label{p3f:thm:v2219-coherent-uv-scalaron}
Assume the \emph{same} parent action supplies an Einstein-normalized,
regular two-sided UV quadratic continuation through signed Lorentzian
$R_J=0$, with local coefficient $Q/U$.  This is an additional
parent/source certificate, not a continuation of the one-sided
$C_7$-normalized IR Legendre chart.  Then
\[
 f_{\rm UV}(R_J)=R_J+\frac{Q}{U}R_J^2
 =R_J+\frac{R_J^2}{6m_{\rm UV}^2},\qquad
 \boxed{m_{\rm UV}^2=\frac{U}{6Q}}.
\]
In the Einstein frame, $\phi=\sqrt{3/2}\,\bar M_{\rm Pl}\ln f_{\rm UV,R}$
has the exact potential
\[
 V_{\rm UV}(\phi)=\frac34\bar M_{\rm Pl}^2m_{\rm UV}^2
 \left(1-e^{-\sqrt{2/3}\phi/\bar M_{\rm Pl}}\right)^2,
\]
with expansion
\[
 V_{\rm UV}=\frac12m_{\rm UV}^2\phi^2
 -\frac{m_{\rm UV}^2}{\sqrt6\bar M_{\rm Pl}}\phi^3
 +\frac{7m_{\rm UV}^2}{36\bar M_{\rm Pl}^2}\phi^4+O(\phi^5).
\]
The positive mass pole about $\phi=0$ supports coherent oscillations
when expansion and backreaction are sufficiently slow.  The terminal
IR scale $C_7$ is \emph{not} a UV turning-point or stopping wall.
\end{theorem}
\begin{proof}
The stated UV two-sided action has $F=1+R_J/(3m_{\rm UV}^2)$.
Legendre--Weyl conversion yields the displayed potential, whose second
derivative at zero is positive $m_{\rm UV}^2$. Its local harmonic
solutions are oscillatory if $H_E/m_{\rm UV}$ and the nonlinear sources
are appropriately small. No condition at the distinct terminal
IR response chart enters this proof.
\end{proof}
\begin{proposition}[Closed-$S^3$ trace and coefficient-work terms]
\label{p3f:prop:v2219-closed-uv-trace}
On $ds_J^2=-dt_J^2+a_J^2d\Omega_3^2$ with $k=+1$,
\[
 R_J=6(\dot H_J+2H_J^2+a_J^{-2}),\quad
 3F(H_J^2+a_J^{-2})=\frac{\rho_J}{\bar M_{\rm Pl}^2}
 +\frac{FR_J-f}{2}-3H_J\dot F.
\]
For a fixed UV coefficient the homogeneous trace is exactly
\begin{equation}\label{p3f:eq:v2219-closed-uv-trace}
\boxed{\ddot R_J+3H_J\dot R_J+m_{\rm UV}^2R_J
 =\frac{m_{\rm UV}^2}{\bar M_{\rm Pl}^2}(\rho_J-3p_J).}
\end{equation}
For a \emph{separately owned} homogeneous coefficient coordinate $s(t)$,
if the parent reduction admits $f_{\rm UV}(R_J,s)=R_J+A(s)R_J^2$ with
$A=D_4(s)/U$, its metric trace includes
\begin{equation}\label{p3f:eq:v2219-variable-A-trace}
\boxed{\ddot R_J+\left(3H_J+2\frac{\dot A}{A}\right)\dot R_J
 +\left(\frac1{6A}+\frac{\ddot A+3H_J\dot A}{A}\right)R_J
 =\frac{\rho_J-3p_J}{6A\bar M_{\rm Pl}^2}.}
\end{equation}
Here the stress tensor includes the coefficient owner's metric variation;
the same coefficient work may not be counted again as a new source.
Neither $s=q_{\rm sp}(t)$ nor $s=q_R(R_J(t))$ has been established by this
equation. Oscillating physical Ricci curvature may be negative while
the spatial sections remain the positive-curvature $S^3$.
\end{proposition}
\begin{proof}
Use $FR-2f+3\Box F=T/\bar M_{\rm Pl}^2$ and
$\Box F=-\ddot F-3H_J\dot F$. For constant $A$,
$F=1+2AR_J$ yields the harmonic equation. For variable $A$,
$\ddot F+3H\dot F=2A(\ddot R_J+3H\dot R_J)
 +4\dot A\dot R_J+2(\ddot A+3H\dot A)R_J$, proving
Eq.~\eqref{p3f:eq:v2219-variable-A-trace}.
\end{proof}
\begin{proposition}[Complete particle Schur vertex and exact energy-bookkeeping port]
\label{p3f:prop:v2219-scalar-particle-vertex}
For $K_P=A_P-B_PD_P^{-1}B_P^\dagger$ (Hessian blocks of the
common action), the UV scalaron/particle vertex is
\[
\boxed{\begin{aligned}
\partial_\phi K_P={}&(A_P)_\phi-(B_P)_\phi D_P^{-1}B_P^\dagger
-B_PD_P^{-1}(B_P^\dagger)_\phi\\
&+B_PD_P^{-1}(D_P)_\phi D_P^{-1}B_P^\dagger.
\end{aligned}}
\]
It is distinct from $\partial_\eta K_P$.  For local real modes with
$L_P=(\dot q^TG(\phi)\dot q-q^TK(\phi)q)/2$ the explicit power
per physical volume deposited into the modes is
\[
\boxed{Q_{\phi\to P}=\frac{\dot\phi}{2}
 \left\langle q^T K_\phi q-\dot q^T G_\phi\dot q\right\rangle,}
\]
up to explicit moving-basis, derivative, metric and retained-boundary
contacts in the same action. Thus a two-sector coarse-grained ledger is
$\dot\rho_\phi+3H_E(\rho_\phi+p_\phi)=-Q_{\phi\to P}$ and
$\dot\rho_P+3H_E(\rho_P+p_P)=Q_{\phi\to P}$ only when no separately
retained boundary-energy reservoir remains. The $K16$ packet must not
be counted a second time.  The actual finite vertex, occupations, and
backreaction remain unexecuted.
\end{proposition}
\begin{remark}[UV/IR action typing]
The parent has distinct reductions
$S_\star\to S_{\rm UV}^{\rm coh}$ and
$S_\star\to S_{\rm IR}^{\rm eff}$, not two additive gravitational actions.
The latter may use retarded low-frequency Schur elimination; the former
must retain the occupied coherent scalaron until it transfers energy.
The original $n_2\mid n_7$ first-passage trajectory proves no claim
about coherent UV oscillation count.
\end{remark}


\subsection{Variational completion and controlled cap comparisons}
\label{p3f:sec:v223-cap-variation}
This section concerns the conditional capped constructor of
Theorem~\ref{p3e:thm:v223-exact-cap-reduction}.  All displayed spacetime
actions are dimensionless $\mathcal I=S/\hbar$.  Native actions with
other tensor, derivative or boundary contributions must retain them.

\begin{theorem}[First-order cap parent with canonical boundary]
\label{p3f:thm:v223-cap-parent}
Let $\zeta$ be an auxiliary coordinate, $r=R_s\sinh\zeta$,
$t=\tanh(\zeta/2)$, $u=ct$, and $T=T_s(t)$ as in
Definition~\ref{p3e:def:v223-cap-data}.  Put $\rho=R_s/2$ and
\[
 W=Pr+B_I(r^2+2urz+z^2)
       +\operatorname{Tr}(Y^{\mathsf T}TY)+2\rho\operatorname{Tr}Y.
\]
On a smooth region with non-null canonical boundary, consider
\[
 \mathcal I_{\rm par}=\frac18\int_M\sqrt{-g}
       [\Phi\{R[g]-r\}+W]\,d^4x
 +\frac14\sum_\partial\epsilon\int_\partial\sqrt{|h|}\,\Phi K\,d^3x
 +\mathcal I_{\rm rest}.
\]
Assume $\mathcal I_{\rm rest}$ has no $\zeta,z,Y,\Phi$ dependence
on this branch.  Fix $h_{ij}$ and $\Phi$ on the boundary.  Classical
elimination gives the complete capped action of
Theorem~\ref{p3e:thm:v223-exact-cap-reduction}, with
\[
 z_*=-ur,\quad Y_*=-\rho T^{-1},\quad R[g]=r,\quad
 \Phi=f_c'(R[g]).
\]
The parent and reduced interior equations correspond in both directions,
with the induced boundary scalar data.
\end{theorem}
\begin{proof}
The $z,Y$ equations are algebraic and their eliminated value is
$W_*=f_c(r)$.  The multiplier equation gives $R=r$.
The $\zeta$ equation and the envelope theorem give
$0=(f_c'(r)-\Phi)r_\zeta$, where $r_\zeta=R_s\cosh\zeta>0$.
Thus $\zeta=\operatorname{arsinh}(R/R_s)$ uniquely and
$\Phi=f_c'(R)$.  Substitution proves equality of bulk and boundary
actions.  The boundary variation of $\Phi R/8$ is cancelled by
$\epsilon\Phi K/4$ with fixed induced metric and fixed $\Phi$,
by the same scalar-tensor calculation as the canonical endpoint theorem
\ref{p3f:thm:v169-complete-gravity-action}.  The remaining variables
are algebraic and produce no integration-by-parts boundary terms.
Every reduced solution therefore lifts by the displayed algebraic maps,
and the envelope identity proves the converse interior equation.
\end{proof}
This is a variational presentation of a declared curvature map.  Its
multiplier does not prove the upstream origin of that map.  In
particular $\zeta$ is not the centered logarithmic phase coordinate
of Remark~\ref{p3d:rem:v223-log-asinh}.  A derived $q(R[g])$ has no
additional hierarchy equation.  If instead a hierarchy is varied
independently or prescribed externally, its Euler or exchange terms
belong to that different action, as required by
Theorem~\ref{p3a:thm:variational-sort-firewall}.

\begin{corollary}[Complete curvature momentum and metric equation]
\label{p3f:cor:v223-cap-metric}
With $H=R^2+2\sigma^2/B^3$, total curvature derivatives are
\[
 \Phi=f_c'=P+2RB+B'H,\qquad
 f_c''=2B+4RB'+B''H-6\sigma^2(B')^2/B^4.
\]
For $T_{\mu\nu}=-(2/\sqrt{-g})\delta\mathcal I_{\rm rest}/\delta g^{\mu\nu}$,
the metric and trace equations are
\[
 f_c'R_{\mu\nu}-\tfrac12 f_cg_{\mu\nu}
 +(g_{\mu\nu}\Box-\nabla_\mu\nabla_\nu)f_c'=4T_{\mu\nu},
 \qquad 3\Box f_c'+f_c'R-2f_c=4T.
\]
The left metric tensor is identically divergence free.  Its boundary
momentum includes the source derivative $2\sigma^2B'/B^3$.
\end{corollary}
\begin{proof}
Differentiate the complete eliminated density twice, including its
source term, and apply the local metric-$f(R)$ variation theorem
\ref{p3f:thm:general-local-fR-variation}.  The coefficient $1/8$ with
the stated stress convention yields the factor four.  Contracted
Bianchi and the scalar commutator identity prove the divergence claim.
\end{proof}
Keeping $B(R)$ in the density while omitting these derivative contacts
would change both the field equation and the canonical endpoint term.
No null, corner, asymptotic or additional spectral-boundary completion
is asserted by this smooth non-null statement.

\begin{proposition}[Regular Legendre chart of the retained cap]
\label{p3f:prop:v223-cap-legendre}
Write $d=1-B_U/B_I$.  If $2B_U>d$, then on the real curvature line
\[
 f_c''\ge 2B_U-d>0,\qquad f_c'(R)\ge P>0\quad(R\ge0).
\]
Let $r(\Phi)$ invert $\Phi=f_c'(r)$ on an interval and
$U_J(\Phi)=\Phi r-f_c(r)$.  Then $U_J'=r$, $U_J''=1/f_c''$.
On $\Phi>0$, the transformation
\[
 g_E=(\Phi/P)g,\quad
 \varphi=\sqrt{3P/8}\log(\Phi/P),\quad
 V_E=P^2U_J/(8\Phi^2)
\]
gives bulk density $PR_E/8-(\nabla\varphi)^2/2-V_E$ and canonical
boundary $\epsilon PK_E/4$, with the full pulled-back remaining action.
\end{proposition}
\begin{proof}
Set $e=B_U/B_I$, $x=R/R_s$, $G=\sqrt{1+x^2}$, $D_0=2+e(G-1)$,
and $H_0=(B_I/B)^2=(1+G)^2/D_0^2$.  Direct differentiation gives
\[
 (H_0)_{xx}=\frac{8d[d-e(1+G)]x^2}{D_0^4G^2}
       +\frac{4d(1+G)}{D_0^3G^3}
 \le\frac d2\left[d+\frac e{G^2}+\frac1{G^3}\right]\le d.
\]
For the middle inequality, use $D_0\ge2$, discard the negative part
of the first numerator, and use $x^2/G^2=1-G^{-2}$.
Also $(R^2B)_{RR}=2B_U+2\Delta/G^3$ and
$(\sigma^2/B^2)_{RR}=(H_0)_{xx}$.  This proves the first bound.
Since $f_c'(0)=P$, integration proves the second.  The inverse function
and conformal transformation formulae prove the remaining assertions;
the normal derivative from the bulk conformal divergence cancels that
from the transformed boundary term.
\end{proof}
The cap anchors of Proposition~\ref{p3d:prop:v223-curvature-clock}
satisfy $2B_U>d$.  This gives a regular scalaron chart, not a proof
of positive scalaron mass on every background, global evolution or
quantum completeness.  It introduces no second hierarchy particle.


\begin{theorem}[Unique positive cap vacuum and scalaron-sign certificate]
\label{p3f:thm:v288-cap-vacuum-root-stability}
On the exact \emph{conditional capped-action branch} of
Theorem~\ref{p3e:thm:v223-exact-cap-reduction}, write
\[
\begin{gathered}
 e=B_U/B_I\in(0,1),\quad d=1-e,\quad p=P/R_s,\quad x=R/R_s,\\
 G=\sqrt{1+x^2},\quad b(x)=B(R)/B_I=\frac{2+e(G-1)}{G+1}.
\end{gathered}
\]
Assume
\[
 \boxed{B_U>2,\qquad pe>8,\qquad p>1,\qquad B_I\le p.}
\]
The source-free metric-$f_c(R)$ vacuum trace equation
\(R f_c'(R)-2f_c(R)=0\) has \emph{exactly one} solution $R_*>0$.
With $C_{\rm cap}:=R_s^2/P$ it satisfies the strict bounds
\[
 \boxed{C_{\rm cap}<R_*<2C_{\rm cap},\qquad
 0<1-\frac{R_*}{2C_{\rm cap}}
 \le\frac{4dB_I}{p^4}.}
\]
The curvature momentum obeys $f_c'(R)>0$, the local susceptibility
obeys $f_c''(R)>0$, and
\[
 f_c'(R)-R f_c''(R)>0\qquad (R\ge0).
\]
Consequently, on the source-free constant-curvature vacuum, the
standard metric scalaron parameter
\(m_s^2=[f_c'(R_*)-R_* f_c''(R_*)]/[3f_c''(R_*)]\)
is strictly positive.
The hierarchy anchors of Proposition~\ref{p3d:prop:v223-curvature-clock},
with $P$ retained as their curvature owner and
$\alpha=\exp(4\pi^2+\delta)$, $\delta=1/(6\sqrt2)$, satisfy all
four inequalities. There
\(C_{\rm cap}=P\alpha^{-7}\); identifying it with the separate
$U$-owned $C_7=U\alpha^{-7}$ additionally requires $P=U$.
This is a branch-conditional action prediction, not a proof of native
common-action occurrence or a cosmological solution history.
\end{theorem}
\begin{proof}
On the cap domain of
Theorem~\ref{p3e:thm:v223-exact-cap-reduction},
$\sigma=B_I R_s$; write the dimensionless complete density as
\[
 F(x):=\frac{f_c(R_sx)}{R_s^2}
 =p x+B_I x^2b(x)-b(x)^{-2},\qquad
 \mathcal E(x):=xF_x-2F.
\]
Put $D=2+e(G-1)$ and $H=(b+xb_x)/b^3$. Direct differentiation and
rationalization give the exact identities
\[
 x^3 b_x=-\frac{2d(G-1)^2}{G},\qquad
 H-1=\frac{d(G-1)^2}{G}\,\varphi(G),
\]
\[
 \varphi(G)=\frac{e(1+e)G^2+e(5-e)G+2}{D^3},
 \qquad
 H_G=\frac{2d(G^2-1)[e(3G^2+4G-1)+2]}{G^2D^4}\ge0.
\]
Thus
\[
 \boxed{\mathcal E(x)=2-px+
 \frac{2d(G-1)^2}{G}[\varphi(G)-B_I].}
\]
For $z=eG\ge0$, the estimates $D\ge1+z$ and $0<e<1$ yield
\[
 0<\varphi(G)\le\frac1e\frac{2z^2+5z+2}{(1+z)^3}
 \le\frac2e,
\]
where the last inequality follows by expansion. Likewise
$3G^2+4G-1\le6G^2$, so
\[
 0\le H_G\le \frac2e\frac{6z^2+2}{(1+z)^4}
 \le\frac4e.
\]
As $0\le dG/dx=x/G\le1$ and
$(G-1)^2/G=G-2+G^{-1}$ is nondecreasing,
\[
 \mathcal E_x\le-p+2H_x\le-p+8/e<0.
\]
In particular the complete trace function is strictly decreasing.
Since $B_I e=B_U>2$, its correction to $2-px$ is negative for
$x>0$. Hence $\mathcal E(2/p)<0$. From $H\ge H(1)=1$ and
$G-1=x^2/(1+G)\le x^2/2$ we also have
\[
 \mathcal E(x)\ge2-px-\frac{dB_I}{2}x^4,
 \quad\text{so}\quad
 \mathcal E(1/p)\ge1-\frac{dB_I}{2p^4}>0.
\]
Strict monotonicity and continuity give exactly one root in
$(1/p,2/p)$. At the root the same lower estimate, combined with
$x_*<2/p$, gives
$px_*\ge2-8dB_I/p^4$, yielding the stated error bound.
The hypotheses imply $2B_U>d$; Proposition~\ref{p3f:prop:v223-cap-legendre}
therefore proves $f_c''\ge2B_U-d>0$ and $f_c'\ge P>0$ on $R\ge0$.
Because $d\mathcal E/dx=(R f_c''-f_c')/R_s$, strict decrease gives
$f_c'-Rf_c''>0$, which proves the scalaron sign.
Finally Proposition~\ref{p3d:prop:v223-curvature-clock} gives
the stated hierarchy values
$p=\alpha^{7/2}$, $B_I=\alpha^{7/2-\delta/2}<p$,
$B_U=\alpha^{1/2+\delta/2}>2$ and
$pe=\alpha^{1/2+\delta}>8$.
\end{proof}

\begin{theorem}[Fixed finite collar and complete scalar error]
\label{p3f:thm:v223-collar-error}
For fixed dimensionless width $\lambda>0$, with Dirichlet data at zero
and natural Neumann data at the far end, put
\[
 k_\lambda(G)=G\tanh(\lambda G),\quad a=\tanh\lambda,\quad
 h=\frac2{1+G},\quad h_\lambda=\frac{2a}{a+k_\lambda(G)},
\]
\[
 B_\lambda=B_U+\Delta h_\lambda,\quad
 D=\sigma^2/B^2,\quad D_\lambda=\sigma^2/B_\lambda^2,
 \quad f_\lambda=PR+R^2B_\lambda-D_\lambda.
\]
This is the matched finite model: both gravity and the normalized source
use $B_\lambda$.  Then
\[
 ah\le h_\lambda\le h,\quad aB\le B_\lambda\le B,\quad
 D\le D_\lambda\le a^{-2}D,
\]
\[
 0\le f_c-f_\lambda\le
 (1-a)R^2(B-B_U)+(a^{-2}-1)D=:\mathcal E_0(R).
\]
At zero curvature,
$f_\lambda''(0)=2B_I-d[1+2\lambda/\sinh(2\lambda)]$.
For $\lambda\ge1$ define
\[
 \kappa_1=(2+4\lambda)e^{-2\lambda},\quad
 \kappa_2=8\lambda(1+\lambda)e^{-2\lambda},\quad
 E_1=\frac{\kappa_1+1-a}{2a},
\]
\[
 E_2=\frac12\left[\frac{2\kappa_1+\kappa_1^2}{a^2}
                +(a^{-2}-1)+\frac{\kappa_2}{a}\right].
\]
\[
 \epsilon_0=\Delta(1-a),\quad
 \epsilon_1=\Delta E_1/R_s,\quad
 \epsilon_2=\Delta(E_1+E_2)/R_s^2.
\]
Globally $|B_\lambda-B|\le\epsilon_0$,
$|B_\lambda'-B'|\le\epsilon_1$,
$|B_\lambda''-B''|\le\epsilon_2$.
On $|R|\le R_*$ let $m=aB_U$, $M_1=\Delta/(2R_s)$,
$M_2=3\Delta/(4R_s^2)$.  Then
\[
 |f_\lambda'-f_c'|\le\mathcal E_1:=
 2R_*\epsilon_0+R_*^2\epsilon_1
 +2\sigma^2\left(\frac{\epsilon_1}{m^3}
                         +\frac{3M_1\epsilon_0}{m^4}\right),
\]
\[
 \begin{split}
 |f_\lambda''-f_c''|\le\mathcal E_2:={}&
 2\epsilon_0+4R_*\epsilon_1+R_*^2\epsilon_2\\
 &+2\sigma^2\left[\frac{\epsilon_2}{m^3}
 +\frac{3M_2\epsilon_0}{m^4}
 +\frac{3(2M_1+\epsilon_1)\epsilon_1}{m^4}
 +\frac{12M_1^2\epsilon_0}{m^5}\right].
 \end{split}
\]
On the same geometry, the bulk and smooth non-null canonical boundary
action values therefore differ by at most
\[
 \frac18\int_M\sqrt{-g}\,\mathcal E_0(R)\,d^4x
 +\frac{\mathcal E_1}{4}\sum_\partial
                   \int_\partial\sqrt{|h|}\,|K|\,d^3x.
\]
\end{theorem}
\begin{proof}
The finite solution is $v(y)=w\cosh(G(\lambda-y))/\cosh(\lambda G)$;
its endpoint response is $k_\lambda w$.  Normalize its conjugate-flux
joining by the reference $k_\lambda(1)=a$.  This gives $h_\lambda$.
Since $a\le\tanh(\lambda G)\le1$ for $G\ge1$, the value bounds
follow directly, and differentiating at $G(0)=1$ gives the stated
zero-curvature Hessian.
Differentiation and $\operatorname{sech}^2z\le4e^{-2z}$ give
$|k_\lambda'-1|\le\kappa_1$, $|k_\lambda''|\le\kappa_2$.
Put $D_1=a+k_\lambda$, $D_2=1+G$.  Since
$a\le D_1/D_2\le1$ and $D_2\ge2$, use
\[
 h_{\lambda,G}=-2ak_\lambda'/D_1^2,\qquad
 h_{\lambda,GG}=4a(k_\lambda')^2/D_1^3-2ak_\lambda''/D_1^2
\]
to obtain $|\partial_G(h_\lambda-h)|\le E_1$ and
$|\partial_G^2(h_\lambda-h)|\le E_2$.
The chain rule with $|G_x|,|G_{xx}|\le1$ gives the $B$ bounds.
Directly $|B'|\le M_1$, $|B''|\le M_2$.
Apply the mean value theorem to $B^{-3},B^{-4}$ above their common
lower bound $m$ and use
Corollary~\ref{p3f:cor:v223-cap-metric}.  For the squared derivative
use $|(B_\lambda')^2-(B')^2|\le(2M_1+\epsilon_1)\epsilon_1$.
This proves $\mathcal E_1,\mathcal E_2$; integrating the density and
canonical boundary coefficient errors gives the last bound.
\end{proof}
If $\mathcal E_2<2B_U-d$, positivity of the cap Hessian persists on
that compact range.  No all-curvature stability claim for every width
follows from the zero-curvature calculation alone.  Width and reference
frequency are held fixed under the derivatives above; a metric-dependent
proper width has additional contacts.  These are action-value bounds
on the same geometry, not a proof that the two boundary-value problems
have identical fixed $\Phi$ data or nearby solutions.  Identifying the
actual native width and finite source ports is a separate physical
comparison.  The width $\lambda$ is not $L=\log\alpha$ or the Hopf
residue.  A tangential operator requires the independent bound of
Theorem~\ref{p1b:thm:v223-root-perturbation} and its actual jets.

\begin{proposition}[An integrated collar has a different determinant]
\label{p3f:prop:v223-collar-determinant}
If the extension is itself integrated as a Gaussian field with fixed
$L^2$ measure, its Dirichlet--Neumann fluctuation eigenvalues are
$((n+1/2)\pi/\lambda)^2+G^2$.  Its convergent relative determinant is
\[
 \frac{\det(-\partial_y^2+G^2)_{DN}}
      {\det(-\partial_y^2+1)_{DN}}
 =\frac{\cosh(\lambda G)}{\cosh\lambda},\qquad
 \lambda\tanh\lambda(G-1)\le
 \log\frac{\cosh(\lambda G)}{\cosh\lambda}\le\lambda(G-1).
\]
An actual multiplicity $m_c$ therefore adds
$m_c\hbar_{\rm eff}\log[\cosh(\lambda G)/\cosh\lambda]/2$
to the dimensionless finite effective action.
\end{proposition}
\begin{proof}
Separate variables on the interval.  The normalized product for
$\cosh$ gives the determinant ratio; its factors differ from one by
$O(n^{-2})$, so it converges.  Integrate
$\partial_G\log\cosh(\lambda G)=\lambda\tanh(\lambda G)$ from $1$
to $G$ to obtain the bounds.
\end{proof}
The extension used as an operator representation does not automatically
introduce such a field integration.  If it is an actual field, classical
exponential convergence of $h_\lambda$ does not control this determinant
as $\lambda\to\infty$.  The finite source volume does not cancel a
different collar-field determinant.  Its actual multiplicity, measure,
other sector determinants and any absolute vacuum normalization must
be evaluated by the native common action; none is subtracted here.
Likewise a change from $r$ to $\zeta$ transports its measure:
$dr=R_s\cosh\zeta\,d\zeta$.  A flat $d\zeta$ constraint integral
would leave the corresponding nonconstant Jacobian.  Classical
auxiliary equivalence is not by itself a quantum measure theorem.


\subsection{Separated-scale Born--Legendre reconstruction and terminal scope}
\label{p3f:app:born-potential-criterion}
Fix \(0<C<U\) and put
\[
 B=CU,
 \qquad
 \mathscr H_{C,U}(R):=R g_{\rm eff}(R)
 =\frac{R^2+B}{U+C}.
\]
The unit primitive \(\mathscr Q_{C,U}=R^2+B\) and the finite interpolation
\(\mathscr H_{C,U}=\mathscr Q_{C,U}/(U+C)\) have different normalization
roles.  At the Paper~III-E Hessian stage,
\(\mathscr H_{C,U}''=2/(U+C)\) leaves the expected affine kernel.  On the
generated same-action branch, however,
Theorem~\ref{p3f:thm:born-coefficient-family-first-response} supplies the
upstream coefficient-family first response
\[
 \mathfrak f_{\rm nl}^{(1)}(R)
 =\frac{R^2-CU}{U+C}
\]
and reduces the compatible kernel to the parallel Einstein line
\(\operatorname{span}\{R\}\).  The bare reconstruction therefore retains the
already owned terminal value but no longer requires a separately imposed
terminal slope.

\begin{theorem}[Separated-scale bare same-action reconstruction]
\label{p3f:thm:bare-born-potential-criterion}
Assume the generated bare undressed two-root response lane of
Theorem~\ref{p3f:thm:born-coefficient-family-first-response} and retain the
complete terminal value
\[
 \overline f(C)=0.
\]
Then compatibility with the generated coefficient-family first response,
\begin{equation}
\label{p3f:eq:bare-first-response-compatibility}
 \nabla_U^{\rm sd}
 \left(\overline f-\mathfrak f_{\rm nl}^{(1)}\right)=0,
\end{equation}
fixes the complete coefficient uniquely as
\begin{equation}
\label{p3f:eq:regular-separated-action}
 \boxed{
 \overline f_{\rm reg}(R)
 =(R-C)+\frac{(R-C)^2}{U+C}.}
\end{equation}
Its terminal slope is a consequence,
\begin{equation}
\label{p3f:eq:bare-born-terminal-derived}
 \boxed{\overline f_{\rm reg,R}(C)=1,}
\end{equation}
and its same-action Legendre potential is
\begin{equation}
\label{p3f:eq:complete-born-legendre}
 \boxed{
 \mathcal V_{\rm reg}(R)
 :=R\overline f_{\rm reg,R}-\overline f_{\rm reg}
 =\frac{R^2+CU}{U+C}
 =R g_{\rm eff}(R).}
\end{equation}
Consequently
\[
 \frac{\mathcal V_{\rm reg}(R)}R
 =\frac1{U+C}\left(R+\frac{CU}{R}\right)
\]
is Born-even under \(R\mapsto CU/R\).  After restoring the independent
action normalization \(U\), the dimensionful coefficient
\(\overline{\mathcal F}_{\rm reg}=U\overline f_{\rm reg}\) obeys
\[
 R\overline{\mathcal F}_{\rm reg}'-
 \overline{\mathcal F}_{\rm reg}
 =U R g_{\rm eff}(R)
 =\frac{U}{U+C}(R^2+CU).
\]
No negative/anti-contact action transport and no independent terminal
\(C^1\) selector are required for this bare result.
\end{theorem}
\begin{proof}
By Theorem~\ref{p3f:thm:born-coefficient-family-first-response},
\(\ker\nabla_U^{\rm sd}=\operatorname{span}\{R\}\), so
\eqref{p3f:eq:bare-first-response-compatibility} is equivalent to
\[
 \overline f(R)
 =\frac{R^2-CU}{U+C}+\lambda R
\]
for one real \(\lambda\).  The terminal value gives
\[
 0=\frac{C^2-CU}{U+C}+\lambda C,
 \qquad
 \lambda=\frac{U-C}{U+C},
\]
which proves \eqref{p3f:eq:regular-separated-action}.  Differentiation gives
\eqref{p3f:eq:bare-born-terminal-derived}, and direct substitution in
\(R\overline f_R-\overline f\) gives
\eqref{p3f:eq:complete-born-legendre}.  The final Born-even identity is the
exact reciprocal Laurent identity of Paper~III-D.
\end{proof}

\begin{proposition}[Affine Hessian family and exact constructor defect]
\label{p3f:prop:complete-born-action-defect}
Every bare coefficient with the same Hessian and terminal value is
\[
 \overline f_\kappa(R)
 =\overline f_{\rm reg}(R)+\kappa(R-C),
\]
with
\[
 \overline f_\kappa'(C)=1+\kappa,
 \qquad
 \mathcal V_\kappa(R)
 =R\overline f_\kappa'-\overline f_\kappa
 =R g_{\rm eff}(R)+\kappa C.
\]
Relative to the generated nonlinear first-response section,
\begin{equation}
\label{p3f:eq:constructor-affine-defect}
 \boxed{
 \nabla_U^{\rm sd}
 \left(\overline f_\kappa-\mathfrak f_{\rm nl}^{(1)}\right)
 =\kappa C\,du_B.}
\end{equation}
After scalarization by the horizontal dual density the same defect is
\begin{equation}
\label{p3f:eq:constructor-affine-defect-scalarized}
 \boxed{\frac{\kappa C}{R}\,du_B.}
\end{equation}
Hence, since \(C>0\), compatibility with the actual coefficient-family
constructor is equivalent to
\begin{equation}
\label{p3f:eq:kappa-zero-constructor}
 \boxed{\kappa=0.}
\end{equation}
The sharp undressed TERM--REN representative
\[
 \overline f_{\rm sharp}(R)
 =R+\frac{R^2+CU}{U+C}-2C
\]
has
\[
 \boxed{
 \overline f_{\rm sharp}-\overline f_{\rm reg}
 =\frac{2C}{U+C}(R-C),
 \qquad
 \kappa_{\rm sharp}=\frac{2C}{U+C}>0.}
\]
It therefore does not belong to the generated bare same-action
coefficient-family member selected here.  This statement does not remove a
separately realized sharp dressed Laurent endpoint or determine a microscopic
junction profile, width, smoothing mechanism, or history.
\end{proposition}
\begin{proof}
The affine family is the kernel of the ordinary second derivative with the
terminal value fixed.  Since
\[
 \overline f_\kappa-\mathfrak f_{\rm nl}^{(1)}
 =\left(\frac{U-C}{U+C}+\kappa\right)R-\kappa C,
\]
the \(R\)-term is parallel while
\(\nabla_U^{\rm sd}(-\kappa C)=\kappa C\,du_B\), proving
\eqref{p3f:eq:constructor-affine-defect}.  Pairing with \(\mu_R\) gives
\eqref{p3f:eq:constructor-affine-defect-scalarized}.  Applying
\(R\partial_R-1\) to \(\kappa(R-C)\) gives the identical Legendre defect
\(\kappa C\).  Direct subtraction proves the sharp coefficient formula.
\end{proof}

\begin{proposition}[Vacuum shifts, Hessian blindness and coefficient-family detection]
\label{p3f:prop:vacuum-shift-symplectic-scope}
For a local density \(\sqrt{-g}\,\overline{\mathcal F}(R)\), an affine
change \(\Delta\overline{\mathcal F}=aR+b\) leaves the second curvature
derivative unchanged but changes the complete metric field equation through
its Einstein and vacuum pieces. Equality of \emph{curvature Hessians} alone
therefore cannot select \(\kappa\) in the preceding proposition. In contrast,
the full metric off-shell presymplectic current generally changes under the
\(aR\) Einstein--Hilbert term; only the constant \(b\) contributes no
metric-derivative presymplectic potential. The current is not being claimed
invariant under the whole affine family. On the generated determinant-dual coefficient line,
however, only the \(aR\) direction is parallel; a nonzero constant \(b\) is
detected by \(\nabla_U^{\rm sd}\).  The constructor closure above therefore
uses information strictly upstream of the Hessian quotient rather than an
extra boundary regularity law.
\end{proposition}
\begin{proof}
The curvature-Hessian statement is immediate.  The variation of
\(\sqrt{-g}(aR+b)\) contributes the ordinary Einstein tensor term
proportional to \(a\) and the metric term proportional to \(b\), while the
curvature-Hessian data are unchanged. The \(aR\) term also carries the usual
nonzero Einstein--Hilbert presymplectic potential and current for general
off-shell metric perturbations, unlike \(b\sqrt{-g}\), which has no derivatives
of the metric variations. By
\eqref{p3f:eq:born-parallel-kernel}, \(aR\) lies in the coefficient-connection
kernel, whereas \(\nabla_U^{\rm sd}b=-b\,du_B\).  This is exactly the
constant mode seen in
\eqref{p3f:eq:constructor-affine-defect}.
\end{proof}

\begin{corollary}[Constructor discharge of the bare affine selector]
\label{p3f:cor:constructor-born-discharge-bare}
On the generated bare same-action two-root branch, the already owned terminal
value \(\overline f(C)=0\) together with the frozen determinant-evaluation
constructor implies
\[
 \boxed{\kappa=0,\qquad
        \overline f_R(C)=1,\qquad
        R\overline f_R-\overline f=R g_{\rm eff}(R).}
\]
Thus the old terminal-slope atom and bare Born--Legendre mismatch are not
independent physical junction inputs on this branch: terminal \(C^1\) is a
derived consequence of the coefficient-family constructor.  No microscopic
transition profile or nonzero intrinsic width is generated here; the generic
GP exact-interface theorem instead retains the sharp interface exactly.  An
actual crossing history remains G0 occurrence data.  The distinct dressed
coefficient-character and action-placement atoms are discharged later by
Theorem~\ref{p3f:thm:v263-laurent-coefficient-character} and
Theorem~\ref{p3f:thm:v263-dressed-born-legendre-descent}.
\end{corollary}

\subparagraph{Fixed-linear-coefficient potential comparison.}
\label{p3f:app:v2187-fixed-linear-potential-comparison}
As a comparison family, \(\overline{\mathcal F}=UR+R^2+\beta\) has
\(\mathcal V=R^2-\beta\).  Equality of its normalized reciprocal endpoint
values at \(0<C<U\) gives
\[
 C-\frac\beta C=U-\frac\beta U
 \quad\Longleftrightarrow\quad \beta=-CU,
\]
and then \(\mathcal V/R=R+CU/R\) is Born-even.  This fixed-linear family is
not the generated bare coefficient-family member merely because it shares a
Hessian or reciprocal carrier.  The actual bare constructor and its
normalization are those of
Corollary~\ref{p3f:cor:v2187-normalized-born-potential-scope}; its affine
defect is tested upstream by
Proposition~\ref{p3f:prop:complete-born-action-defect}.
For a separately claimed dressed/marked potential comparison, evaluate its
actual normalized target with the complete retained word and declared Born
action.  Ordinary same-action variation ancestry alone proves neither
unmarkedness nor a trivial target action.

\subparagraph{Scope.}
This appendix section closes only the bare two-root affine selector.  The main-text
Theorems~\ref{p3f:thm:v263-laurent-coefficient-character} and
\ref{p3f:thm:v263-dressed-born-legendre-descent} now extend the same-action
construction through the depth-two dressing.  The structural interface is
closed by GP exactness and the same-action field equation; only occurrence of
a particular history remains separate under G0.  The earlier Hessian-only
classification remains correct as a statement about
\(\ker D^2\); the new result goes one constructor level upstream and uses the
generated coefficient-family connection, whose kernel is only
\(\operatorname{span}\{R\}\).  The negative anti-contact lift remains valid
Paper~I-C mathematics but is not required for the physical bare
Born--Legendre discharge.

\subsection{same-action complete-curvature susceptibility and the scalar-response firewall}
\label{p3f:sec:v1164-complete-curvature-response}

\begin{definition}[Generated curvature readout versus dynamical response]
\label{p3f:def:v1164-curvature-response-types}
On a licensed common-owner, local Lorentzian configuration domain, retain
exactly the effective position functional
\(S^{\rm common}_{x,U}\) of
Eq.~\eqref{p3:eq:common-position-spectral-action}, including the full
Born-diagonal copy of every pure-metric density and the inherited half
weight of each nonmetric position constituent.  Let \(S_\star\) denote the
claim-complete action actually licensed for the chosen physical response;
using this notation does not assert that the entire exact nonlocal trace
already possesses a certified Lorentzian action.

Let \(Y_{6,u}[g]\) be Paper~VI-A's rooted exceptional representation of the
Levi--Civita curvature operator.  It is a dependent curvature \emph{readout},
not a new field.  Its action-generated susceptibility is the map from
complete physical sources to \(\delta Y_{6,u}\), evaluated after the metric,
source, Ward and constraint equations are reduced on the occurring member.
No additional metric-curvature action or source normalization is included in
this definition.
\end{definition}

\begin{theorem}[Same-action retarded exceptional-curvature susceptibility]
\label{p3f:thm:v1164-exceptional-curvature-susceptibility}
Assume the causal dynamic Schur hypotheses of
Theorem~\ref{p3f:thm:v1144-causal-dynamic-schur} on the physical
source-conditioned quotient.  Denote the retained metric variable by \(h\),
the eliminated interior/source variable by \(z\), and the full linearized
operator by \(\left(\begin{smallmatrix}A&B\\D&C\end{smallmatrix}\right)\).
Suppose \(C\) and the Schur-reduced metric operator possess their physical
retarded inverses.  Then
\begin{equation}
\label{p3f:eq:v1164-full-curvature-susceptibility}
\boxed{
 \delta Y_{6,u}
 = D_gY_{6,u}\,(K_{\rm phys}^{R})^{-1}_{R}
   \left[j_h-B(C^{-1})_Rj_z\right],\quad
 K_{\rm phys}^{R}=A-B(C^{-1})_RD .}
\end{equation}
The differential \(D_gY_{6,u}\) includes the variation of the metric-induced
coframe used in the exceptional curvature representation.  Its scalar,
traceless-Ricci and Weyl projections are dependent projections of the same
response, not three independently normalized actions or three new compact
momentum carriers.  The complete response exists on precisely the causal,
constraint and source domains supplied in the hypotheses; an exceptional-root
frame identity alone does not construct those inverses.
\end{theorem}
\begin{proof}
The source-conditioned Schur theorem gives
\(K_{\rm phys}^{R}h=j_h-B(C^{-1})_Rj_z\).  Apply the physical retarded
inverse and differentiate the already constituted map \(g\mapsto Y_{6,u}[g]\).
Every factor is a derivative, inverse, or evaluation of the same action and
its inherited geometry.  If the coframe depends on \(g\), differentiating
\(Y_{6,u}=\sqrt{12}J_{6,u}[g]\mathcal K[g]\) gives both
\(\sqrt{12}(D_gJ_{6,u})\mathcal K\) and
\(\sqrt{12}J_{6,u}D_g\mathcal K\).  The displayed formula retains both.
\end{proof}

\begin{proposition}[The exact heat Hessian is fixed before local curvature expansion]
\label{p3f:prop:v1164-heat-hessian-no-insertion}
On the unlocalized compact Euclidean Dirac domain put
\(A(g,z)=D_{x,A}(g,z)^2/\Lambda_g^2\) and
\(S_{\rm heat}=\operatorname{Tr}e^{-A}\).  For two field variations
\(h_1,h_2\), where \(A_i=D A[h_i]\) and
\(A_{12}=D^2 A[h_1,h_2]\), one has
\begin{equation}
\label{p3f:eq:v1164-full-heat-hessian}
\boxed{
 D^2S_{\rm heat}[h_1,h_2]
 =\int_0^1\!ds\,\operatorname{Tr}
    (e^{-sA}A_1e^{-(1-s)A}A_2)
   -\operatorname{Tr}(e^{-A}A_{12}).}
\end{equation}
The same Duhamel differentiation with the analytic localizer, the
Born-diagonal metric pullback, and the nonmetric half weight defines the
actual localized/common-action blocks on their licensed domains.  Thus a
local \(a_4\) or \(a_6\) contribution is a controlled derivative-expansion
readout of the full operator, not a new independently weighted response.
\end{proposition}
\begin{proof}
Differentiate \(\operatorname{Tr}e^{-A}\) once using Duhamel's formula
and trace cyclicity, then differentiate again.  The first variation is
\(-\operatorname{Tr}(e^{-A}A_i)\); differentiating its exponential gives the
integral, and differentiating \(A_i\) gives the second term.  The statement
about role weights follows from
Lemma~\ref{p3f:lem:diagonal-eh-full-variation} and
Eq.~\eqref{p3:eq:common-position-spectral-action}.
\end{proof}

\begin{proposition}[Local Dirac metric coefficients and the non-scalar test]
\label{p3f:prop:v1164-dirac-a4-a6-firewall}
For the ordinary torsion-free pure Dirac specialization, the four-derivative
metric coefficient, modulo a divergence and an overall inherited coefficient,
is
\[
 5R^2-8R_{\mu\nu}R^{\mu\nu}
 -7R_{\mu\nu\rho\sigma}R^{\mu\nu\rho\sigma}
 =-18C_{\mu\nu\rho\sigma}C^{\mu\nu\rho\sigma}+11\mathcal G.
\]
Consequently its constant-coefficient contribution has no independent
\(R^2\) dynamical channel modulo the four-dimensional Gauss--Bonnet term;
its first metric variation vanishes on conformally flat FLRW, although the
Weyl Hessian can remain nonzero on perturbations.  The same local
four-derivative terms are background-inert on Einstein/Kottler metrics.

On the same ordinary Dirac specialization, the established conformally flat
six-derivative density, modulo the same boundary/continuation conditions,
can be written
\begin{equation}
\label{p3f:eq:v1164-a6-flrw-test}
 \mathcal L_{6}^{\rm FLRW}
 =\zeta_6\left[-\frac1{10}R\Box R
 +\frac{31}{2160}R^3-\frac{19}{720}Rq_E^2
 +\frac{13}{1080}q_E^3\right],
 \qquad q_E=6\left(\frac{k}{a^2}-\dot H\right).
\end{equation}
Here \(\zeta_6\) is a \emph{placeholder for the coefficient to be evaluated
from the fixed common heat profile, actual internal trace and licensed
Lorentzian sign}, not a new phenomenological parameter.  At constant
\(R=R_*\) with \(|q_E|\gg |R_*|\), the displayed algebraic density has a
nonzero \(13\zeta_6q_E^3/1080\) term, hence cannot be represented by a
one-variable \(f(R)\) density.  Equation~\eqref{p3f:eq:v1164-a6-flrw-test}
is a \emph{restricted heat-descendant diagnostic}, not an additive
same-action completion of the native nonlinear \(f_{2\mid7}\) branch.
To derive full covariant metric field equations one must vary the unprojected
heat coefficient \emph{before} imposing conformal flatness: terms linear in
Weyl curvature may have nonzero first metric variation when Weyl curvature
itself vanishes.  Thus a restricted FLRW lapse derivative is not, by itself,
a certificate of the complete covariant Born--Hopf source equation.
\end{proposition}
\begin{proof}
The first identity is the four-dimensional Weyl/Gauss--Bonnet algebra.
The \(a_6\) formula follows by specializing the ordinary Dirac
six-derivative coefficient to conformally flat curvature and substituting
\(R=6(A_H+B_H)\), \(q_E=6(B_H-A_H)\); equivalently its cubic part is
\(\tfrac25 A_H(-13A_H^2+57A_HB_H+18B_H^2)\).  These are the relative
coefficients of the ordinary Dirac heat kernel (see R.~Mistry,
A.~Pinzul and L.~Rachwa\l, \emph{Eur. Phys. J. C} \textbf{80}, 266 (2020),
DOI:10.1140/epjc/s10052-020-7805-1).  The scalar-only density cannot
produce an independent \(q_E^3\) variation on the full FLRW curvature
coordinate domain.  This says nothing by itself about physical mode
multiplicity or the complete joint action's mixed terms.
\end{proof}

\begin{proposition}[Terminal scalar jets obstruct a naive additive six-derivative completion]
\label{p3f:prop:v1164-terminal-jet-addition-nogo}
Suppose, only as a diagnostic, that the ordinary Dirac six-derivative term
of Eq.~\eqref{p3f:eq:v1164-a6-flrw-test} were appended as an independent
physical functional to
\(S_f=(M_*^2/2)\int\sqrt{-g}\,f_{2\mid7}(R)\) with
\(M_*^2>0\).  Restricting the candidate to the four-dimensional
space-form family \(q_E=0\), \(\nabla R=0\) gives
\[
 f_{\rm diagnostic}(R)
 =f_{2\mid7}(R)+\frac{31\zeta_6}{1080M_*^2}R^3.
\]
At \(R=C_7>0\), the resulting shifts of the three fixed terminal
coefficient jets are
\[
 \boxed{\Delta f(C_7)=\frac{31\zeta_6 C_7^3}{1080M_*^2},\qquad
 \Delta F(C_7)=\frac{31\zeta_6 C_7^2}{360M_*^2},\qquad
 \Delta f_{RR}(C_7)=\frac{31\zeta_6 C_7}{180M_*^2}.}
\]
For a nonzero inherited \(\zeta_6\), an \emph{uncompensated} additive
term therefore contradicts the exact terminal normalization
\(f(C_7)=0\), \(F(C_7)=1\), \(f_{RR}(C_7)=0\) when these are asserted
as the jets of the complete action.  The calculation does not contradict
the existing restricted \(f_{2\mid7}\) theorem: it proves that an action
quotient/subtraction, exact mixed-term cancellation, or a carefully stated
truncation is required before the two descriptions can be combined.
\end{proposition}
\begin{proof}
On the space-form family the derivative term vanishes, \(q_E=0\), and
the cubic density reduces to \(31\zeta_6R^3/2160\).  Equating the
resulting action to \((M_*^2/2)\int\sqrt{-g}\,f_{\rm diagnostic}\)
and taking zero, one and two derivatives proves the three displayed
expressions.  Since \(C_7>0\), none vanishes when \(\zeta_6\ne0\).
\end{proof}

\begin{theorem}[Scalaron port-matching criterion for a generalized curvature response]
\label{p3f:thm:v1164-scalaron-port-matching}
Consider one admissible background and source-conditioned, Ward-reduced
retarded physical chart in which the inherited scalaron perturbation
\(\phi:=\delta F=f_{RR}(R_0)\delta R\) is a valid coordinate
(\(f_{RR}(R_0)\ne0\)), and collect the remaining constrained physical
metric and source coordinates in \(w\).  Write its complete 
\emph{same-action} linearized equations in blocks,
\[
 \begin{pmatrix}H_{\phi\phi}&H_{\phi w}\\
 H_{w\phi}&H_{ww}\end{pmatrix}
 \begin{pmatrix}\phi\\w\end{pmatrix}
 =\begin{pmatrix}j_\phi\\j_w\end{pmatrix}.
\]
If the physical retarded inverse \((H_{ww}^{-1})_R\) exists, then the
unique eliminated scalar equation is
\begin{equation}\label{p3f:eq:v1164-scalaron-schur-test}
 \boxed{
 K_{\phi,\rm eff}^{R}\phi
 =j_{\phi,\rm eff},\quad
 K_{\phi,\rm eff}^{R}:=H_{\phi\phi}
       -H_{\phi w}(H_{ww}^{-1})_RH_{w\phi},\quad
 j_{\phi,\rm eff}:=j_\phi-H_{\phi w}(H_{ww}^{-1})_Rj_w.}
\end{equation}
To claim the exact \emph{same} scalaron propagation operator as the
terminal-normalized \(f_{2\mid7}\) sector on that common physical domain,
it is necessary and sufficient to establish
\(K_{\phi,\rm eff}^{R}=K_{\phi,f}^{R}\), with matching causal
boundary data, variables and source-port normalization.  This does
\emph{not} force \(H_{\phi w}=0\): inherited non-scalar curvature/source
ports can produce the additional effective forcing
\(-H_{\phi w}(H_{ww}^{-1})_Rj_w\) while leaving the scalaron kinetic
operator and mass law unchanged.  A nonzero forcing is a testable output,
not guaranteed by the block algebra alone.

At the terminal point \(f_{RR}(C_7)=0\), the \(\phi\)-coordinate becomes
degenerate and this criterion must be replaced by the original
scalar-curvature constraint and canonical collar/boundary matching, not
by division through \(f_{RR}(C_7)\).
\end{theorem}
\begin{proof}
Solve the second block equation retardedly for
\(w=(H_{ww}^{-1})_R(j_w-H_{w\phi}\phi)\), substitute it into the
first and collect terms.  Equality of the resulting scalar operator
with the inherited \(f_{2\mid7}\) operator gives the matching criterion
when both have a common invertible causal domain; equality of their
retarded inverses implies equality of the operators there.  The
terminal warning follows from the vanishing local Legendre Jacobian.
\end{proof}

\begin{proposition}[Exact scalar-action matching is required before combining sectors]
\label{p3f:prop:v1164-scalar-overlap-condition}
Let \(S_{2\mid7}\) be the existing terminal-normalized Born--Legendre
scalar action on its fixed-depth local-jet domain.  A proposed decomposition
of the \emph{same} completed action as
\(S_\star=S_{2\mid7}+S_{\rm residual}\) is admissible only if the equality
holds as an action-functional identity in common physical units, including
all pure-metric, nonmetric, mixed and boundary constituents and excluding any
already-consumed term from \(S_{\rm residual}\).  If the existing
\(f_{RR}\) and trace-mode mass are claimed to be the \emph{exact} physical
scalar response of the completed action, its constrained scalar Hessian must
agree with the one derived from \(S_{2\mid7}\); otherwise they describe
only a restricted/truncated response.  In particular a free-standing
addition of Eq.~\eqref{p3f:eq:v1164-a6-flrw-test} is not licensed by the
common-heat theorem alone.

The terminal scalaron-memory theorem is unchanged on its declared regular
Legendre collar, but no inference is made that every higher/tidal metric
response vanishes at \(R=C_7\).  For a specified regular finite-order
higher-jet action, source-free matching additionally requires continuity of
the complete canonical boundary momenta appropriate to that order.
\end{proposition}
\begin{proof}
The first claim is the defining condition for functional decomposition of
one action; counting a term twice changes \(D^2S_\star\), and therefore
changes the derived Schur response.  The scalar-mode equality follows by
restricting that same physical Hessian to the already fixed scalar response
port.  The final claims follow from the scoped terminal Schur theorem and
the general higher-jet variational matching rule, respectively.
\end{proof}

\begin{remark}[Result and exact remaining action obligation]
\label{p3f:rem:v1164-generalized-response-status}
The full-curvature encoding, its differential and the conditional causal
same-action response are closed mathematical compositions.  The missing
physical input is the explicit overlap/quotient identity between the
terminal-normalized \(f_{2\mid7}\) coefficient and the \emph{complete}
metric and mixed local/nonlocal heat functional on the selected member,
followed by the physical spectral pole and constraint analysis.  Neither the
unmarked positive curvature readout nor the ordinary heat coefficients alone
implies a \(C_6\) amplitude cap, an exact post-K48 dust first integral, a
healthy additional spin-two pole, or a strong-field bounded continuation.
\end{remark}


\subsection{scalaron matching, positive Schur no-free-force and raw response typing}
\label{p3f:sec:v1165-action-port-closure}
\begin{proposition}[Raw Laurent secant is not the terminal metric Euler response]
\label{p3f:prop:v1165-raw-euler-port}
Use the restored fixed-depth action of
Theorem~\ref{p3f:thm:v263-dressed-born-legendre-descent} and set
\(C=C_7\).  On the post-K16 lower curvature branch \(C<R\le C_2\),
\[
 \mathfrak P_{\rm raw}(R)=\frac{R^2+UC}{U+C},\qquad
 P_{\rm phys}(R)=C+\int_C^R x\chi_7(x)\,dx.
\]
They have identical terminal value but
\[
 \boxed{\mathfrak P'_{\rm raw}(C)=\frac{2C}{U+C}>0,
 \qquad P'_{\rm phys}(C)=C\chi_7(C)=0.}
\]
The inverse-Euler primitive of the raw secant is
\(f_{\rm raw}(R)=R-C+(R-C)^2/(U+C)\), whereas the restored coefficient
has the near-terminal expansion
\[
 \boxed{f_{\rm phys}(C+\delta)
 =\delta+\frac{\delta^3}{3C(U+C)}+O(\delta^4)}
 \qquad(\delta\downarrow0).
\]
Thus the two objects are not interchangeable metric action responses; the
physical Euler port is the fixed-depth-susceptibility integral, not the
undifferentiated Laurent readout.  This proposition supersedes any older
same-action inverse-Euler identification of the raw dressed secant.
\end{proposition}
\begin{proof}
If \(Rf_R-f=P\), then \(Rf_{RR}=P'\).  Apply this to
\(\mathfrak P_{\rm raw}\) and use terminal value and slope to integrate.
On the physical branch
\(\chi_7(C+\delta)=2\delta/[C(U+C)]+O(\delta^2)\).
Integrate twice with \(f(C)=0,f_R(C)=1\) to obtain the cubic start.
\end{proof}

\begin{theorem}[Stationary positive-block no-free-curvature-force test]
\label{p3f:thm:v1165-positive-no-free-force}
Work on a stationary gauge- and constraint-reduced real Hilbert patch of
one same action, where the eliminated block \(C_w=C_w^*>0\) is boundedly
invertible.  Let \(\phi\) denote the nondegenerate inherited scalaron
coordinate and suppose the quadratic Hessian is
\[
 H=\begin{pmatrix}K_f+\Delta K_{\rm dir}&B\\B^*&C_w\end{pmatrix}.
\]
Eliminating \(w\) gives the exact self-energy
\(\Sigma=B C_w^{-1} B^*\succeq0\), so the full scalaron operator
agrees with \(K_f\) if and only if
\[
 \boxed{\Delta K_{\rm dir}=B C_w^{-1}B^*.}
\]
If \(\Delta K_{\rm dir}=0\), exact preservation implies \(B=0\)
and hence the cross-source response \(j_{\phi,{\rm curv}}
=-BC_w^{-1}j_w\) vanishes for every admissible source.  More
quantitatively, for any scalaron test vector \(u\),
\[
 \boxed{\lvert\langle u,BC_w^{-1}j_w\rangle\rvert^2
 \le\langle u,\Sigma u\rangle
 \langle j_w,C_w^{-1}j_w\rangle.}
\]
These positivity implications are not asserted for an indefinite
Lorentzian retarded block; the operator matching identity remains the
necessary same-action test there with all retarded inverses and domains
made explicit.
\end{theorem}
\begin{proof}
Schur elimination gives
\(K_{\phi,{\rm eff}}=K_f+\Delta K_{\rm dir}-\Sigma\).
Since \(C_w>0\),
\(\langle u,\Sigma u\rangle
 =\|C_w^{-1/2}B^*u\|^2\), so \(\Sigma=0\) forces \(B^*=0\).
The displayed bound is Cauchy--Schwarz applied to
\(C_w^{-1/2}B^*u\) and \(C_w^{-1/2}j_w\).
\end{proof}

\begin{proposition}[Protected scalar and joint Dirac heat owners have unequal curvature jets]
\label{p3f:prop:v1165-scalar-dirac-heat-mismatch}
For the ordinary four-dimensional minimal scalar Laplacian, the metric
heat coefficients, modulo divergences, have invariant polynomials
\[
 a_2^{\rm sc}\sim\tfrac16R,\qquad
 a_4^{\rm sc}\sim\tfrac1{360}(5R^2+3C^2-\mathcal G).
\]
For the ordinary four-component torsion-free Dirac operator with
\(E=-R/4\), after the spin trace,
\[
 a_2^{\rm D}\sim-\tfrac13R,\qquad
 a_4^{\rm D}\sim\tfrac1{360}(-18C^2+11\mathcal G).
\]
The common exponential heat \emph{profile} and the Born-diagonal role
normalization do not identify these distinct Laplace-type operator jets.
Any equality of the protected scalar and joint Dirac \emph{physical action}
ports therefore requires the already generated projection/replacement and
coefficient transport, not simple trace multiplicity.  In particular
adding an independent minimal-scalar \(a_4\) correction with the same
Einstein-normalized heat scale \(\Lambda_g^2=96\pi^2U\) would shift
\(f_{RR}\) by \(1/(6\Lambda_g^2)\), spoiling exact terminal decoupling
unless compensated or excluded in the completed action.  This is an
\emph{additive-counterpart diagnostic}, not a second action prescription.
\end{proposition}

\begin{proposition}[Analytic fixed-rank heat trace cannot itself generate the K16 kink]
\label{p3f:prop:v1165-analytic-k16-obstruction}
Suppose an elliptic operator family \(D(s)^2\) is analytic on a common
compact, fixed-domain, trace-class heat patch with uniform analytic trace
control.  Its heat trace is analytic in \(s\).  If a smooth/analytic
constant-curvature test family \(R(s)\) crosses \(C_2\) transversely,
then \(f_{2\mid7}(R(s))\), whose third curvature derivative has a
nonzero finite jump, is not an analytic function of \(s\) at that crossing:
\[
 \boxed{[\partial_s^3 f_{2\mid7}(R(s))]_{s_2}
 =[f_{RRR}]_{C_2}(R'(s_2))^3\ne0.}
\]
Thus the exact restored branch cannot be the restriction of one analytic,
fixed-rank unprojected heat trace across K16.  The existing positive-part
coefficient/depth projection, a change of domain, or an explicitly certified
nonanalytic elimination must enter the action-reduction map.
\end{proposition}
\begin{proof}
Analytic holomorphic-family trace-class functional calculus makes
\(\operatorname{Tr}e^{-D(s)^2/\Lambda_g^2}\) analytic in \(s\).
On each side of the crossing the chain rule applies.  The first and
second curvature jets of \(f\) are continuous; therefore all terms in
its third parameter derivative cancel across the junction except
\([f_{RRR}](R')^3\).  A nonzero third derivative jump contradicts
analyticity of the unprojected trace.
\end{proof}

\begin{remark}[Zero crossings do not evade the analytic heat obstruction]
\label{p3f:rem:v2217-zero-crossing-heat-regularity}
For a finite analytic Hermitian family $D(s)$, the function
$\operatorname{Tr}e^{-D(s)^2/\Lambda_g^2}$ is analytic even when a signed
eigenvalue crosses zero: multiplication, the entire exponential and trace
preserve analyticity without requiring an inverse of $D$. The same holds
on the fixed-domain uniform analytic trace-class patch above. Thus the
cubic signed curve of Paper~I-C does not by itself generate the K16
$C^2$ but not $C^3$ action junction. Integrating its protected low mode
into a regular inverse or log determinant at the zero would instead be
invalid; retain the low mode and eliminate only its gapped complement.
\end{remark}

\begin{remark}[Physical closure still requires the complete action owner]
The preceding results establish exact derivative-typing and necessary
matching obstructions, but do not prove a complete source/metric spectral
reduction, a healthy retarded propagator, new matter abundance, or a
\(C_6\) ceiling.  Exact scalaron preservation with a nonzero new curvature
force requires the common action to generate the nontrivial matched direct
term; it cannot be imposed independently of the source mixing.
\end{remark}


\subsubsection{Neutral-photon Hessian and terminal optical no-mixing firewall}
\label{p3f:sec:v1166-neutral-optical-hessian}
\begin{proposition}[Zero photon--geometric mixed Hessian on a neutral branch]
\label{p3f:prop:v1166-neutral-photon-mixing}
Let $z$ collectively denote the geometric, scalaron and admissible boundary
variables on a regular homogeneous isotropic branch.  Suppose the actual
common action is differentiable to second order about $A=0$, the photon
background and external electromagnetic current vanish, and its retained
gauge-invariant bosonic sector has no term linear in an electromagnetic
fluctuation at any nearby neutral $z$.  Then
\[
 S[z,A]=S[z,0]+\tfrac12\langle A,K_\gamma[z]A\rangle+O(A^3),
 \qquad
 \boxed{H_{\gamma z}=H_{z\gamma}=0\quad(A=0).}
\]
Consequently, wherever the chosen retarded/reduced geometric inverse
exists on the physical constrained subspace, the photon Schur operator is
\[
 \boxed{K_{\gamma,\mathrm{Schur}}
 =H_{\gamma\gamma}-H_{\gamma z}(H_{zz}^{-1})_R H_{z\gamma}
 =H_{\gamma\gamma}.}
\]
In particular an independent scalaron-lapse Dirichlet-to-Neumann kernel
cannot, solely by quadratic photon--scalaron mixing, create a refractive
cone on this neutral branch.
\end{proposition}
\begin{proof}
The absence of the linear gauge term states $\delta_A S[z,0]=0$ on the
admitted neutral family.  Differentiating this identity along any
permitted $z$ variation gives $\delta_z\delta_A S[z,0]=0$.  The displayed
Schur identity follows algebraically.  Interactions of the form
$\delta z\,F^2$ are cubic in fluctuations and do not enter this cross
Hessian at $A=0$.
\end{proof}
\begin{remark}[Scope and unresolved direct-heat contribution]
This is a \emph{mixed-block} theorem, not a proof that the direct
joint-Dirac heat Hessian $H_{\gamma\gamma}$ has the minimal Maxwell
form at all orders.  Curvature-dependent tensorial, preferred-clock,
nonlocal or other directly generated terms in $H_{\gamma\gamma}$ must be
calculated from the \emph{same} common action and its actual internal
trace, localizer and Lorentzian reduction.  In particular, a numerical
six-derivative photon-cone coefficient has not been certified by the
pure-metric $a_6$ calculation and is not installed here as a Born--Hopf
prediction.  A charged or anisotropic background does not satisfy the
hypotheses of this proposition.
\end{remark}


\subsection*{Route-2 joint-parent closure: trace, boundary and signed source response}
\label{p3f:sec:v118-joint-parent-closure}

\begin{theorem}[On-shell conformal trace of the native local odd-field action]
\label{p3f:thm:v118-conformal-trace}
On the source-free common local domain of
Theorem~\ref{p3f:thm:v176-local-bosonic-euler-noether}, assume the scalar and
metric Euler equations, four-dimensional metric-independent internal gauge
pairing, and Hermitian odd field. With
$\alpha=(16\pi^2)^{-1}$ and $X=\operatorname{Tr}_F\Phi^2$,
\[
 \boxed{\kappa_J R=-2\alpha\Lambda_g^2 X.}
\]
This is not a field equation of the complete parent, which has additional
pure-metric, higher-order and state-conditioned owners beyond this local
subfunctional. If a separately generated constant $V_0$ is added to the
local scalar potential on an enlarged claim domain, the right-hand side is
replaced by $2V_0-2\alpha\Lambda_g^2X$.
\end{theorem}
\begin{proof}
The scalar Euler equation gives
$D^2\Phi=R\Phi/6-2\Lambda_g^2\Phi+2\Phi^3$. Write
$Y=\operatorname{Tr}_F(D\Phi)^2$ and $Z=\operatorname{Tr}_F\Phi^4$.
Then
\[
 \Box X=2Y+\frac R3X-4\Lambda_g^2X+4Z.
\]
The trace of the metric equation of
Theorem~\ref{p3f:thm:v176-local-bosonic-euler-noether}, using the
traceless classical four-dimensional Yang--Mills stress, is
\[
 -2\mathscr F R+6\Box\mathscr F
 =-2\alpha Y-4\alpha(Z-2\Lambda_g^2X),
 \qquad \mathscr F=\kappa_J-\alpha X/6.
\]
Substitution cancels the derivative terms, the $RX$ terms, and the quartic
terms, leaving the stated relation. A genuine constant $V_0$ changes the
metric trace by $-4V_0$ and the scalar equation by zero, giving the
conditional extension.
\end{proof}

\begin{proposition}[Hyperbolic odd-field mass Gram and stabilizer]
\label{p3f:prop:v118-hyperbolic-mass-gram}
On an admitted gauge-stable, trace-orthonormal linear real odd support
$\Phi=h^IT_I$ with $2\alpha\operatorname{Tr}_F(T_IT_J)=\delta_{IJ}$,
let $r^2=h^Ih_I$ and $\mathscr F=\kappa_J-r^2/12>0$.
The Einstein-frame field metric is
\[
 G^E_{IJ}=\frac{\kappa_J}{\mathscr F}\delta_{IJ}
 +\frac{\kappa_J}{12\mathscr F^2}h_Ih_J.
\]
For dimension $N\ge2$ it has constant sectional curvature
$-1/(12\kappa_J)$ before gauge quotient; at $N=1$ it is a flat radial
line. If $T_a$ is the antisymmetric generator of the action on this real
support, the gauge-boson mass Gram is
\[
 \boxed{\mathcal M_{ab}^{2,E}
 =G^E(T_ah,T_bh)=\frac{\kappa_J}{\mathscr F}
 \langle T_ah,T_bh\rangle.}
\]
The physical operator is normalized by the already generated gauge kinetic
pairing $Q_\rho$; its kernel is the stabilizer of $h$. No extra scalar or
gauge field is created by this metric.
\end{proposition}
\begin{proof}
The conformal transformation
$g_E=(\mathscr F/\kappa_J)g$ yields the displayed target metric.
With $L^2=12\kappa_J$, the polar substitution
$\rho=L\operatorname{artanh}(r/L)$ changes its line element to
$d\rho^2+L^2\sinh^2(\rho/L)d\Omega_{N-1}^2$, proving its curvature.
An infinitesimal gauge transformation is tangent to a sphere of constant
$r$, so $h\cdot T_ah=0$. The rank-one radial part of $G^E$ therefore
drops out of the gauge mass Gram. Positivity and $Q_\rho$-whitening preserve
the Gram kernel, which is the original stabilizer.
\end{proof}

\begin{proposition}[Conditional Dirichlet scalar--tensor boundary coefficient]
\label{p3f:prop:v118-dirichlet-scalar-ghy}
On an additionally constituted smooth non-null boundary, fix the induced
metric and odd field and use $\varepsilon=n^\mu n_\mu=\pm1$.
The well-posed Dirichlet completion of the *local curvature term*
$\int\sqrt{-g}\,\mathscr F(\Phi)R$ is
\[
 \boxed{S^{\mathrm{Dir}}_\partial
 =2\varepsilon\int_{\partial U}\sqrt{|h|}\,\mathscr F(\Phi)K.}
\]
Relative to the Einstein term already owned by $\kappa_J$, the fixed
scalar-dependent coefficient is
\[
 \boxed{\Delta S^{\mathrm{Dir}}_{\partial,\Phi}
 =-\frac{\varepsilon\alpha}{3}
 \int_{\partial U}\sqrt{|h|}\,K\operatorname{Tr}_F\Phi^2.}
\]
This is a variational necessity on the stated boundary class, not a proof
that the full joint Dirac spectral boundary condition or its heat coefficient
has been matched. The heat $\Box\Phi^2$ divergence and fermionic, null and
corner boundary contributions remain claim-specific.
\end{proposition}
\begin{proof}
The variation of $\sqrt{-g}\,\mathscr F R$ includes the usual
normal derivative of $\delta g$, multiplied by $\mathscr F$. The
variation of $2\varepsilon\sqrt{|h|}\mathscr F K$ cancels it when
$\delta h=\delta\Phi=0$ at the boundary.
Substituting $\mathscr F=\kappa_J-\alpha\operatorname{Tr}_F\Phi^2/6$
fixes the scalar contribution. The conformal Einstein-frame check uses
$\sqrt{|h_E|}K_E=(\mathscr F/\kappa_J)
\sqrt{|h_J|}[K_J+3n(\log\sqrt{\mathscr F/\kappa_J})]$.
\end{proof}

\begin{proposition}[Canonical squared-domain scalar Robin transport]
\label{p3f:prop:v3031-mixed-square-boundary-domain}
On a smooth non-null \emph{Euclidean} normal-parallel frozen product collar
use the canonical first-order Dirac boundary projector of
Remark~\ref{p3f:rem:v175-canonical-dirac-boundary-completion},
anti-Hermitian Clifford multiplication with $\gamma_n^2=-1$,
$\chi=\gamma_n\gamma_g$ and $\Pi_\pm=(1\pm\chi)/2$.
Let $D_{A,\Phi}=D_A+\gamma_g\otimes\Phi$ be the same generated joint
Dirac family, with a grading-even gauge connection and an internally odd,
Hermitian scalar $\Phi$; retain its gauge-covariant normal derivative.
For the canonical bag $\Pi_-\psi|_{\partial U}=0$, the associated
\emph{squared-operator} mixed boundary condition is
\begin{equation}
\label{p3f:eq:v3031-mixed-square-robin}
\boxed{\quad
 \Pi_-\psi|_{\partial U}=0,\qquad
 \Pi_+\bigl(\nabla_n^A-\tfrac12K_E-\Pi_+\Phi\Pi_+\bigr)
 \Pi_+\psi|_{\partial U}=0.\quad}
\end{equation}
In particular the canonical pure-Dirac Robin endomorphism
$S_0=-K_E\Pi_+/2$ acquires the \emph{nonfrozen} increment
\begin{equation}
\label{p3f:eq:v3031-scalar-robin-shift}
\boxed{\delta S=-\Pi_+\Phi\Pi_+.}
\end{equation}
This statement concerns the displayed collar and specified Dirac bag.
If gauge fields, coframe, tangential projector, boundary localizers or
source support vary, their domain transport and source contacts are
additional inputs, not consequences of this fixed-collar identity.
\end{proposition}
\begin{proof}
Squaring the first-order boundary problem requires both
$\Pi_-\psi|_{\partial U}=0$ and
$\Pi_-D_{A,\Phi}\psi|_{\partial U}=0$.
The tangential Clifford terms preserve the inherited bag splitting,
while $\gamma_g$ exchanges its two sectors.  Project the latter equation
onto the normal block and multiply by the inverse Clifford normal
($\gamma_n^{-1}=-\gamma_n$).  The pure-Dirac part is the inherited
$\nabla_n^A-K_E/2$ Robin operator and the new cross-block contributes
$-\Pi_+\Phi\Pi_+$, giving
Eq.~\eqref{p3f:eq:v3031-mixed-square-robin}.
The differential square has the already derived Laplace endomorphism
$E_{A,\Phi}=E_A-\Phi^2-i\gamma^\mu\gamma_gD_\mu\Phi$ in the source's
Hermitian-gamma convention; hence normal scalar gradients are retained
in the heat coefficients and cannot be discarded as a fixed-domain effect.
\end{proof}

\begin{proposition}[Restricted scalar mixed-boundary heat coefficients]
\label{p3f:prop:v3031-mixed-scalar-boundary-heat}
Under Proposition~\ref{p3f:prop:v3031-mixed-square-boundary-domain},
consider the \emph{ordinary unmarked} heat trace with unit localizer,
normal-parallel bag, fixed projector blocks and the already generated
nonmetric Born half-weight \emph{once}.  For the scalar boundary
monomials assume a tangentially constant, internally odd $\Phi$ at the
boundary and set
$X=\operatorname{Tr}_F\Phi^2$,
$\alpha=(16\pi^2)^{-1}$.
The restricted four-dimensional Euclidean scalar response is
\begin{align}
\label{p3f:eq:v3031-mixed-heat-a3}
 I_{E,\partial}^{a_3,\Phi^2}
  &=\frac{\Lambda_g}{2(4\pi)^{3/2}}
    \int_{\partial U}\!\sqrt{h_E}\,X,\\
\label{p3f:eq:v3031-mixed-heat-a4}
 I_{E,\mathrm{bulk}}^{a_4}\big|_{R_EX}
  &=\frac{\alpha}{6}\int_U\!\sqrt{g_E}\,R_EX,
 &I_{E,\partial}^{a_4}\big|_{K_EX}
  &=-\frac{\alpha}{3}\int_{\partial U}\!\sqrt{h_E}\,K_EX.
\end{align}
Incorrectly freezing $S=S_0=-K_E\Pi_+/2$ in the
scalar $a_4$ computation changes the last coefficient to
$+\alpha/3$; the difference in the two \emph{Euclidean spectral}
coefficients is $-2\alpha K_EX/3$.
When only the scalar varies normally on a flat, normal-parallel collar
($K_E=0$, no tangential source gradients), the explicit boundary $a_4$
packet additionally contains
\begin{equation}
\label{p3f:eq:v3031-mixed-normal-scalar-contact}
 I_{E,\partial}^{a_4,nX}
 =\frac{2\alpha}{3}\int_{\partial U}\!\sqrt{h_E}\,n(X),
\end{equation}
with $n$ the selected inward normal.  This is the explicit local heat
boundary term, \emph{before} combining with the interior $\Box X$
divergence and any orientation/continuation transport.  These formulas
are a restricted derivative expansion, not an exact full interacting
boundary action or a marked-heat theorem.
\end{proposition}
\begin{proof}
Use the universal mixed-boundary heat coefficients
\cite{Vassilevich2003HeatKernel} with
$S=-K_E\Pi_+/2-\Pi_+\Phi\Pi_+$ and
$E=E_A-\Phi^2-i\gamma^\mu\gamma_gD_\mu\Phi$.
On the indicated frozen collar, the $\Phi^2$ contribution at order
$a_3$ reduces to the internal/spin trace of
$96\chi E+192S^2$ divided by $384(4\pi)^{3/2}$.
Taking the spin trace and the once-applied Born half-weight gives
Eq.~\eqref{p3f:eq:v3031-mixed-heat-a3}.
For the $K_E\Phi^2$ coefficient the surviving mixed $a_4$
monomials have trace numerator
\[
 \operatorname{tr}(120EK_E+720SE+480S^2K_E+480S^3).
\]
With the scalar Robin increment this yields
$-\alpha K_EX/3$, while substituting $S_0$ gives
$+\alpha K_EX/3$, after the same spin trace, factor $1/360$
and Born half-weight.  The bulk $a_4$ trace from
$60R_EE+180E^2$ gives $+\alpha R_EX/6$.
The normal-only $a_4$ boundary terms contain the inward derivative
of $E$ and the Robin/scalar cross-contact; their finite Clifford trace
with the stated collar guards yields $2\alpha\,n(X)/3$.
No trace-smearing identity for a noncommuting relative-grade insertion
has been used.
\end{proof}

\begin{corollary}[Scalar boundary two-jet and the unclosed physical port]
\label{p3f:cor:v3031-boundary-scalar-two-jet}
For fixed collar data of Proposition~\ref{p3f:prop:v3031-mixed-scalar-boundary-heat}
let $\Phi(u,v)=uX_1+vX_2$ with admissible, internally odd source
operators $X_1,X_2$.  The restricted $a_3$ term has the exact mixed jet
\begin{equation}
\label{p3f:eq:v3031-boundary-two-source-jet}
 \boxed{\left.\partial_u\partial_v I_{E,\partial}^{a_3}
 \right|_{u=v=0}
 =\frac{\Lambda_g}{2(4\pi)^{3/2}}
 \int_{\partial U}\!\sqrt{h_E}\,
 \operatorname{Tr}_F\{X_1,X_2\}.}
\end{equation}
Freezing the squared-domain Robin operator at $\Phi=0$ omits this
source-variation contribution.  In a \emph{same-sign Euclidean}
curvature-only Dirichlet comparison, the bulk coefficient
$+\alpha R_EX/6$ requires $+\alpha K_EX/3$ as its local
Dirichlet completion; hence the restricted spectral curvature packet
with coefficient $-\alpha K_EX/3$ differs by
$-2\alpha K_EX/3$ before any continuation/source-identity map.
The Lorentzian boundary coefficient of
Proposition~\ref{p3f:prop:v118-dirichlet-scalar-ghy} remains valid on
its distinct physical Dirichlet domain.  Its identification with the
Euclidean mixed heat, the $a_3$ source surface term, the normal
$\Box X$ contact, localizer and fermion/boundary variations, and any
marked $1/96$ projection requires a separately certified
\emph{whole-action, continuation and source-domain map}.
No physical Newton-matching status follows from this corollary.
\end{corollary}
\begin{proof}
Differentiate the quadratic expression
$X=\operatorname{Tr}_F\Phi^2$ twice to obtain
$\partial_u\partial_vX=\operatorname{Tr}_F\{X_1,X_2\}$.
The Euclidean comparison coefficient is twice the bulk
$R_EX$ coefficient under the same-sign Dirichlet curvature variation.
Subtract the restricted spectral coefficient in
Eq.~\eqref{p3f:eq:v3031-mixed-heat-a4}.
The distinct boundary and source-domain qualifications follow from
the hypotheses of the two propositions, not from a cross-signature
identification.
\end{proof}

\begin{proposition}[Complete localizer-aware mixed heat Hessian]
\label{p3f:prop:v118-localizer-aware-hessian}
Let $D=D_{A,\Phi}(g)$ be the same selected Euclidean joint Dirac operator,
$A=D^2/\Lambda_g^2$, $L=\varphi_U$ its analytic localizer, and
$S=\operatorname{Tr}(Le^{-A})$. For actual variations $U,V$ including
coframe and contact variation, define
\[
 A_U=\Lambda_g^{-2}\{D,D_U\},\quad
 A_{UV}=\Lambda_g^{-2}(\{D_U,D_V\}+\{D,D_{UV}\}).
\]
Define
\[
\begin{aligned}
 \mathcal W_{UV}={}&\int_{0\le s\le t\le1} ds\,dt
 \bigl(e^{-sA}A_Ue^{-(t-s)A}A_Ve^{-(1-t)A}
       +e^{-sA}A_Ve^{-(t-s)A}A_Ue^{-(1-t)A}\bigr)\\
 &-\int_0^1ds\,e^{-sA}A_{UV}e^{-(1-s)A}.
\end{aligned}
\]
Then the *exact* localizer-aware second derivative is
\[
\boxed{\begin{aligned}
 S_{UV}={}&\operatorname{Tr}(L_{UV}e^{-A})
 -\int_0^1ds\,\operatorname{Tr}(L_Ue^{-sA}A_Ve^{-(1-s)A})\\
 &-\int_0^1ds\,\operatorname{Tr}(L_Ve^{-sA}A_Ue^{-(1-s)A})
 +\operatorname{Tr}(L\mathcal W_{UV}).
\end{aligned}}
\]
It reduces by trace cyclicity to
Proposition~\ref{p3f:prop:v1164-heat-hessian-no-insertion} only on the
unlocalized trace (or an otherwise justified commutation domain). The
Born nonmetric half-weight applies once to the whole source packet; no
independent mixed-Hessian coefficient is permitted.
\end{proposition}
\begin{proof}
The two derivative identities for $A$ follow by differentiating $D^2$.
Twice applying Duhamel's operator-exponential formula produces the two
ordered integrals, the $A_{UV}$ contact and both localizer derivatives.
Trace cyclicity collapses the ordered pair only when the relevant
localizer commutation hypothesis is available.
\end{proof}

\begin{proposition}[Off-shell gauge Ward correction and protected signed low-mode bound]
\label{p3f:prop:v118-ward-and-low-mode}
On a gauge-invariant claim domain, let $\mathcal E=dS$ and let
$\mathcal R_\xi$ be the infinitesimal gauge-orbit tangent. Then
\[
 \boxed{H(U,\mathcal R_\xi)
 +\langle\mathcal E,D_U\mathcal R_\xi\rangle=0.}
\]
Consequently the Hessian annihilates gauge-orbit directions only after
full stationarity is imposed. On the selected compact weak spatial Dirac
block at the Born midpoint, let $P$ be the protected zero-mode projector,
$Q=1-P$, and $g=\sqrt3/R$ the gap of $QD_0Q$. For a self-adjoint
admitted source perturbation $V$ and
$\|QVQ\|+|z|<g$, the complement inverse exists with
\[
 \|(Q(D_0+V-z)Q)^{-1}\|\le(g-\|QVQ\|-|z|)^{-1}.
\]
The exact signed retained low-mode operator is
\[
 \boxed{K_P(z)=P(D_0+V-z)P
 -PVQ(Q(D_0+V-z)Q)^{-1}QVP.}
\]
Its self-energy norm is bounded above by
$\|PVQ\|^2/(g-\|QVQ\|-|z|)$.
This is a finite/gapped spatial-operator result, not a proof of a new
four-dimensional particle or of a regulated determinant phase.
\end{proposition}
\begin{proof}
Differentiate the off-shell gauge identity
$\langle dS,\mathcal R_\xi\rangle=0$ to obtain the first formula.
The native compact gap and the norm bound on the resolvent follow from
Neumann inversion about $QD_0Q$. The standard block-elimination identity
gives the Schur operator and the norm estimate. Background/projector gauge
transport is required when comparing different gauge representatives.
\end{proof}

\begin{remark}[An evaluated signed cubic curve]
\label{p3f:rem:v2217-cubic-dirac-application}
On the same round fundamental doublet, the admitted even-under-inversion
profile $u=\epsilon(q_1^2-1/4)$ has
$\lambda_{\rm low}=9\epsilon^3/(8192R)+O(\epsilon^5)$ by
Paper~I-C, Theorem~\ref{p1b:thm:v2217-nonconstant-cartan-dirac}.
This curve is different from an inversion-odd profile protecting an exact
zero at every amplitude. Small gauge transformations to the actual local
slice preserve its signed eigenvalue. Its quadratic static shift vanishes,
but its signed cubic shift is nonzero. The uninserted even heat is even in
$\epsilon$; the isolated low heat term starts at sixth order, while gapped
modes may respond quadratically. This licenses no separate cubic bosonic
term, occupied new particle, completed anomaly transition or baryon yield.
\end{remark}


\begin{remark}[Route-2 scope and causal source completion]
\label{p3f:rem:v118-parent-first-full-action-frontier}
The local trace, scalar boundary, hyperbolic mass Gram, and signed spatial
Schur results above close their stated mathematical action domains; they do
not give an exact global Lorentzian nonlocal Born parent, a selected
K16 depth-active off-shell joint-Dirac operator, a fermionic determinant or
Pfaffian phase, or an occurring physical state. On an admitted finite
self-adjoint source Hamiltonian with chosen initial density operator
$\rho_i$, the exact causal two-history source functional is
$\mathcal I[q_+,q_-]=\operatorname{Tr}(U[q_+]\rho_iU[q_-]^\dagger)$;
it obeys $\mathcal I[q,q]=1$ and yields retarded commutator response and
nonnegative symmetrized source covariance. This finite identity licenses
no new dynamical coefficient and does not establish the corresponding
unrestricted continuum quantum measure. The inherited $f_{2\mid7}$ scalaron
remains exact on its certified restricted action domain; Route 2 does
not require full-parent scalar Hessian corrections to vanish. K48 structural
$I_{48}=UC_6$ remains separate from a demonstrated Hilbert density.
\end{remark}


\section{Cap curvature jets, coherent UV target, and common-action current tests}
\label{p3f:sec:v229-cap-uv-hessian}

\begin{proposition}[Same-chart cap/UV Hessian separation]
\label{p3f:prop:v229-cap-vs-coherent-uv-hessian}
The signed cap parent of
Theorem~\ref{p3f:thm:v224-reference-cap-parent},
\[
 f_c(R)=PR+B(R)R^2-\frac{B_I^2R_s^2}{B(R)^2},\qquad
 B(R)=B_U+\frac{2(B_I-B_U)}{1+\sqrt{1+(R/R_s)^2}},
\]
has the exact curvature jets
\[
 \boxed{f_c(0)=-R_s^2,\quad f_c'(0)=P,\quad
 \frac12 f_c''(0)=B_I-\frac{B_I-B_U}{2B_I}.}
\]
If this \emph{same local reduced scalar} is identified, after dividing
by its Einstein coefficient $P>0$, with the separate coherent quadratic
target $R-C_7+(Q/U)R^2$, then necessarily
\[
 \boxed{\frac{B_I-(B_I-B_U)/(2B_I)}{P}=\frac QU,
 \qquad R_s^2/P=C_7.}
\]
The cap anchors with the additional comparison $P=U$ do not satisfy
the first equality (the ratio of the two quadratic coefficients is about
$4.71\times10^{49}$).  Consequently the two local scalar responses may
\emph{not} be identified as one same-chart reduced function under those
assumptions.  Their distinct origins do not exclude different reductions
of one more general common action if the corresponding dynamic or
boundary matching map is derived.
\end{proposition}
\begin{proof}
$B(R)$ is even, $B(0)=B_I$, and
$B(R)=B_I-(B_I-B_U)R^2/(4R_s^2)+O(R^4)$.
Expanding $-B_I^2R_s^2/B(R)^2$ adds
$-(B_I-B_U)R^2/(2B_I)$ to the quadratic coefficient.
Normalization by $P$ and coefficient comparison yield the two necessary
identities.  The numerical ratio evaluates the printed cap parameters
under the explicitly stated extra identification $P=U$; it is not a
physical mass measurement.
\end{proof}

\begin{corollary}[Actual curvature-Schur and coherent-oscillation test]
\label{p3f:cor:v229-actual-curvature-schur-test}
Suppose an actually consumed local parent $W(R,z)$ has a regular stationary
elimination $W_z(R,z_*(R))=0$ with invertible $H_{zz}$.  Then
\[
 \boxed{f_{{\rm red},RR}=H_{RR}-H_{Rz}H_{zz}^{-1}H_{zR}.}
\]
A claimed same-chart Einstein-normalized UV quadratic coefficient
$Q/U$ therefore requires this Schur Hessian, after the actual
normalization, to equal $2Q/U$ at the stated expansion point.  For an
exact quadratic on an interval, all higher local curvature jets must
also satisfy the required vanishing identities.  In a dynamical scalaron
claim, the physical oscillation frequency must instead be read from
the full kinetic-normalized retarded/inverse-propagator pole; a static
local Schur Hessian alone does not establish that pole.
\end{corollary}
\begin{proof}
Differentiate $W_z(R,z_*(R))=0$ to get
$z_*'=-H_{zz}^{-1}H_{zR}$, then differentiate
$f_{\rm red}=W(R,z_*(R))$ twice.  The final statements distinguish
local Taylor data from an entire action and its propagating response.
\end{proof}

\begin{remark}[Positive Euclidean auxiliary-block diagnostic]
\label{p3f:rem:v229-positive-schur-diagnostic}
For a restricted \emph{Euclidean} quadratic family with the same retained
coordinate, $D(s)=D_0+sG$ with $D_0>0$, $G\succeq0$ and constant $A,C$,
the Schur return $K(s)=A-CD(s)^{-1}C^\dagger$ satisfies
\[
 K'(s)=CD(s)^{-1}GD(s)^{-1}C^\dagger\succeq0.
\]
Such a simple positive fixed-coupling interpolation cannot explain a
static Hessian larger than its high-frequency counterpart.  This
conditional monotonicity is \emph{not} asserted for the uncomputed
Lorentzian causal parent, moving retained coordinates or boundary
contacts.  Neither a fitted coefficient nor an imposed Starobinsky
identity substitutes for the actual geometric response.
\end{remark}

\begin{remark}[K16 current identification is an equality test]
\label{p3f:rem:v229-k16-current-equality}
The Schur vertex $\partial_\phi K_P$ follows from differentiating the
complete same-action block.  Identifying scalaron work with the
already normalized direct $K16$ boundary packet requires equality of
the \emph{actual} currents on the same state and physical clock,
including all moving-frame, kinetic, connection and boundary terms.
Exclusive descendant ownership forbids counting the same packet twice;
it neither establishes this equality nor excludes independently sourced
bulk pumping.
\end{remark}

\subsection{Exact bare vacuum trace and scalaron stability certificate}
\label{p3f:subsec:v2188-bare-vacuum-trace}

\begin{theorem}[Exact bare two-root vacuum and trace mass]
\label{p3f:thm:v2188-bare-vacuum-trace}
On the already generated \emph{bare, undressed} same-action branch of
Theorem~\ref{p3f:thm:bare-born-potential-criterion}, let
$0<C<U$ and
\[
 f_{\mathrm b}(R)=(R-C)+\frac{(R-C)^2}{U+C}.
\]
Use the four-dimensional source-free, constant-curvature metric trace equation
$R f_{\mathrm b,R}-2f_{\mathrm b}=0$, without an additional independent
cosmological constant. Then
\[
 \boxed{f_{\mathrm b,R}=1+\frac{2(R-C)}{U+C},\quad
 f_{\mathrm b,RR}=\frac2{U+C}>0,\quad
 R f_{\mathrm b,R}-2f_{\mathrm b}
       =\frac{2CU-(U-C)R}{U+C}.}
\]
The trace equation has exactly one algebraic constant-curvature solution,
\[
 \boxed{R_*=\frac{2CU}{U-C},\qquad
         f_{\mathrm b,R}(R_*)=\frac{U+C}{U-C}>0.}
\]
This solution lies in the closed bare interpolation interval $[C,U]$ if
and only if $U\ge3C$, and in its interior if and only if $U>3C$.
The trace-linearized scalaron mass-squared parameter on every admitted
constant-curvature background is independent of $R$:
\[
 \boxed{m_{\mathrm{alg},\mathrm b}^{2}
 =\frac{f_{\mathrm b,R}-Rf_{\mathrm b,RR}}{3f_{\mathrm b,RR}}
 =\frac{U-C}{6}>0.}
\]
In particular, $R=C$ is \emph{not} a vacuum of this bare action:
$(R f_{\mathrm b,R}-2f_{\mathrm b})|_{R=C}=C>0$.
Multiplying the entire bare $f$ by the positive normalizer $U$ leaves
both $R_*$ and the mass ratio above unchanged.
\end{theorem}
\begin{proof}
Differentiate the printed bare coefficient twice. Substitution into
$Rf_R-2f$ gives its displayed affine function, with strictly negative
slope $-(U-C)/(U+C)$; its unique zero is $2CU/(U-C)$.
At $R=C$ the trace residual is $C$, while at $R=U$ it is
$U(3C-U)/(U+C)$, yielding the stated interval test. Inserting
$f_R$ and $f_{RR}$ in the trace-linearized formula of
Theorem~\ref{p3f:thm:trace-linearized-scalaron-mass} cancels the
$R$-dependence and yields $(U-C)/6$. Direct evaluation gives the
positive slope at $R_*$. Common positive action multiplication cancels
from the trace equation and the mass ratio.
\end{proof}

\begin{corollary}[Born-even response does not imply reciprocal vacua]
\label{p3f:cor:v2188-bare-vacuum-born-scope}
The normalized Legendre response of
Corollary~\ref{p3f:cor:v2187-normalized-born-potential-scope}
obeys $g_{\rm eff}(R)=g_{\rm eff}(CU/R)$, but this alone does not map
solutions of the metric vacuum trace equation to solutions. Indeed,
\[
 \boxed{\frac{CU}{R_*}=\frac{U-C}{2}.}
\]
The dual is another vacuum of the same bare action only when it equals
$R_*$, equivalently when $U/C=3+2\sqrt2$. Otherwise the vacuum is
not fixed by reciprocal inversion.
\end{corollary}
\begin{proof}
The first identity is substitution. The affine vacuum trace has a
unique zero. Equality $CU/R_*=R_*$ gives $(U-C)^2=4CU$;
with $U/C>1$, the only admissible root is $3+2\sqrt2$.
\end{proof}

\subparagraph{Physical and constitutional scope.}
This is \textsc{corpus-closure} algebra on the existing bare coefficient and
its already licensed metric trace formula; it introduces neither a new
field nor a curvature/vacuum selector. An algebraic root outside $[C,U]$
is not promoted to a realized vacuum outside the source's action domain.
The mass expression is a trace-linearized parameter; positivity does not
assert existence of the full metric solution, a realized stationary history,
cosmological attraction, or positivity of every coupled perturbation.
For the dressed $2\mid7$ coefficient, the physical $K48$ history, or an
action with a separate $\Lambda_{\rm ren}$ term, the complete relevant
trace equation must be evaluated afresh.


\begin{corollary}[Laurent--Legendre terminal scalar-remnant representative]
\label{p3f:cor:laur17-terminal-ren-representative}
Let \(P_{2\mid7}=RF-f\) be the generated Euler response reconstructed in
Theorem~\ref{p3f:thm:v263-dressed-born-legendre-descent}.  At the rank-seven
endpoint,
\[
 P_{2\mid7}(C_7)=C_7.
\]
For any same-action inverse-Euler representative \(\widetilde f\) satisfying
\((R\partial_R-1)\widetilde f=P_{2\mid7}\), write
\[
 Z_7:=\widetilde f_R(C_7),\qquad F_7:=\widetilde f(C_7).
\]
Then the endpoint Euler equation gives
\begin{equation}
\label{p3f:eq:v266-terminal-euler-jet}
 \boxed{C_7Z_7-F_7=C_7,\qquad F_7=C_7(Z_7-1).}
\end{equation}
Hence the unique affine terminal Einstein continuation with the same one-jet is
\begin{equation}
\label{p3f:eq:v266-terminal-kernel-invariant-intercept}
 \boxed{
 f_{2\mid7}^{\rm term}(R_{\rm g})
 =F_7+Z_7(R_{\rm g}-C_7)
 =Z_7R_{\rm g}-C_7.}
\end{equation}
For \(Z_7>0\), the conventional cosmological-constant coordinate and the
four-dimensional vacuum curvature are
\begin{equation}
\label{p3f:eq:v267-terminal-general-lambda-ds}
 \boxed{\Lambda_{\rm DE}=\frac{C_7}{2Z_7},\qquad
 R_{\rm g,dS}=\frac{2C_7}{Z_7}.}
\end{equation}
The Euler-kernel freedom \(\widetilde f\mapsto\widetilde f+\lambda R\)
changes \(Z_7\) but leaves the zero-order intercept \(-C_7\) invariant.
On the retained terminal-normalized representative \(Z_7=1\),
\begin{equation}
\label{p3f:eq:laur17-postfreeze-total-action}
 \boxed{
 f_{2\mid7}^{\rm term}(R_{\rm g})=R_{\rm g}-C_7
 =R_{\rm g}-2\Lambda_{\rm eff},
 \qquad \Lambda_{\rm eff}=\frac{C_7}{2}.}
\end{equation}

Paper~III-C identifies the exponent seven as the rank drop from the protected
Clifford module to the infrared scalar remnant, and Paper~III-E transports the
same generated hierarchy owner to \(C_7=U\alpha_\sigma^{-7}\).  Therefore
\(-C_7\) is the zero-order gravitational representative of that terminal
Clifford scalar-remnant scale.  No independent \(\mathrm{TERMREN}_7\) datum
is required on this branch, and adding a second subtraction would double
count the same endpoint constituent.  The raw cross-owner normalization is fixed downstream by
\(U_J^{\rm raw}=96U\).  The primitive curvature owner is fixed, while
identification of the raw and scalar action coupling coordinates with one
physical readout still consumes
Proposition~\ref{p3f:prop:v2211-newton-owner-matching}.
Absolute unit realization follows only after that physical match.
The general vacuum curvature is
\(R_{\rm g,dS}=2C_7/Z_7\); on the slope-one representative it reduces to
\(2C_7\).  The terminal sector carries no scalaron
Hessian.
\end{corollary}

\begin{proof}
At \(R=C_7\), the generated coefficient has \(D_2(C_7)=1\) and the rank-seven
Born response has \(\mathcal H_7^{\rm act}(C_7)=C_7\), so
\(P_{2\mid7}(C_7)=C_7\).  Evaluating
\(R\widetilde f_R-\widetilde f=P_{2\mid7}\) at the endpoint proves
Eq.~\eqref{p3f:eq:v266-terminal-euler-jet}.  Affine continuation of that
one-jet gives
\(F_7+Z_7(R_{\rm g}-C_7)=Z_7R_{\rm g}-C_7\), proving that the intercept is
independent of the Euler kernel.  Writing \(Z_7R_{\rm g}-C_7=Z_7(R_{\rm g}-2\Lambda_{\rm DE})\)
gives \(\Lambda_{\rm DE}=C_7/(2Z_7)\); the vacuum trace equation gives
\(R_{\rm g,dS}=2C_7/Z_7\).  Setting \(Z_7=1\) recovers the retained
terminal-normalized representative and \(\Lambda_{\rm eff}=C_7/2\).  The
scalar-remnant interpretation is the composition of Paper~III-C's rank-seven
endpoint theorem with Paper~III-E's same-owner hierarchy inheritance; no new
zero-order constituent is inserted.
\end{proof}

\begin{corollary}[Native cyclic zero-order completion and \(\mathrm{REN}_0\) nonduplication]
\label{p3f:cor:v2118-native-cyclic-ren0-firewall}
On the generated native cyclic terminal branch of
Corollary~\ref{p3f:cor:laur17-terminal-ren-representative}, the
terminal-normalized metric action is already
\[
 \boxed{f_{2\mid7}^{\rm term}(R_{\rm g})=R_{\rm g}-C_7
 =R_{\rm g}-2\Lambda_{\rm eff},\qquad \Lambda_{\rm eff}=C_7/2.}
\]
Consequently an additional independently chosen
\(\mathrm{REN}_0(\Lambda_{\rm ren})\) term is not a datum of this branch.
Appending one would replace the derived zero-order coefficient by
\(-C_7-2\Lambda_{\rm ren}\) and therefore define a different action unless a
compensating terminal shift were generated; no such independent shift occurs
in the native cyclic lineage.  Equivalently, the generic TERM--REN quotient
remains available as an enlarged/local-action comparison atlas, but it is not
a residual native-cyclic cosmological input.

Moreover
\[
 \frac{d^2 f_{2\mid7}^{\rm term}}{dR_{\rm g}^2}=0,
\]
so no nonlinear curvature Hessian survives after the terminal handoff.  This
endpoint statement does not delete the earlier active nonlinear branch or the
fresh-daughter \(R^2\) owner; it states only that those higher-order dynamics
have terminated on the terminal Einstein branch.
\end{corollary}
\begin{proof}
The preceding terminal-intercept corollary derives the affine continuation and
its intercept from the same-action inverse-Euler endpoint data and proves that
no independent \(\mathrm{TERMREN}_7\) constituent is required.  Subtracting
an additional constant changes the complete action-level zero-order invariant
and hence is not a representative change unless accompanied by the exact
TERM--REN compensating shift.  The native cyclic construction supplies no such
second endpoint datum.  Affinity gives the displayed vanishing second
derivative, while the pre-terminal nonlinear action remains unchanged.
\end{proof}

\begin{corollary}[Historical rank-two K16/K48 local-jet classification]
\label{p3f:cor:transported-k16-k48-action-junction}
For the restored fixed-depth Laurent--Legendre action,
\[
 \boxed{[f_{RR}]_{C_2}=0,\qquad
 [f_{RRR}]_{C_2}=D_2'(C_2^+)\chi_7(C_2)>0.}
\]
Thus K16 is a finite third-jet coefficient freeze with no Hessian atom or
jump.  K48 lies on the smooth lower branch and has the finite positive
susceptibility $f_{RR}(C_6)=\chi_7(C_6)$.  At K56,
$f_{RR}\to0^+$ continuously.  No independent microscopic smoothing width is
licensed by any of these statements.
\end{corollary}

\begin{theorem}[Historical rank-two GP-exact matching at the K16 coefficient freeze]
\label{p3f:thm:v266-gp-exact-k16-junction-condition}
Let an actual classical history be piecewise regular across a non-null K16
interface $\Sigma_2$ with $R=C_2$.  For the restored action,
\[
 [F]_{\Sigma_2}=0,\qquad [\chi]_{\Sigma_2}=0,
 \qquad \chi=f_{RR}>0.
\]
In the absence of an independently generated surface stress/curvature layer,
there is therefore no action-generated delta source at K16.  The exact
source-free matching law is ordinary $F$-flux continuity,
\[
 \boxed{[\nabla_nF]_{\Sigma_2}=0,}
\]
and because the same nonzero $\chi_{16}$ occurs on both sides, this is
\[
 \boxed{(\nabla_nR)_+=(\nabla_nR)_-.}
\]
A transverse $C^1$ curvature crossing is consequently allowed without a
surface impulse.  The coefficient freeze is recorded by the finite one-sided
change in $f_{RRR}$, which changes higher time derivatives but does not
refract the curvature rate.
\end{theorem}

\begin{proof}
The preceding theorem gives a $C^2$ metric coefficient at C2.  Hence $F$ is
$C^1$ as a function of $R$ and its normal derivative is
$\nabla_nF=\chi\nabla_nR$ with the same finite nonzero $\chi_{16}$ on both
sides.  The distributional second derivative of $F$ has no delta term when
this flux is continuous.  The equivalence with continuity of $\nabla_nR$
follows immediately.  The finite third-jet statement is the already derived
coefficient-character derivative kink.
\end{proof}

\begin{theorem}[Positive multiplicative kernel of the TERM--REN action]
\label{p3f:thm:term-ren-multiplicative-action-quotient}
On the same classical action domain, write
\[
 S_g[f,\Lambda_{\rm ren},G]
 :=\frac{c_{\rm op}^3}{16\pi G}
   \int\sqrt{-g}\,[f(R)-2\Lambda_{\rm ren}]\,d^4x.
\]
For every \(a>0\), the simultaneous transformation
\begin{equation}
\label{p3f:eq:term-ren-positive-multiplicative-action}
 (f,\Lambda_{\rm ren},G)
 \longmapsto
 (af,a\Lambda_{\rm ren},aG)
\end{equation}
leaves \(S_g\) exactly unchanged.  It is therefore an exact comparison-kernel
action, not permission to rescale one coefficient in isolation.

On the branch \(Z_C=f'(C)>0\), every positive orbit has the unique
representative with terminal slope one.  Its invariant normalized data are
\begin{equation}
\label{p3f:eq:normalized-term-ren-data}
 \boxed{
 \widehat f:=\frac f{Z_C},\quad
 \widehat\Lambda_{\rm ren}:=\frac{\Lambda_{\rm ren}}{Z_C},\quad
 G_{\rm term}:=\frac G{Z_C},\quad
 \widehat{\overline F}_C
 :=\frac{F_C-2\Lambda_{\rm ren}}{Z_C},\quad
 \widehat\gamma_\sigma:=\frac{\gamma_\sigma}{Z_C}.}
\end{equation}
Here \(\widehat f'(C)=1\).  Thus \(Z_C\) is not an independent structural
coordinate of the uncalibrated same-action quotient.  If \(G\) has already
been fixed as an external physical Newton calibration, transformation
\eqref{p3f:eq:term-ren-positive-multiplicative-action} is no longer an
admissible change within that calibrated slice; then terminal slope
normalization is a calibration statement rather than a structural
identification.
\end{theorem}

\begin{proof}
Direct substitution gives
\[
 \frac1{aG}\,[af(R)-2a\Lambda_{\rm ren}]
 =\frac1G\,[f(R)-2\Lambda_{\rm ren}]
\]
pointwise, so the gravitational functional is identical for every metric and
the unchanged matter action leaves the full action identical.  Along an orbit,
\(Z_C\mapsto aZ_C\).  If \(Z_C>0\), choosing \(a=Z_C^{-1}\) gives slope one,
and uniqueness follows because no other positive \(a\) does so.  Each
quantity in \eqref{p3f:eq:normalized-term-ren-data} is unchanged by the
simultaneous action; the Hessian calibration scales with \(f''\), hence
\(\gamma_\sigma\mapsto a\gamma_\sigma\).  This proves the quotient statement.
It is invariant descent in the sense of Paper~II-A,
Theorem~\ref{p2a:thm:geometric-defect-repair-admission}.
\end{proof}

\begin{corollary}[PQTI relative coefficient on the same-action quotient]
\label{p3f:cor:pqti-same-action-invariance}
Assume the imported \((\mathrm{CRP4})+(\mathrm{PQTI})\) branch of
Paper~III-E.  The relative coefficient between its Einstein--Hilbert scalar
insertion and its primitive nonlinear insertion is invariant under the
positive action kernel
\eqref{p3f:eq:term-ren-positive-multiplicative-action}.  Hence on the
terminal-slope-one representative
\[
 \boxed{\widehat\gamma_\sigma
 =16\sqrt{\alpha_\sigma}},
 \qquad
 \boxed{\kappa_U=\frac{R_+}{\Lambda_{\rm EH}^2}=1}.
\]
The common action preserves the generated PQTI normalization.  It still
cannot be used in reverse to derive PQTI or CRP4 from an already assembled
response.  Its contact identification consumes
Theorem~\ref{p3b:thm:v2211-action-contact-certificate}; after that test,
the earlier action coefficient compiler and its second variation provide
the conditional same-owner transport.
\end{corollary}

\begin{proof}
The transformation
\eqref{p3f:eq:term-ren-positive-multiplicative-action} multiplies the complete
gravitational integrand, and therefore both protected insertions, by the same
positive scalar.  Their relative trace coefficient is unchanged.  Equivalently,
\(\gamma_\sigma\) and \(Z_C\) both acquire the factor \(a\), so
\(\widehat\gamma_\sigma=\gamma_\sigma/Z_C\) is invariant.  Paper~III-E
Theorem~\ref{p3e:thm:pqti-quadratic-normalization} evaluates that invariant
as displayed.  The chronology statement follows from the dependency order of
the construction and blocks a circular discharge.
\end{proof}

\begin{corollary}[REN-aware terminal constant-curvature condition]
\label{p3f:cor:term-ren-terminal-vacuum}
If \(R=C\) is a four-dimensional constant-curvature vacuum of the same action,
so that \(R_{\mu\nu}=\tfrac C4g_{\mu\nu}\), \(T_{\mu\nu}=0\), and all
derivatives of \(f'(R)\) vanish, then
\[
 \boxed{CZ_C-2F_C+4\Lambda_{\rm ren}=0}
\]
or, equivalently,
\[
 \boxed{\overline F_C=\frac{CZ_C}{2}}.
\]
Thus the source pair \((F_C,Z_C)=(0,1)\) together with
\(\Lambda_{\rm ren}=0\) is not a terminal constant-curvature vacuum when
\(C>0\).  If the displayed \(G\) is calibrated as the physical Newton
constant at this terminal point, the normalization \(Z_C=1\) is calibration,
not an additional structural selector.
\end{corollary}

\begin{proof}
At a constant-curvature vacuum the derivative term in
Theorem~\ref{p3f:thm:general-local-fR-variation} vanishes.  Substitution of
\(R_{\mu\nu}=(C/4)g_{\mu\nu}\), \(T_{\mu\nu}=0\),
\(f(C)=F_C\), and \(f'(C)=Z_C\) gives
\[
 \left(\frac{CZ_C}{4}-\frac{F_C}{2}
       +\Lambda_{\rm ren}\right)g_{\mu\nu}=0.
\]
Nondegeneracy of the metric and multiplication by four give the first boxed
identity.  Substitution of
\(\overline F_C=F_C-2\Lambda_{\rm ren}\) gives the second.  The source
specialization has residual \(C\), so it fails the vacuum condition whenever
\(C>0\).
\end{proof}

\begin{corollary}[Normalized terminal-vacuum datum]
\label{p3f:cor:normalized-terminal-vacuum-data}
On the positive-slope same-action quotient, the terminal
constant-curvature vacuum condition is exactly
\[
 \boxed{\widehat{\overline F}_C=\frac C2.}
\]
Thus the raw triple \((F_C,Z_C,\Lambda_{\rm ren})\) contributes only one
normalized zero-order endpoint coordinate after the additive and positive
multiplicative action kernels are both removed.  The remaining
\(G_{\rm term}\) and, on the unrestricted MRESP atlas,
\(\widehat\gamma_\sigma\) are invariant quotient coordinates.  On the PQTI
branch Corollary~\ref{p3f:cor:pqti-same-action-invariance} fixes the latter.
\end{corollary}

\begin{proof}
Divide
\(\overline F_C=CZ_C/2\) from
Corollary~\ref{p3f:cor:term-ren-terminal-vacuum} by \(Z_C>0\), and use
equation~\eqref{p3f:eq:normalized-term-ren-data}.
\end{proof}

\begin{theorem}[Trace-linearized scalaron mass]
\label{p3f:thm:trace-linearized-scalaron-mass}
Suppose the four-dimensional trace equation
\[
 f'(R)R-2f(R)+3\Box f'(R)+4\Lambda_{\rm ren}
 =\kappa T,
 \qquad \kappa:=\frac{8\pi G}{c_{\rm op}^4},
\]
holds along a differentiable one-parameter family.  Let \(R_0\) and \(T_0\)
be constant background values, keep \(\Lambda_{\rm ren}\) fixed, and assume
\[
 R_0f'(R_0)-2f(R_0)+4\Lambda_{\rm ren}=\kappa T_0.
\]
Assume only that \(f\) is twice differentiable at \(R_0\), and write
\(\delta R\) and \(\delta T\) for the first variations.  Then the exact
linearized trace relation is
\begin{equation}
\label{p3f:eq:trace-linearized-scalaron}
 3f''(R_0)\Box_0\delta R
 +\bigl(R_0f''(R_0)-f'(R_0)\bigr)\delta R
 =\kappa\,\delta T.
\end{equation}
If \(f''(R_0)\ne0\), then with the paper's Lorentzian \(\Box\) convention
\[
 (\Box_0-m_s^2)\delta R
 =\frac{\kappa}{3f''(R_0)}\,\delta T,
\]
with the derived algebraic mass parameter
\[
 \boxed{m_s^2(R_0)
 =\frac{f'(R_0)-R_0f''(R_0)}{3f''(R_0)}.}
\]
The vacuum homogeneous equation is the specialization
\(T_0=\delta T=0\), not a premise of the sourced formula.
If \(f''(R_0)>0\), this parameter is positive exactly when
\(f'(R_0)-R_0f''(R_0)>0\).  The formula is therefore a consequence of the
trace equation on the stated constant background, not an additional MRESP
postulate.  If \(f''(R_0)=0\), division is forbidden and the first-order
relation is instead
\[
 -f'(R_0)\delta R=\kappa\,\delta T.
\]
For \(f'(R_0)\ne0\) this is algebraic and carries no scalar trace wave; if
also \(f'(R_0)=0\), first order requires \(\delta T=0\) but does not determine
\(\delta R\), so a higher-order analysis is necessary.  Calling \(m_s^2\) a
healthy physical mass additionally requires an actual full-equation
background and the appropriate nondegeneracy and stability signs.
\end{theorem}

\begin{proof}
At the background, \(R_0\) and hence \(f'(R_0)\) are constant.  Therefore
the variation of \(\Box\) acting on the background value is zero and
\[
 \delta\!\left(\Box f'(R)\right)
 =\Box_0\!\left(f''(R_0)\delta R\right)
 =f''(R_0)\Box_0\delta R.
\]
The algebraic variation is
\[
 \delta\!\left(f'(R)R-2f(R)\right)
 =\bigl(R_0f''(R_0)-f'(R_0)\bigr)\delta R.
\]
The cosmological constant is fixed, so its variation vanishes, while the
source varies by \(\kappa\delta T\).  Combining the two
displayed variations proves \eqref{p3f:eq:trace-linearized-scalaron}.
Division by \(3f''(R_0)\) proves the sourced Klein--Gordon form and mass
formula when that denominator is nonzero; setting it to zero gives the stated
algebraic and first-order-degenerate branches.
\end{proof}

\begin{corollary}[Derived smooth-MRESP scalaron endpoint criterion]
\label{p3f:cor:derived-minimal-response-scalaron}
Let \(f\) be any member of the smooth complete minimal-response family imported
from Paper~III-E, and assume \(Z_C=f'(C)>0\).  This corollary concerns that
smooth-MRESP family, not the v2.63 Laurent--Legendre branch whose lower active
Hessian is finite and constant up to the terminal handoff.  On \((C,R_+]\), define the
algebraic mass function by the boxed formula above.  Then
\[
 f''(R)>0,\qquad f''(R)\longrightarrow0^+,\qquad
 f'(R)-Rf''(R)\longrightarrow Z_C>0,
\]
and consequently
\[
 \boxed{m_s^2(R)\longrightarrow+\infty\qquad(R\downarrow C).}
\]
Whenever a constant-curvature vacuum background (or a family of such
backgrounds approaching \(C\)) satisfies the hypotheses of
Theorem~\ref{p3f:thm:trace-linearized-scalaron-mass}, this algebraic function
is the physical trace-linearized scalaron mass (or its limiting family).
The corollary does not assert existence of such a vacuum family; in
particular Paper~III-E's source endpoint is not itself a vacuum when
\(C>0\) and \(\Lambda_{\rm ren}=0\).
\end{corollary}

\begin{proof}
Paper~III-E's complete-family theorem gives the positivity and terminal
decay of \(f''\), while its endpoint reconstruction gives
\(f'(R)\to Z_C\).  Since \(R\to C<\infty\), one has
\(Rf''(R)\to0\).  The numerator therefore tends to \(Z_C>0\), and the
denominator tends to zero through positive values.  The physical-background
qualification is exactly the domain statement of the preceding theorem.
\end{proof}


\begin{theorem}[Conformal Hessian of the complete local curvature density]
\label{p3f:thm:v199-rgas-conformal-hessian}
Fix a four-dimensional Riemannian constant-scalar-curvature background
\(\bar g\), a \(C^2\) density \(F(R)\) on its variation interval, and the
declared action sign in \(\mathcal I[g]=\int\sqrt g\,F(R)\).
Assume a compact background without boundary or compactly supported
variations, so the integrations by parts below are valid.
In the quantum units used here the imported density is
\[
 F_k(R)=Z_{N,k}[f_k(R)-2\Lambda_{{\rm ren},k}],
 \qquad Z_{N,k}=(16\pi G_k)^{-1}.
\]
The cosmological volume term is included.  With \(P=-\bar\Box\)
and \(g=e^{2\sigma}\bar g\), define the Hessian by the quadratic term
\(\tfrac12\int\sqrt{\bar g}\,\sigma\mathcal H_\sigma\sigma\).  Then
\begin{equation}
\label{p3f:eq:v199-rgas-conformal-hessian}
 \mathcal H_\sigma
 =16F-12RF'+12F'P+4F''(R-3P)^2 .
\end{equation}
All density derivatives are evaluated at the background curvature,
at fixed \(k\).
\end{theorem}
\begin{proof}
Apply Paper~I-D, Theorem~\ref{p1b:thm:conformal-hessian}
\cite{Kugelmann2026Paper1D}, with its density \(\Phi=F_k\), at fixed
\(k\). Its Hessian convention, constant-curvature background and
integration-by-parts domain agree with those stated here. In particular,
\(F_k\) includes the cosmological volume term, so the imported formula
retains its contribution to \(16F-12RF'\).
\end{proof}

\begin{corollary}[Vacuum scalaron factor and dimensionless Hessian]
\label{p3f:cor:v199-rgas-scalaron-factor}
If the background is a vacuum of this complete local sector and obeys
\(RF'-2F=0\), with \(F''\ne0\), then
\begin{equation}
\label{p3f:eq:v199-rgas-scalaron-factor}
 \mathcal H_\sigma
 =36F''(P-R/3)(P+m_s^2),\qquad
 m_s^2=\frac{F'-RF''}{3F''}.
\end{equation}
For \(F_k(R)=k^4\varphi_k(r)\), \(r=R/k^2\), and \(P=k^2\zeta\),
\[
 k^{-4}\mathcal H_\sigma
 =16\varphi-12r\varphi'+12\zeta\varphi'
                         +4\varphi''(r-3\zeta)^2.
\]
\end{corollary}
\begin{proof}
The vacuum factorization and dimensionless substitution are the final
parts of Paper~I-D, Theorem~\ref{p1b:thm:conformal-hessian}
\cite{Kugelmann2026Paper1D}, with \(\Phi=F_k\) and \(\mu^2=m_s^2\).
The factorization uses the stated complete-sector equation \(RF'=2F\).
\end{proof}
For \(F=Z_N(f-2\Lambda_{\rm ren})\), the vacuum condition is
\(Rf'-2f+4\Lambda_{\rm ren}=0\).  The mass numerator agrees with the
trace-linearized result, while the background depends on the complete
volume term.  This is the conformal restriction before gauge constraints,
ghosts, Jacobians and other fluctuation blocks; its determinant alone is
not the full gravity trace.  On an Einstein background a transverse
traceless variation has \(\delta R=0\), so \(F''(\delta R)^2\) supplies
no fourth-order TT principal part.  The scalar fourth-order factor does
not determine the tensor anomalous dimension.


\begin{definition}[Imported nonlinear response realization atlas]
\label{p3f:def:nonlinear-response-realization-atlas}
Paper~III-E separates the nonlinear response data into a universally
quantified REG bundle, an exact raw affine endpoint pair, a generated
response-owner bridge, a generated PQTI owner, and
generated/co-realized constituents.  On the enlarged atlas before the
same-action owner theorem is imposed, its unrestricted
MRESP classification has one dimensionless ultraviolet coordinate.  Import
\[
 \boxed{\mathfrak A_{\rm REG}
 =\{\mathrm{REG}_{\sigma}:\sigma_{\rm fin}\in\mathbb R\}},
 \qquad
 \boxed{\mathfrak A_{\rm TERM}
 =\{(F_C,Z_C):F_C\in\mathcal L_{\scal},\ Z_C\in\mathbb R\}},
\]
and that unrestricted minimal susceptibility atlas
\[
\boxed{\mathfrak A_{\rm MRESP}^{\min}
 =\{\gamma_\sigma\in\mathbb R_{>0}\}},\qquad
\chi_{\rm UV}=\frac{\gamma_\sigma}{R_+}.
\]
The symbol \(R_+\) is III-F's name for Paper~III-E's upper curvature root
\(U\); this prevents collision with the spacetime patch \(U\).
Here \(F_C\) is curvature-valued and \(Z_C\) is dimensionless.  Numerical
coordinates in \(\mathbb R^2\) arise only after a curvature unit is fixed.
The source endpoint
\(\mathrm{TERM}_f^0:(F_C,Z_C)=(0,1)\) remains a declared member.  The
octahedral compiler derives
\(\mathrm{REG}_0:\sigma_{\rm fin}=0\) on the generated \(2O\) interface;
unanchored branches retain the affine family.  The visibility process, closed trace, normalized state,
entropy, and determinant-depth coordinate co-realize from the active
composite.  Once the minimal MRESP composite also contains the CRP4 response
ancestry and a frozen
\(\mathrm{REG}_{\sigma}\in\mathfrak A_{\rm REG}^{+}\), owner inheritance
additionally gives the scale interpolation and hierarchy-active state:
\[
 \boxed{\operatorname{CoReal}
 (V_\tau,Z(\tau),\rho_\tau,S(\rho_\tau),q_{\rm sp}(\tau),
   \scal(\tau),\sigma_{\rm act}(\scal)\mid\mathrm{MRESP}_{\min}).}
\]
The ancestry port is generated before numerical evaluation by the
scalar-action coefficient compiler and is not a competing response member.
Within \((\mathrm{MRESP}_{\min})\), Paper~III-E derives the Hessian transport
as ordinary covariant pullback and derives the active-state coupling as the
unique continuous entropy-sufficient, tensor-multiplicative,
octet-normalized effective count.  Neither is a selector.  On the PQTI
subfamily, the primitive response/associated-line morphism and the common
octet trace give
\[
 \boxed{\widehat\gamma_\sigma
 =16\sqrt{\alpha_{\rm hier}(\sigma_{\rm fin})}}.
\]
Thus \(\gamma_\sigma\) is no longer an independent physical coordinate on
that branch, and \(R_+=\Lambda_{\rm EH}^2\).  Dropping the minimal tensor or
entropy signature reopens an enlarged
\(\mathfrak T_+\times\mathfrak K_+\) role not used below.
The III-D/III-E two-root response also adopts the convenient
\(\kappa=2\) response-normalization representative; other
\(g_\kappa=\kappa\lambda_B(W_{\rm scal})\) representatives are absorbed into the
conversion scale and give the same calibrated
\[
 \boxed{g_{\rm eff}(\scal)=\frac{\scal^2+CR_+}{(C+R_+)\scal}}.
\]
Thus \(\scal_{\rm resp}=C+R_+\) is a normalization representative, not a
further physical selector.
\end{definition}

\begin{theorem}[MRESP specialization of the local metric equation]
\label{p3:thm:local-fR-field-equation}
Assume
\[
\begin{gathered}
 \mathsf C_{\rm IR}^{\rm der}+(\mathrm{GRAV\!-\!ASYM})
 +(\mathrm{VDOM}_{x}^{\min})+(\mathrm{JOINT}_x)\\
 +(\mathrm{SOLHIST}_{U}^{\rm phys})
 +\bigl(\mathrm{REG}_{\sigma}\in\mathfrak A_{\rm REG}^{+}\bigr)\\
 +(\mathrm{CRP4})+(\mathrm{PQTI})+(\mathrm{TERM}_f)
 +(\mathrm{MRESP}_{\min})+(\mathrm{REN}_0(\Lambda_{\rm ren})),
\end{gathered}
\]
with the protected Einstein branch inherited from Paper~III-E.  Assume one of
the two explicit action/normalization routes:
\[
\begin{gathered}
 (\mathrm{ACT}_{\rm MRESP}):\\[2pt]
 \begin{aligned}
 (\mathrm{LOR}_{\rm dir}):\quad&
 \text{the normalized Lorentzian response action with physical \(G\)}\\
 &\text{occurs directly};\\[2pt]
 (\mathrm{CONT}_{\rm cl}):\quad&
 \text{the spectral response is continued with fixed field/sign role},\\
 &\text{and its EH coefficient is matched with the III-E primitive}\\
 &\text{curvature actualization}.
 \end{aligned}
\end{gathered}
\]
Let \(U_0\Subset U\) be the interior
unit-localizer domain and restrict metric and matter variations to compact
support in \(U_0\).  Impose the premise-minimal classical MRESP subdomain:
for every admitted variation with support \(K\Subset U_0\), the action (or
compactly supported relative action difference) is meaningful and there are
a neighborhood \(V\supset K\) and an open curvature interval \(I\) such that
the nearby metrics remain nondegenerate,
\[
 (\mathrm{MRESP}_{\rm cl}):\qquad
 \scal[g_\varepsilon](V)\subset I,\qquad
 \mathfrak f_{\rm sp}\in C^1(I),\qquad
 \mathfrak f_{{\rm sp},\scal}(\scal[g])\in C^2(V).
\]
Let \(\mathfrak f_{\rm sp}(\scal)\) be any metric response of
Theorem~\ref{p3e:thm:complete-minimal-smooth-response-family}; the historical
source response is its \((F_C,Z_C)=(0,1)\) specialization.
On the direct branch the following dimensional action is the declared
normalized occurrence.  On a spectral-derived branch it additionally
requires the actual physical action/source projection of
Proposition~\ref{p3f:prop:v2211-newton-owner-matching}, together with the
certified response ports of
Definition~\ref{p3f:def:v2211-certified-response-domain}.  The raw joint
coordinate is \(G_{J,\mathrm{raw}}=c_{\rm op}^3/(\hBorn U_J^{\rm raw})\),
where \(U_J^{\rm raw}=96U\).  The primitive geometric scale
\(\mathcal K_p=6/\Rmom^2\), \(U=\Lambda_{\rm EH}^2\); the further
comparison \(U=(3/\pi)\mathcal K_p\) holds only under the Paper~III-E
\((\mathrm{CELL}_{K\to U})\) interface, as restated below.
Neither geometric comparison itself evaluates that action projection.  The scalar-response
coordinate is separately \(G_{\rm scal}=c_{\rm op}^3/(\hBorn U)\).
On this direct or projection-certified domain the action is
\[
 S_{\rm eff}[g,\psi]
 =\frac{c_{\rm op}^3}{16\pi G_{\rm scal}}
   \int_U d^4x\sqrt{-g}\,
   [\mathfrak f_{\rm sp}(\scal)-2\Lambda_{\rm ren}]
 +S^{\rm norm}_{\rm m,U}[g,\psi].
\]
A numerical comparison \(G_{\rm scal}=G_N^{\rm obs}\) is downstream
metrology only after this physical normalization and the source map have
been established.  A raw coefficient ratio does not establish them.
Here \(S^{\rm norm}_{\rm m,U}\) is the half-weighted position-role matter
functional inherited from
Theorem~\ref{p3:thm:common-position-spectral-origin}; the imported smooth
curvature response belongs to the full Born-diagonal pure-metric channel and
is not a raw unweighted position trace.  Define
\[
 T_{\mu\nu}:=-\frac{2c_{\rm op}}{\sqrt{-g}}
 \frac{\delta S^{\rm norm}_{\rm m,U}}{\delta g^{\mu\nu}}.
\]
The imported response is the integrable Hessian reconstruction
\[
 \mathfrak f_{{\rm sp},\scal}(\scal)
 =Z_C+\int_C^{\scal}h(s)\,ds,
 \qquad
 \mathfrak f_{{\rm sp},\scal\scal}=h,
\]
with \((F_C,Z_C)\) independently specified endpoint data.  No independent
first-variation semigroup transport is assumed; the coefficient response is
the generated restriction of \(dS_\star\).  On the stronger sufficient
subdomain \(h\in C^1(I)\), one has
\(\mathfrak f_{\rm sp}\in C^3(I)\); together with
\(g\in C^4_{\rm loc}\) and \(\scal[g]\in C^2_{\rm loc}\), this is a
convenient sufficient route to \((\mathrm{MRESP}_{\rm cl})\), not an
irreducible premise.  The direct composite condition displayed above is what
the field-equation proof uses.
Writing \(\mathfrak f_{{\rm sp},\scal}=d\mathfrak f_{\rm sp}/d\scal\), metric
variation gives on \(U_0\)
\begin{equation}
\label{p3:eq:master-fR-field-equation}
\boxed{
\mathfrak f_{{\rm sp},\scal}R_{\mu\nu}
-\frac12\mathfrak f_{\rm sp}g_{\mu\nu}
 +\bigl(g_{\mu\nu}\Box-\nabla_\mu\nabla_\nu\bigr)
 \mathfrak f_{{\rm sp},\scal}
 +\Lambda_{\rm ren}g_{\mu\nu}
 =\frac{8\pi G}{c_{\rm op}^4}T_{\mu\nu}.}
\end{equation}
Its trace is
\begin{equation}
\label{p3:eq:master-fR-trace}
 \mathfrak f_{{\rm sp},\scal}\scal-2\mathfrak f_{\rm sp}
 +3\Box\mathfrak f_{{\rm sp},\scal}
 +4\Lambda_{\rm ren}
 =\frac{8\pi G}{c_{\rm op}^4}T.
\end{equation}
On the generated-owner Paper~III-E branch
\[
 R_+=U,\qquad
 C=U\alpha_{\rm hier}(\sigma_{\rm fin})^{-7},
\]
and the imported nonlinear Hessian is
\[
 \mathfrak f_{{\rm sp},\scal\scal}(\scal)
 =w_{\rm nl}(\scal)u(\scal)
  \frac{16\sqrt{\alpha_\sigma}}{R_+}
 =w_{\rm nl}(\scal)u(\scal)
  \frac{16\sqrt{\alpha_\sigma}}{U},
 \qquad
 w_{\rm nl}(\scal)=\frac18\,8^{n_{\rm act}(\scal)}.
\]
Here \(n_{\rm act}\) is the entropy index of the co-realized active state;
it is not Paper~III-D's geometric secant exponent \(n_{\rm sec}\).  The
state, entropy, and determinant-depth coordinate are generated/co-realized
constituents; the scale interpolation is their owner-typed inheritance and no
active-state selector remains.  Their
functional coupling to the Hessian and its tensor transport are derived in
the minimal signature.  PQTI fixes the relative ultraviolet normalization,
and the action compiler gives \(R_+=U=\Lambda_{\rm EH}^2\); Paper~III-E supplies the conditional comparison \(\mathcal K_p=6/\Rmom^2,\qquad U=(3/\pi)\mathcal K_p\quad\text{only under }(\mathrm{CELL}_{K\to U})\).  The convexity and endpoint
properties are proved in Paper~III-E.  These are imports of the frozen
\(\mathrm{REG}_{\sigma}\in\mathfrak A_{\rm REG}^{+}\),
\((\mathrm{CRP4})+(\mathrm{PQTI})+(\mathrm{TERM}_f)
+(\mathrm{MRESP}_{\min})\) branch, not
consequences of \(\Dcl\) or \((\mathrm{CSA})\).
The seven integer scalar hierarchy levels are generated ordered scale
surfaces.  They are not rank-one projector transitions without
\((\mathrm{FLAG})\), and they acquire named cosmological meanings only
through downstream solution predicates.
\end{theorem}

\begin{proof}
The selected branch of \((\mathrm{ACT}_{\rm MRESP})\) supplies the normalized
Lorentzian action, and \((\mathrm{MRESP}_{\rm cl})\) supplies exactly the
support-local regularity of the general theorem.
Apply Theorem~\ref{p3f:thm:general-local-fR-variation} to
\(f=\mathfrak f_{\rm sp}\), with the common matter projection varied in the
Hilbert sense using \eqref{p3f:eq:matter-tensor-convention}.  Compact support
in \(U_0\), where \(\varphi_U=1\), removes
localizer-derivative terms.  The Hessian and calibrated endpoint statements
are direct imports of the \((\mathrm{MRESP}_{\min})\) branch and certify why
this response occurs; once the function and action are given, they are not
mathematical inputs to the variation formula.
\end{proof}

\begin{corollary}[Geometric gauge--gravity split of the reciprocal role]
\label{p3:cor:born-carrier-gauge-gravity-split}
The Born-equivariant ratio submersion of
Theorems~\ref{p3b:thm:relative-readout-admission} and
\ref{p3:thm:born-parity-decomposition} retain the smooth reciprocal product
and present its relative readout together with a diagonal orbit fibre.  On the vector
part, the diagonal fibre is one round three-sphere and has connected spin
isometry
\[
 \Spin(4)_{\rm fib}\cong SU(2)^a_L\times SU(2)^a_R.
\]
Hence the native left--right geometric gauge slot is the isometry algebra of
one round diagonal vector fibre; it is neither the Born-odd quotient tangent
module treated as a Lie algebra nor a duplication of the original momentum
carrier.  The scalar pair supplies only a dimensionless diagonal
\(U(1)_{\rm fib}\) representation slot at this constitutional stage.  A
physical scalar-circle radius or absolute \(U(1)_{\rm fib}\) kinetic
coefficient requires an additional metric calibration and is therefore
supplementary.
\end{corollary}

\begin{proof}
The ratio submersion gives a diagonal \(SU(2)\) fibre over the relative
quotient.  Its round bi-invariant metric has spin-isometry
\(SU(2)_L\times SU(2)_R\).  The scalar role supplies a diagonal circle at the
representation level, but current Paper~III-B retains no physical scalar-circle
metric modulus in the constitutional core.  The statement therefore stops at
the dimensionless \(U(1)\) slot.
\end{proof}

\subsection{Born-spherical gauge--gravity reduction}
\label{p3:sec:born-spherical-spectral-action-reduction}

Paper~III-B's canonical scalar-diagonal constructor supplies a generated
compact seven-dimensional principal-\(SU(2)\) readout role before this
section.  The quotient by itself does not choose a connection, but III-B's
parent-Cartan descent theorem now supplies one unique compact residual
connection from already generated data.  Composed with Paper~III-A's already
generated position chart, that connection has a canonical local pullback.
Paper~II-A supplies hereditary Born covariance for this single-valued unmarked
descendant, while Paper~I-C's homogeneous family is retained only as an
independent uniqueness/classification certificate.
The two theorems below discharge precisely that geometric carrier and the
local pullback of the already generated background connection.  They do not actualize fluctuated
\(SU(2)_L\times SU(2)_R\) gauge potentials, adopt a physical action, choose a
fibre/action normalization, or perform the scalar-curvature reduction.  The
main calculation below remains the conditional seven-dimensional connection
reduction on that reduced interface.  The optional eight-dimensional Born-even comparison
\((\mathrm{BULK}_8)\) is a separate diagnostic and is not silently identified
with that connection space.  The four-dimensional gravitational coefficient
is not rederived here: it is the already normalized protected
Einstein coefficient imported through the
\(\mathsf C_{\rm IR}^{\rm der}
 +(\mathrm{GRAV\!-\!ASYM})\) generated-heat branch.  On that branch,
\((\mathrm{GRAV}_{\circ})\), \((\mathrm{CORN})\),
and \((\mathrm{BD}_{\rm EH})\) are derived process/constituent-identification
certificates, not additional occurrence assumptions.  The vector-fibre
geometry fixes the left--right equality; the absolute seven-dimensional
coefficient additionally requires its connection reduction, compact-cell
saturation when the dimensionless value is quoted, and generator
normalization.  The eight-dimensional factor enters only if
\((\mathrm{MATCH}_{8\to7})\) is separately supplied.

\begin{theorem}[Canonical local pullback of the compact GRED carrier]
\label{p3f:thm:canonical-local-gred-pullback}
Let \(U\subset K_x=S_B^3\times S_t^1\) be any product chart on which
Paper~III-A's generated reciprocal localization is realized, let
\(\iota_U:U\hookrightarrow K_x\) be inclusion, and let
\[
 \phi_U:=j_x\circ\iota_U:U\longrightarrow K_{\rm rel},
\]
where \(j_x\) is Paper~III-B's canonical factor restriction of the ratio map.
Then
\begin{equation}
\label{p3f:eq:canonical-local-gred-pullback}
 \boxed{
 X_{7,U}^{\rm can}
 :=\phi_U^*X_7^{\rm cmp}
 =\{(y,\xi)\in U\times X_7^{\rm cmp}:
      \phi_U(y)=\pi_{\rm rel}(\xi)\}}
\end{equation}
is a canonically generated principal \(SU(2)_{\rm diag}\)-bundle over \(U\).
It has dimension seven, is compatible with restriction and chart overlap, and
is, through the diffeomorphism \(j_x\), exactly the restriction of
\(X_7^{\rm cmp}\) to \(j_x(U)\).  Thus no arbitrary four-dimensional base,
bundle topology, or gluing map remains in the local GRED carrier.
\end{theorem}

\begin{proof}
Paper~III-B proves that \(j_x:K_x\to K_{\rm rel}\) is a canonical
diffeomorphism and that
\(X_7^{\rm cmp}\to K_{\rm rel}\) is a generated principal
\(SU(2)_{\rm diag}\)-bundle.  Pullback along \(\phi_U\) therefore gives a
principal bundle over \(U\), with the displayed fibre product and inherited
right action.  Since \(j_x\) is a diffeomorphism, the map
\[
 (y,\xi)\longmapsto \xi
\]
identifies the pullback with
\(\pi_{\rm rel}^{-1}(j_x(U))\), while its inverse sends
\(\xi\) to \(((j_x|_U)^{-1}\pi_{\rm rel}(\xi),\xi)\).  Pullback functoriality
gives restriction and overlap compatibility, and
\(\dim U+\dim SU(2)=4+3=7\).

The construction is the exact pullback of an already generated bundle along
an already generated chart map.  It is natural, choice-free, and its graph is
its full essential image.  Paper~II-A,
Theorem~\ref{p2a:thm:geometric-defect-repair-admission}, therefore admits it
by exact transport; merely permitting an unrelated local bundle would not
pass that criterion.
\end{proof}

\begin{theorem}[Parent-descended Born connection and homogeneous uniqueness certificate]
\label{p3f:thm:unique-born-invariant-gred-connection}
Identify
\[
 X_7^{\rm cmp}
 \cong G:=SU(2)_x\times SU(2)_p\times U(1)_{\rm rel},
 \qquad
 H:=SU(2)_{\rm diag},
\]
with the principal right action
\((x,p,z)\cdot h=(xh,ph,z)\).  Paper~III-B,
Theorem~\ref{p3b:thm:parent-cartan-scalar-diagonal-descent}, already generates
on this compact principal bundle the unique scalar-diagonal descent of the
parent \(M_8\) Cartan connection:
\begin{equation}
\label{p3f:eq:unique-born-invariant-connection}
 \boxed{\omega_{\rm B}:=\omega_{\rm cmp}
 =\frac12(\theta_x+\theta_p).}
\end{equation}
It is fixed by the descended Born exchange.  Its horizontal curvature is
\begin{equation}
\label{p3f:eq:born-invariant-connection-curvature}
 \boxed{
 \Omega_{\rm B}\bigl((X,-X,0),(Y,-Y,0)\bigr)=-[X,Y].}
\end{equation}
For the local section \(s(r,z)=(r,e,z)\), and with the matrix-valued
convention \(F=dA+A\wedge A\),
\[
 A_{\rm B}=s^*\omega_{\rm B}=\frac12\theta_r,
 \qquad
 F_{\rm B}=-\frac14\theta_r\wedge\theta_r.
\]
Pullback to \(X_{7,U}^{\rm can}\) gives the same distinguished nonflat
background connection on every generated chart.

Independently, Paper~I-C,
Theorem~\ref{p1b:thm:homogeneous}, classifies all left-\(G\)-invariant
connections natural under the unmarked diagonal \(H\)-action as
\[
 \omega_f=f\theta_x+(1-f)\theta_p.
\]
Within that mathematical family, Born fixedness forces \(f=1/2\).  Therefore
the already generated parent-descended connection is also the unique
homogeneous Born-fixed member.  The classification is a target-identification
and uniqueness certificate; it is not the physical selector that generates
\(\omega_{\rm B}\).

This theorem fixes the compact principal background and its generated chart
pullbacks only.  The next theorem functorially extends that background to the
generated \(\Spin(4)_{\rm fib}\) factors and one joint Dirac family.  Arbitrary
independent left/right fluctuation members, a physical action, and all
absolute generator or fibre normalizations remain outside the present
conclusion.
\end{theorem}

\begin{proof}
Paper~III-B,
Theorem~\ref{p3b:thm:parent-cartan-scalar-diagonal-descent}, proves that
\(\pi_{\rm su}(\mathcal A_{M8})\) is basic for the scalar-diagonal quotient and
therefore descends uniquely to \(\omega_{\rm cmp}\).  It also proves the
explicit formula, Born fixedness, curvature, and section representative above.
The generated chart identification with \(X_7^{\rm cmp}\) then pulls the same
connection back to \(X_{7,U}^{\rm can}\); functoriality of connection pullback
gives the local background.

For the independent uniqueness check, Paper~I-C,
Theorem~\ref{p1b:thm:homogeneous} \cite{Kugelmann2026Paper1C}, gives the family
\(f\theta_x+(1-f)\theta_p\).  Born exchange sends \(f\mapsto1-f\), so the fixed
member has \(f=1/2\), exactly the upstream-generated \(\omega_{\rm cmp}\).
Thus no unresolved family member is selected by reciprocity alone: the
single-valued physical connection is already generated upstream, and the
homogeneous calculation certifies its unique position in the larger
mathematical family.
\end{proof}

\begin{remark}[Born inheritance and the homogeneous comparison family]
\label{p3f:rem:born-inheritance-connection-family-firewall}
Paper~II-A, Theorem~\ref{p2a:thm:generated-born-inheritance}, still does not
select a member of an unresolved connection family.  The strengthened ancestry
above avoids that forbidden inference: Paper~III-B first generates one
single-valued connection by basic descent of the parent Cartan connection, and
Born covariance of that generated output is inherited/verified on its
Born-unmarked owner.  The left-invariant \(\omega_f\) family remains a
mathematical classification space.  Paper~I-C then shows that the generated
connection is exactly its unique Born-fixed midpoint.  Hence reciprocity is
hereditary on the generated connection but is not being used as a free family
selector.
\end{remark}


\begin{corollary}[Full restricted holonomy of the canonical Born connection]
\label{p3f:cor:v181-full-born-holonomy}
For the connection of
Theorem~\ref{p3f:thm:unique-born-invariant-gred-connection},
\[
 \boxed{\operatorname{Hol}_0(\omega_{\rm B})=SU(2)_{\rm diag}.}
\]
Hence a single Hopf $U(1)$ may be used as a readout of a chosen Cartan, but it
is not a connection-faithful replacement for the generated principal owner.
\end{corollary}
\begin{proof}
This is the restricted-holonomy conclusion of Paper~I-C,
Theorem~\ref{p1b:thm:homogeneous} \cite{Kugelmann2026Paper1C}, for the
midpoint connection identified in
Theorem~\ref{p3f:thm:unique-born-invariant-gred-connection}.
A connection-faithful reduction to a Hopf circle would have to contain
this holonomy group, which is the full \(SU(2)_{\rm diag}\).
\end{proof}

\begin{corollary}[Protected-projector ancestry of the midpoint connection]
\label{p3f:cor:protected-projector-midpoint-ancestry}
Normalize Paper~I-A's protected coordinates by
$\Gamma_a=2\widehat X_a/\rho_{\rm prot}$, so that
$\{\Gamma_a,\Gamma_b\}=2\delta_{ab}I$ as in
Theorem~\ref{p1:thm:s3one-cliffordization}.
On their coefficient sphere $n\in S^3$, choose the chiral frame
\[
 \Gamma(n)=\begin{pmatrix}0&U(n)\\U(n)^\dagger&0\end{pmatrix},
 \qquad U(n)=n_4I+i n_i\sigma_i\in SU(2).
\]
Both rank-two spectral projectors
$P_\pm=(I\pm\Gamma(n))/2$ have the midpoint projected connection
\[
 V_\pm=\frac{1}{\sqrt2}\begin{pmatrix}\pm U\\I\end{pmatrix},
 \qquad A_\pm=V_\pm^\dagger\mathrm dV_\pm
             =\frac12U^\dagger\mathrm dU,
 \qquad F_\pm=-\frac14(U^\dagger\mathrm dU)\wedge(U^\dagger\mathrm dU).
\]
Thus their connections are equivalent, as in
Theorem~\ref{p2c:thm:protected-projector-midpoint}, and identify explicitly
with the relative $SU(2)$ midpoint of Paper~I-C,
Theorem~\ref{p1b:thm:homogeneous}.  Their restricted holonomy is $SU(2)$.
An oriented Hopf axis can enlarge a marked pointwise commutant without
reducing this connection's holonomy to $U(1)$.
\end{corollary}
\begin{proof}
The displayed frames satisfy $V_\pm^\dagger V_\pm=I$ and
$\Gamma(n)V_\pm=\pm V_\pm$.  Direct differentiation gives the same
$A$ in both frames.  The Maurer--Cartan identity yields the curvature,
whose values span $\mathfrak{su}(2)$; the holonomy claim follows as in
Corollary~\ref{p3f:cor:v181-full-born-holonomy}.
\end{proof}
This is an exact geometric comparison between existing protected and
relative-quotient owners.  Its use as a physical gauge or adiabatic
connection requires the actual action or gapped Hamiltonian and pullback.
The complementary coordinate-sign projectors alone do not produce a
helicity imbalance.  Their common connection does not select a populated
history or replace a physical source normalization.


\begin{theorem}[Canonical fibre-spin extensions and generated joint operator]
\label{p3f:thm:generated-fibre-spin-joint-operator}
\label{p3f:thm:joint-connection-occurrence-discharge}
Let
\[
 P_{\rm B}:=X_{7,U}^{\rm can}\longrightarrow U,
 \qquad H:=SU(2)_{\rm diag},
\]
carry the nonflat connection \(\omega_{\rm B}\) of
Theorem~\ref{p3f:thm:unique-born-invariant-gred-connection}.  Identify its
round fibre with the unit quaternions.  Left and inverse-right multiplication,
\[
 \lambda_h(q)=hq,
 \qquad r_h(q)=qh^{-1},
\]
are orientation-preserving isometries and give the two generated
homomorphisms
\[
 \iota_L:H\longrightarrow SU(2)^a_L,
 \qquad
 \iota_R:H\longrightarrow SU(2)^a_R
\]
inside
\(\Spin(4)_{\rm fib}=SU(2)^a_L\times SU(2)^a_R\).  Both extensions of
structure group therefore occur as rooted descendants,
\begin{equation}
\label{p3f:eq:canonical-fibre-spin-extensions}
 P_\epsilon:=P_{\rm B}\times_{\iota_\epsilon}SU(2)^a_\epsilon,
 \qquad
 \omega_\epsilon:=\operatorname{Ext}_{\iota_\epsilon}(\omega_{\rm B}),
 \qquad \epsilon\in\{L,R\}.
\end{equation}
They satisfy
\begin{equation}
\label{p3f:eq:canonical-fibre-spin-curvature}
 \Omega_\epsilon=d\iota_\epsilon(\Omega_{\rm B})\ne0.
\end{equation}
Their fibre product
\[
 P_{\Spin}:=P_L\times_U P_R\longrightarrow U,
 \qquad \omega_{\Spin}:=(\omega_L,\omega_R),
\]
is consequently a generated principal \(\Spin(4)_{\rm fib}\)-bundle with a
nonflat connection.  The two components are co-realized images of the same
\(\omega_{\rm B}\); the conclusion does not identify them as independently
selected fluctuations.

Let \(W=W_L\oplus W_R\), with the two fundamental half-spin actions supplied
by Theorem~\ref{p3b:thm:diagonal-fibre-spin-ancestry}.  Then
\begin{equation}
\label{p3f:eq:faithful-fibre-spin-representation}
 \rho_W:\Spin(4)_{\rm fib}\longrightarrow U(W)
 \quad\hbox{is faithful}.
\end{equation}
Under the Paper~II-D left-Hopf factor dictionary, these symbols
\(SU(2)^a_L,SU(2)^a_R\) in the multiplication-side definition refer to
\emph{geometric} left and inverse-right actions. Their algebraic
half-spin labels are interchanged:
\begin{equation}\label{p3f:eq:v270-geometric-algebraic-side-dictionary}
\boxed{\rho_W(u_{\rm geom,L},v_{\rm geom,R})
=\rho_{1/2}(v_{\rm geom,R})\big|_{W_L}
\oplus\rho_{1/2}(u_{\rm geom,L})\big|_{W_R}.}
\end{equation}
Thus the geometrically right extension is the algebraic-left
weak extension, while the geometrically left factor contains the
oriented source-Hopf circle. This is a representation dictionary,
not a new connection or a collapse of the full co-realized product.

For every already generated nonzero finite spectator carrier \(K\), the
identity extension \(F:=W\otimes K\) remains faithful, and the associated
connection \(\nabla_{\rm B}^F\) has nonzero represented curvature on both
half-spin blocks.  On a generated local spin chart, put
\begin{equation}
\label{p3f:eq:canonical-joint-dirac-anchor}
 D_{{\rm B},F}
 :=c_g\!\circ\!\left(
       \nabla^{\mathbb S_g}\otimes1_F
       +1_{\mathbb S_g}\otimes\nabla_{\rm B}^F
     \right),
 \qquad
 \sigma_1(D_{{\rm B},F})(\xi)=c_g(\xi)\otimes1_F .
\end{equation}
This is a nontrivial covariantly represented joint datum.

More generally, if \(\mathfrak a\subseteq\mathfrak u(F)\) is any later
generated gauge-operation algebra and \(\mathcal V_{\rm odd}\subseteq
\End(F)\) any later generated self-adjoint odd-field family, their complete
unmarked joint operator family is
\begin{equation}
\label{p3f:eq:generated-joint-dirac-family}
 \mathfrak D_J(U;F)
 :=\left\{
 D_{{\rm B},F}+c_g(a)+\gamma_g\otimes\Phi:
 \begin{array}{l}
 a\in\Omega_c^1(U,\mathfrak a),\\[-1mm]
 \Phi\in\Gamma_c(U,\mathcal V_{\rm odd})
 \end{array}
 \right\},
\end{equation}
with the conventional factor of \(i\) absorbed into the skew-Hermitian
generator convention.  Equation~\eqref{p3f:eq:canonical-joint-dirac-anchor}
is the member \((a,\Phi)=(0,0)\), but it is not the zero-connection member.
Thus the zero-member test is passed by the required nonzero output itself,
and
\begin{equation}
\label{p3f:eq:joint-x-derived}
 \boxed{(\mathrm{JOINT}_x^{\rm der}):=
 \operatorname{CoReal}\bigl(P_{\Spin},\omega_{\Spin},
 \nabla_{\rm B}^F,\mathfrak D_J(U;F)\bigr).}
\end{equation}
No independent \((\mathrm{JOINT}_x^{\rm occ})\) premise remains.  The
theorem co-realizes the complete family and its canonical nonflat anchor; it
does not select \(a\), \(\Phi\), a physical action, a stationary solution, or
an absolute normalization.
\end{theorem}

\begin{proof}
For either \(\epsilon\), realize the extension in
\eqref{p3f:eq:canonical-fibre-spin-extensions} as the quotient of
\(P_{\rm B}\times SU(2)^a_\epsilon\) by
\((p,g)h=(ph,\iota_\epsilon(h)^{-1}g)\).  On the product, the standard form
\begin{equation}
\label{p3f:eq:extension-connection-form}
 \widetilde\omega_\epsilon
 =\operatorname{Ad}_{g^{-1}}
      d\iota_\epsilon(\omega_{\rm B})+g^{-1}dg
\end{equation}
is basic for this action in the connection sense and descends uniquely to
\(P_\epsilon\).  Its curvature is
\(\operatorname{Ad}_{g^{-1}}d\iota_\epsilon(\Omega_{\rm B})\).  Both
\(d\iota_\epsilon\) are injective, so
Theorem~\ref{p3f:thm:unique-born-invariant-gred-connection} and
\eqref{p3f:eq:born-invariant-connection-curvature} prove
\eqref{p3f:eq:canonical-fibre-spin-curvature}.  Product of the two extensions
gives \(P_{\Spin}\) and the displayed product connection.

By the geometric/algebraic dictionary
\eqref{p3f:eq:v270-geometric-algebraic-side-dictionary}, if
\((u_L,u_R)\) acts trivially on \(W_L\oplus W_R\), its restriction to
\(W_L\) gives \(u_R=1\), and its restriction to \(W_R\) gives \(u_L=1\).
Thus \(\rho_W\) is faithful.  GP3 cannot collapse the product to the
diagonal copy of \(H\): the two factor projections and the two marked
half-spin actions are distinct generated operations, and a diagonal collapse
would delete the independent \(SU(2)^a_L\) and \(SU(2)^a_R\) connection
directions required by GP2 no-omission.  Their canonical curvatures are
correlated images of one source, not duplicate representation roles.
Tensoring with the identity on a nonzero
carrier preserves its kernel and preserves every nonzero curvature block.
Functorial associated-bundle transport gives \(\nabla_{\rm B}^F\).
Clifford contraction with the generated local cotangent role gives
\eqref{p3f:eq:canonical-joint-dirac-anchor}; a connection term is zeroth order,
so the principal symbol is unchanged.  In local coordinates it is a finite
sum of terms
\[
 A_\mu^a\rho_F(T_a)[D_g,x^\mu]
 =A_\mu^a\rho_F(T_a)c_g(dx^\mu),
\]
which lies in the covariant one-form module of
Theorem~\ref{p3f:thm:joint-same-datum-relation-closure}.  Faithfulness and
nonzero curvature exclude collapse to the trivial represented connection.

Finally, structure-group extension, associated-bundle transport, direct sum,
identity tensor extension, local cotangent multiplication, Clifford
contraction, graded addition, and formation of the complete unmarked affine
family are precisely the already licensed operations used by the frozen
principal-bundle and local-operator grammar.  Their target sorts are declared
in equations~\eqref{p3f:eq:canonical-fibre-spin-extensions}--
\eqref{p3f:eq:generated-joint-dirac-family} before the outputs are evaluated.
They therefore form a frozen constructor-origin word in the sense of
Definition~\ref{p2a:def:frozen-constructor-origin-certificate}.  The inputs
\(P_{\rm B},\omega_{\rm B}\), the local cotangent/Clifford role, and \(W\)
are generated; later spectators and field families enter only when they have
themselves been generated.  The rooted-image theorem
\ref{p2a:thm:rooted-generated-image-corealization} and GP2 no-idle closure
co-realize every required output of the word.  Because the required anchor
has curvature \(d\rho_F(\Omega_{\Spin})\ne0\), the zero-member corollary
\ref{p2a:cor:zero-member-occurrence-test} cannot preserve the frozen word by
deleting it.  This proves \eqref{p3f:eq:joint-x-derived} and all stated
boundaries.
\end{proof}

\subsection{Evaluated coupled native Einstein--scalar vacua}
\label{p3f:sec:v248-native-constant-vacua}

The isolated quartic minima do not solve the metric equation. The
following calculation solves both native equations on a declared
constant background and retains the independently owned absolute
vacuum coefficient. It does not select or subtract that coefficient.

\begin{theorem}[Complete constant Einstein--scalar equations with the owned vacuum coefficient]
\label{p3f:thm:v248-native-constant-vacua}
On the actual local action domain of
Definition~\ref{p3f:def:v176-local-bosonic-noether-domain}, retain a
fixed action-owned constant density $V_0$ and write
\[
 \mathscr F=\kappa_J-\frac\alpha6\operatorname{Tr}_F\Phi^2,
 \qquad
 \mathcal V=V_0+\alpha\operatorname{Tr}_F(\Phi^4-2L\Phi^2),
 \qquad L=\Lambda_g^2>0.
\]
Here $\kappa_J,\alpha>0$ are the inherited action coefficients. If a
constant term is allocated to a residual source, $V_0$ denotes its
same-owner constant density in this particular vacuum comparison;
setting it to zero is an additional branch specialization, not a
vacuum subtraction. Assume the actual source-free comparison has
zero gauge curvature, $D_\mu\Phi=0$, no other retained source or
metric constituent, and a regular positive-graviton chart
$\mathscr F>0$. Consume the full native scalar equation on its
admitted matrix support, rather than only a selected ray variation.
In particular the actual variation slot must see
$\Phi^3+(R/12-L)\Phi$ faithfully; stationarity of its projection onto
an arbitrary restricted plane does not imply this matrix equation.
The functional throughout is the native ordinary $\operatorname{Tr}_F$.
It is not replaced by a rooted comparison functional without the
separate actual trace-incidence certificate.
Then constant Einstein vacua are exactly
\[
 \boxed{\mathscr F R=2\mathcal V,\qquad
       \Phi^3=(L-R/12)\Phi.}
\]
The zero-scalar solution has $R=2V_0/\kappa_J$.
For a nonzero solution let $n=\operatorname{rank}\Phi$ and write its
nonzero eigenvalues as $\pm v$, $v>0$, with whatever multiplicities
the actual scalar support allows. If
$D_n=6\kappa_J-\alpha nL\ne0$, all its scalar and curvature values
are explicitly
\[
 \boxed{v^2=\frac{6\kappa_JL-V_0}{6\kappa_J-\alpha nL},\qquad
 \mathscr F=\kappa_J-\frac{\alpha n v^2}{6},\qquad
 R=12(L-v^2).}
\]
Only $v^2>0$, $\mathscr F>0$ and the actual allowed spectral patterns
are retained. These algebraic equations are also sufficient for the
stated constant sector when its parallel scalar and Einstein metric
are admitted by the actual bundle/domain.

At $D_n=0$ there is no nonzero rank-$n$ solution unless
$V_0=6\kappa_JL$. If both equalities hold, every admitted
$0<v^2<L$ on that rank obeys the constant equations, with
\[
 \mathscr F=\kappa_J(1-v^2/L),\qquad R=12(L-v^2).
\]
For $g_E=(\mathscr F/\kappa_J)g$ the Einstein coefficient is
$\kappa_J$, and the potential and Einstein scalar curvature on this
exceptional family are exactly
\[
 W_E=\frac{\kappa_J^2\mathcal V}{\mathscr F^2}
          =6\kappa_JL,\qquad R_E=12L.
\]
\end{theorem}
\begin{proof}
The native metric equation at a parallel scalar has no derivative
terms and is $2\mathscr F G_{\mu\nu}=-\mathcal V g_{\mu\nu}$.
On an Einstein metric this is the first boxed equation. The native
scalar equation is
$\alpha R\Phi/3+4\alpha(\Phi^3-L\Phi)=0$, giving the second.
Hermiticity implies every nonzero eigenvalue has square
$v^2=L-R/12>0$, and $\operatorname{Tr}_F\Phi^2=nv^2$.
Combining the scalar equation with the metric trace gives
\[
 2\kappa_JR=4V_0-4\alpha Lnv^2,
 \qquad (6\kappa_J-\alpha nL)v^2=6\kappa_JL-V_0.
\]
Solving proves the nondegenerate formulas and the exceptional
alternative. The converse substitutes them in both full constant
equations. In the exceptional case
$\mathcal V=6\kappa_JL(1-v^2/L)^2$, so the Einstein potential is
constant, and constant Weyl scaling gives
$R_E=\kappa_JR/\mathscr F=12L$.
\end{proof}

\begin{proposition}[Evaluated radial and dormant scalar Hessians at those vacua]
\label{p3f:prop:v248-native-vacuum-masses}
On a nonzero vacuum of the theorem use the Einstein metric above and
the canonical scalar kinetic form inherited from that same Weyl
transport. If the scaling ray of its actual matrix $\Phi$ is retained,
its canonically normalized radial potential curvature is
\[
 \boxed{m_{\rm rad}^2
       =\frac{2v^2(6\kappa_J-\alpha nL)}{3\mathscr F}.}
\]
It vanishes on the exceptional family. If a nonzero admissible
Hermitian scalar tangent $T$ is supported in $\ker\Phi$, so
$\Phi T=T\Phi=0$, its canonically normalized potential curvature is
\[
 \boxed{m_{\rm ker}^2=-\frac{2\kappa_Jv^2}{\mathscr F}<0.}
\]
Thus every such dormant scalar direction is an actual negative
direction of the Einstein potential. At the zero-scalar vacuum the
common canonical potential curvature is
\[
 \boxed{m_0^2=\frac{V_0}{3\kappa_J}-2L.}
\]
These are local potential curvatures on the declared source-free
constant sector. They do not by themselves establish a global
Lorentzian instability with unspecified boundary conditions, or the
spectrum of an action with additional source/gauge/higher terms.
\end{proposition}
\begin{proof}
On a rank-$n$ scaling ray, put $x=v^2$ so
$\mathscr F=\kappa_J-\alpha nx/6$ and
$\mathcal V=V_0+\alpha n(x^2-2Lx)$. Direct differentiation gives
\[
 \frac{dW_E}{dx}
 =\frac{\alpha n\kappa_J^2}{3\mathscr F^3}
       [(6\kappa_J-\alpha nL)x-(6\kappa_JL-V_0)].
\]
At stationarity its second derivative with respect to $v$ is
$4v^2\alpha n\kappa_J^2(6\kappa_J-\alpha nL)/(3\mathscr F^3)$.
The actual radial kinetic metric is
\[
 G_{vv}=\frac{2\alpha\kappa_J n}{\mathscr F}
         +3\kappa_J(\partial_v\log\mathscr F)^2
       =\frac{2\alpha\kappa_J^2n}{\mathscr F^2}.
\]
Division gives $m_{\rm rad}^2$.
For $\Phi T=T\Phi=0$, the first derivatives of $\mathscr F$ and
$\mathcal V$ vanish, whereas
$\mathscr F''=-\alpha\operatorname{Tr}_FT^2/3$ and
$\mathcal V''=-4\alpha L\operatorname{Tr}_FT^2$.
Using $\mathcal V/\mathscr F=R/2$ gives
\[
 W_E''[T,T]
 =\frac{\kappa_J^2}{\mathscr F^2}
       \alpha(R/3-4L)\operatorname{Tr}_FT^2
 =-\frac{4\alpha\kappa_J^2v^2}{\mathscr F^2}
                      \operatorname{Tr}_FT^2.
\]
Its kinetic norm is
$2\alpha\kappa_J\operatorname{Tr}_FT^2/\mathscr F$, proving the
second formula. At $\Phi=0$ the metric is already Einstein and the
quadratic equation has curvature coefficient $R/6-2L$, which with
$R=2V_0/\kappa_J$ gives the last formula.
\end{proof}

\begin{corollary}[All-rank obstruction and global potential minimum for the actual bare constant]
\label{p3f:cor:v248-bare-all-rank-vacuum}
Use the actual bare constant
$V_0^{\rm bare}=12\kappa_JL$ of
Proposition~\ref{p3f:prop:v2216-bare-heat-vacuum}, with its
same-support coefficient ancestry. On the preceding constant sector
there is no nonzero Hermitian Einstein--scalar vacuum of any rank
with $\mathscr F>0$. The only one is
\[
 \boxed{\Phi=0,\qquad R=24L,\qquad m_0^2=2L.}
\]
More generally, as a function on the entire positive-graviton
Hermitian scalar chart, put
$x=\alpha\operatorname{Tr}_F\Phi^2/(6\kappa_J)\in[0,1)$ and
$Y=\operatorname{Tr}_F\Phi^4\ge0$. The exact formal Einstein
potential of this same truncated local packet is
\[
 \boxed{W_E(\Phi)=\frac{12\kappa_JL}{1-x}
                    +\frac{\alpha Y}{(1-x)^2}
       \ge\frac{12\kappa_JL}{1-x}\ge12\kappa_JL.}
\]
Equality at the lower endpoint holds only for $\Phi=0$; the potential
is strictly increasing on every nonzero retained scaling ray and
blows up at $\mathscr F\downarrow0$. This strengthens the earlier
single canonical-ray result without assuming a full-rank Clifford
ansatz or identifying arbitrary scalar supports. The potential proof
needs only that the actual nonzero configuration retains its scaling
variation. Thus the exclusion of nonzero critical points also applies
on any narrower homogeneous linear scalar support with that property,
even when the full matrix classification hypothesis is unavailable.
\end{corollary}
\begin{proof}
For $V_0=12\kappa_JL$, positivity of the formula for $v^2$ would
require $\alpha nL>6\kappa_J$ and give
\[
 v^2=\frac{6\kappa_JL}{\alpha nL-6\kappa_J},\qquad
 \mathscr F=-\frac{6\kappa_J^2}{\alpha nL-6\kappa_J}<0.
\]
The denominator-zero case also fails its required numerator equality.
Thus every nonzero rank is excluded. The origin values follow from
the theorem and proposition. Independently,
$\mathcal V=12\kappa_JL(1-x)+\alpha Y$ and
$\mathscr F=\kappa_J(1-x)$ give the displayed potential identity.
Along a nonzero scaling ray, $x$ grows quadratically and $Y$ quartically,
so both nonconstant nonnegative terms are strictly increasing wherever
$x<1$. The lower bound gives the endpoint divergence.
\end{proof}

\begin{corollary}[Exact promotion condition for an isolated native quartic extremum]
\label{p3f:cor:v248-quartic-vacuum-promotion}
Suppose an admitted nonzero constant matrix obeys the isolated native
quartic extremum condition $\Phi^3=L\Phi$ with rank $n$. It solves the
coupled constant Einstein--scalar equations only if
\[
 \boxed{R=0,\qquad V_0=\alpha nL^2,\qquad
                    \alpha nL<6\kappa_J.}
\]
These conditions are sufficient on the stated constant geometric
and scalar domain. In particular they cannot hold for the actual
bare $V_0=12\kappa_JL$. Isolated scalar minimization therefore does
not supply the native bare coupled vacuum.
\end{corollary}
\begin{proof}
The full scalar equation forces $R\Phi=0$, hence $R=0$.
Then the metric equation requires $\mathcal V=0$; at this quartic
extremum $\mathcal V=V_0-\alpha nL^2$. The strict inequality is exactly
$\mathscr F>0$. For the bare value, the equality would instead force
$\alpha nL=12\kappa_J$, contradicting that inequality. Conversely
these conditions satisfy both constant equations.
\end{proof}

These formulas close the constant coupled vacuum calculation on the
specified local action, including its absolute constant. They do not
close the vacuum of the complete framework. That further claim still
consumes its actual constant allocation, admissible parallel scalar
support, source preparation and residual gauge, normal, boundary and
higher-metric terms. The formal bare origin has $R/L=24$, so it remains outside the
controlled low-curvature expansion, exactly as in the earlier bare
heat test. The all-rank result is not a controlled prediction of that
high-curvature solution or a global vacuum theorem for the full heat
parent. The rooted midpoint scalar has a separately
evaluated spatial covariant-gradient contribution; it is not silently
put on the present parallel-scalar sector.


\section{Source-free active vacuum adjoining the terminal scalaron edge}
\label{p3f:sec:v346-near-terminal-vacuum}
The following statement concerns the \emph{restored one-variable metric
$f(R)$ action} with the same local jet fixed in
Theorem~\ref{p3f:thm:v263-dressed-born-legendre-descent}; its
constant-curvature trace equation is not the complete joint-action
vacuum equation in the presence of other higher-curvature, matter or
nonlocal terms.  Set $C:=C_7>0$ and assume the retained hierarchy
$U>5C$ and $3C<C_2$.  On $C<R\le3C$ one has $D_2=1$,
\[
 \chi_7(R)=\frac{(1+n_7(R))(R^2-C^2)}{(U+C)R^2},
 \qquad 0\le n_7(R)\le1,
\]
and the terminal-normalized metric coefficient is
$ f(R)=R-C+\int_C^R(R-x)\chi_7(x)\,dx$.

\begin{theorem}[Sharp neighborhood for the actual source-free active vacuum]
\label{p3f:thm:v346-near-terminal-active-vacuum}
For $T(R):=Rf'(R)-2f(R)$ one has the exact identity
\[
 \boxed{T(R)=2C-R+\int_C^R(2x-R)\chi_7(x)\,dx.}
\]
There is no source-free constant-curvature root in the terminal sector
$R\le C$ or in the adjoining active interval $C<R\le2C$.
There is exactly one such root $R_*$ in $(2C,3C)$, with explicit bounds
\[
 \boxed{
 \frac{(2-2\log2)C^2}{U+5C}
 \ \le\ R_*-2C\ \le\ \frac{2C^2}{U-5C}.
 }
\]
In particular $R_*=2C+O(C^2/U)$, while $f''(R_*)>0$:
the self-consistent source-free vacuum is on the \emph{active} branch,
not the affine terminal stratum where $f''=0$.
\end{theorem}
\begin{proof}
For $R\le C$, $f_{\rm term}(R)=R-C$ gives
$T(R)=2C-R>0$.  On $C<R\le2C$, both
$2C-R\ge0$ and $2x-R\ge0$ for $x\ge C$,
so $T(R)>0$.  At $R=2C$, boundedness
$1\le1+n_7\le2$ implies
\[
 \frac{(2-2\log2)C^2}{U+C}\le T(2C)
 \le\frac{2C^2}{U+C}.
\]
At $R=3C$, discarding the negative part of the integral gives
$T(3C)\le-C+9C^2/[2(U+C)]<0$.
Differentiation yields
\[
 T'(R)=-1+R\chi_7(R)-\int_C^R\chi_7(x)\,dx,
\]
and, for $2C\le R\le3C$,
$-1-4C/(U+C)\le T'(R)\le-1+6C/(U+C)<0$.
Thus precisely one zero lies in $(2C,3C)$; the mean-value theorem
combined with the two-sided bounds on $T(2C)$ proves the displayed
curvature estimates.
\end{proof}

\begin{corollary}[Positive trace-mode scalaron mass at the active vacuum]
\label{p3f:cor:v346-active-vacuum-scalaron}
On the trace-linearized metric $f(R)$ domain of
Theorem~\ref{p3f:thm:trace-linearized-scalaron-mass}, the scalaron at
$R_*$ satisfies $f'(R_*)>0$, $f''(R_*)>0$, and
\[
 \boxed{m_{\rm sc}^2
 =\frac{f'(R_*)-R_*f''(R_*)}{3f''(R_*)}
 =\frac{4(U+C)}{9[1+n_7(2C)]}
  \left[1+O\!\left(\frac CU\right)\right]>0.}
\]
At $\sigma_{\rm fin}=0$, the formal half-determinant convention gives
$C/U\simeq4.216119310922\cdot10^{-121}$, and
$n_7(2C)\simeq0.00250704$, so $m_{\rm sc}^2/U\simeq0.44333$.
This is not an all-spin or nonlocal ghost-freedom certificate.
\end{corollary}
\begin{proof}
At $2C$, $\chi_7=3[1+n_7(2C)]/[4(U+C)]$.
The location bound gives $R_*=2C+O(C^2/U)$ and hence
$f'(R_*)=1+O(C/U)$, with the same leading susceptibility at $R_*$.
Substitute into the existing trace-mode scalaron mass formula; positivity
uses $U>5C$ and the hierarchy $U/C\gg1$.  This proof is a
near-terminal asymptotic on the certified jet branch, not a derivation of
the complete physical cosmological history.
\end{proof}

\begin{remark}[Affine terminal Einstein root is an extrapolation]
\label{p3f:rem:v346-terminal-affine-extrapolation}
The affine representative $f_{\rm term}(R_g)=R_g-C$ defines the
cosmological coefficient $\Lambda_{\rm eff}=C/2$ and has the formal
Einstein constant-curvature solution $R_g=2C$.  If its \emph{valid
terminal domain} is $R_g\le C$, that solution lies outside this
domain.  It must not be counted as a terminal-branch source-free vacuum
alongside $R_*$.  A distinct physical rank-changing branch and its
background matching require the actual action/constraint and source
certificate; one cannot infer it by analytic continuation of a
terminal curvature formula.
\end{remark}

\section{Evaluated normalization and native root readouts}
\label{p3f:sec:native-source-root-readout}

The even-insertion coefficient in
Corollary~\ref{p3f:cor:cartan-even-source-cubic-error} is already
constructively specified. Theorem~\ref{p1b:thm:native-source-scalar-weight}
now proves that its value for the constant insertion cannot vanish,
and Theorem~\ref{p1b:thm:native-source-root-interval} supplies an explicit
monotone interval. The following evaluates that coefficient in closed
form for the nonconstant mode already present in
Proposition~\ref{p3f:prop:m8td-explicit-nonconstant-zero}.

\begin{proposition}[A nonconstant Born-equivariant zero mode]
\label{p3f:prop:m8td-explicit-nonconstant-zero}
Let $s=q_3$ be an imaginary quaternion coordinate on the unit sphere and
set $u=\epsilon s$. Put $y=1-s^2$, $a=\sqrt3\epsilon/2$ and
\[
 F(y)=\frac{\sinh(a\sqrt y)}{a\sqrt y}
     =\sum_{m\ge0}\frac{(a^2y)^m}{(2m+1)!}.
\]
Then a smooth, unnormalized zero of
\eqref{p3f:eq:m8td-variable-cartan-dirac} is
\begin{equation}
 \phi=F(1-s^2),\qquad
 \alpha=\frac{4R}{\epsilon}F'(1-s^2)\,ds.
 \label{p3f:eq:m8td-explicit-nonconstant-zero}
\end{equation}
It extends at $\epsilon=0$ to $(1,0)$ and is even under $\mathcal U$.
For $|\epsilon|\le1$ it is the unique doublet zero. Its expansion is
$\phi=1+\epsilon^2(1-s^2)/8+O(\epsilon^4)$,
$\alpha=R\epsilon\,ds/2+O(\epsilon^3)$.
\end{proposition}
\begin{proof}
For the positive scalar Laplacian on the round sphere,
\(d^\dagger ds=3s/R^2\), \(|ds|^2=(1-s^2)/R^2\), and
\(dy=-2s\,ds\). The entire defining series of \(F\) obeys
\[
 2yF''(y)+3F'(y)=\frac{3\epsilon^2}{8}F(y).
\]
Since \(d\alpha=0\), the one-form row of
\eqref{p3f:eq:m8td-variable-cartan-dirac} is
\(d\phi+u\alpha/(2R)=0\). For the scalar row,
\[
 d^\dagger\alpha
 =\frac{4s}{\epsilon R}\bigl(3F'+2yF''\bigr)
 =\frac{3\epsilon s}{2R}F,
\]
which cancels its Cartan scalar multiplier. Thus the stated pair is a
zero for \(\epsilon\ne0\). The entire series gives the removable
zero-amplitude extension and both displayed Taylor coefficients.
Inversion sends \(s\) to \(-s\), making the scalar even and the
pulled-back one-form odd. Uniqueness on \(|\epsilon|\le1\) follows from
\eqref{p3f:eq:m8td-nonhomogeneous-gap}. This is the same explicit
representative as Paper~I-C,
Proposition~\ref{p1b:result:an-explicit-inhomogeneous-zero-mode}.
\end{proof}

\begin{proposition}[Closed normalization of the supplied dipole zero]
\label{p3f:prop:native-source-dipole-normalization}
Take \(u=a h\), \(h=q_3\), \(|a|\le1\), put
\(k=\sqrt3a/2\), \(y=1-h^2\), and retain the existing raw mode
\[
 \phi_a=\frac{\sinh(k\sqrt y)}{k\sqrt y},\qquad
 \alpha_a=\frac{4R}{a}\frac{d}{dy}
       \left(\frac{\sinh(k\sqrt y)}{k\sqrt y}\right)dh.
\]
With averages defined by \(\langle f\rangle=V_R^{-1}\int f\,d\mu_R\),
its scalar and one-form squared norms are
\begin{align}
 A_\phi(k)&=\langle\phi_a^2\rangle
     =\frac{I_0(2k)-1}{k^2},
       \label{p3f:eq:native-source-dipole-scalar}\\
 A_\alpha(k)&=\langle|\alpha_a|^2\rangle
     =\frac{3(I_0(2k)+1)}{k^2}-\frac{6I_1(2k)}{k^3}.
       \label{p3f:eq:native-source-dipole-oneform}
\end{align}
The removable values at \(k=0\) are \(A_\phi=1\), \(A_\alpha=0\).
The normalized zero is \((\phi_a,\alpha_a)/\sqrt{V_R(A_\phi+A_\alpha)}\),
and its constant-offset slope is exactly
\begin{equation}
 \kappa_a=\frac{1}{2R}
       \left(1-\frac{4A_\phi(k)}{A_\phi(k)+A_\alpha(k)}\right).
 \label{p3f:eq:native-source-dipole-slope}
\end{equation}
\end{proposition}
\begin{proof}
Put \(h=\cos\chi\), \(r=\sin\chi\).
The normalized coordinate measure is \((2/\pi)r^2d\chi\) and
\(|dh|=r/R\). The raw one-form length is
\[
 |\alpha_a|^2=
 \frac{3(kr\cosh(kr)-\sinh(kr))^2}{k^4r^4}.
\]
The integral representation of \(I_0\) gives the scalar norm directly
\cite{NativeSourceDLMF1032}. For the one-form put
\(K(k)=\int_0^\pi\sinh^2(k\sin\chi)/\sin^2\chi\,d\chi\).
It has \(K(0)=K'(0)=0\),
\(K''(k)=2\pi I_0(2k)\), and therefore
\(kK'(k)-K(k)=\pi k I_1(2k)\). Expansion of the squared numerator
in the preceding formula gives the stated \(A_\alpha\).
All apparent endpoint singularities are removable from the original
hyperbolic series. In particular
\[
 A_\phi=1+k^2/4+k^4/36+O(k^6),\qquad
 A_\alpha=k^2/4+k^4/24+O(k^6).
\]
Normalization and Hellmann--Feynman give the slope. The previously
supplied zero itself is not a new construction here; its evaluated
normalization and response coefficient are the new conclusions.
\end{proof}


The separately declared square has a unique quadratic energy ground and
two local offset roots for each admitted positive target; Paper~III-G, Appendix~
\ref{p3f:app:native-source-square-readout} states the additional operator
choice and proves the bounds. The native signed readout does not supply
a physical terminal-to-rest dictionary.


\subsection{A unique offset after a supplied complete bounded insertion}
\label{p3f:sec:native-shifted-root}

Retain the preceding actual \(D_u,s,P,Q,G_u,g\). Supply one bounded
self-adjoint operator \(V\) on the same \(\mathcal H\), with
\(w=\lVert V\rVert\); its whole operator is retained. Define
\[
 J=\frac1{2R}\operatorname{diag}(-3,1,1,1),\quad j=\lVert J\rVert=\frac3{2R},
 \qquad D(\eta,t)=D_u+\eta J+tV\quad\hbox{on }H^1.
\]
Here \(\eta,t\) are real dimensionless amplitudes; a target \(e\)
is a signed spatial Dirac readout, with units \(R^{-1}\).

\begin{theorem}[Quantified shifted signed root]
\label{p3f:thm:native-shifted-signed-root}
Let \(\delta=gR/18\). If
\begin{equation}
 |e|+w|t|<g/90,
 \label{p3f:eq:native-shifted-root-target}
\end{equation}
there is exactly one \(\eta_*(t;e)\in(-\delta,\delta)\) with
\(\lambda(\eta_*(t;e),t)=e\), where \(\lambda\) is the simple branch
continuing the native zero. Throughout this offset interval,
\begin{equation}
 -\frac3{2R}\le\partial_\eta\lambda(\eta,t)\le-\frac1{5R},\qquad
 |\eta_*(t;e)|\le5R(|e|+w|t|).
 \label{p3f:eq:native-shifted-root-slope}
\end{equation}
The branch and its local inverse are real analytic. For \(e=0\),
the zero is a zero of the full spatial operator, and all other signed
eigenvalues stay outside \((-g/2,g/2)\).
\end{theorem}
\begin{proof}
The scalar-weight and slope premise is
Theorem~\ref{p1b:thm:native-source-scalar-weight}. For the normalized native zero \(s=(\phi,\alpha)\), put
\(F=\lVert\phi\rVert^2\). The lowest one-form Laplacian eigenvalue is
\(3/R^2\); the zero equation gives
\[
 3(1-F)\le\tfrac14 U_0^2(9F+1-F),\qquad
 F\ge\frac{12-U_0^2}{12+8U_0^2}\ge\frac{11}{20}.
\]
For \(W=\eta J+tV\), \(v=j|\eta|+w|t|\), the hypotheses imply
\(v<g/12+g/90=17g/180<g/6\). Resolvent Neumann convergence on
\(|z|=g/2\) preserves the rank-one low Riesz projection. The sole
eigenvalue in that circle has \(|\lambda|\le v\), by normal-resolvent
spectral inclusion. For its unit vector \(s_{\eta,t}\), the old-complement
equation gives
\[
 \lVert Qs_{\eta,t}\rVert\le\frac{v}{g-|\lambda|}
 \le\frac{v}{g-v}<\frac15.
\]
The expectation of the scalar-form projector changes by at most this
angle sine. Thus its new expectation is at least \(7/20\).
Differentiating the simple eigenvalue yields
\(\partial_\eta\lambda=(1-4F_{\eta,t})/(2R)\), giving the bounds.
Also \(|\lambda(0,t)|\le w|t|\). Integrating the negative slope to
\(\eta=\pm\delta\) brackets every target in the stated range, since
\(\delta/(5R)=g/90\). The intermediate-value theorem, strict
monotonicity and the implicit-function theorem give existence,
uniqueness, analyticity and the displacement bound. The same contour
argument excludes any other zero of the complete spatial operator.
\end{proof}

\begin{theorem}[Displacement coefficients and an explicit residual bound]
\label{p3f:thm:native-shift-displacement}
Define real coefficients
\[
 \kappa=\langle s,Js\rangle\le-3/(5R),\quad
 a_0=\langle s,Vs\rangle,\quad
 d_0=\operatorname{Re}\langle Js,G_uVs\rangle,\quad
 c_0=\langle Vs,G_uVs\rangle.
\]
For the zero target,
\begin{equation}
 \eta_*(t;0)=A_1t+A_2t^2+O(t^3),\qquad
 A_1=-a_0/\kappa,\qquad A_2=(2d_0A_1+c_0)/\kappa.
 \label{p3f:eq:native-shift-displacement-coefficients}
\end{equation}
More generally set
\[
 \eta_1=(e-a_0t)/\kappa,\quad
 \eta_2=(2d_0\eta_1t+c_0t^2)/\kappa,\quad
 \eta_{\dagger}=\eta_1+\eta_2,\quad v_{\dagger}=j|\eta_{\dagger}|+w|t|.
\]
If the target hypothesis holds, \(|\eta_{\dagger}|\le\delta\), and
\(v_{\dagger}+|e|<g\), the approximation has the computable bound
\begin{equation}
 |\eta_*-\eta_{\dagger}|\le5R\left(
 2|d_0\eta_2t|+
 \frac{v_{\dagger}^2(v_{\dagger}+|e|)}{g(g-v_{\dagger}-|e|)}\right).
 \label{p3f:eq:native-shift-displacement-remainder}
\end{equation}
All coefficients use the actual moving complementary inverse of the
preceding theorem and the whole supplied insertion.
\end{theorem}
\begin{proof}
The Born involution fixes \(s\), commutes with \(J\), and anticommutes
with \(G_u\), so \(\langle Js,G_uJs\rangle=0\).
The exact scalar Schur function is
\[
 k(\eta,t,z)=\kappa\eta+a_0t-z
 -\langle QWs,(QD_uQ+QWQ-z)^{-1}QWs\rangle.
\]
Its complementary Neumann expansion is
\[
 k=\kappa\eta+a_0t-z-2d_0\eta t-c_0t^2+r,\qquad
 |r|\le\frac{v^2(v+|z|)}{g(g-v-|z|)}.
\]
Cancelling the first two orders gives the stated coefficients. At the
approximation and \(z=e\), the uncancelled polynomial term is
\(-2d_0\eta_2t\). For real \(z\) between the continuing eigenvalue and
\(e\),
\(\partial_z k=-1-\lVert(QD_uQ+QWQ-z)^{-1}QWs\rVert^2\le-1\).
Thus \(|\lambda-e|\le|k(\eta_{\dagger},t,e)|\). The offset slope bound
then gives exactly the displayed root-error estimate. This also
supplies a controlled cubic error for the zero-target expansion.
\end{proof}


If the application additionally licenses only \(\lVert u+\eta\rVert_\infty\le1\),
oddness implies the exact further condition \(U_0+|\eta_*|\le1\).
In particular \(U_0=1\) admits no nonzero constant offset in that subrange.
For the separately declared square \(D(\eta,t)^2\), the local low energy
is \(\lambda(\eta,t)^2\); a positive small energy therefore has two
offset branches whenever both signed targets \(\pm\sqrt E\) satisfy
the target and profile conditions. An actual tagged electron polynomial
requires its own complete-action and source identification.


\subsection{Actual weak carrier and midpoint Morse--Dirac propagation}
\label{p3f:sec:v2183-native-weak-carrier-propagation}

\begin{theorem}[The generated weak doublet is the trivial G15 class with pure fixed isotropy]
\label{p3f:thm:v2183-native-weak-carrier-class}
Let
\[
 SU(2)_{\rm diag}\longrightarrow X_7^{\rm cmp}
 \longrightarrow K_{\rm rel}\cong S^3_{\rm rel}\times S^1_{\rm rel}
\]
be the compact residual principal bundle generated in Paper~III-B, equipped
with the descended Born connection \(\omega_{\rm cmp}\).  In the canonical
presentation
\(X_7^{\rm cmp}\cong SU(2)_x\times SU(2)_p\times U(1)_{\rm rel}\), the map
\[
 s(r,z)=(r,e,z)
\]
is a global principal-bundle section.  Hence the fundamental associated
doublet \(E_{\rm w}\to K_{\rm rel}\) is ordinarily trivial.

Pull it back along the relative map
\[
 r_{M_8}:M_8\longrightarrow K_{\rm rel},
 \qquad (a,b)\longmapsto ab^{-1}
\]
with the corresponding relative scalar coordinate.  Then
\[
 \boxed{r_{M_8}^*E_{\rm w}\cong M_8\times\mathbb C^2,}
\]
so \(c_2=0\), and its restriction to
\(D_6=[M_8,M_8]\cong S^3\times S^3\) has
\[
 \boxed{\epsilon_{\rm IB15}=0.}
\]
In the mixed-Chern coordinate convention
\(c_2=A\mathcal D\), this ordinary class has \(A=0\).  Thus the pair used for
this rank-two comparison is
\[
 \boxed{(A,\epsilon_{\rm IB15})=(0,0).}
\]

Moreover, the actual descended Born action on the exchange fixed locus
\(F=\Delta K\subset M_8\) acts on the fundamental fibre by \(+I_2\).
Consequently this generated weak carrier is a pure determinant-preserving
principal-\(SU(2)\) Born lift, not a mixed fixed-fibre \((1,1)\) lift.
The Paper-I-B mixed coherent counts \(72/120\) therefore do not count the
physical weak carrier.
\end{theorem}
\begin{proof}
Paper~III-B, Theorems~\ref{p3b:thm:canonical-compact-seven-quotient} and
\ref{p3b:thm:parent-cartan-scalar-diagonal-descent}, has generated the
actual residual vertical carrier, induced exchange and inherited connection.
Apply Paper~I-C, Theorems~\ref{p1b:thm:minimal-native-provider} and
\ref{p1b:thm:conditional-residual-descent}, to this diagram.  They give
the global section, ordinary triviality and the explicit determinant-one
gauge \(u(a,b)=b^{-1}\) that trivializes the pulled-back exchange cocycle.
Thus \(c_2=0\), the restriction bit is zero, and the fixed action on
\(\Delta K\) is \(+I_2\).  The mixed Chern coordinate is consequently
\(A=0\).  These conclusions evaluate the generated carrier itself.

On a relative spatial sphere the same induced lift instead has fixed signs
\((+,-)\), by Theorem~\ref{p1b:thm:minimal-spatial-lifts}.  The fixed
set there is \(\{1,-1\}\); on \(M_8\) it is \(\Delta K\), where
the relative coordinate equals one.  The two assertions concern different
quotient maps and are consistent.  This pure principal lift is outside
the mixed \(U(2)\) counting domain.  No generic descent-fibre enumeration
is used.
\end{proof}

\begin{corollary}[Discharge of Paper I-C's provisional native \(SU(2)\) premise]
\label{p3f:cor:v259-native-su2-discharge}
For the compact weak principal role generated here,
\(\mathsf H_{\rm vert}\) of Paper~I-C,
Definition~\ref{p1b:def:conditional-residual-datum}, is proved rather
than imposed.  Its group is the connected derived relative-vertical group,
its action is the actual residual right action after the scalar quotient,
and its exchange is induced from the source.  Its inherited connection is
the parent-descended midpoint.  Consequently the conditional native
weak-sector applications in Paper~I-C apply to this generated datum,
with each application's displayed additional inputs retained.
\end{corollary}
\begin{proof}
The generated relative diagram in Paper~III-B,
Theorem~\ref{p3b:thm:canonical-compact-seven-quotient}, uses the actual
fibres of \((a,b)\mapsto ab^{-1}\), the scalar diagonal and inherited
exchange.  The relative-carrier rigidity and derived-group calculation of
Paper~I-C, Theorem~\ref{p1b:thm:minimal-native-provider}, identify the
remaining group and right action as precisely those in the definition of
\(\mathsf H_{\rm vert}\).  Theorem~\ref{p3b:thm:parent-cartan-scalar-diagonal-descent}
generates its inherited connection.  Theorem~\ref{p3f:thm:v2183-native-weak-carrier-class}
above identifies the generated fundamental weak carrier and its actual
ordinary and equivariant class.  These are exactly the stated native
role data.  The mathematical constructions in I-B were proved without
this later occurrence conclusion, so their application introduces no
circular premise.  This proves the native antecedent, not existence of a
descent for an arbitrary principal bundle or every generic fibre count.
\end{proof}

\begin{corollary}[Minimal Born protection before midpoint specialization]
\label{p3f:cor:v259-minimal-born-protection}
On either exchange-invariant relative scalar slice, every smooth connection
preserving the native induced weak lift has Chern--Simons class \(1/2\),
full \(SU(2)\) holonomy and fundamental reciprocal Dirac index one on any
supplied inversion-invariant spatial metric.  Its spatial kernel is odd
and nonzero.  Simplicity and a numerical gap require their stated analytic
premises.  For the generated parent-descended midpoint and supplied round
radius, Paper~I-C's complete every-isospin spectrum and local response
theorems apply.  Other associated representations here are mathematical
constructions; this corollary does not generate additional matter species.
\end{corollary}
\begin{proof}
Apply Paper~I-C, Theorems~\ref{p1b:thm:minimal-born-cs-holonomy},
\ref{p1b:thm:minimal-born-dirac-index},
\ref{p1b:thm:minimal-all-isospin-dirac} and
\ref{p1b:thm:minimal-born-zero-jets}, on their metric, connection,
spin and norm domains.  Theorem~\ref{p1b:thm:minimal-provider-consumer-closure}
transports the complete consumed data.  The action and source applications
below keep their generated Yang--Mills and source inputs.
\end{proof}

\begin{corollary}[The weak Morse mode is the protected Dirac crossing coordinate]
\label{p3f:cor:v2183-native-weak-morse-dirac}
Restrict the generated weak connection to the retained round
\(S_R^3\) and to the classical pure-gauge constituent of the generated local
Yang--Mills action.  Then the midpoint
\(A_{\rm B}=\Theta/2\) has, modulo infinitesimal gauge transformations,
exactly one negative spatial Yang--Mills mode,
\[
 \lambda_{\rm YM,-}=-\frac2{R^2},
\]
no physical spatial zero mode, and positive physical spectrum beginning at
\(1/R^2\).  The unique negative mode is the homogeneous scalar
\(\mathbf1\subset\mathbf3_{\rm space}\otimes\mathbf3_{\rm weak}\).
By Paper~I-C, Theorem~\ref{p1b:thm:native-morse-dirac-alignment}, this is
also the unique physical homogeneous weak-gauge direction acting nontrivially
on the protected fundamental-doublet Dirac zero.

Along the corresponding Cartan coordinate \(A_f=f\Theta\), its signed
protected eigenvalue \(\lambda\), Chern--Simons coordinate and Yang--Mills
barrier potential obey exactly
\begin{equation}
 \boxed{
 f=\frac12-\frac{R\lambda}{3},\qquad
 N_{\rm CS}=\frac12-\frac{R\lambda}{2}
             +\frac{2R^3\lambda^3}{27},\qquad
 \mathcal E_{\rm YM}=\frac14
 \left(1-\frac{4R^2\lambda^2}{9}\right)^2.}
 \label{p3f:eq:v2183-native-weak-morse-coordinate}
\end{equation}
Thus the same one-dimensional coordinate controls the specified midpoint
Morse direction, the protected Dirac crossing and the unit
Chern--Simons interpolation.
\end{corollary}
\begin{proof}
The spatial Morse statement and homogeneous source selection are
Paper~I-C, Theorems~\ref{p1b:thm:native-midpoint-gauge-morse} and
\ref{p1b:thm:native-morse-dirac-alignment}.  Their connection and trace
normalizations are the same midpoint data already identified in
Theorem~\ref{p3f:thm:unique-born-invariant-gred-connection}.  The coordinate
identities are Corollary~\ref{p1b:cor:native-dirac-cs-barrier-coordinate}.
The generated local Yang--Mills action supplies the physical pure-gauge
action role; no additional gauge-field selector is introduced.
\end{proof}

\begin{corollary}[Native common-heat barrier normalization]
\label{p3f:cor:v2183-native-weak-barrier-energy}
On the source-native common-heat branch, where
Paper~V-A gives \(g_2^2=2\pi^2\), the static homogeneous pure-weak
Yang--Mills barrier energy at the midpoint is
\begin{equation}
 \boxed{
 E_{\rm bar}^{\rm heat}
 =\frac{3\pi^2}{g_2^2R}
 =\frac3{2R}
 =|\lambda(f=0)|=|\lambda(f=1)|.}
 \label{p3f:eq:v2183-native-weak-barrier-energy}
\end{equation}
This is a native matching-datum equality on the stated round pure-gauge
branch.  It is not an observed particle-mass relation and does not identify
\(R^{-1}\) with an experimental renormalization scale.
\end{corollary}
\begin{proof}
Write proper Euclidean time as \(t=R\tau\) in the Paper-I-B reduced
functional.  Its static energy is
\[
 E(f)=\frac{12\pi^2}{g_2^2R}\,4f^2(1-f)^2.
\]
At \(f=1/2\) this is \(3\pi^2/(g_2^2R)\).  Substitute the common-heat
normalization of Corollary~\ref{p5:cor:common-heat-absolute-gauge-normalization}
and the endpoint Dirac values in
Corollary~\ref{p1b:cor:native-dirac-cs-barrier-coordinate}.
\end{proof}

\begin{corollary}[Pure-gauge Lorentzian midpoint instability]
\label{p3f:cor:v2183-native-weak-lorentzian-instability}
On the same round homogeneous restriction of the generated Lorentzian
Yang--Mills constituent, put \(x=f-1/2\) and let \(t\) be proper time.
Then
\begin{equation}
 \boxed{
 \ddot x=\frac2{R^2}x-\frac8{R^2}x^3.}
 \label{p3f:eq:v2183-native-weak-lorentzian-x}
\end{equation}
Hence the pure-gauge midpoint has linear instability rate
\(\sqrt2/R\).  In the signed Dirac coordinate,
\begin{equation}
 \boxed{
 \ddot\lambda=\frac2{R^2}\lambda-\frac89\lambda^3.}
 \label{p3f:eq:v2183-native-weak-lorentzian-lambda}
\end{equation}
The exactly symmetric datum \(x=\dot x=0\) remains at the midpoint, and the
complete joint action may add matter or other Hessian blocks.  This corollary
therefore supplies a generated pure-gauge escape direction, not an occurring
cosmological or baryogenesis history.
\end{corollary}
\begin{proof}
The Euclidean reduced equation of Paper~I-C changes sign in the potential term
under the Lorentzian Yang--Mills action.  In dimensionless time this gives
\(x''=2x-8x^3\); replacing \(\tau\) by proper time \(t=R\tau\) yields
\eqref{p3f:eq:v2183-native-weak-lorentzian-x}.  The relation
\(\lambda=-3x/R\) gives
\eqref{p3f:eq:v2183-native-weak-lorentzian-lambda}.
\end{proof}

\begin{remark}[Downstream firewall]
\label{p3f:rem:v2183-native-weak-propagation-firewall}
The preceding results remove weak-bundle topology, homogeneous weak-connection
selection and the pure-gauge midpoint Morse direction from the unresolved
input list on their stated domains.  They do not construct the actual C01
finite quantization member \(\Gamma_e\), its complete remainder, occupation
and constraint projectors, stationary support or tagged electron source
polynomial.  Likewise, the existence of the unit Chern--Simons interpolation
and its pure-gauge instability does not actualize a Lorentzian charge-changing
history.  A21 still consumes the same-history CP/BL2 kinetic incidence,
collision/washout and spectator response, freeze-out, material entropy and
final yield normalization.
\end{remark}


The exact heat readout and its strictly negative Cartan curvature on this
spatial restriction are given in Paper~III-G, Appendix~\ref{p3f:app:m8td-heat}.
They determine a response of this specified operator, while full vacuum
stability requires the full action. Paper~III-G, Appendix~\ref{p3f:app:m8td-examples}
extends the main explicit nonconstant protected zero with contractions
and exact evolution equations for a supplied homogeneous history. The Born grading includes base
inversion; its index is distinct from four-dimensional Clifford chirality
and from the native GW/clutch index.

\subsection{Parameter comparison and Born connection}

\begin{remark}[Parameter versus Born connection comparison]
\label{p3f:rem:parameter-versus-born-connection-comparison}
The joint-orbit line of Paper I-A carries its unique unitary
connection preserving the derived constant-norm quadratic.  Its
base is the geometric parameter orbit \(\mathcal O\).
The Born connection of this paper has the reciprocal geometric
base and two fibre-spin extensions.  A comparison needs an actual
base map \(f:B\to\mathcal O\), a target representation and a bundle
intertwiner identifying the specified descendant with
\((f^*N_{\mathcal O},f^*\nabla_{N_{\mathcal O}})\).
Common \(M_8\) hosting supplies none of those identifications by
itself.  Both existing nonflat fibre-spin anchors are retained.
The curvature test in
Paper~III-G, Appendix~\ref{p3f:app:born-comparison-curvature-test} excludes
identifying the flat line with an unreduced full \(SU(2)\) spin
descendant, but does not exclude a correctly constituted reduced
or pulled-back comparison.
\end{remark}


\begin{definition}[Normalized Born vector ratio submersion]
\label{p3:def:normalized-born-product-metric}
On the derived carrier \(X_{7,U}^{\rm can}\), the residual
\((\mathrm{GRED}_7)\) branch chooses equal round bi-invariant auxiliary
comparison metrics \(g_x^{\rm aux},g_p^{\rm aux}\) on
\(SU(2)_x\times SU(2)_p\), distinct from the dynamical position metric varied
under \((\mathrm{VDOM}_x^{\min})\), and set
\[
 \mathcal G_{\rm B}^{\rm vec}
 :=\frac12(g_x^{\rm aux}\oplus g_p^{\rm aux}).
\]
For
\[
 \pi_{\rm rel}(x,p)=xp^{-1},
\]
the fibres are diagonal three-spheres
\[
 \pi_{\rm rel}^{-1}(b)=\{(bg,g):g\in SU(2)\}\cong S^3_{\rm diag}.
\]
With the displayed normalization, a vertical tangent \((A,A)\) has the same
norm as \(A\) on either original round three-sphere, so the vector-fibre radius
is
\[
 r_{\rm fib}=\Rmom.
\]
The scalar pair supplies a diagonal \(U(1)\) representation slot but no
physical scalar-circle metric modulus is assigned in the main constitutional
branch.
\end{definition}

\begin{definition}[Displayed connection, fibre average, and Yang--Mills convention]
\label{p3:def:gred-connection-generator-normalization}
Let \(X_7:=X_{7,U}^{\rm can}\xrightarrow{\pi}U\) be the derived principal
carrier, \(\dim X_7=7\), with fibre
\(F=S^3_{\rm diag}\cong SU(2)\).  Let
\(d\mu_F:=\operatorname{Vol}(F)^{-1}d\operatorname{vol}_F\), and for
\(s\in\{L,R\}\) choose a basis \(T_A^{(s)}\) of the corresponding isometry
ideal with Killing fields \(K_A^{(s)}\).  Write
\[
 [T_A^{(s)},T_B^{(s)}]
 =c^{C}_{\ AB}T_C^{(s)},
 \qquad
 A_s=A_s^AT_A^{(s)},
 \qquad
 F_s^A=dA_s^A+\frac12c^A_{\ BC}A_s^B\wedge A_s^C.
\]
The residual physical isometry-connection member is described by the
potentials \(A_s\); the connection metric is defined by the horizontal lifts
\[
\begin{aligned}
 X^H&:=X-\sum_{s,A}A_s^A(X)K_A^{(s)},
 &\widetilde g(X^H,Y^H)&=g_4(X,Y),\\
 \widetilde g(X^H,K_A^{(s)})&=0,
 &\widetilde g|_{TF}&=g_F.
\end{aligned}
\]
The normalized fibre-averaged generator metrics are
\begin{equation}
\label{p3:eq:gred-generator-metric}
 \Gamma_{AB}^{(s)}
 :=\int_F g_F(K_A^{(s)},K_B^{(s)})\,d\mu_F
 =\kappa r_{\rm fib}^2\delta_{AB},
 \qquad s=L,R,
\end{equation}
with vanishing left--right cross average.  Equal \(\kappa\) is forced once the
same invariant basis convention is used on the two isometric ideals; its
absolute value depends on that convention.  We compare the reduced term with
\begin{equation}
\label{p3:eq:gred-ym-convention}
 S_{\rm YM}^{(s)}
 :=-\frac1{4g_s^2}\int_U\sqrt{-g}\,
 \delta_{AB}F_{s,\mu\nu}^AF_s^{B,\mu\nu}\,d^4x.
\end{equation}
The declaration of
\eqref{p3:eq:gred-generator-metric}--\eqref{p3:eq:gred-ym-convention}, including
the generator basis, is \((\mathrm{CAL}_{\rm GRED}(\kappa))\).

For the period-\(2\pi\) quaternion Killing basis,
\(\Gamma_{AB}=r_{\rm fib}^2\delta_{AB}\) and \(\kappa=1\).  If the Paper-II
spin-\(\tfrac12\) generators \(L_A\) are represented geometrically by
\(K_A/2\), their averaged squared length is one quarter as large and
\(\kappa=1/4\).  Conversion between these basis conventions and any
phenomenological standard coupling is a later matching step.  The dependence
of Kaluza--Klein gauge coefficients on the fibre-integrated generator metric
is standard; see \cite{Baptista2024InternalSymmetries}.
\end{definition}

\subparagraph{Optional eight-dimensional diagnostic deferred.}
On \((\mathrm{BULK}_8)\), the Born-even projector gives ordinary
eight-dimensional bulk factor \(1/2\), while the fixed-set term is lower
dimensional; it yields no seven-dimensional coefficient without the separate
\((\mathrm{MATCH}_{8\to7})\) bridge.  The optional construction, proof, and
routing corollary are placed in
Paper~III-G, Appendix~\ref{p3f:app:born-even-fixed-set}.  None is a premise of the
seven-dimensional \((\mathrm{GRED}_7)\) theorem below.

\begin{lemma}[Protected scalar rank and corner weight]
\label{p3:lem:minimal-scalar-trace-multiplicity}
On the physical protected scalar branch,
\[
 \boxed{N_{\rm sc}^{\rm rank}=1},
 \qquad
 \boxed{w_{\rm sc}^{\rm corner}=\frac18}.
\]
The first statement is the rank of the surviving scalar object on
\(\mathsf C_{\rm IR}^{\rm der}\); the second is the ambient
normalized weight of its co-realized rank-one constituent, certified by
\((\mathrm{CORN})\).  They are distinct
typed facts and are not multiplied to define a new multiplicity.
\end{lemma}

\begin{proof}
The protected octet has exactly one scalar infrared object after the
\(O(3)\)-equivariant extraction.  Paper~III-B's co-realization theorem then
derives the ambient rank-one real-channel coefficient \(1/8\) certified by
\((\mathrm{CORN})\).  Current Paper~III-B explicitly
separates this readout from the regular complex-Hilbert \(\Dcl\) trace.
\end{proof}

\begin{lemma}[Totally geodesic diagonal vector fibres]
\label{p3:lem:totally-geodesic-born-odd-fibers}
For any positive constant weights on the two round product factors, every
fibre of \(\pi_{\rm rel}\) is a totally geodesic round three-sphere.  Hence the
O'Neill \(T\)-tensor vanishes.  Equal weights are needed only for the specific
Born-normalized induced-radius convention used later, not for total geodesy.
\end{lemma}

\begin{proof}
For fixed \(b\), the fibre is the graph of the left translation
\(L_b:g\mapsto bg\), an affine isometry of the bi-invariant connection.
Multiplying either factor metric by a positive constant does not change its
Levi--Civita connection.  For tangent fields \(X,Y\) on the graph, the product
covariant derivative is
\[
 (\nabla_XY,\nabla_{(L_b)_*X}(L_b)_*Y)
 =(\nabla_XY,(L_b)_*\nabla_XY),
\]
which is again tangent to the graph.  Thus its second fundamental form and
the O'Neill \(T\)-tensor vanish for arbitrary positive constant weights.
\end{proof}

\begin{proposition}[Geometric left--right equality on the round vector fibre]
\label{p3:prop:round-fibre-left-right-equality}
The two simple isometry ideals of the single round diagonal vector fibre have
identical fibre-averaged Killing length.  Therefore any Yang--Mills
normalization obtained from the same invariant vector-fibre metric obeys
\[
 \boxed{g_L=g_R}.
\]
This equality is geometric and does not by itself determine the common
absolute coefficient.  It consumes neither \((\mathrm{BULK}_8)\),
\((\mathrm{MATCH}_{8\to7})\), compact-cell saturation, the spectral/EH
coefficient, nor total geodesy; those data enter only stronger reduction or
absolute-normalization claims.
\end{proposition}

\begin{proof}
Inversion of the round group fibre exchanges left and right translations and
is an isometry.  The two simple ideals are therefore related by an isometry of
the same invariant metric and have equal normalized Killing length.
\end{proof}

\begin{theorem}[Seven-dimensional Born-spherical connection reduction]
\label{p3:thm:born-spherical-spectral-action-reduction}
Assume
\((\mathrm{GRED}_7)+(\mathrm{CAL}_{\rm GRED}(\kappa))\), and let the
four-dimensional Einstein coefficient entering the normalized fibre reduction
be
\[
 \mathcal C_{\rm EH}:=\frac{f_2\Lambda_{\rm cut}^2}{48}.
\]
For the declared connection metric on \(X_7\to U_4\), the normalized
scalar-curvature reduction contains
\[
 \int_F\scal[\widetilde g]\,d\mu_F
 =\scal[g_4]+\scal_F
 -\frac14\sum_{s=L,R}\Gamma_{AB}^{(s)}
 F_{s,\mu\nu}^AF_s^{B,\mu\nu}
 +\mathcal R_{\rm other},
\]
where \(\mathcal R_{\rm other}\) denotes the other explicitly retained
geometric terms of the declared \((\mathrm{GRED}_7)\) connection member.
The coefficient of the displayed curvature-square term follows from the
horizontal connection metric and \(\Gamma^{(s)}\) alone.  Total geodesy is not
needed for that coefficient; Lemma~\ref{p3:lem:totally-geodesic-born-odd-fibers}
sets the fibre second-fundamental-form contribution inside
\(\mathcal R_{\rm other}\) to zero on the round graph branch.
Comparison with
\eqref{p3:eq:gred-ym-convention} gives
\[
 \boxed{
 \frac1{g_L^2}=\frac1{g_R^2}
 =\kappa r_{\rm fib}^2\mathcal C_{\rm EH}.}
\]
With \(r_{\rm fib}\Lambda_{\rm cut}=1\),
\[
 \boxed{\frac1{g_L^2}=\frac1{g_R^2}=\frac{\kappa f_2}{48}},
 \qquad \alpha_L^{-1}=\alpha_R^{-1}=\frac{\kappa\pi f_2}{12}.
\]
The equality \(g_L=g_R\) is geometric; the absolute value is conditional on
the displayed reduction conventions.  On that compact-cell saturation branch, in particular,
\[
 \kappa=1:\quad g^{-2}=\frac{f_2}{48},
 \qquad
 \kappa=\frac14:\quad g^{-2}=\frac{f_2}{192}.
\]
The first is the period-\(2\pi\) quaternion/integral-generator calibration;
the second is the Paper-II spin-\(\tfrac12\) generator calibration.  No physical
\(U(1)_{\rm fib}\) coefficient is claimed here.  Neither
\((\mathrm{BULK}_8)\) nor its factor \(1/2\) is a premise of this theorem.
\end{theorem}

\begin{proof}
For the connection metric of
Definition~\ref{p3:def:gred-connection-generator-normalization}, the
horizontal-curvature part of the reduced Einstein term is
\[
 -\frac{\mathcal C_{\rm EH}}4
 \int_U\sqrt{-g}\,\Gamma_{AB}^{(s)}
 F_{s,\mu\nu}^AF_s^{B,\mu\nu}\,d^4x
\]
for each ideal, with the same sign/contraction convention as
\eqref{p3:eq:gred-ym-convention}; see
\cite{Baptista2024InternalSymmetries}.  Comparing coefficients and using
\eqref{p3:eq:gred-generator-metric} gives
\(g_s^{-2}=\kappa r_{\rm fib}^2\mathcal C_{\rm EH}\).  Total geodesy removes
the O'Neill \(T\)-tensor but does not determine this horizontal-curvature
coefficient.  Inversion gives the equality of the two averaged generator
metrics.  Compact-cell saturation then gives the displayed \(\kappa f_2/48\)
value.
\end{proof}

\begin{corollary}[Calibrated absolute left--right coefficient and compact-cell condition]
\label{p3:cor:no-independent-lr-gauge-coefficient}
Inside \((\mathrm{GRED}_7)+(\mathrm{CAL}_{\rm GRED}(\kappa))\),
\[
 \boxed{\frac1{g_L^2}=\frac1{g_R^2}
 =\frac{\kappa f_2}{48}(r_{\rm fib}\Lambda_{\rm cut})^2.}
\]
On the \emph{additional} compact-cell saturation branch
\(r_{\rm fib}\Lambda_{\rm cut}=1\), this specializes to
\[
 \boxed{\frac1{g_L^2}=\frac1{g_R^2}=\frac{\kappa f_2}{48}.}
\]
Thus \(f_2/48\) requires both \(\kappa=1\) and compact-cell saturation.
The generator calibration \((\mathrm{CAL}_{\rm GRED}(\kappa))\) fixes
\(\kappa\) but does \emph{not} imply saturation of the dimensionless product
\(r_{\rm fib}\Lambda_{\rm cut}\).
The geometric equality of the left--right coefficients does not require
saturation. Changing their common absolute value can reflect generator
calibration, fibre radius, or spectral cutoff; changing their \emph{relative}
value leaves the symmetric round-fibre branch or adds a left--right
distinguishing datum. The detailed negative case analysis is deferred to
Paper~III-G, Appendix~\ref{p3f:app:lr-diagnostic}.
\end{corollary}

\begin{corollary}[Separation of the connection reduction and bulk diagnostic]
\label{p3:cor:born-spectral-action-closure-status}
Under \((\mathrm{GRED}_7)+(\mathrm{CAL}_{\rm GRED}(\kappa))\), the
seven-dimensional connection reduction supplies the conditional left--right
kinetic normalization.  Independently, \((\mathrm{BULK}_8)\) supplies the
ordinary eight-dimensional factor
\(\beta_{\rm B}^{\rm bulk}=1/2\).  If that diagnostic is used to normalize the
seven-dimensional coefficient, \((\mathrm{MATCH}_{8\to7})\) is additionally
required.  The physical scalar rank
\(N_{\rm sc}^{\rm rank}=1\) and corner weight
\(w_{\rm sc}^{\rm corner}=1/8\) remain separate typed quantities and are not
multiplied into the bulk factor.  Fixed-set terms and later running/matching
remain separate.
\end{corollary}

\begin{remark}[Optional scalar-circle metric branch]
The constitutional scalar-circle role is metric-free.  A physical
\(U(1)_{\rm fib}\) kinetic coefficient therefore requires an additional
scalar-circle metric/connection realization, for example a radius
\(\rho_{\rm fib}>0\), rather than following from the existing role alone.
This is a downstream selector on the scalar-circle realization family.  The
corresponding optional formula and its assumptions are deferred to
Paper~III-G, Appendix~\ref{p3f:app:scalar-circle-metric}.
\end{remark}

\begin{corollary}[Conditional Euclidean quasi-locality]
\label{p3f:cor:euclidean-quasi-locality}
Assume \((\mathrm{QLOC})\) together with all hypotheses of
Theorem~\ref{p3:thm:local_spectral_asymptotics}.  Let \(U\) be a
reciprocal-position chart in the semiclassical derivative-expansion regime.
On \(U_0\), the compressed observables obey the quantified separated-support
estimate, \(TU\oplus0\) is involutive, and the Euclidean effective action has
the displayed local two-derivative leading term together with its explicitly
retained higher-derivative remainder.  No Lorentzian dynamics follows from
this corollary.
\end{corollary}

\begin{proof}
The local reciprocal atlas supplies \(T^*U\), the exact doubled bundle
\(E_U=TU\oplus T^*U\), and the involutive position summand
\(F=TU\oplus0\).  The anchored connection restricts to the Levi--Civita
connection on \(F\), while Theorem~\ref{p3:thm:local_spectral_asymptotics}
provides the Euclidean derivative expansion in the stated protected Einstein
branch on \(U_0\).  Under \((\mathrm{QLOC})\),
equation~\eqref{p3:eq:quantified-qloc} gives rapid operator-norm suppression of
commutators of the declared separated coarse observables at distances large
compared with the declared coarse scale.  This proves precisely the stated
Euclidean conclusion.
\end{proof}

\begin{corollary}[Conditional Lorentzian macrolocality]
\label{p3:cor:d3_derived}
Assume the hypotheses of
Corollary~\ref{p3f:cor:euclidean-quasi-locality}, together with either an
occurring Lorentzian action \((\mathrm{LOR}_{\rm dir})\) or a fixed classical
continuation \((\mathrm{CONT}_{\rm cl})\) that preserves the field identity,
reality/sign convention, and two-derivative coefficient.  Assume also, on a
stated normed field/domain class \(\mathcal X\),
\[
 (\mathrm{IR}_{2\partial}(\varepsilon_*)):\qquad
 \|\mathcal E^{\ge4}\|_{\mathcal X}
 \le\varepsilon_*\|\mathcal E^{\rm EH}\|_{\mathcal X}.
\]
Then \(U_0\) supports a macroscopically local Lorentzian infrared sector in
the sense of Definition~\ref{p3:def:macrolocal_sector}, with the displayed
quantitative remainder.  Any homogeneous cosmological realization remains
downstream and is not a premise.
\end{corollary}

\begin{proof}
The preceding corollary supplies QLOC and involutivity.  The declared
Lorentzian occurrence/continuation supplies condition (iii) of
Definition~\ref{p3:def:macrolocal_sector} without changing the fields or the
two-derivative coefficient, and the norm inequality supplies the claimed
control of the retained higher-derivative variation.  These four antecedents
prove the result.  Without the Lorentzian bridge the conclusion stops at
Euclidean quasi-locality; without the norm bound no ``suppressed'' remainder
claim is licensed.
\end{proof}

\begin{remark}[Forward realization interface deferred]
The local package above may be instantiated by a later theorem that proves the
required hypotheses on a concrete homogeneous or other macroscopic branch.
No such downstream realization is used as a Paper~III-F premise.  The generic
composition statement is recorded in
Paper~III-G, Appendix~\ref{p3f:app:forward-realization}.
\end{remark}


\label{p3:part:summary-export-ledger}


\section[Native chirality, boundary source, and complete-action parity tests]{Native chirality, boundary source, and\protect\linebreak complete-action parity tests}
\label{p3f:sec:v353-native-boundary-tests}
This section is a \emph{discriminator} for the conditional cap in
Paper~III-E, Section~\ref{p3e:sec:v353-schur-boundary}; it does not
identify the existing raw heat trace with a stationary Gaussian cap.
Write $\Gamma=\gamma_g$ on a four-component comparison spinor space.

\begin{proposition}[Transverse projector and inverse-square identity]
\label{p3f:prop:v353-projector-transversality}
Let $P$ be a rank-$r$ orthogonal projector, $M=M^\dagger$ an
involution $M^2=I$, and $L_u=(I+uM)^{-1}$ for real $|u|<1$.
Then
\begin{equation}
\label{p3f:eq:v353-projector-iff}
 \operatorname{Tr}(PL_uPL_u)=\frac{r}{(1-u^2)^2}
 \quad\text{for every }u\text{ in an open interval}
 \quad\Longleftrightarrow\quad PMP=0.
\end{equation}
For the conditional canonical non-null Dirac boundary of
Remark~\ref{p3f:rem:v175-canonical-dirac-boundary-completion},
$\chi=\gamma_n\gamma_g$, $P=\Pi_+=(I+\chi)/2$ and $M=\gamma_g$
satisfy $PMP=0$ on its specified Clifford convention.  Therefore
with $J=R_sP/\sqrt2$, $r=2$ and
$u=cR/[R_s(1+\sqrt{1+(R/R_s)^2})]$ one has the exact
\emph{comparison} identity
\[
 -\frac18\operatorname{Tr}(JL_uJL_u)
 =-\frac{R_s^2}{8(1-u^2)^2}
 =-\frac{\sigma^2}{8B(R)^2}.
\]
This identity alone does not show that the native spectral operator
produces $L_u$ or $J$, nor that a physical non-null boundary occurs.
\end{proposition}
\begin{proof}
Since $L_u=(I-uM)/(1-u^2)$,
$PL_uP=(P-uPMP)/(1-u^2)$ and
\[
 \operatorname{Tr}(PL_uPL_u)=
 \frac{r-2u\operatorname{Tr}(PMP)
      +u^2\operatorname{Tr}[(PMP)^2]}{(1-u^2)^2}.
\]
Equality on an interval forces the polynomial coefficients to vanish.
Self-adjointness gives $\operatorname{Tr}[(PMP)^2]=0$ precisely when
$PMP=0$.  For the canonical boundary $\{\chi,\gamma_g\}=0$.
Insert $\operatorname{Tr}P=2$ and $\sigma=B_IR_s$.
\end{proof}

\begin{theorem}[Interior parity and its inherited reduction obstruction]
\label{p3f:thm:v353-interior-parity}
On an ordinary compact interior Dirac heat domain assume
$\{D_0,\Gamma\}=0$, $[V,\Gamma]=0$, and
$D(\phi)=D_0+\phi\Gamma V$, with common, conjugation-invariant
operator domain, trace and localizer.  Then
\[
 \Gamma D(\phi)\Gamma=-D(-\phi),\qquad
 S_{\rm heat}(R,\phi)=S_{\rm heat}(R,-\phi).
\]
If the curvature variation transports this grading equivariantly,
$\partial_R\partial_\phi S_{\rm heat}(R,0)=0$.
Moreover, any stationary auxiliary elimination or positive auxiliary
integration respecting the same involution on fields and measure
preserves this evenness.  An unmarked odd scalar source cannot thereby
produce a nonzero mixed curvature/odd-source first jet.
The claim does not apply to a boundary domain that fails the stipulated
conjugation invariance or to a source not identified with $\phi$.
\end{theorem}
\begin{proof}
Conjugate the operator, square and take the invariant heat trace.
For $S(R,\phi,y)=S(R,-\phi,Cy)$ with $C^2=I$ and invariant
admissible $y$ domain, change variables $y\mapsto Cy$ in the
minimum (or invariant integral).  The reduced action remains even.
Differentiate at $\phi=0$.
\end{proof}


\begin{corollary}[Signed multiplicity selection rule for a complete finite heat two-jet]
\label{p3f:cor:v3030-signed-heat-twojet}
Consider a finite regular metric/source comparison
\(\mathcal H(h,\lambda)\) on \(E\otimes\mathbb C^2\), with
\(Z=\operatorname{diag}(1,-1)\) in the signed multiplicity,
and let \(S(h,\lambda)=\operatorname{Tr}\Phi(\mathcal H(h,\lambda))\),
where \(\Phi\) is differentiable on the finite real spectrum.
Assume, in the \emph{same fixed source comparison}, that
\[
 \mathcal H(h,0)=A(h)\otimes I_2,
 \qquad
 \left.\partial_\lambda\mathcal H(h,\lambda)\right|_{0}
 =B(h)\otimes Z
\]
for every \(h\) near the comparison value; both \(A\) and \(B\)
are differentiable, and the trace, regulator, source measure, and
metric/source contacts preserve this factorization.
Then the \emph{complete} finite heat response obeys
\[
 \boxed{\left.\partial_\lambda S(h,\lambda)\right|_0=0,
 \qquad
 \left.\partial_h\partial_\lambda S(h,\lambda)\right|_0=0.}
\]
The second identity includes the contribution from
\(\partial_hB(h)\otimes Z\); it is not only a direct-Hessian
cancellation.  A physically distinct signed receiver or
source-state marker may violate the hypotheses.  The rooted
imaginary coherence of Corollary~
\ref{p1:cor:v3030-rooted-imaginary-coherence} supplies a
\emph{candidate marker}, not the required native metric/source
incidence. Neither a phase-to-curvature transport nor a physical
dark-sector response is inferred.
\end{corollary}
\begin{proof}
For a finite Hermitian \(\mathcal H\), the trace spectral
functional has the exact first variation
\(\partial_\lambda\operatorname{Tr}\Phi(\mathcal H)
 =\operatorname{Tr}[\Phi'(\mathcal H)\partial_\lambda\mathcal H]\),
including noncommuting \(A,B\).  At \(\lambda=0\) the right-hand
side factorizes as
\(\operatorname{Tr}_E[\Phi'(A(h))B(h)]
 \operatorname{Tr}_{\mathbb C^2}Z=0\).
This vanishes for \emph{every} admitted \(h\), so differentiation
with respect to \(h\) proves the complete mixed-jet identity.
The latter differentiation also differentiates \(B(h)\), hence
retains the mixed metric/source contact.  A marked metric or measure
outside the stated tensor-factorization hypothesis is not covered.
\end{proof}


\begin{proposition}[Exact cap chiral mixed-jet discriminator]
\label{p3f:prop:v353-chiral-mixed-jet}
In the purely finite comparison representation set
$\Gamma=\mathrm{diag}(I_2,-I_2)$,
$\chi=\bigl(\begin{smallmatrix}0&I_2\\I_2&0\end{smallmatrix}\bigr)$,
$P=(I+\chi)/2$, $\rho=R_s/2$, $J=\sqrt2\rho P$ and
\[
 F(R,q)=-\frac18\operatorname{Tr}
 \bigl[(J+q\Gamma)L_u(J+q\Gamma)L_u\bigr].
\]
Then the exact finite polynomial in the source is
\begin{equation}
\label{p3f:eq:v353-chiral-return}
 F(R,q)=-\frac{\rho^2-2\sqrt2\rho u q+(1+u^2)q^2}
                      {2(1-u^2)^2},
\end{equation}
and hence
\begin{equation}
\label{p3f:eq:v353-mixed-jet}
 \left.\partial_qF\right|_{q=0}=\frac{\sqrt2\rho u}{(1-u^2)^2},
 \qquad
 \left.\partial_R\partial_qF\right|_{R=q=0}
 =\frac{c}{2\sqrt2}.
\end{equation}
For $B_I>B_U$ this mixed jet is nonzero.  Therefore a direct
identification of $q\Gamma$ with the parity-odd source $\phi\Gamma V$
on the domain of Theorem~\ref{p3f:thm:v353-interior-parity}
would be inconsistent.  Other licensed physical source maps must
be compared by their own complete source transport, not ruled out.
\end{proposition}
\begin{proof}
Use the $2+2$ chiral blocks, $L_u=\mathrm{diag}((1+u)^{-1}I_2,
(1-u)^{-1}I_2)$ and $P=(I+\chi)/2$.  Taking the trace gives
Eq.~\eqref{p3f:eq:v353-chiral-return}.
Differentiate, use $u(0)=0$ and $u'(0)=c/(2R_s)$.
\end{proof}

\begin{remark}[Physical boundary and finite grading]
\label{p3f:rem:v353-grading-frontier}
The canonical Robin endomorphism $S=-K_E\Pi_+/2$ in
Remark~\ref{p3f:rem:v175-canonical-dirac-boundary-completion}
is diagonal in the $\Pi_+$ grading and does not \emph{alone}
identify the transverse resolvent $L_u$ or its absolute source norm.
The finite internal grading used in
Proposition~\ref{p3f:prop:v250-zero-scalar-member-fibre}
can preserve evenness of a physical odd scalar source even on a
boundary that preserves that internal grading.  Conversely a boundary
may fail the \emph{spin} grading hypothesis, so interior parity is not
an unconditional assertion about all boundary classes.
On a genuine standard Dirichlet-to-Neumann problem the response also
carries tangential mode dependence: for the flat half-space
$(-\partial_n^2-\Delta_{\parallel}+m^2)\psi=0$ the decaying
symbol is $\sqrt{|\xi|^2+m^2}$, not a curvature-only number.
The cap's two-channel $S_+(R)$ requires a derived finite reduction.
\end{remark}

\begin{proposition}[Observable-relative quadratic source identity]
\label{p3f:prop:v353-observable-quadratic}
For two real stationary source models on a \emph{common physically licensed}
real probe space $q$, allow distinct real auxiliary dimensions and put
\[
 F_a(q)=\ell_a-(j_a+Q_aq)^\mathsf T T_a^{-1}(j_a+Q_aq),
 \quad T_a=T_a^\mathsf T>0,\qquad a=1,2.
\]
The two reduced quadratic polynomials coincide on an open real
probe set if and only if the three reduced quantities coincide:
\[
 \ell_a-j_a^\mathsf T T_a^{-1}j_a,\qquad
 Q_a^\mathsf T T_a^{-1}j_a,\qquad
 Q_a^\mathsf T T_a^{-1}Q_a.
\]
Complex auxiliary systems can be compared after a fixed, actual
realification of their physical source variables and Hermitian forms.  This criterion certifies \emph{physical quadratic response}
without assuming all auxiliary directions are observable; it does
not certify their complete parent identity or higher nonlinear jets.
Quantum-integrated spectators additionally require their measure
and determinant terms.
\end{proposition}
\begin{proof}
Expand each quadratic polynomial into constant, linear and quadratic
coefficients.  Equality on an open real probe domain is equivalent
to equality of these coefficients; no invertibility of $Q_a$ is used.
\end{proof}


\section{Charged and neutral sources and finite internal boundaries}
\label{paper3f:sec:organized-sources}
\subsection{Directed source returns and the finite assembly handoff}
\label{p3f:sec:v1126-directed-source-handoff}
For a directed finite member $A:G_R\to G_L$, dagger completion in the
generated chiral owner means
\[
 D[A]=\begin{pmatrix}0&A\\A^\dagger&0\end{pmatrix},\qquad
 D[A]^2\big|_{G_L}=AA^\dagger.
\]
The physical left-return quotient is therefore the positive Gram $AA^\dagger$.
The family expression $A+A^\dagger$, when source and target coordinates
have additionally been identified, is a different forward response and
needs its own origin certificate. Paper~III-G, Appendix~\ref{p3f:app:v1126-dagger-typing}
gives the typing calculation.

The rooted generated-member handoff
(Corollary~\ref{p3f:cor:v1122-root-generated-fermion-member-handoff}) continues
to consume independently generated finite action members. Paper V-A v3.93
discharges the older primitive-source antecedents on the rooted branch;
Paper V-B v5.85 supplies canonical terminal supports and a conditional
finite spectral reconstruction. A response of the form
$q_s(A_sA_s^\dagger)A_s$ enters this handoff only when its coherent parent,
terminal-role matching, and response weighting are required outputs of a
licensed common source constructor. Polynomial interpolation alone is not
such a constructor and is not an inserted kinetic normalization.
\texttt{FERMION\_SAME\_ACTION\_SVD\_ASSEMBLY} remains
\texttt{OPEN\_PHYSICAL}; primitive \texttt{SrcVal} is not reopened.


\subsection{Source covectors and moving-source boundary response}
\label{p3f:sec:v2109-source-response-transport}
The native terminal deep \(S+T\) law is retained.  The following positive
calculations supply source and response descendants on their declared
quadratic or regular stationary domains; they do not identify the complete
global galaxy response with a chosen scalar interpolation.

\begin{proposition}[Hilbert source pulled back to the lapse]
\label{p3f:prop:v2109-hilbert-lapse-source}
Use \(c=1\), zero shift and
\(ds^2=-N^2dt^2+h_{ij}dx^idx^j\), with fixed matter/source variables in
the metric variation.  In the convention
\(\delta S_m=\tfrac12\int\sqrt{-g}\,T^{\mu\nu}\delta g_{\mu\nu}\,d^4x\),
let \(\rho=T_{\mu\nu}n^\mu n^\nu\), \(n^\mu=N^{-1}\delta^\mu_0\).
Then
\[
 \delta_NS_m=-\int dt\,d^3x\,\sqrt h\,\rho\,\delta N.
\]
Thus the source covector is an existing Hilbert variation, not a new
gravitational-charge coefficient.
\end{proposition}
\begin{proof}
\(\delta g_{00}=-2N\delta N\), \(\sqrt{-g}=N\sqrt h\), and
\(T^{00}=\rho/N^2\).  Substitution gives the displayed result.
\end{proof}

\begin{theorem}[Adjoint Poisson projection of a quadratic bulk source]
\label{p3f:thm:v2109-poisson-source-projection}
For \(k>0\), consider the half-line quadratic form
\[
 S[f;j]=\frac12\int_0^\infty(f'^2+k^2f^2)\,d\xi
       -\int_0^\infty jf\,d\xi,
 \qquad f(0)=f_0,
\]
on a decay/finite-energy domain on which the Dirichlet inverse \(G_D\)
and displayed source pairings exist.  Then
\[
 f(\xi)=e^{-k\xi}f_0+(G_Dj)(\xi),\qquad
 j_\partial=\int_0^\infty e^{-k\xi}j(\xi)\,d\xi,
\]
\[
 S_\partial(f_0;j)=\frac12kf_0^2-j_\partial f_0
                    -\frac12\langle j,G_Dj\rangle,
 \qquad \pi_0=kf_0-j_\partial.
\]
\end{theorem}
\begin{proof}
The Poisson extension solves the homogeneous Euler equation.  Substitute
\(f=e^{-k\xi}f_0+u\) with \(u(0)=0\).  Its quadratic cross term is zero
by integration by parts; minimizing the interior form gives \(u=G_Dj\).
The remaining cross term is the adjoint Poisson source pairing, and
variation in \(f_0\) gives the outward flux \(\pi_0\).
\end{proof}
For linear response, \(j\) is the incremental forcing and the quadratic
operator must contain all variations consumed by that response.  A source
already included in a source-conditioned operator is not added twice.
Neither \(j_\partial=g_b/a_0\) nor exhaustion of the full baryonic source
by one scalar follows from this projection.

\begin{theorem}[Moving-source Schur derivatives and exact residual]
\label{p3f:thm:v2109-moving-source-schur}
Let a twice differentiable self-adjoint block Hessian on fixed retained
spaces be
\[
 H(y)=\begin{pmatrix}A&B^\dagger\\ B&C\end{pmatrix},\quad
 C^{-1}\ \hbox{exists},\quad
 Z=-C^{-1}B,\quad L=\binom I Z,\quad K=A-B^\dagger C^{-1}B.
\]
Derivatives include the actual changing stationary background.  With
\(W=\dot B+\dot C Z\),
\[
 \dot K=L^\dagger\dot HL,\qquad
 \ddot K=L^\dagger\ddot HL-2W^\dagger C^{-1}W.
\]
At a regular base value \(y_0\), put \(Z_0=Z(y_0)\),
\(L_0=\binom I{Z_0}\) and \(\mathcal R(y)=B(y)+C(y)Z_0\).  Exactly,
\begin{equation}
 \label{p3f:eq:v2109-exact-schur-residual}
 K(y)=L_0^\dagger H(y)L_0-\mathcal R(y)^\dagger C(y)^{-1}\mathcal R(y).
\end{equation}
Consequently the complete identity \(K(y)=yG\), for a fixed positive
metric \(G\), is equivalent to
\(L_0^\dagger H(y)L_0-yG=\mathcal R^\dagger C^{-1}\mathcal R\).
\end{theorem}
\begin{proof}
\(HL=(K,0)^T\), and the upper block of \(\dot L\) is zero, so
 differentiating \(L^\dagger HL\) gives the first formula.
Differentiating \(CZ+B=0\) gives \(\dot Z=-C^{-1}W\).
The two lift terms in the second derivative are both
\(-W^\dagger C^{-1}W\).  Completing the square at the fixed lift \(L_0\)
proves \eqref{p3f:eq:v2109-exact-schur-residual}.
\end{proof}
Even for \(C>0\), the term \(L^\dagger\ddot HL\) has no prescribed
sign.  A first-jet match does not prove the full source law.  The physical
port and its kinetic metric must also be transported, as recorded in
Paper~III-G, Appendix~\ref{p3f:app:v2109-quadratic-boundary-details}.

\begin{theorem}[Terminal scalaron decoupling with regular retained mixing]
\label{p3f:thm:v2109-mixed-scalaron-decoupling}
On a fixed or uniformly controlled normed block domain, suppose
\[
 \mathcal H_\chi=\begin{pmatrix}
 A&E^\dagger&B^\dagger\\E&N&D^\dagger\\B&D&C_\chi
 \end{pmatrix},\qquad \|C_\chi^{-1}\|=O(\chi),
\]
all mixed blocks are uniformly bounded, and \(N\) is uniformly invertible
on the reduced quotient.  The full boundary Schur return is
\[
 K_\partial=A-E^\dagger N^{-1}E+O(\chi).
\]
Using the inherited terminal expansion
\(\chi(C_7+\delta)=2\delta/[C_7(U+C_7)]+O(\delta^2)\), the error is
\(O(\delta)\) on this same controlled domain.
\end{theorem}
\begin{proof}
Taking the same-action Schur complement of the scalaron block changes \(A,E,N\) to
\(A-B^\dagger C_\chi^{-1}B\),
\(E-D^\dagger C_\chi^{-1}B\),
\(N-D^\dagger C_\chi^{-1}D\), respectively.  Each change is \(O(\chi)\).
Uniform invertibility and the inverse identity give
\((N-D^\dagger C_\chi^{-1}D)^{-1}=N^{-1}+O(\chi)\).
The remaining Schur elimination proves the result.
\end{proof}
The non-scalaron term \(E^\dagger N^{-1}E\) remains.  This extends the
scalaron block theorem, not complete G28 or uncontrolled continuum
uniformity.


\subsection{Complete normal-current return to a neutral pairing port}
\label{p3f:sec:v2167-cpbl2-current-return}
The following statements specialize the complete action and metric/source
constructors to a licensed quadratic normal branch. They do not select a
K16 Hamiltonian. All earlier eliminated sectors, kinetic metrics,
connections and consumed source jets belong to the same complete parent.

\begin{proposition}[Three-block complete neutral return]
\label{p3f:prop:v2167-complete-neutral-return}
On the canonical carrier $\mathcal K_R\oplus\mathcal K_L\oplus\mathcal K_E$
write
\[
 h=\begin{pmatrix}h_R&T^\dagger&S^\dagger\\
 T&h_L&W\\S&W^\dagger&h_E\end{pmatrix},
 \quad T:R\to L,\quad S:R\to E,\quad W:E\to L.
\]
For a common regular spectral argument $\zeta$, set
$G_E=(\zeta I-h_E)^{-1}$ and
$G_L=(\zeta I-h_L-WG_EW^\dagger)^{-1}$. The complete retained kernel is
\[
 \boxed{K_R=\zeta I-h_R-S^\dagger G_ES
 -(T^\dagger+S^\dagger G_EW^\dagger)
 G_L(T+WG_ES).}
\]
Equivalently, with $\mathsf B=(T^\dagger,S^\dagger)$,
$h_A=\begin{pmatrix}h_L&W\\W^\dagger&h_E\end{pmatrix}$ and
$G_A=(\zeta I-h_A)^{-1}$,
\[
 \Sigma_R=\mathsf B G_A\mathsf B^\dagger,\qquad
 K_R=\zeta I-h_R-\Sigma_R,\qquad
 K_R^{-1}=[(\zeta I-h)^{-1}]_{RR}.
\]
\end{proposition}
\begin{proof}
Eliminate $E$, retaining both direct $S$ blocks and their mixed couplings
to $L$, then eliminate $L$. Direct elimination of the complete rest block
$A=L\oplus E$ gives the second formula and its resolvent identity.
Corollary~\ref{p3f:cor:complete-reduction-reuse} preserves these equalities
and all consumed jets on their common uniquely generated regular branch.
\end{proof}

\begin{corollary}[Actual charged-source return]
\label{p3f:cor:v2167-charged-source-return}
For an actual charged source $X_E=\delta h_E$ with $\mathsf B$ fixed,
\[
 \delta\Sigma_R=\mathsf B G_Aj_EX_Ej_E^\dagger G_A\mathsf B^\dagger.
\]
Both resolvents have the same $\zeta$. On a real regular slice put
$F_E=j_E^\dagger G_A\mathsf B^\dagger$. Then
\[
 \delta\Sigma_R=F_E^\dagger X_EF_E,\qquad
 F_E^\dagger P_aF_E
 =|F_E^\dagger g_a\rangle\langle F_E^\dagger g_a|.
\]
The returned ray is nonzero exactly when $F_E^\dagger g_a\ne0$.
If the source moves the coupling as well, its complete variation is
\[
 \delta\Sigma_R=(\delta\mathsf B)G_A\mathsf B^\dagger
 +\mathsf B G_A(\delta\mathsf B)^\dagger
 +\mathsf B G_A(\delta h_A)G_A\mathsf B^\dagger.
\]
\end{corollary}
\begin{proof}
Differentiate the complete inverse using $\delta G_A=G_A(\delta h_A)G_A$.
On the real slice $G_A^\dagger=G_A$; substituting $P_a=g_ag_a^\dagger$
gives the outer product. A complex retarded resolvent does not obey this
Hermitian replacement. All actual embedding, metric and connection
derivatives remain included in the complete source member.
\end{proof}
Returned rays are a positive packet, not generally an orthogonal projector
triple. On the $S=0$ chain branch, a bounded finite large-$|\zeta|$ expansion
gives
\[
 \delta\Sigma_R=\zeta^{-4}T^\dagger WX_EW^\dagger T+O(\zeta^{-5}).
\]
With $S\ne0$, this is the coefficient of the specifically marked
five-slot word; the complete response also contains direct and mixed
returns. A role marker does not itself license a background potential.

\begin{proposition}[Causal return retains the inherited rest state]
\label{p3f:prop:v2167-causal-inherited-return}
For the clocked quadratic equations with actual rest propagator $V_A$,
\[
 \begin{aligned}
 i\hbar\dot\psi_R(t)={}&h_R(t)\psi_R(t)
 +\mathsf B(t)V_A(t,t_0)\psi_A(t_0)\\
 &-\frac i\hbar\int_{t_0}^t
 \mathsf B(t)V_A(t,s)\mathsf B(s)^\dagger\psi_R(s)\,ds.
 \end{aligned}
\]
The neutral return of the full normal propagator,
$V_R(t,s)=J_R(t)^\dagger V(t,s)J_R(s)$, need not be unitary.
\end{proposition}
\begin{proof}
Variation of constants in the rest equation followed by substitution
gives the identity. The full $V$ is unitary, but its returned subblock can
have outflow. The initial-rest term retains inherited covariances even
when its field first moments vanish.
\end{proof}
The charge law uses total lepton charge on the complete retained carrier.
If the normal owner has quartic or quantum boson--fermion interactions,
its complete many-body response replaces the quadratic specialization.

\begin{corollary}[Family-equivariant canonical retention]
\label{p3f:cor:v2167-family-equivariant-retention}
Suppose the actual full-rank field chart intertwines every family matrix
unit. In compatible fixed module frames $J=J_m\otimes I_G$.
On the uninserted kinetic branch $G_p=G_m\otimes I_G$, let
$G=J^\dagger G_pJ$ and
$\widetilde Z=G_p^{1/2}JG^{-1/2}$. Then
\[
 \widetilde Z^\dagger\widetilde Z=I,\qquad
 \widetilde Z^\dagger(I\otimes A)\widetilde Z=I\otimes A
 \quad(A\in M_3(\mathbb C)).
\]
The metric normalization and moving connection retain the identity family
factor on this equivariant domain.
\end{corollary}
\begin{proof}
Intertwining diagonal units preserves family components; intertwining
off-diagonal units equates their multiplicity maps. Thus $J$ factors.
$G$ and its positive root factor as well, proving both identities.
Differentiation in these fixed module frames preserves factorization of
the connection in Theorem~\ref{p3f:thm:v1148-finite-metric-constructor}.
\end{proof}
For a canonical compression $\pi_Z(A)=Z^\dagger AZ$, multiplication has
defect $-Z^\dagger A(I-ZZ^\dagger)BZ$. Hence preservation of a complete
star algebra requires its reducing-range guard; the short positive
argument and exact diagnostic are in
Lemma~\ref{p1b:lem:v2167-compression-multiplication}.
Faithful pair insertion alone does not prove this guard, even on the same
canonical chart (Proposition~\ref{p1b:prop:v2167-faithful-pair-lossy-current}).
The actual K16 constructor must derive its required intertwining or
evaluate its complete nonzero mixed return. Numerical mode functions are
not a prerequisite for a structural proof, and this conditional result
does not supply that still-unproved handoff as a new input.


\subsection{Exact charged-source return and the free clock comparison}
\label{p3f:sec:v2168-charged-free-pole}
The positive terminal coefficients, already generated by the direct
Clifford/split compiler, are fixed before the source response is calculated:
\[
 M=\diag(\varepsilon_e,\varepsilon_\mu,\varepsilon_\tau)>0,\qquad
 \begin{aligned}
 \varepsilon_e&=E_{40}(1+\sqrt5)^{-1}(24/25)^{-3/2},\\
 \varepsilon_\mu&=(E_{36}/2)(24/25),\\
 \varepsilon_\tau&=E_{36}(2\sqrt6+\sqrt{12})(24/25).
 \end{aligned}
\]
The occurring source is $R_{\rm act}=\sum_r\varepsilon_r\iota_r$,
with $\iota_r=\Phi_{\rm HG}(A_r^{\rm hier})$. Its physical occurrence and
assembly are closed by V-A; no independent source choice is introduced.
The current frame $U=(g_e,g_\mu,c)$ is the generated incidence frame,
not a spectrum-generated mass frame. The primitive momentum level remains
$m_p=1$. A local rest-comparison symbol requires its admitted physical
dictionary; it creates neither a primitive zero-momentum state nor a
position lattice.

On the admitted local uninserted clock branch III-F gives
\[
 (Z_L,D_e,Z_R)=c_F(I,R_{\rm act},I),
\]
so canonical normalization leaves $R_{\rm act}$ unchanged. The common $c_F$ is not a missing scale. In this
section write $R=R_{\rm act}$ for the exact criterion, and use the
explicit $T_*=UMV^\dagger$ only for the stated comparison construction,
with $V=(h_e,c,h_\tau)$ as in V-B,
Theorem~\ref{p6:thm:v2168-right-completion}.
\subsubsection{The exact test, without an unnecessary right-frame choice}
For the canonically normalized left-current sources set
$J=\binom{U}{0}$ and let $G_U(z)=J^\dagger(zI-h_R)^{-1}J$.

\begin{theorem}[Left-return characterization of the free tagged clock claim]
\label{p3f:thm:v2168-left-return-test}
The following statements are equivalent:
\begin{enumerate}
\item The complete free left-current two-point function is
\[
 G_U(z)=\diag_r\frac{z}{z^2-\varepsilon_r^2}.
\]
\item The actual canonically normalized family source satisfies
\[
 \boxed{RR^\dagger=UM^2U^\dagger.}
\]
\item With the already fixed $M$, the comparison readout
$K:=M^{-1}U^\dagger R$ is unitary.
\item There exists a unitary $V_R$, obtained as a readout of $R$, such that
$R=UMV_R^\dagger$.
\end{enumerate}
Under these conditions every tagged left-current channel has poles
$\pm\varepsilon_r$ with residue $1/2$. In the same admitted free local
Clifford chart its dispersion is
$E_r(\mathbf p)^2=c_{\rm op}^2|\mathbf p|^2+\varepsilon_r^2$.
\end{theorem}
\begin{proof}
Since $h_R^2=\diag(RR^\dagger,R^\dagger R)$, multiplication by
$zI+h_R$ gives
\[
 G_U(z)=zU^\dagger(z^2I-RR^\dagger)^{-1}U.
\]
Consequently (2) implies (1). Conversely, equality of these meromorphic
matrix functions implies equality of their $z^{-3}$ coefficients at infinity,
which are $U^\dagger RR^\dagger U$ and $M^2$. This gives (2).
Direct substitution gives
$KK^\dagger=M^{-1}U^\dagger RR^\dagger UM^{-1}$.
As $K$ is square, $KK^\dagger=I$ is equivalent to unitarity. Take
$V_R=K^\dagger$ for (4); the converse follows by multiplication.
The two simple partial fractions of each diagonal entry have residue $1/2$.
If energies coincide, retain their combined spectral projection.
The dispersion follows by squaring the retained Clifford symbol, as in
the geometric completion; the spin and family factors commute.
\end{proof}

\subparagraph{Minimal free-response certificate.}
The equation $D_e=c_FUMV^\dagger$ for V-B's explicit Hopf completion
is sufficient. It is stronger than the free left-current/rest-pole claim
requires. The exact next test is
$RR^\dagger=UM^2U^\dagger$. This removes an unnecessary demand to prove one
particular right frame. It does not authorize changing the actual source or
its right-current and interaction incidences. A right-unitary change is an
equivalent full presentation only when every consumed field, source, kinetic
and interaction block is transformed together. No right frame is chosen as
new physical input here.

The theorem characterizes the complete tagged free two-point function.
It is stronger than merely asking whether the unordered pole set contains
three given values with some nonzero overlap. It is not a necessity theorem
for an interacting, shifted or differently typed physical observable.

\subsubsection{The hierarchy removes four ordered cross returns exactly}
The retained charged directions give, in their common rooted family
presentation,
\[
 \iota_e\in Q_\perp\operatorname{End}(G)Q_\perp,\quad
 \iota_\mu\in Q_\perp\operatorname{End}(G)Q_0,\quad
 \iota_\tau\in Q_0\operatorname{End}(G)Q_\perp.
\]
These are statements about the source packet's hierarchy supports.
They do not identify the raw source arrows with normalized current arrows.

\begin{theorem}[Exact surviving native cross return]
\label{p3f:thm:v2168-hierarchy-source-return}
On these hierarchy supports,
\[
 \iota_e\iota_\mu^\dagger=\iota_\mu\iota_e^\dagger
 =\iota_\mu\iota_\tau^\dagger
 =\iota_\tau\iota_\mu^\dagger=0.
\]
Hence the complete left return, retaining every term, is
\[
 \boxed{\begin{aligned}
 RR^\dagger={}&\varepsilon_e^2\iota_e\iota_e^\dagger
 +\varepsilon_\mu^2\iota_\mu\iota_\mu^\dagger
 +\varepsilon_\tau^2\iota_\tau\iota_\tau^\dagger\\
 &+\varepsilon_e\varepsilon_\tau
   (\iota_e\iota_\tau^\dagger+
    \iota_\tau\iota_e^\dagger).
 \end{aligned}}
\]
The free test of Theorem~\ref{p3f:thm:v2168-left-return-test} is equivalent to
these three owner-block identities:
\[
 \boxed{\begin{aligned}
 \iota_e\iota_\tau^\dagger&=0,\\
 \iota_\tau\iota_\tau^\dagger&=Q_0,\\
 \varepsilon_e^2\iota_e\iota_e^\dagger+
 \varepsilon_\mu^2\iota_\mu\iota_\mu^\dagger
 &=\varepsilon_e^2P_e+\varepsilon_\mu^2P_\mu.
 \end{aligned}}
\]
\end{theorem}
\begin{proof}
In every displayed muon cross product the adjacent right-support projectors
are $Q_0Q_\perp$ or $Q_\perp Q_0$, so the product is zero.
Expansion of the complete source gives the claimed return, without deleting
any term that can survive. Its off-diagonal $Q_\perp$--$Q_0$ corner is exactly
$\varepsilon_e\varepsilon_\tau\iota_e\iota_\tau^\dagger$.
The opposite ordered product occupies the adjoint corner and cannot cancel it.
The diagonal $Q_0$ corner is
$\varepsilon_\tau^2\iota_\tau\iota_\tau^\dagger$; the other two diagonal
terms lie in $Q_\perp\operatorname{End}(G)Q_\perp$.
Compare these three corners with
$UM^2U^\dagger=\varepsilon_e^2P_e+\varepsilon_\mu^2P_\mu+
\varepsilon_\tau^2Q_0$. Positivity of the terminal energies permits cancellation
of their nonzero scalar factors. This proves necessity and sufficiency.
\end{proof}

\subparagraph{Paper II-A's exact role.}
Preterminal relation preservation requires the expanded cross return to be
retained. The four zero products above have now been proved from their typed
supports. Paper II-A does not force the remaining electron--tau product to
vanish: the inputs share the right augmentation plane. Its value must come
from the existing source owner. This is a downstream source-energy comparison,
not a revived charged source-occurrence or parent-contraction gap.

\subsubsection{An explicit finite normal form for the remaining input}
the geometric completion provides a right-plane basis without an additional geometric selector:
\[
 h_e=R_-^\dagger g_e=\frac12(-1-i,1,i)^T,\qquad
 h_\tau=\frac1{2\sqrt6}(2,-1-3i,-1+3i)^T,
\]
where $R_-=\diag(1,i,-i)$. Put
$L=(g_e,g_\mu)$ and $W=(h_e,h_\tau)$.
Both are isometries onto $Q_\perp G$.
Every hierarchy-compatible physical source arrow has the unique form
\[
 \boxed{\iota_e=L A W^\dagger,\qquad
 \iota_\mu=L b c^\dagger,\qquad
 \iota_\tau=c d^\dagger W^\dagger,}
\]
where $A\in M_2(\mathbb C)$ and $b,d\in\mathbb C^2$ are evaluation outputs of the
actual source arrows. They are not new fitted parameters or a source family
to be selected. In the ordered right presentation $\widetilde V=(h_e,h_\tau,c)$,
\[
 U^\dagger R\widetilde V=
 \begin{pmatrix}\varepsilon_e A&\varepsilon_\mu b\\
                 \varepsilon_\tau d^\dagger&0\end{pmatrix},
\]
and therefore
\[
 \boxed{U^\dagger RR^\dagger U=
 \begin{pmatrix}
 \varepsilon_e^2AA^\dagger+\varepsilon_\mu^2bb^\dagger
       &\varepsilon_e\varepsilon_\tau A d\\
 \varepsilon_e\varepsilon_\tau d^\dagger A^\dagger
       &\varepsilon_\tau^2d^\dagger d
 \end{pmatrix}.}
\]
The full necessary and sufficient evaluation test is thus
\[
 \boxed{A d=0,\qquad d^\dagger d=1,\qquad
 \varepsilon_e^2AA^\dagger+\varepsilon_\mu^2bb^\dagger
       =\diag(\varepsilon_e^2,\varepsilon_\mu^2).}
\]
This formulation needs no assumption that the mass-source range and a
normalized current range are equal. If the test holds, $d$ is nonzero and
$Ad=0$ forces $\operatorname{rank}A\le1$.
The positive rank-two plane return then forces $A$ and $b$ to have nonzero,
independent left images, so $\operatorname{rank}A=1$.

\subsubsection{Where the existing construction supplies these inputs}
Let $J_{\rm HG}$ be Paper I-C's rooted hierarchy/family unitary and
$\Phi_{\rm HG}(X)=J_{\rm HG}^\dagger XJ_{\rm HG}$ its $*$-isomorphism.
For the source owner's already retained hierarchy arrows,
\[
 \iota_r=\Phi_{\rm HG}(A_r^{\rm hier}).
\]
The exact evaluation inputs are
\[
 A=L^\dagger\Phi_{\rm HG}(A_e^{\rm hier})W,\quad
 b=L^\dagger\Phi_{\rm HG}(A_\mu^{\rm hier})c,\quad
 d=W^\dagger\Phi_{\rm HG}(A_\tau^{\rm hier})^\dagger c.
\]
In particular the first surviving cross-return target is
\[
 \boxed{\Phi_{\rm HG}
   \bigl(A_e^{\rm hier}(A_\tau^{\rm hier})^\dagger\bigr)
       =\iota_e\iota_\tau^\dagger\stackrel{?}{=}0.}
\]
The same source owner supplies the two diagonal returns. Since
$\Phi_{\rm HG}$ is a $*$-isomorphism,
\[
 \Phi_{\rm HG}(XY^\dagger)
   =\Phi_{\rm HG}(X)\Phi_{\rm HG}(Y)^\dagger.
\]
Consequently proving the source-return identity upstream on the rooted
hierarchy arrows automatically transports it to the physical family source;
no additional family embedding is needed. The signed triality
channel/family $*$-isomorphism similarly preserves these returns. Paper II-A
then retains the generated identity in the complete action word.


\subsubsection{Observable-minimal charged return: the exact new evaluation frontier}
\label{p3f:sec:v267-one-observable-charged-test}
On the actual canonical finite charged chart, \emph{if constituted}, let
\[
 M_s=Z_{L,s}^{-1/2}\bigl(\delta^2S_\star/
 \delta\bar\psi_{L,s}\delta\psi_{R,s}\bigr)Z_{R,s}^{-1/2}
\]
be the complete mixed chiral action derivative in sector $s$.
Set $B_s=M_sM_s^\dagger$ and retain independently constituted
rooted left projections $P_{s,r}$ and target coefficients
$\lambda_{s,r}>0$ in
$D_s=\sum_r\lambda_{s,r}P_{s,r}$.  All assertions in this
subsection are on these actual \emph{same-source} blocks, not on
unconsumed comparison matrices.

\begin{theorem}[Two scalar contractions certify a complete finite charged left return]
\label{p3f:thm:v267-two-scalar-return-defect}
For the included mutually orthogonal finite charged sectors, let
$w_s>0$ be their actually inherited multiplicities and write
\[
 Q_4=\sum_sw_s\operatorname{Tr}(B_s^2),\quad
 Q_{\rm root}=\sum_sw_s\operatorname{Tr}(B_sD_s),\quad
 Q_*=\sum_sw_s\operatorname{Tr}(D_s^2).
\]
Then the exact nonnegative action comparison is
\[
 \boxed{Q_4-2Q_{\rm root}+Q_*
 =\sum_sw_s\|B_s-D_s\|_{\rm HS}^2\geq0,}
\]
with equality if and only if $B_s=D_s$ in \emph{every} included sector.
For one sector the alternative single positive-time scalar test is
\[
 \boxed{\operatorname{Tr}e^{-\tau B_s}
 -\operatorname{Tr}e^{-\tau D_s}
 +\tau\operatorname{Tr}(e^{-\tau D_s}(B_s-D_s))=0,}
\]
for any $\tau>0$, by Theorem~\ref{p2a:thm:v267-heat-bregman-converse}.
\end{theorem}
\begin{proof}
Expand $\operatorname{Tr}((B_s-D_s)^2)$ and use Hermiticity
and positive sector multiplicities.  The heat test is the same-source
application of the strict matrix-convex trace defect of Paper~II-A.
\end{proof}

\begin{corollary}[Native two-contact test for the free charged-lepton left return]
\label{p3f:cor:v3043-native-three-role-left-acceptance}
On the fixed common rooted left charged-lepton chart of
Theorem~\ref{p5:thm:v3218-two-contact-hurwitz-return}, let the
independently calculated, canonically normalized \emph{actual}
chiral jet be $N=Z_L^{-1/2}D_{\rm ch}^{\rm act}Z_R^{-1/2}$ and
$B_{\rm act}=NN^\dagger$.  With finite squared coefficients
$E,M,T>0$, rooted projectors $P_e,P_\mu,P_\tau$ and the \emph{geometric
comparison return}
$K_*=E I_G+(T-E)P+\frac98(M-E)CC^\dagger$, the already proved
Theorem~\ref{p3f:thm:v267-two-scalar-return-defect} specializes to
\[
 \boxed{\Delta_L=\operatorname{Tr}(B_{\rm act}^2)
  -2\operatorname{Tr}(B_{\rm act}K_*)+(E^2+M^2+T^2)
  =\|B_{\rm act}-K_*\|_{\rm HS}^2\ge0.}
\]
Here
$\operatorname{Tr}(B_{\rm act}K_*)=
E\operatorname{Tr}B_{\rm act}+(T-E)\operatorname{Tr}(PB_{\rm act})
+\frac98(M-E)\operatorname{Tr}(CC^\dagger B_{\rm act})$.
If and only if the independently calculated action return has
$\Delta_L=0$, its finite left free resolvent agrees with
\[
 \boxed{z(z^2-B_{\rm act})^{-1}=\sum_{r=e,\mu,\tau}
        \frac{z}{z^2-(m_r^{\rm fin})^2}P_r}
\]
on the common resolvent domain.  Neither the source-jet equality,
its physically occurring background/member, the stronger directed
$\mathsf{Cert}_{\rm hier}$, nor the interacting-pole identification
is established by this comparison corollary.
\end{corollary}
\begin{proof}
Use the orthogonal projector decomposition of $K_*$ supplied by
Paper~V-A, $\operatorname{Tr}(K_*^2)=E^2+M^2+T^2$,
and specialize the existing finite-dimensional squared-norm
identity.  Equality gives $B_{\rm act}=K_*$ and then its spectral
resolvent; the converse follows by equality of the resolvents and
their inverse pencils away from poles.
\end{proof}

\begin{corollary}[One unmarked complete-source heat value on the anchored domain]
\label{p3f:cor:v267-anchored-heat-action-test}
Assume additionally the physical source has the \emph{independently
verified} rooted first moments
$\langle g_{s,r},B_sg_{s,r}\rangle=\lambda_{s,r}$
on a complete orthonormal left frame in every included sector.
Then at a single $\tau>0$,
\[
 \boxed{\sum_sw_s\operatorname{Tr}e^{-\tau B_s}
 =\sum_sw_s\sum_r e^{-\tau\lambda_{s,r}}
 \ \Longleftrightarrow\ \forall s\; B_s=D_s.}
\]
The same implication follows from the normalized determinant equality
$\det B_s=\prod_r\lambda_{s,r}$ for each sector if the source is
already anchored into individually unit-normalized rank-one right arrows.
Neither route requires choosing a right-family member.
\end{corollary}
\begin{proof}
Each sector's heat excess is nonnegative by
Theorem~\ref{p2a:thm:v267-rooted-one-heat}; positive multiplicities
exclude cancellation.  The determinant alternative is
Proposition~\ref{p2a:prop:v267-normalized-determinant}.
\end{proof}

\begin{proposition}[Conditional extraction from the native exact Euclidean scalar heat]
\label{p3f:prop:v267-exact-heat-jet-normalization}
In the exact factorized Euclidean scalar branch of
Theorem~\ref{p3f:thm:v250-exact-parallel-scalar-heat}, suppose its
\emph{actual} finite self-adjoint factor has specifically the square
chiral form
$\Phi_M=\bigl(\begin{smallmatrix}0&M\\M^\dagger&0\end{smallmatrix}\bigr)$
for the \emph{same} $3\times3$ action source $M$.
With $B=MM^\dagger$, $\tau=s^2/L$ and
$S(a,s)=Z_g(a)\operatorname{Tr}_F e^{-s^2\Phi_M^2/L}$,
\[
 \boxed{\frac{S(a,s)}{2Z_g(a)}=\operatorname{Tr}e^{-\tau B},\qquad
 \operatorname{Tr}(B^2)=\frac{L^2}{24Z_g(a)}
 \left.\frac{\partial^4 S(a,s)}{\partial s^4}\right|_{s=0}.}
\]
For a complete sector direct sum, exact multiplicities and every
extra neutral/bosonic contribution must be separated before using
these as charged-sector invariants. The native raw scalar branch does
\emph{not} by itself identify $\Phi_M$ with the physical direct-mass
charged source, generate a rooted mark, or prove a Lorentzian pole.
\end{proposition}
\begin{proof}
The two chiral blocks of $\Phi_M^2$ are $MM^\dagger$ and
$M^\dagger M$; square finite sources have matching spectra, so their
complete heat traces coincide.  Expanding
$e^{-s^2\Phi_M^2/L}$ at order $s^4$ gives
$\partial_s^4S|_0=12Z_g\operatorname{Tr}_F\Phi_M^4/L^2
=24Z_g\operatorname{Tr}(B^2)/L^2$.
\end{proof}

\paragraph{What is now closed, and what remains.}
The mathematical free left-return test is reduced to a scalar heat
value if the anchored same-action first moments are proved, or to one
heat value plus one rooted mark without those moments.  Conversely,
a source that only shares the hierarchy \emph{corners} or the complete
unmarked heat spectrum can fail the rooted test: a unitary rotation of
$g_e,g_\mu$ within $Q_\perp$, applied to the normalized comparison
packet, preserves all three corner types and the entire spectrum but
changes its current-rooted left return.  Thus
Theorem~\ref{p3f:thm:v2168-hierarchy-source-return} cannot be replaced
by spectral matching.  The independent hierarchy-origin certificate
of Definition~\ref{p5:def:v228-hierarchy-origin}, the actual rooted row
moments, and a same-source full heat evaluation remain to be proved
before any conditional charged-source claim becomes unconditional.

\begin{corollary}[Conditional inverse of the occurring action corners]
\label{p3f:cor:v2191-action-corner-inverse}
On the existing uninserted local action branch, with $c_F>0$, the
same-action jet $(Z_L,D_e,Z_R)=c_F(I,R_{
m act},I)$ gives
\[
 \widehat R=Z_L^{-1/2}D_eZ_R^{-1/2}=R_{
m act}.
\]
For the fixed positive coefficients and the three hierarchy corner types,
an independently evaluated complete $R_{
m act}$ recovers its directed
source arrows uniquely:
\[
 \boxed{\begin{aligned}
 \iota_e&=\varepsilon_e^{-1}Q_\perp R_{
m act}Q_\perp,\\
 \iota_\mu&=\varepsilon_\mu^{-1}Q_\perp R_{
m act}Q_0,\\
 \iota_\tau&=\varepsilon_\tau^{-1}Q_0R_{
m act}Q_\perp.
 \end{aligned}}
\]
The hierarchy values are the corresponding corners of
$J_{
m HG}R_{
m act}J_{
m HG}^\dagger$, divided by their coefficients.
This is a lossless inverse on the existing three-corner source domain.
\end{corollary}
\begin{proof}
The two kinetic factors contribute $c_F^{-1}$ and cancel the factor in
$D_e$. Compress the existing source sum by the left and right projectors.
Disjoint corners kill the other two summands; the positive coefficient
of the remaining summand is nonzero. Transport back by the already fixed
unitary $J_{
m HG}$ preserves products and projector corners.
\end{proof}

This inverse is not an evaluation when $R_{
m act}$ is supplied only as
its symbolic sum over unevaluated hierarchy arrows. Evaluate the existing
same-action zero-order coefficient with its complete inherited marks, or
prove the needed actual return contractions directly. No additional source
occurrence, source selector or dynamics axiom is introduced. In particular
the electron-only test of Corollary~\ref{p6:cor:v2191-electron-defect} does
not demand a unique full right-arrow reconstruction. The normalized Hopf
current relation remains a different target sort unless a licensed
same-action physical coefficient word proves their equality.


\begin{proposition}[Minimal same-action chiral certificate for one marked two-pole response]
\label{p3f:prop:v239-same-action-electron-pencil}
Let $Z_L,Z_R>0$ and $D_{\rm ch}$ be the \emph{independently evaluated}
kinetic and mixed zero-order Hessian coefficients of one licensed
uninserted local action, all on their inherited chiral owner domains.
Define $M=Z_L^{-1/2}D_{\rm ch}Z_R^{-1/2}$ and let $g_e$ be the
specified unit left source in its canonical metric. For the marked
mass-only chiral resolvent, exact poles $\pm\varepsilon_e$ with
residues $1/2$ are equivalent to
\[
 \boxed{MM^\dagger g_e=\varepsilon_e^2g_e.}
\]
Writing $v=Z_L^{-1/2}g_e$, this is the same-action generalized pencil
\begin{equation}\label{p3f:eq:v239-electron-generalized-pencil}
 \boxed{D_{\rm ch}Z_R^{-1}D_{\rm ch}^\dagger v
          =\varepsilon_e^2Z_Lv,
 \qquad v^\dagger Z_Lv=1.}
\end{equation}
In particular it does not require the full three-family return identity,
individually reconstructed right arrows or an a posteriori choice of
an action-source coefficient.
\end{proposition}
\begin{proof}
Multiplying $MM^\dagger g_e=\varepsilon_e^2g_e$ by
$Z_L^{1/2}$ yields the displayed generalized eigenvalue equation.
Conversely the equation conjugated by $Z_L^{-1/2}$ yields the
canonical return. The chiral block resolvent and its spectral measure
give the two equally weighted poles when the positive return fixes the
unit source, as in Corollary~\ref{p6:cor:v2191-electron-defect}.
\end{proof}
\noindent\textit{Constitutive boundary.}
The equation tests an actual Hessian and does not generate one.
The geometric negative-Hopf reducing ray is an upstream sufficient
candidate, but the action must still show that its normalized mixed
jet consumes that ray with the required scale. The full interacting pole,
its stationary support and its residue remain separate physical readouts.

\begin{proposition}[Three right-Gram contractions certify the actual electron return]
\label{p3f:prop:v240-action-source-gram-test}
On the uninserted local branch of
Corollary~\ref{p3f:cor:v2191-action-corner-inverse}, the common
kinetic factor $c_F$ cancels and the canonical chiral coefficient is
$T_\star=R_{\rm act}$; it does not generate an independent exterior
normalization or right-family frame. Suppose, \emph{in addition},
that the actual mixed source arrows have independently certified
left-line supports
\[
 \iota_e=g_er_e^\dagger,\quad
 \iota_\mu=g_\mu r_\mu^\dagger,\quad
 \iota_\tau=cr_\tau^\dagger,
\]
where the right vectors are outputs of the same action rather than
chosen representatives. For $R_{\rm act}=e\iota_e+m\iota_\mu+t\iota_\tau$,
\begin{equation}\label{p3f:eq:v240-three-right-gram-defect}
\boxed{\begin{aligned}
 \|(R_{\rm act}R_{\rm act}^\dagger-e^2I)g_e\|^2
   ={}&e^4(\|r_e\|^2-1)^2\\
    &+e^2m^2|r_\mu^\dagger r_e|^2
     +e^2t^2|r_\tau^\dagger r_e|^2.
\end{aligned}}
\end{equation}
Thus on this certified left-support branch the actual electron
mass-only response is exact if and only if
\[
 \boxed{\|r_e\|^2=1,\qquad
 r_\mu^\dagger r_e=0,\qquad r_\tau^\dagger r_e=0.}
\]
The polar-normalized geometric Hopf triple of
Theorem~\ref{p6:thm:v240-canonical-polar-hopf-triple} satisfies these
contractions, but its existence does not prove that the mixed physical
Hessian produces this triple. Without certified left supports use the
full generalized pencil
\eqref{p3f:eq:v239-electron-generalized-pencil} instead.
\end{proposition}
\begin{proof}
The left orthonormality yields
$R_{\rm act}^\dagger g_e=e r_e$ and hence
\[
 R_{\rm act}R_{\rm act}^\dagger g_e
 =e^2\|r_e\|^2g_e
 +em(r_\mu^\dagger r_e)g_\mu
 +et(r_\tau^\dagger r_e)c.
\]
Subtract $e^2g_e$ and square orthogonally. The local common
kinetic cancellation is Corollary~\ref{p3f:cor:v2191-action-corner-inverse}.
\end{proof}
\begin{remark}[Comparison-frame placement remains an action datum]
A relative comparison frame with factor
$\theta=\sqrt{24/25}$ can replace a same-frame mass block $T$
by $\theta T$. If $TT^\dagger g_e=e^2g_e$, then with the
\emph{original fixed} terminal $e$ the marked defect becomes
$e^4(\theta^2-1)^2=e^4/625$. Thus common kinetic normalization
alone does not determine the physical chiral comparison-leg placement.
\end{remark}

\begin{proposition}[Exact two-jet electron defect without assuming a reducing block]
\label{p3f:prop:v245-electron-actual-block-defect}
Let $M=Z_L^{-1/2}D_{\rm ch}Z_R^{-1/2}$ be the \emph{actually
evaluated} mixed chiral coefficient of
Proposition~\ref{p3f:prop:v239-same-action-electron-pencil}.
In a left unit frame beginning with the tagged $g_e$ and an arbitrary
right unit frame, write
\[
 M=\begin{pmatrix}a&b\\c&D\end{pmatrix},\qquad
 a\in\mathbb C,\quad b\in\mathbb C^{1\times2},\quad
 c\in\mathbb C^{2\times1},\quad D\in\mathbb C^{2\times2}.
\]
For $e=\varepsilon_e>0$ its exact tagged defect is
\begin{equation}\label{p3f:eq:v245-electron-block-defect}
 \boxed{\mathcal D_e
 =\bigl(|a|^2+\|b\|^2-e^2\bigr)^2
 +\|\overline a\,c+Db^\dagger\|^2.}
\end{equation}
Thus the two-pole mass-only response $z/(z^2-e^2)$ holds if and only if
both displayed nonnegative summands vanish. Equivalently it suffices
and is necessary to evaluate the three tagged action contractions
\[
 \boxed{d_0=\langle g_e,MM^\dagger g_e\rangle-e^2=0,\quad
 d_1=\langle g_\mu,MM^\dagger g_e\rangle=0,\quad
 d_2=\langle c_0,MM^\dagger g_e\rangle=0,}
\]
where $(g_e,g_\mu,c_0)$ is the fixed rooted left role frame.
The geometric Hopf $1\oplus2$ form has $a=e,b=c=0$ and therefore
passes this certificate for arbitrary lower block $D$, but \emph{it
must not be assigned to the action-generated $M$ without evaluating its
actual mixed jet}. On the unchanged three-corner charged source grammar,
the singlet--singlet entry additionally vanishes:
$Q_0MQ_0=0$ (the bottom-right entry in the rooted right frame).
\end{proposition}
\begin{proof}
Direct multiplication gives
$MM^\dagger g_e=(|a|^2+\|b\|^2,\,\overline a c+Db^\dagger)^T$.
Orthogonality of the left frame gives the squared defect. The already
proved tagged resolvent equivalence in
Proposition~\ref{p3f:prop:v239-same-action-electron-pencil}
then gives the iff claim. Three-corner support implies
$Q_0MQ_0=0$ only when the \emph{actual} $M$ is the normalized
three-corner source in the same rooted frames.
\end{proof}

\begin{remark}[New geometric returns do not replace the mixed action jet]
\label{p3f:rem:v242-native-provider-action-scope}
Corollary~\ref{p1:cor:v242-joint-phase-vector-lift} computes the
signed lift of the existing joint Hopf phase with fixed rooted spinor
frames, and Proposition~\ref{p1:prop:v242-complete-transverse-return}
evaluates its complete transverse return. The scalar family embedding of
Theorem~\ref{p1b:thm:v242-family-quadratic-embedding} also fixes a
concrete finite cubic and its raw register-exclusion normalization.
These are provider calculations; none identifies their output with the
actual mixed chiral derivative. The independently evaluated
\(D_{\rm ch},Z_L,Z_R\) still enter
Proposition~\ref{p3f:prop:v239-same-action-electron-pencil}, or the
three right-Gram tests above when their left supports are certified.
No new physical source-occurrence token is introduced, and no source
insertion, stationary support or interacting pole is inferred from the
geometric coefficient alone.
\end{remark}


\begin{proposition}[A result-independent constructor certificate]
\label{p3f:prop:v2168-source-arrow-certificate}
If the actual source arrows satisfy
\[
 \boxed{\iota_r\iota_s^\dagger=\delta_{rs}P_r,
 \qquad(P_e,P_\mu,P_\tau)=(P_e,P_\mu,Q_0),}
\]
then the free tagged clock test holds for every positive coefficient tuple
in the role-linear word. Conversely, if the word $R(x)=\sum_r x_r\iota_r$
satisfies $R(x)R(x)^\dagger=\sum_r x_r^2P_r$ for every positive formal
coefficient tuple $x$, then the displayed arrow certificate holds.
\end{proposition}
\begin{proof}
Expansion proves sufficiency. For necessity, coefficient comparison in the
polynomial identity $R(x)R(x)^\dagger=\sum_r x_r^2P_r$ gives
$\iota_r\iota_r^\dagger=P_r$ and
$\iota_r\iota_s^\dagger+\iota_s\iota_r^\dagger=0$ for $r\ne s$.
Each rank-one diagonal return implies $\iota_r=u_rv_r^\dagger$ with
$u_r=Ue_r$ and $\|v_r\|=1$, absorbing its phase into $v_r$.
Thus the off-diagonal coefficient is
$\langle v_r,v_s\rangle u_ru_s^\dagger$ plus its adjoint. Contracting on
the left with $u_r^\dagger$ and on the right with $u_s$ proves
$\langle v_r,v_s\rangle=0$. This proves the certificate.
\end{proof}

The coefficient variation is an algebraic test of the result-independent
constructor, not a claim that the physical terminal masses vary freely.
For the single fixed physical tuple, the three block tests above are the
weaker exact necessary and sufficient condition.
The normalized geometric arrows $B_r$ constructed in the geometric completion satisfy this
certificate exactly. The supplied source does not establish
$\iota_r=B_r$, or this return certificate for the physical $\mathsf{SrcVal}$
arrows. Current-ray normalization alone is not that statement.

\subsubsection{The finite free clock poles and their accessible source weights}
Define the explicit comparison descendant
\[
 T_*:=\sum_r\varepsilon_rB_r=UMV^\dagger,
 \qquad h_*:=\begin{pmatrix}0&T_*\\T_*^\dagger&0\end{pmatrix}.
\]
It uses the existing terminal energies and the completed geometric relation.
It is not declared to be the actual physical source merely because it exists.

\begin{theorem}[Exact free rest pole and source certificate]
\label{p3f:thm:v2168-free-clock-poles}
On an admitted stationary finite clock branch whose complete quadratic
rest-symbol jet is $D_*(z)=zI-h_*$, one has
\[
 \Spec(h_*)=\{\pm\varepsilon_e,\pm\varepsilon_\mu,
                    \pm\varepsilon_\tau\}.
\]
Writing $u_r=Ue_r$, $v_r=Ve_r$,
\[
 q_{r,\pm}=\frac1{\sqrt2}\binom{u_r}{\pm v_r},\qquad
 h_*q_{r,\pm}=\pm\varepsilon_rq_{r,\pm}.
\]
For the normalized left-current probe $j_r=(u_r,0)^T$,
\[
 j_r^\dagger D_*(z)^{-1}j_s
 =\delta_{rs}\frac{z}{z^2-\varepsilon_r^2},\qquad
 \operatorname*{Res}_{z=\pm\varepsilon_r}
 j_r^\dagger D_*(z)^{-1}j_r=\frac12.
\]
Thus every role is seen by its corresponding left-current channel. If several
energies coincide, combine their projections; the resolvent poles remain
first order and channel weights remain as displayed.
\end{theorem}
\begin{proof}
With $W=\diag(U,V)$,
$h_*=W\left(\begin{smallmatrix}0&M\\M&0\end{smallmatrix}\right)W^\dagger$.
The two-component vector $(e_r,\pm e_r)/\sqrt2$ gives the claimed eigenpair.
Its overlap with $(e_r,0)$ has squared norm $1/2$. Summing the two resolvent
terms gives $\tfrac12[(z-\varepsilon_r)^{-1}+
(z+\varepsilon_r)^{-1}]=z/(z^2-\varepsilon_r^2)$.
\end{proof}

The stationary clock action is
\[
 S_*[\psi]=\int dt\,\psi^\dagger(i\hbar\partial_t-h_*)\psi.
\]
On its finite CAR Fock space, fill the negative-energy modes to form the
free ground state. Its energy is $-\sum_r\varepsilon_r$ per retained spin
copy. Addition in a positive mode and removal from a negative mode each cost
$\varepsilon_r$, by
$[d\Gamma(h_*),a^\dagger(q)]=a^\dagger(h_*q)$.
This gives a dynamical free one-particle interpretation on the stated
stationary clock/constraint domain, not only a matrix determinant.

\begin{corollary}[Rest-mass dictionary on the admitted free local Clifford branch]
\label{p3f:cor:v2168-free-clifford-dispersion}
Retain the same source and the already generated local Lorentzian Clifford
first symbol. For a local cotangent readout $\mathbf p$, its free Hamiltonian
symbol is, in a Weyl presentation,
\[
 H_*(\mathbf p)=
 \begin{pmatrix}
 -c_{\rm op}(\boldsymbol\sigma\cdot\mathbf p)\otimes I_G&I_2\otimes T_*\\
 I_2\otimes T_*^\dagger&c_{\rm op}(\boldsymbol\sigma\cdot\mathbf p)\otimes I_G
 \end{pmatrix}.
\]
On role $r$ its energies obey
\[
 E_r(\mathbf p)^2=c_{\rm op}^2|\mathbf p|^2+\varepsilon_r^2,
 \qquad \mathfrak m^{\rm free,rest}_r c_{\rm op}^2=\varepsilon_r.
\]
\end{corollary}
\begin{proof}
The spin and family factors commute, so the cross blocks of $H_*^2$ vanish.
The Pauli identity gives
$(\boldsymbol\sigma\cdot\mathbf p)^2=|\mathbf p|^2I_2$, while
$T_*T_*^\dagger=UM^2U^\dagger$ and
$T_*^\dagger T_*=VM^2V^\dagger$. This proves the displayed dispersion.
Its invariant free rest-energy coefficient is $\varepsilon_r$, with the
operational speed already fixed by the same first symbol.
\end{proof}

This is a theorem on the admitted free local symbol and source certificate.
It does not insert a new cotangent carrier, identify a compact harmonic
energy with this local rest energy, or establish a charged interacting
asymptotic particle. The primitive momentum resolution remains unchanged.

The derived finite free ratio on this comparison branch is
\[
 \frac{\varepsilon_\tau}{\varepsilon_\mu}
 =4\sqrt6+4\sqrt3\simeq16.7261622014.
\]
It is a consequence of the already generated terminal coefficients; it is
not yet the complete interacting physical mass-ratio prediction.

\subsubsection{The certificate retains nontrivial kinetic metrics}
\begin{theorem}[Same-action congruence certificate]
\label{p3f:thm:v2168-kinetic-congruence}
Let $C_L,C_R$ be invertible, time-independent actual coordinate-to-role maps inherited from one
complete stationary clock action, and $c_F>0$ its common coefficient. Suppose
its complete rest-symbol jet is
\[
 Z_L=c_F C_L^\dagger C_L,\quad
 D=c_F C_L^\dagger M C_R,\quad
 Z_R=c_F C_R^\dagger C_R.
\]
Then, for $\mathcal C=\diag(C_L,C_R)$,
\[
 \mathcal D(z)=\begin{pmatrix}zZ_L&-D\\-D^\dagger&zZ_R\end{pmatrix}
 =c_F\mathcal C^\dagger
 \begin{pmatrix}zI&-M\\-M&zI\end{pmatrix}\mathcal C.
\]
Its isolated pole positions are $\pm\varepsilon_r$. The canonically
normalized chiral block has those singular values as a derived comparison
identity. For $j=\sqrt{c_F}\mathcal C^\dagger(e_r,0)^T$, the source response
and its two residues equal those in Theorem~\ref{p3f:thm:v2168-free-clock-poles}.
\end{theorem}
\begin{proof}
The displayed congruence is direct multiplication. Its determinant differs
from $\prod_r(z^2-\varepsilon_r^2)$ by the nonzero constant
$c_F^6|\det\mathcal C|^2$. Taking the inverse and transforming the full source
gives the identical scalar response. Alternatively, if
$C_L=U_L(C_L^\dagger C_L)^{1/2}$ and
$C_R=U_R(C_R^\dagger C_R)^{1/2}$ are their polar decompositions, then
\[
 Z_L^{-1/2}DZ_R^{-1/2}=U_L^\dagger M U_R,
\]
where the common coefficient cancels. This proves the spectral comparison
without changing any already generated terminal energy.
\end{proof}

The theorem does not authorize choosing $Z_L,Z_R$ to force a desired mass.
All three blocks must be extracted together from the same actual action.
On a time-dependent frame the induced connection terms must also be retained;
the stationary formula does not discard them.
For the existing uninserted local branch the source already fixes
$Z_L=Z_R=c_FI$; the candidate above takes $C_L=U^\dagger$ and
$C_R=V^\dagger$, so the remaining finite identity is simply
\[
 \boxed{D_e=c_FUMV^\dagger.}
\]
More generally the three displayed congruence identities are the finite
action test. An energy-dependent Schur-reduced jet must instead retain its
full derivative and cannot be replaced by this affine certificate.


For an inserted constant kinetic chart the same criterion is applied to
$\widehat R=Z_L^{-1/2}D_eZ_R^{-1/2}$ and the actual normalized current
frame of that action. The source transforms with the same kinetic map.
The current-evaluation Gram is not automatically the propagation kinetic
block. An energy-dependent reduced jet retains the full inverse propagator,
source and derivative norm of Theorem~\ref{p3f:thm:v1140-finite-pole-source-transport}.
The three owner-block equalities above have not been evaluated for the
actual source exporter. They characterize the stronger complete free
response, not every interacting mass prediction.


\subsection{Actual component-band retention of the internal family coefficients}
\label{p3f:sec:v2168-component-family-retention}
\subsubsection{Applying the actual component-band retention constructor}
The III-F reference component-band branch normalizes the clock/spin form
and spatial half-density before any energy projection. A component band
acts on the external field components with the same inclusion for every
family component. The spherical product-band calculation likewise uses
its complete inherited embedding before instantaneous energy projection.
Its half-density calculation supplies the external normalization; its
invariant scalar mass-sector display is not used to project the family
carrier here. The relation-free core/family composition extends these
external constructions as $I_G$. This keeps the external mode/spin role
distinct from $G$.

\begin{proposition}[Family retention on the native component-band chart]
\label{p3f:prop:v2168-component-band-retention}
On this uninserted chart write the actual external inclusion as $J_m$.
Its family extension is $J=J_m\otimes I_G$, and the parent kinetic form is
$G_p=G_{p,m}\otimes I_G$. Then
\[
 G=J^\dagger G_pJ=G_m\otimes I_G,\qquad
 Z=JG^{-1/2}=Z_m\otimes I_G.
\]
Consequently every inherited coefficient $K_m\otimes A_G$ has the canonical
slot pullback
\[
 \boxed{Z^\dagger(K_m\otimes A_G)Z
       =(Z_m^\dagger K_mZ_m)\otimes A_G.}
\]
For an internal identity coefficient, the normalized family return is
exactly $A_G$, independent of the admitted external band index.
Different external charged and neutral inclusions give their actual
external overlap but retain the identity family coefficient between them.
Changing to the existing current tags therefore gives $C=U^\dagger W_\nu$.
\end{proposition}
\begin{proof}
Clock hermitization, spatial half-density and a component-band inclusion
act on external field slots and extend as $I_G$. The tensor-product
identities give $G$ and its positive root, hence $Z$ and the slot formula.
For the parent-metric isometry
$\widetilde Z=G_p^{1/2}JG^{-1/2}$,
\[
 \widetilde Z^\dagger\widetilde Z=I,\qquad
 \widetilde Z^\dagger(I\otimes A_G)\widetilde Z=I\otimes A_G.
\]
This is the actual factorization antecedent of III-F's v2.167 retention
corollary, established here by the component-band construction rather than
by faithful pair transport. Every derivative that preserves this external
chart has the same family factor. The cross-slot formula follows by
substitution of the two inclusions.
\end{proof}

The common readout across the admitted component-band family is its
\emph{internal coefficient}. Paper II-A can reuse that common readout
without choosing an external band. The external overlaps and full responses
can depend on the band and remain an image family. This does not prove that
a family-dependent mass spectral projection or an arbitrary constraint
quotient retains all current units. Additional reductions require their
actual reducing/intertwining certificate or the full generated return.
Nor does this establish the separate K16 normal-action handoff of A12.
The source Gram comparison $R_{\rm act}R_{\rm act}^\dagger=UM^2U^\dagger$
of the earlier free route is not inferred from current retention.

A component band can carry nonzero spatial kinetic energy. Retention of
the same internal current therefore does not assert a pole at
$\varepsilon_r$ on every such band. The terminal-energy kernel test below
consumes the claimed same-owner rest-comparison chart. If the source chart
instead reads a nonrest energy, its target must first be transported by the
already certified energy/rest dictionary. If that dictionary has not been
proved, it is part of the first chart-identification obligation. No
zero-momentum state is prepared or assumed by this coefficient calculation.


\subsection{Source-mode protection as an optional complete-action route}
\label{p3f:sec:v2169-pole-protection}

The direct physical-support polynomial in IV uses the complete $H$ and
requires no free-reference terminal equality. A separate sufficient route is
available when an actual source mode $O_0$ has already been generated with
$[H_0,O_0]=\varepsilon O_0$ and the consumed source/state overlap is nonzero.
For the complete inherited split $H=H_0+V$, the additional action identity
$[V,O_0]=0$ then gives $[H,O_0]=\varepsilon O_0$ by linearity.
This protects the pole without requiring the whole free two-point function.
It does not supply either antecedent from the terminal coefficient.

More generally, on the actual faithful operator chart let $P_\varepsilon$
be the full reference-frequency projection and $Q=I-P_\varepsilon$.
At a regular complement the exact endpoint condition is a nonzero
source-visible kernel of
\[
 P_\varepsilon\mathcal D P_\varepsilon+
 P_\varepsilon\mathcal DQ(\varepsilon-Q\mathcal LQ)^{-1}
 Q\mathcal DP_\varepsilon,
 \quad \mathcal D=\operatorname{ad}_V,\quad\mathcal L=\operatorname{ad}_H.
\]
On a physical nonfaithful support use its actual response quotient.
Paper~III-G, Appendix~\ref{p3f:app:v2169-pole-protection} constructs the complete mode,
all mixed first and second shifts, and the exact obstruction.
The formal mark $H_0+tV$ is evaluated at $t=1$ and creates no fitting
parameter. Zero first weak shift on definite-occupation reference eigenstates
leaves the complete second and higher returns to evaluate; two zero early
coefficients do not prove endpoint protection. A singular complement uses
IV's full source polynomial without this inverse.


\subsection{Native tensor packet and reference-boundary source comparison}
\label{p3f:sec:v224-native-source-comparison}
This section integrates the complete gravity/source bridge on its stated
domains. The ordinary common heat, a retained primitive readout, and the
conditional cap parent are three specified functional words. Equal scalar
coefficients do not identify them. No term below is added a second time
to a native packet already represented in other variables.

\begin{proposition}[Native local tensor coefficients]
\label{p3f:prop:v224-native-tensor}
On the ordinary torsion-free pure-metric Dirac specialization, with the
native internal rank 48, protected weight $1/8$ and the established local
clock, the four-derivative polynomial packet is
\[
 \mathcal L_{g,J}^{\le4}=\kappa_JR-V_0+\beta C^2+\eta\mathcal E_4,
 \qquad \kappa_J=\frac{\Lambda_g^2}{8\pi^2}=12U,
 \quad V_0=\frac{3\Lambda_g^4}{2\pi^2}=12\kappa_J\Lambda_g^2,
\]
\[
 \beta=\frac3{160\pi^2},\qquad \eta=-\frac{11}{960\pi^2},
 \quad \mathcal E_4=R_{\mu\nu\rho\sigma}^2
                          -4R_{\mu\nu}^2+R^2.
\]
The raw Einstein comparison is $8\kappa_J=96U$.
\end{proposition}
\begin{proof}
Before internal trace and protected weight, the Euclidean spin traces
are $a_0=4$, $a_2=-R_E/3$ and
$a_4=(-18C_E^2+11\mathcal E_{4,E})/360$, with common
$(4\pi)^{-2}$ and powers $\Lambda_g^4,\Lambda_g^2,1$.
Multiply by $48/8$ and apply the established negative Lorentzian
polynomial clock. This gives the coefficients directly. The nonmetric
position half-weight is not applied again to this pure-metric word.
\end{proof}
The statement inherits its controlled local expansion and actual boundary
problem. A same-metric pure $f(R)$ does not reproduce its off-shell tensor
Hessian: on a flat transverse-traceless jet $R^{(1)}=0$ while
$(C^{(1)})^2=2(\operatorname{Ric}^{(1)})^2\ne0$, modulo the Euler
boundary term. This does not infer a physical extra pole by extrapolating
the truncation outside its domain.

\begin{theorem}[Perturbative metric transport with sources]
\label{p3f:thm:v224-stress-transport}
Let $S_0=\int\sqrt{-g}\,\kappa(R-2\Lambda_0)+S_m$,
$\delta S_m=-\tfrac12\int\sqrt{-g}T_{\mu\nu}\delta g^{\mu\nu}$,
and define
\[
 H_{\mu\nu}=\Lambda_0g_{\mu\nu}
      +\frac{T_{\mu\nu}-Tg_{\mu\nu}/2}{2\kappa},
 \qquad H=4\Lambda_0-\frac T{2\kappa}.
\]
For constant $a,b$, introduce a formal first-order parameter $\varepsilon$
and the perturbatively invertible contravariant metric map
\[
 g^{\mu\nu}=\widetilde g^{\mu\nu}+\varepsilon Y^{\mu\nu},\qquad
 Y^{\mu\nu}=-\frac1\kappa
 \left[b(R^{\mu\nu}+H^{\mu\nu})
       -(a+b/2)(R+H)g^{\mu\nu}\right].
\]
All right-hand quantities are evaluated at unshifted fields to this order.
Then, with the whole source action, observables and boundary transported,
\[
 S_0+\varepsilon\int\sqrt{-g}(aR^2+b\operatorname{Ric}^2)
 \simeq S_0+\varepsilon\int\sqrt{-g}
                  (aH^2+bH_{\mu\nu}H^{\mu\nu})+O(\varepsilon^2).
\]
\end{theorem}
\begin{proof}
The leading metric variation is
$E_{\mu\nu}=\kappa(G_{\mu\nu}+\Lambda_0g_{\mu\nu})-T_{\mu\nu}/2$.
Its trace gives
$R_{\mu\nu}-H_{\mu\nu}=(E_{\mu\nu}-Eg_{\mu\nu}/2)/\kappa$.
Polarizing $P(A)=a(\operatorname{tr}A)^2+bA_{\mu\nu}A^{\mu\nu}$
therefore gives $P(\operatorname{Ric})-P(H)=-E_{\mu\nu}Y^{\mu\nu}$.
The first-order change of the complete $S_0$ cancels this difference.
The quadratic correction's own change is second order. This is an
off-shell functional identity after a change of variables, not the
insertion of vacuum equations into a fixed matter action.
\end{proof}
For Proposition~\ref{p3f:prop:v224-native-tensor}, the identity
$C^2=\mathcal E_4+2\operatorname{Ric}^2-2R^2/3$ gives
\[
 a=-\frac1{80\pi^2},\qquad b=\frac3{80\pi^2},\qquad
 \beta+\eta=\frac7{960\pi^2}.
\]
The residual Euler term retains its boundary/topological completion.
At $\kappa=12U$, $V_0=2\kappa\Lambda_0$, the source descendant is
\begin{equation}
 \Delta\mathcal L=
 \frac{3T_{\mu\nu}T^{\mu\nu}-T^2+2V_0T-4V_0^2}
                  {46080\pi^2U^2}.
 \label{p3f:eq:v224-stress-descendant}
\end{equation}
Indeed polarization also gives
$4(4a+b)\Lambda_0^2-(4a+b)\Lambda_0T/\kappa
 +(aT^2+bT_{\mu\nu}T^{\mu\nu})/(4\kappa^2)$.
Thus two co-consumed sectors have a fixed cross term proportional to
$2[3T^{(1)}_{\mu\nu}T^{(2)\mu\nu}-T^{(1)}T^{(2)}]$.
Conserved traceless stresses can still have a nonzero tensor contraction;
a trace-only scalar response cannot erase that observable-facing test.
Improved stresses retain all their curvature contacts. Higher orders,
boundary work and any quantum Jacobian are separate retained terms.
The bare value $\Lambda_0=6\Lambda_g^2$ is not set to zero by this map;
its putative vacuum $R=24\Lambda_g^2$ is outside the small-curvature
expansion. No absolute vacuum cancellation follows.

\begin{proposition}[Complete retained/primitive normalization dictionary]
\label{p3f:prop:v224-normalization-transport}
The retained readout of Proposition~\ref{p1b:prop:v224-retained-primitive}
has pure-metric $\kappa_{\rm rp}=\kappa_J/96=U/8$.
If a complete parity-even packet is uniformly multiplied by
$s=1/96$, canonical normalization transports its coefficients as
\[
 \kappa\mapsto s\kappa,\quad g_i^2\mapsto g_i^2/s,\quad
 h\mapsto\widehat h=\sqrt{s}\,h,\quad
 \lambda\mapsto\lambda/s.
\]
A retained degree-one Yukawa coefficient becomes $y/\sqrt{s}$;
a direct constant fermion mass and a scalar quadratic mass remain fixed.
For a cap written as $PR+BR^2-\sigma^2/B^2$, its complete scaling is
\[
 P\mapsto sP,\quad B\mapsto sB,\quad\sigma^2\mapsto s^3\sigma^2.
\]
All curvature derivatives and the canonical endpoint scale by the same
$s$; the vacuum/Einstein ratio stays fixed.
\end{proposition}
\begin{proof}
The first coefficient follows from the pure-metric trace identity.
Substitute canonically normalized fields into the multiplied kinetic,
quartic and degree-one terms. Substitute the displayed cap coefficients
into its full density and differentiate its equality.
\end{proof}
The native marked word also retains its graded trace. An arbitrary
interacting or boundary packet need not be uniformly rescaled. This
dictionary does not license a gravity-only division while leaving gauge
and scalar coefficients fixed, nor a quantum equivalence at fixed
operational $\hbar$. The original action's consumption of the marked
readout remains to be identified.

\begin{proposition}[Metric-reparametrization invariance of source exchange]
\label{p3f:prop:v3036-newton-exchange-invariance}
Let a gauge-fixed, invertible quadratic metric-source packet on its
physical nonzero-mode quotient be
\[
 S^{(2)}[h;J]=\tfrac12\langle h,Kh\rangle-\langle h,J\rangle,
 \qquad K=K^\dagger,\quad K^{-1}\text{ defined on the quotient}.
\]
For any bounded invertible field-coordinate map $A$, transport the whole
quadratic packet by $h'=Ah$ and the same source by
$K'=A^{-\dagger}KA^{-1}$, $J'=A^{-\dagger}J$.  Then
\begin{equation}\label{p3f:eq:v3036-canonical-source-exchange}
 \boxed{\langle J',(K')^{-1}J'\rangle
       =\langle J,K^{-1}J\rangle.}
\end{equation}
In particular, putting $h'=\sqrt{96}\,h$ rewrites the raw tensor
coefficient $12U$ as $U/8$ but also changes its canonical source vertex
by $1/\sqrt{96}$; it does not alter the measured exchange amplitude.
This statement is classical quadratic covariance on its stated quotient,
not a quantum path-integral/Jacobian equivalence certificate.
\end{proposition}
\begin{proof}
The inverse kinetic operator is
$(K')^{-1}=A K^{-1}A^\dagger$. Substitution of
$J'=A^{-\dagger}J$ cancels both field-coordinate maps. For
$A=\sqrt{96}\,I$ one has $K'=K/96$ and $J'=J/\sqrt{96}$.
\end{proof}

\begin{proposition}[Leading tensor pole under regular gapped stationary reduction]
\label{p3f:prop:v3036-gapped-schur-pole}
On a source-native local transverse-traceless weak-field quotient suppose,
for momentum $k$ near zero,
\[
 K_{hh}(k)=12U\,k^2 I+O(k^4),\quad
 K_{h\xi}(k)=O(k^2),\quad
 K_{\xi\xi}(k)=C_0+O(k^2),\qquad C_0^{-1}\ \text{bounded}.
\]
Then its exact stationary Schur complement obeys
\begin{equation}\label{p3f:eq:v3036-gapped-schur-residue}
 \boxed{K_{hh}^{\rm eff}(k)=
 K_{hh}-K_{h\xi}K_{\xi\xi}^{-1}K_{\xi h}
 =12U\,k^2 I+O(k^4).}
\end{equation}
Hence this restricted regular-gapped reduction cannot change the
leading massless tensor kinetic residue from $12U$ to $U/8$.
The statement concerns the measured force residue only if the physical
source vertex is also held fixed on this quotient (in particular no
independent interior-source incidence $J_\xi$). It does not exclude
additional massless modes, $O(k)$ or nonanalytic mixing, boundary/nonlocal
response or a separately derived complete marked action map.
\end{proposition}
\begin{proof}
Invertibility of $C_0$ gives $K_{\xi\xi}^{-1}=C_0^{-1}+O(k^2)$.
Both mixing factors are $O(k^2)$, so their Schur product is $O(k^4)$.
A physical source-source exchange additionally uses the fully transported
$J_h-K_{h\xi}K_{\xi\xi}^{-1}J_\xi$, not this kinetic identity alone.
\end{proof}

\begin{remark}[Status of the primitive Newton owner]
\label{p3f:rem:v3036-newton-owner-status}
The existing primitive insertion $P_+/48$ produces the 
\emph{pure-metric} coefficient $U/8$ exactly. The native ordinary heat
still yields $12U$, and neither an invertible metric-coordinate change
nor the regular gapped reduction above identifies the physical
canonically sourced exchange with the primitive readout. Physical
$G_N=G_{\rm scal}$ requires the already specified complete-action
metric, matter, gauge, boundary and measure matching of
Proposition~\ref{p3f:prop:v2211-newton-owner-matching}; it is not
promoted to closed by the two mathematical exclusions.
\end{remark}

\begin{theorem}[Reference-boundary origin of the complete cap parent]
\label{p3f:thm:v224-reference-cap-parent}
Fix $R_s,\ell>0$, $B_I>B_U>0$, $\Delta=B_I-B_U$,
$c=\sqrt{\Delta/B_I}$ and a constant Hilbert--Schmidt unit matrix $M$.
For an auxiliary curvature $r$, set
\[
 x=r/R_s,\quad G=\sqrt{1+x^2},\quad k=\ell G,\quad
 t=\frac{x}{1+G},\quad u_{\rm op}(r)=ct,\quad
 B(r)=B_U+\frac{2\Delta}{1+G}.
\]
Supply the positive normal square $\mathcal A_r=k^2I$ and fixed datum
\[
 B_0=b_0M,\qquad b_0^2=\frac{4\Delta R_s^2}{\ell},\qquad
 Q_0=b_0e^{-\ell y}M,\quad C_0=2\Delta R_s^2.
\]
Then the translation $Q=Q_0-F$, $F(0)=0$, forces
\begin{equation}
 j_r=(\mathcal A_r-\mathcal A_0)Q_0
     =\frac{2\ell^{3/2}\sqrt\Delta}{R_s}r^2e^{-\ell y}M,
 \label{p3f:eq:v224-forced-source}
\end{equation}
and the full off-shell identity is
\[
 B_Ur^2+\mathscr E_r[Q]-C_0
      =B_Ir^2+\mathscr E_D[F;j_r],
 \qquad \mathscr E_D=\tfrac12\|F\|_{L_r}^2
                                  -\operatorname{Re}\langle j_r,F\rangle.
\]
Its stationary value is exactly $r^2B(r)$.

With $\rho=R_s/2$, $T(u)=M_+(u)\otimes M_-(u)$,
$M_\pm(u)=\left(\begin{smallmatrix}1&\pm u\\\pm u&1\end{smallmatrix}\right)$,
and $W_Y=\operatorname{Tr}(Y^\dagger TY)+2\rho\operatorname{ReTr}Y$,
the first-order parent has bulk bracket
\[
 \Phi\{R[g]-r\}+Pr+B_Ur^2+\mathscr E_r[Q]-C_0
                 +W_Y[Y;u_{\rm op}(r)],\qquad Q(0)=B_0.
\]
With the existing factor $1/8$, boundary $\epsilon\Phi K/4$,
and a remainder independent of these auxiliary variables, it reduces
to the complete $f_c(R)=PR+R^2B(R)-\sigma^2/B(R)^2$,
$\sigma=B_IR_s$, and boundary $\epsilon f_c'(R)K/4$.
\end{theorem}
\begin{proof}
Theorem~\ref{p1b:thm:v224-poisson-defect} fixes the source, its sign,
profile and operator ordering. Direct integration gives
\[
 \mathscr E_r[Q_0]=\frac{b_0^2(\ell^2+k^2)}{4\ell}
                 =C_0+\Delta r^2,\qquad
 \min\mathscr E_r=\frac{b_0^2k}{2}=2\Delta R_s^2G.
\]
The difference is the complete source centre
$-J(r)^2/[4\ell(k+\ell)^2]=-\Delta r^2t^2$.
Rationalizing $G-1$ proves
$B_Ur^2+2\Delta R_s^2(G-1)=r^2B(r)$.
The positive four-port kernel has
$\operatorname{Tr}T^{-1}=4/(1-u^2)^2$, so its unique stationary
source value is $-R_s^2/(1-u^2)^2=-\sigma^2/B(r)^2$.
The multiplier equation gives $r=R[g]$, and the complete envelope
gives $\Phi=f_c'(r)$. The canonical boundary calculation is precisely
that of Theorem~\ref{p3f:thm:v223-cap-parent} with fixed $h_{ij},\Phi$.
All derivatives of $k,j,u$ occur in this same functional identity.
\end{proof}
The $R^2$ source is now a consequence of supplied normal and fixed-boundary
data, rather than an independent forcing ansatz. Its norm still encodes
the supplied susceptibility $\Delta=b_0^2\ell/(4R_s^2)$.
The nonzero datum is not a scalar VEV or a selected physical vacuum.
Here $r=0$ is an analytic reference, not the positive Clifford endpoint.
The mandatory $-C_0$ is an absolute term in this presentation; the theorem
does not authorize subtracting the native gravitational vacuum.

For the fixed-reference source profile, there is also a restricted
uniqueness check. Suppose its coefficient is a finite polynomial in $r$
with fixed matrix coefficients, and its negative source centre must
preserve the assigned primitive's constant, Einstein and infrared
quadratic slots without reallocating other sectors. Positivity of the
first nonzero squared norm forces its constant and linear source
coefficients to vanish. If its highest degree is $p$, the large-$r$
centre is $-R_s^2\|M_p\|_{\rm HS}^2r^{2p-2}/(4\ell^3)$.
A finite nonzero change of ultraviolet quadratic coefficient therefore
forces $p=2$ and $\|M_2\|_{\rm HS}^2=4\ell^3\Delta/R_s^2$.
This is uniqueness in that supplied profile/polynomial class; the
reference-defect theorem no longer needs a polynomial premise for its
forward implication. The complete inverse-square source still contributes
its own curvature derivatives. For example $(r^2B(r))''=2B_U+2\Delta/G^3$
is the gravity primitive's second derivative, not the complete $f_c''$.

\begin{proposition}[Off-shell form transport to the finite parent]
\label{p3f:prop:v224-collar-finite-transport}
Let $h_0=e^{-\ell y}M$ and
\[
 \psi_k=L_r^{-1}h_0
   =\frac{e^{-\ell y}-e^{-ky}}{k^2-\ell^2}M,\qquad
 q_0=\langle h_0,\psi_k\rangle=\frac1{2\ell(k+\ell)^2}.
\]
The formula has its regular $k=\ell$ limit. Put
$a_k=2\sqrt{B_I\ell}(k+\ell)$ and decompose the real sourced form
direction as $F=-a_k z\psi_k+F_\perp$, with
$\operatorname{Re}\langle\psi_k,L_rF_\perp\rangle=0$. Then
\[
 \mathscr E_D[F;j_r]=B_I(z^2+2u_{\rm op}(r)rz)
                                      +\tfrac12\|F_\perp\|_{L_r}^2.
\]
\end{proposition}
\begin{proof}
The Green image follows by differentiating its exponentials. Their
integrals give $q_0$. The complement is also source-orthogonal because
$L_r\psi_k=h_0$. Substitute into the whole form and use
$a_k^2q_0/2=B_I$, $Ja_kq_0=2B_Iu_{\rm op}r$.
\end{proof}
This gives the existing $B_I(r^2+2urz+z^2)$ parent before elimination,
with a displayed positive complement. A complex presentation may retain
the imaginary sourced direction in that real orthogonal complement.
Other consumed sources or observables on the complement remain present.
The moving form coordinates retain their full parameter jets and
determinant; a finite affine translation alone proves no continuum
measure or physical state.

The translated stationary outward flux is
$\pi_F=-b_0(k-\ell)M$; hence, with its continuous value at $r=0$,
\[
 \frac{-\operatorname{ReTr}(M^\dagger\pi_F)}
                   {2\sqrt{\ell B_I}\,r}=ct=u_{\rm op}(r).
\]
This is the translated normal-response flux, distinct from total unforced
flux $kB_0$ and from canonical metric momentum. The four-port kernel in
the theorem consumes $u_{\rm op}(r)$ as a geometric coefficient. Replacing
it by an independently varied off-shell $F$ flux changes the $F$ equation
through $W_Y$ and requires a new coupled reduction. Equality of stationary
fluxes alone never licenses that substitution. The existing shared-interior
rule is preserved, not bypassed.

\begin{proposition}[Companion and tangential response tests]
\label{p3f:prop:v224-companion-descendant}
For the rooted full-rank internal source
$M_0=K_F/\sqrt{48}$, $K_F^\dagger K_F=I$, its fixed-boundary energy is
$E_*=b_0^2\operatorname{Tr}\sqrt{\mathcal A}/96$.
At $\mathcal A=k^2I$, no nonzero positive perturbation on this same
internal fibre is invisible. For $\|V\|<k^2$,
\[
 E_*(k^2I+V)-\tfrac12b_0^2k
 =\frac{b_0^2}{4k\,48}\operatorname{Tr}V
 -\frac{b_0^2}{16k^3\,48}\operatorname{Tr}V^2+O(\|V\|^3).
\]
If an actual positive Euclidean realization also has tangential kinetic form
\[
 \tfrac12\int_0^\infty\operatorname{Tr}
                    (D_\alpha Q)^\dagger D^\alpha Q\,dy,
\]
its slow-variation reduction through two tangential derivatives, at fixed
$b_0,\ell,R_s$ and unit $M$, contains
\[
 \mathcal L_\partial^{(2)}=
 \frac{b_0^2}{4k}\operatorname{Tr}(D_\alpha M)^\dagger D^\alpha M
             +\frac{b_0^2}{8k^3}(\nabla k)^2.
\]
For parallel $M$ the latter coefficient of $(\nabla r)^2$ is
$\Delta(r/R_s)^2/(2\ell^2G^5)$.
\end{proposition}
\begin{proof}
Full rank and Proposition~\ref{p1b:prop:v224-root-source-jets} prove
the positive-support statement. Expand the matrix root at the scalar
baseline to second order. For the derivative term insert
$Q_*=b_0e^{-ky}M$ into the tangential form. Unit norm gives
$\operatorname{ReTr}M^\dagger D_\alpha M=0$; use
$\int e^{-2ky}dy=1/(2k)$ and
$\int y^2e^{-2ky}dy=1/(4k^3)$. The correction to $Q_*$ first affects
the next derivative order by zeroth-order stationarity.
\end{proof}
A nonzero traceless companion can have zero first contact and nonzero
second contact. For example
$\sqrt{\left(\begin{smallmatrix}25&24\\24&25\end{smallmatrix}\right)}
=\left(\begin{smallmatrix}4&3\\3&4\end{smallmatrix}\right)$ changes
the $e_1$ root response from 5 to 4 despite a zero diagonal contact of
$V=\mathcal A-25I$. A separate invariant sourced-orthogonal block can
preserve the scalar return only with its actual embedding and own action.
The tangential result is conditional on the physical kinetic form,
positive gap and controlled derivative expansion. Its clock, continuation
and boundary work must be inherited. A fibrewise analytic representation
does not itself postulate this term or a physical extra dimension.

\subsubsection{Boundary port and universal-weight tests}
\label{p3f:subsec:v224-boundary-port-tests}
There is also an independent boundary-domain test. For the same distributed
source on a frozen product collar, Dirichlet elimination gives
$E_D=-J^2/[4\ell(k+\ell)^2]$. Allowing its Neumann endpoint to relax
adds $-J^2/[2k(k+\ell)^2]$. The resulting coefficient profiles are
$H_D=H$, $H_N=H/G$, where $H=2/(1+G)$. Thus a block-diagonal mixed
response with normalized positive Neumann source population $p$ has
\[
 H_p=(1-p)H+pH/G,\qquad H_{1/2}=1/G.
\]
The target $H$ on an interval forces $p=0$. This conclusion concerns
that mixed response problem. An actual separately typed Dirichlet
interior Schur port need not inherit a freely relaxing Neumann endpoint
from another heat role; the action must establish which port is consumed.

On the common fibre let $\Gamma=\Gamma_{\rm rel}$ and $\chi$ be
the boundary involution with $\{\Gamma,\chi\}=0$.
For $P=(I+\Gamma)/2$ and $\Pi_D=(I-\chi)/2$,
\[
 P\Pi_DP=P/2,\qquad
 V_{D\leftarrow P}=\sqrt2\,\Pi_DP|_{\operatorname{ran}P}
\]
is an isometry onto the Dirichlet range. The overlap fixes its polar
normalization without a fitted angle. The insertion $P/48$ has trace
two but balanced Neumann population; $\Pi_D/48$ has trace two and
pure Dirichlet support. Transporting the whole operator, domain and
source by a unitary preserves their relative population. Wiring only
the source through this polar map changes its incidence and requires
the actual physical source constructor and all its derivatives.

The scalar-dependent extension of this local boundary test is established,
on a more restrictive ordinary-heat collar domain, by
Propositions~\ref{p3f:prop:v3031-mixed-square-boundary-domain}
and~\ref{p3f:prop:v3031-mixed-scalar-boundary-heat};
Corollary~\ref{p3f:cor:v3031-boundary-scalar-two-jet}
keeps the physical source/continuation comparison open.

The inherited mixed heat coefficient supplies a further restricted
check. For a normal-parallel fixed-block insertion $Z$ of trace two,
with zero normal smearing derivative and Robin
$S=-K_\partial\Pi_N/2$, its local boundary contribution is
$\operatorname{Tr}[Z(2K_\partial I+12S)]/6$.
Putting $p=\operatorname{Tr}(Z\Pi_N)/2$ yields
\[
 a_2^{\rm bulk}=-R/6,\qquad
 a_2^\partial=\tfrac23(1-3p)K_\partial.
\]
The boundary-to-bulk coefficient ratio is $12p-4$. The canonical
Einstein/GHY ratio two requires $p=1/2$, whereas this same frozen
nonlinear response has the target $H$ only at $p=0$. Thus one universal
frozen weighting on this particular mixed domain cannot repair both
comparisons. This does not exclude the declared cap parent, whose
Einstein and nonlinear response roles retain their own canonical
endpoint. A whole-action map must license those roles together; a
universal factor 96 changes neither ratio nor source population.

\subsubsection{Evaluated limits of normalization and mixing routes}
\label{p3f:subsec:v224-route-tests}
The following tests reject specified mechanisms, not every completion.
They accompany the positive vacuum bound already established in
Proposition~\ref{p3f:prop:v2216-bare-heat-vacuum} and
Theorem~\ref{p3f:thm:v2216-positive-profile-vacuum-bound}.

The native unit-preserving scalar Laplacian and the retained finite LC
square have a marked constant kernel. A line-preserving injective
comparison cannot identify either full operator with the strictly
positive cap normal square. Selecting a nonzero band, shifting the
operator or taking a source-conditioned Schur response is a different
word requiring its actual map. The source-support criterion in
Proposition~\ref{p1b:prop:v224-heat-cap} gives a sharper positive-weight
test of a proposed single cap; equality of scalar endpoint values is
insufficient.

On the normalized core/complement constants of the native nearest-shell
walk, the comparison $zI+I-T_I$, $z>0$, has block
\[
 \begin{pmatrix}z+1&-1/2\\-1/2&z+1/4\end{pmatrix}.
\]
Its Schur kernels are $Q_H=z(4z+5)/(4z+1)$ and
$Q_Y=z(4z+5)/[4(z+1)]$. Elimination with sources
$-j_Hb-j_Yy$ transports $j_{H,\rm eff}=j_H+2j_Y/(4z+1)$
and the centre $-j_Y^2/[2\gamma(z+1/4)]$, where $\gamma$ multiplies
the quadratic form. The normalized common-mode coefficients obey
$24\cdot5=96\cdot(5/4)=120$. Thus this actual finite comparison
does not produce $1/96$; the retained primitive readout is a different
operation and is not excluded by this result.

On the local doublet attenuation branch, $F(h)=\kappa_J-h^2/12$.
Requiring $F_0=\kappa_J/96$ gives
$h_0^2=95\kappa_J/8$. Its Einstein scalar metric is
$Z_E=F_0\kappa_J/F^2$; differentiating the gauge mass function
proportional to $h^2F_0/F$ along canonical scalar coordinates yields
the complete squared coupling sums
$\sum_i\kappa_{iW}^2=\sum_i\kappa_{iZ}^2=\kappa_J/F_0=96$.
Hence attenuation and gauge--scalar response cannot be selected separately.

On a parallel source-free odd ray $\Phi=vK_F$, the native local equation
is $\Box v-Rv/6+2\Lambda_g^2v-2v^3=0$.
Its nonzero constant solution has $v^2=\Lambda_g^2-R/12$, with
$v'=-1/(24v)$ and $v''=-1/(576v^3)$, rather than $v=aR$.
With $C=1/(16\pi^2)$ and
$\delta_\Phi S_{\rm extra}=\int\sqrt{-g}\operatorname{Tr}
(\mathcal J\delta\Phi)$, an exact affine branch $v=aR$ requires
\[
 \mathcal J_{\rm req}=C[-2a\Box R-4\Lambda_g^2aR
                                +aR^2/3+4a^3R^3]K_F.
\]
This follows by substituting the branch into the complete scalar Euler
equation. Moving $K_F$ adds its covariant derivatives. With the retained
bare vacuum, scalar stationarity on this nonzero positive-quartic branch
requires $R<12\Lambda_g^2$, while the joint constant metric equation
with positive Einstein coefficient requires $R>12\Lambda_g^2$.
The tested bare constant-background mechanism therefore fails; an
interacting or nonconstant completion must supply its current, stress
and stabilizing second variations together.

For a literal intrinsic lift $C_f=i[\mathscr D,f]$ on the round Dirac
domain,
\[
 \mathscr D+sC_f=e^{-isf}\mathscr D e^{isf}.
\]
The full unmarked heat spectrum is independent of constant $s$.
For the lowest spinor projector $P$, the native multiplication images
give $PC_fPC_fP=9P/16$, $PC_f(1-P)C_fP=3P/16$.
The isolated square has $m_0^2+9s^2/16$, while the full corner has
$m_0^2+3s^2/4$. The second-order pole shifts from the opposite retained
sign and first omitted band are respectively $3r_p/16$ and
$-3r_p/16$. Their cancellation is also forced by the unitary identity.
A fixed projector readout differs from a transported projector. Thus an
isolated pencil is not by itself a full heat-action proof.

On a round product collar, the lowest tangential Dirac square is
$3R_3/8$. Identifying it with $2\ell^2x$ for $x=r/R_s$ requires the
specific comparison $R_3=r$ and $\ell^2=3R_s/16$.
The affine completion of Proposition~\ref{p1b:prop:v224-normal-completion}
then gives the candidate normal square. If $R_s^2=UC_7$ is separately
licensed, its gap is $\ell=(\sqrt3/4)(UC_7)^{1/4}$.
This is not an identity between spatial round curvature and general
four-dimensional Ricci curvature. The native band, added odd channel
and scale remain to be identified in the actual action.

Finally, on the explicit common $1/96$, $P=U$ comparison, the bare heat
constant inside the cap bracket is $V_{\rm heat}^{(f)}=1152\pi^2U^2$.
With $R_s=U\alpha^{-7/2}$, $\alpha>1$ and
$0<\Delta<B_I=\alpha^{7/2-\delta/2}<\alpha^{7/2}$,
\begin{equation}
 \frac{C_0}{V_{\rm heat}^{(f)}}=
 \frac{\Delta\alpha^{-7}}{576\pi^2}
       <\frac{\alpha^{-7/2}}{576\pi^2}<1.
 \label{p3f:eq:v224-reference-vacuum-test}
\end{equation}
This exact coefficient mismatch prevents assigning the reference energy
as the sole cancellation of the native heat vacuum. Multiplicity scales
complete words and cannot supply that missing absolute energy.

The remaining common-action comparison must therefore establish the
actual positive normal band and scale, nonzero fixed boundary datum and
norm, curvature/source and normalization maps, all tensor, derivative,
source-orthogonal and boundary companions, physical measure and one
absolute energy allocation. I/II co-realization transports a licensed
constructor with every required relation; it does not infer this missing
action equality backwards from its scalar stationary value.


\subsection{Charged origin certificate and the smaller cyclic source test}
\label{p3f:sec:v228-charged-character-test}
\noindent\textbf{Current source scope.}
The noncentral charged handoff of
Theorem~\ref{p5:thm:v397-charged-preterminal-srcval-realization}
consumes independently generated hierarchy row arrows. Its current domain
is the certified packet of Definition~\ref{p5:def:v228-hierarchy-origin}.
The native uninserted kinetic normalization and family-preserving component
inclusion remain valid; they transport an actual source and do not create
its input arrows. The explicit rooted primitive identity-family source
retains its separate origin proof.
\begin{corollary}[A derived right electron pair]
\label{p3f:cor:v228-derived-electron-pair}
On an admitted canonical finite source chart let $T$ be the actual mass-only
chiral coefficient and $J_L$ an isometric electron embedding. If
$(TT^\dagger-e^2I)J_L=0$, $e>0$, then
\[
 J_R=e^{-1}T^\dagger J_L,\qquad
 J_R^\dagger J_R=I,\qquad TJ_R=eJ_L.
\]
Thus $h_T=\left(\begin{smallmatrix}0&T\\T^\dagger&0\end{smallmatrix}\right)$
has the paired block
$\left(\begin{smallmatrix}0&eI\\eI&0\end{smallmatrix}\right)$
and left source resolvent $z/(z^2-e^2)$ on this embedding.
If the actual hierarchy word belongs to the source class in
Theorem~\ref{p1b:thm:v228-character-defect}, its positive character
formula is necessary and sufficient for the scalar marked claim.
Corollary~\ref{p1b:cor:v228-operator-character-defect} supplies the
complete test with finite noncommuting outer coefficients.
\end{corollary}
\begin{proof}
Insert $TT^\dagger J_L=e^2J_L$ in $J_R^\dagger J_R$ and $TJ_R$.
The pair is isometric and invariant under the Hermitian $h_T$, hence
reducing. Inverting its two-by-two block proves the resolvent formula.
No independent right-family isometry is introduced.
\end{proof}
An actual source in the rank-one fibre of
Proposition~\ref{p1b:prop:v228-opposite-odd-fibre} has the complete
three-role free left return without a member selection. Its source-domain
identification is still required. Weak/scalar terms, stationary support
and physical pole/rest interpretation are separate consumed data.


\subsection{Parent Cartan origin of connection and boundary response}
\label{p3f:sec:v187-parent-cartan-boundary}

\begin{theorem}[Vector Born connection as the descended parent restriction]
\label{p3f:thm:v187-born-connection-parent-restriction}
Let $\mathcal A_{M8},\mathcal E_{M8}$ be the complete $K$-valued Cartan split
of Paper~III-B.  The scalar-diagonal descent theorem~\ref{p3b:thm:parent-cartan-scalar-diagonal-descent} first generates
$\omega_{\rm cmp}$ on $X_7^{\rm cmp}$.  Restricting that generated connection
to the vector $SU(2)$ factor and then to the canonical labelled-factor section
gives
\[
 A_B=E_B=\frac12\theta_r,
 \qquad
 \boxed{F_B=-E_B\wedge E_B=-\frac14\theta_r\wedge\theta_r.}
\]
Thus the Born vector connection and its curvature are not selected locally
from an unresolved family and are not a separate bridge out of the relative
geometry: they are the uniquely descended vector component of the $M_8$
Maurer--Cartan split.
\end{theorem}

\begin{theorem}[Boundary-germ action descent and quotient--Schur ordering]
\label{p3f:thm:v187-boundary-germ-schur-quotient}
Suppose the regular-boundary constructor generates the local oriented collar
and trace $T:V\to J_\partial$ from the parent normal of Paper~III-C, and let
$H=d^2S$ be the same-action Hessian.  Put $I=\ker T$.  On the regular branch
where the restriction $H_{II}$ is invertible, stationarity canonically selects
the $H$-orthogonal complement of $I$ and generates the Schur descendant
\[
 \boxed{K_\partial
 =H_{\partial\partial}-H_{\partial I}H_{II}^{-1}H_{I\partial}.}
\]
Every projectable Ward/constraint descendant is retained with it.  Any
terminal rank-changing quotient must be applied only after this elimination,
unless the mixed memory block vanishes.  Consequently the local structural
G25 response and structural G28 endpoint interface descend from the same
parent collar and action; an actual endpoint crossing remains G0 history.
\end{theorem}

\begin{remark}[Global dynamics guardrail]
Paper~I-C, Theorem~\ref{p1b:thm:g12-double},
Proposition~\ref{p1b:prop:g12-reduction}, and
Proposition~\ref{p1b:prop:relational-two-forms}
\cite{Kugelmann2026Paper1C}, construct and compare the compact-double and
reciprocal branches.  With the normalized common maximal torus they agree
there, but globally their rank and exterior-derivative data differ; in
particular there is neither a global equality of the two forms nor a closed
two-form \(B_0\) whose addition turns the compact-double form into
\(\Omega\).  This proves mathematical existence and a sharp separation of
the branches, but does not make the compact double a physical constructor
here: the present source-derived chain supplies neither a physical
normalization nor a licensed action/moment-map occurrence for it.  The
quasi-Hamiltonian completion therefore remains separately typed and unused in
the physical III-F dynamics until those missing bridges are proved.
\end{remark}

\subsection{Generated same-action boundary owner}
\label{p3f:sec:v183-generated-boundary-owner}
\begin{theorem}[Same-action boundary-owner generation]
\label{p3f:thm:v183-generated-boundary-owner}
Let $S_\star$ be a claim-complete physical action, $q_\star$ a stationary or
source-conditioned configuration, and suppose the physical boundary trace and
split
\[
 T_{q_\star}\mathcal C_{\rm stat}=J_\partial\oplus J_I
\]
are already declared parts of that same action datum.  On a domain where the
interior Hessian block $H_{II}$ is invertible, stationary Schur reduction generates
\[
 S_\star^\partial,
 \qquad
 \Theta_\partial=dS_\star^\partial,
 \qquad
 K_\partial=H_{\partial\partial}-H_{\partial I}H_{II}^{-1}H_{I\partial},
\]
together with every projectable Ward/constraint descendant and every licensed
source derivative such as $V^\partial(\eta)=\nabla_\eta K_\partial$.
Therefore
\[
 \boxed{\mathcal P_\partial(S_\star,q_\star)=
 (S_\star^\partial,J_\partial,\Theta_\partial,K_\partial,V^\partial,
 \mathcal W_\partial,\mathcal C_\partial)}
\]
is a generated joint owner in the sense of Paper~II-A.  No second occurrence
is required for these descendants.

The theorem does not generate the physical trace/split itself, nor any
cross-sort map from an independently typed geometric carrier into
$J_\partial$.
\end{theorem}


\subsection{typed $q=5$ paired-complement transport gate}
\label{p3f:sec:v1113-q5-complement-transport}

The upstream geometry now fixes the unique $q=5$ rotation plane, its exact
trimer realization and the Born-odd sign line.  The remaining action-level
question is whether the paired complement of the actual finite fermion source
uses that generated transport.

\begin{definition}[Actual $q=5$ complement-transport occurrence target]
\label{p3f:def:v1113-q5-complement-transport-gate}
On a typed finite left/right complement write
\[
 QD_{\rm fin}Q
 =D_Q
 =\Delta
 \begin{pmatrix}
 0&U_Q\\
 U_Q^\dagger&0
 \end{pmatrix},
 \qquad \Delta>0.
\]
Historically, the v1.113 G15 half-step branch introduced the conditional
same-action identification
\[
 \boxed{U_Q=\Psi_u^{-1}R_u\Psi_u=X_G}
\]
for one rooted orientation, with $U_Q=X_G^\dagger$ on the completely reversed
rooted presentation.
\end{definition}

\begin{proposition}[Conditional Feshbach specialization]
\label{p3f:prop:v1113-q5-feshbach-specialization}
Assume the occurrence target above and let the retained-to-complement coupling
be chirality even,
\[
 B=\begin{pmatrix}A_L&0\\0&A_R\end{pmatrix}.
\]
Then stationary Schur reduction of the invertible complement yields
\[
 -BD_Q^{-1}B^\dagger
 =-\frac1\Delta
 \begin{pmatrix}
 0&A_LX_GA_R^\dagger\\
 A_RX_G^\dagger A_L^\dagger&0
 \end{pmatrix}.
\]
Thus the descendant is automatically a zero-order odd
$L\leftrightarrow R$ source, and its normalized reciprocal/orientation-odd
family direction is the already generated
$C_{2,G}=-i(X_G-X_G^\dagger)$.
\end{proposition}

\begin{proof}
The inverse of
$D_Q=\Delta\left(\begin{smallmatrix}0&X_G\\X_G^\dagger&0\end{smallmatrix}\right)$
is
$\Delta^{-1}\left(\begin{smallmatrix}0&X_G\\X_G^\dagger&0\end{smallmatrix}\right)$
because $X_G$ is unitary.  Direct block multiplication gives the displayed
Schur/Feshbach correction.  The odd family direction is then the antisymmetric
Hermitian projection of the reciprocal pair.
\end{proof}

\begin{remark}[No occurrence from covariance alone]
\label{p3f:rem:v1113-occurrence-firewall}
Born covariance fixes how an already occurring transport transforms; it does
not insert that transport into $QD_{\rm fin}Q$.  In particular a Born-covariant
source may carry its Born parity in another tensor factor while its family
factor remains $I_G$.  Therefore the v1.113 equality $U_Q=X_G$ was a conditional same-action
source-incidence hypothesis rather than a consequence of covariance.  Release
v1.117 retires it as a live selector: the actual nonzero fermion-source/member
occurrence is the upstream G0 boundary, while the generated q=5 transport is
retained only as supporting structural geometry.
\end{remark}


\subsection{safeguard synchronization of the fermion-source boundary}
\label{p3f:sec:v1117-g15-safeguard-sync}
The q=5/trimer transport, its rooted reversal, and the conditional Feshbach
calculation above remain valid mathematical statements on their stated domains.
They do not own the physical occurrence of a nonzero fermion source.  In the
current framework the structural G15 branch is closed upstream and the actual
sector-specific nonzero source/member is typed as
\texttt{G0\_FERMION\_SOURCE\_MEMBER}.  Hence
$U_Q=X_G$ is no longer an active source gate and is not a premise of any
downstream mass, mixing, CP, RG, pole, QCD, or cosmology readout.

The dependency direction is therefore
\[
\begin{gathered}
 \text{generated source/Dirac owner family}
 +\texttt{G0\_FERMION\_SOURCE\_MEMBER}\\
 \longrightarrow\text{invariant readouts}
 \longrightarrow\texttt{G16}_{\rm RG},
\end{gathered}
\]
with pole certification added only for a claim that needs a pole mass.  No
scalar vacuum expectation value and no q=5 occurrence selector is inserted in
this chain.


\subsubsection*{Finite fermion gap and transport scope}
\label{p3f:sec:v1118-fermion-gap-transport}

\begin{theorem}[Finite source-to-operator gap on a source-preserving branch]
\label{p3f:thm:v1118-finite-source-operator-gap}
Let an actually occurring finite odd source be represented by
\[
 \Phi_T=\begin{pmatrix}0&T\\T^\dagger&0\end{pmatrix},
 \qquad m_T:=\sigma_{\min}(T)>0.
\]
On a Euclidean local branch for which the same-action Dirac square has the
standard source-covariant form and the source is covariantly constant,
$D_\mu\Phi_T=0$, its cross term vanishes and
\[
 D_{A,\Phi_T}^{\,2}=D_A^2+\Phi_T^2.
\]
Consequently
\[
 \boxed{D_{A,\Phi_T}^{\,2}\ge m_T^2 I},\qquad
 \boxed{0\notin\sigma(D_{A,\Phi_T})}.
\]
This is a finite/local operator-gap theorem.  It does not identify $m_T$ with
an interacting pole mass or prove a whole charged-sector Hamiltonian gap.
\end{theorem}
\begin{proof}
$\Phi_T^2=\operatorname{diag}(TT^\dagger,T^\dagger T)\ge m_T^2I$.
On the stated source-preserving Euclidean branch $D_A^2\ge0$, so the displayed
operator inequality follows.
\end{proof}

\begin{proposition}[Schur/Feshbach transmits a gap but is not first genesis]
\label{p3f:prop:v1118-feshbach-propagation-not-genesis}
For an invertible paired complement
\[
 D_Q=\Delta\begin{pmatrix}0&U\\U^\dagger&0\end{pmatrix},\qquad
 \Delta\ne0,\quad U:\mathbb C^r\to\mathbb C^r\text{ unitary},
\]
and retained couplings $A_L:\mathbb C^r\to\mathbb C^{m}$ and
$A_R:\mathbb C^r\to\mathbb C^{n}$, the odd Schur edge is
\[
 \boxed{T_{\rm Fesh}=-\Delta^{-1}A_LUA_R^\dagger}.
\]
If $A_L,A_R$ are both \emph{full column rank on that same paired
$r$-channel complement}, then $T_{\rm Fesh}$ has rank $r$ and
\[
 \boxed{\sigma_{\min}^{+}(T_{\rm Fesh})\ge
 |\Delta|^{-1}\sigma_{\min}(A_L)\sigma_{\min}(A_R)>0,}
\]
where $\sigma_{\min}^{+}$ denotes the smallest \emph{nonzero} singular
value. This is not a lower bound on all singular values when the retained
spaces are larger than the paired complement. Full \emph{row} rank alone
or misaligned coupling ranges does not imply a positive gap. If one
instead seeks $\sigma_{\min}(T_{\rm Fesh})>0$ on a prescribed retained
sector, its domain and codomain must be identified with a compatible
full-rank restriction of the paired range.
Because $D_Q^{-1}$ is already required, this construction propagates or
repackages an existing invertible complement; it is not a theorem of first
source/gap genesis from a source-free owner.
\end{proposition}
\begin{proof}
On the paired $r$-channel complement the Schur inverse is defined because
$U$ is unitary and $\Delta\ne0$. With full-column-rank couplings, take
reduced polar decompositions $A_L=V_LL_L$ and $A_R=V_RL_R$, where
$V_L,V_R$ are isometries and $L_L,L_R$ are invertible positive
$r\times r$ matrices. Hence the nonzero singular values of
$A_LUA_R^\dagger=V_L(L_LUL_R)V_R^\dagger$ are those of $L_LUL_R$.
For each unit $x\in\mathbb C^r$,
$\|L_LUL_Rx\|\ge\sigma_{\min}(A_L)\sigma_{\min}(A_R)$.
This proves the rank and positive nonzero-singular-value bound; excess
orthogonal retained directions carry zero singular values.
\end{proof}

\begin{proposition}[Finite odd-source compression has odd coupling legs]
\label{p3f:prop:primitive-odd-complement-grading}
Let $D$ be odd under $\Gamma$ and let orthogonal projectors $P,Q$
commute with $\Gamma$.  Define $B=PDQ$, and let $\Gamma_P,\Gamma_Q$
be the restricted gradings. Then
\[
 \Gamma_P B\Gamma_Q=-B.
\]
A grading-even retained--complement coupling, written in
chirality order as $\mathrm{diag}(A_L,A_R)$, cannot be obtained
from such a compression of $D$ unless it vanishes.  This
restricts direct compressions of the primitive finite odd Dirac
member; it does not forbid a separately licensed even kinetic
coupling, a distinct total grading, or a grading-odd complement.
\end{proposition}
\begin{proof}
$\Gamma_P B\Gamma_Q=P\Gamma D\Gamma Q=-PDQ=-B$.
\end{proof}


\begin{corollary}[Balanced Schur return]
\label{p3f:cor:v1118-balanced-schur-return}
If $D_i=E_i\widehat D_i$ and
$B_{ji}=\sqrt{E_iE_j}\,\widehat B_{ji}$ on an invertible eliminated block,
then
\[
 \boxed{-B_{ji}D_i^{-1}B_{ji}^\dagger
 =-E_j\widehat B_{ji}\widehat D_i^{-1}\widehat B_{ji}^\dagger.}
\]
Thus balanced midpoint weights are response coefficients whose Schur return
lands at the endpoint scale; they need not be particle carriers.
\end{corollary}

\begin{remark}[Source, running operator and pole are different-scale objects]
\label{p3f:rem:v1118-source-running-pole-scale-separation}
The finite source member, its RG/threshold-transported representative, the
interacting inverse two-point response and a physical pole are distinct typed
objects.  The framework does not impose an equality between them.  Their
relations require the generated transport/matching and pole certificates.
This is the fermionic specialization of Paper~II-A's separated-scale kernel
elimination and scale-matching firewall.
\end{remark}


\subsubsection*{Rooted Feshbach family-return factorization}
\label{p3f:sec:v1119-rooted-feshbach-family-return}

\begin{proposition}[Left family return of a paired Feshbach descendant]
\label{p3f:prop:v1119-left-family-return}
On the domain of Proposition~\ref{p3f:prop:v1118-feshbach-propagation-not-genesis},
put
\[
 R_s=-\Delta_s^{-1}A_{L,s}U_sA_{R,s}^\dagger,
 \qquad G_{R,s}:=A_{R,s}^\dagger A_{R,s}.
\]
Then the positive left-family return is
\[
 \boxed{
 H_s:=R_sR_s^\dagger
 =|\Delta_s|^{-2}A_{L,s}U_sG_{R,s}U_s^\dagger A_{L,s}^\dagger .}
\]
Consequently the right polar unitary of $A_{R,s}$ is presentation data and
drops out of $H_s$; its positive Gram $G_{R,s}$ is the physical right-input
data seen by the left quotient.
\end{proposition}

\begin{corollary}[Orientation-visibility criterion]
\label{p3f:cor:v1119-orientation-visibility}
If $G_{R,s}=\lambda_s I$, or more generally if
$[U_s,G_{R,s}]=0$, then the rooted orientation carried by $U_s$ disappears
from the middle factor $U_sG_{R,s}U_s^\dagger$.  Hence a nontrivial
orientation-dependent family frame requires a same-action noncommuting
rooted Gram/transport pair.  This is a necessary condition, not a sufficient
member-selection theorem.
\end{corollary}

\begin{remark}[Family response is not one scalar contrast]
\label{p3f:rem:v1119-family-response-not-one-scalar}
A response of the form $Q_\perp+rQ_0$ supplies only one positive real
contrast parameter per sector.  Once two three-point spectra are fixed,
generic quark flavor has four real relative-frame coordinates.  Therefore
one such ratio in the up sector and one in the down sector cannot by
themselves generate generic CKM data.  Additional rooted noncirculant
same-action response directions are required.
\end{remark}

\subsubsection*{Antiprism frame handoff and same-action scope}
\label{p3f:sec:v1120-antiprism-frame-handoff}

\begin{remark}[Canonical noncentral frame input without occurrence promotion]
\label{p3f:rem:v1120-antiprism-frame-input}
Paper~I-C supplies the full-rank odd antiprism primitive $\mathcal D_A$
and its exact paired-Feshbach realization.  On any physical branch that
separately proves occurrence of the corresponding polar blocks, the left
return of Proposition~\ref{p3f:prop:v1119-left-family-return} may consume
those blocks as a noncentral family-frame input.  The mathematical I-B theorem
does not by itself establish their occurrence in the finite fermion action.
\end{remark}

\begin{corollary}[Physically sufficient positive-return target]
\label{p3f:cor:v1120-positive-return-target}
For charged masses and left-handed family mixing, once an action-owned sector
singular-value matrix $\Sigma_s$ and left frame $U_s$ are co-owned by one
sector source, the physically sufficient positive return is uniquely
\[
 \boxed{H_s=U_s\Sigma_s^2U_s^\dagger.}
\]
The right polar unitary of a representative $R_s$ is not required by these
readouts unless a later interaction explicitly consumes it.  This uniqueness
statement does not prove the required same-action co-ownership.
\end{corollary}


\subsection*{Rank-seven local-response normalization}
\begin{corollary}[N7JET closure]
\label{p3f:cor:v1124-n7jet-closure}
The chain
\[
 n_7^{\rm sec}
 \longrightarrow \chi_7^{\rm jet}
 \longrightarrow \chi_{2\mid7}=D_2\chi_7^{\rm jet}
 \longrightarrow (F,f,P)
\]
is a same-action generated chain with no extra switching function, response
owner, or phenomenological damping coefficient.  The normalized unresolved
trace $\nu_7$ remains a distinct visibility observable.  Accordingly
\[
 \boxed{\mathrm{N7JET}=\texttt{CLOSED\_ANALYTIC}.}
\]
The identification $f_{RR}\propto\nu_7$ does not follow from the rank-seven response.
\end{corollary}


\subsection*{Root-state typing and terminal scalaron Schur decoupling}

\begin{theorem}[QROOT cannot descend from the classical fresh-root datum alone]
\label{p3f:thm:v1125-qroot-typing-no-go}
On the classical-only fresh-root sub-signature, the conditions
\[
 H_0=\dot F_0=0,\qquad \rho_0=p_0=0
\]
do not by themselves imply
\(\langle T^{\rm exc}_{\mu\nu}\rangle_0=0\).
The reason is typed: the former data belong to the generated classical
homogeneous Cauchy/source role, whereas
\(T^{\rm exc}_{\mu\nu}[\omega]\) consumes a quantum state and a licensed
quantum/effective-action constructor.  The claim-complete-action theorem does
not create that state, and the bare/effective-action firewall forbids deriving
an occurring quantum member from the classical action merely because its
classical source vanishes.  Consequently the classical tuple by itself supplies no QROOT state map.  This is a classical-only typing no-go, not an obstruction to closure by an already licensed quantum constructor co-realized with the fresh-root branch.
\end{theorem}

\begin{proof}
Definition~\ref{p3f:def:v2105-vacuum-excitation-source-split} makes the
excitation source explicitly state dependent.  Paper~III-A's
bare/effective-action firewall states that a generated microscopic action does
not by itself generate a quantum effective action, measure, vacuum, or
occurring stationary member.  Hence the classical root tuple contains no
map whose target is the quantum state sort.  Inferring a particular
\(\omega\) from the vanishing classical fluid source would therefore add a
cross-sort selector absent from the frozen signature.  This is precisely the missing cross-sort step isolated by
Definition~\ref{p3f:def:v2105-qroot}; it cannot be supplied by the classical
tuple alone.
\end{proof}

\begin{theorem}[Fresh-root neutral-constructor closure of QROOT]
\label{p3f:thm:v1128-qroot-neutral-constructor-closure}
On Paper~VI-A's occurring native complete-cap cyclic lineage, restrict QROOT
to the generated primordial scalar/tensor perturbation source family.  The
fresh daughter is a new branch origin with no inherited parent quantum state
or parent-anchored quantum correlation and no independently generated daughter
state/preparation/boundary-process insertion.  Paper~IV's licensed Fock
constructor therefore supplies its neutral degree-zero unit state.  The
same-action root perturbation Hessian is strictly positive and its
vacuum-subtracted excitation operator is
\[
 H_{{\rm exc},0}=\sum_A\hbar\omega_{A0}N_A\succeq0,
\]
so the neutral state has zero excitation.  Since the excitation Hilbert source
of Definition~\ref{p3f:def:v2105-vacuum-excitation-source-split} is obtained
from the same claim-complete perturbation action, Paper~IV's positive-family
ground-source compression gives
\[
 \boxed{\langle T_{\mu\nu}^{\rm exc}\rangle_0=0.}
\]
Hence on this native cyclic perturbation branch
\[
 \boxed{\mathrm{QROOT}=\texttt{CLOSED\_ANALYTIC}.}
\]
This result neither fixes \(\mathrm{REN}_0\) nor infers ground support for an
auxiliary ambient parent Hilbert/Fock extension.  On the native cyclic lineage
Paper~VI-A closes the physically retained parent requirement at the complete
cap, where the parent-relative operational algebra is trivial; the stronger
ambient parent-ground statement is claim-scoped and is not a QROOT dependency.
No repeated vacuum rule is introduced at later spectral frontiers.
\end{theorem}

\begin{proof}
The no-insertion fresh-root chronology removes every nonneutral state
preparation word while retaining the already licensed Fock unit.  Paper~IV's
neutral branch-origin theorem therefore gives the degree-zero state, and
Paper~VI-A's positive root oscillator theorem identifies its number operators
and strict excitation gap.  The resulting state is ground-supported with zero
excitation.  The same finite positive family persists on a neighbourhood of
the root, so the ground compression of every covered off-shell source
variation vanishes.  This is exactly the source statement in QROOT.  The
classical-only theorem above remains true because the quantum constructor is
an additional licensed word, not an inference from \(\rho_0=p_0=0\) alone.
\end{proof}

\begin{theorem}[Terminal scalaron Schur-memory decoupling in the Legendre field]
\label{p3f:thm:v1125-terminal-scalaron-schur-decoupling}
On the restored active branch let \(F=f_R\), \(\chi=f_{RR}>0\), and use the
Legendre scalar-tensor representative
\[
 S[g,F]\supset\int\sqrt{-g}\,[F R-\mathcal V(F)],
 \qquad \mathcal V(F)=R(F)F-f(R(F)).
\]
Then
\[
 \boxed{\mathcal V''(F)=\frac{dR}{dF}=\chi^{-1}.}
\]
On the regular terminal collar of the generated same-action branch,
\(F\to1\) and \(\chi\to0^+\).  The following operator-norm conclusion
is on a fixed retained block, or on a normed realization with uniformly
controlled mixed blocks and inverse bounds; finite differential
coefficients alone do not imply uniformity over an unbounded refinement.  The metric--Legendre mixed Hessian block is
finite there because the only principal mixed owner is the term \(FR\),
linear in the finite field \(F\), together with the already regular boundary
completion.  Hence the scalaron interior block has inverse norm
\(\|C_F^{-1}\|=O(\chi)\), while \(\|B_F\|=O(1)\), and therefore
\[
 \boxed{
 B_F^{*}C_F^{-1}B_F\longrightarrow0
 \qquad(R\downarrow C_7).}
\]
Thus, for the terminally stiff scalaron block, same-action stationary Schur
reduction commutes with the terminal quotient and no finite scalaron Schur memory
survives.
This closes the scalaron-memory clause of G28 on the regular collar; complete
packet matching, common-parent lift completeness, nonempty handoff support,
and actual crossing occurrence remain separately typed.
\end{theorem}

\begin{proof}
The first identity is the ordinary Legendre derivative:
\(\mathcal V'(F)=R(F)\).  The restored N7JET theorem gives
\(\chi\to0^+\), hence \(\mathcal V''\to+\infty\).  In the regular
\((g,F)\) collar chart the mixed second variation of \(FR\) and its generated
boundary completion has finite coefficients at \(F=1\).  The positive
potential stiffness therefore dominates the scalaron interior block, so its
inverse is \(O(\chi)\).  The displayed Schur estimate follows from
\(\|B^*C^{-1}B\|\le\|B\|^2\|C^{-1}\|\).  Applying
Theorem~\ref{p3f:thm:v167-terminal-schur-order} gives the quotient statement.
\end{proof}
\begin{proposition}[Terminal Legendre coefficient Schur slope]
\label{p3f:prop:v245-terminal-legendre-schur-slope}
On the restored branch of
Theorem~\ref{p3f:thm:v263-dressed-born-legendre-descent}, fix
$C=C_7$ and $D=U+C$. Let $\mathcal L_{\rm Leg}(R,F)
=\kappa(FR-\mathcal V(F))$ with fixed $\kappa>0$ and
$F=f_R(R)$, and take \emph{coefficient-chart} variations in $R,F$.
On $R>C$ the Hessian and its stationary $F$ Schur reduction are
\[
 \boxed{H_{R,F}=\kappa\begin{pmatrix}0&1\\1&-\chi^{-1}
 \end{pmatrix},\qquad K_R=\kappa\chi_{2\mid7}(R).}
\]
Because $D_2=1$ on the terminal collar,
\[
 \boxed{K_R(C^+)=0,\quad K_R'(C^+)
 =\frac{2\kappa}{C(U+C)},\quad
 f_{RRR}(C^+)=\frac{2}{C(U+C)}.}
\]
This is the exact stationary \emph{scalar coefficient} response of the
already specified Legendre action. It is not by itself the full
metric/boundary Hessian of $\int\sqrt{-g}\,f(R)$, which contains
measure, tensor and curvature-pullback contacts.
\end{proposition}
\begin{proof}
At a stationary $F$, $\mathcal V'(F)=R$ and
$\mathcal V''(F)=\chi^{-1}$. Schur elimination of the lower-right
pivot $-\kappa\chi^{-1}$ gives
$K_R=-\kappa(1)(-\chi)(1)=\kappa\chi$.
From Eq.~\eqref{p3f:eq:v263-n7-local-jet},
$\chi(R)=2(R-C)/(C(U+C))+O((R-C)^2)$ on the lower
$D_2=1$ collar, giving the two right endpoint derivatives.
\end{proof}


\section{Preterminal finite-source and charge-graded boundary response}
\label{p3f:sec:v1131-preterminal-source-chronology}
The positive-return handoff does not authorize reconstruction of a finite
source from independently terminalized scalar descendants.  The finite member
must first exist in its action owner.

\begin{proposition}[Action first, readout second]
\label{p3f:prop:v1131-action-before-terminal-readout}
Let a rooted constructor independently generate a finite chiral member
$D_s:F_{R,s}\to F_{L,s}$ as a required output of the occurring action, and
suppose $D_s$ contains owner-resolved channels with a common receiver.  Then
the uninserted jet and the left positive return consume the complete $D_s$,
including every generated cross term.  Any terminal map
$q(D_s)$ used only for masses, support dimensions, singular values or other
readouts is downstream:
\[
 \boxed{D_s\longrightarrow D_sD_s^\dagger\longrightarrow q(D_s).}
\]
Conversely, a list of separate terminal values $q(D_{s,r})$ does not determine
$D_s$ unless an injective reconstruction theorem is supplied.
\end{proposition}

\begin{proof}
Apply the existing root-generated fermion-member handoff and the positive
return target to the complete finite member.  Paper II-A's preterminal
relation-preservation theorem forbids commuting a noninjective terminal
readout through a common-parent contraction without a relation-faithfulness
certificate.
\end{proof}

\begin{remark}[Scope]
\label{p3f:rem:v1131-preterminal-scope}
This proposition fixes chronology only.  It does not generate the still-open
charged noncentral finite member; that occurrence remains a Paper-V source
claim.
\end{remark}

\subsection{charged preterminal source-value handoff}
\label{p3f:sec:v1132-charged-preterminal-handoff}
On the independently certified hierarchy-arrow domain
$\mathsf{Cert}_{\rm hier}$, Paper~V-A,
Theorem~\ref{p5:thm:v397-charged-preterminal-srcval-realization},
assembles the charged preterminal finite odd member itself.
The member-handoff theorem applies on that same domain; assembly does not
supply the independent hierarchy-arrow origin certificate.

\begin{corollary}[Charged preterminal member enters the local jet without a second occurrence]
\label{p3f:cor:v1132-charged-preterminal-member-handoff}
Assume $\mathsf{Cert}_{\rm hier}$ of
Definition~\ref{p5:def:v228-hierarchy-origin}, and let $D_{F,s}^{\rm ch}$
be the charged finite odd source value of
Theorem~\ref{p5:thm:v397-charged-preterminal-srcval-realization}
on its occurring rooted action domain.
Then the local clocked uninserted fermion jet is
\[
 \boxed{
 \mathfrak J_F(D_{F,s}^{\rm ch})
 = (\operatorname{Owner},\operatorname{Occ},
 c_FI_L,c_FD_{F,s}^{\rm ch},c_FI_R),
 \qquad
 \widehat D_{F,s}^{\rm ch}=D_{F,s}^{\rm ch}.}
\]
No duplicate source-occurrence token or wavefunction-normalization coefficient
is introduced.
\end{corollary}

\begin{proof}
The member is the required output of a licensed finite geometric constructor
in the same finite odd source owner consumed by
Corollary~\ref{p3f:cor:v1122-root-generated-fermion-member-handoff}.  Apply
that corollary and Theorem~\ref{p3f:thm:v177-local-clocked-fermion-jet}.
\end{proof}

\begin{corollary}[Action chronology after charged closure]
\label{p3f:cor:v1132-charged-action-chronology}
On the same certified-packet domain, the action consumes the complete
mass-labelled incidence member before any pole, RG, terminal-support or other noninjective comparison readout.  The
finite charged masses carried by the member are not regenerated from the
spectrum of a later family return.  Pole and running-mass statements retain
their existing downstream certificates.
\end{corollary}


\subsection{direct-mass source handoff and spectral non-regression}
\label{p3f:sec:v1133-direct-mass-source-handoff}

\begin{corollary}[Finite source coefficients consume, but do not regenerate, terminal masses]
\label{p3f:cor:v1133-direct-mass-source-handoff}
Let a finite fermion source coefficient be generated downstream from a
role-tagged packet $((m_r^{\rm fin},\iota_r))_r$, for example
\[
 R_s^{\rm ch}=\sum_r m_{s,r}^{\rm fin}\iota_{s,r}.
\]
The local uninserted action handoff consumes this coefficient together with
its generated kinetic/current data.  Canonical normalization and later
positive returns may expose spectral information of the action coefficient,
but they do not redefine the already generated $m_{s,r}^{\rm fin}$.  Any
identification of a later eigenvalue or singular value with a terminal mass is
a separately typed comparison/pole theorem.
\end{corollary}

\begin{remark}[Meaning of ``Yukawa'' in the local jet]
Within the current direct-mass branch, ``Dirac/Yukawa member'' names the
zeroth-order chiral interaction/source coordinate carried by the finite
operator slot.  It is not a statement that native masses are generated as
singular values of $vY/\sqrt2$.
\end{remark}


\subsection{successor-cell Schur/Cayley response}
\label{p3f:sec:v1134-successor-cell-schur}

\begin{theorem}[Successor-cell Schur/Cayley response]
\label{p3f:thm:v1134-successor-cell-cayley}
Assume the physical boundary hypotheses of the source-conditioned Schur
descendant theorem and let a generated successor cell carry normalized residual
$0\le\rho\le1$.  Suppose the positive boundary spectral root on the declared
response block is $G>0$ and exact stationary gluing gives the two-boundary
Gaussian family.  Then the one-boundary reduced response is
\[
 \boxed{K_\partial(\rho)=G\frac{1-\rho}{1+\rho}}.
\]
With
\[
 \eta:=G^{-1}K_\partial,
\]
\[
 \boxed{\eta=\frac{1-\rho}{1+\rho}}.
\]
If the cell thresholds are $b_j<b_{j+1}$ and
$\rho=(b_{j+1}-K)/(b_{j+1}-b_j)$, then
\[
 \boxed{\eta_j(K)=
 \frac{K-b_j}{2b_{j+1}-b_j-K}}.
\]
For $b_j=8j$ this is
\[
 \boxed{\eta_j(K)=\frac{K-8j}{8(j+2)-K}}.
\]
\end{theorem}

\begin{corollary}[Hyperbolic boundary flow]
\label{p3f:cor:v1134-hyperbolic-flow}
Writing $\rho=e^{-\ell}$ gives
\[
 \boxed{\eta=\tanh\frac{\ell}{2}},
 \qquad
 \boxed{2\frac{d\eta}{d\ell}=1-\eta^2}.
\]
For a fixed normalized boundary datum $q$,
\[
 S_\partial^{(2)}[q;\rho]
 =\eta(\rho)S_\partial^{(2)}[q;0].
\]
Hence the normalized quadratic boundary-response fraction is exactly $\eta$.
\end{corollary}

\begin{remark}[Stationary-descendant firewall]
\label{p3f:rem:v1134-stationary-descendant-firewall}
The Schur descendant is evaluated after stationary interior substitution.
When an independently typed physical theorem licenses this reduced object, it
is a representation of the same retained interior mode.  Evolving a separately
normalized duplicate in addition to the exact parent would double-count the
same occurrence unless a second owner were independently generated.  The exact
parent remains the physical background history when its Cauchy evolution is
retained.
\end{remark}


\subsection{nested internal boundary response}
\label{p3f:sec:v1135-nested-internal-boundary-response}

\begin{theorem}[Nested stationary boundary descent]
\label{p3f:thm:v1135-nested-stationary-boundary-descent}
Let one occurring stationary action have a frozen tangent decomposition
\[
 T_{\bar q}\mathcal C_{\rm stat}=J_P\oplus J_B\oplus J_I,
\]
where $J_I$ is the stationary interior, $J_B$ is a retained nonparticle
boundary/background block and $J_P$ is a separately declared retained block.
Assume all required inverse blocks exist on the declared gauge/constraint
quotient.  Then Schur reduction is associative:
\[
 \boxed{
 \operatorname{Schur}_{J_B}\!\left(
   \operatorname{Schur}_{J_I}D^2S_\star\right)
 =
 \operatorname{Schur}_{J_B\oplus J_I}D^2S_\star.}
\]
Consequently the quadratic response on $J_P$ is independent of whether the
true interior and the remaining boundary/background variables are eliminated
in one step or in the displayed nested order.
\end{theorem}

\begin{proof}
Order the Hessian as $J_P\oplus J_B\oplus J_I$,
\[
H=\begin{pmatrix}A&B&C\\B^\dagger&D&E\\C^\dagger&E^\dagger&F\end{pmatrix}.
\]
Eliminating $J_I$ gives
\[
 \widehat A=A-CF^{-1}C^\dagger,\quad
 \widehat B=B-CF^{-1}E^\dagger,\quad
 \widehat D=D-EF^{-1}E^\dagger.
\]
The nested response is
\[
 K_P=\widehat A-\widehat B\widehat D^{-1}\widehat B^\dagger.
\]
The standard block-inverse formula for
$\begin{psmallmatrix}D&E\\E^\dagger&F\end{psmallmatrix}$ gives identically
\[
 K_P=A-\begin{pmatrix}B&C\end{pmatrix}
 \begin{pmatrix}D&E\\E^\dagger&F\end{pmatrix}^{-1}
 \binom{B^\dagger}{C^\dagger},
\]
which is the direct Schur complement by $J_B\oplus J_I$.
\end{proof}

\begin{corollary}[Internal particle-response realization]
\label{p3f:cor:v1135-internal-particle-response}
If $J_P$ is independently licensed as a physical particle/one-particle field
carrier, then
\[
 \boxed{K_P=\operatorname{Schur}_{J_B\oplus J_I}D^2S_\star}
\]
is its unique quadratic internal same-action boundary response.  On a regular
branch,
\[
 \boxed{K_P^{-1}=(H^{-1})_{PP}},
\]
so isolated zeros of $K_P$ are full-action poles visible through that retained
particle channel.  This statement identifies the response carried by an
already constituted particle carrier; it does not assign particle semantics to
an arbitrary spectral block or complete boundary-energy packet.
\end{corollary}

\begin{corollary}[Outer-history variation of an internal response]
\label{p3f:cor:v1135-history-variation-internal-response}
Let $\eta$ be a genuine source/history/boundary argument of the same reduced
action.  Wherever the preceding reductions are differentiable,
\[
 \boxed{V_{\eta PP}:=D_\eta K_P=D_\eta D_P^2S_\star^\partial.}
\]
Thus motion of an outer/history boundary changes the internal particle response
through the mixed third action variation.  A subsequent Bogoliubov or kinetic
occupation statement remains state/history dependent.
\end{corollary}

\subsection{fixed-datum support preservation}
\label{p3f:sec:v1136-fixed-datum-support}

\begin{corollary}[Support preservation for a fixed successor-cell datum]
\label{p3f:cor:v1136-fixed-datum-support}
On the domain of Theorem~\ref{p3f:thm:v1134-successor-cell-cayley}, let the
normalized endpoint quadratic form be represented on the declared boundary
tangent by $K_0$, and let $P$ be an owner-preserving projector onto a retained
physical response block.  If
\[
 K_0=PK_0P,
\]
then the complete successor-cell family satisfies
\[
 \boxed{K_\partial(\rho)=\eta(\rho)K_0
       =P K_\partial(\rho)P.}
\]
Hence the cell cannot generate an orthogonal response target from a fixed
boundary datum.  If $P$ is an independently constituted particle block, the
whole fixed-datum cell is an internal particle-boundary response on that block.
\end{corollary}

\begin{proof}
Corollary~\ref{p3f:cor:v1134-hyperbolic-flow} gives scalar multiplication of
the same normalized quadratic datum throughout the cell.  Multiplying
$K_0=PK_0P$ by the scalar $\eta(\rho)$ gives the displayed identity.  No
additional Schur block, target sort or normalization is created.  When the
particle block itself is obtained after stationary Schur reduction, associativity
of Theorem~\ref{p3f:thm:v1135-nested-stationary-boundary-descent} shows that
the same statement is independent of the order in which the true interior and
retained nonparticle/background block are eliminated.
\end{proof}

\begin{corollary}[Direct-cell current has no independent realization factor]
\label{p3f:cor:v1136-direct-cell-no-efficiency}
If the fixed datum of the preceding corollary has a generated physical energy
normalization $\mathcal I_\partial$ and completion coordinate $\eta(t)$, then
the direct same-owner packet current is
\[
 \boxed{Q_\partial=\mathcal I_\partial\dot\eta.}
\]
An additional factor multiplying this direct packet is not generated by the
successor-cell/Schur construction.  State-dependent Bogoliubov pumping of
other comparison branches remains a distinct dynamical readout and is not
identified with this same-owner packet current by this corollary.
\end{corollary}

\subsection{charge-graded internal response and history no-go}
\label{p3f:sec:v1137-charge-graded-internal-response}

\begin{theorem}[Exact-charge preservation under stationary Schur descent]
\label{p3f:thm:v1137-charge-preserving-schur}
Let a regular stationary quadratic action have tangent decomposition
\[
 J_P\oplus J_R,
 \qquad
 H=D^2S_\star=
 \begin{pmatrix}A&B\\B^\dagger&C\end{pmatrix},
\]
where $J_P$ is a retained physical carrier and $C$ is invertible on the
stated gauge/constraint quotient.  Let an exact additive charge act block
diagonally,
\[
 Q=Q_P\oplus Q_R,
\]
and assume $[Q,H]=0$.  Then the retained stationary response
\[
 K_P=A-BC^{-1}B^\dagger
\]
obeys
\[
 \boxed{[Q_P,K_P]=0.}
\]
Thus stationary Schur descent cannot manufacture a charge-violating
quadratic block absent from the parent Hessian.
\end{theorem}

\begin{proof}
The commutator hypothesis gives
$Q_PA=AQ_P$, $Q_PB=BQ_R$, $Q_RB^\dagger=B^\dagger Q_P$ and
$Q_RC=CQ_R$.  Hence $Q_RC^{-1}=C^{-1}Q_R$, and therefore
\[
\begin{aligned}
 Q_PK_P
 &=AQ_P-BC^{-1}Q_RB^\dagger\\
 &=AQ_P-BC^{-1}B^\dagger Q_P
 =K_PQ_P.
\end{aligned}
\]
\end{proof}

\begin{corollary}[History differentiation preserves an exact charge sector]
\label{p3f:cor:v1137-charge-preserving-history-variation}
Let $K_P(\eta)$ be the internal response of
Corollary~\ref{p3f:cor:v1135-history-variation-internal-response}, with
$Q_P$ independent of the genuine history/source coordinate $\eta$.  If
$[Q_P,K_P(\eta)]=0$ throughout an interval, then for every differentiability
order admitted on that interval,
\[
 \boxed{[Q_P,D_\eta^nK_P]=0\qquad(n\ge1).}
\]
Equivalently every nonzero-charge projection of $D_\eta^nK_P$ vanishes.
A changing boundary profile can modulate an already occurring charged
interaction, but cannot create a missing charge sector from a charge-preserving
parent action.
\end{corollary}

\begin{corollary}[Neutral \texorpdfstring{$B-L$}{B-L} no-go on the current Dirac branch]
\label{p3f:cor:v1137-bl2-dirac-branch-no-go}
On the active Paper~V-A Dirac-neutrino branch with no constant odd/Majorana
offset, let $Q_{B-L}$ be the generated $B-L$ charge and let $J_\nu$ be the
already constituted neutral one-particle carrier.  Then
\[
 [Q_{B-L},D^2S_\star]=0
\]
and hence the internal neutral response satisfies
\[
 \boxed{\Pi_{\Delta(B-L)=\pm2}K_\nu=0},
\qquad
 \boxed{\Pi_{\Delta(B-L)=\pm2}D_\eta^nK_\nu=0\quad(n\ge1).}
\]
This is a branch statement, not an impossibility theorem for the framework:
a separately generated charge-two action owner would change the parent
Hessian and lie outside the hypothesis of the corollary.
\end{corollary}

\subsection{co-realized charge-two neutral boundary response}
\label{p3f:sec:v1139-bl2-boundary-pair-import}

\begin{corollary}[Charge-two K16 boundary block on the co-realized pair branch]
\label{p3f:cor:v1139-bl2-boundary-pair-block}
On the K16 neutral branch of Paper~V-A,
Theorem~\ref{p5:thm:v3105-k16-bl2-pair-corealization}, the occurring physical
boundary quadratic owner contains the nonzero pair cell
\[
 \bigl(\Delta_{BL2}^{\rm geom},
       (\Delta_{BL2}^{\rm geom})^\dagger\bigr)
\]
with charge weights $+2$ and $-2$.  Therefore the boundary quadratic response
has a nonzero charge-two projection,
\[
 \boxed{
 \Pi_{\Delta(B-L)=\pm2}K_{\nu,\partial}\ne0}
\]
on that branch.  Corollary~\ref{p3f:cor:v1137-bl2-dirac-branch-no-go}
remains exact on the Dirac-only comparison branch, but its hypothesis
$[Q_{B-L},D^2S_\star]=0$ is not satisfied by the co-realized pair owner.
\end{corollary}

\begin{corollary}[Fixed-owner transport of the BL2 pair cell]
\label{p3f:cor:v1139-bl2-fixed-owner-transport}
Let $K_{0}^{BL2}$ denote the normalized charge-two component of the fixed K16
boundary datum.  On the same successor-cell domain as
Corollary~\ref{p3f:cor:v1136-fixed-datum-support},
\[
 \boxed{
 K_{\partial}^{BL2}(\rho)=\eta(\rho)K_{0}^{BL2}.}
\]
No independent BL2 realization efficiency is generated.  The statement
transports an already occurring pair block; it does not derive a CP-odd phase
history or a net charge expectation.
\end{corollary}

\subsection{fixed-owner K48 H response and normalization jet}
\label{p3f:sec:v1141-k48-h-response-jet}

\begin{corollary}[Homogeneous K40--K48 fixed-owner specialization]
\label{p3f:cor:v1141-k48-fixed-owner-specialization}
Let the normalized nonzero K48 endpoint boundary datum be independently owned
by the homogeneous $H$ modality.  For the successor cell $40\le K\le48$,
\[
 \rho_{40\to48}=\frac{48-K}{8},
 \qquad
 \eta_{40\to48}=\frac{1-\rho_{40\to48}}{1+\rho_{40\to48}}
 =\frac{K-40}{56-K}.
\]
Fixed-owner inheritance then gives
\[
 \boxed{K_{\partial,40\to48}^{H}(K)
 =\eta_{40\to48}(K)K_{0,48}.}
\]
No second $H$ owner and no realization-efficiency coefficient are generated by
this transport.
\end{corollary}

\begin{definition}[K48 homogeneous H-boundary response jet]
\label{p3f:def:v1141-k48-h-response-jet}
Let $S_{\partial,40\to48}^{H,\rm red}$ be the complete same-action reduced
boundary action after all licensed stationary interior Schur reductions and after
co-consumption of the complete terminal quadratic relation.  Define
\[
 \boxed{
 \mathfrak J_{48}^{H}:=
 \frac{-2u^\mu u^\nu\,
 \delta S_{\partial,40\to48}^{H,\rm red}/
 (\sqrt{-g}\,\delta g^{\mu\nu})}{I_{48}}\bigg|_{K48}.}
\]
The dimensionless K48 background normalization is
\[
 \boxed{\gamma_{48}:=\mathfrak J_{48}^{H}.}
\]
\end{definition}

\begin{remark}[Only the response jet remains open]
\label{p3f:rem:v1141-k48-normalization-firewall}
Generated co-consumption closes occurrence of the terminal relation in the
$H$-owned action, but not the numerical value of
$\mathfrak J_{48}^{H}$.  Neither the raw scalar corner $-1/4$, the positive edge
metric $1/4$, a moment ratio $3/4$, the absolute semigroup coefficient $1/7$,
nor a Cayley formula alone determines $\gamma_{48}$.  The chain
$1-1/4\to3/4\to4/3$ becomes a theorem only if the complete reduced action
independently proves that the normalized $H$ Schur self-energy equals $1/4$.
The corresponding source pullback generally yields an action jet such as
$-(1/4)\,\partial_sK_H$, not automatically a finite shift $-(1/4)K_H$.
\end{remark}

\subsection{K48 homogeneous Schur decoupling}
\label{p3f:sec:v1142-k48-homogeneous-schur-closure}

\begin{theorem}[Unbroken homogeneous K48 Schur self-energy vanishes]
\label{p3f:thm:v1142-k48-h-t5-schur-zero}
Work on the physical post-K16 homogeneous K48 branch after the licensed
stationary interior reductions already consumed by
Corollary~\ref{p3f:cor:v1141-k48-fixed-owner-specialization}.  Let $h$ denote
the retained real $H$ coordinate and $v\in T_5$ the unmarked terminal
transverse coordinate.  On the unbroken $C_3$ branch the complete quadratic
block has the invariant form
\[
 \mathcal H_{H\oplus T_5}^{(2)}
 =\begin{pmatrix}
 A_H & B_{H,T}\\
 B_{T,H} & C_T
 \end{pmatrix},
 \qquad
 B_{H,T}\in\operatorname{Hom}_{C_3}(T_5,\mathbb R).
\]
Paper~III-C, Corollary~\ref{p3c:cor:v193-k48-h-t5-mixed-zero}, gives
$B_{H,T}=0$.  Hence on every regular stationary quotient,
\[
 \boxed{
 \Sigma_{H,48}^{\rm Schur}
 :=B_{H,T}C_T^{-1}B_{T,H}=0,
 \qquad
 K_{H,48}^{\rm red}=A_H.}
\]
If $C_T$ is singular, the compatible stationary quotient requires
$B_{H,T}(\ker C_T)=0$, which is automatic here; the induced quotient therefore
has the same retained quadratic form.  No Moore--Penrose inverse or deleted
critical mode can manufacture a nonzero Schur correction from the vanishing
mixed block.
\end{theorem}

\begin{corollary}[The K48 edge quarter is not a homogeneous Schur decrement]
\label{p3f:cor:v1142-k48-quarter-no-schur-shift}
On the branch of Theorem~\ref{p3f:thm:v1142-k48-h-t5-schur-zero},
\[
 \boxed{\Sigma_{H,48}^{\rm Schur}/A_H=0.}
\]
Accordingly the geometric identity
$G_{48}^{\rm edge}=I_{T_5}/4$ cannot yield
$K_H^{\rm red}=(3/4)A_H$ or an inverse-response factor $4/3$.
The nonlinear terminal pullback may still generate a source--$H$--$H$ jet such
as $-(1/4)\,\partial_sK_H$; this is a higher variation and not an additive
second-variation self-energy.
\end{corollary}

\begin{remark}[Branch scope]
\label{p3f:rem:v1142-k48-c3-branch-scope}
The preceding theorem is authoritative for the current unbroken homogeneous
$C_3$ branch.  A separately generated symmetry-breaking critical branch with a
nonzero transverse stationary member would require its own occurrence,
stability, kinetic-support and Hilbert-variation theorems; no such branch is
inserted by this result.
\end{remark}

\subsection{causal dynamic source--Schur reduction}
\label{p3f:sec:v1144-dynamic-source-schur}

\begin{theorem}[Causal source-conditioned dynamic Schur reduction]
\label{p3f:thm:v1144-causal-dynamic-schur}
Let one occurring reduced action/source branch have linearized dynamical equations
\[
 \begin{pmatrix}
 A^{\rm dyn}&B^{\rm dyn}\\
 D^{\rm dyn}&C^{\rm dyn}
 \end{pmatrix}
 \binom{w}{z}
 =\binom{j_w}{j_z},
\]
on its declared gauge/constraint quotient.  Assume that the source/interior block
$C^{\rm dyn}$ has the physical retarded inverse $(C^{\rm dyn})^{-1}_{R}$ on the
stated causal domain.  Eliminating $z$ gives the unique causal scalar descendant
\[
 \boxed{
 K_R^{\rm dyn}
 =A^{\rm dyn}-B^{\rm dyn}(C^{\rm dyn})^{-1}_{R}D^{\rm dyn},}
\]
\[
 \boxed{
 K_R^{\rm dyn}w
 =j_w-B^{\rm dyn}(C^{\rm dyn})^{-1}_{R}j_z.}
\]
For a variationally symmetric frozen-frequency Hessian one has
$D^{\rm dyn}=(B^{\rm dyn})^*$; in a causal first-order transport representation
$C^{\rm dyn}$ need not be self-adjoint.  The response kernel and transformed
source are both descendants of the same action/source system.  No relaxation
time, sound speed, viscosity, damping coefficient or fitted nonlocal kernel is
created by the Schur step.
\end{theorem}

\begin{corollary}[Collisionless kinetic source block on the active-stress scalar leaf]
\label{p3f:cor:v1144-collisionless-source-schur}
On an admitted local collisionless kinetic branch with canonical Hamiltonian
$H_a$ and occupation $f_a$, set $\Phi_N=\log N=-w/2$ and
\[
 D_a:=\partial_{\Phi_N}H_a,
 \qquad Q_a:=\partial_{\Phi_N}^2H_a.
\]
Let
\[
 \bar{\mathscr L}_a:=\partial_t+\{\,\cdot\,,\bar H_a\}_{x,p}
\]
be the background Liouville operator and let
$(\bar{\mathscr L}_a^{-1})_R$ denote its retarded inverse on the admitted
non-Casimir transport subspace.  The equal-time scalar matter Hessian is
\[
 \chi_{\rm kin}
 =\frac{\kappa_g}{2N\sqrt h}\sum_a\int \bar f_a Q_a\,d^3p.
\]
Linearized collisionless transport gives
\[
 \bar{\mathscr L}_a\,\delta f_a
 =\frac12\{\bar f_a,D_aw\}_{x,p}+q_a^{\rm ind}.
\]
After retarded elimination of $\delta f_a$, the exact causal kinetic source
block on a constant-$F_\star$ Einstein tangent is
\[
 \boxed{
 \mathcal V_{\rm kin}^{R}
 =M_{\chi_{\rm kin}/F_\star}
 -\mathcal B\,\bar{\mathscr L}_{R}^{-1}\mathcal C,}
\]
where
\[
 (\mathcal Cw)_a=\{\bar f_a,D_aw\},
 \qquad
 \mathcal Bg=\frac{\kappa_g}{2F_\star N\sqrt h}
 \sum_a\int D_ag_a\,d^3p.
\]
The independent source column is
\[
 \boxed{
 J_{\rm env}^{R}
 =\frac{\kappa_g}{F_\star N\sqrt h}
 \sum_a\int D_a\,\delta f_a^{\rm ind}\,d^3p.}
\]
All pressure, dispersion and scalar anisotropic-stress responses are momentum
moments of the same retarded $\delta f_a$; they do not require a separately
chosen fluid closure on this branch.
\end{corollary}

\begin{corollary}[Frozen, relaxed and high-frequency kinetic limits]
\label{p3f:cor:v1144-kinetic-limits}
For fixed canonical occupations the transport variable is not eliminated and
\[
 \boxed{\mathcal V_{\rm kin}^{\rm frozen}=\chi_{\rm kin}/F_\star.}
\]
For a stationary collisionless family $\bar f_a=\mathcal F_a(\bar H_a)$, the
zero-frequency retarded limit on the invertible non-Casimir transport subspace is
\[
 \boxed{\mathcal V_{\rm kin}^{R}(0)=\chi_{\rm rel}/F_\star,}
\]
\[
 \chi_{\rm rel}
 =\frac{\kappa_g}{2N\sqrt h}\sum_a\int
 \left[\mathcal F_a(H_a)Q_a+\mathcal F_a'(H_a)D_a^2\right]d^3p.
\]
The second term need not be positive.  For frequencies large compared with the
kinetic transport frequencies,
\[
 \boxed{\mathcal V_{\rm kin}^{R}(\omega)\to\chi_{\rm kin}/F_\star.}
\]
Thus the fixed-source static Hessian and the freely relaxed zero-frequency
response are different tangent problems of one action.
\end{corollary}

\begin{remark}[Scope]
\label{p3f:rem:v1144-dynamic-source-scope}
The collisionless corollaries are branch theorems, not a universal source model.
The constituted finite interacting source block is now supplied by
Corollary~\ref{p3f:cor:v1145-finite-quantum-source-schur}.  Collisional,
open-system, unbounded-field and other source domains retain their own
generated dynamical blocks.  The general Schur theorem remains applicable
once their actual source blocks and source columns are supplied.
\end{remark}


\subsection{Matched native sources: trace, target geometry and response selection}
\label{p3f:sec:v2202-matched-native-sources}

This section records consequences of the existing source constructors;
it selects no scalar section, Yukawa law, state or history. Work in
$\hbar=c=1$, $M^2=U/4$, $C=C_7$, $F=f_R$ and $\chi=f_{RR}$.
The active gravitational representative has $\chi>0$ and
$P(F)=Fr(F)-f(r(F))$. In the displayed native local heat units put
$\alpha_\Phi=(16\pi^2)^{-1}$,
$X=\operatorname{Tr}_F\Phi^2$ and
$V_\Phi=\alpha_\Phi\operatorname{Tr}_F(\Phi^4-2\Lambda_g^2\Phi^2)$.

\subparagraph{Common-action domain.}
The results involving both $F$ and $\Phi$ consume an exact, claim-complete
overlap in common physical units, as required by
Proposition~\ref{p3f:prop:v1164-scalar-overlap-condition}. On that domain
the displayed gravity--scalar constituent is
\begin{equation}
\label{p3f:eq:v2202-matched-packet}
 S_{g\Phi}=\int\sqrt{-g}\left[
 \frac{M^2}{2}(FR-P(F))
 -\alpha_\Phi\operatorname{Tr}_F(D\Phi)^2
 -\frac{\alpha_\Phi}{6}RX-V_\Phi\right]d^4x .
\end{equation}
Every actual residual source, metric and boundary variation is retained.
If the actual overlap transports the scalar packet by a positive weight,
$\alpha_\Phi$ in all formulas denotes its transported action coefficient,
with every field and source map transported consistently. Its native local
value is not an independently adjustable weight or an automatic
raw-to-physical normalization identity.
For a canonical non-null Dirichlet boundary, fix its induced metric,
$F$ and $\Phi$ and retain
$2\sum_\partial\epsilon\int\sqrt{|h|}\,\mathcal A K$, where
$\mathcal A=M^2F/2-\alpha_\Phi X/6$.
This is the completion already specified in
Section~\ref{p3f:sec:v169-minimal-complete-action}.
The source's $R(g)$ is varied before restriction; substituting the bare
$r(F)$ into it is not this matched action.

The raw common-heat Einstein coefficient is $\kappa_J=12U$, whereas
$M^2/2=U/8$; their ratio is $96$. One must not append
$\kappa_JR$ to~\eqref{p3f:eq:v2202-matched-packet} or infer full action
overlap from this coefficient ratio. The generated gravitational heat
branch fixes $\Lambda_g^2=96\pi^2U$, not $\Lambda_g^2=U$.
A uniform positive transport of the entire scalar packet, followed by
its canonical field rescaling, preserves its curvature-to-kinetic and
quadratic-to-kinetic ratios. Its quartic and mass-map coefficients must
also transport. This observation supplies neither an overlap proof nor
a new adjustable weight or heat scale.

\begin{corollary}[Inherited improved trace on the matched scalaron port]
\label{p3f:cor:v2202-improved-trace-handoff}
On the matched domain let $S_r$ have no independently inserted $F$
coupling, and define its complete variations by
$\delta_\Phi S_r=-\int\sqrt{-g}\,\operatorname{Tr}_F(J\delta\Phi)$
and $T_r=g^{\mu\nu}T^r_{\mu\nu}$. The existing complete improved trace
and the nonlinear gravity trace compose to
\begin{align}
 T_\Phi&=4\alpha_\Phi\Lambda_g^2X+
                       \operatorname{Tr}_F(\Phi J),
 \label{p3f:eq:v2202-improved-trace}\\
 \boxed{M^2(3\Box F+FR-2f)
   =4\alpha_\Phi\Lambda_g^2X+
      \operatorname{Tr}_F(\Phi J)+T_r.}
 \label{p3f:eq:v2202-matched-trace-law}
\end{align}
\end{corollary}
\begin{proof}
Use the matter identity of
Proposition~\ref{p3f:prop:native-source-curvature-trace}, and the
gravity-independent handoff in Paper~VI-A,
Proposition~\ref{p8:prop:v2189-c06-classical-scalar-sign}.
The scalar Euler equation cancels its kinetic, curvature and quartic
trace terms, leaving
$\operatorname{Tr}_F(\Phi V_{\Phi,\Phi})-4V_\Phi
 +\operatorname{Tr}_F(\Phi J)$.
Taking the metric trace of the matched $f(R)$ owner gives the result.
\end{proof}
This cancels terms in the trace, not their energy, pressure or metric
response. It is a handoff of the existing matter identity, not a new
trace mechanism or reuse of the local-EH curvature equation.
Functional residual insertions retain their higher jets.
The section $\Phi=0$ remains admissible where $J(0)=0$; its stability
does not follow from that statement.

\begin{corollary}[Actual mass-map Euler selector]
\label{p3f:cor:v2202-mass-map-selector}
For an actually retained classical Dirac mass slot
$\mathcal L_D\supset-\bar\psi\mathsf M(\Phi)\psi$, with its complete
on-shell metric and scalar variations,
\begin{equation}
\label{p3f:eq:v2202-mass-map-selector}
 \boxed{T_D+\operatorname{Tr}_F(\Phi J_D)
 =\bar\psi\bigl(D\mathsf M_\Phi[\Phi]-\mathsf M(\Phi)\bigr)\psi.}
\end{equation}
A constant direct mass contributes $-m\bar\psi\psi$; a homogeneous
degree-$q$ slot contributes $(q-1)\bar\psi\mathsf M\psi$.
\end{corollary}
\begin{proof}
The on-shell Dirac trace is $T_D=-\bar\psi\mathsf M\psi$ in the
printed $(-+++)$ convention. Variation of the same mass term gives
$\operatorname{Tr}_F(\Phi J_D)=\bar\psi D\mathsf M_\Phi[\Phi]\psi$.
Euler's identity gives the homogeneous case.
\end{proof}
Degree-one cancellation is conditional on that actual mass map.
The native direct-mass ``Yukawa'' terminology does not select it
(Remark following Corollary~\ref{p3f:cor:v1133-direct-mass-source-handoff}).
The combined trace does not determine a body's total gravitational
charge or remove its energy. Quantum anomalies consume their own owner.

\begin{theorem}[Matched hyperbolic target and terminal induced geometry]
\label{p3f:thm:v2202-hyperbolic-target}
On an admitted $n\ge1$ dimensional gauge-stable real linear support
$\Phi=h^IT_I$ with $2\alpha_\Phi\operatorname{Tr}_F(T_IT_J)=\delta_{IJ}$,
put $\rho_h^2=h^Ih_I$, $\Omega=F-\rho_h^2/(6M^2)>0$ and
$\psi=\sqrt{3/2}M\log\Omega$. Before gauge quotient the displayed
two-derivative constituent has Einstein variables $g_E=\Omega g$ and
\begin{align}
 d\ell^2&=d\psi^2+\Omega^{-1}\delta_{IJ}dh^Idh^J,
 \qquad K=-\frac1{6M^2},
 \label{p3f:eq:v2202-hyperbolic-target}\\
 W_E&=\frac{M^2P(F)/2+V_\Phi(h)}{\Omega^2},\qquad
 F=\Omega+\frac{\rho_h^2}{6M^2}.
 \label{p3f:eq:v2202-hyperbolic-potential}
\end{align}
For $n$ retained real scalar coordinates the active target has dimension
$n+1$ and scalar curvature $-n(n+1)/(6M^2)$.
With no retained scalar support, $n=0$, the active field space is the
single canonical scalaron line and has no intrinsic sectional curvature.
On the separately licensed terminal quotient $F=1$, its induced metric is
\begin{equation}
\label{p3f:eq:v2202-terminal-target}
 G^{\rm term}_{IJ}=\frac{\delta_{IJ}}\Omega
       +\frac{h_Ih_J}{6M^2\Omega^2},\qquad
 \Omega=1-\frac{\rho_h^2}{6M^2}.
\end{equation}
For $n\ge2$ this is a totally geodesic $H^n$ with the same sectional
curvature; its scalar curvature is $-n(n-1)/(6M^2)$.
For $n=1$ it is a radial line and has no intrinsic sectional curvature.
The terminal metric is the existing local-EH scalar target of
Proposition~\ref{p3f:prop:v118-hyperbolic-mass-gram} after its actual
Einstein coefficient is transported to $M^2/2$; the active $n+1$ target and
its embedding retain the independent nonlinear $F$ coordinate.
\end{theorem}
\begin{proof}
The trace normalization gives kinetic term $-(Dh)^2/2$ and curvature
term $-R\rho_h^2/12$. Weyl transport of $M^2\Omega R/2$, with its
endpoint, gives $-(\partial\psi)^2/2$ and transports the scalar kinetic
term by $\Omega^{-1}$. The warp $e^{-b\psi/2}$, $b=\sqrt{2/3}/M$,
has radial curvature $-f''/f$ and tangential curvature $-(f'/f)^2$,
both $-b^2/4$.
Set $y=\sqrt\Omega$ and $x^I=h^I/(\sqrt6M)$.
The line element is $6M^2(dy^2+dx^Idx_I)/y^2$, while $F=1$ is
$y^2+x^Ix_I=1$, a hemisphere orthogonal to its conformal boundary.
It is totally geodesic and gives the induced metric above.
\end{proof}
Residual velocity insertions retain the complete source-corrected
principal matrix of Proposition~\ref{p3f:prop:complete-source-cauchy-matrix}.
Gauge quotient, stable physical normal modes and their occupations are
not supplied by this target metric. Terminal scalaron removal reduces
the field dimension; it does not delete an independently retained
native scalar support.

\begin{theorem}[Homogeneous free-band source selection]
\label{p3f:thm:v2202-free-band-selection}
On the minimal free Dirac specialization, retain the complete invariant
band of Section~\ref{p3f:sec:v169-minimal-complete-action}.
At fixed Einstein metric, $\mu=m e^{-\phi/c_\phi}$,
$c_\phi=\sqrt6M$, and
\begin{equation}
\label{p3f:eq:v2202-free-band-vertices}
 Q_\phi=-c_\phi^{-1}d\Gamma(\mu\beta_D)
                  +\varepsilon_{,\phi}I,
 \qquad \partial_\phi^k\mu=(-c_\phi^{-1})^k\mu .
\end{equation}
Every homogeneous mass insertion preserves the $D_\gamma^2$
eigenspaces. On each sign pair, $p=|\lambda|/a_E$,
$E^2=p^2+\mu^2$, its energy-axis insertion is
\begin{equation}
\label{p3f:eq:v2202-sign-pair-vertex}
 Q_\phi=-\frac{\mu^2}{c_\phi E}\tau_3
          +\frac{\mu p}{c_\phi E}\tau_1
          +\varepsilon_{,\phi}I.
\end{equation}
The free pair word preserves $n_a-n_b$ and fermion parity; its
homogeneous transverse mass insertion vanishes for $m=0$ or $p=0$.
\end{theorem}
\begin{proof}
The exponential is the existing exact Weyl-transported mass.
$\{D_\gamma,\beta_D\}=0$ implies $[D_\gamma^2,\beta_D]=0$;
the invariant band retains this identity. Rotating
$p\sigma_1+\mu\sigma_3$ to its energy axis gives the displayed insertion.
Both the occupation term and $a^\dagger b^\dagger+ba$ commute with
$n_a-n_b$ and parity.
\end{proof}
Complete spatial-isometry multiplets preserve that symmetry; equivariant
ports have Schur form $I_{\rm irrep}\otimes A_{\rm multiplicity}$, not
an assumed identity on internal multiplicities. Inhomogeneous insertions
have their actual spatial jets and need not preserve $|\lambda|$.
The source band is not the microscopic label $m_p=1$.
Central derivatives, the moving energy-frame connection and the exact
source work of Section~\ref{p3f:sec:v2201-derived-d10} remain retained.
These selection rules create neither a density matrix nor an asymmetry.

\begin{corollary}[Scoped terminal scalar and tensor diagnostics]
\label{p3f:cor:v2202-terminal-diagnostics}
At the affine terminal quotient, on a $\Phi=0$, $J(0)=0$ background,
the complete residual trace $T_r=-\rho_r+3p_r$ gives
$R=2C+(\rho_r-3p_r)/M^2$. The local canonical scalar equation therefore
has mass coefficient
\begin{equation}
\label{p3f:eq:v2202-terminal-local-mass}
 m_{h,\rm loc}^2=\frac C3-2\Lambda_g^2
                     +\frac{\rho_r-3p_r}{6M^2}.
\end{equation}
Covariant spatial eigenvalues and the complete residual scalar
contact/memory are additional terms in its actual linear operator.
For an admitted zero spatial mode with no such extra terms, the zero
of this coefficient is
$\rho_r-3p_r=12M^2\Lambda_g^2-2M^2C$.
On a homogeneous isotropic two-derivative tensor restriction with no
additional tensor principal or memory terms, the same coefficient
$M^2\Omega$ gives
\begin{equation}
\label{p3f:eq:v2202-tensor-friction}
 \ddot\gamma_{ij}+
 \left(3H+\frac{\dot\Omega}{\Omega}\right)\dot\gamma_{ij}
 +\frac{k_T^2}{a^2}\gamma_{ij}=0,
\end{equation}
where $k_T^2$ denotes the actual short-wavelength/closed-harmonic tensor
eigenvalue on that restriction.
\end{corollary}
\begin{proof}
Trace $G_{\mu\nu}+(C/2)g_{\mu\nu}=M^{-2}T^r_{\mu\nu}$ and
linearize the native scalar equation at $\Phi=0$. The tensor velocity
and spatial-gradient terms both carry $a^3M^2\Omega/8$; differentiating
its velocity coefficient gives the stated friction.
\end{proof}
Radiation's vanishing classical trace is not vanishing gravitational
stress. With the same-unit heat scale, the empty local coefficient is
$m_{h,\rm loc}^2/U=C/(3U)-192\pi^2$ and the critical trace divided by
$U^2$ is $288\pi^2-C/(2U)$. This is a local branch-stability diagnostic
conditional on the overlap, not an exhibited late cosmological
instability or a controlled ultraviolet growth rate.
The tensor equation's geometric-optics propagation factor is
$d_L^{\rm GW}/d_L^{\rm EM}=\sqrt{\Omega_0/\Omega_z}$ only with compatible
source/detector and photon-clock maps; it is not an observation fit.
On the unextended terminal $\Phi=0$ branch, $\Omega=1$ and this factor
is unity. No result here selects a stable spectator, a nonzero scalar
background, a new scalaron, or an observed reservoir abundance.
\begin{theorem}[Bare terminal radial obstruction on an actually matched packet]
\label{p3f:thm:v2211-terminal-radial-obstruction}
On the terminal \(F=1\) restriction of the matched packet, suppose the
displayed two-derivative gravity--scalar constituent is the complete
potential for the claim: no residual source, gauge background or
quantum/effective correction changes its radial energy.  Let
\(\Phi=rT\) be an actually admitted Hermitian ray,
\(\operatorname{Tr}_FT^2=1\), \(c_4=\operatorname{Tr}_FT^4\in(0,1]\).
Use the \emph{actual projected} coefficients \(M^2>0\),
\(\alpha_\Phi>0\) and \(\Lambda_g^2\), not merely their raw names.
For \(x=\alpha_\Phi r^2/(3M^2)\in[0,1)\), its Einstein potential is
\[
 W(x)=\frac{M^2C/2-6M^2\Lambda_g^2x+
                   9c_4M^4x^2/\alpha_\Phi}{(1-x)^2},
\]
and obeys the exact derivative test
\begin{equation}
 \boxed{(1-x)^3W'(x)
 =M^2(C-6\Lambda_g^2)
       +(18c_4M^4/\alpha_\Phi-6M^2\Lambda_g^2)x.}
 \label{p3f:eq:v2211-terminal-radial-test}
\end{equation}
If the constant coefficient is strictly negative and the linear
coefficient is nonpositive, there is no nonzero radial stationary point
in this admitted positive-graviton chart.  The origin then has negative
radial mass squared.

In particular, \emph{if the actual overlap proves}
\[
 M^2=U/4,\quad \alpha_\Phi=(16\pi^2)^{-1},\quad
 \Lambda_g^2=96\pi^2U,\quad 0<C<U,
\]
the test becomes
\[
 (1-x)^3W'(x)=M^2C-2304\pi^2M^4
                 +288\pi^2M^4(c_4-8)x<0.
\]
The positive native quartic therefore does not by itself certify a
physical broken vacuum or a stable terminal spectator.
\end{theorem}
\begin{proof}
At \(F=1\), the potential numerator in
\eqref{p3f:eq:v2202-hyperbolic-potential} is
\(M^2C/2+\alpha_\Phi(c_4r^4-2\Lambda_g^2r^2)\), while
\(\Omega=1-\alpha_\Phi r^2/(3M^2)\).  Substitute \(r^2=3M^2x/\alpha_\Phi\)
and differentiate the resulting quotient to obtain
\eqref{p3f:eq:v2211-terminal-radial-test}.  A strictly negative
\(W'(x)\) excludes a stationary point at \(r>0\), since \(dx/dr>0\)
there.  At the origin, the canonically normalized radial coordinate is
\(h=\sqrt{2\alpha_\Phi}\,r+O(r^3)\); the quadratic potential coefficient
gives \(m_{\rm rad}^2=C/3-2\Lambda_g^2<0\).
Hermiticity and \(\operatorname{Tr}_FT^2=1\) imply \(0<c_4\le1\)
by summing fourth powers of its real eigenvalues.
The displayed native substitutions make the linear coefficient negative,
and \(C<U=4M^2\) makes the constant coefficient negative.  This proves
the conditional native specialization.
\end{proof}
If the actual physical projection changes these coefficients, if the ray
is removed by a valid constraint, or if a retained source or quantum
effective potential contributes, the complete radial derivative must be
evaluated instead.  This obstruction concerns the specified bare matched
constituent; it neither excludes every completed interacting vacuum nor
supplies a compensating field, source or adjustable normalization.


\section{Fixed-current source identifiability and higher heat response}
\label{paper3f:sec:v280-current-heat}
\subsection{The right source overlap with fixed current rays}
\label{p3f:sec:v249-fixed-current-right-word}

The remaining charged-source proof cannot be replaced by re-identifying the
already known current rays. The following comparison holds those rays fixed
and, in addition, grants unit normalization of every mass-source arrow. It
isolates the surviving right-word datum by evaluating the complete tagged
response, rather than only giving a test for that response.

\begin{proposition}[Explicit finite source actions with the same current and single-role data]
\label{p3f:prop:v249-fixed-current-source-nonimplication}
Use the established left frame $(g_e,g_\mu,c)$ and right frame
$(h_e,h_\tau,c)$ of the rooted negative-Hopf construction, with all vectors
and their phases held fixed. Keep the already generated positive coefficients
$e,m,t$ unchanged. For $0\le\theta<\pi/2$ put
\[
 r_\theta=\cos\theta\,h_\tau+\sin\theta\,h_e,\qquad
 \iota_e=g_eh_e^\dagger,\quad
 \iota_\mu=g_\mu c^\dagger,\quad
 \iota_{\tau,\theta}=cr_\theta^\dagger,
 \qquad M_\theta=e\iota_e+m\iota_\mu+t\iota_{\tau,\theta}.
\]
These concrete finite source words have the same following data:
\begin{enumerate}
\item the native three hierarchy corners
$Q_\perp\operatorname{End}(G)Q_\perp$,
$Q_\perp\operatorname{End}(G)Q_0$,
$Q_0\operatorname{End}(G)Q_\perp$;
\item the three role coefficients, rank-one arrow singular spectra and
unit left returns
$\iota_r\iota_r^\dagger=P_r$, where
$(P_e,P_\mu,P_\tau)=(g_eg_e^\dagger,g_\mu g_\mu^\dagger,Q_0)$;
\item every scalar heat trace of each \emph{separate} weighted chiral arrow,
the fixed left-current map $W_e=(g_e,g_\mu,c)$ and its current Gram,
and the uninserted kinetic pairing $c_FI_L\oplus c_FI_R$;
\item the complete second source moment
$\operatorname{Tr}(M_\theta M_\theta^\dagger)=e^2+m^2+t^2$.
\end{enumerate}
They are realized by the explicit self-adjoint mass-only Hamiltonians
\[
 h_\theta=\begin{pmatrix}0&M_\theta\\M_\theta^\dagger&0\end{pmatrix},
 \qquad
 S_\theta[\Psi]=c_F\int\Psi^\dagger(i\hbar\partial_t-h_\theta)\Psi\,dt,
 \qquad c_F>0.
\]
The current source, clock and any fixed preparation may be held identical.
For a source $M$ put $B=MM^\dagger$ and define
$\mathcal D_e:=\|(B-e^2I)g_e\|^2$; use the subscript $\theta$
for the comparison words. Nevertheless their exact tagged electron resolvents are
\begin{equation}
 \boxed{G_{e,\theta}(z)
 =\frac{z(z^2-t^2)}
 {(z^2-e^2)(z^2-t^2)-e^2t^2\sin^2\theta},\qquad
 \mathcal D_{e,\theta}=e^2t^2\sin^2\theta.}
 \label{p3f:eq:v249-right-overlap-electron-return}
\end{equation}
In particular $G_{e,0}(z)=z/(z^2-e^2)$, while every
$0<\theta<\pi/2$ has a different response.
Thus even the listed data, which are stronger than presently proved for the
actual mass-source arrows, do not imply the electron mass-only closure.
\end{proposition}
\begin{proof}
Both $r_\theta$ and $h_e$ are unit vectors in $Q_\perp G$, so every
support, rank and single-arrow return in the list follows by multiplication.
For a weighted rank-one arrow of weight $a$, the six-dimensional chiral
square has spectrum $(a^2,a^2,0,0,0,0)$, proving equality of all separate
heat traces. The current map is unchanged by construction; its role Gram
is $I$, independently of $r_\theta$. Canonical normalization cancels $c_F$.
In the fixed left role frame the complete positive return is
\[
 B_\theta=M_\theta M_\theta^\dagger
 =\begin{pmatrix}
 e^2&0&et\sin\theta\\
 0&m^2&0\\
 et\sin\theta&0&t^2
 \end{pmatrix}.
\]
Its trace is fixed, but its action on $g_e$ contains the displayed nonzero
tau component, giving $\mathcal D_{e,\theta}$.
The left block of the chiral resolvent is $z(z^2I-B_\theta)^{-1}$.
Inverting its electron--tau block proves the rational formula. The
finite CAR anticommutator of this linear source equals its one-particle
resolvent, so no change of state is used to create the difference.
\end{proof}

\begin{proposition}[Complete tagged response from the mixed fourth source moment on this stronger branch]
\label{p3f:prop:v249-fourth-moment-right-overlap}
Suppose, independently, that the \emph{actual} three-corner source has the
unit, current-aligned arrows
\[
 \begin{gathered}
 \iota_e=g_er_e^\dagger,\qquad
 \iota_\mu=g_\mu r_\mu^\dagger,\qquad
 \iota_\tau=cr_\tau^\dagger,\\
 \|r_e\|=\|r_\mu\|=\|r_\tau\|=1,
 \quad r_e,r_\tau\in Q_\perp G,\quad r_\mu\in Q_0G.
 \end{gathered}
\]
Let $M=e\iota_e+m\iota_\mu+t\iota_\tau$, $B=MM^\dagger$, and retain
its actual complete mixed scalar contraction
\[
 \Delta_4:=\operatorname{Tr}B^2-(e^4+m^4+t^4).
\]
Then the full tagged electron mass-only response is exactly
\[
 \boxed{0\le\Delta_4\le2e^2t^2,\qquad
 \mathcal D_e=\frac{\Delta_4}{2},\qquad
 G_e(z)=\frac{z(z^2-t^2)}
 {(z^2-e^2)(z^2-t^2)-\Delta_4/2}.}
\]
For $0<\Delta_4<2e^2t^2$ it has four simple source-visible poles
$\pm\sqrt{\lambda_-},\pm\sqrt{\lambda_+}$, with
\[
 \lambda_\pm=\frac{e^2+t^2
       \pm\sqrt{(e^2-t^2)^2+2\Delta_4}}2,
\]
\[
 \boxed{\operatorname*{Res}_{z=\pm\sqrt{\lambda_-}}G_e
       =\frac{t^2-\lambda_-}{2(\lambda_+-\lambda_-)},\qquad
 \operatorname*{Res}_{z=\pm\sqrt{\lambda_+}}G_e
       =\frac{\lambda_+-t^2}{2(\lambda_+-\lambda_-)}.}
\]
All four residues are positive. In this independently certified branch,
$\Delta_4=0$ proves not only the electron return but the complete three-role
left-return identity
\[
 \boxed{B=e^2P_e+m^2P_\mu+t^2P_\tau.}
\]
It does not choose a preferred right frame.
\end{proposition}
\begin{proof}
Set $\gamma=r_e^\dagger r_\tau$. The muon mixed terms vanish by the
orthogonal right supports. The electron--tau block of $B$ is
$\left(\begin{smallmatrix}e^2&et\gamma\\et\bar\gamma&t^2\end{smallmatrix}\right)$.
Consequently
$\Delta_4=2e^2t^2|\gamma|^2$ and
$\mathcal D_e=e^2t^2|\gamma|^2$. Cauchy--Schwarz gives the range.
Direct inversion gives $G_e$; its characteristic polynomial gives
$\lambda_\pm$. In the strict range both eigenvalues are positive and
$\lambda_-<t^2<\lambda_+$; partial fractions give the stated positive
residues. Vanishing $\Delta_4$ forces the only remaining mixed return to
vanish, while all three diagonal returns are already the independently
certified unit ones. This proves the full left-return identity.
\end{proof}

\begin{remark}[What is, and is not, proved about the rooted source]
\label{p3f:rem:v249-right-word-origin-scope}
These propositions establish an exact finite information boundary and
response on their printed domains. The comparison actions are concrete
members of the predeclared finite source grammar, not a claim that every
$\theta$ is a GP2-generated physical output of the frozen rooted parent.
Their complete fourth source moment, determinant and full-source heat trace
are \emph{not} equal. Therefore they do not refute a possible upstream
constructor, nor an independently evaluated complete mixed action identity.

The new comparison keeps the derived current lines and all unit individual
source returns fixed. The v2.246 scalar-isospectral comparison instead
rotated the source's left lines while preserving its full unordered
spectrum. The two statements isolate different missing information.
Here a fourth \emph{complete-source} contraction would determine the tagged
response on the stronger unit-line branch. An action coefficient such as
$b_Y$ may include additional chiral, color, sector and comparison factors;
its identification with the printed $\operatorname{Tr}B^2$ must retain those
actual factors before this result is used.

For the native physical source, the licensed directed hierarchy word,
its left-support/unit-return certificates, or an equivalent actual mixed
return calculation still has to be supplied. The polar Hopf construction
and the derived current frame do not supply that action-incidence assertion.
No interacting particle pole, source preparation or continuum claim follows
from this mass-only calculation.
\end{remark}


\subsection{The complete correction budget and the first evaluated higher heat term}
\label{p3f:sec:v249-native-vacuum-budget}

\begin{theorem}[An exact same-owner trace budget for low native curvature]
\label{p3f:thm:v249-native-vacuum-correction-budget}
Use the actual native local scalar--metric action, with
$\kappa_J,\alpha>0$, $L=\Lambda_g^2>0$,
\[
 X=\operatorname{Tr}_F\Phi^2,\qquad
 \mathscr F=\kappa_J-\alpha X/6,\qquad
 \mathcal V_0=12\kappa_JL+
                   \alpha\operatorname{Tr}_F(\Phi^4-2L\Phi^2).
\]
The constant is the owned bare heat value of
Proposition~\ref{p3f:prop:v2216-bare-heat-vacuum}; it has not been
subtracted. Let an actually co-consumed remainder supply complete
action-normalized metric trace $T_\Delta$ and scalar insertion
$J_\Delta$, including any retained source, measure, nonlocal or higher
term exactly once. Thus its scalar variation is
$-\int\sqrt{-g}\operatorname{Tr}_F(J_\Delta\delta\Phi)$, and its
metric variation defines its Hilbert trace in the same pairing.
Assume these insertions are defined on the compared solution domain,
and set
\[
 \mathcal I_\Delta:=T_\Delta+\operatorname{Tr}_F(\Phi J_\Delta).
\]
On the full native metric and scalar mean equations, the exact
pointwise improved trace is
\[
 \boxed{2\kappa_JR
       =24L(\kappa_J+\mathscr F)-\mathcal I_\Delta,
 \qquad
 \frac RL=12\left(1+\frac{\mathscr F}{\kappa_J}\right)
                    -\frac{\mathcal I_\Delta}{2\kappa_JL}.}
\]
Consequently on any portion with $\mathscr F\ge f\kappa_J>0$
and $R\le\epsilon L$, one necessarily has
\[
 \boxed{\mathcal I_\Delta\ge
            2\kappa_JL[12(1+f)-\epsilon].}
\]
Equivalently, a certified bound
$|\mathcal I_\Delta|\le2\kappa_JL\eta$ implies
$R/L\ge12(1+f)-\eta$ there. A low-curvature target with
$\epsilon<12(1+f)-\eta$ is therefore impossible on that domain.
This identity also applies to a changing native scalar and metric:
the native covariant-gradient terms cancel in the improved trace.
It does not replace a complete remainder insertion by a potential
value or a single scalar heat trace. A common positive whole-action
normalization leaves every dimensionless inequality unchanged.
\end{theorem}
\begin{proof}
The full native improved trace, with the constant retained, is
\[
 2\kappa_JR=4(12\kappa_JL)-4\alpha LX
            -T_\Delta-\operatorname{Tr}_F(\Phi J_\Delta).
\]
It follows by the same metric trace and scalar-Euler cancellation as
Proposition~\ref{p3f:prop:native-source-curvature-trace}. In
particular the quartic and native kinetic contributions cancel, so
parallel or frozen fields have not been assumed. Substitute
$\alpha X=6(\kappa_J-\mathscr F)$ to obtain the exact boxed
identity. Inserting the two stated inequalities proves the budget.
For an actual common normalization $z>0$, the coefficients
$\kappa_J,\alpha,\mathscr F,T_\Delta,J_\Delta$ all acquire the
same factor $z$, while $L$ and the curvature coordinate do not;
their displayed ratios are unchanged.
\end{proof}

\begin{proposition}[What a controlled constant local remainder must supply]
\label{p3f:prop:v249-weighted-remainder-budget}
For a declared constant-scalar-curvature and parallel-scalar
comparison, suppose the actual local remainder is
$\Delta S=-\int\sqrt{-g}\,W(R,\Phi)$ with its true fixed
coefficients. Compact interior tests or the actual endpoint rule
are used. Then
\[
 \boxed{\mathcal I_\Delta
       =2RW_{,R}+\operatorname{Tr}_F(\Phi W_{,\Phi})-4W.}
\]
Off the constant-curvature/parallel-scalar comparison the metric
trace also contains $6\Box W_{,R}$; other derivative or tensor
terms retain their actual insertions. If
$W=\sum_dW_d$ is an actually supplied sum with
\[
 W_d(e^{2s}R,e^s\Phi)=e^{ds}W_d(R,\Phi)
\]
and the sum and its displayed first derivatives converge on the
comparison, then
\[
 \boxed{\mathcal I_\Delta=\sum_d(d-4)W_d.}
\]
Thus quartic weighted-degree-four terms cannot change this constant
improved-trace equation. A genuinely controlled bound
$\sum_d|d-4|\,|W_d|\le2\kappa_JL\eta$ obeys the previous theorem's
exclusion. Smallness of $W$ alone, without these metric and scalar
derivative bounds, is insufficient.

For an isolated weighted-degree-six correction the necessary
low-curvature size is explicitly
\[
 \boxed{|W_6|\ge\kappa_JL[12(1+f)-\epsilon].}
\]
At small $\epsilon$ this is at least the order of the entire bare
vacuum density $12\kappa_JL$, even as the positive-graviton margin
$f$ approaches zero. If instead only an owned constant correction
$\delta V$ is available and $\Phi=0$, the exact required value is
\[
 \boxed{\delta V=-12\kappa_JL+\kappa_JR/2.}
\]
This is a required independently derived cancellation, not a license
to choose or subtract that constant.
\end{proposition}
\begin{proof}
Variation of $-W$ gives
$T_\Delta=2RW_{,R}-4W+6\Box W_{,R}$ and
$J_\Delta=W_{,\Phi}$. The derivative term vanishes on the stated
constant comparison. Differentiating the displayed weighted scaling
at $s=0$ gives
$2RW_{d,R}+\operatorname{Tr}_F(\Phi W_{d,\Phi})=dW_d$.
Termwise summation proves the trace identity and the error bound.
For degree six, $|\mathcal I_\Delta|=2|W_6|$, so the previous
budget proves the stated size. At $\Phi=0$, a constant remainder has
$\mathcal I_\Delta=-4\delta V$ and $\mathscr F=\kappa_J$;
substitution in the exact trace proves its required cancellation.
\end{proof}

\begin{theorem}[The actual first raw higher heat term still has no low-curvature bare vacuum]
\label{p3f:thm:v249-evaluated-raw-a6-vacuum}
Use the actual unprojected ordinary torsion-free Dirac pure-metric
heat coefficients with their fixed exponential profile, common trace
support and finite-local Lorentzian clock domain, as in
Proposition~\ref{p3f:prop:v2212-dirac-a6-einstein}. Consider the
specified raw truncation through $a_6$, with $\Phi=0$, no source,
and a four-dimensional maximally symmetric interior metric. Retain
the same bare $a_0$ and Einstein $a_2$ coefficients and all
pure-metric $a_4$ terms. Then its exact constant-curvature Euler
equation is
\[
 \boxed{r+\frac{31}{181440}r^3=24,\qquad r=R/L.}
\]
It has exactly one real root, with
\[
 \boxed{22<r<23,\qquad r\simeq22.1446183528.}
\]
Thus this actually evaluated first higher raw heat term shifts the
bare $r=24$ stationary candidate but does not produce a controlled
$|R|/L\ll1$ vacuum. The root remains outside the small-curvature
domain; it is a stationary value of this stated truncation, not a
controlled solution or no-go for the full heat parent.
\end{theorem}
\begin{proof}
The cited same-owner evaluation gives the interior Euler
representative, before specializing the metric,
\[
 P_{6,E}\simeq_\delta\frac{31}{2160}R^3
                -\frac16RC^2+\frac1{15}J_C,
 \qquad S_{6,L}^{\rm raw}=-z\int\sqrt{-g}\,P_6,
 \qquad z=\frac{\kappa_J}{84L^2}>0.
\]
On the maximally symmetric jet $C=0$. Both the Weyl-square and
cubic-Weyl constituents have zero first variation there. The pure
$a_4$ terms also have zero interior Euler tensor there: they reduce
to constant-coefficient $R^2$, Weyl-square and Euler-density
constituents plus divergences; the $R^2$ tensor vanishes on a
constant-curvature Einstein metric, the Weyl-square tensor vanishes
at $C=0$, and Euler/divergence terms have no compact-interior Euler
variation in four dimensions. Their boundary terms are not being
deleted from a separate boundary claim. The full compared first
variation is therefore that of
\[
 \int\sqrt{-g}\left[
    \kappa_JR-12\kappa_JL
                 -\frac{31\kappa_J}{181440L^2}R^3\right].
\]
Its trace equation is
$2\kappa_JR=48\kappa_JL-
31\kappa_JR^3/(90720L^2)$, yielding the cubic. The derivative of
$r+31r^3/181440-24$ is strictly positive for every real $r$.
Its values at $22$ and $23$ are respectively negative and positive,
which proves uniqueness and the exact bracket. Ordinary bisection
gives the printed approximation.
\end{proof}

The new result fixes the sign and size of one actual missing
higher-order calculation. It does not assume that all later terms
are small at $r\simeq22$. A full resummation, a genuinely large
same-owner source/effective correction, a different licensed physical
branch or a different claimed observable can leave this truncation's
domain. Any proposed low-curvature completion nevertheless has to
meet the complete metric--scalar trace budget on the domain where
that decomposition is actually valid.


\subsection{The printed action's mixed fermion jet and the field/member distinction}
\label{p3f:sec:v250-field-versus-member}

\begin{proposition}[Exact mixed jet before and after regular classical boson elimination]
\label{p3f:prop:v250-boson-elimination-fermion-jet}
On an admitted complete finite coefficient chart, or a supplied common
differentiable action domain where the corresponding reduction is licensed,
write the printed common action in its uninserted fermion form
\[
 S[B,\bar\psi,\psi;Y]
 =S_B[B;Y]+c_F\langle\bar\psi,\mathscr D(B;Y)\psi\rangle,
 \qquad c_F>0.
\]
Here $B$ denotes exactly the retained bosonic fields, $Y$ the fixed
finite-owner member when the chart uses such a member, and $S_B$ includes
the actual spectral trace, its role weights and any retained purely
bosonic companions. The Dirac operator has no fermion dependence in this
printed word. Its exact mixed fermion derivative is consequently
\[
 \boxed{\frac{\delta^2S}{\delta\bar\psi_L\,\delta\psi_R}
       =c_F\mathscr D_{LR}(B;Y),}
\]
with the inherited graded derivative convention and the required
zero-order coefficient extraction when the source sort calls for it.
The bosonic heat term contributes zero to this derivative, independently
of whether its admissible profile is expanded locally.

Suppose $B_0(Y)$ is an actual stationary bosonic background at zero
fermion fields and a regular classical stationary elimination of $B$
is supplied on that same chart, with its actual gauge, constraint and
boundary treatment. Then the reduced fermion two-jet is still exactly
\[
 \boxed{\frac{\delta^2S_{\rm red}}
 {\delta\bar\psi_L\,\delta\psi_R}\bigg|_{\psi=0}
       =c_F\mathscr D_{LR}(B_0(Y);Y).}
\]
Thus the uninserted canonical kinetic cancellation leaves the finite
member already present in this evaluated Dirac coefficient. The background
can affect that coefficient through $B_0(Y)$, but regular classical boson
elimination supplies no additional quadratic source word.
\end{proposition}
\begin{proof}
$S_B$ contains neither fermion variable. Differentiating the printed
bilinear therefore gives the first identity directly. At zero fermion
fields each mixed boson--fermion Hessian block vanishes: differentiating
the bilinear once in a fermion variable leaves one fermion factor.
The exact stationary Schur formula consequently has zero cross-block
correction to the fermion two-jet. Equivalently the stationary bosonic
shift starts in fermion degree two. Its contribution to $S_B$ starts in
degree four because $DS_B(B_0)=0$, and its change of the fermion bilinear
also starts in degree four. The degree-two coefficient is precisely the
one displayed. This concerns regular classical stationary elimination;
quantum integration, a nonzero fermion background or a singular reduction
retains its actual measure, state and reduction problem.
\end{proof}

\begin{proposition}[What the native scalar variation varies]
\label{p3f:prop:v250-native-scalar-variation-scope}
On the no-constant-odd-offset coordinate-linear family
$\Phi_Y(H)=\Theta_Y(Y)(H)$ of
Theorem~\ref{p5:thm:v186-generated-yukawa-dirac-owner-family}, fixing the
finite member $Y$ and varying the retained scalar field $H$ gives
\[
 \boxed{\delta_HS_Y
 =D_\Phi S\,[\Theta_Y(Y)(\delta H)].}
\]
It gives the Euler equation on those admitted field directions, not an
Euler equation for $Y$. If a broader matrix-scalar variational domain is
actually licensed, its full scalar Euler equation instead acts on that
broader field; it still needs an independently proved identification of
its selected background with the charged preterminal mass-source target.
Promoting constant owner-family coordinates $Y$ to variational variables,
or promoting them to additional spacetime fields, is a different
variational problem unless that promotion has its own native constructor.
\end{proposition}
\begin{proof}
The first statement is the exact chain rule for the printed field map.
Its tangent at fixed $Y$ contains $\delta H$ and not $\delta Y$.
The owner-family bijection enumerates available source members; a
bijection does not enlarge the tangent space of a field variational
problem. Broader admitted scalar variations may be used only on their
own declared domain. The charged direct-source target and its incidence
are distinct entries of the action, so an equality to a chosen stationary
scalar background requires its actual source correspondence.
\end{proof}

\begin{proposition}[The zero-scalar fibre does not become a unique member selector]
\label{p3f:prop:v250-zero-scalar-member-fibre}
For this particular no-offset family, restrict to the source-free printed
spectral word and only companions that share its grading invariance and
have no independent $Y$ dependence on the $H=\psi=0$ fibre,
with no independent scalar source. Suppose the complete retained
Dirac domain, localizer/window and gauge data preserve the internal
grading $\Gamma_F$, the source-free base operator commutes with it, and
$\Gamma_F\Phi_Y(H)\Gamma_F=-\Phi_Y(H)$. Then the actual bosonic spectral
word is even under $H\mapsto-H$. Hence at $H=\psi=0$ the scalar first
variation vanishes for every $Y$, and the action and its metric/gauge
variations there are independent of $Y$.

If a metric/gauge critical configuration of that zero-scalar word is
independently supplied, it therefore extends to the entire available
$Y$ family at $H=\psi=0$. Even a formal enlargement allowing constant
$\delta Y$ would have $\delta_YS=0$ on this fibre. Thus existence of
such a zero-scalar stationary configuration cannot select a noncentral
charged member. This statement does not assert existence of an exact
full-heat metric critical configuration.
\end{proposition}
\begin{proof}
On the stated complete domain the constant unitary internal grading
conjugates $\mathscr D(H)$ to $\mathscr D(-H)$, and hence conjugates
its square and admissible profile in the same way. The retained trace
and localizer are invariant under that conjugation. Thus the bosonic
functional is even, and its derivative at zero scalar is zero.
Coordinate linearity gives $\Phi_Y(0)=0$ for every $Y$; the source-free
operator, action and metric/gauge derivatives on that fibre consequently
contain no $Y$. The fermionic term and its first variation vanish at
zero fermion fields. These facts prove the conditional extension and
also its formally enlarged constant-$Y$ version.
\end{proof}

These statements execute the printed action's variation; they do not add a
source-generation rule. The primitive member $Y_*=(I_G,I_G,I_G,I_G)$ is
already distinguished by its separate rooted triality source word, not by
extremizing over the family. The charged hierarchy still requires its actual
preterminal directed source constructor or an equivalent actual mixed-return
identity. The result does not prohibit a future native constructor or a
proved claim-specific compatibility condition from determining that source.


\subsection{Exact full-heat scale and scalar variations}
\label{p3f:sec:v250-exact-heat-scale-tests}

The preceding $a_6$ calculation is a test of a stated truncation.
The next result differentiates the exact unlocalized compact
Euclidean heat functional already defined in
Proposition~\ref{p3f:prop:v1164-heat-hessian-no-insertion}. It does
not continue the full trace to Lorentzian signature or create a new
physical round-space branch.

\begin{theorem}[The exact ordinary heat trace has no unconstrained homothetic stationary metric]
\label{p3f:thm:v250-exact-heat-homothety}
Let $D_0$ be the actual ordinary self-adjoint geometric Dirac operator
on a closed compact Riemannian spin parent in the unlocalized
Euclidean heat domain, and put $L=\Lambda_g^2>0$. Retain its fixed
trace prescription, fixed positive cutoff and the common transported
spinor/domain identification. Suppose the configuration domain admits
the homothetic metric variation $g_a=a^2g_0$, $a>0$, without a
fixed-volume constraint or other independent scale condition. With
the standard unitary density transport one has $D_a=a^{-1}D_0$ and
\[
 \boxed{S(a)=\operatorname{Tr}e^{-D_0^2/(a^2L)},\qquad
 \frac{dS}{d\log a}=2\operatorname{Tr}(A_ae^{-A_a})>0,
 \quad A_a=D_0^2/(a^2L).}
\]
Strictness holds whenever the retained carrier contains a nonzero
Dirac eigenvalue; in particular it holds for the complete compact
elliptic carrier. Therefore this exact ordinary, unnormalized pure
heat functional has no finite-scale critical metric on that admitted
unconstrained homothetic branch. A positive constant action or
protected trace multiplier preserves this conclusion.

The same identity holds on an actual fixed finite reducing spectral
window, with strictness precisely when that window sees a nonzero
eigenvalue. It is not asserted across a moving sharp band or a change
of boundary domain.
\end{theorem}
\begin{proof}
Constant metric scaling multiplies the inverse coframe and the
geometric Dirac operator by $a^{-1}$ under the common unitary
density identification. Hence $D_a^2=a^{-2}D_0^2$. Write the actual
squared eigenvalues of $D_0$ as $\lambda_j^2\ge0$, including
multiplicities. On a compact $a$ interval bounded away from zero,
both the heat sum and its differentiated sum converge uniformly:
$x e^{-cx}\le C e^{-cx/2}$ for every fixed $c>0$, and the latter
heat operator is trace class. Thus termwise differentiation is
licensed and gives
\[
 \frac{dS}{d\log a}
   =2\sum_j\frac{\lambda_j^2}{a^2L}
                    e^{-\lambda_j^2/(a^2L)}.
\]
Every term is nonnegative and a retained positive eigenvalue makes
the sum strictly positive. The geometric metric derivative cannot
vanish on all allowed variations because this one is nonzero.
In a fixed finite reducing window the same calculation is a finite
sum. It includes the actual transported volume through the spectrum;
there is no additional volume division.
\end{proof}

\begin{proposition}[An explicit source-native compact product test with momentum fixed]
\label{p3f:prop:v250-native-product-heat-scale}
Use the actual ordinary-spin product spectrum and heat trace of
Paper~I-C, Eqs.~\eqref{p1b:der:g08:eq-G08-6}--
\eqref{p1b:der:g08:eq-G08-9}. Hold the momentum radius $R_p$, torus
metric $G$, spin structure $\sigma$ and heat time $t>0$ fixed, and
vary only the round position-factor radius $a=R_x$. The exact
declared spectral readout is
\[
 Z_{\rm spin}(t;a)
   =4\mathcal S_a(t)\mathcal S_{R_p}(t)\Theta_{G,\sigma}(t),
 \quad
 \mathcal S_a(t)=2\sum_{k\ge0}(k+1)(k+2)
                  e^{-t(k+3/2)^2/a^2}.
\]
Its actual derivative is
\[
 \boxed{a\partial_a Z_{\rm spin}
   =\frac{16t}{a^2}\mathcal S_{R_p}(t)\Theta_{G,\sigma}(t)
       \sum_{k\ge0}(k+1)(k+2)(k+3/2)^2
                   e^{-t(k+3/2)^2/a^2}>0.}
\]
Fixed nonempty complete finite product shell windows obey the same
strict formula with their actual finite sums. An arbitrary fixed
reducing multiindex window instead retains its complete finite
positive derivative sum, without assuming product factorization.
No homothety of the
microscopic momentum radius has been used. This is an exact test
of the declared eight-dimensional geometric readout, not an
identification of that readout with a four-dimensional physical
metric action.
\end{proposition}
\begin{proof}
The native product square is the sum of the commuting squared sphere
Dirac operators and fixed shifted torus Laplacian. Its actual heat
trace therefore factorizes as printed. Differentiate the position
factor only; its complete two-sign multiplicity is already included
in $\mathcal S_a$, so no additional sign or Clifford multiplicity is
inserted. Uniform convergence follows from the Gaussian spectral
tail. Every factor and every term in the derivative is positive.
\end{proof}

\begin{theorem}[The exact native raw joint heat has no nonzero parallel scalar stationary point]
\label{p3f:thm:v250-exact-parallel-scalar-heat}
On an admitted even-dimensional compact unlocalized Euclidean joint
Dirac domain, use the actual generated operator word of
Theorem~\ref{p3f:thm:v175-native-local-scalar-action},
\[
 D_{a,s}=a^{-1}D_0\otimes I+\gamma_g\otimes s\Phi_0,
 \qquad \Phi_0=\Phi_0^\dagger,
\]
with a finite internal factor, zero added gauge curvature and a
parallel $\Phi_0$. Require its actual tensor-product domain and
Clifford grading, so $\{D_0,\gamma_g\}=0$,
$\gamma_g^2=I$, and $[D_0,\Phi_0]=0$ on the common core. Keep the
ordinary raw trace, fixed cutoff $L$ and any actual constant common
positive multiplier. The geometric carrier contains a nonzero
Dirac mode, as automatically holds on the complete
positive-dimensional compact elliptic parent. Then
\[
 \boxed{\begin{aligned}
 D_{a,s}^2&=a^{-2}D_0^2\otimes I+I\otimes s^2\Phi_0^2,\\
 S(a,s)&=Z_g(a)\operatorname{Tr}_F e^{-s^2\Phi_0^2/L},\qquad
 Z_g(a)=\operatorname{Tr}_g e^{-D_0^2/(a^2L)}.
 \end{aligned}}
\]
The two exact derivatives are
\[
 \boxed{\partial_{\log a}S
   =2\operatorname{Tr}_g(A_ae^{-A_a})
                    \operatorname{Tr}_F e^{-s^2\Phi_0^2/L}>0,}
\]
\[
 \boxed{\partial_sS
  =-\frac{2s}{L}Z_g(a)
             \operatorname{Tr}_F(\Phi_0^2e^{-s^2\Phi_0^2/L})<0
  \quad(s>0,\ \Phi_0\ne0).}
\]
Thus any admitted nonzero parallel scalar configuration has a
nonvanishing radial variation of this exact raw heat functional.
If the full parallel scalar linear support is varied, its only
finite scalar stationary point is $\Phi=0$. The metric variation
there still does not vanish on the admitted homothetic branch.

This is an exact Euclidean raw-parent statement. It does not replace
the physical common-action pure-metric/nonmetric role scalarizer
by one uniform raw weight, assert an exact global Lorentzian heat
action, or apply to a rooted midpoint scalar whose covariant spatial
derivative is nonzero.
\end{theorem}
\begin{proof}
The actual Clifford anticommutation and parallel internal factor
remove the mixed square term. The two surviving positive operators
commute on different tensor factors, so their heat operator and
ordinary trace factorize exactly. The first derivative is the
preceding theorem. Finite internal functional calculus gives the
second derivative. A nonzero Hermitian matrix has a positive squared
eigenvalue, and every exponential weight is positive, so its stated
trace is strictly positive. The radial direction at any nonzero
matrix therefore prevents stationarity. At zero all scalar first
variations vanish, proving the remaining assertion.
\end{proof}

\begin{remark}[What has and has not been resummed]
\label{p3f:rem:v250-full-heat-scope}
These are full spectral sums, not an extrapolation of the $a_6$
root. In particular a stationary point of a finite curvature or
scalar polynomial truncation does not establish a stationary point
of the exact raw Euclidean heat function on these same admitted
variations. In the scalar factor even the ordinary quartic
approximation $1-x+x^2/2$ has a spurious stationary value at $x=1$,
whereas $e^{-x}$ has none. The native local quartic action remains
its correct finite-jet coefficient on its own domain.

The cutoff, trace normalization and variation space matter. If
$L=L(a)$ instead of the fixed native comparison cutoff, the scale
derivative acquires the factor
$1+\tfrac12\partial_{\log a}\log L$; a cutoff with $L\propto a^{-2}$
makes this particular spectral ratio constant. Dividing the trace
by volume introduces a different derivative. Moving projectors,
localizers, boundaries and scalar/metric constraints retain their
actual contacts. None of these changes is licensed merely to cancel
the nonzero derivative proved above. Most decisively, the
finite-local clock real-form theorem explicitly does not supply a
global continuation of the exact nonlocal heat functional. No
Lorentzian vacuum or universal physical no-vacuum claim is inferred
from this Euclidean result.
\end{remark}


\section{Protected native photon cone and full-source Hessian test}
\label{p3f:sec:v358-photon-cone-certificate}
\begin{proposition}[Native cone and complete-action frontier]
\label{p3f:prop:v358-native-photon-cone-scope}
On the stationary, parity-even source-free native Yang--Mills branch of
Theorem~\ref{p3f:thm:joint-gauge-principal-symbol-inheritance},
Corollary~\ref{p3f:cor:v244-transverse-gauge-propagation} already
establishes the transverse wave operator
$K_{\gamma,T}=\partial_t^2+c_{\rm op}^2\Delta_1^{\rm co}$
and its common metric characteristic cone.  The represented massless
photon eigendirection inherits the cone on the domain of the
corresponding Paper~V electroweak export.  The neutral mixed-Hessian
vanishing of Proposition~\ref{p3f:prop:v1166-neutral-photon-mixing}
identifies its Schur block with the \emph{direct} photon Hessian on the
stated neutral branch; it does not identify the complete direct block
with the retained local Yang--Mills constituent.  Full physical optical
transport requires a separate same-action source/cone certificate.
\end{proposition}
\begin{proof}
The cited native corollary specifies the stationary coexact transverse
operator.  A vanishing mixed block removes the mixed Schur correction
but imposes no additional identity on the direct diagonal block.
Thus neither the protected native cone nor the stated full-action
frontier is strengthened beyond its original domain.
\end{proof}

\begin{proposition}[Exact finite heat-trace two-jet]
\label{p3f:prop:v358-photon-heat-hessian}
On a fixed finite covariant source window and fixed positive cutoff
$\Lambda$, let $D=D^*$ be twice differentiable in real source
coordinates $u,v$ and $S(D)=\frac12\operatorname{Tr}
\exp(-D^2/\Lambda^2)$.  At the base source put
\[
\begin{gathered}
 X=D^2/\Lambda^2,\quad X_u=\{D,D_u\}/\Lambda^2,\\
 X_v=\{D,D_v\}/\Lambda^2,\quad X_{uv}=(D_uD_v+D_vD_u+\{D,D_{uv}\})/\Lambda^2.
\end{gathered}
\]
In an eigenbasis $X|r\rangle=\lambda_r|r\rangle$, define
$\phi(a,b)=\int_0^1e^{-(1-s)a-sb}\,ds$.  Then
\begin{equation}
\label{p3f:eq:v358-photon-heat-two-jet}
 d_vd_uS=\frac12\sum_{r,s}\phi(\lambda_r,\lambda_s)
 (X_u)_{rs}(X_v)_{sr}
 -\frac12\sum_r e^{-\lambda_r}(X_{uv})_{rr}.
\end{equation}
The formula is continuous through spectral degeneracy.  Its contact
term, including $D_{uv}$, cannot in general be omitted.
\end{proposition}
\begin{proof}
Differentiate $e^{-X}$ using the Duhamel integral and cyclicity
of the finite trace; then diagonalize $X$.  The integral definition
of $\phi$ is continuous at $a=b$, with $\phi(a,a)=e^{-a}$.
This is a specialization of the exact finite heat two-jet already
proved in Paper~III-B, Proposition~\ref{p3b:prop:v245-exact-finite-heat-two-jet}.
\end{proof}

\begin{corollary}[Pure-gauge Ward cancellation and regulator test]
\label{p3f:cor:v358-photon-gauge-orbit}
If the retained window is invariant under a unitary source orbit
$D(u)=U(u)D_0U(u)^*$ with $U(u)=e^{uT}$, $T^*=-T$, then
$D_u=[T,D_0]$, $D_{uu}=[T,[T,D_0]]$, and the complete quadratic
heat response along the gauge orbit is zero.  A non-invariant
projection can spoil this cancellation: if
$D_0=\operatorname{diag}(1,2)$ is rotated through angle $u$ and
compressed onto its first coordinate, one may take
$D_P(u)=1+\sin^2u$, for which
\[
 \left.\frac{d^2}{du^2}\left(\frac12e^{-D_P(u)^2}\right)
 \right|_{u=0}=-2e^{-1}\ne0.
\]
This is an example of a \emph{bad} window, not an assertion that the
framework's own finite cutoff is noncovariant.
\end{corollary}
\begin{proof}
On the invariant window the heat trace is constant under conjugation;
differentiate twice, including the nonzero orbit acceleration.
The compressed counterexample follows from two elementary derivatives.
\end{proof}

\begin{remark}[Execution contract for a source-selected electromagnetic Hessian]
\label{p3f:rem:v340-s3lit-em-source-execution}
The finite two-jet in
Proposition~\ref{p3f:prop:v358-photon-heat-hessian} and the Ward
criterion in Corollary~\ref{p3f:cor:v358-photon-gauge-orbit}
are the required \emph{same-action} tests, not a license to
replace the physical electromagnetic current by a graph-coordinate
commutator or an auxiliary Wilson/overlap twist.  Before claiming a
numerical physical electromagnetic response, retain on one
source-selected window: (i) the occurring $U(1)_Q$-charged physical
Hamiltonian and source field; (ii) all first and second source jets,
including the contact/acceleration term; (iii) the associated
transverse physical current and local kinetic normalization;
(iv) a gauge-invariant regulator and zero response along every
retained pure-gauge orbit; and (v) the actual charged-state spatial
form factor and declared local/global/infrared chart.  The temporal
current-response bounds of
Corollary~\ref{p12:cor:v340-s3lit-current-envelope}
apply only when its positivity/gap hypotheses are met.  The
physical source packet and common Thomson/ultraviolet limits are
not evaluated here and remain open action-level obligations.
\end{remark}


\begin{proposition}[Observable-oriented local optical-cone criterion]
\label{p3f:prop:v358-local-photon-cone-test}
Assume an actually selected, local, homogeneous and isotropic,
parity-even, gauge-invariant two-derivative transverse source action
has positive coefficients $Z_E,Z_B$ and quadratic part
\[
 S_T^{(2)}=\frac12\int dt\,d\mu_h
 \bigl(Z_E|\dot a_T|^2-c_{\rm op}^2Z_B|da_T|^2\bigr).
\]
Its transverse characteristic speed equals the native operational
speed exactly when $Z_E=Z_B$, since
$c_\gamma^2=c_{\rm op}^2Z_B/Z_E$.  The native Yang--Mills branch
already satisfies this criterion.  For the complete nonlocal direct
action this test is only applicable after deriving the appropriate
principal symbol and physical Lorentzian transport; minimal Maxwell
form is \emph{not} a necessary premise for a shared null cone.
\end{proposition}
\begin{proof}
The transverse Euler equation has principal part
$Z_E\partial_t^2+c_{\rm op}^2Z_B\Delta_1^{\rm co}$.
Its characteristic equation yields the displayed ratio; positivity
selects the stated positive speed.
\end{proof}


\section[The complete action and the polyhedral connector-return frontier]{The complete action and the polyhedral\protect\linebreak connector-return frontier}
\label{p3f:sec:v1264-octahedral-icosahedral-action-scope}
\begin{remark}[No force follows from a finite stabilizer alone]
\label{p3f:rem:v1264-polyhedral-mixed-source-scope}
The full local octahedral response stabilizer and the five-core
icosahedral incidence of Paper I-B are exact finite geometry.  The
balanced $E_{36}$ source of Paper III-C is an actual finite operator
word, but the color-adjoint three-face field, the charged $u,d$
source slots and the three selected $2I$ core-pair edges have
independent physical owner types.  An $A_5$-invariant tensor
or compatible $C_3$ projectors do not themselves produce a term
in the complete action.  For a dynamically generated connector
source $J$ and positive Hessian $K$ on a common owner, a Schur
reduction produces the conditional return $-J^\dagger K^{-1}J/2$;
its sign follows from positivity, but neither its coefficient
nor the occurrence of $J$ is fixed by symmetry.  Ground-state
stabilization must be distinguished from a higher excitation
and from an open-channel resonance.  The existing fermion action,
its source jets and all previous Ward/kinetic hypotheses are unchanged.
\end{remark}


\subsection{Marked source Schur jet and horizontal Clifford-to-cap comparison}
\label{p3f:sec:et-source-port}

\begin{proposition}[Reference Schur two-jet of the declared cap source]
\label{p3f:prop:et-cap-source-schur}
In the \emph{conditional} cap parent of
Definition~\ref{p3e:def:v223-cap-data}, retain its sole source-density
factor $1/8$.  Write
\[
 \mathcal L_s(R,Y)=\frac18\Bigl[\operatorname{Tr}
 (Y^{\mathsf T}T_s(R)Y)+2\rho\operatorname{Tr}Y\Bigr],
 \quad \rho=\frac{R_s}{2},\quad
 X=\begin{pmatrix}0&1\\1&0\end{pmatrix},
\]
$u(R)=c\,(R/R_s)/(1+\sqrt{1+(R/R_s)^2})$,
$T_s=(I_2+uX)\otimes(I_2-uX)$, and
$A=X\otimes I_2-I_2\otimes X$.  At the \emph{stationary marked}
source $R=0$, $Y_*=-\rho I_4$, the full $R\oplus\mathbb R^{16}$
source-density Hessian has blocks
\[
 \boxed{H_{YY}=\tfrac14 I_{16},\qquad
 H_{RY}=-\tfrac{c}{16}\operatorname{vec}(A)^\mathsf T,\qquad
 H_{RR}^{(s)}=0.}
\]
It follows that
\[
 \boxed{H_{RY}H_{YY}^{-1}H_{YR}=\frac{c^2}{8},\qquad
 \left.\frac{d^2\mathcal L_{s,*}}{dR^2}\right|_0=-\frac{c^2}{8}.}
\]
The mixed source row lies in the nine-dimensional symmetric trace-free
subspace of $\operatorname{Mat}_4(\mathbb R)$, orthogonal at the
reference to $\mathbb RI_4\oplus\Lambda^2\mathbb R^4$.
\end{proposition}
\begin{proof}
At $R=0$, $T_s=I_4$, $u'=c/(2R_s)$ and
$T_s'=cA/(2R_s)$.  The source quadratic gives
$H_{YY}=I_{16}/4$ and, because $Y_*=-\rho I_4$,
$H_{RY}=-\rho\,\operatorname{vec}(T_s')^\mathsf T/4
=-c\operatorname{vec}(A)^\mathsf T/16$.
Here $T_s''=-2(u')^2(X\otimes X)$ has zero trace, so
$H_{RR}^{(s)}=0$.  Finally $A=A^\mathsf T$,
$\operatorname{Tr}A=0$, and $\operatorname{Tr}A^2=8$ give the
stated Schur value.  This is a second jet of the \emph{declared cap
chart}, not a prior derivation of its input parameter $c$.
\end{proof}

\begin{proposition}[Marked horizontal isometry into the cap source ports]
\label{p3f:prop:et-marked-port-isometry}
For a fixed oriented horizontal plane $H_u\cong\mathbb R^2$, set
$J(v)=v_1\sigma_x+v_2\sigma_z$, with the real symmetric Pauli pair,
and define
\[
 \iota(v):=\frac1{\sqrt2}
 \bigl(J(v)\otimes I_2-I_2\otimes J(v)\bigr)
 \in S_0^2(\mathbb R^4).
\]
The sixteen-port per-coordinate bilinear form
$\mathfrak t_{16}(M)=\operatorname{Tr}_{\mathbb R^4}(M)/16$
(note that $\mathfrak t_{16}(I_4)=1/4$, so this is \emph{not}
a normalized tracial state on $M_4$) satisfies
\[
 \boxed{\mathfrak t_{16}(\iota(v)\iota(w))
 =\frac14\langle v,w\rangle_{H_u}.}
\]
Thus it matches the $\tau_8$ Clifford source-edge Gram in
Proposition~\ref{p3c:prop:et-shell-edge-gram}.  The map is equivariant
for the \emph{marked} horizontal Spin(2) action by paired real
half-angle conjugation, and $A=\sqrt2\,\iota(1,0)$.
It transports the complete exceptional unnormalized quadratic shell
returns $(4/3,1,2)$, while their coherent linear shell sums vanish.
This is a comparison isometry, not a physical map from the native
Dirac/heat source to the gravitational cap ports.
\end{proposition}
\begin{proof}
$\operatorname{Tr}J(v)=0$ and
$\operatorname{Tr}J(v)J(w)=2\langle v,w\rangle$.
Expanding the tensor factors gives
$\operatorname{Tr}\iota(v)\iota(w)=4\langle v,w\rangle$,
proving the Gram equality.  Simultaneous conjugation with the real
half-angle rotation preserves both tensor factors and intertwines the
marked horizontal rotation (with its orientation convention).
The shell result follows from the raw tight frames, and antipodality
forces the first moment to vanish.
\end{proof}

\begin{remark}[Native source selection is not supplied by an isometry]
\label{p3f:rem:et-native-selection-firewall}
The source-free, unmarked restricted boundary-heat packet at $\Phi=0$
can have a vanishing first curvature--odd-source mixed jet, while the
marked stationary cap source above has a nonzero mixed jet.
Nor does the unbroken terminal $C_3$ branch admit a generic unmarked
linear $E_{48}$ source.  A physical equality requires an actual
same-action marked incidence, the boundary Hessian and complete
metric/localizer derivatives, the reciprocal determinant-line
transport, and the correct energy-level rebase.  An available isometry
and the conditional cap Schur value do not discharge these inputs.
\end{remark}


\begin{remark}[Same-action test for charged-family origin and visibility]
\label{p3f:rem:v1271-charged-hessian-origin-gate}
The mixed chiral coefficient in a charged sector $s\in\{u,d,e\}$ is a
physical source assertion only if the underlying fermion action and
kinetic factors are independently fixed on one occurring action member:
\[
 M_{s}^{\rm act}=Z_{L,s}^{-1/2}
 \frac{\delta^2S_\star}{\delta\bar\psi_{L,s}\,\delta\psi_{R,s}}
 Z_{R,s}^{-1/2}.
\]
An equality of this coefficient to a geometric hierarchy/Hopf word is a
\emph{test of the independently constructed action}; assigning that word
to the input Yukawa member and then differentiating is not a proof of its
origin.  Proposition~\ref{p3f:prop:v250-native-scalar-variation-scope}
varies $H$ at fixed $Y$ and supplies no missing charged-arrow selector.
Kinetic, connection, boundary and mixed terms are retained according to
their actual action domain; no bosonic stationary reduction is used to
invent a quadratic family term.  The $u,d,e$ insertions need separate,
owner-resolved signed certificates.  A common free left-current response
on $R\mapsto RU$ does not select an unmarked right frame: a stronger
right-sensitive physical claim requires an independently occurring
right-channel source/mark and its same-action response.  This remark
records an acceptance gate, not an unconditional charged-source theorem.
\end{remark}


\section{Schur reduction, boundary momentum, and matched current}
\label{paper3f:sec:organized-boundary}
\subsection{Exact finite pole and source transport through a boundary}
\label{p3f:sec:v1140-finite-pole-source-transport}
The following same-action boundary theorem retains the full energy-dependent
kinetic normalization and the source in both blocks.  It is consumed by the
finite-to-continuum transport layer of Paper~VII-C; the analytic pole domain
is stated here rather than inferred from a terminal mass readout.

\begin{theorem}[Finite isolated pole and residue under exact Schur elimination]
\label{p3f:thm:v1140-finite-pole-source-transport}
Fix one finite physical action owner and its declared gauge/constraint
quotient.  Let its inverse propagator be real analytic and Hermitian on a
real energy interval containing $z_*$, with blocks
\[
 D(z)=\begin{pmatrix}A(z)&B(z)^\dagger\\B(z)&C(z)\end{pmatrix},
 \qquad \det C(z)\ne0.
\]
The block formula is written on the real interval; off that interval use its
holomorphic continuation, not a pointwise complex-energy adjoint.  Define
\begin{equation}
\label{p3f:eq:v1140-schur-pole-lift}
 S(z)=A(z)-B(z)^\dagger C(z)^{-1}B(z),\qquad
 L(z)=\begin{pmatrix}I\\-C(z)^{-1}B(z)\end{pmatrix}.
\end{equation}
Then $L(z)$ identifies the kernels of $S(z)$ and $D(z)$ and
\begin{equation}
\label{p3f:eq:v1140-schur-pole-determinant}
 D(z)L(z)=\begin{pmatrix}S(z)\\0\end{pmatrix},\qquad
 \det D(z)=\det C(z)\det S(z).
\end{equation}
Suppose $\ker S(z_*)=\mathbb C u$ with $u^\dagger u=1$ and
\[
 a:=u^\dagger S'(z_*)u\ne0.
\]
The local inverse has a simple pole and its residue is
\begin{equation}
\label{p3f:eq:v1140-pole-residue-lift}
 \boxed{\operatorname*{Res}_{z=z_*}D(z)^{-1}
   =\frac{L_*uu^\dagger L_*^\dagger}{a},\qquad
 a=(L_*u)^\dagger D'(z_*)(L_*u),}
\end{equation}
where $L_*=L(z_*)$.  For an analytic retained source
$j(z)=(j_R(z),j_C(z))$, with all quantities in the following expression
evaluated at $z_*$, the source residue is
\begin{equation}
\label{p3f:eq:v1140-retained-source-residue}
 \boxed{\operatorname*{Res}_{z=z_*}\bigl(j^\dagger D^{-1}j\bigr)
 =\frac{|u^\dagger j_{\rm eff}|^2}{a},\qquad
 j_{\rm eff}=j_R-B^\dagger C^{-1}j_C.}
\end{equation}
The adjoint source is continued from the same real analytic source chart.
A source orthogonal to this effective null channel has zero residue in that
channel even though the inverse propagator has the displayed pole.
\end{theorem}

\begin{proof}
Block multiplication gives $DL=(S,0)^{\mathsf T}$, while the second block of
$D(v,w)^{\mathsf T}=0$ forces $w=-C^{-1}Bv$.  Block triangular elimination
gives the determinant identity and, wherever $S$ is invertible,
\[
 D^{-1}=LS^{-1}L^\dagger+
 \begin{pmatrix}0&0\\0&C^{-1}\end{pmatrix}.
\]
These real-interval identities extend to the corresponding holomorphic
germs.  The simple Hermitian zero of $S$ has eigenvalue derivative
$u^\dagger S'u=a$; its inverse therefore has residue $uu^\dagger/a$.
The complement inverse is analytic, so the displayed block inverse gives
\eqref{p3f:eq:v1140-pole-residue-lift}.  Finally $L^\dagger DL=S$.
After differentiation and contraction with $u$ at $z_*$, the terms containing
$L'$ vanish because $D(z_*)L_*u=0$, giving the identity for $a$.
Contract the block inverse with the full source and use
$L_*^\dagger j=j_R-B^\dagger C^{-1}j_C$ to obtain
\eqref{p3f:eq:v1140-retained-source-residue}.
\end{proof}


\begin{corollary}[Same-sector lowering and source weight on a positive Hamiltonian pivot]
\label{p3f:cor:schur-positive-gap-weight}
Specialize Theorem~\ref{p3f:thm:v1140-finite-pole-source-transport}
to the canonical Hamiltonian inverse $zI-H$, with a fixed
energy-independent Hermitian owner
$H=\left(\begin{smallmatrix}A&B\\B^\dagger&C\end{smallmatrix}\right)$.
If $\delta I\preceq C-EI\preceq DI$ with $\delta>0$, and a normalized
retained null direction $u$ has
$s=u^\dagger B(C-EI)^{-1}B^\dagger u$, then
\[
 \boxed{\frac{\delta}{\delta+s}\le
 \frac{1}{1+u^\dagger B(C-EI)^{-2}B^\dagger u}
 \le\frac{D}{D+s}.}
\]
The middle expression is the retained population of the normalized full
lift; at a simple pole it multiplies $uu^\dagger$ in the retained
residue.  For degeneracy use the complete nullspace matrix residue and
Loewner bounds of Theorem~\ref{p1b:thm:response-positive-gap-weight}.
\end{corollary}
\begin{proof}
The Hamiltonian derivative norm in
\eqref{p3f:eq:v1140-pole-residue-lift} is
$1+u^\dagger B(C-EI)^{-2}B^\dagger u$.
Apply the positive-gap spectral inequalities in
Theorem~\ref{p1b:thm:response-positive-gap-weight} to that same return.
\end{proof}
A direct-block Rayleigh lowering and its retained residue cannot be
specified independently for this fixed sector.  Full-source overlaps
still use \eqref{p3f:eq:v1140-retained-source-residue}; the inequality does
not set those overlaps.  Generic energy-dependent kinetic terms or an
indefinite complementary Dirac block retain the general derivative norm
and do not inherit this positive-gap estimate.

\begin{corollary}[Grouping-independent finite pole and source transport]
\label{p3f:cor:v1140-pole-transport-composition}
For a fixed final retained block, any two exact nested Schur elimination
orders with invertible eliminated blocks on the same analytic neighborhood
produce the same isolated pole positions and full source residues.
\end{corollary}
\begin{proof}
Both orders complete the same finite quadratic form with the same full source.
Successive substitution is associative; each result is the corresponding
block-inverse/source identity of Theorem~
\ref{p3f:thm:v1140-finite-pole-source-transport}.  Hence their meromorphic
source functions, and therefore their residues, agree.
\end{proof}

The derivative $a$ includes every inherited energy-dependent self-energy and
kinetic term; replacing it by $1$ requires a separately justified coordinate
normalization.  This is a finite transport theorem on an isolated analytic
branch.  It neither establishes a sharp charged pole in a continuum with
massless quanta nor identifies a terminal geometric finite mass with a pole
mass.  At a singular complement the eliminated modes must be retained in a
larger chart.  Exact block elimination of the owned inverse propagator is the
constructor used here; a truncated effective action is not substituted for it.

\subsection{Second source jet of the same-action boundary reduction}
\label{p3f:sec:v199-second-schur-jet}
Work on a fixed, normalized boundary/interior trivialization after the
declared gauge quotient.  Assume the occurring stationary family exists
smoothly and its finite Hermitian Hessian has twice differentiable blocks
\[
 H(y)=\begin{pmatrix}A&B^\dagger\\ B&C\end{pmatrix},\qquad C(y)>0,
 \qquad K=A-B^\dagger C^{-1}B.
\]
The positivity of the interior block is the elimination hypothesis; the
full Hessian may have a critical boundary mode.  No inverse of the full
Hessian is needed for the following extension of
Theorem~\ref{p3f:thm:v176-explicit-source-jet-schur}.

\begin{theorem}[Exact first and second boundary Schur jets]
\label{p3f:thm:v199-second-schur-jet}
Put \(Z=-C^{-1}B\), \(L=(I,Z)^{\mathsf T}\), and
\(R_{\rm load}=B'+C'Z\), where primes denote the actual source derivative.
Then
\begin{equation}
 K'=L^\dagger H'L,\qquad
 K''=L^\dagger H''L-2R_{\rm load}^\dagger C^{-1}R_{\rm load}.
 \label{p3f:eq:v199-second-schur-jet}
\end{equation}
For a fixed normalized critical boundary vector \(e\), the coefficient of
\(y^2\) in \(\langle e,K(y)e\rangle\) at zero is
\begin{equation}
 c_2=\tfrac12\langle Le,H''Le\rangle
       -\|C^{-1/2}R_{\rm load}e\|^2,
 \label{p3f:eq:v199-critical-second-coefficient}
\end{equation}
with all right-hand quantities evaluated at zero.
\end{theorem}
\begin{proof}
Apply the fixed-lift Schur identity of Paper~I-D,
Theorem~\ref{p1b:thm:all-order-schur}
\cite{Kugelmann2026Paper1D}, at each source value \(y\).
For increment \(\epsilon\), its residual is
\(\epsilon R_{\rm load}+o(\epsilon)\); its first two Taylor coefficients
are therefore \(L^\dagger H'L\) and
\(\tfrac12L^\dagger H''L-R_{\rm load}^\dagger C^{-1}R_{\rm load}\).
Contracting the latter with the fixed vector \(e\) gives \(c_2\).
All derivatives are of the actual moving source family.
\end{proof}

\begin{corollary}[Moving simple critical eigenvalue]
\label{p3f:cor:v199-moving-critical-eigenvalue}
If \(K(0)e=0\) is simple and
\(Q_eK(0)Q_e>0\) on \(\operatorname{ran}Q_e\), where
\(Q_e=I-ee^\dagger\), the moving eigenvalue instead satisfies
\begin{equation}
 \lambda''(0)=\langle e,K''(0)e\rangle
 -2\langle Q_eK'(0)e,(Q_eK(0)Q_e)^{-1}Q_eK'(0)e\rangle.
 \label{p3f:eq:v199-moving-eigenvalue}
\end{equation}
\end{corollary}
\begin{proof}
Use the simple-eigenvalue differentiation formula following Paper~I-D,
Theorem~\ref{p1b:thm:schur-jets} \cite{Kugelmann2026Paper1D}, at
\(\lambda=0\). The other eigenvalues are positive by the stated
\(Q_eK(0)Q_e>0\), so their sum is the displayed negative
reduced-resolvent term.
\end{proof}

Where the full stationary Hessian is invertible, the same action itself
computes the missing fourth-vertex contraction.  With
\(v=\bar q'\), \(a=\bar q''\), and \(H=S_{qq}(\bar q(y),y)\),
\begin{align}
 v&=-H^{-1}S_{qy},\qquad
 a=-H^{-1}\bigl(S_{qqq}[v,v]+2S_{qqy}v+S_{qyy}\bigr),\nonumber\\
 H''&=S_{qqyy}+2S_{qqqy}[v]+S_{qqqq}[v,v]+S_{qqq}[a].
 \label{p3f:eq:v199-fourth-vertex-contraction}
\end{align}
Twice differentiating stationarity and then the Hessian proves these
identities.  At a critical full Hessian, use the explicit interior reduction of
Corollary~\ref{p3f:cor:v202-critical-source-vertices}; only the
nonsingular interior block is inverted.  An unbounded-operator extension additionally
requires common form domains and differentiable inverses.

These formulas turn the stronger historical local-\(L\)/G26-GLOBAL calculation into an explicit comparison of
background curvature and interior relaxation.  They do not assert their
cancellation.  Paper~VI-A applies the trace-force part in
Theorem~\ref{p8:thm:v199-trace-force-vertices} and retains the complete
metric--matter boundary reduction.  Paper~III-G, Appendix~\ref{p3f:app:v199-affine-schur}
shows that, in the additional affine-Hessian class, one loading-vector test
is sufficient for an all-order linear response; affinity of an external
source in the action does not establish that class.

\subsubsection{Higher boundary vertices using only the interior inverse}
\label{p3f:sec:v202-higher-boundary-vertices}
Work on a fixed finite regular constraint quotient and smooth
trivialization, with the actual kinetic and source normalization.
Use real coordinates, realifying complex fields. Let $S(w,z)$ be the
actual $C^4$ action, where $w$ contains retained variables and any external
controls, $S_z=0$ at the base point, and $C=S_{zz}$ is nonsingular.
The full Hessian may be critical. Interior stationarity gives the local
branch $z_*(w)$ and
\[
 \Phi(w)=S(w,z_*(w)),\quad B=S_{zw},\quad Z=-C^{-1}B,
 \quad La=(a,Za),\quad T_j=D^jS.
\]

\begin{theorem}[Exact cubic and quartic boundary reduction]
\label{p3f:thm:v202-boundary-higher-vertices}
Define $T_{3,I}[v_1,v_2](\zeta)=T_3[(0,\zeta),v_1,v_2]$. Then
\begin{align}
 D^3\Phi[a,b,c]&=T_3[La,Lb,Lc],\nonumber\\
 D^4\Phi[a,b,c,d]&=T_4[La,Lb,Lc,Ld]\nonumber\\
 &\quad-\langle T_{3,I}[La,Lb],C^{-1}T_{3,I}[Lc,Ld]\rangle\nonumber\\
 &\quad-\langle T_{3,I}[La,Lc],C^{-1}T_{3,I}[Lb,Ld]\rangle\nonumber\\
 &\quad-\langle T_{3,I}[La,Ld],C^{-1}T_{3,I}[Lb,Lc]\rangle.
 \label{p3f:eq:v202-boundary-fourth-vertex}
\end{align}
Nonsingularity of $C$ suffices for these identities. If $C>0$, their
repeated-direction specialization is
\[
 D^4\Phi[a,a,a,a]=T_4[La,La,La,La]
 -3\|C^{-1/2}T_{3,I}[La,La]\|^2.
\]
\end{theorem}
\begin{proof}
Apply Paper~I-D, Theorem~\ref{p1b:thm:schur-jets}
\cite{Kugelmann2026Paper1D}, to the same-action stationary branch
\(S_z(w,z_*(w))=0\) and its actual tensors \(T_3,T_4\).
Its harmonic lift is the \(L\) defined here, so it gives all three
interior contractions. Only \(C\) must be nonsingular; \(C>0\)
is used for the final norm-square form.
\end{proof}

\begin{corollary}[Source derivatives at a critical full Hessian]
\label{p3f:cor:v202-critical-source-vertices}
Include the source in $w=(u,y)$ before interior Schur reduction. On an actual
boundary background $u=\bar u(y)$, set
$t=(\bar u',1)$, $a_y=(\bar u'',0)$,
$\widetilde e=(e,0)$, $\widetilde f=(f,0)$ for fixed boundary test vectors.
For $K_{ef}(y)=\Phi_{uu}(\bar u(y),y)[e,f]$,
\begin{align}
 K'_{ef}&=D^3\Phi[\widetilde e,\widetilde f,t],\nonumber\\
 K''_{ef}&=T_4[L\widetilde e,L\widetilde f,Lt,Lt]
 +T_3[L\widetilde e,L\widetilde f,La_y]\nonumber\\
 &\quad-\langle T_{3,I}[L\widetilde e,L\widetilde f],
                      C^{-1}T_{3,I}[Lt,Lt]\rangle\nonumber\\
 &\quad-2\langle T_{3,I}[L\widetilde e,Lt],
                       C^{-1}T_{3,I}[L\widetilde f,Lt]\rangle.
 \label{p3f:eq:v202-critical-source-fourth}
\end{align}
\end{corollary}
\begin{proof}
Differentiate $K_{ef}$ along $(\bar u(y),y)$ twice, then substitute
Theorem~\ref{p3f:thm:v202-boundary-higher-vertices}.
\end{proof}
The actual $\bar u'$ and $\bar u''$ remain inputs from the source
background. A moving simple eigenvalue still includes the boundary
relaxation term of Corollary~\ref{p3f:cor:v199-moving-critical-eigenvalue}.
A varying physical kinetic metric must be normalized with its own jets
before interpreting this coordinate Hessian as a physical eigenvalue.

\begin{theorem}[Exact all-order same-action Schur residual]
\label{p3f:thm:v202-all-order-schur-residual}
Let $H(y)=\left(\begin{smallmatrix}A&B^\dagger\\B&C\end{smallmatrix}\right)$
be the actual Hermitian Hessian family, with $C(y)$ nonsingular, and
$K(y)=A-B^\dagger C^{-1}B$. Put
\[
 Z_0=-C(0)^{-1}B(0),\quad L_0=\binom I{Z_0},\quad
 R(y)=B(y)+C(y)Z_0.
\]
Then
\begin{equation}
 K(y)=L_0^\dagger H(y)L_0-R(y)^\dagger C(y)^{-1}R(y).
 \label{p3f:eq:v202-fixed-lift-schur}
\end{equation}
If $H$ is differentiable at zero, this is equivalently
\begin{align}
 K(y)&=K(0)+yK'(0)+\mathcal E(y),\nonumber\\
 \mathcal E(y)&=L_0^\dagger[H(y)-H(0)-yH'(0)]L_0
              -R(y)^\dagger C(y)^{-1}R(y).
 \label{p3f:eq:v202-schur-residual}
\end{align}
For a scalar physical channel with the actual normalization $K(0)=0$,
$K'(0)=1$, the R1 law $K(y)=y$ on the interval is equivalent to
$\mathcal E(y)\equiv0$ there.
\end{theorem}
\begin{proof}
The two identities are Paper~I-D,
Theorem~\ref{p1b:thm:all-order-schur} \cite{Kugelmann2026Paper1D},
with \(R_0=R\), applied to this actual same-action Hessian throughout
its interior invertibility interval. For the physical normalization
\(K(0)=0\), \(K'(0)=1\), subtracting \(y\) leaves precisely
\(\mathcal E(y)\), proving the stated equivalence for the terminal response.
\end{proof}
If $C>0$ and the lifted term is independently proved affine, exact
linearity reduces to $R(y)=0$. Neither premise is inferred from a source
being affine in the action on a moving background. This closes the
algebraic compiler, while the actual native residual cancellation, R2
return positivity and occurring unbounded depth remain obligations only for that stronger local-\(L\)/G26-GLOBAL construction; they do not reopen native G26-DEEP.
Paper~III-G, Appendix~\ref{p3f:app:v202-vertex-diagnostics} checks the exchange factor
and shows why a vanishing second correction alone is insufficient.

\subsection{Generated Laurent coefficient character and dressed Born--Legendre descent}
\label{p3f:sec:v263-laurent-legendre-descent}
The preceding same-action constructor results distinguish the generated
coefficient-family tangent from the local configuration tangent.  This
separation closes the two remaining mathematical pieces of the depth-two
Laurent dressing without identifying the physical secant clock with an
additive process parameter.

\begin{theorem}[Generated positive coefficient character and Laurent pullback]
\label{p3f:thm:v263-laurent-coefficient-character}
On the generated PQTI coefficient line, let
\[
 Q_\sigma:=8\sqrt{\alpha_\sigma}>1,
 \qquad
 A_Q:=Q_\sigma\,\operatorname{id}.
\]
The positive functional calculus of the one-dimensional coefficient line
constructs the continuous one-parameter representation
\begin{equation}
\label{p3f:eq:v263-coefficient-character}
 \boxed{
 \rho_Q(s):=\exp\!\bigl(s\log A_Q\bigr)
 =Q_\sigma^s\,\operatorname{id},
 \qquad s\in\mathbb R.}
\end{equation}
It is the unique continuous positive representation satisfying
\[
 \rho_Q(s+t)=\rho_Q(s)\rho_Q(t),\qquad
 \rho_Q(0)=\operatorname{id},\qquad
 \rho_Q(1)=A_Q.
\]
If the generated depth-two Laurent exponent of Paper~III-D is denoted
\(n_2^+(R)\), its coefficient pullback is therefore
\begin{equation}
\label{p3f:eq:v263-D2-pullback}
 \boxed{
 D_2(R)\,\operatorname{id}
 :=\rho_Q\!\left(n_2^+(R)\right)
 =Q_\sigma^{\,n_2^+(R)}\operatorname{id}.}
\end{equation}
The multiplicative law belongs to the abstract exponent parameter \(s\).
It neither asserts nor requires additivity of the physical secant coordinate
\(n_2^+\) under concatenation of logarithmic intervals.
\end{theorem}

\begin{proof}
Every positive automorphism of a one-dimensional real coefficient line is
multiplication by one element of \(\mathbb R_{>0}\), so its real logarithm is
canonical.  Equation~\eqref{p3f:eq:v263-coefficient-character} therefore
defines a continuous positive representation and the exponential law gives
its multiplicativity.  Conversely, if \(\rho(s)=r(s)\operatorname{id}\) is
continuous and multiplicative, then \(\ell(s):=\log r(s)\) is a continuous
additive function and hence \(\ell(s)=cs\).  The condition
\(\rho(1)=A_Q\) gives \(c=\log Q_\sigma\), proving uniqueness.  Paper~III-D
constructs \(n_2^+\) as a scalar Laurent exponent, so evaluation of the
already generated representation at that scalar gives
Eq.~\eqref{p3f:eq:v263-D2-pullback}.  Paper~III-D,
Remark~\ref{p3d:rem:v263-exponent-character-pullback}, keeps this evaluation
separate from physical secant concatenation.
\end{proof}

\begin{theorem}[Historical depth-two fixed-depth local-jet Born--Legendre comparison]
\label{p3f:thm:v263-dressed-born-legendre-descent}
Put \(C:=C_7\).  Work on the certified action/contact and local-excess
branch of Definition~\ref{p3f:def:v2211-certified-response-domain}.
In particular the actual fixed-depth excess passes
Theorem~\ref{p3d:thm:v2211-excess-character-certificate}; a diagonal
secant value alone is insufficient.  Import the exact rank-seven secant readout
\[
 \mathscr H_7(R)=\frac{R^2+UC}{U+C},\qquad
 n_7(R)=\Pi_+(n_7^{\rm L}(R)),
\]
and the Paper~III-D local susceptibility
\begin{equation}
\label{p3f:eq:v263-n7-local-jet}
 \chi_7(R):=
 \frac{(1+n_7(R))(R^2-C^2)}{(U+C)R^2},\qquad \chi_7(C):=0.
\end{equation}
Let $D_2(R)=Q_\sigma^{n_2^+(R)}$ be the independently generated depth-two
coefficient character.  Since $d_{\rm cfg}$ does not differentiate the
scheme/depth coefficient, the same-action physical local Hessian is
\begin{equation}
\label{p3f:eq:v263-dressed-positive-response}
 \boxed{\chi_{2\mid7}(R):=D_2(R)\chi_7(R).}
\end{equation}
The unique terminal-normalized one-variable metric coefficient having this
local Hessian is
\begin{align}
\label{p3f:eq:v263-dressed-legendre-action}
 \bar f^{\rm L}_{2\mid7}(R)
 &=\boxed{R-C+\int_C^R(R-x)\chi_{2\mid7}(x)\,dx},\\
\label{p3f:eq:v263-dressed-euler-response}
 F(R)&=\boxed{1+\int_C^R\chi_{2\mid7}(x)\,dx},\\
 P_{2\mid7}(R):=RF-f
 &=\boxed{C+\int_C^R x\chi_{2\mid7}(x)\,dx}.
\end{align}
Hence
\[
 \boxed{f(C)=0,\quad F(C)=1,\quad f_{RR}=\chi_{2\mid7}>0\ (R>C),
 \quad f_{RR}(C)=0.}
\]
The Euler identity $\mathcal Ef=P_{2\mid7}$ and
$P_{2\mid7}'=R\chi_{2\mid7}$ hold identically.

On $C<R\le C_2$, $D_2=1$ and therefore $f_{RR}=\chi_7(R)$; the old
constant lower Hessian is recovered only in the high-curvature limit
$n_7\to1$, $C/R\to0$.  At K16, $D_2(C_2)=1$ on both sides, so $f$, $F$,
and $f_{RR}$ are continuous.  The coefficient freeze produces the finite
third-jet step
\begin{equation}
\label{p3f:eq:v263-k16-third-jet-step}
 \boxed{[f_{RRR}]_{C_2}
 =D_2'(C_2^+)\chi_7(C_2)>0}
\end{equation}
(upper minus lower).  At the terminal point the susceptibility vanishes
continuously, and the terminal Einstein representative $f_{\rm term}=R_g-C$
matches value, slope, and zero scalaron Hessian.
\end{theorem}

\begin{proof}
On its explicit excess-character certificate, Paper~III-D obtains
\(\chi_7\) by taking the local configuration derivative of the actual
Laurent excess at fixed coefficient/depth.  That certificate is exactly
the zero-defect differential test and diagonal normalization for the
printed germ; it is not inferred by subtracting a constant from a total
character.  Corollary~\ref{p3f:cor:generated-action-response-variation}
proves that this local derivative and the scheme/depth differential are
distinct tangent sorts.  The generated coefficient character can therefore
multiply the local Hessian after depth evaluation, but its depth derivative
is not a second local configuration derivative.  This gives
$\chi_{2\mid7}=D_2\chi_7$ with no additional response profile.

Twice integrating the supplied continuous Hessian with terminal data
$f(C)=0$, $F(C)=1$ gives the displayed unique coefficient.  Direct
differentiation proves the Euler identities.  Positivity follows from
$D_2>0$ and $\chi_7>0$.  At C2 the coefficient is continuous and equals one;
below C2 its background derivative is zero, while above it is finite and
positive, proving the third-jet formula.  Paper~III-D gives
$\chi_7(C)=0$, which proves the $C^2$ terminal match.
\end{proof}

\begin{corollary}[Discharge of the Laurent--Legendre structural G24 subfrontier]
\label{p3f:cor:v263-g24-dressing-placement-discharge}
On the certified action/contact and local-excess PQTI/Laurent branch,
\[
\boxed{\begin{gathered}
 \mathrm{DCHAR}_{2}^{\rm coeff},\quad\mathrm{ACT}_{2\mid7}^{\rm L},\\
 \mathrm{JUNC}_{2\mid7}^{\rm str,C^2},\quad C_7\text{ terminal scalar-remnant intercept}\\
 \text{are derived descendants}.
\end{gathered}}
\]
The coefficient-character, supplied-Hessian integrations and junction
identities are proved; physical owner realization additionally consumes
the two actual port certificates in
Definition~\ref{p3f:def:v2211-certified-response-domain}.
This corollary does not declare either unevaluated defect to vanish.
There is no independent free \(\mathrm{JUNC}_{2\mid7}^{\rm micro}\)
profile/width and no independent \(\mathrm{TERMREN}_7\) zero-order datum on
this branch.  A physical crossing remains G0/history content subject to
Theorem~\ref{p3f:thm:v266-gp-exact-k16-junction-condition}; the restored junction is a smooth-Hessian/finite-third-jet interface.  The raw G24 coefficient relation is closed as \(U_J^{\rm raw}=96U\), while the primitive momentum-curvature owner is fixed, while its EH identification is conditional: \(\mathcal K_p=6/\Rmom^2,\qquad U=(3/\pi)\mathcal K_p\quad\text{only under }(\mathrm{CELL}_{K\to U})\).  Absolute numerical unit realization is metrological rather than a structural premise.  This corollary does not identify
a nonadditive physical secant trajectory with the additive exponent parameter.
\end{corollary}


\subsection{source-conditioned same-action boundary descent}
\label{p3f:sec:v167-boundary-schur-corealization}

\begin{theorem}[Source-conditioned boundary Schur descendant]
\label{p3f:thm:v167-source-conditioned-boundary-schur}
Let \(S_\star[q;y]\) be one occurring action with a supplied source datum
\(y\), and let \(\bar q(y)\) be a stationary solution.  Suppose its frozen
configuration tangent is split by an actual boundary-domain construction as
\[
 T_{\bar q(y)}\mathcal C_{\rm stat}=J_\partial\oplus J_I,
\]
and write the same-action Hessian as
\[
 H_y=D_q^2S_\star\big|_{\bar q(y)}
 =\begin{pmatrix}
 H_{\partial\partial}(y)&H_{\partial I}(y)\\
 H_{I\partial}(y)&H_{II}(y)
 \end{pmatrix}.
\]
If \(H_{II}(y)\) is invertible on the declared gauge/constraint quotient,
then eliminating the interior fluctuation gives the unique quadratic
boundary descendant
\begin{equation}
 \label{p3f:eq:v167-boundary-schur}
 \boxed{
 K_{\partial,y}
 =H_{\partial\partial}(y)
  -H_{\partial I}(y)H_{II}(y)^{-1}H_{I\partial}(y).}
\end{equation}
It is the Hessian of the reduced boundary action obtained by stationary
interior substitution.  Hence its source dependence and any nonlocality
carried by the interior Green operator are generated by the same action;
neither requires a second boundary action or an independently fitted kernel.
This co-realization conclusion is conditional on the displayed physical
boundary split and trace map.  The abstract equality of two representation
types cannot create those data.
\end{theorem}

\begin{proof}
For a boundary fluctuation \(u_\partial\), stationarity of the quadratic form
in the interior variable gives
\[
 u_I=-H_{II}^{-1}H_{I\partial}u_\partial.
\]
Substitution into \(\frac12\langle(u_\partial,u_I),
H_y(u_\partial,u_I)\rangle\) yields
\(\frac12\langle u_\partial,K_{\partial,y}u_\partial\rangle\), proving
\eqref{p3f:eq:v167-boundary-schur}.  Equivalently, the implicit-function
theorem solves the interior stationary equation and two boundary variations
of the resulting reduced action give the same expression.  Every block,
inverse, and contraction therefore descends from \(D_q^2S_\star\) on the
source-conditioned solution.  The construction presupposes the stated split;
it does not generate a cross-sort trace from finite kinematics.
\end{proof}

\begin{corollary}[Regular dual boundary response is a restriction]
\label{p3f:cor:regular-dual-boundary-restriction}
On a finite-dimensional regular coordinate patch of the preceding
gauge/constraint quotient, assume also that \(H_y\) is invertible.
The local Legendre dual of the stationary reduced boundary action
is the full local dual action restricted to zero dual interior source.
Its Hessian is
\[
 \boxed{K_{\partial,y}^{-1}=(H_y^{-1})_{\partial\partial}.}
\]
The same statement holds on a Banach patch where the corresponding
implicit and inverse function theorems and bounded block inverses apply.
\end{corollary}
\begin{proof}
Apply Paper~I-C,
Theorem~\ref{p1b:thm:partial-elimination-legendre-duality}, to the
actual stationary action and its supplied split. Solving
\(H_y(u_\partial,u_I)^T=(p_\partial,0)^T\) retains the same
interior response as stationary substitution. No positive-definiteness
premise is needed for these regular local maps.
\end{proof}
A second Schur elimination of the inverse would instead give
\(H_{\partial\partial}^{-1}\), when that additional inverse exists.
At a critical boundary mode \(K_{\partial,y}^{-1}\) need not exist.
The source-to-boundary law, critical R1/R2 behavior and physical trace/history remain tasks only for the stronger local-\(L\)/G26-GLOBAL route identified by the parent theorem; native terminal G26-DEEP is already closed in the separate \(S+T\) modality.


\begin{corollary}[Boundary third variation is the K16 response analogue]
\label{p3f:cor:v167-boundary-third-variation}
Let \(\eta\) be a genuine source or boundary argument of the same reduced
action \(S^\partial_{\star}\).  Wherever the preceding inverse and
differentiation are defined,
\[
 \boxed{
 V^\partial(\eta)
 :=\nabla_\eta K_\partial
 =D_\eta D_{u_\partial}^2S^\partial_{\star}.}
\]
Thus the source-to-boundary-Hessian response is the boundary counterpart of
\(V_{16}=\delta^3S_\star\).  Its chain-rule expansion is fixed by the parent
second and third variations together with \(H_{II}^{-1}\).  A prescribed
linear critical lift, such as \(\mu^2(y)=y\) to all orders, is not implied by
the existence of this vertex; it remains a calculation of the actual action,
solution, boundary trace, and higher variations.
\end{corollary}

\begin{theorem}[Explicit first source jet of the stationary Schur descendant]
\label{p3f:thm:v176-explicit-source-jet-schur}
Under the hypotheses of
Theorem~\ref{p3f:thm:v167-source-conditioned-boundary-schur}, assume the
boundary/interior trivialization is fixed in the source parameter near
\(y=0\), the stationary family is differentiable, and the full Hessian is
invertible on the declared gauge/constraint quotient.  Write
\[
 J:=S_{qy}|_0,\qquad H:=S_{qq}|_0,\qquad
 W:=S_{qqy}|_0,\qquad V:=S_{qqq}|_0.
\]
Differentiating stationarity gives
\begin{equation}
 \label{p3f:eq:v176-stationary-source-velocity}
 \boxed{\dot{\bar q}_0=-H^{-1}J,}
\end{equation}
and therefore the first source jet of the same-action Hessian is
\begin{equation}
 \label{p3f:eq:v176-hessian-source-jet}
 \boxed{\dot H_0=W-V[H^{-1}J].}
\end{equation}
If
\[
 H_y=\begin{pmatrix}A_y&B_y\\ D_y&C_y\end{pmatrix},\qquad
 K_y=A_y-B_yC_y^{-1}D_y,
\]
then
\begin{equation}
 \label{p3f:eq:v176-schur-source-jet}
 \boxed{
 \dot K_0=
 \dot A_0-\dot B_0C_0^{-1}D_0
 -B_0C_0^{-1}\dot D_0
 +B_0C_0^{-1}\dot C_0C_0^{-1}D_0.}
\end{equation}
Thus the leading source response of the reduced boundary stiffness is a fixed
contraction of the parent second/third variations and the interior Green
operator.  No independently fitted boundary coefficient is introduced.
\end{theorem}

\begin{proof}
Differentiate
\(D_qS_\star(\bar q(y);y)=0\) at \(y=0\):
\[
 H\dot{\bar q}_0+J=0,
\]
which gives \eqref{p3f:eq:v176-stationary-source-velocity}.  Differentiating
\(H_y=D_q^2S_\star|_{\bar q(y)}\) then gives
\(\dot H_0=W+V[\dot{\bar q}_0]\), hence
\eqref{p3f:eq:v176-hessian-source-jet}.  Finally differentiate
\(K_y=A_y-B_yC_y^{-1}D_y\) and use
\(
\frac{d}{dy}C_y^{-1}=-C_y^{-1}\dot C_yC_y^{-1}
\), which yields \eqref{p3f:eq:v176-schur-source-jet}.
\end{proof}

\begin{remark}[Affine source insertion does not imply affine stationary stiffness]
\label{p3f:rem:v176-affine-source-firewall}
Even if the parent action is exactly affine in the source,
\(S[q;y]=S_0[q]+yS_1[q]\), the stationary point \(\bar q(y)\) generally moves
with \(y\).  The term
\(-V[H^{-1}J]\) in \eqref{p3f:eq:v176-hessian-source-jet} therefore feeds
pre-existing cubic vertices back into the Hessian and can generate quadratic
and higher source jets.  Source affinity, boundary support, or block-additive
matter sourcing alone does not prove an all-order law such as
\(P_0K_yP_0=yP_0\).  Paper~VI-A records explicit counterexamples and the
basis-relative G26 no-go in its appendix.
\end{remark}


\subsection{Noether-constrained terminal boundary handoff}
\label{p3f:sec:v167-noether-terminal-boundary-handoff}

\begin{theorem}[Ward descent through the same-action Schur reduction]
\label{p3f:thm:v167-ward-schur-descent}
Under the hypotheses of
Theorem~\ref{p3f:thm:v167-source-conditioned-boundary-schur}, let
\(R_A=(R_{A,\partial},R_{A,I})\) be an infinitesimal generated symmetry
direction of the stationary configuration for which the on-shell same-action
Ward identity gives
\[
 H_yR_A=0.
\]
If \(H_{II}(y)\) is invertible on the declared gauge/constraint quotient,
then the reduced boundary Hessian satisfies
\begin{equation}
 \label{p3f:eq:v167-ward-schur-null}
 \boxed{K_{\partial,y}R_{A,\partial}=0.}
\end{equation}
Hence stationary interior Schur reduction cannot destroy a Noether/Ward null
direction carried by the physical boundary component of the same generated
symmetry.  The conclusion is a relation descendant of the same action; it is
not an independent boundary conservation premise.
\end{theorem}

\begin{proof}
Writing \(H_yR_A=0\) in blocks gives
\[
 H_{\partial\partial}R_{A,\partial}+H_{\partial I}R_{A,I}=0,
 \qquad
 H_{I\partial}R_{A,\partial}+H_{II}R_{A,I}=0.
\]
The second equation yields
\[
 R_{A,I}=-H_{II}^{-1}H_{I\partial}R_{A,\partial}.
\]
Substitution into the first gives
\[
 \bigl(H_{\partial\partial}-H_{\partial I}H_{II}^{-1}H_{I\partial}\bigr)
 R_{A,\partial}=0,
\]
which is exactly \eqref{p3f:eq:v167-ward-schur-null}.  No extra charge,
current, or boundary action is introduced.
\end{proof}

\begin{theorem}[Eliminate-before-quotient rule for a rank-changing terminal mode]
\label{p3f:thm:v167-terminal-schur-order}
Let \(\varepsilon>0\) parametrize an active-side family of stationary
same-action Hessians on an actually constituted split
\(J_\partial\oplus J_I\), and write
\[
 H_\varepsilon=
 \begin{pmatrix}
  A_\varepsilon&B_\varepsilon^{*}\\
  B_\varepsilon&C_\varepsilon
 \end{pmatrix},
 \qquad C_\varepsilon^{-1}\ \text{defined on the reduced interior quotient}.
\]
Suppose a terminal role quotient removes an interior mode contained in
\(J_I\), while the active-side reduced boundary Hessian
\[
 K_{\partial,\varepsilon}
 =A_\varepsilon-B_\varepsilon^{*}C_\varepsilon^{-1}B_\varepsilon
\]
has a finite limit.  Then the quadratic boundary descendant obtained from the
same action by stationary Schur reduction before the terminal quotient is
\begin{equation}
 \label{p3f:eq:v167-terminal-schur-limit}
 \boxed{
 K_{\partial,T}
 :=\lim_{\varepsilon\downarrow0}
 \left(A_\varepsilon-B_\varepsilon^{*}C_\varepsilon^{-1}B_\varepsilon\right).}
\end{equation}
A quotient-first deletion that keeps only the limiting boundary block
\(A_T:=\lim_{\varepsilon\downarrow0}A_\varepsilon\) is equivalent to
\eqref{p3f:eq:v167-terminal-schur-limit} only if
\begin{equation}
 \label{p3f:eq:v167-schur-commutation-condition}
 \boxed{
 \lim_{\varepsilon\downarrow0}
 B_\varepsilon^{*}C_\varepsilon^{-1}B_\varepsilon=0,}
\end{equation}
or if an independently proved identity gives the same equality.  Thus the
terminal disappearance of an interior coordinate does not by itself license
deletion of its finite Schur memory.
\end{theorem}

\begin{proof}
For every \(\varepsilon>0\),
Theorem~\ref{p3f:thm:v167-source-conditioned-boundary-schur} fixes the unique
quadratic boundary descendant as \(K_{\partial,\varepsilon}\).  Taking the
stated finite active-side limit gives
\eqref{p3f:eq:v167-terminal-schur-limit}.  Quotient-first removal discards the
mixed term and therefore gives \(A_T\).  Their difference is exactly
\[
 A_T-K_{\partial,T}
 =\lim_{\varepsilon\downarrow0}B_\varepsilon^{*}C_\varepsilon^{-1}B_\varepsilon,
\]
whenever the displayed limits exist.  Equality is therefore equivalent to
\eqref{p3f:eq:v167-schur-commutation-condition}.  In particular, a vanishing
interior quadratic block does not imply a vanishing memory term when the mixed
block vanishes more slowly than its inverse-weighted contraction.
\end{proof}

\begin{corollary}[Terminal scalaron load is not yet a boundary stress]
\label{p3f:cor:v167-terminal-load-owner-firewall}
Let \(\mathcal L_C\) denote the active-side terminal restoring-load diagnostic
of the nonlinear scalaron branch in Paper~VI-A.  The same-action boundary
Schur calculus can consume \(\mathcal L_C\) only after an owner-preserving
source or trace map places that scalaron-side datum in the boundary/source
sort of \(S^\partial_\star\).  Without such a map, neither
Theorem~\ref{p3f:thm:v167-ward-schur-descent} nor
Theorem~\ref{p3f:thm:v167-terminal-schur-order} identifies
\(\mathcal L_C\) with a boundary stress, contact force, or RAR acceleration.
Release v1.87 separately supplies the local owner-preserving
\(B_{48}	o J_\partial\) map from the generated parent collar and PCAD.  This
corollary therefore no longer leaves a local G25 port open; it retains only
the firewall that the scalaron-side diagnostic \(\mathcal L_C\) is not thereby
identified with a physical boundary source.  Noether/Ward descent constrains
the resulting response but does not by itself prove an all-order critical law
such as \(\mu^2(y)=y\).
\end{corollary}

\begin{remark}[Exact scope of the terminal boundary result]
\label{p3f:rem:v167-terminal-boundary-scope}
The positive theorem here is an order-of-operations and relation-preservation
result for a boundary split that already exists physically.  The theorem
itself supplies no new boundary carrier, global ADM/Brown--York/Iyer--Wald
charge, terminal crossing history, or RAR critical exponent.  Release v1.87
separately generates the local \(B_{48}\)-to-boundary parent relation, so that
local port is no longer a downstream G26 premise.  For the stronger G26-GLOBAL/local-\(L\) route, the critical source law,
unbounded physical realization, and return/history data remain analytic,
action, and occurrence calculations rather than consequences of GP2 or
Noether alone.
\end{remark}


\section{Native Schur residue reconstruction and full metric compatibility}
\label{paper3f:sec:v280-native-schur-metric}
\subsection{Complete native response and residue reconstruction}
\label{p3f:sec:v226-response-reconstruction}
\begin{theorem}[Coefficient-complete native Schur response]
\label{p3f:thm:v226-response-reconstruction}
Let an actual same-action elimination supply a finite response $M(z)$ on
an admitted domain, and suppose its retained native fibre lies in
$\mathcal A_D$ (or in the tensor extension of
Corollary~\ref{p1b:cor:v226-source-jets}). Then
\[
 M(z)=\sum_{j=1}^{137}m_j(z)D_j,
 \qquad m_j(z)=\operatorname{Tr}(E_jM(z)).
\]
Holomorphic derivatives and residues at admitted isolated poles obey the
same inverse. Source and metric derivatives obey it in a fixed frame, or
with the common transported-frame covariant derivatives. Thus all native
operator, residue and consumed-jet matching is coefficient-complete.
\end{theorem}
\begin{proof}
Apply the complex-linear inverse of
Theorem~\ref{p1b:thm:v226-dual-coefficients} pointwise. In finite dimension
it commutes with differentiation and contour integration, including the
residue contour. The source-jet and tensor statements are the cited corollary.
\end{proof}
The source-loaded response is still the complete same-parent object,
including $BC^{-1}B^*$, moving incidence, contacts and boundary terms
where they occur. The coordinate algebra does not delete an odd source
leg whose completed return is even. A gapped or positive complement still
requires its actual domain and inverse. The spectral coefficient error bound
of Proposition~\ref{p2d:prop:v226-coefficient-spectral-error} applies only
after the common kinetic normalization has been supplied. These results
provide exact evaluation tests rather than an uncomputed electron member,
new occupied field, selected physical pole, or completed cosmic history.


\subsection{Finite exact certification of the complete response handoff}
\label{p3f:sec:v227-finite-handoff}
\begin{corollary}[The complete Schur handoff needs only finitely many energies]
\label{p3f:cor:v227-finite-handoff}
Two actual finite Hermitian complete Schur words with the same retained
and consumed source target, common kinetic normalization and complementary
sizes at most $m$ agree meromorphically once their full matrices agree at
$m+1$ distinct upper-half-plane energies. All consumed order-$k$ source
jets agree meromorphically after the enlarged tests of
Corollary~\ref{p1b:cor:v227-response-jets}. Equality then passes to their
admitted Green responses and contour poles/residues.
\end{corollary}
\begin{proof}
Use Theorem~\ref{p1b:thm:v227-finite-response}; use the common inverse
on its regular domain and integrate around common admitted pole contours.
For a native target, the preceding 137-coefficient inverse supplies every
sample matrix. Larger targets retain their actual complete frames.
\end{proof}
The existing source-cyclic, moment and cofactor routes are alternative
finite certificates on their printed domains; the present result does not
reopen their derivations. A coupled complementary zero must remain in a
real-zero calculation, although it obstructs none of these nonreal tests
(Corollary~\ref{p1b:cor:v227-zero-mode-comparison}). Equality of a kernel
without its actual source incidence is insufficient. In particular the
uncomputed external/dual/remainder source returns in the current charged
handoff are not replaced by numerical sample choices or by dark-sector
retirement.


\begin{proposition}[An exact scalar-isospectral ambiguity of the charged source grammar]
\label{p3f:prop:v246-charged-marked-isospectral-ambiguity}
Fix the rooted role and right frames
$U=(g_e,g_\mu,c)$ and $V^\sharp=(h_e,h_\tau,c)$, and positive
coefficients $e\ne m$, $t>0$. On the algebraic three-corner source domain put
\[
\begin{aligned}
 M_0&=eB_e+mB_\mu+tB_\tau, &M_\theta&=L_\theta M_0,\\
 L_\theta g_e&=\cos\theta\,g_e+\sin\theta\,g_\mu,\\
 L_\theta g_\mu&=-\sin\theta\,g_e+\cos\theta\,g_\mu,\quad L_\theta c=c.
\end{aligned}
\]
Then the three arrows $L_\theta B_r$ retain the same hierarchy corner
supports, rank one and Hilbert--Schmidt norm one. Every unordered singular
value and every scalar heat trace of the finite chiral operator
$h_\theta=\left(\begin{smallmatrix}0&M_\theta\\M_\theta^\dagger&0\end{smallmatrix}\right)$
is independent of $\theta$:
\[
 \operatorname{Tr}e^{-\tau h_\theta^2}
   =2(e^{-\tau e^2}+e^{-\tau m^2}+e^{-\tau t^2})
 \quad(\tau>0).
\]
Nevertheless the \emph{fixed rooted} left electron source
$j_e=(g_e,0)^T$ has the exact resolvent
\[
\boxed{\begin{gathered}
 j_e^\dagger(zI-h_\theta)^{-1}j_e
 =\cos^2\theta\frac{z}{z^2-e^2}+\sin^2\theta\frac{z}{z^2-m^2},\\
 \|(M_\theta M_\theta^\dagger-e^2I)g_e\|^2=(m^2-e^2)^2\sin^2\theta.
\end{gathered}}
\]
Thus three-corner membership, unit source-arrow norms, the entire scalar
spectrum and all uninserted scalar heat moments together do not imply
the marked electron response.
\end{proposition}
\begin{proof}
$L_\theta$ commutes with $Q_0$ and $Q_\perp$, so left multiplication
preserves all three source corners, ranks and norms. The polar Hopf
identities give $M_0M_0^\dagger=e^2P_e+m^2P_\mu+t^2Q_0$;
left multiplication gives its conjugate by $L_\theta$, while
$M_\theta^\dagger M_\theta=M_0^\dagger M_0$. These prove the heat formula.
The left compression of the chiral resolvent is
$z(z^2I-M_\theta M_\theta^\dagger)^{-1}$.
The fixed $g_e$ has squared overlaps $\cos^2\theta$ and $\sin^2\theta$
with the first two eigenlines. Spectral decomposition gives both displayed
identities.
\end{proof}
\begin{remark}[Exact scope of this obstruction]
\label{p3f:rem:v246-isospectral-owner-scope}
This is an insufficiency theorem for the stated algebraic data. It does not
assert that the frozen rooted constructor licenses every $L_\theta$, or
that two different complete GP2-generated parents have been constructed.
Nor does it assert equality of the complete heat functional with arbitrary
current, scalar-gradient, source or boundary insertions. Those mixed data
can distinguish the words. A simultaneous transformation of every source
and interaction is a presentation change; here the rooted probe is held
fixed, so the displayed difference is observable on the stated free chart.
To determine the physical word one must evaluate the independently generated
hierarchy arrows, rather than promote this diagnostic family to a physical
choice. In particular $\mathsf{Cert}_{\rm hier}$ is not replaced by scalar
heat-trace equality.
\end{remark}

\begin{proposition}[A one-contraction electron test on an independently certified spectral floor]
\label{p3f:prop:v246-electron-spectral-floor-test}
Let $B=MM^\dagger$ be the actual normalized action return and let
$g_e$ be the actual normalized tagged left source. If the same action has
\emph{independently} established $e^2I\le B\le\Lambda I$, then with
$d_0=g_e^\dagger Bg_e-e^2$ one has
\[
 \boxed{0\le\|(B-e^2I)g_e\|^2\le(\Lambda-e^2)d_0.}
\]
Consequently the exact electron mass-only response holds if and only if the
single contraction $d_0$ vanishes. The lower spectral bound is indispensable:
a mean equal to an interior spectral value need not make its probe an
eigenvector.
\end{proposition}
\begin{proof}
For $A=B-e^2I$, spectral calculus gives $0\le A^2\le(\Lambda-e^2)A$.
Contract with $g_e$. If $d_0=0$, positivity gives $A^{1/2}g_e=0$ and hence
$Ag_e=0$; the converse is immediate. The existing tagged chiral-resolvent
theorem then applies. If the bound is unavailable, use all three actual
contractions of Proposition~\ref{p3f:prop:v245-electron-actual-block-defect}.
\end{proof}

\subsection{Full metric and causal compatibility of a gravity parent}
\label{p3f:sec:v246-gravity-parent-compatibility}

\begin{proposition}[The endpoint cap mismatch is an interior metric Hessian defect]
\label{p3f:prop:v246-cap-metric-hessian-defect}
Let $f_+=\bar f^{\rm L}_{2\mid7}$ on its admitted lower collar, and
let $f_c$ satisfy the hypotheses of
Theorem~\ref{p3f:thm:v224-reference-cap-parent}, with $R_s^2=UC$.
Choose positive density normalizations $\nu_+,\nu_c$ and constants $a,b$
so the scalar densities
\[
 \Phi_+=\nu_+f_+,\qquad \Phi_c=\nu_cf_c+a+bR
\]
agree in value and first derivative at $C$. Let $g_r$ be a smooth
family of constant-scalar-curvature metrics on one fixed interior patch,
with $R[g_r]=r\downarrow C$ and regular limit $g_C$. Assume the
actual compactly supported metric test variations and pairings are
transported smoothly along this family. For
$I_\bullet[g]=\int\sqrt{|g|}\,\Phi_\bullet(R[g])$, one has
\[
 \boxed{\lim_{r\downarrow C}(D^2 I_c-D^2 I_+)_{g_r}[h,k]
 =k_C\int\sqrt{|g_C|}\,\delta R_C[h]\,\delta R_C[k],
 \qquad k_C=\nu_cf_c''(C)>0.}
\]
The right side has a nonzero rank-one fourth-order principal symbol
at every nonzero covector $\xi$:
\[
 k_C\bigl(\xi^\mu\xi^\nu h_{\mu\nu}-\xi^2
                 g_C^{\mu\nu}h_{\mu\nu}\bigr)^2.
\]
The irrelevant common Fourier sign in $\delta R$ disappears on squaring.
Compact-interior boundary contributions cannot cancel this defect.
An invertible derivative-free metric field change preserves its
nonzero principal rank. Positive whole-action rescaling preserves it.
\end{proposition}
\begin{proof}
Write $\Delta\Phi=\Phi_c-\Phi_+$. Twice differentiating
$\sqrt{|g|}\Delta\Phi(R[g])$ gives one term
$\sqrt{|g|}\Delta\Phi''\delta R[h]\delta R[k]$;
every other term contains $\Delta\Phi$ or $\Delta\Phi'$.
The latter vanish in the limit by endpoint alignment, whereas
$\Delta\Phi''\to\nu_cf_c''(C)$ because $f_+''(C)=0$.
The asserted regularity permits passing the limit through compact
test integrals. The two-derivative principal symbol of scalar-curvature
linearization is, up to its common sign,
$L_\xi(h)=\xi^\mu\xi^\nu h_{\mu\nu}-\xi^2\operatorname{tr}_g h$.
In dimension four $L_\xi$ is a nonzero linear form for every nonzero
$\xi$, also for a null Lorentzian covector. Thus $k_CL_\xi^*L_\xi$
has rank one. Under a derivative-free field chart with invertible
Jacobian $J$, the fourth-order symbol transforms by $J^*(\cdot)J$.
The additional chain-rule term pairs the action Euler covector with
the pointwise second derivative of the chart and has no fourth-order
principal part. Rank is therefore preserved. Boundary variations
vanish on the stated compact-interior tests.
\end{proof}

This comparison does not require an extension of the one-sided active
action below $C$, nor an inverse of its singular terminal Legendre
chart. It compares the regular interior limits. Unknown non-scalar,
shared-interior, source or measure terms could supply an opposite
principal coefficient; such a cancellation is an additional actual
mixed-variation identity, not a boundary adjustment or coordinate
normalization of these two scalar sectors.

\begin{proposition}[Sharp scaling required for a terminal scalaron cancellation]
\label{p3f:prop:v246-singular-scalaron-cancellation}
Consider a one-dimensional retained curvature port and a finite positive
scalaron block on $\chi>0$. Suppose
\[
 C_F^{-1}=\chi(T_0+o(1)),\qquad T_0>0,\qquad
 B_F=\chi^{-p}(B_0+o(1)),\quad B_0\ne0,
\]
and the direct coefficient tends to $k_C>0$. If its stationary Schur
coefficient tends to zero, then necessarily
\[
 \boxed{p=\tfrac12,\qquad B_0^*T_0B_0=k_C.}
\]
More generally, the bound $\|C_F^{-1}\|\le M\chi$ and convergence
$B_F^*C_F^{-1}B_F\to k_C>0$ imply
$\liminf\sqrt\chi\,\|B_F\|\ge\sqrt{k_C/M}$.
\end{proposition}
\begin{proof}
The return is
$\chi^{1-2p}(B_0^*T_0B_0+o(1))$. Its leading coefficient is strictly
positive. A finite nonzero limit therefore requires $1-2p=0$ and
the stated coefficient equality. For the general assertion use
$B_F^*C_F^{-1}B_F\le M\chi\|B_F\|^2$ and take the lower limit.
\end{proof}

The threshold is sharp as a block-algebra statement. The diagnostic
quadratic form
\[
 \frac12\begin{pmatrix}u&z\end{pmatrix}
 \begin{pmatrix}k_C+\chi&\sqrt{k_C/\chi}\\
                 \sqrt{k_C/\chi}&1/\chi\end{pmatrix}
 \binom uz
 =\frac12\chi u^2+\frac12(z/\sqrt\chi+\sqrt{k_C}u)^2
\]
has Schur coefficient exactly $\chi$. The chart $z=\sqrt\chi\,y$
makes its displayed quadratic polynomial regular, but its Jacobian
degenerates at $\chi=0$; a linear source $-j_z z$ becomes
$-\sqrt\chi j_z y$ and must be transported. This comparison neither
proves that such a singular mixed block occurs nor supplies a native
source or additional physical field. It identifies precisely one
escape from the bounded-block hypotheses of
the bounded-block Schur-decoupling hypothesis.

\begin{proposition}[Nonlocal causal memory is not a one-history action Hessian]
\label{p3f:prop:v246-causal-hessian-integrability}
Let $\Gamma[q]$ be a twice continuously differentiable scalar
functional of one real history with compact-interior test variations.
Suppose its Hessian has distribution kernel $K(t,s)$ and is retarded,
so $\operatorname{supp}K\subseteq\{t\ge s\}$. Then
\[
 \boxed{\operatorname{supp}K\subseteq\{t=s\}.}
\]
In particular a nonzero memory kernel away from the time diagonal in
Theorem~\ref{p3f:thm:v2216-source-response} cannot itself be this Hessian.
\end{proposition}
\begin{proof}
Symmetry of the second Fr\'echet differential gives $K(t,s)=K(s,t)$
distributionally. Retarded support and its transposed support then
intersect only on the time diagonal. A nonzero off-diagonal retarded
term contradicts that intersection.
\end{proof}

This obstruction concerns the integrated causal response with its
fixed preparation. A single local action for explicit metric and
source degrees of freedom is consistent with retarded solutions after
initial data are imposed. An effective action that has integrated out
those degrees of freedom needs its actual doubled-history or other
licensed causal construction. Neither operation is an equality between
the nonzero retarded kernel and the Hessian of an instantaneous
single-history $f(R)$ functional.


\begin{theorem}[An exact causal parent for the admitted finite source]
\label{p3f:thm:v246-finite-source-influence-parent}
On the actual finite source domain of
Theorem~\ref{p3f:thm:v2216-source-response}, let $H[q](t)=H(t,q(t))$ be
its complete time-local Hermitian generator, including its central contact,
in the fixed inherited clock and frame. Work on a finite interval
$[t_i,t_f]$, with histories in $C([t_i,t_f],O)$ in the uniform norm.
Assume $H$ and its first two control derivatives are jointly continuous
on a compact tube about the base history. Let $\rho_i$ be its supplied
positive trace-one initial
density, held fixed under the stated source variation. Write $U_q$ for
its unitary propagator from $t_i$ to $t_f$, with
$i\hbar\partial_tU_q=H[q](t)U_q$. Define
\[
 Z[q_+,q_-]=\operatorname{Tr}(U_{q_+}\rho_iU_{q_-}^\dagger),\qquad
 \Gamma[q_+,q_-]=-i\hbar\log Z[q_+,q_-].
\]
There is a unique continuous logarithm with $\log Z[q,q]=0$ on a
sufficiently small simply connected neighborhood of each diagonal
history. Thus $\Gamma[q,q]=0$, and on that neighborhood
\[
\begin{gathered}
 Z[q_-,q_+]=Z[q_+,q_-]^*,\quad |Z|\le1,\\
 \Gamma[q_-,q_+]=-\Gamma[q_+,q_-]^*,\\
 \operatorname{Im}\Gamma=-\hbar\log|Z|\ge0.
\end{gathered}
\]
Put
$q_\pm=q_c\pm q_\Delta/2$ and
$B=\partial_qH$, $Q=\partial_q^2H$ on the admitted local control
member, with $\langle B(t)\rangle=\operatorname{Tr}(\rho_i B_H(t))$
and the analogous convention for $Q$. Its exact diagonal physical derivatives are
\[
 \boxed{\left.\frac{\delta\Gamma}{\delta q_\Delta(t)}\right|_{q_\Delta=0}
       =-\langle B(t)\rangle,}
\]
\[
 \boxed{\left.\frac{\delta^2\Gamma}
 {\delta q_\Delta(t)\delta q_c(s)}\right|_{q_\Delta=0}
 =-\langle Q(t)\rangle\delta(t-s)
 +\frac{i}{\hbar}\Theta(t-s)
          \langle[B_H(t),B_H(s)]\rangle.}
\]
In the manuscript's $\hbar=1$ source units this is exactly the
complete causal source-action kernel already evaluated there. It is
a mixed Hessian of the doubled-history parent, whose transposed
branch block has advanced support; hence it obeys the symmetry of a
full second differential without making a single block time symmetric.
\end{theorem}
\begin{proof}
Finite-dimensional unitary evolution and smooth dependence on the
supplied coefficients give continuous $Z$ and its stated variations.
On the diagonal unitarity and trace normalization give $Z[q,q]=1$;
continuity therefore gives a convex product neighborhood with
$|Z-1|<1/2$, on which the principal logarithm has the stated regularity
and normalization. Swapping the two histories
conjugates the trace. Also
$Z=\operatorname{Tr}(\rho_iU_{q_-}^\dagger U_{q_+})$ is a convex
combination of diagonal matrix elements of a unitary in an eigenbasis
of $\rho_i$, so $|Z|\le1$. The normalized logarithm then gives the
stated conjugation and nonnegative imaginary part. Duhamel's identity
\[
 \delta U_q(t_f,t_i)
 =-\frac{i}{\hbar}\int_{t_i}^{t_f}
 U_q(t_f,t)\delta H(t)U_q(t,t_i)\,dt
\]
gives $\delta_+\Gamma=-\langle\delta H\rangle$ and
$\delta_-\Gamma=+\langle\delta H\rangle$ on the diagonal.
The difference-coordinate derivative is their half difference.
Differentiation in the common coordinate varies the actual unitary
history and the insertion. The former is the retarded commutator
term and the latter is the displayed $Q$ contact, by the same
Duhamel identity. This gives the boxed mixed kernel. Transposition
interchanges both time arguments and branch indices of the complete
second differential, proving the final assertion.
\end{proof}

This constructs a source-sector causal parent from the actual finite
generator and supplied preparation; it introduces no new state,
vacuum reset or continuum measure. Its extension through zeros of
$Z$, and its identification with a complete rooted gravity influence
functional, are not asserted. If a separately generated local metric
action $S_g$ is consumed on the same branch, the candidate combined
causal functional is $S_g[g_+]-S_g[g_-]+\Gamma[g_+,g_-]$ after the
actual source-to-metric pullback. The common-action firewall still
requires the native source incidence, normalization and boundary
contacts of that combination; writing it does not certify them.

\begin{corollary}[Exact finite mean-source Euler parent on a constituted common branch]
\label{p3f:cor:v246-finite-mean-source-euler-parent}
Suppose $S_0[q]$ is an actually constituted classical action on the
same source-control chart, with its inherited source incidence,
clock, physical normalization and boundary conditions. Work on a common
differentiable history subspace on which $S_0$ is $C^2$ and the source
functional above restricts continuously (for example sufficiently smooth
histories with the actual endpoint rule). Retain the actual $H[q]$ and fixed initial $\rho_i$ of the theorem, and assume
their source contribution has not already been included in $S_0$.
Then the locally defined doubled functional
\[
 \boxed{\mathcal I[q_c,q_\Delta]
 =S_0[q_c+q_\Delta/2]-S_0[q_c-q_\Delta/2]
                   +\Gamma[q_c+q_\Delta/2,q_c-q_\Delta/2]}
\]
generates the exact finite-source mean equation and its causal
linearization:
\[
 \left.\frac{\delta\mathcal I}{\delta q_\Delta(t)}\right|_0
 =\frac{\delta S_0}{\delta q(t)}-\langle B(t)\rangle,
 \qquad
 \left.\mathcal I_{\Delta c}\right|_0=D^2S_0-K_q^{\rm force}.
\]
In addition $\rho$ obeys the exact finite equation
$i\hbar\dot\rho=[H[q],\rho]$ with its supplied initial condition.
These are exact differential identities. Setting the mean Euler
expression to zero is the declared stationarity problem; existence or
occurrence of its self-consistent solution is not proved here.
No continuum measure or semiclassical expansion of the finite source
trace enters these identities. The classical control remains a
classical mean field; a quantum measure for that control is a separate
construction.
\end{corollary}
\begin{proof}
The difference-coordinate derivative of the classical difference is
the first variation of $S_0$; its common-coordinate derivative is the
full second variation with the actual boundary rule. Add the two
identities of the theorem. The density equation follows by
differentiating $\rho(t)=U_q(t,t_i)\rho_iU_q(t,t_i)^\dagger$.
\end{proof}

For the native source example one substitutes the already consumed
\[
 H_\eta[q]=d\Gamma\!\left(P_m[D_\gamma+e^qm\beta_D]P_m\right)
                      +c(q)I,
\]
on its fixed coherent frame and complete finite band, with the
actual inherited clock. Variations of the central $c(q)I$ contribute
to both $B$ and $Q$. A moving projector, frame, source marking or
preparation has its own contacts and must either be transported in
$H,\rho_i$ or retained as an independent source coordinate. The
statement with fixed initial $\rho_i$ does not license dropping such
contacts in a different physical preparation. The corollary closes
the finite causal integration step once its common-branch antecedents
are supplied; it does not select a previously unspecified $S_0$.


\section{Complete metric work and the common boundary parent}
\label{p3f:sec:v2203-workports}
The Hilbert definition already used in this paper makes stress the covector
conjugate to physical metric strain. In this section use $c=1$ and the
convention
\[
 \delta_gS_s=\tfrac12\int\sqrt{-g}\,T^{\mu\nu}\delta g_{\mu\nu}.
\]
The dimensional convention multiplies the Hilbert
conversion by $c_{\rm op}$ as in
\eqref{p3f:eq:matter-tensor-convention}. The following results complete the
work of an already occurring action. They do not constitute another stress
principle or identify a spectral carrier with a physical Hamiltonian.

\begin{proposition}[Work conservation with retained source marks]
\label{p3f:prop:v2203-workport-ward}
On the same-owner covariant domain of
Theorem~\ref{p3f:thm:v166-native-local-action-invariance}, suppose the
source action and all its marks admit the stated integrations by parts.
Define $\mathscr W_\nu$ by the complete field and external-mark Lie variation
$\delta_{\rm fields}S_s=\int\sqrt{-g}\,\mathscr W_\nu\zeta^\nu$
for an interior-supported vector field $\zeta$. Then
\[
 \nabla_\mu T^\mu{}_{\nu}=\mathscr W_\nu.
\]
On the dynamical field equations and without external work,
$\mathscr W_\nu=0$. For external scalar marks $y^A$ their term is
$\mathcal E_A\partial_\nu y^A$. Tensor, spinor and gauge marks consume
their full Lie-action adjoint, including derivative terms.
\end{proposition}
\begin{proof}
Insert $\delta g_{\mu\nu}=2\nabla_{(\mu}\zeta_{\nu)}$ in the
Hilbert variation, integrate once and add the complete field variation.
Covariance and arbitrary supported $\zeta$ give the identity. This is the
source-work specialization of the native Euler--Noether packet, with its
retained marks kept in the variation.
\end{proof}
For a smooth timelike boundary, outward spacelike normal $n$, induced
metric $\gamma$ and a scalar boundary mark $b$, the convention
$\delta S_\partial=\tfrac12\int\sqrt{|\gamma|}\tau^{ij}\delta\gamma_{ij}
+\int\sqrt{|\gamma|}\pi_b\delta b$ gives the on-bulk-equation identity
$D_i\tau^i{}_j=n_\mu T^\mu{}_{\nu}e^\nu{}_j+\pi_bD_jb$.
Normal displacement, corners and actual endpoints have their own terms.
The gravitational Euler derivative remains gravitationally allocated;
the variation of a total action cannot also be added as a second source.

\subsection{Hierarchy coefficients and moving characteristic data}
\label{p3f:subsec:v2203-workport-hierarchy}
Inherited $U,\alpha,C_j$ and grading weights are fixed in a physical metric
variation wherever the native tangent fixes them. An actual coefficient,
projector, source, measure, localizer or embedding map is differentiated
where it occurs. A constituted diagonal comparison family is different
from a local variation at fixed depth, as in
Proposition~\ref{p3e:prop:laur17-active-hessian}.
For example the constant four-dimensional bulk homothety
$g_s=e^{2s}g$ of $I=\int\sqrt{-g}f(R)$ gives
\[
 I'_0=\int\sqrt{-g}(4f-2RF),\qquad
 I''_0=\int\sqrt{-g}(16f-12RF+4R^2\chi),
 \quad F=f_R,\quad\chi=f_{RR}.
\]
This follows from $R_s=e^{-2s}R$ and uses the complete native running
$f$, with its inherited hierarchy. Nonconstant strain also retains the
derivative and endpoint terms. This calculation stays gravity-owned.

On a regular fixed-pairing chart let $H_{\rm car}=\sum_i h_iP_i$.
The full energy pullback, if this carrier actually occurs as the energy
insertion, is
\[
 \delta E=\operatorname{Tr}((\delta\sigma)H_{\rm car})
 +\sum_i\delta h_i\operatorname{Tr}(\sigma P_i)
 +\sum_i h_i\operatorname{Tr}(\sigma\delta P_i).
\]
The Hilbert contact fixes the independent state; its actual response is a
separate $\delta\sigma$ term. If the inherited $h_i$ are fixed and an
orthogonal projector orbit has skew-adjoint generator $A$, then
\[
 \delta H_{\rm car}=[A,H_{\rm car}],\qquad
 \delta E=\operatorname{Tr}([H_{\rm car},\sigma]A).
\]
Thus projector-only first work selects coherences between unequal grades.
For a two-line rotation with rate $r$ and diagonal populations $n_i,n_j$,
$\delta^2E=2(h_j-h_i)(n_i-n_j)r^2$. A commuting state has zero first
work without necessarily having zero second work. An uninserted spectral
trace is constant on the pure unitary orbit. The balanced hierarchy
weight $\sqrt{h_ih_j}$ therefore cannot replace the orbit vertex
$h_i-h_j$; their native distinction is
Theorem~\ref{p3c:thm:v179-balanced-hierarchy-response}.
Likewise a normalized current $J_D=Q_D\Omega_H$ has the complete tangent
\[
 \delta J_D=(\delta Q_D)\Omega_H+Q_D\delta\Omega_H.
\]
The moving amplitude cannot be discarded merely because the measure is
normalized. A dependent Euler characteristic witness does not replace
this existing current or its source jets.

\begin{proposition}[Symmetry-preserving work and the first broken jet]
\label{p3f:prop:v2203-workport-symmetry-onset}
Retain the unbroken homogeneous $C_3$ branch, the trivial $H$ role and
the irreducible real $T_5$ plane of
Theorem~\ref{p3f:thm:v1142-k48-h-t5-schur-zero}. If an admitted
scalar metric/source family preserves this branch and its equivariant
ports, its mixed block $B_{H,T}(g)$ is identically zero. Every defined
scalar-family jet of that particular Schur self-energy is therefore zero.
At a regular base point $B_0=0$, two admitted symmetry-breaking tangents
instead have
\[
 \delta_{12}Q_0=\delta_{12}A_0
 -\delta_1B_0^\dagger C_0^{-1}\delta_2B_0
 -\delta_2B_0^\dagger C_0^{-1}\delta_1B_0.
\]
For one tangent the relaxation correction is
$-2\delta B_0^\dagger C_0^{-1}\delta B_0$, negative semidefinite when
$C_0>0$; no such sign follows on an indefinite quotient.
\end{proposition}
\begin{proof}
The $C_3$ invariant map from its irreducible plane to its trivial line
vanishes at every point of the preserving family. For a broken tangent
differentiate $Q=A-B^\dagger C^{-1}B$ twice at $B=0$. Terms containing
undifferentiated $B$ vanish and leave the displayed terms.
\end{proof}
Invariant metric--$T$--$T$ vertices can still be nonzero and consume the
actual covariance. $C_3$ does not imply full $O(2)$ parity. Nor does a
frozen eigenvalue prove dynamical retention: the native terminal edge
eigenvector has a nonzero transverse source vertex, as in
Proposition~\ref{p3c:prop:v1101-terminal-edge-all-orders}. A moving support
needs its actual connection and $i\hbar\dot P=[H_{\rm full},P]$.

\begin{theorem}[Complete stationary boundary work]
\label{p3f:thm:v2203-complete-boundary-work}
Let $S_\star[g,b,z;y]$ be the occurring complete action, with actual source
marks $y$ and a constituted boundary/interior split. On a $C^3$ regular
constraint/gauge-reduced chart suppose $S_z=0$ has a differentiable
stationary solution $z_\star(g,b;y)$ and the relevant interior inverse
exists. Set $S_{\rm red}=S_\star[g,b,z_\star;y]$. For an admitted strain
$s$ and boundary tangents $u,v$ fixed in this chart,
\[
 D_gS_{\rm red}=(D_gS_\star)|_{z_\star},\qquad
 \boxed{D_b^2(D_sS_{\rm red})[u,v]
       =D_s(D_b^2S_{\rm red})[u,v].}
\]
This relation consumes all actual pairing, source, port, projector,
localizer, collar and boundary variations. In fixed reference coordinates
write $\widetilde K=\beta^\dagger M_\partial Q\beta$, where $\beta$ is
the physical port, $M_\partial$ the boundary pairing and $Q$ the
same-action Schur operator. Its complete metric tangent is
\begin{align*}
 \delta\widetilde K={}&(\delta\beta)^\dagger M_\partial Q\beta
 +\beta^\dagger(\delta M_\partial)Q\beta\\
 &+\beta^\dagger M_\partial(\delta Q)\beta
 +\beta^\dagger M_\partial Q(\delta\beta).
\end{align*}
Here the dagger is that of fixed reference coordinates; a changing
physical adjoint is thereby retained in $M_\partial$.
\end{theorem}
\begin{proof}
The chain term $S_zD_gz_\star$ vanishes in the stationary envelope.
Three derivatives commute on the assumed chart. For
$H=\left(\begin{smallmatrix}A&B^\dagger\\B&C\end{smallmatrix}\right)$,
$Q=A-B^\dagger C^{-1}B$ and $L=(I,-C^{-1}B)^T$, the native Schur
identity $\delta Q=L^\dagger(\delta H)L$ supplies its quadratic
realization. Differentiate the actual port and pairing product to obtain
the four displayed terms. The changing stationary background is already
part of $\delta H$, including the native source correction
$\dot H=W-V[H^{-1}J]$ of
Theorem~\ref{p3f:thm:v176-explicit-source-jet-schur}.
\end{proof}
This envelope applies to the action for which the eliminated variables
are stationary. If they are stationary for $S_g+S_s$ but not for $S_s$
alone, differentiating $S_s[g,z_\star(g)]$ also gives
$S_{s,z}D_gz_\star$. Source Hilbert work is first defined by partial
metric variation at fixed independent fields, then evaluated. A total
Schur kernel alone is not a source tensor. Moving quotient/projector
transport belongs to the same chain rule, and a differentiable common
domain is a hypothesis. Rank changes require the separate limit and
eliminate-before-quotient rule of
Theorem~\ref{p3f:thm:v167-terminal-schur-order}. Causal response uses the
actual dynamics or a licensed two-history influence action; varying a
single-history quadratic retarded word would symmetrize its kernel.

\subsection{Source-centre work and noncommuting collar variation}
\label{p3f:subsec:v2203-workport-collar}
Theorem~\ref{p3f:thm:v2109-poisson-source-projection} already reduces the
owned half-line form $\tfrac12\int_0^\infty(f'^2+k^2f^2)d\xi
-\int_0^\infty jf\,d\xi$, $f(0)=b$, $k>0$. For the admitted
fixed-collar source $j=Je^{-\ell\xi}$, $\ell>0$, its complete descendant
and comparison derivative are
\[
 S_{\rm red}=\frac12kb^2-\frac{Jb}{k+\ell}
 -\frac{J^2}{4\ell(k+\ell)^2},
\]
\[
 \partial_{k^2}S_{\rm red}
 =\frac{b^2}{4k}+\frac{Jb}{2k(k+\ell)^2}
 +\frac{J^2}{4k\ell(k+\ell)^3}.
\]
Indeed $f=be^{-k\xi}+J(e^{-\ell\xi}-e^{-k\xi})/(k^2-\ell^2)$;
its $k=\ell$ limit and the reduced expression are regular.
At $b=0$, $J\ne0$, the complete work is positive although the boundary
Hessian is always $k$. The source-linear and centre terms therefore
cannot be reconstructed from the boundary Hessian alone. The comparison
$k^2$ is not automatically four-dimensional Ricci curvature: an actual
metric tangent also varies normal maps, measure, source and ports and
consumes the inherited Lorentzian lapse normalization.

For a fixed positive product-collar realization $\mathcal N=P^{1/2}$,
with $P\ge p_*I>0$ and a bounded differentiable perturbation on the
stated fixed domain, $X=\delta\mathcal N$ obeys
\[
 \mathcal N X+X\mathcal N=\delta P,
 \qquad X=\int_0^\infty e^{-t\mathcal N}(\delta P)
                         e^{-t\mathcal N}\,dt.
\]
In its frozen eigenframe $X_{mn}=(\delta P)_{mn}/(\lambda_m+\lambda_n)$.
Only a commuting tangent gives $(2\mathcal N)^{-1}\delta P$.
The second mixed derivative obeys
$\mathcal N X_{12}+X_{12}\mathcal N=P_{12}-X_1X_2-X_2X_1$.
For bounded positive $\delta P$, $X\ge0$ and
$\|X\|\le\|\delta P\|/(2\sqrt{p_*})$. These statements follow by
differentiating $\mathcal N^2=P$ and integrating the Sylvester equation;
they are boundary-vertex bounds. General metric Laplacian tangents are
unbounded on $L^2$ and require their actual form/Sobolev-domain estimates.
A causal moving collar additionally consumes its sectorial domain and
physical port, rather than this frozen self-adjoint specialization.

\begin{proposition}[Two cosmological work ports and formation]
\label{p3f:prop:v2203-two-work-ports}
Retain $N$, $\xi=\log a$ and $V=2\pi^2a^3$ before the proper-time
restriction. On the complete source domain of
Proposition~\ref{p8:prop:v2185-c06-e6-functional-ports}, let
$E_s=\langle\mathcal O_N\rangle$ and
$W_s=N^{-1}\langle\mathcal O_\xi\rangle$. Then
\[
 u_s=E_s/V,\qquad p_s=-W_s/(3V),\qquad
 \Theta_s=\frac{\kappa}{V}(E_s+W_s),\quad
 \Xi_s=\frac{\kappa}{V}(E_s-W_s/3),
\]
where $\Theta_s=\kappa(u_s-3p_s)$ and $\Xi_s=\kappa(u_s+p_s)$.
In a proper-time coherent chart suppose an already allocated instantaneous
sector has energy insertion $H_D(a,z;t)$ and the state obeys
$\dot\sigma=-i[\widehat H,\sigma]/\hbar$ for the full generator. Put
\[
 \mathcal Q_D=V^{-1}\left\langle(\partial_tH_D)_{a,z}
 +\dot z^I\partial_IH_D+\frac{i}{\hbar}[\widehat H,H_D]\right\rangle,
 \qquad p_D=-\langle\partial_\xi H_D\rangle/(3V).
\]
Provided the displayed trace derivatives exist,
\[
 \dot u_D+3H(u_D+p_D)=\mathcal Q_D,\qquad
 E_D(t)=E_D(t_i)+\int_{t_i}^tV(\mathcal Q_D-3Hp_D)\,ds.
\]
Equivalently $\dot\Theta_D+3\dot\Xi_D+12H\Xi_D=4\kappa \mathcal Q_D$.
\end{proposition}
\begin{proof}
The first identities substitute the complete lapse/scale insertions.
For the instantaneous specialization trace cyclicity gives
$d\langle X\rangle/dt=\langle\dot X+(i/\hbar)[\widehat H,X]\rangle$.
Split the scale derivative from other work in $H_D$ and use
$\dot V=3HV$. Finally $\Theta_D+3\Xi_D=4\kappa u_D$ gives the
trace--enthalpy form.
\end{proof}
Exact coldness requires $W_D=0$, positive $E_D$, isotropy and zero flux.
A trace result alone cannot establish it: a conformal radiation block has
zero $\Theta$ with positive density and enthalpy. The three homogeneous
strains $(\delta\log N,\delta\xi)=(v,0),(v,-v),(v,v)$ pair respectively
with $-NvV u$, $-NvV(u+3p)$ and $-NvV(u-3p)$; hence lapse, active
static and Weyl work are different contractions. State evolution is not
differentiated against $a$ in the fixed-independent-state scale contact.
Interaction energy and centres are allocated once. A projected
$P_DHP_D$ alone does not establish an owned Hilbert stress split, and
velocity-dependent words need the complete functional Euler insertions.

\begin{corollary}[Finite physical pressure exclusions]
\label{p3f:cor:v2203-workport-pressure-bounds}
For complete Hermitian insertions $A=\mathcal O_N$ and
$B=N^{-1}\mathcal O_\xi$ on a finite physical Hilbert space and a positive
trace-one state $\sigma$, exact coldness $\operatorname{Tr}(\sigma B)=0$
requires $0\in[\lambda_{\min}(B),\lambda_{\max}(B)]$. For each real
$\lambda$ its energy obeys
\[
 \lambda_{\min}(A+\lambda B)\le E_D\le
 \lambda_{\max}(A+\lambda B).
\]
If $B$ is semidefinite, coldness forces state support into $\ker B$;
if definite it is impossible on that packet.
\end{corollary}
\begin{proof}
The expectation of a Hermitian finite matrix lies in its spectral
interval. Under coldness the expectation of $A+\lambda B$ is $E_D$.
For $B\ge0$, zero expectation implies
$B^{1/2}\sigma B^{1/2}=0$, hence support in its kernel; use $-B$ for
the other sign. These tests exclude states without claiming their
native history reaches any surviving state.
\end{proof}
Finite oscillator mode number is not finite Hilbert-space dimension;
an infinite Fock realization needs its own operator/form-domain version
and cannot borrow these finite spectral endpoints without a proof.

\subsection{Euler pullback, proper-time energy and the existing heat jet}
\label{p3f:subsec:v2203-workport-euler}
Theorems~\ref{p1:thm:v2203-weighted-euler-work} and
\ref{p1:thm:v2203-lift-transgression} give the weighted Euler and lifted
geometric derivative on their fixed-topology domains. Their physical
pullback retains the actual embedding $X$ and scalar weight:
$\delta h=X^*(\delta g)+X^*(\mathcal L_{\delta X}g)$ and
$\delta w=\delta(w\circ X)$. Normal displacement, moving domains and
endpoints must therefore be kept. A two-dimensional base strain does not
itself identify a three-dimensional FLRW volume strain.

If an actual reduction of the existing action supplies the proper-time
word
\[
 S_E=-\int dt\,N\varepsilon_7\int_{S^3}J_E,
 \qquad \int_{S^3}J_E=2,
\]
with uniform $\varepsilon_7(a,y)$, its homogeneous ports are
\[
 E_E=2\varepsilon_7,\qquad u_E=\frac{2\varepsilon_7}{V},\qquad
 p_E=-\frac{2\partial_\xi\varepsilon_7}{3V}.
\]
This action form is conditional. A topological term independent of the
metric and without a proper-time coupling
has zero Hilbert work. Constant $\varepsilon_7$ gives zero homogeneous
pressure and $a^{-3}$ dilution. With
$\nu_E=4\pi\kappa u_E/C_7$, imposing $\nu_E=2$ at every volume by
$\varepsilon_7=VC_7/(4\pi\kappa)$ instead gives
$u_E=C_7/(2\pi\kappa)$ and $p_E=-u_E$. Topology fixes neither the
dimensionful energy per charge nor a formation/observation history.
The inherited curvature Euler summand evaluated in the preceding
action results must retain its actual coefficient and companions; it
does not automatically realize this constant-energy target form.

\begin{proposition}[Complete Euler-marked heat third jet]
\label{p3f:prop:v2203-euler-heat-third-jet}
On the native localizer-aware parent $S=\operatorname{Tr}(Le^{-A})$,
$A=D^2/\Lambda_g^2$, let $E$ be a licensed source/metric derivative
whose Euler readout is being examined and $H$ a homogeneous metric
derivative. Assume a common $C^3$ trace/Duhamel domain. On a fixed-scale
chart the operator contacts are
\begin{align*}
 A_X&=\Lambda_g^{-2}\{D,D_X\},\\
 A_{XY}&=\Lambda_g^{-2}(\{D_X,D_Y\}+\{D,D_{XY}\}),\\
 A_{EHH}&=\Lambda_g^{-2}
 (2\{D_H,D_{EH}\}+\{D_E,D_{HH}\}+\{D,D_{EHH}\}).
\end{align*}
Write $I_k(B_1,\ldots,B_k)$ for the ordered simplex integral
\[
 \int_{t_i\ge0,\,\sum_{i=0}^kt_i=1}
 e^{-t_0A}B_1e^{-t_1A}\cdots B_ke^{-t_kA}\,dt_1\cdots dt_k
\]
and set $Z=e^{-A}$. Then
\begin{align*}
 Z_X={}&-I_1(A_X),\\
 Z_{XY}={}&I_2(A_X,A_Y)+I_2(A_Y,A_X)-I_1(A_{XY}),\\
 Z_{EHH}={}&-2I_3(A_E,A_H,A_H)-2I_3(A_H,A_E,A_H)\\
 &-2I_3(A_H,A_H,A_E)\\
 &+2\bigl(I_2(A_{EH},A_H)+I_2(A_H,A_{EH})\bigr)\\
 &+I_2(A_E,A_{HH})+I_2(A_{HH},A_E)-I_1(A_{EHH}),
\end{align*}
and the complete insertion is
\[
 \boxed{\begin{aligned}
 S_{EHH}=\operatorname{Tr}\bigl(&L_{EHH}Z+2L_{EH}Z_H+L_{HH}Z_E\\
 &+L_EZ_{HH}+2L_HZ_{EH}+LZ_{EHH}\bigr).
 \end{aligned}}
\]
\end{proposition}
\begin{proof}
Differentiate $D^2$ three times and apply Duhamel's formula. The
three orders of the singleton $E$ among two $H$ insertions each occur
twice. The $EH$ and $H$ partitions each occur twice and the $E$ and
$HH$ partitions once. The product rule for $LZ$ gives the six displayed
localizer terms. This extends
Proposition~\ref{p3f:prop:v118-localizer-aware-hessian} without invoking
commutation of the localizer with the heat factors.
\end{proof}
Actual variation of $\Lambda_g$ adds its product-rule contacts.
The Born nonmetric half-weight applies once to the complete source
packet. Stationary background, Schur and physical port derivatives are
taken after the parent insertion and retain the complete work theorem
above. An auxiliary Euler observable with no independently varied
action argument creates no $D_E$ or $L_E$; if it reads an actual metric
tangent, that tangent is the existing metric variation rather than an
added term. Thus the third-jet formula fixes the required calculation
where occurrence is present, without asserting a new Euler incidence,
stable mode, populated state or numerical abundance.


\subsection{Inherited Euler work, boundary momentum and complete source accounting}
\label{p3f:sec:v2203-euler-boundary-action}

This section composes the existing local gravity action, its canonical
endpoint, the exactly matched native scalar packet and the retained finite
source. It derives action and stress identities on their stated domains;
it does not select a new scalar support, source state, history or dark
abundance. Use $\hbar=c=1$, $M^2=U/4$, $C=C_7$ and
$\kappa=M^{-2}$. The running branch has $F=f_R$ and $\chi=f_{RR}>0$.
The terminal affine quotient is a separately licensed stratum, not a
regular inverse-$F$ chart at $\chi=0$.

\subsubsection{The same susceptibility fixes the anchored action work}
\begin{proposition}[Terminal anchored work identity]
\label{p3f:prop:v2203-anchored-work}
For the inherited representative with $f(C)=0$, $F(C)=1$ and
$P(r)=rF(r)-f(r)$, define $\mathcal J=P-CF$. Then
\begin{equation}
\label{p3f:eq:v2203-anchored-work}
 \boxed{\mathcal J(r)=\int_C^r(s-C)\chi(s)\,\mathrm ds,
 \qquad \mathrm d\mathcal J=(r-C)\,\mathrm dF,
 \qquad \mathcal J(C)=0.}
\end{equation}
On each regular active chart, $\mathcal J_F=r(F)-C$. In particular,
$\mathcal J\ge0$ for $r\ge C$ on the positive-susceptibility branch.
The complete gravity constituent can consequently be written
\begin{equation}
\label{p3f:eq:v2203-anchored-action}
 S_g=\frac{M^2}{2}\int\sqrt{-g}
       [F(R_g-C)-\mathcal J(F)]\,\mathrm d^4x
 +M^2\sum_\partial\epsilon\int\sqrt{|h|}\,FK_\partial\,\mathrm d^3x.
\end{equation}
Its bare auxiliary equation is $R_g=C+\mathcal J_F=r(F)$ when the
actual source has no independent $F$ variation.
\end{proposition}
\begin{proof}
The inherited integrals give
$F=1+\int_C^r\chi(s)\,\mathrm ds$ and
$P=C+\int_C^r s\chi(s)\,\mathrm ds$. Subtract $CF$ and differentiate.
Division by $F_r=\chi$ is licensed only on the regular active chart.
Substitution of $P=CF+\mathcal J$ in
Eq.~\eqref{p3f:eq:v169-jordan-action} gives the displayed action.
\end{proof}
The anchored work, curvature coefficient and endpoint therefore have one
normalization. A source curvature slot retains its complete variation
before restriction, as in
Theorem~\ref{p3f:thm:v169-complete-gravity-action}. This proposition does
not replace that source's $R_g$ by the bare $r(F)$. On the terminal affine
quotient $F=1$, $\mathcal J=0$, the gravity term is $R_g-C$; its empty
Einstein vacuum has $R_g=2C$. The hierarchy anchor $C$ is not that
vacuum's spacetime Ricci scalar. The work function $\mathcal J$ is also
distinct from a tangent Euler class or the mass-map homogeneity identity.

\begin{proposition}[Source-normalized geometric remainder and canonical work]
\label{p3f:prop:v2203-source-normalized-remainder}
On the complete bare $f(R)$ metric action, retain the full on-shell
homogeneous source $\rho_s,p_s$, assumed separately conserved on this
specialization. Let dots denote proper time. Comparing its same-action
metric equations to the fixed Einstein coefficient $M^2$, define
$\rho_{\rm tot}=\rho_s+\rho_\Lambda+\rho_\Delta$ and
$p_{\rm tot}=p_s-\rho_\Lambda+p_\Delta$, where
$\rho_\Lambda=M^2C/2$. Then
\begin{align}
\label{p3f:eq:v2203-source-normalized-remainder}
 \rho_\Delta&=\frac{M^2}{F}
      \left(\frac{\mathcal J}{2}-3H\dot F\right)
       -\frac{F-1}{F}\rho_s,\\
 p_\Delta&=\frac{M^2}{F}
      \left(\ddot F+2H\dot F-\frac{\mathcal J}{2}\right)
       -\frac{F-1}{F}p_s.\nonumber
\end{align}
The derivative density term is the canonical scalar work
$p_F\dot F/(FV)$, with $p_F=-3M^2VH$ from
Eq.~\eqref{p3f:eq:v169-flrw-momenta}. The full remainder is conserved,
$\dot\rho_\Delta+3H(\rho_\Delta+p_\Delta)=0$.
\end{proposition}
\begin{proof}
The two metric equations give the curvature contributions
$\rho_{\rm curv}=M^2(P/2-3H\dot F)/F$ and
$p_{\rm curv}=M^2(\ddot F+2H\dot F-P/2)/F$, together with
$\rho_s/F,p_s/F$. Insert $P=CF+\mathcal J$ and subtract the
unchanged comparison source and terminal vacuum. The momentum formula
proves the work identity. The metric Ward law and separately conserved
source give conservation of the full comparison remainder.
\end{proof}
If one keeps only the curvature-fluid convention instead, its actual
exchange is
$\dot\rho_{\rm curv}+3H(\rho_{\rm curv}+p_{\rm curv})
=\rho_s\dot F/F^2$. The source-rescaling terms are therefore necessary.
Exact $p_\Delta=0$ imposes the history condition
$\ddot F+2H\dot F=\mathcal J/2+(F-1)p_s/M^2$;
it is not a second adjustable constitutive law. This decomposition does
not create an additional independently evolving fluid. An actual residual
action or source with extra slots retains them before the comparison.

\subsubsection{The Euler contribution is already in the curvature action}
On the bare branch put $\mathcal A=M^2F/2$. On the genuinely retained,
exactly matched real scalar support of
Theorem~\ref{p3f:thm:v2202-hyperbolic-target}, put
\begin{equation}
\label{p3f:eq:v2203-curvature-owner}
 \rho_h^2=h^Ih_I=2\alpha_\Phi\operatorname{Tr}_F\Phi^2,
 \qquad \Omega=F-\frac{\rho_h^2}{6M^2}>0,
 \qquad \mathcal A=\frac{M^2\Omega}{2}
                 =\frac{M^2F}{2}-\frac{\rho_h^2}{12}.
\end{equation}
All scalar formulas consume the common-action overlap of
Proposition~\ref{p3f:prop:v1164-scalar-overlap-condition}, including its
transported normalization. The existing word is
$\int\sqrt{-g}\,\mathcal A R_g$ with the inherited canonical
$2\mathcal A K_\partial$ completion. No Euler coupling is appended.

\begin{theorem}[Inherited Hopf--Euler action incidence]
\label{p3f:thm:v2203-euler-incidence}
On a physical position-space Hopf slice, with spatially homogeneous
$\mathcal A(t)$ and fibre length $\ell(t)$, write
\begin{equation}
\label{p3f:eq:v2203-hopf-curvature}
 h_3=h_b+\ell^2\eta_0^2,\quad
 \int_{\rm fibre}\eta_0=2\pi,\quad
 \mathcal F=\mathrm d\eta_0,\quad
 R_3=R_2-\frac{\ell^2}{4}\mathcal F_{ij}\mathcal F^{ij}.
\end{equation}
The intrinsic spatial part of the same curvature action splits into
\begin{align}
\label{p3f:eq:v2203-euler-split}
 S_E&=8\pi^2\int N\ell\mathcal A Q_E\,\mathrm dt,
 &Q_E&=\frac1{4\pi}\int_{S^2}R_2\,\mathrm dA_b=2,\\
 S_{\mathcal F}&=-\frac\pi2\int N\ell^3\mathcal A
       \left(\int_{S^2}\mathcal F_{ij}\mathcal F^{ij}\,\mathrm dA_b\right)
       \mathrm dt.\nonumber
\end{align}
For the round physical sphere, $b=a/2$, $\ell=a$ and $V=2\pi^2a^3$,
\begin{equation}
\label{p3f:eq:v2203-round-euler-split}
 R_2=\frac8{a^2},\quad \mathcal F^2=\frac8{a^4},\quad
 L_E=16\pi^2Na\mathcal A,\quad
 L_{\mathcal F}=-4\pi^2Na\mathcal A,\quad
 L_{R_3}=12\pi^2Na\mathcal A.
\end{equation}
Both terms are inherited action constituents and must be retained together
with the temporal curvature and endpoint contributions.
\end{theorem}
\begin{proof}
The Hopf metric of Sections~\ref{p1:der:riemannian-submersion}
and~\ref{p1:subsec:coordinates-metric} gives the circle-submersion
curvature identity for spatially constant $\ell$. Its volume is
$\ell\eta_0\wedge\mathrm dA_b$; fibre integration gives $2\pi\ell$.
The oriented base Euler form is $e_h=R_2\mathrm dA_b/(4\pi)$, with
integral $Q_E=\chi(S^2)=2$. This gives the Euler summand and its
connection companion. For a round base of radius $b$,
$R_2=2/b^2$, $\mathcal F^2=1/(2b^4)$ and area $4\pi b^2$.
Set $b=a/2$, $\ell=a$ to obtain Eq.~\eqref{p3f:eq:v2203-round-euler-split},
including $R_3=6/a^2$.
\end{proof}

Already at $h=0$, the Euler term is $8\pi^2M^2NaF$. It is gravitational
metric work whose lapse sign and constraint are inherited from $S_g$;
it is not an independently additive positive matter density. Its scalaron
dependence lies in $F$, not in varying the integer $Q_E$. The geometric
Euler witness does not replace the native primitive source current
$J_D=Q_D\Omega_H$, and this physical position Hopf slice is not an
unproved identification with the intrinsic momentum source owner.

\begin{proposition}[Allocated scalar insertion and anisotropic work]
\label{p3f:prop:v2203-euler-scalar-work}
For the matched scalar constituent $\mathcal A_h=-\rho_h^2/12$, the
round uneliminated Euler term gives
\begin{equation}
\label{p3f:eq:v2203-euler-scalar-vertex}
 L_{E,h}=-\frac{4\pi^2}{3}Na\rho_h^2,
 \quad E_{E,h}:=-\partial_NL_{E,h}=\frac{4\pi^2}{3}a\rho_h^2,
 \quad
 -\partial_N\partial_{h^I}\partial_{h^J}L_{E,h}
       =\frac{8\pi^2}{3}a\delta_{IJ}.
\end{equation}
For independent base and fibre strains with fixed independent scalar data,
$V=8\pi^2b^2\ell$, the signed constituent energies
$E_E\propto\ell$ and $E_{\mathcal F}\propto\ell^3/b^2$ have
\begin{equation}
\label{p3f:eq:v2203-anisotropic-work}
 (p_b,p_\ell)_E=(0,-u_E),\qquad
 (p_b,p_\ell)_{\mathcal F}=(u_{\mathcal F},-3u_{\mathcal F}).
\end{equation}
At roundness $u_{\mathcal F}=-u_E/4$, so their sum has equal
pressures $p_b=p_\ell=-u_{R_3}/3$.
\end{proposition}
\begin{proof}
Substitute $\mathcal A_h$ in the inherited action split. For the
anisotropic work use
$2Vp_b=-\partial_{\log b}E$ and
$Vp_\ell=-\partial_{\log\ell}E$ at fixed independent fields; the
powers of $b,\ell$ prove the claim.
\end{proof}
At $h=0$ its classical energy and first scalar derivative vanish, while the
quadratic insertion remains on a retained support. A source-free scalaron
covariance is not thereby a native-$h$ covariance. The insertion precedes
stationary elimination and the full-parent completion. An isotropic strain
of the Euler term alone gives average $p_E=-u_E/3$; that does not make it
a separately isotropic fluid or the complete scalar equation of state.

\subsubsection{Temporal completion and the active boundary square}
For $ds^2=-N^2\mathrm dt^2+a^2\mathrm d\Omega_3^2$, let
$D_t=N^{-1}\mathrm d/\mathrm dt$ and $H=D_ta/a$. The completed curvature
word of Eq.~\eqref{p3f:eq:v2203-curvature-owner} reduces exactly to
\begin{equation}
\label{p3f:eq:v2203-completed-curvature}
 L_{\mathcal A R}=2\pi^2\left[
 -\frac{6a\mathcal A\dot a^2}{N}
 -\frac{6a^2\dot a\dot{\mathcal A}}{N}+6Na\mathcal A\right],
\end{equation}
where dots in this displayed lapse-preserving Lagrangian denote coordinate
time derivatives. The endpoint cancels the discarded total derivative,
as in Theorem~\ref{p3f:thm:v169-flrw-lapse-and-phase-action}.

\begin{proposition}[The complete nonminimal conformal constituent]
\label{p3f:prop:v2203-conformal-completion}
In a homogeneous coherent scalar frame, adding the canonical scalar kinetic
term to the $-\rho_h^2R_g/12$ constituent gives
\begin{equation}
\label{p3f:eq:v2203-conformal-completion}
 L_{\rm kin+NM}=\pi^2\left[
       \frac aN\left|\frac{\mathrm d(ah)}{\mathrm dt}\right|^2
       -Na|h|^2\right].
\end{equation}
Thus $v=ah$ and $\mathrm d\eta=N\mathrm dt/a$ give
$S_{\rm kin+NM}=\pi^2\int(|v'|^2-|v|^2)\,\mathrm d\eta$.
On the linearized native scalar around $h=0$, its actual negative quadratic
potential yields, for a position scalar harmonic of degree $j$,
\begin{equation}
\label{p3f:eq:v2203-native-conformal-mode}
 \boxed{v_j''+[(j+1)^2-2\Lambda_g^2a^2]v_j=0.}
\end{equation}
\end{proposition}
\begin{proof}
The scalar part of Eq.~\eqref{p3f:eq:v2203-completed-curvature} is
$\pi^2[a\rho_h^2\dot a^2/N+
2a^2\dot a\,h_I\dot h^I/N-Na\rho_h^2]$.
Its canonical kinetic term is $\pi^2a^3|\dot h|^2/N$; the three
velocity terms form $a|\mathrm d(ah)/\mathrm dt|^2/N$.
The native quadratic potential is $-\Lambda_g^2\rho_h^2$, whose
action contribution is $2\pi^2Na\Lambda_g^2|v|^2$.
The scalar harmonic gradient contributes $j(j+2)$, giving the result.
\end{proof}
The harmonic normalization can remove the common spatial volume factor
without changing the equation. The Euler curvature contribution is already
included: appending it again changes this same equation. Temporal gauge or
moving-frame slots retain their complete covariant derivatives and contacts.

Define the final-slice canonical curvature stress by
$\tau_{\partial,ij}=-2(\sqrt h)^{-1}\delta W/\delta h^{ij}$, with fixed
independent scalar endpoint data, and let
$\tau_\partial=h^{ij}\tau_{\partial,ij}$. Differentiating the completed
Lagrangian gives
\begin{equation}
\label{p3f:eq:v2203-active-boundary-trace}
 \boxed{\tau_\partial=\frac{ap_a}{V}
 =-12\mathcal A H-6D_t\mathcal A,
 \qquad \mathcal B_J=-\tau_\partial
 =3M^2(2\Omega H_J+D_t\Omega).}
\end{equation}
The lower endpoint has opposite orientation. At $h=0$ this is the bare
trace $-3M^2(2FH+D_tF)$; the scalar curvature allocation adds
$\rho_h^2H+h_ID_th^I$ to that trace. A source with additional metric
velocity slots retains its own endpoint momentum. The formula is the full
canonical trace only when those additional momenta vanish. Its dimension
is $M^2H$, rather than the bulk density dimension $M^2H^2$.

\begin{theorem}[Active density as a matched boundary square]
\label{p3f:thm:v2203-active-boundary-square}
On the regular matched $\Omega>0$ action and canonical endpoint domain,
put $T_J=|D_th|^2/2$, $W_J=M^2P/2+V_h$ and $K_J=k/a_J^2$.
Let $\rho_r$ be the complete instantaneous homogeneous residual source
density. For the displayed canonical kinetic constituent and residual
source without additional metric velocity terms, lapse variation before
fixing $N$ yields
\begin{equation}
\label{p3f:eq:v2203-matched-lapse}
 3M^2\Omega(H_J^2+K_J)+3M^2H_JD_t\Omega
       =W_J+T_J+\rho_r.
\end{equation}
Consequently
\begin{equation}
\label{p3f:eq:v2203-active-boundary-square}
 \boxed{\rho_r+T_J+V_h
 =\frac{\mathcal B_J^2}{12M^2\Omega}+3M^2\Omega K_J
 -\frac{3M^2(D_t\Omega)^2}{4\Omega}-\frac{M^2P}{2}.}
\end{equation}
Under $g_E=\Omega g_J$ and the same homogeneous physical preparation,
\begin{equation}
\label{p3f:eq:v2203-boundary-frame}
 \mathcal B_J=\Omega^{3/2}\mathcal B_E,\quad
 \mathcal B_E=6M^2H_E,\quad K_E=K_J/\Omega,
 \quad \psi=\sqrt{\frac32}M\log\Omega,
\end{equation}
and the constraint becomes
\begin{equation}
\label{p3f:eq:v2203-boundary-einstein-constraint}
 3M^2(H_E^2+K_E)
 =\frac12\left(\frac{\mathrm d\psi}{\mathrm dt_E}\right)^2
  +\frac1{2\Omega}|D_{t_E}h|^2
  +\frac{W_J+\rho_r}{\Omega^2}.
\end{equation}
The bare specialization sets $h=T_J=V_h=0$, $\Omega=F$.
\end{theorem}
\begin{proof}
Vary the completed curvature, scalar kinetic and potential terms and the
same source in the lapse to obtain Eq.~\eqref{p3f:eq:v2203-matched-lapse}.
The boundary square expands to
$3M^2\Omega H_J^2+3M^2H_JD_t\Omega+
3M^2(D_t\Omega)^2/(4\Omega)$; subtract its last term and the
potential to obtain the square identity. The homogeneous Weyl map gives
$H_E=\Omega^{-1/2}[H_J+D_t\Omega/(2\Omega)]$, proving the
boundary transport. Divide the lapse equation by $\Omega^2$ and use
$\tfrac12(\mathrm d\psi/\mathrm dt_E)^2=
3M^2(D_t\Omega)^2/(4\Omega^3)$ to obtain its Einstein form.
\end{proof}
This is the constraint realization of the hyperbolic kinetic target in
Theorem~\ref{p3f:thm:v2202-hyperbolic-target}. The canonical scalar work
already occurs in the boundary square; adding it a second time duplicates
the same action. On homogeneous Einstein variables the two spatial
constituents individually become
$L_E=8\pi^2M^2N_Ea_E$ and
$L_{\mathcal F}=-2\pi^2M^2N_Ea_E$. Their direct target derivatives
at fixed Einstein geometry vanish, while the temporal Weyl contacts supply
the scalar kinetic term and the full potential remains active. Nonuniform
conformal weights retain their spatial derivative ports. This does not
erase the allocated Jordan scalar insertion of
Eq.~\eqref{p3f:eq:v2203-euler-scalar-vertex}.

\subsubsection{Complete scalar trace and the cold-source restriction}
\begin{theorem}[Isolated classical native scalar cannot be a positive cold source]
\label{p3f:thm:v2203-isolated-native-scalar-cold-obstruction}
On the exactly matched classical on-shell native scalar owner, retain its
complete improved stress, positive real support and $\Lambda_g^2\ge0$.
If the residual scalar force is $J=0$, its homogeneous isotropic trace is
\begin{equation}
\label{p3f:eq:v2203-isolated-native-scalar-cold-obstruction}
 T_h=-u_h+3p_h=2\Lambda_g^2\rho_h^2\ge0,
 \qquad \boxed{u_h>0\ \Longrightarrow\ p_h\ge u_h/3>0.}
\end{equation}
It therefore cannot supply positive pressureless stress. At
$\Lambda_g^2=0$, the isolated classical identity is $p_h=u_h/3$.
\end{theorem}
\begin{proof}
Corollary~\ref{p3f:cor:v2202-improved-trace-handoff} gives
$T_h=4\alpha_\Phi\Lambda_g^2\operatorname{Tr}_F\Phi^2+
\operatorname{Tr}_F(\Phi J)$.
Canonical normalization gives $\rho_h^2=2\alpha_\Phi
\operatorname{Tr}_F\Phi^2\ge0$. Set $J=0$ and solve the trace relation
in the printed $(-+++)$ convention for $p_h$.
\end{proof}
The curvature, kinetic and quartic terms cancel in this on-shell trace,
not in their energy or variations. If an actual exchanged source is
present, positive cold stress in this scalar allocation instead requires
\begin{equation}
\label{p3f:eq:v2203-native-scalar-exchange-test}
 \operatorname{Tr}_F(\Phi J)=-u_h-2\Lambda_g^2\rho_h^2.
\end{equation}
This is a necessary test of the inherited source constructor, not permission
to choose a force. The actual mass-map identity of
Corollary~\ref{p3f:cor:v2202-mass-map-selector} supplies
$T_{h+D}=2\Lambda_g^2\rho_h^2+(q-1)\bar\psi\mathsf M\psi$
for a homogeneous degree-$q$ mass slot. An actually supplied degree-one
map has the same combined classical obstruction in the absence of other
trace terms. The theorem does not remove quantum anomalies, retained
subtractions, additional exchange, nonlocal terms or differently allocated
mixed stress. Those must consume their own complete action derivatives.

\subsubsection{First metric variation, reduction and causal source mean}
\begin{proposition}[Stationary reduction preserves the parent mean]
\label{p3f:prop:v2203-stationary-envelope}
Let $S[g,b,y]$ be one complete differentiable action on a regular stationary
reduction, with retained metric and boundary data $g,b$. If
$D_yS[g,b,y_*]=0$ including the consumed endpoints, then
\begin{equation}
\label{p3f:eq:v2203-stationary-envelope}
 W[g,b]=S[g,b,y_*(g,b)]\quad\Longrightarrow\quad
 D_gW=(D_gS)_{y_*}.
\end{equation}
The next derivative is its inherited Schur response, not another mean.
\end{proposition}
\begin{proof}
The complete chain rule gives $D_gW=D_gS+D_yS\,D_gy_*$; stationarity
annuls the last term. Differentiating the stationarity equation gives the
same-action Hessian of Corollary~\ref{p3f:cor:complete-reduction-reuse}.
\end{proof}
For the finite quadratic word $S=S_0+y^TCy/2+J_y^Ty$, $C=C^T$
invertible, the correction is $-J_y^TC^{-1}J_y/2$.
Its first metric correction vanishes at the unexcited classical point
$J_y=y_*=0$ even when its response Hessian does not vanish.

An actually admitted positive fixed-rank Euclidean Gaussian integration is
a different operation. With inherited measure factor $\mu$, mean
$m=-C^{-1}J_y$ and covariance $\Sigma=\hbar C^{-1}$, its effective
action has complete variation
\begin{equation}
\label{p3f:eq:v2203-quantum-envelope}
 \delta W_G=\delta S_0+\delta J_y^Tm+\tfrac12m^T\delta Cm
       +\tfrac12\operatorname{Tr}(\delta C\,\Sigma)
       -\hbar\delta\log\mu.
\end{equation}
The covariance term is the same derivative as
$\hbar\log\det C/2$, rather than an additional quantum mean. Classical
stationary reduction does not automatically supply this integration,
its causal contour or preparation. In particular the formula does not
reset a propagated Lorentzian covariance to $\hbar C^{-1}$. Constituent,
measure and endpoint matching remain necessary before comparing it to
the determinant already allocated to $f$ or the retained source.

\begin{proposition}[The complete two-history source gives both mean ports]
\label{p3f:prop:v2203-causal-mean-ports}
For the same normalized initial density $\sigma_i$ and admitted complete
unitary generator, set
\begin{equation}
\label{p3f:eq:v2203-causal-source}
 \mathcal I[x_+,x_-]=\operatorname{Tr}(U[x_+]\sigma_iU[x_-]^\dagger),
 \quad\Gamma_s=-i\hbar\log\mathcal I,\quad
 x_c=(x_++x_-)/2,\quad x_d=x_+-x_-.
\end{equation}
Continue the logarithm near equal histories. On the inherited instantaneous
homogeneous fixed-frame word $\widehat H=N\widehat H_0(a)$,
$\lambda=\log a$, its complete mean ports are
\begin{equation}
\label{p3f:eq:v2203-causal-mean-ports}
 \left.\frac{\delta\Gamma_s}{\delta N_d}\right|_{x_d=0}=-V\rho_s,
 \qquad
 \left.\frac{\delta\Gamma_s}{\delta\lambda_d}\right|_{x_d=0}
       =3NVp_s.
\end{equation}
Thus the active metric direction
$(\delta\log N,\delta\lambda)=(\varphi,-\varphi)$ reads
$-NV(\rho_s+3p_s)$, rather than the improved trace $-\rho_s+3p_s$.
\end{proposition}
\begin{proof}
The complete unitary variation is
$\delta U=-(i/\hbar)\int U(t_f,t)\delta\widehat H(t)
U(t,t_i)\,\mathrm dt$. At equal histories trace cyclicity gives
$\delta\Gamma_s=-\int\operatorname{Tr}(\sigma\delta\widehat H[x_d])
\,\mathrm dt$. Its lapse and scale derivatives are precisely the source
work of Theorem~\ref{p3f:thm:v169-finite-source-work}, including the
central energy and its actual scale derivative. Combine them for the
active metric direction.
\end{proof}
Although $\Gamma_s[x,x]=0$, its difference-direction derivative gives
the physical mean. Derivative-dependent or moving source words retain the
functional insertion and endpoint operator of
Theorem~\ref{p4:thm:v4106-functional-hilbert-response}, with every
normalization, projector, preparation and frame contact. A subsequent
common-history derivative gives causal response; a connected second
difference-history cumulant gives noise. Neither is another first metric
variation. Nonzero field covariance or connected noise alone therefore
does not establish an extra mean energy.

\begin{proposition}[Homogeneous collisionless driven memory cancels]
\label{p3f:prop:v2203-collisionless-memory-cancellation}
On the admitted constant-$F_\star$ collisionless tangent of
Corollary~\ref{p3f:cor:v1144-collisionless-source-schur}, suppose the
homogeneous isotropic occupation $\bar f(q)$ and its homogeneous scalar
port $D(q)$ depend on the same spatial kinetic Casimir $q$. For $w=w(t)$,
\begin{equation}
\label{p3f:eq:v2203-collisionless-memory-cancellation}
 \mathcal C_0w=\{\bar f(q),D(q)w(t)\}_{x,p}=0,
 \qquad \mathcal B\bar{\mathscr L}_R^{-1}\mathcal C_0=0.
\end{equation}
The equal-time contact, original occupation stress and independent initial
occupation perturbation columns remain.
\end{proposition}
Here $\mathcal B$ denotes the inherited kinetic return operator, rather
than the scalar endpoint momentum used elsewhere in this section.
\begin{proof}
Functions of the same canonical Casimir Poisson commute, also in curved
spatial coordinates. Substitute into the inherited kinetic memory block.
\end{proof}
This is a homogeneous collisionless specialization; inhomogeneous ports
need not have the same cancellation.

\begin{proposition}[Exact frozen free-sign-pair response with its contact]
\label{p3f:prop:v2203-frozen-dirac-active-kernel}
On one admitted complete free Dirac sign pair of
Eq.~\eqref{p3f:eq:v169-explicit-pair-block}, use $\hbar=1$,
$H_0=p\sigma_1+m\sigma_3$, $p=|\lambda|/a$, $m>0$ and
$E=\sqrt{p^2+m^2}$. For the active metric leaf
$(N,a)\mapsto(e^\varphi N,e^{-\varphi}a)$ on a frozen member, the
noncentral first and second generator jets are
\begin{equation}
\label{p3f:eq:v2203-frozen-dirac-active-jets}
 D=2p\sigma_1+m\sigma_3,\qquad Q=4p\sigma_1+m\sigma_3.
\end{equation}
The same prescribed centre has
$H_{\rm cen}(\varphi)=e^\varphi\varepsilon_0(ae^{-\varphi})$, with
jets $\varepsilon_0-a\varepsilon_0'$ and
$\varepsilon_0-a\varepsilon_0'+a^2\varepsilon_0''$. They remain in
the mean and contact, cancelling only from commutators.
For the stationary normalized pair state
$\bar\sigma=(I+zH_0/E)/2$, $-1\le z\le1$, its covector kernel is
\begin{align}
\label{p3f:eq:v2203-frozen-dirac-active-kernel}
 \Pi^R(t)&=2z\frac{m^2p^2}{E^2}\Theta(t)\sin(2Et),\\
 \Pi^R(\omega)&=\frac{4zE(m^2p^2/E^2)}
       {4E^2-(\omega+i0)^2},\qquad
 \Pi^R(0)=z\frac{m^2p^2}{E^3}.\nonumber
\end{align}
The retarded transform is $\int_0^\infty e^{i(\omega+i0)t}\Pi^R(t)
\,\mathrm dt$. Without the central term, the full zero-frequency
contact plus memory is
\begin{equation}
\label{p3f:eq:v2203-frozen-dirac-contact-memory}
 \langle Q\rangle+\Pi^R(0)
 =z\frac{4p^4+6p^2m^2+m^4}{E^3}
 =\left.\partial_\varphi^2[zE(\varphi)]\right|_{\varphi=0}.
\end{equation}
\end{proposition}
\begin{proof}
In the energy basis $D=d_\parallel\tau_3+d_\perp\tau_1$, with
$d_\parallel=(2p^2+m^2)/E$ and $d_\perp^2=m^2p^2/E^2$.
Use $e^{iE\tau_3t}\tau_1e^{-iE\tau_3t}
=\tau_1\cos(2Et)-\tau_2\sin(2Et)$ in
$-i\Theta(t)\operatorname{Tr}(\bar\sigma[D(t),D(0)])$.
Retarded integration gives the frequency formula. Finally
$\langle Q\rangle=z(4p^2+m^2)/E$; adding the kernel gives the
same instantaneous eigenenergy curvature.
\end{proof}
The transition poles are $\pm2E$, with no zero pole. In a late frozen
massive member $\Pi^R(0)/V=O(a^{-5})$, whereas retained rest-energy
mean can be $O(a^{-3})$. This is not a uniform bound on an arbitrary
nonstationary two-time kernel: earlier driving can persist in its state.
The actual force density still consumes its inherited Hilbert/Riesz,
measure, projector and frame contacts. No source state, mass or central
allocation is selected by the frozen response calculation.

\begin{corollary}[No second mean on the completed minimal domain]
\label{p3f:cor:v2203-no-second-minimal-mean}
For the specifically completed minimal gravity constituent and the same
complete retained finite source, reduced on its actual causal and canonical
endpoint domain, exact constituent matching gives
\begin{equation}
\label{p3f:eq:v2203-no-second-minimal-mean}
 \Gamma_{\rm ex}:=\Gamma_{\rm red}
       -(S_g[x_+]-S_g[x_-])-\Gamma_s=0,
 \qquad \rho_{\rm ex}=p_{\rm ex}=0.
\end{equation}
\end{corollary}
\begin{proof}
The reduced minimal action is exactly the sum being subtracted.
Proposition~\ref{p3f:prop:v2203-causal-mean-ports} verifies that its
source subtraction includes the complete mean, including central
allocation. The Euler constituent of $S_g$ is already in that sum.
\end{proof}
This corollary concerns the supplied minimal reduction, not every completed
parent. A genuinely nonzero overlap-matched parent remainder must be
derived before its first metric variation is called an additional mean.
It may not be appended by retaining the same Euler constituent,
determinant, causal mean or endpoint a second time.

\subsubsection{Terminal boundary budget and the separate cold condition}
On the separately licensed affine gravity stratum $F=1$, $D_tF=0$,
$P=C$, define the bare gravitational scale $\mathcal B=6M^2H$.
It equals the full canonical stress scale only if the actual sources have
no additional metric endpoint momentum; it is not the general matched
$\Omega$ scale of Eq.~\eqref{p3f:eq:v2203-active-boundary-trace}.

\begin{theorem}[Terminal residual boundary budget]
\label{p3f:thm:v2203-terminal-boundary-budget}
Let $\rho_s,p_s$ be the already counted comparison source, and
$\rho_{\rm ex},p_{\rm ex}$ a genuinely admitted additional mean on the
same terminal metric and causal domain. Put
$\rho_\Lambda=M^2C/2$ and $\mathcal B_\infty^2=6M^4C$.
The full lapse and scale work imply
\begin{align}
\label{p3f:eq:v2203-terminal-boundary-budget}
 \rho_{\rm ex}
 &=\frac{\mathcal B^2-\mathcal B_\infty^2}{12M^2}
         +3M^2K-\rho_s,\\
 p_{\rm ex}
 &=-\frac{\dot{\mathcal B}}3-\frac{\mathcal B^2}{12M^2}
       -M^2K+\rho_\Lambda-p_s.\nonumber
\end{align}
Here dots denote proper time. Positive density and zero pressure are
separate readouts; exact cold additional stress requires
\begin{equation}
\label{p3f:eq:v2203-terminal-cold-trajectory}
 \boxed{\dot{\mathcal B}=3\rho_\Lambda
       -\frac{\mathcal B^2}{4M^2}-3M^2K-3p_s.}
\end{equation}
On an admitted expanding affine vacuum with $a\to\infty$ and
vanishing nonvacuum density, $\mathcal B\to\mathcal B_\infty>0$.
At finite closed-vacuum radius, however,
$\mathcal B^2=\mathcal B_\infty^2-36M^4K$; the asymptotic scale
is not a pointwise lower bound.
\end{theorem}
\begin{proof}
The complete terminal metric equations are
$3M^2(H^2+K)=\rho_\Lambda+\rho_s+\rho_{\rm ex}$ and
$-M^2(2\dot H+3H^2+K)=-\rho_\Lambda+p_s+p_{\rm ex}$.
Insert $H=\mathcal B/(6M^2)$, then set $p_{\rm ex}=0$ for the
cold condition. The source-free asymptotic constraint gives
$H_\infty^2=C/6$, and the finite-radius vacuum equation retains $K$.
\end{proof}
These are geometric constraint readouts, not a source production law or
a selected transition to the terminal stratum. The canonical endpoint has
no compactly supported interior variation that could add a volume source.
The exact RAR normalization and the proposed abundance's separate boundary
test consume this same terminal budget in Paper~VI-B.

If the comparison source is separately conserved, define
\begin{equation}
\label{p3f:eq:v2203-conserved-cold-charge}
 \mathscr D=a^3\left[H^2+K-\frac C6-\frac{\kappa\rho_s}{3}\right].
\end{equation}
Then $E_{\rm ex}=3M^2V_\gamma\mathscr D$ and
$\dot{\mathscr D}=-\kappa Ha^3p_{\rm ex}$. On an expanding interval,
exact cold additional stress is equivalent to constant $\mathscr D$;
the integration constant and observation slice remain actual history data.
The independent trace test is
$R_g=2C-\kappa T_s-\kappa T_{\rm ex}$.
For the isolated classical native scalar its nonnegative trace agrees with
Theorem~\ref{p3f:thm:v2203-isolated-native-scalar-cold-obstruction}, rather
than supplying a cold negative trace. Anomalies, retained subtractions and
actual exchanges require their own complete trace terms.

\subsubsection{An existing source can retain cold comoving energy}
The local geometric-duration restriction of
Theorem~\ref{p8:thm:v2203-local-duration} concerns an additional
final-local geometric remainder after the original source is counted.
It does not delete the already counted source's state and history energy.

\begin{theorem}[Cold survival on the complete finite local source]
\label{p3f:thm:v2203-local-source-cold-tail}
On the inherited fixed-rank complete finite CAR source of
Eq.~\eqref{p3f:eq:v169-finite-free-source}, work in a fixed coherent
frame and proper time with
\begin{equation}
\label{p3f:eq:v2203-local-source-cold-generator}
 \widehat H_0(t)=H_m+a(t)^{-1}\widehat D+\varepsilon_0(a(t))I,
 \quad H_m=d\Gamma(m\beta_D),\quad
 \widehat D=d\Gamma(\mathcal D_\gamma),\quad m>0.
\end{equation}
All operators are restricted to the same inherited complete band. Suppose
$a\to\infty$, $\int_{t_i}^{\infty}a^{-1}\,\mathrm dt<\infty$,
$\varepsilon_0\to\varepsilon_\infty$ and
$a\varepsilon_0'(a)\to0$. For any supplied normalized initial density,
the $H_m$ interaction-picture state has a limit $\sigma_\infty$, and
\begin{equation}
\label{p3f:eq:v2203-local-source-cold-tail}
 \boxed{V\rho_s\longrightarrow
 E_\infty=\operatorname{Tr}(\sigma_\infty H_m)+\varepsilon_\infty,
 \qquad Vp_s\longrightarrow0.}
\end{equation}
If its physically allocated $E_\infty>0$, then
$a^3\rho_s\to E_\infty/V_\gamma$ and $p_s/\rho_s\to0$.
An exact local FLRW interpretation additionally consumes homogeneous
energy, isotropic stress and vanishing spatial flux as in
Theorem~\ref{p3f:thm:v169-finite-source-work}.
\end{theorem}
\begin{proof}
Remove the central phase for state transport only, keeping its energy
and metric work. In the $H_m$ interaction picture the perturbation is
the unitary conjugate of $a^{-1}\widehat D$. Its norm is integrable
on the finite CAR Hilbert space, so the interaction propagator is Cauchy
in operator norm and converges to a unitary $W_\infty$. Thus
$\sigma_\infty=W_\infty\sigma_iW_\infty^\dagger$.
The noncentral spatial energy is bounded by $\|\widehat D\|/a$;
the mass energy converges and the retained centre supplies
$\varepsilon_\infty$. The exact scale work gives
$3Vp_s=a^{-1}\operatorname{Tr}(\sigma\widehat D)
-a\varepsilon_0'(a)$, whose terms vanish. Divide by the positive
limiting energy for the ratio assertion.
\end{proof}
This is conditional survival of an already owned source, without WKB or a
postulated dust current. It prepares no state and selects no new dark owner.
Positive populations in a chosen spectral basis do not by themselves prove
positive $E_\infty$ when the inherited central allocation is retained.
A finite response gap is compatible with surviving population energy;
a zero-frequency linear pole is not a universal requirement for cold stress.

\begin{proposition}[The late energy is an actual history output]
\label{p3f:prop:v2203-cold-energy-history}
For fixed initial density, mass, carrier and endpoint allocation, put
$M_{\rm out}=W_\infty^\dagger H_mW_\infty$. Then
\begin{equation}
\label{p3f:eq:v2203-cold-energy-history}
 E_\infty=\operatorname{Tr}(\sigma_iM_{\rm out})+\varepsilon_\infty.
\end{equation}
For a compactly supported variation of an admitted history coordinate $q$,
let $D_q(t)=\partial_q\widehat H_0(t)$ and
$D_q^H=U(t,t_i)^\dagger D_qU(t,t_i)$. Its complete fixed-endpoint
history derivative is
\begin{equation}
\label{p3f:eq:v2203-cold-energy-history-derivative}
 \boxed{\delta E_\infty=i\int_{t_i}^{\infty}
   \operatorname{Tr}\bigl(\sigma_i[D_q^H(t),M_{\rm out}]\bigr)
    \delta q(t)\,\mathrm dt.}
\end{equation}
\end{proposition}
\begin{proof}
The energy formula is the cyclic form of
Eq.~\eqref{p3f:eq:v2203-local-source-cold-tail}. Unitary variation gives
$\delta W_\infty=-iW_\infty\int D_q^H\delta q\,\mathrm dt$;
differentiate $W_\infty^\dagger H_mW_\infty$ to obtain the commutator.
\end{proof}
Initial-state, mass, band, moving-frame or endpoint variations add their
actual inherited terms. No new amplitude factor occurs. A finite observation
epoch requires the exact mean
$E_*=\operatorname{Tr}(\sigma(t_*)\widehat H_0(a_*))$
and the complete scale work $p_s(t_*)$. Substituting $E_\infty$ for
$E_*$ requires a controlled late-tail error; $p_s/\rho_s\to0$ does not
prove exact $p_s(t_*)=0$.
For example the proof bounds the unmeasured remaining propagation by
$\|\widehat D\|\int_{t_*}^{\infty}a^{-1}\,\mathrm dt$, and the energy
error also retains $\|\widehat D\|/a_*$ and
$|\varepsilon_0(a_*)-\varepsilon_\infty|$. The actual history must
evaluate these terms. A separately conserved cold source has a ratio to
the fixed vacuum that decreases as $a^{-3}$; maintaining a constant
nonzero ratio instead requires actual exchange $Q=3H\rho_c$, with its
compensating work in the complete action.

\begin{proposition}[Frozen late mean and energy noise are different jets]
\label{p3f:prop:v2203-frozen-late-influence}
For the frozen generator $H_\infty=H_m+\varepsilon_\infty I$ and the
limiting interaction-picture state of
Theorem~\ref{p3f:thm:v2203-local-source-cold-tail}, the exact lapse-only
influence and its expansion are
\begin{align}
\label{p3f:eq:v2203-frozen-late-influence}
 \mathcal I_\infty&=\operatorname{Tr}
       (\sigma_\infty e^{-iH_\infty\Delta T/\hbar}),
 &\Delta T&=\int(N_+-N_-)\,\mathrm dt,\\
 \Gamma_{s,\infty}&=-E_\infty\Delta T
       +\frac{i}{2\hbar}\operatorname{Var}_{\sigma_\infty}(H_\infty)
              (\Delta T)^2+O((\Delta T)^3).\nonumber
\end{align}
Its frozen mean gives $\rho_s=E_\infty/V$, $p_s=0$. The nonnegative
variance is a separate fluctuation, not another mean density.
\end{proposition}
\begin{proof}
The two lapse histories commute because they multiply the same fixed
generator. Expand the logarithm of their characteristic function through
its first two cumulants. A central energy shifts only the first cumulant,
while cancelling from the variance and commutator response.
\end{proof}
This exact frozen identity asserts neither convergence of arbitrary
deformed early histories nor a complete inhomogeneous dust action.

The finite-CAR norm argument does not apply to the infinite Fock Hilbert
space of a finite-mode bosonic Gaussian interface. Its symplectic covariance
and actual physical work require their own domain and late-history
certificate; a constrained propagation generator is not automatically a
separately counted source energy. The conditional finite-matrix alternative
is Theorem~\ref{p4:thm:v2203-gaussian-cold-tail}; its positive limiting
physical source matrix and integrable history are additional actual
certificates. Euler characteristic $2$, quantum
covariance purity and the proposed abundance normalization are distinct
readouts. The remaining cosmological closure must identify the retained
source or derive a nonzero matched parent remainder, generate its state
and history, and evaluate both mean ports at the actual observation epoch.


\section{Forward common-action and scalar--metric source execution}
\label{paper3f:sec:v280-forward-action}
\subsection{Forward execution of the native local joint packet}
\label{p3f:sec:v251-forward-native-packet}

The starting domain is the intersection of the generated local Einstein,
Yang--Mills and odd-scalar real-form domains, the clock-hermitized uninserted
fermion branch, and an actually licensed finite quantization chart where it
is used.  Compactly supported interior variations are retained.  Canonical
non-null boundary data can be appended only on their separately proved
domain.  There is no assertion that all analytic and boundary domains have
one greatest member.

Write $\alpha_\Phi=(16\pi^2)^{-1}$, and retain the actual common connection,
finite trace, source member and scalar family.  In dimensionless local
action units the already constituted bosonic action is
\begin{equation}
\label{p3f:eq:v251-forward-native-action}
 S_{\rm loc}=\int\sqrt{-g}\left[
 \mathscr F(\Phi)R-\frac14Q_\rho(F_{\mu\nu},F^{\mu\nu})
 -\alpha_\Phi\operatorname{Tr}_F(D_\mu\Phi D^\mu\Phi)-V(\Phi)
 \right]d^4x,
\end{equation}
\[
 \mathscr F=\kappa_J-\frac{\alpha_\Phi}{6}\operatorname{Tr}_F\Phi^2,
 \qquad
 V=\alpha_\Phi\operatorname{Tr}_F
          (\Phi^4-2\Lambda_g^2\Phi^2).
\]
Its physical multiplier is $\hbar$ on the compatible spherical/source chart.
The fermion word is the actual ordered bilinear of the same Dirac datum,
with one positive $c_F$ multiplying both kinetic and zero-order blocks.
All inherited vertices and central terms are retained.  Pure-metric vacuum,
curvature-squared and exact nonlocal constituents are not silently included
in this two-derivative packet; each must be added through its actual owner
and claim domain if the claim consumes it.

\begin{proposition}[Explicit primitive scalar and fermion member packet]
\label{p3f:prop:v251-primitive-canonical-scalar-packet}
On the no-offset coordinate-linear Dirac-neutrino branch, the independently
generated primitive $Y_*=(I_G,I_G,I_G,I_G)$ has
\[
 \boxed{a_*=b_*=24,\quad
 H_c=\frac{\sqrt3}{\pi}H,\quad
 Y_{s,c}=\frac\pi{\sqrt3}I_G,\quad
 \lambda_* =\frac{\pi^2}{3}.}
\]
Consequently its native scalar--Einstein restriction is exactly
\begin{equation}
\label{p3f:eq:v251-primitive-canonical-local-action}
 \int\sqrt{-g}\left[
 \left(\kappa_J-\frac{H_c^\dagger H_c}{6}\right)R
 -(D_\mu H_c)^\dagger D^\mu H_c
 +2\Lambda_g^2H_c^\dagger H_c
 -\frac{\pi^2}{3}(H_c^\dagger H_c)^2
 \right]d^4x .
\end{equation}
The Yang--Mills term keeps its inherited $Q_\rho$ and the weak family
vertex remains $I_G$.  The canonically normalized local fermion member is
the original $D_{Y_*}$; no new Yukawa member or positive wavefunction
insertion is selected by this evaluation.
\end{proposition}
\begin{proof}
The generated trace invariants are
$a_Y=3\operatorname{tr}Y_u^\dagger Y_u+
3\operatorname{tr}Y_d^\dagger Y_d+
\operatorname{tr}Y_e^\dagger Y_e+
\operatorname{tr}Y_\nu^\dagger Y_\nu$, with the analogous squares for
$b_Y$.  Each $\operatorname{tr}I_G=3$, so both evaluate to
$9+9+3+3=24$.  Substitute in the already constituted canonical formulas
\[
 H_c=\sqrt{a_Y/(8\pi^2)}H,\quad
 Y_{s,c}=\sqrt{8\pi^2/a_Y}Y_s,\quad
 \lambda_Y=8\pi^2 b_Y/a_Y^2.
\]
Moreover $\operatorname{Tr}_F\Phi_Y^2=2a_YH^\dagger H$ implies
$\alpha_\Phi\operatorname{Tr}_F\Phi_Y^2=H_c^\dagger H_c$.
This proves every scalar/curvature coefficient at once.  The gauge action
and its representation are family blind, giving $I_G$ in the family
vertex.  The uninserted fermion jet is
$(Z_L,D,Z_R)=c_F(I,D_{Y_*},I)$, whose canonical normalization cancels
$c_F$ exactly.  No field equation or target mass is used.
\end{proof}

The primitive normalized finite internal chiral operator
$h_*=(\begin{smallmatrix}0&I_G\\I_G&0\end{smallmatrix})$ therefore has
the explicit tagged left response
\[
 \boxed{P_L(zI-h_*)^{-1}P_L=\frac z{z^2-1}I_G,\qquad
         \operatorname*{Res}_{z=\pm1}G_L=\frac12 I_G.}
\]
This is an internal/free comparison response, not a statement that the
compact geometric Dirac operator has a zero-spatial-momentum mode or that
these poles are named physical spacetime rest masses.

\begin{proposition}[Same-packet scalar and mixed jets]
\label{p3f:prop:v251-native-scalar-mixed-jets}
Let $\Phi=\sum_a\phi_aT_a$ in actual orthonormal real Hermitian coordinates,
$\operatorname{Tr}_F(T_aT_b)=\delta_{ab}$.  The local coefficient jets are
\[
 \mathscr F_{,a}=-\frac{\alpha_\Phi}{3}\phi_a,\qquad
 \mathscr F_{,ab}=-\frac{\alpha_\Phi}{3}\delta_{ab},
\]
\[
 V_{,a}=4\alpha_\Phi\operatorname{Tr}_F
         [T_a(\Phi^3-\Lambda_g^2\Phi)],
\]
\[
 V_{,ab}=4\alpha_\Phi\operatorname{Tr}_F
 \left[T_a\left(\Phi^2T_b+\Phi T_b\Phi+T_b\Phi^2
                         -\Lambda_g^2T_b\right)\right].
\]
For any actual commuting control coordinates $q^i$ in the same fermionic
word, the ordered bilinear coefficients are
\[
 [\bar\psi_a\psi_b]S=c_F\mathcal D_{ab},\quad
 [\bar\psi_a\psi_b]\partial_iS=c_F\mathcal D_{ab,i},\quad
 [\bar\psi_a\psi_b]\partial_i\partial_jS=c_F\mathcal D_{ab,ij}.
\]
Thus the mixed connection, scalar and metric vertices are derivatives of
the same member, not freely chosen companion maps.
\end{proposition}
\begin{proof}
Differentiate $\operatorname{Tr}\Phi^2$ and the cyclic polynomial trace.
For the quartic, $D(\Phi^3)[T_b]=\Phi^2T_b+\Phi T_b\Phi+T_b\Phi^2$;
cyclicity proves the displayed symmetric Hessian without assuming that
$\Phi$ and $T_b$ commute.  Ordered Grassmann coefficient extraction
commutes with ordinary coefficient differentiation, proving the last
three formulas.
\end{proof}

\paragraph{Forward finite quantization.}
On an actually finite CAR or finite joint-occupation chart, use the actual
embedding $J$, parent kinetic form $G_p$, skew connection $B_p$ and
energy $K_p$.  The complete forward map is
\[
 G=J^\dagger G_pJ,\qquad K=J^\dagger K_pJ,
\]
\[
 B=J^\dagger B_pJ+\tfrac12
       (J^\dagger G_p\dot J-\dot J^\dagger G_pJ),
 \quad S=G^{1/2},\quad W=S^{-1},\quad Z=JW,
\]
\[
 e=WKW,\qquad
 \gamma=WBW+\tfrac12(W\dot S-\dot SW),
\]
\begin{equation}
\label{p3f:eq:v251-forward-finite-generator}
 \boxed{\widehat H=d\Gamma(e)-i\hbar d\Gamma(\gamma)
       +\widehat V_{\rm inherited,can}+\varepsilon_0I.}
\end{equation}
Every inherited creation slot receives $Z^*$ and annihilation slot $Z$.
The positive-root jets are the unique Sylvester solutions
$SS_a+S_aS=G_a$ and
$SS_{ab}+S_{ab}S=G_{ab}-S_aS_b-S_bS_a$; inverse, embedding, constraint,
occupation and vertex slots are differentiated too.  For the native scalar
quadratic word $L_\Phi=\dot z^TM\dot z/2-\dot z^TCz+
z^TDz/2-z^TKz/2-U(z)$ this gives
\[
 H_\Phi=\tfrac12(p+Cz)^TM^{-1}(p+Cz)
                    +\tfrac12z^T(K-D)z+U(z).
\]
The actual ordering and occupation projection remain part of this map.
Finite mode rank alone is not a finite bosonic occupation space.  This
constructor computes $H$ after its embedding and retained chart are
identified; it does not provide an unspecified physical chart by itself.

\begin{theorem}[Forward causal source integration and Ward coherence]
\label{p3f:thm:v251-forward-causal-native-packet}
Restrict the same native action to its compatible homogeneous spherical
chart with $q=(A,B,\phi_a)$, $V_3=\Omega_zAB^2$, and its actual smooth
coordinate-only proper generator $h(q)$ in a flat canonical source frame.
Retain the actual initial density $\sigma_i$, every central term and the
lapse $N$ until variation.  With $H[N,q]=Nh(q)$ define
\[
 Z[q_+,q_-]=\operatorname{Tr}
  (U_{q_+}\sigma_iU_{q_-}^\dagger),\qquad
 \Gamma=-i\hbar\operatorname{Log}Z
\]
on the local normalized logarithm branch containing the diagonal.
The action
\[
 \mathcal I=\hbar S_{g\Phi}[N_+,q_+]
            -\hbar S_{g\Phi}[N_-,q_-]+\Gamma
\]
has the exact source covectors
\[
 F_N=-\langle h\rangle,\qquad F_i=-N\langle h_{,i}\rangle,
\]
and complete mixed source kernel
\[
 \Gamma_{\Delta a,c b}(t,s)
 =-\langle H_{,ab}(t)\rangle\delta(t-s)
 +\frac i\hbar\Theta(t-s)
       \langle[H_{,a}(t),H_{,b}(s)]\rangle.
\]
In particular $H_{,Ni}=h_{,i}$ is a required lapse/source contact.
The action-normalized ports are
\[
 \rho_r=\frac{\langle h\rangle}{\hbar V_3},\quad
 p_{z,r}=-\frac{A\langle h_{,A}\rangle}{\hbar V_3},\quad
 p_{\perp,r}=-\frac{B\langle h_{,B}\rangle}{2\hbar V_3},\quad
 J_a=\frac{\langle h_{,\phi_a}\rangle}{\hbar V_3}.
\]
These satisfy the exact work identity and its differentiated contact
identity.  For a compact control perturbation $(n,u^i)$ the latter is
\[
 \sum_i(\delta F_i\dot q^i+F_i\dot u^i)
       -n\dot F_N-N\partial_t\delta F_N=0.
\]
The same action gives the coupled curvature insertion
\begin{equation}
\label{p3f:eq:v251-forward-curvature-insertion}
 \boxed{R_*=-\frac{2\alpha_\Phi\Lambda_g^2}{\kappa_J}
             \operatorname{Tr}_F\Phi^2
       -\frac{\operatorname{Tr}(\sigma\mathsf T)}
                    {2\kappa_J\hbar V_3},\qquad
 \mathsf T=-h-Ah_{,A}-Bh_{,B}+\sum_a\phi_a h_{,\phi_a}.}
\end{equation}
No independent scalar charge, pressure coefficient, central-energy
subtraction, or fitted curvature cap occurs in these equations.
\end{theorem}
\begin{proof}
The finite propagator is unitary and the diagonal trace equals one, so
the normalized logarithm exists on a sufficiently small history
neighborhood.  Duhamel differentiation gives $F_a=-\langle H_{,a}\rangle$
and the contact-plus-retarded-commutator kernel, including every mixed
derivative.  Cyclicity and $i\hbar\dot\sigma=[Nh,\sigma]$ give
$d\langle h\rangle/dt=\sum_i\langle h_{,i}\rangle\dot q^i$.
Therefore $\sum_iF_i\dot q^i-N\dot F_N=0$.  Dividing by the actual
volume and differentiating that volume gives
\[
 D_\tau\rho_r+H_A(\rho_r+p_{z,r})
 +2H_B(\rho_r+p_{\perp,r})=\sum_aJ_aD_\tau\phi_a.
\]
The mixed kernel is its actual differentiated source response.  Trace
the native metric equation and use the same scalar Euler equation:
\[
 2\kappa_JR=4V-\operatorname{Tr}(\Phi V_{,\Phi})
                 -T_r-\operatorname{Tr}(\Phi J),
 \qquad
 4V-\operatorname{Tr}(\Phi V_{,\Phi})
       =-4\alpha_\Phi\Lambda_g^2\operatorname{Tr}\Phi^2.
\]
Substitute the four finite-source ports.  Their combination is exactly
$\operatorname{Tr}(\sigma\mathsf T)/(\hbar V_3)$, proving the boxed
curvature insertion.  The native common-action Cauchy equivalence retains
the lapse constraint and source state, so no independently prescribed
source balance is needed on this chart.
\end{proof}

\paragraph{Actual scope of the forward result.}
The primitive branch now has explicit inherited scalar/Yukawa coefficients,
finite internal response, source quantization map and coupled metric/scalar
source insertions.  A genuinely consumed family-dependent transport in the
complete $h$, or an appropriate noninvariant actual source, is necessary to
obtain a noncentral family response from otherwise family-blind data.  This
is an execution of one native branch, not a proof that all rooted transport
operators are physically absent.  Charged hierarchy and neutral rank targets
must be compared to its actual complete response; they are not inserted
into $h$ to define it.  The local functional does not identify the global
conormal cap, restored response, a $1/96$ projection, a physical vacuum or
global/continuum dynamics outside its stated domain.


\subsection{Forward scalar--metric and finite-source integration}
\label{p3f:sec:v251-forward-native-gravity-source}

Throughout this subsection put $L=\Lambda_g^2>0$, let $\kappa>0$ be
the actual raw joint Einstein coefficient and retain $\hbar$ in finite
source ports. The native ordinary-trace scalar action has
\[
 \mathscr F=\kappa-\frac\alpha6\operatorname{Tr}\Phi^2,
 \qquad V=12\kappa L+\alpha\operatorname{Tr}(\Phi^4-2L\Phi^2),
 \qquad\alpha=(16\pi^2)^{-1}.
\]
The bare constant has been included once. A finite remainder source below
does not include a second copy of that same constant. The domain is the
common local scalar--metric and compatible coordinate-only finite-source
domain of the native influence/Cauchy theorems. No scalar cap or restored
$f(R)$ action is substituted for this packet. The actual raw coupling
$G_{J,\mathrm{raw}}=G_{\mathrm{scal}}/96$ is retained as a cross-owner
dictionary, rather than changed by a sectorwise normalization.

\begin{proposition}[The exact finite-source operator for the bare trace budget]
\label{p3f:prop:v251-forward-source-trace-operator}
Let the complete retained remainder Hamiltonian be
$h=h(A,B,\phi_1,\ldots,\phi_d)$ on the actual flat finite source frame,
with state $\sigma$, volume $V_3=\Omega_zAB^2$, and native source ports.
Define the geometric scaling derivation and Hermitian operator
\[
 \mathscr D_0=A\partial_A+B\partial_B-\sum_a\phi_a\partial_{\phi_a},
 \qquad \mathsf K_h=-h-\mathscr D_0h.
\]
Then the complete improved source trace is
\[
 \mathcal I_h=T_h+\sum_a\phi_aJ_a
          =\frac{\operatorname{Tr}(\sigma\mathsf K_h)}{\hbar V_3}.
\]
On the native metric and scalar mean equations with no omitted remainder,
\[
 \boxed{\frac RL=12\left(1+\frac{\mathscr F}{\kappa}\right)
       -\frac{\operatorname{Tr}(\sigma\mathsf K_h)}{2\kappa L\hbar V_3}.}
\]
If the state eigenvalues are $p_1\ge\cdots\ge p_n$ and the eigenvalues
of $\mathsf K_h$ are $k_1\ge\cdots\ge k_n$, its expectation lies sharply
between $\sum_jp_jk_{n+1-j}$ and $\sum_jp_jk_j$. Consequently any point
with $\mathscr F\ge f\kappa>0$ and $R\le\epsilon L$ requires
\[
 \sum_jp_jk_j\ge
 2\kappa L\hbar V_3\,[12(1+f)-\epsilon].
\]
This is a necessary pointwise feasibility test, not a proof that its
extremizing state occurs or obeys all other field equations.
\end{proposition}
\begin{proof}
The actual native ports are
$\rho_h=\operatorname{Tr}(\sigma h)/(\hbar V_3)$,
$p_{z,h}=-A\operatorname{Tr}(\sigma h_{,A})/(\hbar V_3)$,
$p_{\perp,h}=-B\operatorname{Tr}(\sigma h_{,B})/(2\hbar V_3)$ and
$J_a=\operatorname{Tr}(\sigma h_{,a})/(\hbar V_3)$.
Insert them into $T_h=-\rho_h+p_{z,h}+2p_{\perp,h}$ to get the first
identity. Substitution into the native bare improved-trace identity gives
the boxed equation. The expectation bounds follow from the finite
Hermitian trace rearrangement inequality, and the final inequality follows
by using the upper bound on the required expectation.
\end{proof}

\begin{proposition}[Conformal source homogeneity and a positive-excitation sign test]
\label{p3f:prop:v251-forward-source-trace-sign}
On one actual fixed finite source chart the following implications hold.
\begin{enumerate}
\item If
$h(\lambda A,\lambda B,\lambda^{-1}\phi)=\lambda^{-1}h(A,B,\phi)$
in its fixed canonical frame, then $\mathsf K_h=0$ for every state.
Such a source cannot change the bare improved trace.
\item Let a free positive one-particle source have
$\omega=(K+M^2)^{1/2}>0$, where $K,M^2\succeq0$,
$\mathscr D_0K=-2K$ and $\mathscr D_0M^2=0$ on the compared chart.
Then
\[
 \boxed{\mathsf K_\omega=-\omega-\mathscr D_0\omega
       =-2\int_0^\infty e^{-t\omega}M^2e^{-t\omega}\,dt\preceq0.}
\]
The same sign holds for its positive free excitation second quantization
on an admitted fixed finite fermionic or bounded bosonic occupation space.
It does not require $[K,M^2]=0$. Such excitations alone cannot supply the
positive trace cancellation necessary for low bare curvature.
\end{enumerate}
\end{proposition}
\begin{proof}
Differentiate the homogeneity identity at $\lambda=1$.
For the second statement differentiating $\omega^2=K+M^2$ gives the
Sylvester equation
$\omega(\mathscr D_0\omega)+(\mathscr D_0\omega)\omega=-2K$.
Therefore
$\omega\mathsf K_\omega+\mathsf K_\omega\omega=-2M^2$.
Positive $\omega$ makes its Lyapunov inverse the displayed convergent
positive integral. Second quantization on free excitation sectors
preserves the operator order and commutes with the fixed-chart derivative.
The trace budget proves the last assertion.
\end{proof}

\begin{remark}[Absolute central energy is indispensable]
Write the actual source as $h=c(q)I+h_{\mathrm{exc}}$ only when that
split has been derived, retaining
$\mathsf K_c=-c-\mathscr D_0c$. It is invisible to fixed-source energy gaps
but not to this trace. An actually generated constant local vacuum density
$c=\hbar V_3\delta V$, with $\delta V$ independent of the compared
fields and metric coordinates, gives $\mathcal I_c=-4\delta V$.
For a general coordinate-dependent density the exact result is instead
$\mathcal I_c=-4\delta V-\mathscr D_0\delta V$. At $\Phi=0$,
low curvature would require
$\delta V=-12\kappa L+\kappa R/2$; this is a required cancellation,
not permission to insert it.
For an admitted finite filled fermionic ground with symmetric one-particle
levels $\pm\omega$, the ground center is instead $c_0=-\operatorname{Tr}\omega$.
Its scaling trace is $-\operatorname{Tr}\mathsf K_\omega\ge0$.
It can therefore contribute in the needed sign, whereas the positive
excitation contribution has the opposite sign. Its magnitude, physical
state, volume scaling, pressure and history remain to be evaluated from
the actual complete Hamiltonian. A pointwise trace cancellation alone
does not make the state a constant-curvature vacuum. No vacuum energy is
silently removed by a change of representative.
\end{remark}

\subsection{Executed primitive-member scalar--metric branch}
\label{p3f:sec:v251-forward-primitive-scalar-metric}
For the actual rooted primitive member $Y_*=(I,I,I,I)$, the native
canonical doublet has $a_Y=b_Y=24$,
$H_c=\sqrt3H/\pi$ and $\lambda=\pi^2/3$. With $x=H_c^\dagger H_c$,
the actual degree-at-most-four scalar--metric packet is
\[
 S=\int\sqrt{-g}\left[
 (\kappa-x/6)R-(D_\mu H_c)^\dagger D^\mu H_c
 -V_0+2Lx-\lambda x^2\right],
\quad \kappa=\frac{L}{8\pi^2},\quad V_0=12\kappa L.
\]
The last coefficient relation uses the printed common raw heat scales.
On a source-free parallel constant comparison,
\[
 (\kappa-x/6)R=2(V_0-2Lx+\lambda x^2),\qquad
 [R/6-2L+2\lambda x]H_c=0.
\]
The zero branch is $x=0$, $R=24L$. The formal nonzero branch is
\[
 x=L/\pi^2,\qquad R=8L,\qquad
 \mathscr F=-L/(24\pi^2)<0,
\]
so it is excluded from the positive-graviton chart. This is an explicit
specialization of the existing all-rank native obstruction, rather than a
new controlled high-curvature vacuum prediction. A common positive
normalization changes none of these dimensionless conclusions.

\begin{theorem}[The complete Einstein scalar metric of the primitive packet]
\label{p3f:thm:v251-forward-higgs-hyperbolic-target}
Write the canonical complex doublet as four real fields $q^a$, so
$x=\tfrac12\sum_a(q^a)^2$ and its kinetic term is
$-\tfrac12\delta_{ab}Dq^aDq^b$. On $\mathscr F>0$ set
$g_E=(\mathscr F/\kappa)g$. The complete transformed scalar metric is
\[
 G_{ab}=\frac\kappa{\mathscr F}\delta_{ab}
           +\frac{3\kappa}{\mathscr F^2}\mathscr F_{,a}\mathscr F_{,b},
 \qquad\mathscr F=\kappa-|q|^2/12.
\]
It is the complete hyperbolic metric of radius $\beta=\sqrt{12\kappa}$.
In radial coordinates $r=|q|=\beta\tanh(\varphi/\beta)$,
\[
 ds_G^2=d\varphi^2+\beta^2\sinh^2(\varphi/\beta)d\Omega_3^2.
\]
The bare Einstein potential is exactly
\[
 \boxed{W(\varphi)=12\kappa L\cosh^2(\varphi/\beta)
              +36\lambda\kappa^2\sinh^4(\varphi/\beta),}
\]
and its covariant Hessian obeys $\nabla_G^2W\succeq2L G$.
For the actual primitive coefficients its entire dimensionless shape is
\[
 \frac{W}{12\kappa L}=\cosh^2u+\tfrac18\sinh^4u,
 \quad u=\varphi/\sqrt{12\kappa}.
\]
\end{theorem}
\begin{proof}
The Weyl transformation of $\sqrt{-g}\mathscr F R$ contributes
$-3\kappa(\partial\log\mathscr F)^2/2$ to the scalar kinetic form.
Combining this with the transformed original kinetic term gives $G$,
including the metric--scalar mixing contribution. Its radial eigenvalue
is $\kappa^2/\mathscr F^2$ and its tangential eigenvalue is
$\kappa/\mathscr F$. Substitution of $r=\beta\tanh(\varphi/\beta)$
gives the hyperbolic metric and makes the boundary $\mathscr F=0$ an
infinite metric distance. The potential is
$W=\kappa^2(V_0-Lr^2+\lambda r^4/4)/\mathscr F^2$.
The stated bare $V_0$ gives the displayed formula.
Put $u=\varphi/\beta$. The radial Hessian eigenvalue is
\[
 W''=2L(\cosh^2u+\sinh^2u)
          +36\lambda\kappa\sinh^2u\cosh^2u
          +12\lambda\kappa\sinh^4u\ge2L.
\]
The tangential eigenvalue is
$W'\coth u/\beta=2L\cosh^2u+
12\lambda\kappa\sinh^2u\cosh^2u\ge2L$, with its regular origin
limit $2L$. These are the full covariant Hessian eigenvalues.
\end{proof}

\begin{corollary}[Source-force bound and positive-graviton margin]
\label{p3f:cor:v257-native-force-margin}
Import precisely the exact primitive scalar chart of
Theorem~\ref{p3f:thm:v251-forward-higgs-hyperbolic-target}, with
$L>0$, $\kappa=L/(8\pi^2)$, $\beta=\sqrt{12\kappa}$,
$\mathscr F=\kappa-|q|^2/12>0$ and $\rho=d_G(q,0)$.
Its bare potential satisfies
\[
 W(q)-W(0)\ge Lr^2,\qquad |\nabla_GW(q)|\ge2Lr.
\]
If an actually licensed extra covector $J$ on this same scalar chart
satisfies $dW(q_*)+J(q_*)=0$ and $|J(q_*)|_{G^{-1}}\le\varepsilon$,
then
\[
 \rho_*\le\frac{\varepsilon}{2L},\qquad
 \mathscr F(q_*)\ge\kappa\operatorname{sech}^2
 \!\left(\frac{\varepsilon}{2L\sqrt{12\kappa}}\right)>0.
\]
For $|dO|_{G^{-1}}\le K$ along the radial geodesic,
$|O(q_*)-O(0)|\le K\varepsilon/(2L)$.
If $J=dV$ and globally $\nabla_G^2V\succeq-\delta G$ with $\delta<2L$,
then $W+V$ has a unique critical point and minimum, with
\[
 \rho_*\le\frac{|dV(0)|_{G^{-1}}}{2L-\delta}.
\]
\end{corollary}
\begin{proof}
Along a unit-speed radial geodesic, the source's Hessian bound gives
$(W\circ\gamma)''\ge2L$ and $(W\circ\gamma)'(0)=0$. Integrate twice.
At stationarity the gradient norm equals the source norm. Use
$\mathscr F=\kappa\operatorname{sech}^2(\rho/\beta)$ and integrate $dO$.
For the final claim, $\mu=2L-\delta>0$ is a global convexity constant.
Along geodesics from zero,
$(W+V)(q)\ge(W+V)(0)-|dV(0)|\rho+\mu \rho^2/2$.
Completeness and finite dimension give compact bounded closed balls,
hence existence of a minimum. Strict geodesic convexity gives uniqueness;
the integrated first derivative at the minimum gives the distance bound.
\end{proof}
This is an exact polynomial-branch result. Source bounds must be proved
in the inherited metric. Corrections changing that metric or conformal
coefficient require a new calculation. The bare equilibrium $R=24L$
does not lie in a controlled low-curvature heat-expansion regime.

\begin{remark}[An equation-level stationary witness on the exact bare restriction]
\label{p3f:rem:v257-bare-equation-witness}
For the exact source-free primitive polynomial restriction, set $q=0$,
gauge curvature zero and retained matter fields zero. With $W_0=12\kappa L$,
the Einstein equation is $\kappa G_{\mu\nu}=-(W_0/2)g_{\mu\nu}$.
The closed de Sitter chart
\[
 \ell=(2L)^{-1/2},\qquad a(t)=\ell\cosh(t/\ell),\qquad
 H(t)=\ell^{-1}\tanh(t/\ell)
\]
has $H^2+a^{-2}=2L=W_0/(6\kappa)$ and $R=24L$; the scalar equation
holds since $dW(0)=0$. A finite time interval with compactly supported
variations gives a nonempty Euler--Lagrange zero locus of this exact
restriction. Membership in the complete physical stationary family also
requires the admitted configuration role, retained joint GP2 conditions
and actual endpoint rule. The equation-level witness alone does not
establish those additional conditions, an actual history, or control of
the original heat expansion at $R=24L$.
\end{remark}


\begin{remark}[No additional scalaron has been inserted]
Here $\mathscr F$ is the composite scalar coefficient of this same
doublet, not an independent Legendre coordinate. The field $\varphi$ is
the canonical Higgs radial field after its actual metric mixing. Introducing
an independent restored-$f(R)$ scalaron requires a separately established
complete action overlap. Gauge angular directions above remain subject to
their actual gauge constraints; they are not four automatically independent
physical scalar particles.
\end{remark}

\begin{theorem}[A global source-free radial future for the exact polynomial packet]
\label{p3f:thm:v251-forward-bare-radial-global-future}
Regard the supplied degree-at-most-four scalar--metric action as an exact
classical polynomial functional on its positive-graviton chart. Restrict
to its gauge-compatible source-free real doublet
$H_c=(0,h/\sqrt2)$, with zero gauge curvature/current and no omitted
source or higher term. Use closed Einstein-frame FLRW proper time.
Let $W$ be the potential above extended evenly to $\varphi\in\mathbb R$,
let $W_0=12\kappa L$, and impose its Friedmann constraint at $t_i$.
If
\[
 a_i>\sqrt{6\kappa/W_0}=1/\sqrt{2L},\qquad H_i>0,
 \qquad E_i=\tfrac12\dot\varphi_i^2+W(\varphi_i)<\infty,
\]
then its solution exists for every $t\ge t_i$ and remains expanding.
Moreover
\[
 \mathscr F(t)\ge\kappa W_0/E_i>0,\quad
 \varphi(t)\to0,\quad\dot\varphi(t)\to0,\quad
 H(t)\to\sqrt{2L},\quad a(t)\to\infty.
\]
The corresponding Jordan proper-time interval is also infinite, since
$dt_J=\sqrt{\kappa/\mathscr F}\,dt_E$ and the conformal factor remains
bounded above and below by positive constants on this solution.
This is a global theorem for the stated exact polynomial restriction,
not a controlled low-curvature theorem for the original heat functional.
\end{theorem}
\begin{proof}
The actual Einstein-frame equations are
\[
 H^2+a^{-2}=E/(6\kappa),\qquad
 \dot H=a^{-2}-\dot\varphi^2/(4\kappa),\qquad
 \ddot\varphi+3H\dot\varphi+W'(\varphi)=0.
\]
Thus $\dot E=-3H\dot\varphi^2$. On any expanding interval,
$a\ge a_i$, $E\le E_i$ and $E\ge W_0$. The constraint implies
\[
 H^2\ge W_0/(6\kappa)-a_i^{-2}=:H_{\min}^2>0.
\]
Continuity excludes a first zero of $H$, so expansion persists. Also
$H\le\sqrt{E_i/(6\kappa)}=:H_{\max}$. Coercivity and
$W\ge W_0\cosh^2(\varphi/\beta)$ bound $\varphi$, its velocity and
$\mathscr F=\kappa\operatorname{sech}^2(\varphi/\beta)$ as stated.
Use $b=1/a$ instead of $a$. The variables $(b,H,\varphi,p)$, with
$p=\dot\varphi$, remain in a compact set of the smooth system
\[
 \dot b=-Hb,\quad \dot H=b^2-p^2/(4\kappa),\quad
 \dot\varphi=p,\quad\dot p=-3Hp-W'(\varphi).
\]
It is smooth also at $b=0$, so no finite-time endpoint exists. This
proves global future continuation and $b\le b_ie^{-H_{\min}(t-t_i)}$.
Energy dissipation gives $\int_{t_i}^\infty p^2dt<\infty$.
The compact bounds give bounded $\dot p$ and $\ddot p$.
If $p$ failed to tend to zero, bounded $\dot p$ would produce disjoint
intervals of fixed positive width on which $p^2$ is bounded below,
contradicting its finite integral. If $\dot p$ failed to tend to zero,
bounded $\ddot p$ would produce fixed-width intervals on which $\dot p$
has one sign and fixed magnitude, contradicting $p\to0$. Thus both
$p\to0$ and $\dot p\to0$. The equation gives $W'(\varphi)\to0$.
Since the strongly convex potential has the unique critical point zero,
$\varphi\to0$. The constraint now gives $H\to\sqrt{W_0/(6\kappa)}$.
\end{proof}

\begin{remark}[What this forward construction does and does not close]
The global trajectory approaches $R_E=24L$, exactly the high-curvature
bare constant solution already identified by the native vacuum calculation.
The original local heat expansion is not controlled there. The theorem
therefore supplies a real coercive continuation mechanism on an explicitly
declared classical truncation; it does not turn that truncation into a
physical low-curvature cosmology. The barrier is at $\mathscr F=0$, where
the actual Einstein potential diverges, and is unrelated to the finite-height
terminal susceptibility endpoint of a different restored-$f(R)$ action.
Populated finite sources, actual central vacuum energy, nonzero gauge
fields, normal/boundary terms and higher metric channels must be co-consumed
before this theorem can be extended to those claims.
\end{remark}


\subsection{Regular curvature polarization and the actual Jordan boundary momentum}
\label{p3f:sec:v260-curvature-canonical-transport}
This subsection works on the \emph{declared} regular local metric-$f_c(R)$
parent of Theorem~\ref{p3f:thm:v223-cap-parent} and its
first-order Jordan presentation, with fixed $h_{ij},\Phi$ endpoint data.
It is not a proof that the full native common action selects that cap.
Write $P_E$ for the Einstein coefficient in $f_c$, reserving $P_c$
for the curvature polarization coordinate below.
\begin{corollary}[Local curvature canonical exchange with endpoint and Hamiltonian]
\label{p3f:cor:v260-canonical-exchange}
\label{p3f:cor:v229-canonical-exchange}
Let $\Phi=f_c'(r)$ on an interval with $\chi(r)=f_c''(r)\ne0$,
and let $\eta\ne0$ be a constant conversion factor. At any regular
canonical chart with scalar momentum $p_\Phi$, define
\[
 P_c=\eta r,\qquad Q=-\frac{\chi(r)}{\eta}p_\Phi.
\]
Then $r=P_c/\eta$, $p_\Phi=-\eta Q/\chi(P_c/\eta)$, and
\begin{equation}
\label{p3f:eq:v260-canonical-exchange}
 \boxed{p_\Phi\,d\Phi=P_c\,dQ-d(P_cQ),\qquad
  dp_\Phi\wedge d\Phi=dP_c\wedge dQ.}
\end{equation}
For closed-FLRW gravity with $B_g=M^2V_\gamma$, $k=1$, and
$U_J(\Phi)=r\Phi-f_c(r)$, the exact pulled-back gravitational
Hamiltonian is
\begin{equation}
\label{p3f:eq:v260-exchanged-hamiltonian}
\begin{aligned}
\mathcal H_g'={}&
\frac{\eta p_aQ}{3B_ga^2\chi(P_c/\eta)}
+\frac{\eta^2\Phi(P_c)Q^2}{3B_ga^3\chi(P_c/\eta)^2}
-3B_gka\Phi(P_c)\\
&+\frac{B_ga^3}{2}\left[
\frac{P_c}{\eta}\Phi(P_c)-f_c(P_c/\eta)\right],\qquad
\Phi(P_c)=f_c'(P_c/\eta).
\end{aligned}
\end{equation}
The full source Hamiltonian and canonical/source momenta must be pulled
back from the \emph{same} action. The phase action changes by
$-[P_cQ]_{t_i}^{t_f}$, so changing the fixed endpoint polarization
requires the compensating Legendre boundary term. No unchanged $P_c^2$
energy or boundary flux follows from the exchange.
\end{corollary}
\begin{proof}
Since $d\Phi=\chi\,dr=(\chi/\eta)\,dP_c$,
$p_\Phi d\Phi=-QdP_c=P_c dQ-d(P_cQ)$; exterior differentiation
proves the symplectic statement. The inverse exists by the
inverse-function theorem when $\eta\chi\ne0$. Substitute
$p_\Phi=-\eta Q/\chi$ and $\Phi=f_c'(P_c/\eta)$ into the
existing Jordan Hamiltonian in
Eq.~\eqref{p3f:eq:v169-flrw-hamiltonian}; the displayed mixed,
quadratic, curvature, and potential terms follow separately.
The phase-action endpoint claim is the integrated one-form identity.
No transformation of an omitted velocity-dependent source is implied.
\end{proof}

\begin{proposition}[Jordan lapse constraint and nonuniqueness of boundary flux at fixed curvature]
\label{p3f:prop:v260-jordan-flux}
\label{p3f:prop:v229-jordan-flux}
On the same closed-FLRW Jordan action with a velocity-independent
homogeneous source, retain the lapse until variation. The gravity
momenta and Hamiltonian constraint are
\begin{align}
 p_a&=-6B_gFa\dot a/N-3B_ga^2\dot F/N,
 &p_F&=-3B_ga^2\dot a/N=-B_ga^3K,\\
 3F(H_\Sigma^2+a^{-2})&=\frac12U_J(F)-3H_\Sigma D_tF
                              +\frac{\rho_s}{M^2},
 &K&=3H_\Sigma .
\end{align}
On the sourceless complete cap $f_c(r)=P_Er+r^2B(r)-\sigma^2/B(r)^2$,
with $B(r)=B_U+2\Delta/(1+\sqrt{1+(r/R_s)^2})$,
$\Delta=B_I-B_U>0$, $\sigma=B_IR_s$,
$B_U,P_E,R_s>0$, there exists a positive constant-curvature root
$r_*>0$ with $r_*f_c'(r_*)-2f_c(r_*)=0$. If the regularity
condition of Proposition~\ref{p3f:prop:v223-cap-legendre}
holds, $F_*=f_c'(r_*)>0$ and the vacuum equations admit the closed
de Sitter chart
\[
  H_*^2=r_*/12,\qquad a(t)=H_*^{-1}\cosh(H_*t),\qquad
  K(t)=3H_*\tanh(H_*t),\qquad F(t)=F_* .
\]
At any $t_0>0$, the two slices $\pm t_0$ have identical induced metric
and identical $r_*,F_*$, yet opposite nonzero $p_F$.
Thus neither the scalar curvature nor its Legendre field and the
induced metric alone determine the Jordan scalar boundary momentum.
In particular the fixed-source complete conormal norm of
Theorem~\ref{p3e:thm:v229-conormal-cap} is not identified with
$p_F$, $p_F^2$, or the metric canonical flux by this calculation.
\end{proposition}
\begin{proof}
Differentiating the lapse-unfixed Jordan minisuperspace Lagrangian
\eqref{p3f:eq:v169-flrw-lagrangian} gives the momenta. Lapse variation
gives the displayed constraint; these are the \emph{same-action}
scalar and gravitational boundary data. For the cap,
$f_c(0)=-R_s^2$, hence $[rf_c'-2f_c]_{r=0}=2R_s^2>0$.
As $r\to+\infty$,
\[
 B(r)=B_U+\frac{2\Delta R_s}{r}+O(r^{-2}),\qquad
 f_c(r)=B_Ur^2+(P_E+2\Delta R_s)r+O(1),
\]
so $rf_c'(r)-2f_c(r)=-(P_E+2\Delta R_s)r+O(1)\to-\infty$.
Continuity yields $r_*>0$. For constant $F_*$, the vacuum
metric equation becomes
$(F_*r_*/4-f_c(r_*)/2)g_{\mu\nu}=0$ on an Einstein metric;
the root makes it zero. Closed de Sitter has
$R_{\mu\nu}=3H_*^2g_{\mu\nu}$ and
$H_\Sigma^2+a^{-2}=H_*^2$.
The Jordan constraint holds because
$U_J(F_*)=r_*F_*-f_c(r_*)=r_*F_*/2$.
Finally $a(-t_0)=a(t_0)$, whereas $K(-t_0)=-K(t_0)$;
the momentum formula proves the counterexample.
The usual contracted-Bianchi identity plus the scalar/source Euler
identity propagates the lapse constraint within this consistent
sourceless reduction; no arbitrary extra curvature cap is imposed.
\end{proof}
\begin{corollary}[Conormal norm and absolute-amplitude firewall]
\label{p3f:cor:v260-conormal-jordan-distinction}
For the separately constituted positive two-channel conormal problem
of Theorem~\ref{p3e:thm:v229-conormal-cap}, the constrained minimum is
\[
 \boxed{\mathcal E_{\partial}(r;b)=
 \frac{b^2\gamma_N^2d_0r^2}{2(1+G(r))},\quad
 G(r)=\sqrt{1+(r/R_s)^2}.}
\]
It equals the target $2\Delta r^2/(1+G(r))$ precisely when
$b^2\gamma_N^2d_0=4\Delta$. This fixes the required magnitude
\emph{within the matched source class}; it does not determine the
physical incoming amplitude $b$, actual source/operator identity,
or the occurrence of that cap in the common action.
\end{corollary}
\begin{proof}
For $p_0+p_1=j_r$ and inverse blocks
$\bar K^{-1},(G\bar K)^{-1}$, variation gives
$p_0=j_r/(1+G)$ and $p_1=Gj_r/(1+G)$.
The minimum is
$\langle j_r,\bar K^{-1}j_r\rangle/[2(1+G)]$.
Putting $j_r=b\gamma_NrJ_1$ and transporting the frozen trace
as in Theorem~\ref{p3e:thm:v229-conormal-cap} gives the boxed norm.
Comparison with the target fixes the displayed necessary and
sufficient scalar equation. The preceding proposition proves that
Jordan canonical momentum is not fixed by this curvature-only value.
\end{proof}


\begin{theorem}[Common finite action jets and compressed observables without subgroup ceiling]
\label{p3f:thm:v275-joint-positive-action-jets}
Work on the existing compact $SU(2)$ comparison factor, with
Haar probability $\sigma$ and a \emph{supplied} tetrahedral
core $H$. Fix a finite-dimensional field-parameter domain
$\theta$, a finite jet order $s$, a parameter point $\theta_0$,
and finitely many continuous action integrands
$L_j(q,\theta)$ whose parameter derivatives through order $s$
are jointly continuous on $G\times U$ for a neighborhood $U$
of $\theta_0$. There is a single positive
finite rule on words in the two icosahedral completions of $H$
that reproduces all prescribed integrated action jets exactly:
\[
 \left.\partial_\theta^\alpha\int L_j(q,\theta)
                       \,d\sigma(q)\right|_{\theta_0}
 =\left.\partial_\theta^\alpha
             \sum_iw_iL_j(q_i,\theta)\right|_{\theta_0}
 \quad (|\alpha|\le s),
\]
provided differentiation under the integral is licensed by the
stated compact-domain regularity.

If $L_j(q,\theta)=\sum_{a=1}^{d_j}c_{ja}(\theta)f_{ja}(q)$
has a finite separated expansion, the same rule reproduces the
entire finite action for all $\theta$ in its domain, not only
the specified jets. Any finite list of identities involving
these exact integrated coefficient functions is preserved.

Additionally, let $a_1,\dots,a_t$ be polynomial scalar
observable symbols with degrees at most $d_a$, and consider
the compressed multiplication operators
$C_{a_j}^{(m)}=P_mM_{a_j}|_{\mathcal P_m}$. If a single positive
rule is exact on $\mathcal P_L$ with
\[
 \boxed{L\ge\max\{2m+d_a,\ r+2,\ d_{\rm action}\}},
\]
where $r$ is the prescribed endpoint-POVM moment degree and
$d_{\rm action}$ bounds the polynomial action-jet integrands,
then that same rule reconstructs the exact matrices
\[
 \boxed{C_{a_j}^{(m)}
       =(S_m)^*\operatorname{diag}(a_j(q_i))S_m}
\]
and their finite-dimensional spectra and compressed-matrix
identities, simultaneously with the action jets and the
readout moments. The evaluation $S_m$ is isometric and the
node count is at least $\dim\mathcal P_m$.
No result here asserts a physical occurrence, all-resolution
exactness by one fixed finite set, local gauge Ward identities
beyond the tested span, or full interacting continuum dynamics.
\end{theorem}
\begin{proof}
The finitely many parameter jets form a finite-dimensional
real span of continuous scalar functions; apply
Theorem~\ref{p1b:thm:v275-dense-twin-word-cubature}
(or Corollary~\ref{p2e:cor:v275-joint-observable-exactification})
to that span. Uniform regularity on the compact factor permits
interchange of integration and the stated derivatives, and
exactness applies to every jet coefficient separately. A
finite separated expansion places every $L_j(\cdot,\theta)$
in one fixed finite-dimensional span, proving full-parameter
exactness.

For the polynomial assertion, $\overline f a_jg$ has degree
at most $2m+d_a$ whenever $f,g\in\mathcal P_m$.
By cubature exactness,
\[
 \langle f,C_{a_j}^{(m)}g\rangle_{L^2(\sigma)}
  =\sum_iw_i\overline{f(q_i)}a_j(q_i)g(q_i)
  =\langle S_mf,\operatorname{diag}(a_j(q_i))S_mg\rangle.
\]
The degree-$2m$ subcase gives $S_m^*S_m=I$; hence the
operator matrix identity holds exactly. The terms $r+2$ and
$d_{\rm action}$ guarantee the other specified equalities.
Injectivity of $S_m$ yields the node lower bound. Notice that
products of compressed operators are not generally equal to
compressions of products without additional off-band control;
only actual compressed-matrix identities are exported.
\end{proof}


\section{Complete metric work, source stress, and homogeneous closure}
\label{paper3f:sec:organized-metric}
\subsection{Completed closed-FLRW density and pressure variation}
\label{p3f:sec:v2205-closed-two-port-stress}

Use $\hbar=c=1$ and the already completed homogeneous coherent-scalar
domain. The controlling action constituents are
Eqs.~\eqref{p3f:eq:v2203-completed-curvature} and
\eqref{p3f:eq:v2203-matched-lapse}. No spatial wall or independent source
is appended. Let $V_\gamma=2\pi^2$, $V=V_\gamma a^3$, $C=C_7$,
$\rho_\Lambda=M^2C/2$, and
\[
 \mathcal A=\frac{M^2}{2}\Omega,\qquad
 \Omega=F-\frac{|h|^2}{6M^2}>0,\qquad
 W=\frac{M^2P(F)}2+V_h(h).
\]
Here $h$ is the retained native material scalar, not a second scalaron;
$P$ is the same-action Legendre potential. The potential above has no
explicit scale dependence at fixed independent $F,h$ on this homogeneous
domain. Any further gradient, mixed, moving-support or metric-dependent
source term retains its complete functional work in $S_s$. The source census
$S_s$ excludes the scalar energy and nonminimal term already displayed here;
every retained constituent is allocated exactly once.

\begin{theorem}[Lapse and spatial work of the completed matched action]
\label{p3f:thm:v2205-closed-matched-work}
With canonical endpoints fixed, the inherited action has the lapse-preserving
representative
\begin{equation}
\label{p3f:eq:v2205-lapse-preserving-action}
 S=V_\gamma\int dt\left[
 -\frac{3M^2\Omega a\dot a^2}{N}
 -\frac{3M^2a^2\dot a\dot\Omega}{N}
 +3M^2Na\Omega
 +\frac{a^3|\mathcal D_t h|^2}{2N}-Na^3W\right]+S_s.
\end{equation}
Dots here denote coordinate derivatives, and $\mathcal D_t h$ is the
coordinate covariant derivative with its retained temporal connection.
Define the complete source ports by
\[
 \frac{\delta S_s}{\delta N}=-V_\gamma a^3\rho_s,\qquad
 \frac{\delta S_s}{\delta a}=3V_\gamma Na^2p_s.
\]
Vary before setting $N=1$. With subsequent dots denoting proper time,
$H=\dot a/a$, $K=a^{-2}$ and $T_h=|D_th|^2/2$, the two metric equations are
\begin{align}
\label{p3f:eq:v2205-closed-matched-equations}
 3M^2\Omega(H^2+K)
 &=\rho_s+T_h+W-3M^2H\dot\Omega,\\
 -M^2\Omega(2\dot H+3H^2+K)
 &=p_s+T_h-W+M^2(\ddot\Omega+2H\dot\Omega).\nonumber
\end{align}
Both ports are fixed by the same action; a density calculation alone is
not a pressure calculation.
\end{theorem}
\begin{proof}
Insert $\mathcal A=M^2\Omega/2$ into
Eq.~\eqref{p3f:eq:v2203-completed-curvature}, add the canonical scalar
kinetic and potential terms, and retain the complete $S_s$.
The lapse derivative gives the first equation. The scale Euler derivative
$\partial_aL-d(\partial_{\dot a}L)/dt$ at $N=1$, divided by
$3V_\gamma a^2$, is
\[
 M^2\Omega(2\dot H+3H^2+K)
 +M^2(\ddot\Omega+2H\dot\Omega)+T_h-W+p_s.
\]
Its vanishing gives the second equation. Fixed endpoints remove only the
canonical total derivative, not either interior equation.
\end{proof}

\begin{corollary}[Fixed-Einstein-coefficient response stress]
\label{p3f:cor:v2205-fixed-coefficient-response}
Define $S_{\rm resp}$ as the difference between this completed action and
the already counted $S_s$ plus the completed reference gravity
$M^2(R_g-C)/2$. This is a decomposition of one action, not an added
fluid. Its homogeneous Hilbert ports are
\begin{align}
\label{p3f:eq:v2205-direct-response-ports}
 \rho_{\rm resp}&=T_h+W-\rho_\Lambda-3M^2H\dot\Omega
                -3M^2(\Omega-1)(H^2+K),\\
 p_{\rm resp}&=T_h-W+\rho_\Lambda
       +M^2(\ddot\Omega+2H\dot\Omega)
       +M^2(\Omega-1)(2\dot H+3H^2+K).\nonumber
\end{align}
On a solution of Eq.~\eqref{p3f:eq:v2205-closed-matched-equations}, these
are equivalently the explicit same-source-normalized expressions
\begin{align}
\label{p3f:eq:v2205-explicit-response-stress}
 \boxed{\rho_{\rm resp}
 =\frac{T_h+W-\rho_\Lambda\Omega-3M^2H\dot\Omega
                -(\Omega-1)\rho_s}{\Omega},}\\
 \boxed{p_{\rm resp}
 =\frac{T_h-W+\rho_\Lambda\Omega
       +M^2(\ddot\Omega+2H\dot\Omega)
                -(\Omega-1)p_s}{\Omega}.}\nonumber
\end{align}
Thus exact zero pressure on this domain is precisely
\begin{equation}
\label{p3f:eq:v2205-response-cold-condition}
 M^2(\ddot\Omega+2H\dot\Omega)
 =W-\rho_\Lambda\Omega-T_h+(\Omega-1)p_s.
\end{equation}
For the bare $h=0$ branch, $\Omega=F$ and
$W-\rho_\Lambda F=M^2\mathcal J/2$, recovering
Eq.~\eqref{p8:eq:v2203-local-source-normalized-remainder} exactly.
On the separately licensed bare affine terminal stratum,
$F=1$, $P=C$, $h=0$, both response ports vanish.
This statement concerns that stratum, not the substitution $R_g=C$:
the terminal vacuum solution instead has $R_g=2C$.
\end{corollary}
\begin{proof}
Vary the difference action in $N,a$ to get the direct ports. Use the two
full metric equations to eliminate the explicit curvature terms and collect
the same counted source once. Set the second numerator to zero for the
cold condition. The stated bare and affine specializations then follow.
\end{proof}

\begin{remark}[Action scope, canonical endpoint and retained sources]
\label{p3f:rem:v2205-response-action-scope}
The formulas evaluate the displayed completed local action and its actual
fields and source. They do not manufacture a nonlocal or mixed boundary
constituent outside it. A further admitted constituent contributes its own
complete two ports before any comparison-source subtraction. The canonical
three-dimensional endpoint stress has dimension $M^2H$, whereas the bulk
ports above have dimension $M^2H^2$. A temporal endpoint term alone has no
compactly supported interior variation and is not an extra cosmic fluid.
Stationary or causal elimination reuses the parent mean as in
Propositions~\ref{p3f:prop:v2203-stationary-envelope} and
\ref{p3f:prop:v2203-causal-mean-ports}; it does not duplicate it.
In particular the current no-second-minimal-mean statement, the isolated
native-scalar cold obstruction and the K48 material-dust scope are preserved.
\end{remark}


\subsection{Mixed metric response and the complete source-loaded energy}
\label{p3f:sec:v2206-mixed-loaded-energy}

This section executes the mixed metric derivative behind the two-port
stress of v2.205. It retains the complete stationary work of
Theorem~\ref{p3f:thm:v2203-complete-boundary-work}, with fixed reference
coordinates and all changing pairing, source, port, projector and endpoint
contacts. Euclidean spectral derivatives and physical Lorentzian energy
are kept on their separately licensed domains.

\begin{proposition}[Complete mixed heat derivative]
\label{p3f:prop:v2206-mixed-heat-jet}
On the same regular trace domain as
Proposition~\ref{p3f:prop:v1164-heat-hessian-no-insertion}, set
$X=D_{x,A}^2/\Lambda_g^2$ and $\mathcal F(X)=\operatorname{Tr}e^{-X}$.
Let $s$ be an admitted metric direction and $b$ a fixed boundary tangent.
Write $X_s,X_b,X_{sb},X_{bb},X_{sbb}$ for the complete pulled-back
derivatives of $X$. Then
\begin{equation}
\label{p3f:eq:v2206-mixed-heat-jet}
 \boxed{D_sD_b^2S_{\rm heat}
 =\mathcal F^{(3)}[X_s,X_b,X_b]
 +2\mathcal F^{(2)}[X_{sb},X_b]
 +\mathcal F^{(2)}[X_s,X_{bb}]
 +\mathcal F^{(1)}[X_{sbb}].}
\end{equation}
Here
\[
 \mathcal F^{(1)}[V]=-\operatorname{Tr}(e^{-X}V),\qquad
 \mathcal F^{(2)}[V,W]
 =\int_0^1\operatorname{Tr}(e^{-tX}Ve^{-(1-t)X}W)\,dt.
\]
For three labelled directions $V_1,V_2,V_3$, even when two coincide,
\begin{equation}
\label{p3f:eq:v2206-third-heat-ordered}
 \mathcal F^{(3)}[V_1,V_2,V_3]
 =-\sum_{\pi\in S_3}\int_{\Delta_3}
 \operatorname{Tr}(e^{-t_0X}V_{\pi(1)}e^{-t_1X}
 V_{\pi(2)}e^{-t_2X}V_{\pi(3)}e^{-t_3X})\,d^3t,
\end{equation}
where $t_i\ge0$, $\sum_{i=0}^3t_i=1$ and the simplex has volume $1/6$.
Actual localizer or trace-weight derivatives obey the same product rule;
they are not deleted by using this unlocalized presentation. Likewise
all variations of the inherited cutoff and metric/field embedding remain
inside the $X$ jets. No local $a_6$ addition or new response coefficient
is introduced by this identity.
\end{proposition}
\begin{proof}
The second composite derivative is
$\mathcal F^{(2)}[X_b,X_b]+\mathcal F^{(1)}[X_{bb}]$.
Differentiate it in $s$ to obtain the four terms. Iterated Duhamel
variation of $e^{-X}$ gives the ordered simplex formula. In a commuting
scalar block its six simplex terms reduce to
$-e^{-X}V_1V_2V_3$, verifying both multiplicity and sign.
\end{proof}

\begin{corollary}[Metric derivative of the actual Schur block]
\label{p3f:cor:v2206-metric-schur-jet}
For the fixed-reference quadratic blocks
$Q=A-B^\dagger C^{-1}B$ of the physical reduced action,
\begin{equation}
\label{p3f:eq:v2206-metric-schur-jet}
 \boxed{Q_s=A_s-B_s^\dagger C^{-1}B-B^\dagger C^{-1}B_s
             +B^\dagger C^{-1}C_sC^{-1}B.}
\end{equation}
For the completed port $\widetilde Q=T^\dagger M_\partial QT$, its
derivative additionally retains $T_s$, $(M_\partial)_s$ and $T_s^\dagger$.
If the background itself is the stationary $q_\star(s)$ of the full
action, its Hessian derivative is
$H_s=S_{qqs}-S_{qqq}[H^{-1}S_{qs}]$ wherever the full quotient inverse
exists. This is the metric specialization of the existing source-jet
identity, not differentiation on a frozen unrelated background.
\end{corollary}
\begin{proof}
Differentiate $C^{-1}$ and the physical port product. Differentiate
$S_q(q_\star(s),s)=0$ for the background velocity, then the Hessian.
This also gives $Q_s=L^\dagger H_sL$, $L=(I,-C^{-1}B)^T$.
\end{proof}

\subsubsection{The lapse derivative of a static loaded response}
\begin{theorem}[Stored energy, source work and net induced energy]
\label{p3f:thm:v2206-loaded-static-energy}
Suppose an already admitted static physical reduction, with its canonical
source marks and endpoints retained, has the Routh action
\[
 S_{\rm stat}=-\int dt\,\nu\left[
 E_0+\tfrac12b^\dagger\mathsf A b-\operatorname{Re}(j^\dagger b)\right],
 \qquad \mathsf A=\mathsf A^\dagger\succ0.
\]
$\nu$ is a uniform clock lapse and the relative spatial lapse, source
marks and physical spatial geometry are held fixed in its variation.
$\mathsf A,j,E_0$ therefore do not acquire a new $\nu$ dependence on
this static branch. They include the actual measure and port normalization
in fixed coordinates. This is a domain predicate on the parent action,
not a new oscillator or a claim about an arbitrary FLRW causal kernel.

Stationary elimination gives $b_\star=\mathsf A^{-1}j$ and
\begin{equation}
\label{p3f:eq:v2206-loaded-static-energy}
 \boxed{E_{\rm red}=-\frac{\delta S_{\rm red}}{\delta\nu}
 =E_0-\tfrac12j^\dagger\mathsf A^{-1}j.}
\end{equation}
Relative to the same unloaded reference, the stored field and source
interaction contributions at this point are
\[
 E_{\rm field}=\tfrac12j^\dagger\mathsf A^{-1}j\ge0,\qquad
 E_{\rm int}=-j^\dagger\mathsf A^{-1}j=-2E_{\rm field},\qquad
 \Delta E_{\rm red}=-E_{\rm field}.
\]
For a scale or shape strain $s$, their complete net work is
\begin{equation}
\label{p3f:eq:v2206-loaded-energy-strain}
 D_sE_{\rm red}=D_sE_0
 -\operatorname{Re}(j_s^\dagger\mathsf A^{-1}j)
 +\tfrac12j^\dagger\mathsf A^{-1}\mathsf A_s\mathsf A^{-1}j.
\end{equation}
Keeping only $E_{\rm field}$ changes the complete energy census and is
not a derivation of an additional positive mean.
\end{theorem}
\begin{proof}
Complete the positive quadratic square, or apply the stationary envelope.
The relative metric data are static, so the proper-time Hamiltonian is
multiplied by the uniform clock lapse once. Its lapse derivative gives the
displayed energy, including the already allocated $E_0$. Evaluate the
quadratic and linear pieces at $b_\star$ and differentiate the inverse.
The static global lapse variation is distinct from the active combined
strain $\delta\log N=-\delta\log a$.
\end{proof}

\begin{corollary}[No positive excess from passive global relaxation]
\label{p3f:cor:v2206-passive-relaxation-sign}
More generally, at fixed source marks let a complete static response
energy $\mathcal U(b;y)$ have an admitted reference $b=0$. If the
physical relaxed member is a global minimizer, then
\[
 \mathcal U(b_\star;y)-\mathcal U(0;y)\le0.
\]
Strict convexity and an attained minimum are a sufficient uniqueness
predicate. A positive response stiffness by itself is not a positive
extra energy. This statement is about the induced change relative to
the already allocated reference, not the sign of its absolute energy.
A metastable stationary branch, nonpassive preparation, dynamical work,
source-dependent baseline or different complete action must retain its
own mean and does not follow from this corollary.
\end{corollary}
\begin{proof}
Compare the minimum with the admitted reference configuration. No claim
that every native stationary boundary is such a minimum is made.
\end{proof}

\begin{remark}[The possible positive contribution has a different work test]
\label{p3f:rem:v2206-positive-loaded-excitation}
An already admitted positive kinetic block would instead have the exact
square decomposition
\[
 \mathcal H=\tfrac12\pi^\dagger\mathsf M^{-1}\pi
 +\tfrac12(b-b_\star)^\dagger\mathsf A(b-b_\star)
 +E_0-\tfrac12j^\dagger\mathsf A^{-1}j,\qquad \mathsf M\succ0.
\]
Its actually populated kinetic/displacement energy can be positive.
This does not promote the static constraint variable to a dynamical
mode: $\mathsf M$, the physical state, finite-band/order contacts and
central allocation must already be supplied by the parent. A quantum
reference ground energy is retained once; it is not an adjustable
abundance coefficient. Formation is tested by the existing complete
source work, not inferred from a positive Hessian or field covariance.
The physical $S+T$ branch and the actual cosmic history are consequently
not excluded by the static sign calculation above.
\end{remark}


\subsection{Homogeneous closure requires temporal metric jets and source work}
\label{p3f:sec:v2207-homogeneous-closure-work}

The complete common action remains the owner. The following results sharpen
the explicit overlap obligation of
Proposition~\ref{p3f:prop:v1164-scalar-overlap-condition}; they do not append a
new completion to it. A static force, its entire stationary source family,
and a finite carrier energy do not specify the physical homogeneous lapse
and scale derivatives. The causal quadratic calculation below applies only
to an already admitted physical dynamical port, not to a mode manufactured
from a static constraint.

\subsubsection{An explicit obstruction to inference from static data}
\begin{proposition}[Temporal stress is not identifiable from stationary response alone]
\label{p3f:prop:v2207-static-temporal-identifiability}
On an admitted observer/source branch, let $u^\mu$ be its unit timelike
velocity, $q$ an already retained scalar coordinate, and $b$ an advected
number density. For the purpose of testing identifiability only, consider
the covariant diagnostic functional
\begin{equation}
\label{p3f:eq:v2207-temporal-diagnostic}
 \Delta S_Z=\frac12\int\sqrt{-g}\,Z(b,q)
                  (u^\mu\nabla_\mu q)^2\,d^4x.
\end{equation}
It and every derivative tangent to the strictly stationary family vanish
where $u^\mu\nabla_\mu q=0$. Thus stationary action, source, boundary and
spatial-metric jets cannot distinguish these diagnostics.

On a homogeneous history with fixed advected census, $b\propto a^{-3}$,
its lapse-unfixed reduction is
\[
 \Delta L_Z=\frac{VZ(b,q)}{2N}\dot q^{\,2},\qquad V=2\pi^2a^3.
\]
At $N=1$, with the independent scalar data held fixed under metric variation,
the physical insertions are
\begin{equation}
\label{p3f:eq:v2207-temporal-diagnostic-stress}
 \boxed{\Delta\rho_Z=\frac12 Z\dot q^{\,2},\qquad
 \Delta p_Z=\frac12(Z-bZ_b)\dot q^{\,2}.}
\end{equation}
A positive constant $Z$ gives a stiff kinetic insertion; a positive
$Z$ proportional to $b$ gives zero scale pressure. Both are invisible to
the strictly static data.
\end{proposition}
\begin{proof}
The proper-time derivative vanishes identically throughout the stationary
family, so its squared functional and all derivatives restricted to that
family vanish. In the homogeneous restriction,
$-V^{-1}\partial_N\Delta L_Z=Z\dot q^2/(2N^2)$ and
$\partial_{\log a}b=-3b$. Applying the completed lapse and scale ports
gives Eq.~\eqref{p3f:eq:v2207-temporal-diagnostic-stress}.
\end{proof}

The symbol $Z$ is an identifiability witness, not an available Born--Hopf
parameter, an admitted extra action or an abundance fit. The full heat
action and its continuation must select or exclude such temporal jets.
In particular, an uncompensated diagnostic that changes the inherited
scalaron Hessian fails
Theorem~\ref{p3f:thm:v1164-scalaron-port-matching}. Its invisibility on the
static family does not make it invisible to that Lorentzian matching test.
This construction proves an inference limitation; it does not prove two
different completions of the native action are allowed.

\subsubsection{Two physical ports on an admitted dynamical reduction}
\begin{theorem}[Quadratic response energy, pressure and complete work]
\label{p3f:thm:v2207-quadratic-history-ports}
Suppose the existing action, canonical completion and physical quotient
actually supply a real finite dynamical response $q$, represented in a
fixed physical chart by
\begin{equation}
\label{p3f:eq:v2207-quadratic-lapse-action}
 W=\int dt\left[\frac{\dot q^T\mathsf M(a,t)\dot q}{2N}
 -\frac N2q^T\mathsf K(a,t)q+Nj(a,t)^Tq-NE_0(a,t)\right].
\end{equation}
The matrices and drive are generated data; $mathsf M,mathsf K$ are
symmetric. All independent source, measure, projector and metric contacts
are included in their partial derivatives. No unshown lapse derivative
is discarded: Eq.~\eqref{p3f:eq:v2207-quadratic-lapse-action} is an
assumption about the admitted complete reduced form. For a retarded
nonlocal reduction use the inherited two-history ports instead.

Let $\lambda=\log a$, set $N=1$ after variation, and subtract the same
reference constituents once. Then
\begin{align}
\label{p3f:eq:v2207-quadratic-energy-pressure}
 \boxed{E=V\rho=E_0+\frac12\dot q^T\mathsf M\dot q
          +\frac12q^T\mathsf Kq-j^Tq,}\\
 \boxed{3Vp=-E_{0,\lambda}
       +\frac12\dot q^T\mathsf M_{,\lambda}\dot q
       -\frac12q^T\mathsf K_{,\lambda}q+j_{,\lambda}^Tq.}\nonumber
\end{align}
The equation of motion and exact proper-time energy derivative are
\begin{align}
\label{p3f:eq:v2207-quadratic-source-work}
 \frac{d}{dt}(\mathsf M\dot q)+\mathsf Kq&=j,\\
 \dot E&=\dot E_0-\frac12\dot q^T\dot{\mathsf M}\dot q
             +\frac12q^T\dot{\mathsf K}q-\dot j^Tq.\nonumber
\end{align}
Splitting each total coefficient derivative into its metric scale part
$H\partial_\lambda$ and its explicitly changing source/background part
$\partial_t^{\rm exp}$ gives
\begin{equation}
\label{p3f:eq:v2207-quadratic-ward-work}
 \boxed{\dot E+3HVp=\partial_t^{\rm exp}E_0
 -\frac12\dot q^T(\partial_t^{\rm exp}\mathsf M)\dot q
 +\frac12q^T(\partial_t^{\rm exp}\mathsf K)q
 -(\partial_t^{\rm exp}j)^Tq.}
\end{equation}
The remaining source sectors carry the compensating work.
\end{theorem}
\begin{proof}
Take $-\partial_N W$ and $\partial_\lambda W$ before setting the lapse.
Differentiate $E$, insert the equation of motion, and cancel the terms
linear in $\dot q$. The metric-scale contributions are exactly $-3HVp$.
This also shows why field energy alone is not the completed insertion.
\end{proof}

For quantum use the same physical energy and pressure insertions, inherited
ordering and central allocation are evaluated in the one actually evolved
state. A time-dependent canonical generator must first be transported to
this physical chart; its chart-connection term is not additional energy,
as in Corollary~\ref{p8:cor:v2203-chart-work-firewall}.

\subsubsection{Executed causal loading, with the interaction retained}
\begin{theorem}[Monotone frozen loading cannot give a positive completed increment]
\label{p3f:thm:v2207-monotone-loading-sign}
On a supplied frozen positive one-mode dynamical block, put
$\mathsf M=m>0$, $\mathsf K=k=m\omega^2>0$ and let $j(t)$ rise from
zero to $j_f$ during $[0,T]$. Hold the reference $E_0$ fixed and start with
zero coherent response increment $q(0)=\dot q(0)=0$. This leaves any
already allocated covariance/vacuum baseline intact. For $t\ge T$,
\begin{equation}
\label{p3f:eq:v2207-forced-final-energy}
 \boxed{\Delta E=-\frac{j_f^2}{2k}
 +\frac1{2k}\left|\int_0^T\dot j(s)e^{-i\omega s}\,ds\right|^2.}
\end{equation}
The first term is the stationary binding decrement and the second is the
positive oscillation above the loaded minimum. If $j_f\ne0$ and
$j(t)/j_f$ is nondecreasing, then $\Delta E\le0$.
The same inequality holds for an admitted frozen positive finite block
driven as $j(t)=f(t)j_f$, with $f(0)=0$, $f(T)=1$, $\dot f\ge0$.
\end{theorem}
\begin{proof}
The retarded solution is
$q(t)=\int_0^t\sin[\omega(t-s)]j(s)\,ds/(m\omega)$.
Integrating by parts at and after the handoff gives the displacement and
velocity relative to $q_*=j_f/k$ as the two real quadratures of
$k^{-1}\int_0^T\dot j(s)e^{-i\omega s}ds$. Completing the physical
energy square gives Eq.~\eqref{p3f:eq:v2207-forced-final-energy}.
Monotonicity bounds the integral modulus by $|j_f|$. A fixed positive
finite block has an $\mathsf M$-orthonormal normal-mode basis; applying
the same bound mode by mode proves the final statement.
\end{proof}

For a linear ramp $j=j_ft/T$ on $[0,T]$ the result is exactly
\begin{equation}
\label{p3f:eq:v2207-ramp-energy}
 \boxed{\Delta E_{\rm ramp}=-\frac{j_f^2}{2k}
 \left[1-\operatorname{sinc}^2\!\left(\frac{\omega T}{2}\right)\right],
 \qquad \operatorname{sinc}x=\frac{\sin x}{x}.}
\end{equation}
The sudden monotone limit gives zero completed increment, although the
oscillation above the loaded minimum is positive. The adiabatic limit gives
the negative binding decrement. Neither permits adding that positive
oscillation while discarding the stationary interaction.

For comparison, a supplied pulse $j=j_0$ during $0<t<T$, with zero drive
afterwards, gives the positive coherent energy
\begin{equation}
\label{p3f:eq:v2207-pulse-energy}
 \boxed{\Delta E_{\rm pulse}=\frac{2j_0^2}{k}
                   \sin^2\!\left(\frac{\omega T}{2}\right)\ge0.}
\end{equation}
For a general finite pulse the exact expression is
$\Delta E=|\int_0^Tj(s)e^{-i\omega s}ds|^2/(2m)$.
This is energy supplied by the actual source work, not a free occupation.
The compensated source loses the corresponding work when its full action
is included. Histories containing an earlier pulse can therefore have the
same terminal static operator and drive but different retained energies.
Frozen coefficients, positive physical kinetic block, admitted forcing and
initial coherent data are essential to these tests. Background variation,
damping/memory, nonlinear response, nonmonotone modal source history and
initial excitations require the complete existing evolution instead.


\subsection{Explicit applications of the geometric provider identities}
\label{p3f:sec:v215-geometric-applications}
\subsubsection{Comparison with the specified relative-sphere operator}
\label{p3f:sec:v215-geometric-connection-comparisons}

\begin{proposition}[Exact bundle and operator comparisons]
\label{p3f:prop:v215-geometric-connection-comparisons}
The relative section connection \(A_{\rm B}=\Theta/2\) in
Paper~III-F, Theorem~\ref{p3f:thm:unique-born-invariant-gred-connection},
has the following exact comparisons with the preceding geometry.
Its fundamental rank-two representative is the Levi--Civita spin
connection of the same round relative sphere, after matching the
marking and Clifford convention of
Proposition~\ref{p1b:prop:providers-sphere-dictionary}.
Its curvature is nonzero and its restricted holonomy is \(SU(2)\).
With the first spin factor kept Levi--Civita, its rank-four twisted
spin contraction is exactly \(B_R\), including the scalar zero line
and the perturbation
\((f-1/2)\operatorname{diag}(-3,1,1,1)/R\) in
Theorem~\ref{p3f:thm:m8td-hodge-dirac}.

The two protected coordinate-sign projectors of Paper~II-C,
Theorem~\ref{p2c:thm:protected-projector-midpoint}, and Paper~III-F,
Corollary~\ref{p3f:cor:protected-projector-midpoint-ancestry}, are the
constant-midpoint cases of
Theorem~\ref{p1b:thm:providers-profile-projectors}.  Their projected
connections are thus exactly intertwined.  The globally flat tangent
dilation of Theorem~\ref{p1b:thm:providers-flat-dilation} gives the same
Levi--Civita connection only after the displayed orthogonal tangent
projection and then the rank-three spin lift.

These are comparisons with an already specified connection.  The flat
Bismut endpoint is distinct from its curved midpoint by curvature and
holonomy; the tangent-torsion deformation is distinct from its twisting
deformation by Theorem~\ref{p1b:thm:providers-torsion-obstruction}.
None of these comparisons selects a member of a larger family.
\end{proposition}
\begin{proof}
The source theorem already supplies the section matrix, radius and
fundamental twisting representation.  Apply the spin-lift and Hodge
computations of Proposition~\ref{p1b:prop:providers-sphere-dictionary}
in the matched marked frame.  The protected midpoint frames are exactly
\((U,I)^t/\sqrt2\) and \((-U,I)^t/\sqrt2\); differentiating gives
the same matrix connection for both.  The dilation comparison is its
proved Euclidean normal-bundle construction followed by projection.
Curvature transforms by conjugation under a unitary gauge equivalence,
so the flat Bismut connection cannot equal the curved midpoint under
such a map.  The obstruction theorem proves the remaining operator
distinction.  Each comparison therefore uses an explicit representation,
bundle isomorphism or projection, rather than equality of one spectral
number or a new family-member choice.
\end{proof}

\subsubsection{The singular continuation of the existing boundary relation}
\label{p3f:sec:geo-singular-boundary-application}

\begin{corollary}[Same-action boundary relation at a singular interior block]
\label{p3f:cor:geo-singular-boundary-application}
Retain the actual action, stationary solution, boundary-domain split and
finite gauge/constraint quotient of
Theorem~\ref{p3f:thm:v167-source-conditioned-boundary-schur}. Write its
frozen Hessian and actual source as
\[
 H=\begin{pmatrix}A&B\\B^*&C\end{pmatrix},\qquad
 j=(j_\partial,j_I),\quad N=\ker C,\quad S=A-BC^+B^*.
\]
At a singular interior block the exact boundary descendant is the relation
\[
 Sx+Bn=j_\partial-BC^+j_I,\qquad
 P_NB^*x=P_Nj_I,\qquad n\in N.
\]
It is the stationary relation of the constrained reduced quadratic
function in Paper I-C, Lemma~\ref{p1b:lem:geo-constrained-quadratic}.
It reduces to an unconstrained boundary Hessian only when \(BN=0\),
and still requires \(P_Nj_I=0\). Positivity of this actual frozen
Hessian implies the coupling condition; positivity of a different
kinematical metric does not imply it.

If the actual source and trace use the same supplied linear map \(L\),
so that \(Hq=L^*h\) and the readout is \(Lq\), then the null directions
are invisible precisely when \(L\ker H=0\). In that case all such
sources are compatible and the unique readout is \(LH^+L^*h\).
When \(H\succeq0\), \(S\succ0\), and \(L=(L_\partial,0)\), it equals
\(L_\partial S^{-1}L_\partial^*h\).
\end{corollary}
\begin{proof}
Apply the stated I-B lemma to this Hessian and the actual split. Every
block is a second variation of the same action, so its range constraint
and null variable belong to its stationary relation. Apply
Theorem~\ref{p1b:thm:geo-null-output} to the supplied source/trace map
for the last statements. No additional boundary operator is introduced.
\end{proof}

\begin{corollary}[Crossing tests for the same response]
\label{p3f:cor:geo-boundary-crossing-tests}
For an analytic actual interior pencil whose compression of
\(-C'(z_0)\) to \(\ker C(z_0)\) is positive definite, the residue and
the independent source constraint are exactly those of Paper I-C,
Theorem~\ref{p1b:thm:geo-interior-laurent}. A vanishing interior Schur
residue does not remove a source component in its uncoupled nullspace.
If the actual full pencil has the form \(K-zM\), \(M\succ0\), its
readout residues are the positive weights of
Theorem~\ref{p1b:thm:geo-full-pencil-residues}. An internal Schur pole
is therefore not itself a pole of the full response. Pencils outside
that positive-metric hypothesis retain their actual derivative pairing,
as in Theorem~\ref{p3f:thm:v1140-finite-pole-source-transport}.

The regular source jets of
Theorem~\ref{p3f:thm:v176-explicit-source-jet-schur} retain their printed
invertibility hypotheses. At a singular stationary point their necessary
replacement compatibility conditions are those of
Proposition~\ref{p1b:prop:geo-singular-stationary-jets}; they do not
establish a nonlinear branch by themselves. Even a positive family with
a decoupled null mode at the crossing can have a discontinuous regular
Schur limit. Its exact continuity condition is
Theorem~\ref{p1b:thm:geo-positive-rank-change}. Constant-rank
pseudoinverse jets cannot replace that test at a rank change.
\end{corollary}
\begin{proof}
Each assertion is the stated I-B result applied to the actual response
matrix or stationary action. The hypotheses specify which reduction and
which inverse are being differentiated. They leave the actual source,
trace, preparation, contacts, and any retarded initial-value domain with
their existing owners.
\end{proof}

\subparagraph{The first-order comparison remains on its own carrier.}
The twenty-five-dimensional first-order realization of Paper I-A,
Theorem~\ref{p1:thm:v214-first-order-retained-closure}, is an exact instance
of Paper I-C, Proposition~\ref{p1b:prop:geo-bipartite-resolvent}: the
paired form space has dimension ten and
\(\Lambda=(3/\Rmom^2)P_4+(4/\Rmom^2)P_6\). Its retained first-order
resolvent on those forms is \(z(\Lambda-z^2I)^{-1}\), holomorphic at
zero, despite the internal \(\Lambda/z\) Schur coordinate. The original
constant line supplies the sole zero pole, and the spinor block retains
its existing eigenvalues. This calculation does not enlarge the defined
physical level-one carrier or give the ten auxiliary partners a particle
interpretation. Applying it to a physical boundary response requires
that response's actual action-to-operator and source/trace maps.

\subparagraph{A closing positive transition gap.}
Whenever an actual stationary finite passive source response has the
positive strict-gap spectral form of Paper I-C,
Proposition~\ref{p1b:prop:geo-positive-gap-scaling}, that proposition
determines its closing-gap rates. Its spectral contribution can remain
uniformly bounded only with the corresponding actual transition-weight
scaling. This applies to the same state and source port; its direct,
preparation, degenerate-level and constraint contacts remain separate.
It supplies no continuation of the low-frequency Taylor window through
a nonzero coupled pole at zero.

\subsection{The geometric Cartan profile on an existing scalar--gauge action}
\label{p3f:sec:v215-cartan-action-composition}
The following composition evaluates a restriction of the existing gauge
action. It does not choose a connection from the comparison families of
Paper I-C or identify an internal sphere with a cosmological spatial
sphere. Its domain is a common regular metric realization on which the
covariant local action of
Theorem~\ref{p3f:thm:v174-generated-joint-gauge-phase-space-action}
and the supplied positive scalar-dependent gauge coefficient are both
defined. For the scalaron coefficient of Paper V-C,
Theorem~\ref{p7:thm:v164-scalar-gauge-source}, the scalar coordinate
is the same canonical Einstein field only when its Einstein coefficient
agrees with that in the conformal transport of this paper.

\begin{proposition}[Exact homogeneous restriction and its energy exchange]
\label{p3f:prop:v215-cartan-scalar-gauge-restriction}
On that common domain put
\[
 ds^2=a(\eta)^2(-d\eta^2+d\Omega_3^2),\quad
 A_0=0,\quad A=f(\eta)\Theta,\quad
 C=12\pi^2/g_2^2,\quad V(f)=4f^2(1-f)^2,
\]
with fixed matching scale and homogeneous $Z(\phi_c)>0$.
The geometric connection is $A=g_2W_c^aT_a$, so the canonical
$-ZF_c^2/4$ term is $-ZF_A^{a,2}/(4g_2^2)$. Its restriction is
\[
 S_f=C\int Z(f'^2-V(f))\,d\eta.
\]
If the actual remaining action supplies
$j=(2C)^{-1}\delta S_{\rm other}/\delta f$, its projected equation is
\[
 \boxed{(Zf')'+4Zf(1-f)(1-2f)=j.}
\]
For $j=0$ this is a consistent homogeneous isotropic truncation: its
temporal Gauss equation vanishes and its spatial equation has only
the displayed invariant component. With a general source it is a full
homogeneous solution only if the omitted spatial and Gauss components
also vanish. Its conformal energy and proper energy density obey
\[
 E_c=CZ(f'^2+V),\qquad
 E_c'=CZ'(V-f'^2)+2Cjf',\qquad
 \rho_f=\frac{6Z}{g_2^2a^4}(f'^2+V),\quad p_f=\rho_f/3.
\]
For $j=0$ and a canonical scalar with potential $\mathcal U$,
\begin{align*}
 \dot\rho_f+4H\rho_f
   &=\frac{6\dot Z}{g_2^2a^4}(V-f'^2),\\
 \ddot\phi_c+3H\dot\phi_c+\mathcal U_{,\phi_c}
   &=\frac{6Z_{,\phi_c}}{g_2^2a^4}(f'^2-V).
\end{align*}
Thus exchange is fixed by the same action and has opposite signs in
the scalar and gauge energy balances. It is not a prescribed positive
friction coefficient.
\end{proposition}
\begin{proof}
The curvature is $F=d\eta\wedge f'\Theta+f(f-1)\Theta^2$.
The round-sphere contractions already normalized in Paper I-C give
$\|\Theta\|^2=12\pi^2$ and $\|\Theta^2\|^2=48\pi^2$.
The four-dimensional conformal factors cancel in the gauge action,
giving $S_f$. Its Euler equation and differentiation of $E_c$ give
the first identities. Homogeneity and the locked rotation/color
invariance leave one spatial equation, while the homogeneous Gauss
contraction vanishes. Dividing the proper energy $E_c/a$ by
$2\pi^2a^3$ gives $\rho_f$; the isotropic gauge stress has zero trace,
so $p_f=\rho_f/3$. Use $dt=a\,d\eta$ in the energy identity.
Variation of the same $Z(\phi_c)$ in $S_f$, together with the canonical
scalar action, gives the final scalar equation and the opposite work.
\end{proof}

\begin{corollary}[The represented lowest coexact equation]
\label{p3f:cor:v215-cartan-n2-identification}
Near the flat vacuum put $v=\sqrt{2CZ}\,f$. The linear equation is
\[
 v''+\left[4-\frac{(\sqrt Z)''}{\sqrt Z}\right]v
                         =\sqrt{2C}\,j/\sqrt Z.
\]
This is the $n=2$ member of the coexact round-sphere scalar equations
for this actual gauge action, restricted to its Cartan combination.
Near $f=1/2$ the corresponding constant $4$ is instead $-2$, the
already present negative Hessian direction. For $j=0$, each constant
$f=0,1,1/2$ with zero velocity remains an exact solution for every
smooth positive $Z$. A varying scalar coefficient therefore does not
force a nonzero coherent classical profile from these exact data.

For the existing protected singlet coordinate
$u=a\lambda_0=3/2-3f$, the same projected equation becomes
\[
 u''+(Z'/Z)u'-2u+\tfrac89u^3=-3j/Z.
\]
Its radius dependence is explicit. It is not obtained by treating
$\lambda_0$ as the same coordinate at every radius.
\end{corollary}
\begin{proof}
Linearize the preceding Euler equation and remove its first derivative
by multiplication by $\sqrt Z$. At the midpoint use
$4f(1-f)(1-2f)=-2(f-1/2)+8(f-1/2)^3$.
The constant solutions follow directly. Substitute $f=1/2-u/3$ for
the last formula, using the protected eigenvalue already calculated
in Paper I-C on its spin--twist carrier.
\end{proof}

The flat dilation and projector constructions of Paper I-C give
representations of this same Cartan connection. They do not change its
coupling or select a realized $f(\eta)$. In particular an unweighted
Euclidean kink is not automatically critical after a varying $Z$ is
inserted: its weighted Euler residual contains $Z'f_*'$. The abstract
weighted-profile result of Paper I-C states the exact condition.


\section{Homogeneous shear two-jet and curvature/source transport firewall}
\label{bh:berger:v117:p3f:section}
The following results are local tensor-action diagnostics and
first-order ordinary-packet comparisons. They do not identify the
complete physical Laurent/cap action by asserting that a scalar
$f(R)$ functional already contains the ordinary Dirac Weyl term.

\begin{proposition}[Five-component homogeneous shape two-jet]
\label{bh:berger:v117:p3f:twojet}
Let $\vartheta_i=\sigma_i/2$ be the fixed unit-$S^3$ coframe, and
$B(\eta)\in\operatorname{Sym}^2_0(\mathbb R^3)$ a small homogeneous
position-space shear. In conformal time choose
$g=a(\eta)^2[-d\eta^2+\vartheta^{\mathsf T}e^{2B}\vartheta]$.
For compactly supported interior variations its geometric two-jet is
\[
\begin{split}
 R[g]-R[g_{\mathrm{FLRW}}]
 &=a^{-2}[\operatorname{tr}(B'^2)-8\operatorname{tr}(B^2)]+O(B^3),\\
 C_{\mu\nu\rho\sigma}C^{\mu\nu\rho\sigma}
 &=a^{-4}[2\operatorname{tr}(B''-8B)^2
             -72\operatorname{tr}(B'^2)]+O(B^3).
\end{split}
\]
The Euler/Gauss--Bonnet interior metric variation is a boundary term.
On the \emph{separately specified comparison functional}
$\int\sqrt{-g}[\kappa_0f(R)+\beta C^2]$, its shape term is
\[
\begin{split}
 S_B^{(2)}={}&2\pi^2\kappa_0\int d\eta\,a^2f_R
      [\operatorname{tr}(B'^2)-8\operatorname{tr}(B^2)]\\
 &+4\pi^2\beta\int d\eta
      [\operatorname{tr}(B''-8B)^2-36\operatorname{tr}(B'^2)].
\end{split}
\]
The Weyl equation has the common operator
$(\partial_\eta^2+4)(\partial_\eta^2+16)$ on the five
homogeneous shear components; it does not itself select a marked
icosahedral orientation. An uncancelled Weyl term distinguishes this
action from pure metric $f(R)$ on a fixed off-FLRW metric chart.
\end{proposition}
\begin{proof}
Use the Milnor curvature for the volume-preserving shape and the
traceless homogeneous extrinsic curvature. In the linearized
orthonormal frame with round brackets
$[e_i,e_j]=2\epsilon_{ijk}e_k$, the electric and magnetic Weyl
pieces are linear in $B''-8B$ and $B'$. Their squared Lorentzian
contraction gives the displayed trace identity. Independent axial
and triaxial shear tangents, followed by polarization and rotation
invariance, determine the entire quadratic form. Multiplying by
the round volume $2\pi^2$ gives the action. Integrating the Weyl
term by parts gives the factored operator; Euler is topological
in the interior with its boundary completion retained.
\end{proof}

\begin{proposition}[Ordinary native curvature/source transport on a classical Maxwell branch]
\label{bh:berger:v117:p3f:transport}
On the \emph{ordinary} torsion-free Dirac packet of
Proposition~\ref{p3f:prop:v224-native-tensor}, the coefficients are
$\kappa_J=12U$, $\beta_J=3/(160\pi^2)$ and
$\eta_J=-11/(960\pi^2)$, with
\[
 \beta_J C^2+\eta_J\mathcal E_4
 =-\frac{R^2}{80\pi^2}+\frac{3}{80\pi^2}
    R_{\mu\nu}R^{\mu\nu}+\frac{7}{960\pi^2}\mathcal E_4.
\]
Under the \emph{complete first-order} field/source/boundary metric
transformation of Theorem~\ref{p3f:thm:v224-stress-transport}, the
Ricci-squared terms are transported to the inherited source functional
$P(H)$, with $P(\mathrm{Ric})-P(H)=-E_{\mu\nu}Y^{\mu\nu}$.
On a classical leading Einstein--Maxwell solution with parallel
fields, $T=0$, $T_{\mu\nu}T^{\mu\nu}=4\rho^2$ and
\[
 P(H)=-\frac{\Lambda_0^2}{20\pi^2}
       +\frac{\rho^2}{3840\pi^2U^2}.
\]
The Berger coordinate shifts to first order by
\[
 \Delta s:=s_g-s_{\widetilde g}
 =-\frac{b_R\rho}{\kappa_J^2}+O(b_R^2),
 \qquad b_R=\frac3{80\pi^2}.
\]
The source-dependent contact obeys
$\Delta\mathcal L_{\gamma}=-\rho\Delta s$ after subtracting
the source-free vacuum term, and the near-round gauge splitting
chart satisfies $\Delta r=(8/5)\Delta s$ to the same order.
\end{proposition}
\begin{proof}
Use the four-dimensional identity
$C^2=\mathcal E_4+2\operatorname{Ric}^2-2R^2/3$ and apply the
already established stress-transport theorem to the complete
ordinary coefficient owner. Its change of metric variables cancels
the first-order Ricci-quadratic remainder against the variation of
the full leading field/source action. Contract the displayed Maxwell
stress $T^\mu{}_{\nu}=\operatorname{diag}(-\rho,\rho,\rho,-\rho)$:
$T=0$ and $\operatorname{tr}(T^2)=4\rho^2$.
The source-dependent traceless part of the metric transformation
is $-b_R T/\kappa_J^2$, producing $\Delta s$; the displayed
contact and gauge-chart relation follow by the same first-order
transformation, not from adding two versions of the correction.
\end{proof}
\begin{remark}[Three gravitational-action owners must remain separated]
\label{bh:berger:v117:p3f:ownership}
The ordinary $\kappa_J=12U$ native packet, its \emph{whole-packet}
$1/96$ retained normalization, and the separate active Laurent/cap
scalar/tensor owner are not interchangeable gravitational coefficient
choices. The ordinary transport is perturbative and carries the
metric, clocks, matter fields, physical source measures, localizers,
detector operators and boundary data together. Do not add the
transported stress contact on top of its untransported Weyl term;
do not infer an extra physical high-frequency pole from a truncated
fourth-order action. The complete quantum composite stress contact
and the full physical action identification remain open.
\end{remark}
\begin{corollary}[Holonomy obstruction to a globally parallel marked shear]
\label{bh:berger:v117:p3f:holonomy}
On the already generated Born midpoint connection with full
restricted holonomy $SU(2)$, the associated real spin-two bundle
with fibre $\operatorname{Sym}^2_0(\mathbb R^3)$ has no nonzero
globally parallel section over a connected domain realizing that
holonomy. This does not identify that associated bundle with the
physical spatial tensor bundle without a source/transport map.
\end{corollary}
\begin{proof}
A parallel value is fixed by every holonomy rotation; the only
$SO(3)$-fixed real symmetric matrices are scalar matrices. Trace
zero forces the fixed matrix to be zero.
\end{proof}


\section{Causal evolution, conserved returns, and transverse response}
\label{paper3f:sec:organized-causal}
\subsection{Native spherical action and the coupled local Cauchy system}
\label{p3f:sec:native-spherical-cauchy}

The action of Definition~\ref{p3f:def:v176-local-bosonic-noether-domain}
already owns the scalar--metric coefficients. The following calculation
evaluates its homogeneous spherical restriction and retains a complete finite
source supplied by the metric constructor. It adds no interaction or regulator.
The classical/source partition is fixed before evaluation, so each term is
counted once. All gauge, symmetry, occupation and constraint marks consumed
by that restriction remain required.

\subsubsection{Complete finite-source principal matrix}
\begin{proposition}[Source-corrected Cauchy matrix and action energy]
\label{p3f:prop:complete-source-cauchy-matrix}
On a licensed finite first-jet constraint/gauge chart, use the complete
classical Lagrangian \(L_{\rm cl}(q,v)\) and Hermitian finite generator
\(H(q,v)\), in common physical action units, with \(v=\dot q\) and the
full source density matrix \(\sigma\). Retain central, interaction,
embedding, projection and moving-frame terms. For a closed autonomous member,
\begin{align}
 P_i&=L_{{\rm cl},v_i}-\operatorname{Tr}(\sigma H_{,v_i}),&
 \dot P_i&=L_{{\rm cl},q_i}-\operatorname{Tr}(\sigma H_{,q_i}),&
 \dot\sigma&=-\frac{i}{\hbar}[H,\sigma].
 \label{p3f:eq:complete-source-cauchy-momenta}
\end{align}
The exact acceleration equation is \(\mathsf M_{ij}\dot v_j=b_i\), where
\begin{equation}
 \boxed{\mathsf M_{ij}=L_{{\rm cl},v_iv_j}
 -\operatorname{Tr}(\sigma H_{,v_iv_j}),}
 \label{p3f:eq:complete-source-cauchy-matrix}
\end{equation}
\begin{equation}
 b_i=L_{{\rm cl},q_i}-L_{{\rm cl},v_iq_j}v_j
 -\operatorname{Tr}\!\left\{\sigma\left[
 H_{,q_i}-H_{,v_iq_j}v_j-\frac{i}{\hbar}[H,H_{,v_i}]
 \right]\right\}.
 \label{p3f:eq:complete-source-cauchy-force}
\end{equation}
Explicit time dependence adds
\(-L_{{\rm cl},v_it}+\operatorname{Tr}(\sigma H_{,v_it})\) to \(b_i\).
On \(\det\mathsf M\ne0\), the normal form is
\(\dot q=v\), \(\dot v=\mathsf M^{-1}b\), with the displayed full state
equation. Its conserved autonomous action energy is
\begin{equation}
 \boxed{\mathcal H_{\rm tot}
 =v_iL_{{\rm cl},v_i}-L_{\rm cl}
 +\operatorname{Tr}\{\sigma(H-v_iH_{,v_i})\}.}
 \label{p3f:eq:complete-source-action-energy}
\end{equation}
\end{proposition}
\begin{proof}
Paper IV's complete functional insertion retains the actual \(H_{,v_i}\).
Use the exact finite-state identity
\[
 \frac{d}{dt}\operatorname{Tr}(\sigma X)
 =\operatorname{Tr}\!\left[\sigma
 \left(\dot X+\frac{i}{\hbar}[H,X]\right)\right]
\]
in the momentum equation and collect accelerations. This gives both the
source Hessian correction and the commutator force. Direct differentiation
of \(\mathcal H_{\rm tot}\), followed by the momentum and state equations,
gives zero. No one-body factorization is used. Independently retained
open-source forces/work must instead be included in their actual equations.
\end{proof}
Positive normalization of the finite source kinetic form does not itself
prove invertibility of this coupled matrix. The lapse constraint is varied
before fixing the lapse; conservation of this energy is not a positive
coercive estimate on a gravitational constraint surface.

\subsubsection{Spherical reduction with the native coefficients}
Write \(\alpha_\Phi=(16\pi^2)^{-1}\), the \(\alpha\) printed in the
native scalar-action packet; this local notation distinguishes it from other
series readouts. Put
\[
 \mathscr F=\kappa_J-\frac{\alpha_\Phi}{6}\operatorname{Tr}_F\Phi^2,
 \qquad V=\alpha_\Phi\operatorname{Tr}_F
 (\Phi^4-2\Lambda_g^2\Phi^2).
\]
The already owned \(\kappa_J=12\Lambda_{\rm EH}^2>0\) does not identify
\(\Lambda_{\rm EH}\) with \(\Lambda_g\).
The identity~\eqref{p3f:eq:v267-joint-eh-newton-readout} fixes physical
action normalization: \(S_{\rm phys}=\hbar S_{\rm loc}\). Below the
printed native scalar--metric equations use proper-time \(c=1\), and
physical residual stress/scalar insertions are divided by \(\hbar\).
Equivalently evaluate this native normal form in \(\hbar=c=1\).
In Proposition~\ref{p3f:prop:complete-source-cauchy-matrix}, the physical
classical member is \(L_{\rm cl}=\hbar L_{g\Phi}\).


\begin{corollary}[Native positive-graviton scalar coordinate bound]
\label{p3f:cor:v2174-scalar-coordinate-bound}
On the admitted native chart $\mathscr F(\Phi)>0$, use the already owned
$\alpha_\Phi=(16\pi^2)^{-1}$ and $\kappa_J=12\Lambda_{\rm EH}^2$.
Then
\begin{equation}
 \operatorname{Tr}_F\Phi^2
   <\frac{6\kappa_J}{\alpha_\Phi}
   =1152\pi^2\Lambda_{\rm EH}^2.
 \label{p3f:eq:v2174-scalar-coordinate-bound}
\end{equation}
For an admitted $\Phi=rT$ with $r>0$, $\operatorname{Tr}_FT^2=1$ and
$\Lambda_{\rm EH}>0$, this is $r<24\sqrt2\pi\Lambda_{\rm EH}$.
\end{corollary}
\begin{proof}
Rearrange the inherited inequality
$\kappa_J-\alpha_\Phi\operatorname{Tr}_F\Phi^2/6>0$ and substitute
the two coefficients. The normalized ray has $\operatorname{Tr}_F\Phi^2=r^2$.
\end{proof}
This bounds a coordinate on its chart, not its momentum or occupation
capacity. It does not identify $\Lambda_{\rm EH}$ with $\Lambda_g$.
The native mixed Legendre relation is evaluated in Paper VI-A,
Corollary~\ref{p8:cor:v2174-native-turning-momentum}; the scoped capacity
comparisons are in Paper~III-G, Appendix~\ref{p3f:app:v2174-native-capacity-domain}.


\begin{proposition}[Native spherical action and exact kinetic determinant]
\label{p3f:prop:native-spherical-kinetic-determinant}
On an admitted homogeneous scalar restriction with \(D_i\Phi=0\), retain
the lapse, and after its variation fix a regular temporal/spatial gauge.
For the same axial integration/localizer datum \(L_z\), set
\(\Omega_z=4\pi L_z\), \(V_3=\Omega_z AB^2\), and
\[
 ds^2=-N^2dt^2+A^2dz^2+B^2d\Omega_2^2.
\]
The retained interior integration by parts gives
\begin{equation}
\begin{split}
 L_{g\Phi}=\Omega_z\bigg[&2NA\mathscr F
 -\frac{4\mathscr FB\dot A\dot B+2\mathscr FA\dot B^2
 +(2B^2\dot A+4AB\dot B)\dot{\mathscr F}}{N}\\
 &+\frac{\alpha_\Phi AB^2}{N}\operatorname{Tr}_F\dot\Phi^2
 -NAB^2V\bigg].
\end{split}
\label{p3f:eq:native-spherical-action}
\end{equation}
Temporal-gauge representatives use the actual \(D_t\Phi\).
In orthonormal real retained scalar coordinates \(\phi_a\), its velocity
Hessian satisfies
\begin{gather}
 M_{gg}=\frac{\Omega_z}{N}
 \begin{pmatrix}0&-4\mathscr FB\\-4\mathscr FB&-4\mathscr FA\end{pmatrix},
 \quad M_{A,a}=-\frac{2\Omega_zB^2}{N}\mathscr F_{,a},
 \quad M_{B,a}=-\frac{4\Omega_zAB}{N}\mathscr F_{,a},\nonumber\\
 M_{ab}=\frac{2\alpha_\Phi\Omega_z AB^2}{N}\delta_{ab},
 \label{p3f:eq:native-spherical-kinetic-blocks}\\
 \boxed{\det M_{g\Phi}
 =-16\left(\frac{\Omega_z}{N}\right)^2 B^2\mathscr F\kappa_J
 \left(\frac{2\alpha_\Phi\Omega_z AB^2}{N}\right)^{d_\Phi}.}
 \label{p3f:eq:native-spherical-kinetic-determinant}
\end{gather}
Thus this block is regular for \(A,B,N,\mathscr F,\kappa_J>0\).
The complete source-corrected matrix uses
\eqref{p3f:eq:complete-source-cauchy-matrix}; a coordinate-only proper
source plus its affine connection term has \(H_{,vv}=0\).
\end{proposition}
\begin{proof}
The spatial curvature is \(2/B^2\), and the extrinsic kinetic combination
is \(-4\dot A\dot B/(N^2AB)-2\dot B^2/(N^2B^2)\).
The \(\mathscr F R\) integration by parts supplies
\(-2\Omega_z(B^2\dot A+2AB\dot B)\dot{\mathscr F}/N\).
Together these give the action. Differentiate its quadratic velocity word.
The scalar Schur complement factors the determinant with
\[
 \mathscr F+\frac{3}{2\alpha_\Phi}|\nabla\mathscr F|^2
 =\mathscr F+\frac{\alpha_\Phi}{6}\operatorname{Tr}_F\Phi^2
 =\kappa_J,
\]
since \(\mathscr F_{,a}=-\alpha_\Phi\phi_a/3\).
This is the spherical form of the native ADM trace-block identity, with no
change of its coefficient. Transported endpoint terms remain in any
boundary claim.
\end{proof}

\subsubsection{Trace and source-work closure}
Let \(T_r=-\rho_r+p_{z,r}+2p_{\perp,r}\) and \(J\) be the complete
\emph{action-normalized} residual trace and scalar insertion, including
central energy and mass-owner work. The native scalar equation is
\[
 2\alpha_\Phi D_\mu D^\mu\Phi
 -\frac{\alpha_\Phi}{3}R\Phi-V_{,\Phi}-J=0.
\]
\begin{proposition}[Native curvature trace with the residual source retained]
\label{p3f:prop:native-source-curvature-trace}
On the common first-jet action equations,
\begin{equation}
 \boxed{2\kappa_JR=4V-\operatorname{Tr}_F(\Phi V_{,\Phi})
 -T_r-\operatorname{Tr}_F(\Phi J)
 =-4\alpha_\Phi\Lambda_g^2\operatorname{Tr}_F\Phi^2
 -T_r-\operatorname{Tr}_F(\Phi J).}
 \label{p3f:eq:native-source-curvature-trace}
\end{equation}
\end{proposition}
\begin{proof}
Trace the native metric equation
\(2\mathscr FG_{\mu\nu}+2(g_{\mu\nu}\Box-\nabla_\mu\nabla_\nu)
\mathscr F=T_{\mu\nu}^{\Phi,\min}+T^r_{\mu\nu}\).
Use
\(\Box\mathscr F=-\alpha_\Phi
\operatorname{Tr}_F(\Phi D^2\Phi+D_\mu\Phi D^\mu\Phi)/3\)
and the scalar equation. The scalar kinetic trace cancels; the coefficient
of \(R\) is \(-2\kappa_J\). The quartic cancellation uses
\(V_{,\Phi}=4\alpha_\Phi(\Phi^3-\Lambda_g^2\Phi)\).
Cyclicity suffices even for noncommuting Hermitian fields and insertions.
\end{proof}
For coordinate-only residual ports this gives an algebraic \(R_*\).
When residual functional insertions consume higher derivatives it is an
implicit relation in that complete higher-jet problem, not a first-order
closure. No prescribed curvature cap is used.

\subsubsection{Explicit constrained normal form}
\begin{theorem}[Local spherical first-jet Cauchy construction]
\label{p3f:thm:native-spherical-local-cauchy}
Use the licensed scalar--metric member on a proved compatible homogeneous
gauge/source restriction and a smooth coordinate-only proper source
\(h(A,B,\phi)\) in its actual flat canonical source frame. Retain the full
finite state, constraints and central energy. Its action-normalized ports are
\[
 \rho_r=\frac{\operatorname{Tr}(\sigma h)}{\hbar V_3},\quad
 p_{z,r}=-\frac A{\hbar V_3}\operatorname{Tr}(\sigma h_{,A}),\quad
 p_{\perp,r}=-\frac B{2\hbar V_3}\operatorname{Tr}(\sigma h_{,B}),\quad
 J_a=\frac{\operatorname{Tr}(\sigma h_{,\phi_a})}{\hbar V_3}.
\]
With \(N=1\), \(u=\dot A/A\), \(v=\dot B/B\), \(\theta=u+2v\), and
\(\Pi=D_t\Phi\), define \(R_*\) by
\eqref{p3f:eq:native-source-curvature-trace}, and set
\begin{align}
 \dot\Phi&=\Pi,&
 \dot\Pi&=-\theta\Pi-\frac{R_*}{6}\Phi
 -\frac{V_{,\Phi}+J}{2\alpha_\Phi},
 \label{p3f:eq:native-spherical-scalar-flow}\\
 f_1&=-\frac{\alpha_\Phi}{3}\operatorname{Tr}_F(\Phi\Pi),&
 f_2&=-\frac{\alpha_\Phi}{3}
 \operatorname{Tr}_F(\Pi^2+\Phi\dot\Pi).\nonumber
\end{align}
Write \(\rho,p_z,p_\perp\) for the residual ports plus
\(\rho_\Phi=\alpha_\Phi\operatorname{Tr}_F\Pi^2+V\) and
\(p_\Phi=\alpha_\Phi\operatorname{Tr}_F\Pi^2-V\). Then
\begin{align}
 Y:=\ddot B/B&=-\frac{p_z+2f_2+4vf_1}{4\mathscr F}
 -\frac{v^2+B^{-2}}2,\nonumber\\
 X:=\ddot A/A&=-\frac{p_\perp+2f_2+2(u+v)f_1}{2\mathscr F}-Y-uv,
 \label{p3f:eq:native-spherical-spatial-flow}\\
 \dot A&=uA,\quad\dot B=vB,\quad\dot u=X-u^2,\quad\dot v=Y-v^2,
 \quad\dot\sigma=-i[h,\sigma]/\hbar.\nonumber
\end{align}
The lapse expression, varied before fixing \(N\), is
\begin{equation}
 \boxed{\mathcal C=2\mathscr F(2uv+v^2+B^{-2})+2\theta f_1-\rho.}
 \label{p3f:eq:native-spherical-lapse-constraint}
\end{equation}
For the rest-work defect
\[
 \mathcal W=\dot\rho_r+u(\rho_r+p_{z,r})
 +2v(\rho_r+p_{\perp,r})-\operatorname{Tr}_F(J\Pi),
\]
the exact off-constraint identities are
\begin{equation}
 \boxed{R_{\rm geom}=R_*+\frac{\mathcal C}{2\mathscr F},\qquad
 \dot{\mathcal C}=-\left(\theta-\frac{f_1}{2\mathscr F}\right)
 \mathcal C-\mathcal W.}
 \label{p3f:eq:native-spherical-constraint-propagation}
\end{equation}
The coordinate-only finite source has \(\mathcal W=0\), so compatible
\(\mathcal C(t_0)=0\) is preserved. On an open licensed fixed-rank chart
with \(A,B,\mathscr F>0\), every compatible initial packet has a unique
local solution in the chosen regular gauge. Its full density-matrix
spectrum, positivity and trace are preserved.
\end{theorem}
\begin{proof}
The scalar equation and the two spatial metric equations give the displayed
normal form. Direct contraction gives
\(R_{\rm geom}=2X+4Y+4uv+2(v^2+B^{-2})\).
Differentiate \(\mathcal C\), using the scalar equation and the residual
work identity; this yields both off-constraint formulas. The coefficient
\(f_1/(2\mathscr F)\) is required because \(R_*\) substitutes the trace
before imposing the lapse constraint. Paper IV,
Proposition~\ref{p4:prop:coordinate-source-work}, gives \(\mathcal W=0\)
exactly for the retained coordinate-only source. Thus the geometrical
curvature equals \(R_*\) on the constraint and the original equations hold.
All coefficients are locally Lipschitz on the stated open finite chart;
the local ODE existence and uniqueness theorem applies. The state equation
has a unitary propagator along that solution, proving the state assertions.
All other consumed symmetry/Gauss constraints belong to the proved compatible
restriction, rather than being inferred from the lapse identity alone.
\end{proof}
This supplies the local analytic certificate for the existing solution word
on this member. It does not choose a new occurring root, prove a globally
matched region, or restore higher/tidal sectors by naming their parent.

\subsubsection{Canonical matching and the higher-jet handoff}
\begin{corollary}[Regular first-jet canonical matching]
\label{p3f:cor:native-spherical-canonical-matching}
On a noncharacteristic first-jet matching surface with no independent surface
action, retain the induced fields and full canonical boundary data. The
native scalar--metric contributions are
\[
 \pi^{ij}=\sqrt h[\mathscr F(K^{ij}-Kh^{ij})-h^{ij}D_\perp\mathscr F],
 \qquad p_\Phi=2\alpha_\Phi\sqrt h(D_\perp\Phi+K\Phi/3),
\]
and
\[
 \pi-\tfrac12\operatorname{Tr}_F(\Phi p_\Phi)
 =-2\kappa_J\sqrt h K.
\]
For the coordinate-only/affine-connection source member, matching the full
state and frame/boundary jets makes its canonical contribution
\(-\operatorname{Tr}(\sigma H_{,v})/\hbar\) agree. Equality of the complete
momenta then matches \(K_{ij}\) and \(D_\perp\Phi\) for
\(\mathscr F,\kappa_J>0\). A separately owned surface term contributes its
actual jump; characteristic/rank-changing surfaces retain their own interface.
\end{corollary}
\begin{proof}
The trace combination fixes \(K\). The scalar momentum then fixes
\(D_\perp\Phi\) and hence \(D_\perp\mathscr F\); the traceless metric
momentum fixes the traceless \(K_{ij}\) through \(\mathscr F\ne0\).
Subtract only the already matched complete source contribution. A source
with additional velocity jets instead uses its complete Legendre matrix.
\end{proof}

For a specified finite order-\(r\) local member, the highest-order Euler
block on its actual constrained support has
\[
 \mathsf M^{(r)}_{ij}
 =L_{{\rm cl},q_i^{(r)}q_j^{(r)}}
 -\operatorname{Tr}(\sigma H_{,q_i^{(r)}q_j^{(r)}}),\qquad
 (-1)^r\mathsf M^{(r)}q^{(2r)}.
\]
Here \(L_{\rm cl},H\) use common physical units and the source order is
at most \(r\); derivatives of the state occur below order \(2r\).
Mixed-order fields retain their lower-order equations as well. An invertible
actual principal block permits retaining the lower jets as a finite normal
form, after the complete higher-order constraints have been discharged.
A genuinely nonlocal member consumes its retarded problem instead.

\begin{proposition}[Spherical quadratic metric principal block]
\label{p3f:prop:spherical-quadratic-principal-block}
On the common local real-form domain of
Theorem~\ref{p3f:thm:finite-local-bosonic-clock-real-form}, take the inherited
constant-coefficient metric sector
\(S_{(4)}=\int\sqrt{-g}(aR^2+b\operatorname{Ric}^2+c\operatorname{Rm}^2)\).
In native action units and proper-time gauge its acceleration Hessian is
\begin{gather}
 \boxed{M^{(2)}_{g,(4)}=\Omega_z AB^2
 \begin{pmatrix}
 (8a+4b+8c)/A^2&(16a+4b)/(AB)\\
 (16a+4b)/(AB)&(32a+12b+16c)/B^2
 \end{pmatrix}}\notag\\
 \boxed{\det M^{(2)}_{g,(4)}=32\Omega_z^2B^2(b+4c)(3a+b+c).}
 \label{p3f:eq:spherical-quadratic-principal-block}
\end{gather}
The factors are the scalar and Weyl channels:
\[
\begin{gathered}
 aR^2+b\operatorname{Ric}^2+c\operatorname{Rm}^2=q_sR^2+q_wC^2+q_GG_4,\\
 q_s=a+(b+c)/3,\quad q_w=(b+4c)/2,\quad q_G=-b/2-c.
\end{gathered}
\]
Thus the determinant is \(192\Omega_z^2B^2q_sq_w\).
\end{proposition}
\begin{proof}
With \(X=\ddot A/A\), \(Y=\ddot B/B\), the homogeneous quadratic
acceleration parts of the curvature contractions are respectively
\((2X+4Y)^2\), \(2X^2+4XY+6Y^2\), and \(4X^2+8Y^2\).
Differentiate the reduced density \(\Omega_z AB^2\) times their combination
twice with respect to \(\ddot A,\ddot B\). The displayed invariant
decomposition and determinant follow by expansion.
\end{proof}
The \(R^2\) and \(C^2\) channels individually have metric-acceleration
rank one, and \(G_4\) rank zero; its boundary contribution remains retained.
These are raw principal blocks, with lapse/gauge variation still preceding
gauge fixing and any acceleration-dependent source correction retained.
Critical factors require the full relation, not a generic inverse. No
numerical coefficient or higher/nonlocal strong-field domain is selected.


\subsection{Characteristic geometry of the complete finite source action}
\label{p3f:sec:v2190-native-clock-evolution}

The complete source momenta and Cauchy matrix of
Proposition~\ref{p3f:prop:complete-source-cauchy-matrix} already retain
backreaction, velocity-dependent source insertions and the full density
matrix.  The result below gives those equations one geometric phase-action
presentation.  It does not add a quantum--classical model, a new state
carrier, a vacuum allocation or a second physical clock.  The existing
licensed action/source partition and its common physical action units
control every term.  The complete generator comes from
Theorem~\ref{p3f:thm:v1148-finite-metric-constructor}, including its actual
interaction, ordering, central and moving-frame terms.

\begin{lemma}[Inherited finite-state orbit form]
\label{p3f:lem:v2190-source-orbit-form}
For the actual finite source state \(\sigma_0\), let
\(\mathcal O_{\sigma_0}=\{U\sigma_0U^\dagger:U\in U(N)\}\)
be its existing unitary-orbit presentation.  Its spectrum is retained from
the state, not selected by this construction.  The full matrix orbit is
used as a mathematical presentation, not as a license for every unitary
or state to be physically admissible.  All existing grading, support and
constraint marks remain on the actual equation graph.  For Hermitian \(A,B\), set
\[
 \delta_A\sigma=-\frac{i}{\hbar}[A,\sigma],\qquad
 \omega_\sigma(\delta_A\sigma,\delta_B\sigma)
      =-\frac{i}{\hbar}\operatorname{Tr}(\sigma[A,B]).
\]
This is a real, closed, nondegenerate form on the orbit, including the
zero-dimensional orbit of a central state.  On a local unitary section it
has the primitive
\[
 \vartheta_{\mathcal O}
   =i\hbar\operatorname{Tr}(\sigma_0U^\dagger dU),\qquad
 \omega=-d\vartheta_{\mathcal O}.
\]
It has the normalization of the canonical finite-source kinetic form on a
unitary frame; there is no additional action weight.  It re-presents the
existing source equation rather than postulating a new state action.
For an energy functional
\(f_A(\sigma)=\operatorname{Tr}(\sigma A)\), its Hamiltonian field is
\(\delta_A\sigma\) in the convention \(\iota_{X_f}\omega=df\).
\end{lemma}
\begin{proof}
Adding to \(A\) an operator commuting with \(\sigma\) changes neither the
tangent nor the trace pairing, so the form is well defined.  If the pairing
vanishes for all Hermitian \(B\), nondegeneracy of the ordinary matrix trace
pairing forces \([A,\sigma]=0\).  The form is real because the commutator
is anti-Hermitian.  On \(U(N)\), the Maurer--Cartan identity for
\(M=U^\dagger dU\) gives \(-d\vartheta_{\mathcal O}
=i\hbar\operatorname{Tr}(\sigma_0M\wedge M)\), which is the displayed
pairing and is basic for the stabilizer quotient.  It is therefore closed.
Only local primitives are asserted on the orbit.

The normalized source kinetic expression
\((i\hbar/2)\operatorname{Tr}[\sigma_0(U^\dagger dU-dU^\dagger U)]\)
equals \(i\hbar\operatorname{Tr}(\sigma_0U^\dagger dU)\) by unitarity.
This also fixes the factor.  Finally cyclicity gives
\(df_A(\delta_B\sigma)=-(i/\hbar)\operatorname{Tr}(\sigma[A,B])\).
The orbit is a re-presentation of the existing state equation, not an
independent quantization or an assumption of global orbit coordinates.
\end{proof}

\begin{theorem}[Joint source--geometry characteristic flow]
\label{p3f:thm:v2190-joint-characteristic-flow}
Use the complete licensed first-jet source action of
Proposition~\ref{p3f:prop:complete-source-cauchy-matrix}, after its actual
regular constraint/gauge treatment.  Let \(q,v\) be the retained classical
coordinates and velocities, and \(H(q,v,t)=H^\dagger\) the complete finite
source generator in the same action units as \(L_{\rm cl}(q,v,t)\).
Assume the coefficients are \(C^3\) locally and
\[
 P_i=L_{{\rm cl},v_i}-\operatorname{Tr}(\sigma H_{,v_i}),\qquad
 \mathsf M_{ij}=L_{{\rm cl},v_iv_j}
                 -\operatorname{Tr}(\sigma H_{,v_iv_j})
\]
is invertible on the claimed chart.  Write
\(v=v(q,P,\sigma,t)\) for its unique local inverse.  On
the local matrix presentation \(T^*Q\times\mathcal O_{\sigma_0}\), put
\begin{align*}
 \mathcal E(q,P,\sigma,t)
 &=P_iv_i-L_{\rm cl}(q,v,t)+\operatorname{Tr}(\sigma H(q,v,t)),\\
 \vartheta_{\rm tot}&=P_i\,dq^i+\vartheta_{\mathcal O},\qquad
 \Omega_{\rm tot}=dq^i\wedge dP_i+\omega,\qquad
 \alpha_{\rm tot}=\vartheta_{\rm tot}-\mathcal E\,dt.
\end{align*}
The following is the exact envelope differential, with \(v\) the same
retained inverse in every term:
\[
 \boxed{d\mathcal E
 =v_i\,dP_i+(-L_{{\rm cl},q_i}+\operatorname{Tr}(\sigma H_{,q_i}))\,dq^i
 +\operatorname{Tr}(H\,d\sigma)
 +(-L_{{\rm cl},t}+\operatorname{Tr}(\sigma H_{,t}))\,dt.}
\]
Consequently the clock-normalized characteristic field of
\(-d\alpha_{\rm tot}\) is exactly
\[
 \boxed{\dot q=v,\qquad
 \dot P_i=L_{{\rm cl},q_i}-\operatorname{Tr}(\sigma H_{,q_i}),\qquad
 \dot\sigma=-\frac{i}{\hbar}[H(q,v,t),\sigma].}
\]
It equals the complete source-corrected Cauchy system, including its
commutator force.  This identification is derived within the existing
complete action/evolution target; no second \(\mathsf{DYN}\) generator
selector is required for that same word.  Any further restriction to a
physical support or constraint branch retains its proved invariance or
the actual transported-projector equation; the full matrix orbit is not
promoted to an independent physical carrier.
\end{theorem}
\begin{proof}
The implicit-function theorem applies to the momentum map because its
velocity derivative is \(\mathsf M\).  Differentiate \(\mathcal E\)
before substituting the inverse.  All \(dv\) terms cancel by
\(P=L_{{\rm cl},v}-\operatorname{Tr}(\sigma H_{,v})\), proving the envelope
formula.  In particular its derivative in a state-orbit direction is
\(\operatorname{Tr}(H\,d\sigma)\), not the derivative obtained by
incorrectly holding \(v\) fixed in the energy expression
\(v_iL_{{\rm cl},v_i}-L_{\rm cl}
 +\operatorname{Tr}[\sigma(H-v_iH_{,v_i})]\).
Lemma~\ref{p3f:lem:v2190-source-orbit-form} and the canonical cotangent
form therefore give exactly the three displayed equations.  The phase
action is the regular Legendre presentation of the same retained action;
Paper II-A Theorem~\ref{p2a:thm:v2190-clocked-characteristic} identifies
its characteristic and variational relations.

Differentiate the momentum equation and use
\[
 \frac{d}{dt}\operatorname{Tr}(\sigma X)
  =\operatorname{Tr}\!\left[\sigma\left(\dot X+
                   \frac{i}{\hbar}[H,X]\right)\right].
\]
One obtains \(\mathsf M\dot v=b\), with
\begin{align*}
 b_i={}&L_{{\rm cl},q_i}-L_{{\rm cl},v_iq_j}v_j-L_{{\rm cl},v_it}\\
 &-\operatorname{Tr}\!\left\{\sigma\left[
 H_{,q_i}-H_{,v_iq_j}v_j-H_{,v_it}
                   -\frac{i}{\hbar}[H,H_{,v_i}]\right]\right\},
\end{align*}
which is the existing complete Cauchy formula with its time-dependent
terms.  Thus the identification is an equality of complete equations,
not a declaration that an unrelated kinematic clock follows them.
\end{proof}

\begin{corollary}[Local nonlinear graph, source spectrum and action work]
\label{p3f:cor:v2190-local-interacting-graph}
On the regular chart above, each retained initial tuple gives a unique
local joint solution.  Its finite source spectrum, positivity and trace
are preserved, and
\[
 \frac{d\mathcal E}{dt}
   =-L_{{\rm cl},t}+\operatorname{Tr}(\sigma H_{,t}).
\]
Thus an autonomous closed member conserves the complete action energy.
Every existing field, interaction, central and moving-frame contribution
remains in \(H\) and its jets.  The initial state is not required to be
stationary or Gibbs.  The source-normalization metric being positive does
not by itself prove invertibility of \(\mathsf M\); critical branches
retain the complete relation and its constraints.
\end{corollary}
\begin{proof}
The characteristic field is locally Lipschitz on the regular chart, so the
local Cauchy theorem of Paper II-A applies.  Along the solution, the source
equation has the unitary propagator of
Theorem~\ref{p4:thm:v2190-complete-finite-propagator}, preserving its spectrum.
Antisymmetry of \(\Omega_{\rm tot}\) and the envelope time derivative give
the energy-work identity.  No energy coercivity on a gravitational
constraint surface or unrestricted continuum well-posedness follows.
\end{proof}

\begin{corollary}[Complete finite gauge moment and constraint propagation]
\label{p3f:cor:v2190-complete-gauge-constraint}
Suppose the actual finite word carries an already licensed joint gauge
action, with configuration infinitesimal \(Y_\xi\), its cotangent lift,
and a fixed anti-Hermitian source representation \(\rho(\xi)\).  Suppose
the complete action is invariant under that same action.  The moment is
\[
 \mu^\xi=P(Y_\xi)+\operatorname{Tr}(\sigma\,i\hbar\rho(\xi)).
\]
For the gauge-invariant complete energy \(\mathcal E\),
\(\{\mu^\xi,\mathcal E\}=0\).  If the actual equivariant moment components
satisfy \(\{\mu_a,\mu_b\}=c_{ab}{}^c\mu_c\), the total generator
\(\mathcal E+\lambda^a\mu_a\) gives
\(\dot\mu_a=\lambda^b c_{ab}{}^c\mu_c\).  Its zero constraint surface
is therefore preserved.  The source charge is retained, not replaced by
the pure-gauge Gauss moment.
\end{corollary}
\begin{proof}
Cotangent-lift preservation of \(\vartheta_{\rm can}\) gives moment
\(P(Y_\xi)\).  The orbit lemma applied to
\(A=i\hbar\rho(\xi)\) gives the source tangent
\([\rho(\xi),\sigma]\) and the displayed source moment.  Their sum
generates the joint action.  Invariance of the complete action descends
through its equivariant regular Legendre transform to invariance of
\(\mathcal E\).  Hamilton's equation for \(\mu\) and the retained moment
algebra give the result.  No gauge action on an arbitrary truncated band
is inferred.  Moving representation/projector charts require their actual
transport and Ward jets; the local general covariance application remains
the existing same-owner certificate where consumed.
\end{proof}

The local bosonic scalar and gauge Euler equations are already derived in
Theorem~\ref{p3f:thm:v176-local-bosonic-euler-noether}.  A restricted scalar
variation space consumes its projected Euler covector; other retained
source sectors contribute their actual variations.  The present result
assembles their complete finite source/backreaction evolution and its
clocked characteristic geometry.  An actual separately retained operation
can be compared by Theorem~\ref{p2a:thm:v2190-clock-evolution-closure}; the
Hopf provenance map alone has no asserted map into this joint phase role.
The conditional polynomial continuation and the finite comparison test
are recorded in Paper~III-G, Appendix~\ref{p3f:app:v2190-native-continuation}.


\subsection{The causal Ricci-history map and complete physical Hopf response}
\label{p3f:sec:v2216-causal-hopf-map}

The round homogeneous geometry supplies a complete map to its Hopf mode
operator.  It is a map of metric or Ricci histories with the inherited
Cauchy data, rather than an instantaneous prescription $a^{-2}=cR$.
The construction below uses an actually retained scalar action and its
source domain.  It selects neither a new field nor a state, clock, Gaussian
measure or nonlinear response action.

\begin{proposition}[Independent lapse and scale contacts of curvature]
\label{p3f:prop:v2216-curvature-contacts}
For $ds^2=-N^2dt^2+a^2d\Omega_3^2$, put
$D=N^{-1}d/dt$, $H=D\log a$.  For independent logarithmic variations
$s=\delta\log a$, $n=\delta\log N$,
\begin{align}
 R&=6(DH+2H^2+a^{-2}),\nonumber\\
 \delta R&=6[D^2s+4HDs-HDn-2n(DH+2H^2)-2a^{-2}s].
 \label{p3f:eq:v2216-curvature-pullback}
\end{align}
For the full mixed contact write $\lambda=\log a$, $\nu=\log N$,
\[
 B=\ddot\lambda+2\dot\lambda^2-\dot\lambda\dot\nu,\quad
 B_i=\ddot s_i+(4\dot\lambda-\dot\nu)\dot s_i
                  -\dot\lambda\dot n_i,\quad
 B_{12}=4\dot s_1\dot s_2-\dot s_1\dot n_2-\dot s_2\dot n_1.
\]
For commuting additive variations of these logarithmic coordinates,
\begin{equation}
 \frac{\delta_{12}R}{6}
 =N^{-2}(B_{12}-2n_1B_2-2n_2B_1+4n_1n_2B)
                              +4a^{-2}s_1s_2.
 \label{p3f:eq:v2216-ricci-second-contact}
\end{equation}
\end{proposition}
\begin{proof}
$\delta D=-nD$, $\delta H=Ds-nH$, and
$\delta(DH)=D^2s-HDn-2nDH$ give
\eqref{p3f:eq:v2216-curvature-pullback}.  The exact coordinate expression
$R/6=e^{-2\nu}B+e^{-2\lambda}$ gives
\eqref{p3f:eq:v2216-ricci-second-contact} by two product rules.
In particular fixing a scalar curvature value does not fix a spatial
radius: at $N=1,H=0$, any positive $a$ and
$DH=R/6-a^{-2}$ realize that value as a smooth local metric jet.
\end{proof}

\begin{theorem}[Fixed-Cauchy causal Ricci inverse]
\label{p3f:thm:v2216-causal-ricci-inverse}
After retaining the independent lapse equation, use conformal time $\eta$.
For supplied $a_i>0,a_i'$ the geometric map
$\mathcal R[a]=6(a''a^{-3}+a^{-2})$ has the inverse initial-value problem
\begin{equation}
 a''=\frac{R}{6}a^3-a,\qquad a(\eta_i)=a_i,\quad a'(\eta_i)=a_i'.
 \label{p3f:eq:v2216-ricci-inverse}
\end{equation}
For every continuous $R$ this has a unique local positive solution.
On any fixed regular finite interval on which that solution remains
positive and bounded, it depends smoothly on $R$.  For fixed initial data,
\begin{align}
 \alpha_r:=D_Ra[r]&=(J_a^{-1})_{\mathrm{ret}}r,\quad
 J_a\alpha=6a^{-3}[\alpha''-(3a''/a+2)\alpha],\nonumber\\
 \gamma_{rs}:=D_R^2a[r,s]
 &=-(J_a^{-1})_{\mathrm{ret}}\mathcal B_a[\alpha_r,\alpha_s],
 \label{p3f:eq:v2216-ricci-inverse-jets}\\
 \mathcal B_a[\alpha,\beta]
 &=6a^{-4}[-3(\alpha\beta''+\beta\alpha'')
                         +(12a''/a+6)\alpha\beta].\nonumber
\end{align}
Both inverse contacts have zero Cauchy data at $\eta_i$.
\end{theorem}
\begin{proof}
Equation \eqref{p3f:eq:v2216-ricci-inverse} is the defining curvature
identity solved for $a''$.  Its right side is continuous in time and
smooth in $a$, so local existence, uniqueness and positivity follow
from the initial-value theorem and $a_i>0$.  On a regular compact
interval, differentiable dependence, or the initial-value implicit
function theorem from $C^2$ to $C^0$, gives smooth dependence.
Differentiating $\mathcal R$ gives $J_a$ and $\mathcal B_a$.
The first inverse is equivalently the unique solution of
\[
 \alpha_r''-[(R/2)a^2-1]\alpha_r=(a^3/6)r,
 \qquad \alpha_r(\eta_i)=\alpha_r'(\eta_i)=0.
\]
This defines its retarded Green operator without a spectral positivity
assumption.  Twice differentiating $\mathcal R[a[R]]=R$ yields
$J_a\gamma_{rs}+\mathcal B_a[\alpha_r,\alpha_s]=0$, proving
the second contact.
\end{proof}

On the actual root-generated branch the supplied Cauchy values are those
of that root.  If they vary in a comparison, their true first and second
derivatives add the corresponding initial-data solutions.  The conformal
chart is imposed after lapse variation, so the inverse does not remove
the Hamiltonian constraint.  Only histories obeying the actual source
and constraint equations belong to the physical solution branch.

\begin{theorem}[Canonical four-dimensional physical Hopf mode]
\label{p3f:thm:v2216-physical-hopf-mode}
Take an actually retained real scalar constituent
\[
 S_\varphi=-\frac12\int\sqrt{-g}
     [(\nabla\varphi)^2+(m^2+\xi R)\varphi^2]\,d^4x
\]
with its inherited coefficients and allowed support.  On a charge-one
Hopf restriction the full unit-radius spatial eigenvalue is $4j^2-1$,
including the vertical unit of
Paper~III-C, Theorem~\ref{p3c:thm:v2216-round-hopf-completion}.
For $e=N/a$ and $u_j=a\varphi_j$, its action is, with transported endpoint,
\begin{equation}
 S_j=\frac12\int dt[e^{-1}\dot u_j^2-e\Omega_j^2u_j^2],\qquad
 \Omega_j^2=4j^2+a^2m^2+(\xi-1/6)a^2R.
 \label{p3f:eq:v2216-physical-hopf-frequency}
\end{equation}
For $d=\xi-1/6$, fixed $m^2,\xi$, its Ricci-history contacts are
\begin{align}
 D_R\Omega_j^2[r]&=2a(m^2+dR)\alpha_r+da^2r,\nonumber\\
 D_R^2\Omega_j^2[r,s]
 &=2(m^2+dR)(\alpha_r\alpha_s+a\gamma_{rs})
                          +2da(\alpha_r s+\alpha_s r).
 \label{p3f:eq:v2216-hopf-frequency-jets}
\end{align}
Actual variations of the mass map, coupling, projector or source frame
retain their additional contacts.
\end{theorem}
\begin{proof}
The covariant mode action has kinetic term $a^2\dot\varphi_j^2/(2e)$
and potential $-ea^2(4j^2-1+a^2m^2+\xi a^2R)\varphi_j^2/2$.
Substitute $\varphi_j=u_j/a$.  Integrating its cross term gives the
endpoint $-[\dot a\,u_j^2/(2ea)]$ and the bulk pump
$+e(D_e^2a/a)u_j^2/2$, $D_e=e^{-1}d/dt$.
The lapse-unfixed identity $a^2R=6(D_e^2a/a+1)$ proves
\eqref{p3f:eq:v2216-physical-hopf-frequency}.
The endpoint vanishes only under the actual endpoint rule; otherwise
it is transported.  Differentiating this frequency and inserting
\eqref{p3f:eq:v2216-ricci-inverse-jets} proves the two contacts.
\end{proof}

The native odd-field zero-background scalar has the already derived
$\xi=1/6$, $m^2=-2\Lambda_g^2$, on its admitted canonical scalar
support.  Thus $\Omega_j^2=4j^2-2\Lambda_g^2a^2$ and its nonzero
history vertex is $-4\Lambda_g^2a\alpha_r$.  Conformal coupling does
not make its negative potential Hessian vanish or guarantee positivity.
For the separate massless conformal specialization, $\Omega_j^2=4j^2$
and all these frequency vertices vanish.  Its phase action still
contains $p_j\dot u_j-(N/a)H_{j,\mathrm{conf}}$, so a retained isotropic
state has proper energy $\langle H_{j,\mathrm{conf}}\rangle/a$ and
radiation pressure.  Zero frequency vertex does not imply zero metric
work, nor does it discard a licensed quantum measure or anomaly.

\begin{proposition}[Complete action and causal-source composition]
\label{p3f:prop:v2216-source-chain-rule}
Let $\widehat\Gamma[a,R]$ denote an actually licensed source functional
in the conformal chart, including its true field/state transport,
normalization, projector and endpoint terms.  Its pullback
$\Gamma[R]=\widehat\Gamma[a[R],R]$ obeys
\begin{align}
 D\Gamma[r]&=D\widehat\Gamma[(\alpha_r,r)],\nonumber\\
 D^2\Gamma[r,s]
 &=D^2\widehat\Gamma[(\alpha_r,r),(\alpha_s,s)]
                         +D\widehat\Gamma[(\gamma_{rs},0)].
 \label{p3f:eq:v2216-complete-source-chain-rule}
\end{align}
In particular the second inverse contact cannot be omitted merely
because the independent $(a,R)$ frequency is affine in $R$.
For an admitted canonical scalar with fixed initial covariance, let
$G_j^{\mathrm{ret}}$ be its retarded oscillator Green function and
$H_j(\eta,v)=\langle\{u_j(\eta),u_j(v)\}\rangle$ its transported
anticommutator.  Writing $A=(J_a^{-1})_{\mathrm{ret}}$, one has
\begin{align}
 V_{j,R}(v,s)&=da(v)^2\delta(v-s)
              +2a(v)[m^2+dR(v)]A(v,s),\nonumber\\
 \frac{\delta\langle u_j(\eta)^2\rangle}{\delta R(s)}
 &=-\int_{\eta_i}^{\eta}
        G_j^{\mathrm{ret}}(\eta,v)H_j(\eta,v)V_{j,R}(v,s)\,dv.
 \label{p3f:eq:v2216-causal-scalar-composition}
\end{align}
For an admitted finite Dirac source whose actual conformal generator
is $\widehat H_\eta(a)$, put $B=\partial_a\widehat H_\eta$.
At fixed initial density its observable response composes as
\begin{align}
 K_{B,R}(\eta,s)&=\int_{\eta_i}^{\eta}K_a(\eta,v)A(v,s)\,dv,
 \nonumber\\
 K_a(\eta,v)&=\langle\partial_aB(\eta)\rangle\delta(\eta-v)
  -i\Theta(\eta-v)\langle[B(\eta),B(v)]\rangle .
 \label{p3f:eq:v2216-causal-dirac-composition}
\end{align}
The inherited central energy contact remains in $\partial_aB$.
These retarded observable responses are not identified with a symmetric
one-history variational Hessian.
\end{proposition}
\begin{proof}
The first two identities are the complete Fr\'echet chain rule for
$R\mapsto(a[R],R)$, including its second derivative.  Varying the
oscillator equation with fixed initial covariance gives
$\delta u_j(\eta)=-\int G_j^{\mathrm{ret}}(\eta,v)
\delta\Omega_j^2(v)u_j(v)\,dv$; the anticommutator gives the stated
variance response.  Insert the proved frequency vertex to obtain
\eqref{p3f:eq:v2216-causal-scalar-composition}.  For the finite Dirac
system, Duhamel variation of its actual unitary propagator gives the
commutator in $K_a$ and direct differentiation of the insertion gives
its contact.  Substitution of $\delta a=A\delta R$ proves
\eqref{p3f:eq:v2216-causal-dirac-composition}.  The triangular supports
of both retarded operators prove causality.  Fixed initial preparation
does not mean repeated instantaneous ground-state resetting.
\end{proof}

\begin{proposition}[The complete two metric-work ports]
\label{p3f:prop:v2216-complete-metric-work}
Temporarily represent an actual source action by independent controls
$(a,e,R)$, and define its complete Euler insertions by
\[
 \delta S_s=\int Ndt[\mathcal A s+\mathcal E\delta\log e
                                      +\mathcal J\delta R].
\]
With $V=2\pi^2a^3$, the physical source ports are
\begin{equation}
 -V\rho_s=\mathcal E+L_n^*\mathcal J,\qquad
 3Vp_s=\mathcal A-\mathcal E+L_s^*\mathcal J,
 \label{p3f:eq:v2216-complete-metric-work}
\end{equation}
where, on the proper-time pairing $Ndt$,
\[
 L_s^*=6(D^2-4HD-4DH-2a^{-2}),\qquad
 L_n^*=6(HD-DH-4H^2).
\]
All moving-endpoint terms are retained instead of discarded.
\end{proposition}
\begin{proof}
$\delta\log e=n-s$ and
\eqref{p3f:eq:v2216-curvature-pullback} give
$\delta R=L_ss+L_nn$, with
$L_s=6(D^2+4HD-2a^{-2})$ and
$L_n=-6[HD+2(DH+2H^2)]$.
For compact variations $D^*=-D$ on $Ndt$; integration by parts gives
the displayed adjoints.  Comparing with
$\delta S_s=\int Ndt\,V(-\rho_s n+3p_s s)$ proves
\eqref{p3f:eq:v2216-complete-metric-work}.  A metric-dependent field
chart contributes its field-Euler term and second contacts until those
are removed on their actual solution domain.
\end{proof}

\begin{corollary}[Complete stress tests and exact moving-range execution]
\label{p3f:cor:v2217-source-test-execution}
On a complete finite physical source with
$H(t)=\sum_\alpha f_\alpha(t)H_\alpha$, choose the real Hermitian span
$W_0$ of $I$ and every actual energy, pressure, scalar-force and contact
test consumed by the readout. Suppose their time-dependent insertions lie
in this span, or enlarge it by their actual coefficient matrices. The
iteration $W_{j+1}=W_j+\sum_\alpha i[H_\alpha,W_j]$ stabilizes in at most
$d^2-\dim W_0$ strict increases. On its least invariant space all those
expectations obey the exact closed finite adjoint evolution with inherited
initial tests. No uniqueness of an unconsumed microscopic state is required.

For a differentiable retained isometry $J(t)$, put $Q=I-JJ^*$ and use
the actual covariant frame derivative $D_t$. Evolution from every vector
in $\operatorname{ran}J(t_i)$ stays in $\operatorname{ran}J(t)$ precisely
when
\[
 Q(HJ-i\hbar D_tJ)=0
\]
almost everywhere on the regular interval. The retained generator is
$J^*HJ-i\hbar J^*D_tJ$.
\end{corollary}
\begin{proof}
The increasing spaces are subspaces of the $d^2$-dimensional real
Hermitian space. Stabilization is exactly invariance; every invariant
space containing $W_0$ contains each iterate. Heisenberg's equation in
a basis of the limiting space gives the closed expectation system, as in
Paper~IV, Theorem~\ref{p4:v2171:thm:flow}, now seeded by the complete stress
tests. For the second assertion substitute $\psi=Jx$ into
$i\hbar D_t\psi=H\psi$. Projection with $J^*$ gives the retained generator;
projection with $Q$ gives the residual. Vanishing makes the reconstructed
solution solve the complete equation. Conversely invariance for every
initial retained vector and uniqueness of its invertible retained
propagator force that residual to vanish on the entire retained space.
\end{proof}
If the residual is nonzero, a stationary Schur complement alone does not
remove its coupling, memory or moving-frame contact. Actual stress tests
can be evaluated independently of a different charge-asymmetry readout;
they retain the same complete action, state and time history.


This closes the local fixed-Cauchy geometry-to-source response map on
regular round homogeneous slices.  It does not supply a new physical
action constructor, make a retarded response into an instantaneous
$f_{RR}(R)$, or determine general four-dimensional geometry from one
Ricci history.  Those identifications retain their separate complete
action, source and constraint hypotheses.


\subsection{Exact causal completion of the native spherical source branch}
\label{p3f:sec:v247-native-causal-completion}

The regular local Cauchy construction already exists in
Theorem~\ref{p3f:thm:native-spherical-local-cauchy}. The next result
identifies its complete finite-source insertions with derivatives of
the exact influence parent, including the lapse contacts. It does not
infer a full gravity overlap from an Einstein coefficient.

\begin{theorem}[Native lapse, scale and scalar influence ports with their Ward identity]
\label{p3f:thm:v247-native-influence-ward}
Use exactly the licensed spherical scalar--metric action
\eqref{p3f:eq:native-spherical-action}, its physical multiplier $\hbar$,
and the compatible coordinate-only proper finite source
$h(q)=h(A,B,\phi)$ of
Theorem~\ref{p3f:thm:native-spherical-local-cauchy}, on its actual flat
canonical source frame. Put $q=(A,B,\phi_a)$,
$V_3=\Omega_zAB^2$, and retain lapse $N>0$ before variation.
In the source influence parent of
Theorem~\ref{p3f:thm:v246-finite-source-influence-parent} use the actual
coordinate-time generator
\[
 H[N,q]=Nh(q),\qquad
 i\hbar\dot\sigma=[Nh,\sigma],\qquad
 E=\operatorname{Tr}(\sigma h).
\]
All central, mass-owner and retained interaction terms of $h$ are
included. The exact diagonal difference-coordinate covectors are
\[
 \boxed{F_N:=\Gamma_{N_\Delta}=-E,\qquad
 F_i:=\Gamma_{q_\Delta^i}=-N\operatorname{Tr}(\sigma h_{,i}).}
\]
Their complete action-normalized native ports are therefore
\[
 \boxed{\rho_r=\frac E{\hbar V_3},\quad
 p_{z,r}=-\frac{A\operatorname{Tr}(\sigma h_{,A})}{\hbar V_3},\quad
 p_{\perp,r}=-\frac{B\operatorname{Tr}(\sigma h_{,B})}{2\hbar V_3},\quad
 J_a=\frac{\operatorname{Tr}(\sigma h_{,\phi_a})}{\hbar V_3}.}
\]
For every control history on this source domain, without imposing a
metric equation, these covectors obey the exact reparametrization
Ward identity
\[
 \boxed{\sum_iF_i\dot q^i-N\dot F_N=0.}
\]
Equivalently, with $D_\tau=N^{-1}\partial_t$,
$H_A=D_\tau\log A$, $H_B=D_\tau\log B$,
\[
 \boxed{D_\tau\rho_r+H_A(\rho_r+p_{z,r})
             +2H_B(\rho_r+p_{\perp,r})
       =\sum_aJ_aD_\tau\phi_a.}
\]
For $y^a=(N,q^i)$, let $B_a=H_{,a}$ and $Q_{ab}=H_{,ab}$.
The full retarded source kernel is
\[
 \Gamma_{\Delta a,c b}(t,s)
 =-\langle Q_{ab}(t)\rangle\delta(t-s)
 +\frac i\hbar\Theta(t-s)\langle[B_a(t),B_b(s)]\rangle,
\]
with the actual Heisenberg transport understood in the commutator.
In particular
\[
 B_N=h,\quad B_i=Nh_{,i},\qquad
 Q_{NN}=0,\quad Q_{Ni}=h_{,i},\quad Q_{ij}=Nh_{,ij}.
\]
For every compact control perturbation $(n,u^i)$ its response obeys
the differentiated Ward identity
\[
 \boxed{\sum_i\bigl(\delta F_i\dot q^i+F_i\dot u^i\bigr)
              -n\dot F_N-N\partial_t\delta F_N=0.}
\]
\end{theorem}
\begin{proof}
The first covectors are the difference-coordinate derivative of the
finite trace with $H=Nh$. Comparing their coefficient with
\[
 \delta\Gamma\big|_{\Delta=0}
 =\hbar\int NV_3\left[-\rho_r\frac{\delta N_\Delta}{N}
 +p_{z,r}\frac{\delta A_\Delta}{A}
 +2p_{\perp,r}\frac{\delta B_\Delta}{B}
 -\sum_aJ_a\delta\phi^a_\Delta\right]dt
\]
gives all four ports with their absolute signs and volume factors.
Cyclicity of the full finite trace gives
$\dot E=\sum_i\operatorname{Tr}(\sigma h_{,i})\dot q^i$;
the state commutator contributes
$\operatorname{Tr}([Nh,\sigma]h)=0$. This proves the first Ward
identity, and differentiating $V_3=\Omega_zAB^2$ gives the work law.
It also follows from source invariance under
$\delta q^i=\epsilon\dot q^i$,
$\delta N=\partial_t(\epsilon N)$ with compact $\epsilon$:
$\delta H=\partial_t(\epsilon H)$ is a pure change of clock and
leaves the endpoint propagator unchanged. An integration by parts
gives the displayed Ward identity. Duhamel differentiation gives the
kernel, with $Q_{Ni}=h_{,i}$ before any lapse choice. Differentiating
the already proved Ward identity proves its linearized form. No
commuting approximation, one-body factorization or source reset is
used. A central $c(q)I$ drops out of the commutator but contributes
$c,c_{,i},c_{,ij}$ to the stated covectors and contacts.
\end{proof}

\begin{corollary}[The causal influence equations and the native Cauchy system are equivalent]
\label{p3f:cor:v247-native-causal-cauchy-equivalence}
On the same compatible restriction define
\[
 \mathcal I=\hbar S_{g\Phi}[N_+,q_+]
            -\hbar S_{g\Phi}[N_-,q_-]+\Gamma[N_+,q_+;N_-,q_-],
\]
where $S_{g\Phi}=\int L_{g\Phi}dt$ is precisely the native action
above with its inherited endpoint rule, and the source partition
counts no source contribution twice. Its physical mean equations
$\mathcal I_{\Delta}=0$ are exactly the native metric and scalar
equations with the four source ports in the theorem, including the
lapse equation. After lapse variation and the actual compatible
constraint reduction, $N=1$ gives the already derived native local
Cauchy system and its source state equation. Conversely every solution
of that system with its initial lapse constraint and supplied state
solves these causal mean equations. The lapse constraint propagation
is the existing native propagation identity, now owned by this same
influence functional.
\end{corollary}
\begin{proof}
The first difference-coordinate derivative of the classical
difference is $\hbar\delta S_{g\Phi}$, while the theorem evaluates
the source derivative. Therefore their sum is the original native
Euler system with the source insertion used in the Cauchy
construction, with no missing lapse/scale mixed contact. The source
work identity is precisely its Ward antecedent. The cited Cauchy
theorem then gives the reduced evolution and constraint propagation.
Conversely its construction reconstructs the complete Euler system
on the constraint. Uniqueness of the finite propagator identifies its
state with the state inside the influence functional.
\end{proof}

This is an equivalence on the actual constituted local branch. It
does not supply the separate global cap/source incidence, whole-action
$1/96$ projection or continuum measure. In particular the cap action
is not substituted for $S_{g\Phi}$ in this corollary.

\begin{proposition}[Causal elimination adds no independent memory data and has a quantitative continuation bound]
\label{p3f:prop:v247-native-causal-continuation}
Fix a regular proper-time native Cauchy chart and write its classical
state as $x=(q,v)$, with its smooth source-corrected acceleration
field $\dot x=f(x,\sigma)$. For the coordinate-only source under
consideration the native kinetic determinant is nonzero on
$A,B,\mathscr F,\kappa_J>0$. Suppose the coefficients are sufficiently
smooth to be locally Lipschitz on this chart. The source spectrum,
positivity and trace are preserved. Eliminating $\sigma$ with the
fixed initial state gives a causal history equation for $x$ that is
locally equivalent to this finite Cauchy system; it requires no
additional memory initial data.

More quantitatively, let two solutions remain on a compact regular
chart where
\[
 |f(x,\sigma)-f(\widetilde x,\widetilde\sigma)|
 \le L_f(|x-\widetilde x|+\|\sigma-\widetilde\sigma\|_1),\qquad
 \|h(x)-h(\widetilde x)\|\le L_h|x-\widetilde x|.
\]
Then, for $t\ge t_i$ in that common chart,
\[
 \boxed{|x(t)-\widetilde x(t)|+\|\sigma(t)-\widetilde\sigma(t)\|_1
 \le\bigl(|x_i-\widetilde x_i|+\|\sigma_i-\widetilde\sigma_i\|_1\bigr)
 e^{(L_f+2L_h/\hbar)(t-t_i)}.}
\]
If a solution on a finite maximal time interval keeps $x$ in a compact
subset of the regular chart, it extends past that endpoint. Thus a
finite source density by itself cannot terminate this regular local
evolution. A noncompact geometric/velocity escape, degeneration of the
kinetic chart, or failure of the licensed source/domain remains a
separate possible obstruction.
\end{proposition}
\begin{proof}
Finite unitary conjugation preserves the source spectrum, hence its
positive trace-one orbit is compact. Variation of constants in the
two state equations, using trace-norm invariance under conjugation
and $\|[A,\sigma]\|_1\le2\|A\|\|\sigma\|_1$, gives
\[
 \|\sigma(t)-\widetilde\sigma(t)\|_1
 \le\|\sigma_i-\widetilde\sigma_i\|_1
       +\frac2\hbar\int_{t_i}^t\|h(x(s))-h(\widetilde x(s))\|ds.
\]
Combining this with the integral equation for $x$ and applying
Gronwall gives the displayed bound. With common initial state the
causal source is uniquely $U_x\sigma_iU_x^\dagger$; inserting it or
retaining its finite state equation are inverse descriptions of the
same solutions, which proves the memory-data assertion. On a compact
regular classical chart times the compact source orbit the full
vector field is bounded and locally Lipschitz, so a finite endpoint
has a limit and the local existence theorem extends it.
\end{proof}

\begin{proposition}[State-spectrum bounds for the native curvature source]
\label{p3f:prop:v247-native-curvature-orbit-bounds}
On this same native first-jet branch define the complete Hermitian
trace insertion
\[
 \mathsf T(q)=-h-Ah_{,A}-Bh_{,B}
                       +\sum_a\phi_a h_{,\phi_a},\qquad
 X=\operatorname{Tr}_F\Phi^2.
\]
The actual curvature trace is
\[
 \boxed{R_*=-\frac{2\alpha_\Phi\Lambda_g^2X}{\kappa_J}
            -\frac{\operatorname{Tr}(\sigma\mathsf T)}
                         {2\kappa_J\hbar V_3}.}
\]
Let $p_1\ge\cdots\ge p_d\ge0$ be the inherited, conserved eigenvalues
of $\sigma$, and $t_1(q)\ge\cdots\ge t_d(q)$ the eigenvalues of
$\mathsf T$. Put $R_0=-2\alpha_\Phi\Lambda_g^2X/\kappa_J$. Then
\[
 \boxed{R_0-\frac{\sum_jp_jt_j}{2\kappa_J\hbar V_3}
 \le R_*\le
 R_0-\frac{\sum_jp_jt_{d+1-j}}{2\kappa_J\hbar V_3}.}
\]
These bounds are sharp over the full unitary orbit of the supplied
density spectrum at fixed $q$. They do not assert that a single
actual history reaches either extreme. In particular $R_*$ is
uniformly bounded whenever the actual classical history stays in a
compact regular source chart with $V_3$ bounded away from zero.
\end{proposition}
\begin{proof}
Substitute the exact source ports into
\eqref{p3f:eq:native-source-curvature-trace}:
$T_r+\sum_a\phi_aJ_a=\operatorname{Tr}(\sigma\mathsf T)/(\hbar V_3)$.
This gives $R_*$. In a descending eigenbasis of $\mathsf T$, the sum
of any first $k$ diagonal entries of $\sigma$ is
$\sum_j\alpha_jp_j$, with $0\le\alpha_j\le1$ and
$\sum_j\alpha_j=k$, so it is at most $\sum_{j\le k}p_j$.
Summation by parts against the descending $t_j$ gives
$\operatorname{Tr}(\sigma\mathsf T)\le\sum_jp_jt_j$.
Apply the same argument to $-\mathsf T$ for the lower bound.
Aligned and oppositely aligned density eigenbases attain the two
bounds on its full orbit. The curvature inequalities follow because
$2\kappa_J\hbar V_3>0$. Smooth finite source coefficients and positive
volume on a compact chart give the last assertion.
\end{proof}

This finite native curvature envelope uses the actual source trace
and its transported state. It is not a universal invariant curvature
cap: global compactness, the actual history and any branch transition
must still be proved by their own equations. It supplies an exact
testable bound without replacing the observable by a preferred
standard mechanism.


\section{Complete shared response and causal finite-source execution}
\label{p3f:sec:v2216-complete-response}

The following results strengthen the existing same-action descent on its
actual domain.  They do not identify the raw joint heat, a finite source
germ, and the complete physical effective action.  That identification
still consumes the complete value, slope, source, measure and boundary
matching of Paper~III-D,
Theorem~\ref{p3d:thm:v2216-complete-action-jets}.  In particular its
local-to-global obstruction,
Theorem~\ref{p3d:thm:v2216-exjet-global-obstruction}, is preserved.

\subsubsection{Combining actions before eliminating shared fields}
\begin{theorem}[Shared-interior variance and terminal compatibility]
\label{p3f:thm:v2216-shared-response}
At a common background, let a finite family of actual co-consumed action
constituents have
real quadratic forms
\[
 S_i^{(2)}(u,z)=\tfrac12\langle u,A_i u\rangle+
 \langle u,B_i z\rangle+\tfrac12\langle z,C_i z\rangle,
 \qquad C_i=C_i^*>0.
\]
Here $z$ is the same interior variable in every constituent.  Put
$K_i=A_i-B_iC_i^{-1}B_i^*$, $L_i=C_i^{-1}B_i^*$,
$C_\Sigma=\sum_iC_i$, and
$\bar L=C_\Sigma^{-1}\sum_i C_iL_i$.  Reduction of the sum obeys
\begin{equation}
 \boxed{K_\Sigma=\sum_iK_i+
 \sum_i(L_i-\bar L)^*C_i(L_i-\bar L).}
 \label{p3f:eq:v2216-shared-variance}
\end{equation}
Equality $K_\Sigma=\sum_iK_i$ holds exactly when all $L_i$ agree.
For independent interior variables $z_i$, there is no variance term.
If $K_i(C)\succeq0$ at the one-dimensional terminal curvature port,
then $K_\Sigma(C)=0$ exactly when every $K_i(C)=0$ and all $L_i(C)$
agree.  If the blocks are $C^1$ near $C$ with uniformly positive
interior gaps and the $L_i(C)$ agree, the variance is
$O((R-C)^2)$ and
$K_\Sigma'(C)=\sum_iK_i'(C)$.
\end{theorem}
\begin{proof}
Complete each square:
\[
 S_i^{(2)}=\tfrac12\langle u,K_i u\rangle+
 \tfrac12\|C_i^{1/2}(z+L_i u)\|^2.
\]
The common stationary field is $z_*=-\bar L u$; substitution proves
\eqref{p3f:eq:v2216-shared-variance}.  No commutativity of the $C_i$
is used.  A sum of the displayed nonnegative squared norms vanishes
for all $u$ precisely when $L_i=\bar L$ for all $i$.  With independent
$z_i$ each square vanishes separately.  The terminal assertion follows
by applying the same positivity to every summand.  Under the last
hypotheses, inversion is $C^1$ and uniformly bounded, so
$L_i-\bar L=O(R-C)$, proving the quadratic remainder and derivative
identity.
\end{proof}

This is the positive-interior specialization of the common-action
consumption result in Paper~III-A,
Theorem~\ref{p3a:thm:v2216-common-action-consumption}.  A signed
Lorentzian scalar envelope is not made positive by a change of notation.
For example $S_i=c_i(z-a_i u)^2/2+k_i u^2/2$, $c_i>0$, gives
\[
 K_\Sigma=k_1+k_2+\frac{c_1c_2}{c_1+c_2}(a_1-a_2)^2.
\]
Thus individually null responses need not remain null after sharing.
The native linear terminal slope
$\chi_7(R)=2(R-C)/[C(U+C)]+O((R-C)^2)$ cannot arise solely from
a regular mismatch whose endpoint values agree.  A closing interior gap
or a nonsmooth incidence leaves the theorem's domain and requires its
actual source scaling: for the toy pair
$c(R-C)(z\mp au)^2/2$, a source $-j_z z$ gives
$z_*=j_z/[2c(R-C)]$, so bounded transport requires
$j_z=O(R-C)$.  The toy is not an additional generated action.

\subsubsection{Executed causal kernel and its stationary derivative expansion}
\begin{theorem}[Complete finite-source metric response]
\label{p3f:thm:v2216-source-response}
On the actual fixed coherent frame and complete finite Dirac band, set
$q=\log a$, $V_\gamma=2\pi^2$, and
\[
 H_\eta[q]=d\Gamma\!\left(P_m[D_\gamma+am\beta_D]P_m\right)
       +c(q)I,\qquad c(q)=a\varepsilon_0(a).
\]
Propagate the supplied positive trace-one density by this complete
generator.  With $B=\partial_qH_\eta$, $Q=\partial_q^2H_\eta$,
the force response to $\alpha=\delta q$ is
\begin{align}
 \delta\langle B(\eta)\rangle
 &=\int K_q^{\rm force}(\eta,\xi)\alpha(\xi)\,d\xi,\\
 K_q^{\rm force}(\eta,\xi)
 &=\langle Q(\eta)\rangle\delta(\eta-\xi)
 -i\Theta(\eta-\xi)\operatorname{Tr}
   (\sigma_i[B_H(\eta),B_H(\xi)]).
 \label{p3f:eq:v2216-causal-source-kernel}
\end{align}
The Lorentzian source-action mean is $-\langle B\rangle$, and its
mixed causal action kernel is $-K_q^{\rm force}$.  The actual trace is
$T=-\langle B\rangle/(V_\gamma a^4)$.  With fixed initial scale and
velocity, a curvature perturbation $\rho_R=\delta R$ induces
\[
 \alpha=\mathcal L_q^{-1,\rm ret}\rho_R,\qquad
 \mathcal L_q\alpha=6a^{-2}
 [\alpha''+2(a'/a)\alpha'-2(1+a''/a)\alpha],
\]
and its fully executed trace response is
\begin{equation}
 \delta T=\frac{-\int K_q^{\rm force}(\eta,\xi)\alpha(\xi)d\xi
                   +4\langle B\rangle\alpha}{V_\gamma a^4}.
 \label{p3f:eq:v2216-curvature-source-response}
\end{equation}
This is the pointwise FLRW trace when the complete-multiplet isotropic,
zero-flux embedding is supplied; otherwise it is its averaged homogeneous
port.  Actual moving frames, bands, masses, measures and preparation
insertions are retained whenever they vary.
\end{theorem}
\begin{proof}
Multiplying the proper generator by the conformal lapse gives $H_\eta$.
Duhamel differentiation of its finite unitary propagator gives the
commutator in \eqref{p3f:eq:v2216-causal-source-kernel}; differentiating
the insertion itself gives $Q$.  The two-history identity
$\delta\Gamma_s=-\int\langle\delta H_\eta\rangle$ fixes the action
sign.  The conformal $q$ direction varies both lapse and scale, hence
$\langle B\rangle=aV(\rho-3p)$, rather than the different active
lapse/scale port $-NV(\rho+3p)$.  Differentiate
$R=6e^{-2q}(1+q''+(q')^2)$ to obtain $\mathcal L_q$.  Its nonzero
leading coefficient gives the unique retarded inverse on each regular
$a>0$ interval.  Differentiation of the trace volume gives the last
contact $4\langle B\rangle\alpha$.
\end{proof}

\begin{proposition}[Positive spectral moments, exact nullity and error]
\label{p3f:prop:v2216-source-spectral-moments}
Freeze a regular member with its actual stationary simple ground
$|0\rangle$.  Let $\Delta_r=E_r-E_0\ge\Delta_*>0$ and
$b_r=\langle r|B|0\rangle$.  On the actual finite band,
\[
 K_q^{\rm mem}(\omega)=-2\sum_r
 \frac{\Delta_r|b_r|^2}{\Delta_r^2-(\omega+i0)^2}.
\]
For $|\omega|\le\Omega<\Delta_*$ its expansion through
$\omega^{2N}$ has error at most
\[
 2\sum_r\frac{|b_r|^2}{\Delta_r}
 \frac{(\Omega^2/\Delta_r^2)^{N+1}}{1-\Omega^2/\Delta_r^2}.
\]
The source-action fourth-derivative coefficient is
\begin{equation}
 \boxed{\Sigma_{q,4}=2\sum_r|b_r|^2/\Delta_r^5\ge0.}
 \label{p3f:eq:v2216-positive-source-fourth}
\end{equation}
It vanishes exactly when the insertion has no transition from the
populated ground support to positive-gap states.  Independent sectors
in their product ground add positive weights; opposite fermion/boson
ground-energy signs do not cancel them.  A central instantaneous
allocation changes the mean and contact, but contributes no
$\omega^4$ memory coefficient.
\end{proposition}
\begin{proof}
An energy resolution of the retarded commutator gives
$-2\Theta(s)\sum_r|b_r|^2\sin(\Delta_rs)$ and the displayed transform.
The geometric-series remainder of each denominator gives the estimate.
Reverse the force sign for the action coefficient.  All denominators
and weights in \eqref{p3f:eq:v2216-positive-source-fourth} are positive,
giving the null test.  On a product ground each sector excites a
different tensor factor, so squared matrix elements add, even for equal
gaps.  Central matrices commute and supply only a delta contact.
\end{proof}

For the native paired block
$h_\eta=\ell\sigma_3+am\sigma_1$,
$\epsilon^2=\ell^2+a^2m^2$, the ground transition has
$|b|^2=a^2m^2\ell^2/\epsilon^2$ and gap $2\epsilon$.  A canonically
real stable conformal scalar has
$\Omega_n^2=n^2+a^2m_{s,n}^2>0$, insertion
$a^2m_{s,n}^2u^2$, matrix element squared
$a^4m_{s,n}^4/(2\Omega_n^2)$ and gap $2\Omega_n$.  Thus, with every
actual multiplicity counted once,
\begin{equation}
 \Sigma_{4,D}=\sum_\ell d_\ell
 \frac{a^2m_\ell^2\ell^2}{16\epsilon_\ell^7},\qquad
 \Sigma_{4,s}=\sum_n d_n\frac{a^4m_{s,n}^4}{32\Omega_n^7}.
 \label{p3f:eq:v2216-native-source-moments}
\end{equation}
Massless conformal blocks have exactly zero insertion with nonzero
finite transition gap.  If a covariant local $\nu f(R)$ matching has
additionally been proved in the same density, clock, lapse, source and
boundary units, its homogeneous highest-derivative coefficient is
$36\nu V_\gamma f_{RR}$.  The necessary induced coefficients are then
\[
 \chi_D^{(4)}=\frac1{576\nu V_\gamma}
  \sum_\ell d_\ell\frac{a^2m_\ell^2\ell^2}{\epsilon_\ell^7},
 \qquad
 \chi_s^{(4)}=\frac1{1152\nu V_\gamma}
  \sum_n d_n\frac{a^4m_{s,n}^4}{\Omega_n^7}.
\]
Indeed $\delta R=6a^{-2}(\alpha''-2\alpha)$ on a frozen member,
so the $f_{RR}(\delta R)^2$ Hessian has that coefficient.  This
one homogeneous coefficient does not establish the covariant matching.

\begin{corollary}[Full-pole obstruction and explicit-matter interpretation]
\label{p3f:cor:v2216-source-full-poles}
If the direct gravitational terms and central contacts are finite-order
local metric functionals, a finite stationary ground source with a
nonzero transition weight cannot be absorbed exactly into an affine
Einstein action on an open frequency interval.  Keeping matter explicit
is consistent with $\chi_g(C)=0$ and a nonzero matter response.  If it
is legitimately integrated out, its nonlocal polarization remains; it
does not by itself restore an independent gravitational Legendre field.
\end{corollary}
\begin{proof}
The local frozen kernels are polynomials.  The source kernel has poles
at $\omega=\pm\Delta_r$ with nonzero residues.  An identity on an open
subgap interval extends analytically and would force every residue to
vanish; repeated-gap ground weights add positively.  This is impossible
under the stated hypothesis.  With explicit matter the two operators
occupy different blocks of the coupled equations.  Eliminating that
block retains its rational kernel, rather than a new bare gravitational
coordinate.  A same-action shared reduction can change the conclusion
only by its actual operator and source matching identity.
\end{proof}

For a claimed locally matched total terminal Hessian the exact necessary
fourth-order cancellation is
\[
 \Sigma_{{\rm direct},4}+\Sigma_{{\rm source},4}
 +\Sigma_{{\rm measure},4}-\Sigma_{{\rm elimination},4}=0,
\]
with every term allocated once.  The full shared condition is
$K=A-BC_x^{-1}D$ together with
$j_{\rm red}=j_g-BC_x^{-1}j_x$ and the endpoint contacts.
A gapped $C_x$ can coexist with $K=0$; a gap alone implies no cancellation.
Nor does a matter gap make the curvature map local: on a static closed
member
\[
 \mathcal L_q^{-1,\rm ret}
 =\frac{a_0^2}{6}(D_\eta^2-2)^{-1}_{\rm ret},\qquad
 (D_\eta^2-2)^{-1}_{\rm ret}(s)
 =\Theta(s)\frac{\sinh(\sqrt2s)}{\sqrt2}.
\]
The growing causal mode is not a stationary Fourier approximation.
In a late constant-$H$ expanding flat limit the proper-time map is
$6D_t(D_t+4H)$ and its inverse has a zero-frequency pole.  Constraints
and actual source incidence must be reduced before calling such an
inverse a local susceptibility.  A finite-time derivative expansion
also retains its switch-on transient.  QROOT selects the ground only
on its licensed fresh-root family and then transports it; it supplies
no repeated instantaneous vacuum reset or independent absolute center.

\subsubsection{Cold energy with bounded variation instead of whole-state convergence}
\begin{theorem}[Finite-source cold tail under bounded scale variation]
\label{p3f:thm:v2216-cold-tail-bv}
Let the complete finite-dimensional source in its fixed coherent frame
have
\[
 H(t)=H_m+s(t)\widehat D+\varepsilon(t)I,
 \quad H_m=H_m^*,\quad\widehat D=\widehat D^*,\quad s=a^{-1}>0.
\]
Suppose $s,\varepsilon$ are locally absolutely continuous,
$s\to0$, $\varepsilon\to\varepsilon_\infty$, and
$\operatorname{Var}_{[t_i,\infty)}s=\int_{t_i}^\infty|\dot s|dt<\infty$.
For $i\dot\sigma=[H,\sigma]$, $\sigma\ge0$, $\operatorname{Tr}\sigma=1$,
the actual energy $E=\operatorname{Tr}(\sigma H)$ has a finite limit and
\begin{align}
 \frac d{dt}(E-\varepsilon)&=\dot s\operatorname{Tr}(\sigma\widehat D),\\
 E_\infty-\varepsilon_\infty
 &=\operatorname{Tr}(\sigma_i(H_m+s_i\widehat D))
  +\int_{t_i}^\infty\dot s\operatorname{Tr}(\sigma\widehat D)dt,\\
 |(E_\infty-\varepsilon_\infty)-(E-\varepsilon)|
 &\le\|\widehat D\|\operatorname{Var}_{[t,\infty)}s.
 \label{p3f:eq:v2216-cold-tail-error}
\end{align}
For nondecreasing $a\to\infty$ this gives
$|E_\infty-E|\le\|\widehat D\|/a+|\varepsilon-\varepsilon_\infty|$;
no condition $\int a^{-1}dt<\infty$ is needed.

If additionally $\varepsilon=\varepsilon_0(a)$ is $C^1$ on the actual
scale range and $a\varepsilon_0'(a)\to0$, the inherited metric ports obey
\[
 \rho_s=E/(V_\gamma a^3),\qquad
 3Vp_s=s\operatorname{Tr}(\sigma\widehat D)-a\varepsilon_0'(a)\to0.
\]
If the physically allocated $E_\infty>0$, then
$a^3\rho_s\to E_\infty/V_\gamma$ and $p_s/\rho_s\to0$.
\end{theorem}
\begin{proof}
Finite Hermitian evolution preserves the density spectrum, positivity
and trace, hence $|\operatorname{Tr}(\sigma\widehat D)|\le\|\widehat D\|$.
Cyclicity kills the state derivative in
$d\operatorname{Tr}(\sigma(H-\varepsilon I))/dt$.  Integrability of
$|\dot s|$ proves the energy limit and its tail bound.  For monotone
$s\downarrow0$ the remaining variation equals $s$.  Finally
$\partial_aH=-a^{-2}\widehat D+\varepsilon_0'I$ gives the pressure
identity.  For positive $E_\infty$, eventually $E>E_\infty/2$ and
$p_s/\rho_s=Vp_s/E\to0$.
\end{proof}

\begin{corollary}[Finite-epoch error and positivity certificate]
\label{p3f:cor:v2216-cold-tail-finite-errors}
Put
\[
\begin{aligned}
 b_E&=\|\widehat D\|\operatorname{Var}_{[t,\infty)}s+|\varepsilon-\varepsilon_\infty|,\\
 b_p&=(s\|\widehat D\|+|a\varepsilon_0'|)/3.
\end{aligned}
\]
Then
\[
 |V_\gamma a^3\rho_s-E_\infty|\le b_E,\quad |Vp_s|\le b_p,
 \qquad
 |p_s/\rho_s|\le\frac{b_p}{E_\infty-b_E}\quad(E_\infty>b_E).
\]
An actual epoch with $E(t)>b_E(t)$ certifies $E_\infty>0$.
\end{corollary}
\begin{proof}
Apply the preceding tail inequality and $E\ge E_\infty-b_E$.
\end{proof}
The theorem does not prove convergence of the whole interaction-picture
state, select the initial state, change its center, create massive matter
from $H_m=0$, or establish inhomogeneous dust or its abundance.  The
complete isotropic and zero-flux FLRW certificate is still consumed.
Unbounded bosonic Fock operators require their own domain and moment
control; bounded finite CAR estimates are not silently extended to them.

\begin{corollary}[Monotone massive sign-population tails]
\label{p3f:cor:v2217-sign-population-tail}
Restrict the preceding finite-source theorem to a fixed complete free CAR
band with one-body blocks
$h_r=(\lambda_r/a)\sigma_1+m\sigma_3$, $m>0$ constant,
degeneracy $g_r$, no omitted interaction and nondecreasing $a\to\infty$.
For arbitrary inherited finite-Fock correlations the instantaneous sign
populations have limits and satisfy
\[
 |n_{r\pm}(\infty)-n_{r\pm}(t)|
 \le\frac{g_r}{2}\arctan\frac{|\lambda_r|}{ma(t)}.
\]
On any finite monotone subinterval one likewise has
$|n_{r\pm}(t_f)-n_{r\pm}(t_i)|\le
g_r[\theta_r(t_i)-\theta_r(t_f)]/2$.
\end{corollary}
\begin{proof}
The one-body covariance $C_r$ obeys $0\le C_r\le I$ and
$i\hbar\dot C_r=[h_r,C_r]$ under the complete free Hamiltonian.
For the instantaneous sign projector $P_{r\pm}$,
$\dot n_{r\pm}=\operatorname{Tr}(C_r\dot P_{r\pm})$, because
$[h_r,P_{r\pm}]=0$. With
$\theta_r=\arctan(|\lambda_r|/(ma))$, the eigenvalues of
$\dot P_{r\pm}$ are $\pm|\dot\theta_r|/2$, each $g_r$-fold. Positivity
of $C_r$ bounds the trace by $g_r|\dot\theta_r|/2$.
Its remaining integral is $g_r\theta_r(t)/2$, proving convergence and error.
\end{proof}
This adds population control to the broader bounded-variation energy
theorem; it does not require convergence of the whole state and does not
extend a finite CAR norm bound to unbounded bosonic Fock operators.

\subsection{Actual coupled evolution and terminal conservation}
\label{p3f:sec:v2216-coupled-terminal}

These results consume the already required exact common-action matching
domain; they do not manufacture it.  Write $M^2=2\nu$, $C=C_7$,
$A_0=C(U+C)$, $p=f_A'(r)>1$, and
$P(p)=pr(p)-f_A(r(p))$.  On an admitted canonical native scalar ray use
\begin{equation}
 S_A=\int\sqrt{-g}\left[
 \nu\{pR-P(p)\}-\frac{h^2R}{12}
 -\frac12(\nabla h)^2-V(h)\right]+S_s,
 \quad V=\lambda h^4-\Lambda_g^2h^2.
 \label{p3f:eq:v2216-coupled-action}
\end{equation}
Every central, source and boundary term is retained.  A frozen external
depth or role is not varied as a new field.  The terminal quotient fixes
$p=1$ and removes only its independent variation; it retains $h$ and the
source.  It is a marked variational stratum, rather than a curvature rule
that reactivates the discarded coordinate at $R_g=2C$.

\begin{theorem}[Coupled normal form and exact lapse preservation]
\label{p3f:thm:v2216-coupled-normal-form}
For an actual homogeneous isotropic, zero-flux, fixed-rank
coordinate-only finite source let its proper generator be $H_s(a,h)$.
Set $V_3=V_\gamma a^3$ and
\[
 \rho_s=\frac{\operatorname{Tr}(\sigma H_s)}{V_3},\quad
 p_s=-\frac{a\operatorname{Tr}(\sigma\partial_aH_s)}{3V_3},\quad
 J=\frac{\operatorname{Tr}(\sigma\partial_hH_s)}{V_3},\quad
 \dot\sigma=-i[H_s,\sigma].
\]
Define $\Theta=hV'-4V+hJ-\rho_s+3p_s$,
$\mathcal A=\nu p-h^2/12$, $v=\dot p$, $w=\dot h$.
The complete homogeneous equations are
\begin{align}
 \dot a&=aH,&\dot H&=r(p)/6-2H^2-a^{-2},\\
 \dot p&=v,&\dot v&=-3Hv+[2P-pr-\Theta/(2\nu)]/3,\\
 \dot h&=w,&\dot w&=-3Hw-rh/6-V'-J,
 \qquad\dot\sigma=-i[H_s,\sigma].
 \label{p3f:eq:v2216-coupled-normal-form}
\end{align}
The lapse constraint and its exact off-constraint identity are
\begin{align}
 \mathcal C&=6\mathcal A(H^2+a^{-2})-\nu P
 -\tfrac12w^2-V-\rho_s+6H(\nu v-hw/6),\\
 \boxed{\dot{\mathcal C}=-4H\mathcal C.}
 \label{p3f:eq:v2216-constraint-preservation}
\end{align}
This is a locally unique evolution on each regular open active chart
with $\mathcal A>0$ and smooth supplied source.  The finite source
realization consumes the actual local isotropic embedding, not just
integrated variations.
\end{theorem}
\begin{proof}
Cyclicity gives the source work law
$\dot\rho_s+3H(\rho_s+p_s)=Jw$.
Variation of $p$ gives $R=r(p)$ and hence the displayed geometric
equation.  Scalar variation gives the $w$ equation.  The improved trace
of the complete metric equation gives the $v$ equation with the stated
$\Theta$; lapse variation gives $\mathcal C$.
For the constraint certificate use
\[
 \dot{\mathcal A}=\nu v-hw/6,\quad \dot P=rv,
 \quad\frac d{dt}(\tfrac12w^2+V+\rho_s)
 =-3H(w^2+\rho_s+p_s)-rhw/6.
\]
Differentiate $\mathcal C$, substitute these identities and
\eqref{p3f:eq:v2216-coupled-normal-form}, and collect the coefficients
of $H^2+a^{-2}$, $P$, $w^2$, $V$, $\rho_s$ and
$H\dot{\mathcal A}$.  The result is exactly $-4H\mathcal C$;
the $Jw$ terms cancel before collecting.  Thus initial $\mathcal C=0$
is preserved.  Smoothness of $r,P,H_s$ on an open active chart and the
ordinary finite-dimensional local existence theorem give uniqueness.
\end{proof}

\begin{theorem}[Terminal momentum budget and the required trace]
\label{p3f:thm:v2216-terminal-momentum-budget}
Put $B=2f_A-pr$.  Along the actual coupled history,
\begin{equation}
 \frac d{dt}(a^3\dot p)=-\frac{a^3}{3}
       \left(B+\frac{\Theta}{2\nu}\right).
 \label{p3f:eq:v2216-momentum-balance}
\end{equation}
An impulse-free tangent arrival $\dot p(t_C)=0$ requires
\[
 \int_{t_i}^{t_C}a^3(B+\Theta/(2\nu))dt=3a_i^3\dot p_i.
\]
For incoming $\dot p_i<0$, define
$N_\Theta=\int a^3[-\Theta/(2\nu)]_+dt$ and $B_-=[-B]_+$.
Then necessarily
\[
 N_\Theta\ge\max\{0,\,3a_i^3|\dot p_i|-\int a^3 B_-dt\}.
\]
The actual native scalar obeys $hV'-4V=2\Lambda_g^2h^2\ge0$.
An isolated such scalar cannot brake the descending momentum wherever
$B\ge0$.  On a lower interval with $B\ge-C$, stopping it without
another negative trace requires
$\int a^3dt\ge-3a_i^3\dot p_i/C$.
\end{theorem}
\begin{proof}
Multiply the $v$ equation by $a^3$, using $2P-pr=-B$.
Integration gives the endpoint equality.  Its negative part is at most
$\int a^3B_-dt+N_\Theta$; the positive part can only increase the
required negative part, proving the bound.  Direct differentiation of
the quartic gives the native trace sign.  With no additional trace,
$d(a^3\dot p)/dt\le a^3C/3$ on the stated lower interval, giving the
duration inequality.
\end{proof}
The associated scalaron work identity is also exact: if
$\mathcal W_p=B/3$ and $\mathcal E_p=\dot p^2/2+\mathcal W$, then
\[
 \dot{\mathcal E}_p=-3H\dot p^2-\Theta\dot p/(6\nu).
\]
On a descending path a negative $\Theta$ extracts scalaron energy;
positive energy density alone does not determine that sign.  For an
actual Dirac pair $H=(\ell/a)\sigma_3+m\sigma_1$ with covariance
Bloch vector $\boldsymbol s$, the noncentral trace is $-ms_1/V_3$.
If $m=m(h)$, the combined trace entering this theorem is
$(hm'-m)s_1/V_3$, with all central derivatives retained.  A degree-one
mass cancels this combined insertion.  At $a=\ell=m=1$ the two pure
states $\boldsymbol s_\pm=\pm(1,0,-1)/\sqrt2$ have equal energy and
opposite traces.  Thus no energy or occupation ledger alone reconstructs
the braking trace.  Equal populations on a complete degeneracy multiplet
preserve the isotropic integrated witness; the actual prepared state
must still be evolved.

\begin{theorem}[Local compatible terminal handoff]
\label{p3f:thm:v2216-local-terminal-handoff}
The lower active inverse has
$p-1\sim(r-C)^2/A_0$ and $P(1)=P'(1^+)=C$.
Choose complete endpoint data $(a_C,H_C,h_C,w_C,\sigma_C)$ with
$a_C>0$, $\mathcal A_C=\nu-h_C^2/12>0$ satisfying the endpoint lapse
constraint.  Evaluate its actual source ports and put
\[
 \kappa=\tfrac13(C-\Theta_C/(2\nu)).
\]
If $\kappa>0$, an incoming active solution exists locally on $t<0$,
with $p(0)=1$, $\dot p(0)=0$, and
\[
 p(t)=1+\kappa t^2/2+o(t^2),\qquad
 r(t)=C+\sqrt{A_0\kappa/2}(-t)+o(|t|).
\]
It joins without an impulse to the terminal action with $p=1$,
retained $h$, and identical source and canonical boundary data.  The
metric and fields are $C^1$ across the join, with bounded steps
\begin{equation}
 R_+=2C-\Theta_C/(2\nu),\qquad [R]=3\kappa,
 \quad[\dot H]=\kappa/2,\quad[\ddot h]=-h_C\kappa/2.
 \label{p3f:eq:v2216-handoff-jumps}
\end{equation}
Conversely a twice-differentiable tangent arrival from $p>1$ cannot
have $\kappa<0$.  Thus $\Theta_C\le2\nu C$ is necessary; strict
inequality is locally sufficient on the stated compatible domain.
\end{theorem}
\begin{proof}
Extend the normal-form vector field continuously to $p<1$ by taking
$r=C$ and $P=Cp$ there.  Peano existence gives a local solution at
the endpoint.  Its $p$ acceleration is $\kappa>0$, hence both short
punctured sides have $p>1$ and the incoming side does not use the
extension.  The constraint identity preserves $\mathcal C=0$.
The terminal scalar--metric system is regular for $\mathcal A_C>0$
on the supplied source chart and starts from the same lapse data.
Its improved trace gives $R_+$, and the geometric and scalar equations
give the remaining jumps in \eqref{p3f:eq:v2216-handoff-jumps}.
The common metric momentum depends on $\mathcal A,H,\dot{\mathcal A}$
and the scalar momentum on $w+Hh$; both agree at $\dot p=0$.
The retained non-null boundary completions have opposite normals and
cancel on the matched data.  Since the metric, $p$, and $h$ are $C^1$,
their second derivatives have no delta layer.  The future $p$ variation
is absent by the terminal role quotient; retaining $R=P'(1)=C$ in
that open region would define a different variational problem.
Finally an incoming tangent minimum of $p-1\ge0$ has
$\ddot p(0)=\kappa\ge0$, proving necessity.
\end{proof}
This is local existence for compatible data, rather than selection or
arrival of the cyclic history.  For an isolated native scalar it requires
$h_C^2\le\nu C/\Lambda_g^2$.  Requiring continuous curvature imposes
$\kappa=0$, after which higher source jets decide arrival.  An endpoint
functional cannot delete an active bulk Euler equation on an open future
region.  The role quotient, complete action matching and actual incoming
history remain distinct consumed operations.

\begin{theorem}[Full nonminimal corner-energy identity]
\label{p3f:thm:v2216-corner-energy}
At $p=1$ keep the induced $a,h$ continuous, but allow $H,w$ to jump.
Let the terminal side have $\dot p_+=0$, and match the common canonical
momenta of the complete action.  Then
\[
 H_+=H_-+\dot p_-/2,\qquad
 w_+=w_--h\dot p_-/2,
\]
and the two complete lapse constraints necessarily require
\begin{equation}
 \boxed{\rho_{s,+}-\rho_{s,-}
 =\tfrac{3\nu}{2}\dot p_-^2
 =\tfrac{3M^2}{4}\dot p_-^2.}
 \label{p3f:eq:v2216-corner-energy}
\end{equation}
The native conformal scalar does not absorb this missing energy through
its canonical velocity jump.  A bounded continuous coordinate-only
finite-source generator with continuous $a,h$ cannot supply the positive
zero-duration density jump when $\dot p_-\ne0$.
\end{theorem}
\begin{proof}
Suppress the common positive volume factors.  The momenta are
$\pi_a=-12\mathcal A H-6\nu\dot p+hw$ and
$\pi_h=w+Hh$.  The second gives $[w]=-h[H]$; the first then uses
$12\mathcal A+h^2=12\nu$ to give $[H]=\dot p_-/2$.
Substitute both into the difference of the complete constraints at
$P(1)=C$.  All terms containing $h,w$ and their crossproducts cancel,
leaving \eqref{p3f:eq:v2216-corner-energy}.  A bounded continuous
finite generator has a continuous unitary propagator and density;
continuous $a,h$ then make $\operatorname{Tr}(\sigma H_s)/V_3$
continuous, which cannot produce that jump.
\end{proof}
The identity is a necessary accounting condition, not an interface
constructor.  Singular or velocity-dependent source/interface terms
retain their actual momenta and define a different problem.  A
finite-time transfer must be evolved before arrival and satisfy
\eqref{p3f:eq:v2216-momentum-balance}; susceptibility exhaustion alone
does not supply it.

\subsection{Finite locality and quantitative action-matching control}
\label{p3f:sec:v2216-quantitative-control}
\begin{theorem}[Finite-frontier locality certificate]
\label{p3f:thm:v2216-finite-locality}
Let $L\ge0$ be self-adjoint with compact resolvent and satisfy the
smooth heat estimate
$\|\mathbf1_Ee^{-sL}\mathbf1_F\|\le e^{-d(E,F)^2/(4s)}$.
Let $P_F=\mathbf1_{[0,\Lambda_F^2]}(L)$ have finite range and let
$\lambda_+$ be the least omitted eigenvalue.  Put
$Q=e^{-tL/\Lambda_c^2}$, $Q_F=P_FQ$, $t,\Lambda_c>0$, and
$A_{a,F}=Q_FM_aQ_F$.  For bounded multiplication functions $a,b$
with disjoint supports $E,F$,
\begin{equation}
 \|[A_{a,F},A_{b,F}]\|\le
 2\|a\|_\infty\|b\|_\infty
 [e^{-\Lambda_c^2d(E,F)^2/(8t)}
       +e^{-2t\lambda_+/\Lambda_c^2}].
 \label{p3f:eq:v2216-finite-locality}
\end{equation}
The second term is zero if there is no omitted spectrum.
\end{theorem}
\begin{proof}
$P_F$ commutes with $Q$, and $Q_F$ is a self-adjoint contraction.
The exact tail is
$\|Q^2-Q_F^2\|=e^{-2t\lambda_+/\Lambda_c^2}$.
Consequently the norm of $\mathbf1_EQ_F^2\mathbf1_F$ is bounded
by the sum in brackets.  Factor
\[
 [A_{a,F},A_{b,F}]
 =Q_FM_aQ_F^2M_bQ_F-Q_FM_bQ_F^2M_aQ_F.
\]
Each term passes through one separated-support block and each exterior
$Q_F$ has norm at most one.  The triangle inequality proves the bound.
\end{proof}
For tolerance $0<\epsilon<1$, it suffices that
\[
 \Lambda_cd(E,F)\ge\sqrt{8t\log(4/\epsilon)},\qquad
 \lambda_+/\Lambda_c^2\ge\log(4/\epsilon)/(2t),
\]
since each exponential is at most $\epsilon/4$.  On a round scalar
band $\bigoplus_{j=0}^kE_j$,
$\lambda_+=(k+1)(k+3)/a^2$ and
$d_k=(k+1)(k+2)(2k+3)/6$, so the omitted-tail error is exactly
$\exp[-2t(k+1)(k+3)/(a\Lambda_c)^2]$.
This combines the smooth heat criterion with the upstream exact spectral
tail; it does not identify a sharp frontier with the action scale,
promote a scalar frontier to a spinor cutoff, or establish a Lorentzian
two-derivative approximation.

\begin{theorem}[Action residuals, vacuum roots and trace stability]
\label{p3f:thm:v2216-action-residual-control}
On $[C,U]$, let $f_0''=\chi_0$ and let an independently calculated
physical action obey
$f_{\rm act}''=\chi_0+r$,
$f_{\rm act}(C)-f_0(C)=\delta f_C$,
$f'_{\rm act}(C)-f'_0(C)=\delta F_C$, with continuous $r$.
For $F=f'$ and $B=2f-RF$,
\begin{align}
 \delta f(R)&=\delta f_C+\delta F_C(R-C)
                 +\int_C^R(R-s)r(s)ds,\\
 \delta F(R)&=\delta F_C+\int_C^Rr(s)ds,\\
 \delta B(R)&=2\delta f_C+\delta F_C(R-2C)
                 +\int_C^R(R-2s)r(s)ds.
 \label{p3f:eq:v2216-action-residuals}
\end{align}
Thus $\|r\|_\infty\le\epsilon_\chi$ bounds
$|\delta f|\le|\delta f_C|+|\delta F_C|(R-C)
+\epsilon_\chi(R-C)^2/2$.
If additionally $r$ has bounded variation and
$N_0=F_0-R\chi_0\ge m>0$, then $N_{\rm act}>0$ whenever
\[
 |\delta F_C|+C|r(C)|+\int_{(C,U]}s\,d|r|(s)<m.
\]
For a compact $I\subset(C,U)$, suppose $B_0'\ge m_0>0$,
$B_{\rm act}'\ge m_1>0$, and both zeros lie in $I$.  If
$\sup_I|\delta B|\le\epsilon_B$, their locations satisfy
\[
 |R_{\rm act}-R_0|\le\epsilon_B/\max(m_0,m_1).
\]
Existence of the physical zero follows if the reference endpoint values
of $B_0$ straddle $[-\epsilon_B,\epsilon_B]$ strictly.
\end{theorem}
\begin{proof}
Integrate $\delta f''=r$ twice with its actual endpoint value and
slope; substitution in $2\delta f-R\delta F$ gives
\eqref{p3f:eq:v2216-action-residuals}.  The uniform estimate follows
by integrating the positive kernel $R-s$.  Stieltjes integration by
parts gives
\[
 N_{\rm act}-N_0=\delta F_C-Cr(C)-\int_{(C,R]}s\,dr(s),
\]
and its variation bound proves positivity.  At the physical zero,
$|B_0|\le\epsilon_B$, giving the root bound with $m_0$ by the
mean-value theorem.  At the reference zero the same argument with
$B_{\rm act}$ gives the bound with $m_1$; both hold simultaneously.
The endpoint signs and continuity give existence, and the strictly
positive derivative gives uniqueness.
\end{proof}
Trace-mass stability additionally consumes $\chi_{\rm act}>0$ and
a background solving the complete equations.  A Hessian comparison
alone cannot control a tiny vacuum curvature: an arbitrary endpoint
value shifts $B$ by $2\delta f_C$ while leaving every Hessian unchanged.
The source's terminal intercept is therefore not re-derived by matching
only its susceptibility.

For an actual positive transported action coefficient $\mathcal U(b)$,
define the residual one-form $\eta_U=d\log\mathcal U+d\Phi$.
Integration along an admitted path gives exactly
\[
 \frac{\mathcal U(b)}{\mathcal U(b_0)}
 =\exp[-\Phi(b)+\Phi(b_0)+\int_{b_0}^b\eta_U].
\]
An independently licensed fourth-root energy inherits correction
$\exp(\tfrac14\int\eta_U)$.  This records an evaluated residual,
rather than setting it to zero.  The whole-action normalization must
likewise transport its gauge and source terms; the raw/scalar factor
$96$ is not a gravity-only repair.


\section{Retained-response local dynamics: inherited blocks and geometric work}
\label{p3f:sec:v2209-retained-local-dynamics}

The following results use the complete common-action domain of
Section~\ref{p3f:sec:v2202-matched-native-sources}, its physical constraints
and its already admitted source state. They introduce no new scalar
species, source efficiency, cutoff or independent energy allocation.
They distinguish a local kinetic-transfer identity from net capture and
from the subsequent homogeneous or spatial readout of the same state.

\subsubsection{The positive square is the inherited quadratic parent}
\begin{proposition}[Frozen Schur square with physical kinetic pairing]
\label{p3f:prop:v2209-inherited-schur-square}
On a fixed physical constrained chart let the actual real quadratic
potential in retained coordinates $b$ and interior coordinates $x$ be
\begin{equation}
 U_2=\tfrac12b^TAb-b^TGx+\tfrac12x^TCx,
 \qquad A=A^T,\quad C=C^T>0.
 \label{p3f:eq:v2209-quadratic-parent}
\end{equation}
Then, with $K=A-GC^{-1}G^T$, its exact decomposition is
\begin{equation}
 \boxed{U_2=\tfrac12(x-C^{-1}G^Tb)^TC(x-C^{-1}G^Tb)
                         +\tfrac12b^TKb.}
 \label{p3f:eq:v2209-inherited-schur-square}
\end{equation}
For an actual frozen interior kinetic matrix $\mathsf M_x>0$, define
$\widetilde x=\mathsf M_x^{1/2}x$,
$\widetilde G=G\mathsf M_x^{-1/2}$ and
$\mathsf\Omega^2=\mathsf M_x^{-1/2}C\mathsf M_x^{-1/2}$.
The whitened square has centre $\mathsf\Omega^{-2}\widetilde G^Tb$ and
$A=K+\widetilde G\mathsf\Omega^{-2}\widetilde G^T$.
If $K\ge0$ the whole quadratic potential is nonnegative.
\end{proposition}
\begin{proof}
Expansion of the square gives the mixed term $-b^TGx$ and subtracts
$b^TGC^{-1}G^Tb/2$ from its retained block. The linear kinetic change
gives the stated congruences. Positivity follows from the two displayed
quadratic forms. No term can be chosen independently of $A,C,G$.
\end{proof}
Whitening this block does not erase actual mixed kinetic contacts or
moving-frame terms. Those remain in the complete dynamics. In particular,
a constant-frequency extension of this frozen square to all amplitudes
is not supplied by the identity; nor does it turn its response coordinates
into an additional microscopic energy carrier.

\begin{corollary}[Analytic regular reduction and the cubic critical point]
\label{p3f:cor:v2209-analytic-criticality}
A finite real-analytic parent $U(b,x)$ on an open regular chart, with a
stationary interior solution and invertible $U_{,xx}$ near the
origin, has a real-analytic stationary envelope. Consequently a reduced
ray law $\kappa|b|^3$, $\kappa\ne0$, at $b=0$ cannot be obtained on all
these hypotheses. Its derivation requires the actual critical interior,
nonanalytic domain/pullback or a nonuniform limiting reduction.
\end{corollary}
\begin{proof}
The analytic implicit-function theorem gives analytic $x_*(b)$ from
$U_{,x}=0$; composition makes $U(b,x_*(b))$ analytic.
The third derivative of $|b|^3$ jumps at zero.
\end{proof}
This criterion does not reopen a separately proved static critical RAR
reduction. It prevents replacing its critical mechanism by an unspecified
regular analytic finite block.

\subsubsection{Native normal-to-tangent kinetic transfer}
\begin{theorem}[Geometric work and conditional local tangent loading]
\label{p3f:thm:v2209-geometric-transfer}
Restrict Theorem~\ref{p3f:thm:v2202-hyperbolic-target} to an actually
retained homogeneous real line $h$, with a consistent restriction of all
forces and constraints. Put
\begin{equation}
 \ell_\Phi^2=6M^2,\qquad \sigma=\tfrac12\log F,\qquad
 u=\frac{h}{\ell_\Phi\sqrt F},\qquad D=1-u^2>0.
 \label{p3f:eq:v2209-adapted-target}
\end{equation}
Assume its complete principal is the displayed two-derivative target;
residual sources act through proper-time force covectors $J_\sigma,J_u$.
With $V_E=2\pi^2a_E^3$, $H_E=\dot a_E/a_E$, and dots denoting Einstein
proper time, this homogeneous constituent is
\begin{equation}
 L_{g\Phi}=V_E(T_\sigma+T_u-W_E),\qquad
 T_\sigma=\frac{\ell_\Phi^2\dot\sigma^2}{2D},\quad
 T_u=\frac{\ell_\Phi^2\dot u^2}{2D^2}.
 \label{p3f:eq:v2209-adapted-kinetics}
\end{equation}
The exact kinetic work identities are
\begin{align}
 \dot T_\sigma+6H_ET_\sigma
   &=-(W_{E,\sigma}+J_\sigma)\dot\sigma-Q_g,\nonumber\\
 \dot T_u+6H_ET_u&=-(W_{E,u}+J_u)\dot u+Q_g,\nonumber\\
 \boxed{Q_g=\frac{\ell_\Phi^2u\dot u\dot\sigma^2}{D^2}.}
 \label{p3f:eq:v2209-geometric-work}
\end{align}
At a regular initial point $t_i$ with $\dot u_i=0$ and
$(W_{E,u}+J_u)_i=0$, the tangent has
\begin{equation}
 \ddot u_i=u_i\dot\sigma_i^2,\qquad
 \boxed{T_u(t)=\frac{\ell_\Phi^2u_i^2\dot\sigma_i^4}{2D_i^2}
                      (t-t_i)^2+O((t-t_i)^3).}
 \label{p3f:eq:v2209-local-positive-loading}
\end{equation}
Thus it is locally loaded if $u_i\dot\sigma_i\ne0$, assuming sufficient
smoothness for the displayed Taylor remainder.
\end{theorem}
\begin{proof}
Use $\Omega=F D$, $h=\ell_\Phi e^\sigma u$ and
$\psi=\ell_\Phi(\sigma+\log D/2)$ in the native target.
The cross terms cancel, giving
$d\ell^2=\ell_\Phi^2(d\sigma^2/D+du^2/D^2)$.
The $\sigma$ Euler equation is
$d(V_E\ell_\Phi^2\dot\sigma/D)/dt=-V_E(W_{E,\sigma}+J_\sigma)$.
Multiply it by $\dot\sigma/V_E$ and use $\dot V_E=3H_EV_E$;
the derivative of $D^{-1}$ gives $-Q_g$ in $\dot T_\sigma$.
The $u$ equation gives the opposite geometric work.
At $\dot u_i=0$ it reads
$\ell_\Phi^2\ddot u_i/D_i^2-
 \ell_\Phi^2u_i\dot\sigma_i^2/D_i^2+(W_{E,u}+J_u)_i=0$.
The initial balance and Taylor expansion prove the last equation.
\end{proof}
This is reversible connection-mediated transfer, not a proof of trapping,
positive completed stress or terminal survival. When the actual force
vanishes on $u=0$, uniqueness preserves the classical $u=\dot u=0$
branch. A covariance production claim on that branch must use an actually
admitted state and its complete pump. Additional source acceleration,
velocity and frame terms instead consume the complete principal matrix
of Proposition~\ref{p3f:prop:complete-source-cauchy-matrix}; the coefficient
above cannot then be retained while discarding those terms.

\subsubsection{Complete normal turnover and the physical history}
\begin{proposition}[Matched normal balance with the residual scalar port]
\label{p3f:prop:v2209-complete-normal-turnover}
On the same common-action restriction let the residual action $S_r$ have
no separately inserted $F$ dependence. Retain its complete Jordan metric
trace $T_r=-\Theta_r$ and scalar force
$\delta_hS_r=-\int\sqrt{-g_J}\,J_h\delta h$.
For the native scalar potential, whose quadratic part is
$-\Lambda_g^2h^2$ and quartic part is homogeneous of degree four, put
\begin{equation}
 B(F)=2f(r(F))-r(F)F,\qquad
 \Theta_{\rm complete}=\Theta_r-2\Lambda_g^2h^2-hJ_h.
 \label{p3f:eq:v2209-complete-normal-trace}
\end{equation}
If the actual source has no normal velocity-momentum correction, then
\begin{equation}
 \boxed{P_\sigma=\frac{V_E\ell_\Phi^2\dot\sigma}{D},\qquad
 \dot P_\sigma=-\frac{V_E}{F^2D^2}
                    (M^2B-\Theta_{\rm complete}).}
 \label{p3f:eq:v2209-canonical-turnover}
\end{equation}
Its canonical momentum turnover satisfies
$M^2B=\Theta_{\rm complete}$; a maximum additionally requires the
corresponding sign crossing of $\dot P_\sigma$ along the actual history.
\end{proposition}
\begin{proof}
At fixed $u$, $\partial_\sigma F=2F$, $\partial_\sigma h=h$ and
$\partial_\sigma\Omega=2\Omega$. Differentiating
$W_E=[M^2P(F)/2+V_\Phi(h)]/(F^2D^2)$, with $P_{,F}=r(F)$,
gives $(M^2B+2\Lambda_g^2h^2)/(F^2D^2)$.
At fixed $g_E$, $\delta_\sigma g_J=-2g_J\delta\sigma$;
the complete residual variation is
$\delta_\sigma S_r=\int\sqrt{-g_J}(\Theta_r-hJ_h)\delta\sigma$.
Hence $J_\sigma=(-\Theta_r+hJ_h)/(F^2D^2)$, and the normal Euler
equation in the preceding proof gives the result.
\end{proof}

For an actually retained degree-$q$ Dirac mass map $\mathsf M(h)$,
Corollary~\ref{p3f:cor:v2202-mass-map-selector} implies that its
contribution to $\Theta_r-hJ_h$ is
$(1-q)\bar\psi\mathsf M(h)\psi$. Degree one therefore cancels this
combined normal force contribution, not its energy, pressure or metric
response. No degree-one map is selected by the word ``Yukawa'' alone.

For a velocity-dependent finite source with density $\sigma_{\rm src}$,
the momentum is instead
$P_i=L_{{\rm cl},v_i}-\operatorname{Tr}(\sigma_{\rm src}H_{,v_i})$
and the complete principal is
$L_{{\rm cl},v_iv_j}-\operatorname{Tr}(\sigma_{\rm src}H_{,v_iv_j})$.
Functional insertions retain their endpoints and contacts. Such a source
must use Eq.~\eqref{p3f:eq:complete-source-cauchy-momenta}, rather than
the uncorrected momentum of Eq.~\eqref{p3f:eq:v2209-canonical-turnover}.
Even on its displayed domain, canonical momentum turnover is neither
$\ddot\sigma=0$ nor a positive retained-energy increment. The latter
uses the complete pressure/work identity of
Eq.~\eqref{p3f:eq:v2207-quadratic-ward-work} or the corresponding full
source insertion. A physical epoch requires the inherited solution and
its clock conversion; no hierarchy midpoint is selected by this balance.


\subsection{Physical transverse response and complete boundary variation}
\label{p3f:sec:v2210-lit-native-response}

The results below sharpen the existing common-action restriction and its
physical response tests.  They use the native target of
Theorem~\ref{p3f:thm:v2202-hyperbolic-target}, not an added spectator,
and the actual boundary/source word of
Proposition~\ref{p3f:prop:v1164-scalar-overlap-condition}.  They do not
select a populated state, a moving wall, a microscopic depth projection,
or a new action coefficient.  A result for a supplied regular restriction
is retained as such; its physical domain is part of its certificate.

\subsubsection{The physical scalar support and its complete transverse operator}

\begin{proposition}[Surviving support and the native curvature test]
\label{p3f:prop:v2210-lit-support-curvature}
Suppose an actual gauge/constraint reduction retains a consistently
restricted real scalar line in
Theorem~\ref{p3f:thm:v2209-geometric-transfer}.  With
$\ell_\Phi^2=6M^2$, $u=\tanh r$, set
$\eta_\Phi=\ell_\Phi r$ and
$\varphi_\parallel=\ell_\Phi\sigma$.  Its displayed principal metric is
\begin{equation}
\label{p3f:eq:v2210-lit-fermi-target}
 d\ell^2=d\eta_\Phi^2+
 \cosh^2(\eta_\Phi/\ell_\Phi)d\varphi_\parallel^2,
 \qquad K=-\ell_\Phi^{-2}=-\frac1{6M^2}.
\end{equation}
The two-dimensional target scalar curvature is $-2/\ell_\Phi^2$.
If no physical scalar direction survives the actual quotient, the active
target is only the scalaron line.  It then has no intrinsic sectional
curvature and no transverse entropy direction.  A static response
coordinate, or a scalar removed by the physical gauge quotient, cannot
replace the surviving line required here.
\end{proposition}
\begin{proof}
Since $du=(1-u^2)dr$, the metric in
Eq.~\eqref{p3f:eq:v2209-adapted-kinetics} becomes
$\ell_\Phi^2[dr^2+\cosh^2r\,d\sigma^2]$.
For $d\eta^2+f(\eta)^2d\varphi^2$ the Gaussian curvature is
$-f''/f$; $f=\cosh(\eta/\ell_\Phi)$ gives the displayed value.
The absence of a two-plane when the quotient leaves only a line proves
the final statements.  Retention of that line is a hypothesis, rather
than a consequence of the coordinate change.
\end{proof}

\begin{theorem}[Native geodesic transverse operator with complete source contacts]
\label{p3f:thm:v2210-lit-transverse-operator}
Work on the claim-complete regular Einstein-frame restriction of the
preceding proposition and on a round closed spatial metric
$a_E^2\gamma$.  Dots denote Einstein proper time.  Let an actual background
have $u=\dot u=0$, with its complete transverse force vanishing along
that history.  Require that the displayed target is the retained
two-derivative scalar constituent; every additional principal, mixed,
metric, state and functional source insertion is retained in the full
linear blocks below.  In particular, vanishing of the background force
does not assert vanishing of its linear contact.

At $u=0$ the normalized transverse perturbation is
$Q_\perp=\delta\eta_\Phi=\delta h/\sqrt F$.  The intrinsic target/potential
operator acting on it is
\begin{equation}
\label{p3f:eq:v2210-lit-geodesic-operator}
 \boxed{\mathcal D_{\perp,0}=
 \partial_t^2+3H_E\partial_t-a_E^{-2}\Delta_\gamma
 +\frac{W_{E,uu}}{\ell_\Phi^2}-\dot\sigma^2.}
\end{equation}
Equivalently, its curvature mass contribution is
$-\dot\varphi_\parallel^2/\ell_\Phi^2$.
Here the potential derivative is at fixed $\sigma$ on the actual
background.  The geodesic has zero turn rate.  On the regular source-free
domain of the covariant super-Hubble two-field Einstein entropy-mass
formula, the expression also contains the covariant potential Hessian
and the turn contribution;
the curvature term is $\dot\varphi^2K$
\cite{LitDYNRenauxPetelTurzynski2016}.

Let $z$ collect all other retained physical perturbations and source/state
variables after the actual constraint treatment.  If the full causal
linearization has the admitted form
\begin{equation}
\label{p3f:eq:v2210-lit-complete-transverse-block}
 \begin{pmatrix}
 \mathcal D_{\perp,0}+\Pi_{\perp\perp}^{R}&B^{R}\\
 D^{R}&C^{R}
 \end{pmatrix}
 \binom{Q_\perp}{z}=\binom{j_\perp}{j_z},
\end{equation}
with $(C^{R})^{-1}_{R}$ on its declared domain, then its exact retained
operator and drive are
\begin{equation}
\label{p3f:eq:v2210-lit-complete-transverse-schur}
 \boxed{\mathcal D_{\perp,\rm phys}^{R}
 =\mathcal D_{\perp,0}+\Pi_{\perp\perp}^{R}
 -B^{R}(C^{R})^{-1}_{R}D^{R},\qquad
 j_{\perp,\rm eff}=j_\perp-B^{R}(C^{R})^{-1}_{R}j_z.}
\end{equation}
The contact $\Pi^R$ includes the complete variation of the residual
transverse force and all actually consumed local derivative contacts.
An instantaneous scalar-force contact $J_{u,u}$, with other variables
fixed, contributes $J_{u,u}/\ell_\Phi^2$ to this operator.
Other source velocities, acceleration insertions and memory must be
retained as operators; they are not replaced by that one coefficient.
\end{theorem}
\begin{proof}
Expand Eq.~\eqref{p3f:eq:v2210-lit-fermi-target} at
$\eta_\Phi=0$.  Its parallel kinetic term contributes
$\dot\varphi_\parallel^2Q_\perp^2/(2\ell_\Phi^2)$ to the quadratic
Lagrangian, whereas the potential contributes
$-W_{E,\eta_\Phi\eta_\Phi}Q_\perp^2/2$.
At zero transverse coordinate,
$W_{E,\eta_\Phi\eta_\Phi}=W_{E,uu}/\ell_\Phi^2$.
The transverse kinetic and spatial-gradient terms give
$\partial_t^2+3H_E\partial_t-a_E^{-2}\Delta_\gamma$.
The first-order force vanishes by hypothesis, so these quadratic terms
give Eq.~\eqref{p3f:eq:v2210-lit-geodesic-operator}.  This directly
derives its curvature sign and normalization without importing a new
target metric.  Since $J_{\eta_\Phi}=J_u/\ell_\Phi$ at zero and
$\delta u=Q_\perp/\ell_\Phi$, its fixed-background derivative has the
stated normalization.  All other derivatives belong to the full
linearization.  Solving the second row of
Eq.~\eqref{p3f:eq:v2210-lit-complete-transverse-block} for $z$ gives
the Schur expression and effective drive.  The hypothesis that these
are the complete physical blocks is essential.
\end{proof}

\begin{corollary}[Evaluated native potential and frozen metric-trace contact]
\label{p3f:cor:v2210-lit-native-potential-contact}
On the matched native scalar line use
$h=\ell_\Phi\sqrt F\,u$, $F=e^{2\sigma}$ and the existing potential
$V_\Phi(h)=-\Lambda_g^2h^2+V_4(h)$ with $V_4$ homogeneous of degree
four.  Its displayed Einstein potential satisfies, exactly at $u=0$,
\begin{equation}
\label{p3f:eq:v2210-lit-native-potential-mass}
 \boxed{\frac{W_{E,uu}}{\ell_\Phi^2}
 =\frac{P(F)}{3F^2}-\frac{2\Lambda_g^2}{F},\qquad
 m_{\perp,0}^2=\frac{P(F)}{3F^2}-\frac{2\Lambda_g^2}{F}
                       -\dot\sigma^2.}
\end{equation}
If an admitted residual source has no explicit Jordan $h$ dependence,
is on its source Euler domain, and its declared source-field transport
has vanishing first $u$ jet at $u=0$, its local metric-trace contact is
\begin{equation}
\label{p3f:eq:v2210-lit-frozen-trace-contact}
 \boxed{\frac{J_{u,u}^{\rm met}}{\ell_\Phi^2}
 =-\frac{T_{r,E}}{6M^2}
 =\frac{\rho_{r,E}-3p_{r,E}}{6M^2},\qquad
 T_{r,E}=F^{-2}T_{r,J}.}
\end{equation}
For a fixed-source-field metric partial Hessian the same formula is the
metric contribution only.  A nonzero transported source-field first
jet supplies additional mixed and quadratic contacts.  Those contacts,
explicit scalar couplings, preparation/constraint contacts, derivative
sources and dynamical relaxation remain in the complete blocks of
Eq.~\eqref{p3f:eq:v2210-lit-complete-transverse-block}; this local
contact is not their full response.  In the separately licensed affine
terminal quotient $F=1$, $P=C$ and with no omitted contacts, the
potential plus local trace coefficient agrees with
Eq.~\eqref{p3f:eq:v2202-terminal-local-mass}.  Removing the scalaron
on that quotient does not carry an active $\dot\sigma^2$ pump into it.
\end{corollary}
\begin{proof}
At fixed $F$ the native potential is
$W_E=[M^2P(F)/2+V_\Phi(\ell_\Phi\sqrt F\,u)]/
[F^2(1-u^2)^2]$.  Since $(1-u^2)^{-2}=1+2u^2+O(u^4)$,
its quadratic coefficient is
$M^2P(F)/F^2-\Lambda_g^2\ell_\Phi^2/F$.
Twice that coefficient divided by $\ell_\Phi^2=6M^2$ gives the
first identity; the quartic has zero second derivative at zero.
For the stated metric-only source direction,
$g_J=[F(1-u^2)]^{-1}g_E$ has first $u$ derivative zero and second
derivative $2g_J$ at zero.  The on-source-Euler metric variation is
$\delta S_r=\tfrac12\int\sqrt{-g_J}T_r^{\mu\nu}\delta g_{J,\mu\nu}$.
Its second $u$ variation is therefore
$\int\sqrt{-g_E}F^{-2}T_{r,J}(\delta u)^2$; terms linear in second
source-field jets vanish on the source Euler domain, while the first
source-field jets vanish by the stated transport hypothesis.  Comparing with the convention
$\delta S_r=-\int\sqrt{-g_E}J_u\delta u$ gives
$J_{u,u}^{\rm met}=-T_{r,E}$.  The terminal comparison follows by
substitution in its own quotient domain.
\end{proof}

\begin{corollary}[Closed-curvature correction to the kinetic instability test]
\label{p3f:cor:v2210-lit-closed-kinetic-test}
On the preceding background suppose the complete restriction has the
Einstein equation with constant coefficient $M^2$, with the displayed
parallel scalar and a residual Hilbert source counted once.  Its
Raychaudhuri equation is
\[
 -2M^2(\dot H_E-k/a_E^2)
 =\dot\varphi_\parallel^2+\rho_r+p_r,\qquad k=1.
\]
For $H_E\ne0$ define
$\epsilon_{\rm kin}=\dot\varphi_\parallel^2/(2M^2H_E^2)$ and
$\epsilon_H=-\dot H_E/H_E^2$.  Then
\begin{equation}
\label{p3f:eq:v2210-lit-closed-kinetic-epsilon}
 \boxed{\epsilon_{\rm kin}=
 \epsilon_H+\frac{k}{a_E^2H_E^2}
 -\frac{\rho_r+p_r}{2M^2H_E^2},\qquad
 -\dot\sigma^2=-\frac{\epsilon_{\rm kin}H_E^2}{3}.}
\end{equation}
The replacement by $-\epsilon_HH_E^2/3$ additionally requires negligible
closed-curvature and residual-enthalpy corrections.  At a turning point
$H_E=0$ the finite expression $-\dot\sigma^2$ is used directly.
\end{corollary}
\begin{proof}
The complete Einstein lapse and scale equations give the displayed
Raychaudhuri combination.  Divide it by $2M^2H_E^2$ and use
$\dot\varphi_\parallel^2=\ell_\Phi^2\dot\sigma^2$ and
$\ell_\Phi^2=6M^2$.  No slow-roll approximation enters this identity.
\end{proof}

\begin{corollary}[Canonical finite-harmonic pump and its source conjugation]
\label{p3f:cor:v2210-lit-canonical-pump}
For an actually admitted scalar harmonic
$-\Delta_\gamma Y_L=\lambda_LY_L$,
$\lambda_L=L(L+2)$, put $X_L=a_E^{3/2}Q_{\perp,L}$.
Conjugating the complete physical equation, rather than only its
potential, gives the intrinsic canonical frequency
\begin{equation}
\label{p3f:eq:v2210-lit-canonical-frequency}
 \boxed{\omega_{L,0}^2=
 \frac{\lambda_L}{a_E^2}+
 \frac{W_{E,uu}}{\ell_\Phi^2}-\dot\sigma^2
 -\frac32\dot H_E-\frac94H_E^2.}
\end{equation}
Every full contact or memory operator $\mathcal P^R$ is conjugated as
$a_E^{3/2}\mathcal P^Ra_E^{-3/2}$ and the effective drive as
$a_E^{3/2}j_{\perp,\rm eff}$.  Only harmonics in the actual physical
support/incidence are used.  This change of variables creates no state,
additional physical energy, independent oscillation frequency or
microscopic carrier.
\end{corollary}
\begin{proof}
Differentiate $Q=a_E^{-3/2}X$ twice.  The first derivative of $X$
cancels against $3H_E\dot Q$, leaving
$\ddot X-(3\dot H_E/2+9H_E^2/4)X$ in that combination.
Operator conjugation keeps the actual order and endpoints of every
source kernel.  This is a canonical presentation of the same response,
not a second energy allocation.
\end{proof}

\begin{proposition}[Geometric exchange cancels in the constituent work]
\label{p3f:prop:v2210-lit-transfer-energy}
On the displayed two-derivative homogeneous restriction, with
$W_E=W_E(\sigma,u)$ and the actual residual force covectors, set
$\rho_{g\Phi}=T_\sigma+T_u+W_E$ and
$p_{g\Phi}=T_\sigma+T_u-W_E$.  The exact identity is
\begin{equation}
\label{p3f:eq:v2210-lit-transfer-work}
 \boxed{\dot\rho_{g\Phi}+3H_E(\rho_{g\Phi}+p_{g\Phi})
 =-J_\sigma\dot\sigma-J_u\dot u.}
\end{equation}
The residual source supplies the compensating work in the complete
same-action Ward identity.  If the potential has admitted explicit time
dependence, its explicit derivative is also retained on the right.
Positive instantaneous $Q_g$ or increasing $T_u$ therefore does not by
itself establish a positive reference-subtracted completed reservoir.
\end{proposition}
\begin{proof}
Add the two equations of
Eq.~\eqref{p3f:eq:v2209-geometric-work}; $Q_g$ cancels.
Add $\dot W_E=W_{E,\sigma}\dot\sigma+W_{E,u}\dot u$.
Since $\rho_{g\Phi}+p_{g\Phi}=2(T_\sigma+T_u)$ this gives the
identity.  The source Ward work is inherited from the same action.
The frozen binding and pulse identities of
Theorem~\ref{p3f:thm:v2207-monotone-loading-sign} continue to govern
their stated complete energy increments; no isolated positive term
overrides them.
\end{proof}

\subsubsection{Boundary incidence and metric/source derivatives of the same parent}

\begin{theorem}[Complete finite boundary pullback and moving-contact derivatives]
\label{p3f:thm:v2210-lit-boundary-jets}
On an actual regular gauge/constraint quotient let $S(b,x;m)$ be the
claim-complete stationary parent, including its physical boundary and
source terms.  Here $m$ denotes already retained metric/source controls
and the initial preparation is fixed in variation.  Transport any
moving domain and pairing to the same physical chart, retaining its
embedding, measure and frame derivatives.  Suppose $S_x=0$ defines a
local interior branch with $C=S_{xx}$ invertible.  Put
$\Gamma(b;m)=S(b,x_*(b;m);m)$ and
\[
 A=S_{bb},\quad B=S_{bx},\quad D=S_{xb},\qquad
 K=\Gamma_{bb}=A-BC^{-1}D.
\]
For the actual physical incidence $b=\beta(q;m)$, its complete
second variation is
\begin{equation}
\label{p3f:eq:v2210-lit-nonlinear-boundary-pullback}
 \boxed{D_q^2(\Gamma\circ\beta)[v,w]
 =K[D_q\beta v,D_q\beta w]
 +\Gamma_b[D_q^2\beta[v,w]].}
\end{equation}
Thus the simple pullback $(D_q\beta)^*K(D_q\beta)$ is valid when
the incidence is linear, or when the complete stationary force
$\Gamma_b$ annihilates its second derivative; it is not a general rule
for a nonlinear, nonstationary boundary map.

For a linear physical incidence in the transported pairing, its complete
control derivative is
\begin{align}
\label{p3f:eq:v2210-lit-schur-control-jets}
 K_{,m}&=A_{,m}-B_{,m}C^{-1}D-BC^{-1}D_{,m}\notag\\
 &\quad+BC^{-1}C_{,m}C^{-1}D,\notag\\
 (\beta^*K\beta)_{,m}&=\beta_{,m}^*K\beta+\beta^*K_{,m}\beta+
 \beta^*K\beta_{,m}.\nonumber
\end{align}
The derivatives are total derivatives along the actual stationary
family.  For a nonlinear incidence the derivative of the force term in
Eq.~\eqref{p3f:eq:v2210-lit-nonlinear-boundary-pullback} is also retained.
\end{theorem}
\begin{proof}
Interior stationarity gives the envelope identity
$d\Gamma=S_b\,db+S_m\,dm$ and
$x_{*,b}=-C^{-1}D$.  Differentiating gives the Schur Hessian.
The ordinary second-order chain rule for $\Gamma\circ\beta$ gives
the force/contact term.  Finally
$(C^{-1})_{,m}=-C^{-1}C_{,m}C^{-1}$ gives the first control identity,
and the product rule gives the second.  Transporting to the actual
physical chart is what ensures that boundary motion and pairing
variation belong to these derivatives.  No physical boundary or worldtube
is constructed by the algebra alone.  This distinction is also required
by the covariant boundary variational treatment
\cite{LitDYNHarlowWu2020}.
\end{proof}

The smooth non-null $FK$ completion of
Theorem~\ref{p3f:thm:v169-complete-gravity-action} already supplies its
particular action boundary \cite{DyerHinterbichler2009BoundaryTerms}.
The theorem above neither extends it silently to null/corner data nor
replaces the actual moving $K48$ incidence.  Lapse and scale stresses
are first variations of the complete $\Gamma$; differentiating a
stationary Hessian alone does not recover its zero-order action or all
temporal metric insertions.


\begin{theorem}[Same-action positive-remainder criterion for grade-four susceptibility]
\label{p3f:thm:v286-same-action-grade-four-survival}
Fix an already occurring \emph{complete} action and one regular
finite physical configuration/constraint quotient at a stationary
background. In the same transported metric-curvature tangent $u$
and interior controls $z\in\mathbb R^m$, suppose the full Hessian
splits as $H_\star=H_4+H_{\rm rem}$, where
\begin{enumerate}[label=\textup{(\roman*)}]
\item $H_4$ is the exact affine finite-heat Hessian of
Theorem~\ref{p3b:thm:v286-finite-heat-incidence-gram}, with
operator tangents $B_4,C_{4,i}$ arising from a generated preterminal
$P_4$ insertion into that \emph{same} action;
\item $H_{\rm rem}$ is positive semidefinite on the actual joint
metric/interior quotient, \emph{after} including density, localizer,
connection, boundary, operator-square and source contacts; and
\item the complete interior block $(H_\star)_{zz}$ is positive
definite, while $\{C_{4,i}\}$ are independent on the retained
$H_4$ interior quotient.
\end{enumerate}
Then the actual stationary reduced curvature Hessian satisfies
\begin{equation}
\label{p3f:eq:v286-positive-remainder-bound}
\boxed{\chi_\star:=(H_\star)_{uu}
 -(H_\star)_{uz}(H_\star)_{zz}^{-1}(H_\star)_{zu}
 \ \ge\ \chi_4
 =\operatorname{dist}_{\mathcal G_{A_4}}^2
 (B_4,\operatorname{span}_{\mathbb R}\{C_{4,i}\})\ge0.}
\end{equation}
In particular $B_4$ transverse to the actual eliminated-source span
is a sufficient certificate for a nonzero \emph{same-action}
curvature susceptibility, \emph{once} (i)--(iii) have been verified.
For a physical metric variation $q\mapsto u(q)$, the factor
$(D_qu)^2$ accompanies this bound when the pullback is affine or
the complete stationary force annihilates $D_q^2u$; otherwise the
force/contact term of
Theorem~\ref{p3f:thm:v2210-lit-boundary-jets} must be kept.
This criterion certifies no physical $P_4$ insertion, positivity of a
generic Lorentzian Hessian, or universal equality of $\chi_\star$
with a separately supplied $f_{RR}$.
\end{theorem}
\begin{proof}
The actual reduced stationary Hessian is the Schur complement of
$H_\star$ by the same-action boundary/constraint theorem
\ref{p3f:thm:v2210-lit-boundary-jets}. For a unit metric tangent,
its variational characterization is
\[
 \chi_\star=\min_z\bigl[H_4((1,z),(1,z))
 +H_{\rm rem}((1,z),(1,z))\bigr].
\]
By (ii) this is at least $\min_z H_4((1,z),(1,z))$, which is exactly
$\chi_4$ by Theorem~\ref{p3b:thm:v286-finite-heat-incidence-gram}.
Transversality is equivalent to $\chi_4>0$. The physical pullback
qualification is the full second-order chain rule in the cited
boundary theorem. None of these purely mathematical inequalities
constructs its physical incidence hypotheses.
\end{proof}

\begin{theorem}[Three exact action-level coefficients for one curvature susceptibility]
\label{p3f:thm:v287-three-jet-frr-consistency}
Fix a single four-dimensional metric response density
$Z\sqrt{|g|}[f(R)-2\Lambda_{\rm ren}]$ with $Z>0$, $f\in C^2(I)$,
and fixed $\Lambda_{\rm ren}$.  The following three \emph{different}
local differential tests all recover the \emph{same} function
$\chi(R)=f_{RR}(R)$ on their respective domains:
\begin{enumerate}[label=\textup{(\roman*)}]
\item On an admitted Lorentzian closed-FLRW jet chart,
with $L_Z=2\pi^2ZNa^3[f(R)-2\Lambda_{\rm ren}]$,
\begin{equation}\label{p3f:eq:v287-flrw-chi-witness}
\boxed{\frac{N^3}{72\pi^2Za}\frac{\partial^2L_Z}
 {\partial(\ddot a)^2}=\chi(R).}
\end{equation}
\item At a constant-curvature Lorentzian background $R_0\in I$
which satisfies the \emph{complete} background equations,
the coefficient of $\Box_0\delta R$ in the linearized pure gravitational
trace residual
$\mathcal T_Z=Z[Rf_R-2f+4\Lambda_{\rm ren}+3\Box f_R]$ is
$3Z\chi(R_0)$. Thus, writing $[\Box\delta R]\,\delta\mathcal T_Z$
for this coefficient rather than an operator division,
\begin{equation}\label{p3f:eq:v287-trace-chi-witness}
\boxed{\frac1{3Z}[\Box_0\delta R]\,\delta\mathcal T_Z
       =\chi(R_0).}
\end{equation}
This test retains the separately varied source trace and does not
manufacture a constant-curvature background.
\item On the \emph{separately admitted} Riemannian constant-curvature
comparison domain with the same $f$ and $Z$, the conformal Hessian
$\mathcal H_\sigma$ of
Theorem~\ref{p3f:thm:v199-rgas-conformal-hessian}, for
$P=-\bar\Box$, has highest $P^2$ coefficient $36Z\chi(R_0)$:
\begin{equation}\label{p3f:eq:v287-conformal-chi-witness}
\boxed{\frac1{36Z}[P^2]\,\mathcal H_\sigma=\chi(R_0).}
\end{equation}
\end{enumerate}
These are \emph{coefficient extractions}, not assertions that a given
finite heat block, a Wick rotation, a stationary Schur quotient, or an
observed scalaron has been constructed. In particular a positive
$\chi$ proves positivity of the displayed \emph{local} coefficients
for $N,a,Z>0$ but not the positivity of a full Lorentzian gravitational
Hessian. If a proposed common-action realization gives different
normalized values from (i)--(iii) after its respective contacts and
normalizations are included, it fails that common-$f(R)$ realization
test on the overlapping mathematical domain.
\end{theorem}
\begin{proof}
The FLRW identity is
Theorem~\ref{p3e:thm:v287-s3-acceleration-hessian};
$-2\Lambda_{\rm ren}$ has zero curvature Hessian.
The constant-background trace variation is
Theorem~\ref{p3f:thm:trace-linearized-scalaron-mass} multiplied by $Z$;
the fixed cosmological term affects the zero-order background condition
but not the coefficient of $\Box_0\delta R$.
The conformal Hessian in
Theorem~\ref{p3f:thm:v199-rgas-conformal-hessian} contains
$4Z\chi(R_0)(R_0-3P)^2$, whose $P^2$ coefficient is
$36Z\chi(R_0)$. All operations use the same normalized function,
though their charts and occurrence hypotheses are distinct.
\end{proof}


\begin{proposition}[Finite Gaussian determinant jets when the integration is licensed]
\label{p3f:prop:v2210-lit-gaussian-jets}
If the existing quantization and measure actually license a real finite
positive Gaussian interior block $C(m)>0$, write its un-subtracted
determinant coordinate $w_C=\tfrac12\log\det C$ in the units fixed by
the actual measure.  Its relative contribution is
$W_C=w_C-w_{C,\rm ref}$, with all measure factors kept.  The control
derivatives of the un-subtracted coordinate are
\begin{equation}
\label{p3f:eq:v2210-lit-gaussian-control-jets}
 \boxed{D_iw_C=\tfrac12\operatorname{Tr}(C^{-1}C_i),\qquad
 D_iD_jw_C=\tfrac12\operatorname{Tr}
 (C^{-1}C_{ij}-C^{-1}C_jC^{-1}C_i).}
\end{equation}
Every derivative of the actual measure or reference is retained as well.
These are derivatives of that already licensed integration, not an
additional independently weighted action.  A continuum elliptic
Neumann--Dirichlet determinant reduction additionally consumes its actual
operator, boundary conditions, zero-mode treatment and regularization
\cite{LitDYNBarvinskyNesterov2006}; finite matrix algebra does not prove
those hypotheses or a Lorentzian continuation.
\end{proposition}
\begin{proof}
The finite positive Gaussian integral is proportional to
$(\det C)^{-1/2}$ in its declared measure.  Jacobi's determinant identity
and the inverse derivative give the displayed formulas.  Constants,
measure factors and the reference cancel only where their identical
allocation is actually proved.  A singular endpoint invalidates the
positive-invertible Gaussian hypothesis and requires the retained full
relation or a separately controlled limit.
\end{proof}

\begin{corollary}[Exact scalar response requires both operator and source matching]
\label{p3f:cor:v2210-lit-overlap-test}
On an admitted common invertible domain let the complete scalar/interior
operator be
$\bigl(\begin{smallmatrix}K_f+\Delta K_{\rm dir}&B\\D&C\end{smallmatrix}\bigr)$,
where $K_f$ is the inherited $f_{2\mid7}$ scalar operator.  Its scalar
operator equals $K_f$ precisely when
\begin{equation}
\label{p3f:eq:v2210-lit-overlap-source-test}
 \boxed{\Delta K_{\rm dir}=BC^{-1}D.}
\end{equation}
Equality of the responses for all admitted sources further requires
equality of the effective source incidence
$j_f=j_\perp-BC^{-1}j_x$ in the same physical pairing.  Coincidence of
an Einstein coefficient, a determinant multiplicity or a pole location
proves neither equality.  Any licensed determinant jets, mixed metric
jets and boundary pullback contacts belong to the same test.
\end{corollary}
\begin{proof}
Eliminate the actual interior block before comparison.  The retained
equation is
$(K_f+\Delta K_{\rm dir}-BC^{-1}D)Q=
j_\perp-BC^{-1}j_x$.  Equality of operators gives the first criterion;
on the common invertible domain equality of every sourced solution also
requires the same effective drive.  This is the action-overlap
requirement of Proposition~\ref{p3f:prop:v1164-scalar-overlap-condition}
with the source incidence made explicit.  It is consistent with the
ordinary Dirac higher-curvature diagnostic
\cite{LitDYNMistryPinzulRachwal2020} and licenses no additive completion.
\end{proof}

\subsubsection{Regularity, causality and the remaining native execution}

\begin{proposition}[Regular analytic elimination cannot create the exact K16 kink]
\label{p3f:prop:v2210-lit-analytic-projection}
On the common analytic fixed-rank trace domain of
Proposition~\ref{p3f:prop:v1165-analytic-k16-obstruction}, suppose the
subsequent incidence, stationary interior branch and coefficient
transport are analytic, with an invertible interior Hessian and uniform
analytic trace control.  Any finite sequence of those regular analytic
eliminations, pullbacks and positive constant coefficient transports
has analytic scalar restriction.  It therefore cannot reproduce the
nonzero third-jet jump of
Section~\ref{p3f:sec:v2201-derived-d14} across a transverse analytic
$R=C_2$ crossing.  The actual depth projection, domain change, critical
elimination or other certified nonanalytic map must be exhibited in
the native parent-to-action word.
\end{proposition}
\begin{proof}
The initial heat trace is analytic on the stated domain.  The analytic
implicit-function theorem gives an analytic stationary interior map
when its Hessian is invertible.  Composition, inversion on that regular
domain and constant transport preserve analyticity.  Repeating these
steps preserves it.  Along a transverse crossing the existing scalar
coefficient instead has
$[\partial_s^3f(R(s))]=[f_{RRR}](R')^3\ne0$, contradicting analytic
restriction.  The source-free $F$-flux matching theorem is unchanged;
this result locates the missing microscopic reduction, not a new
surface impulse.
\end{proof}

\begin{proposition}[Causal support is retained by the admitted Schur composition]
\label{p3f:prop:v2210-lit-causal-support}
On one common physical initial-value domain, if the mixed source blocks
are local or retarded and the actual inverse $C_R^{-1}$ is retarded,
then $BC_R^{-1}D$ has retarded support.  If the resulting scalar
operator possesses its admitted retarded inverse, the response in
Eq.~\eqref{p3f:eq:v2210-lit-complete-transverse-schur} depends only on
the inherited initial data and past source.  This statement does not
construct the inverse, prove an unrestricted continuum Cauchy theorem,
or certify the interacting characteristic/front speed.
\end{proposition}
\begin{proof}
Composition convolves the admitted kernels with time ordering
$t\ge t_1\ge t_2\ge s$; local derivative contacts preserve support.
Hence no future source enters the composed kernel.  Application of
the assumed retarded scalar inverse gives the claimed response.
Existing finite Cauchy and Volterra results continue to supply inverses
only on their proved regular charts.  Euclidean positivity alone is not
the missing causal certificate.
\end{proof}

\begin{proposition}[Static metric integrability test for a scalar envelope]
\label{p3f:prop:v2210-lit-static-integrability}
The native static metric readout $g_i=c^2\nabla_i\log N$ is locally a
gradient.  If an optional scalar reduction proposes
$\mathbf g_{\rm trial}=\nu(y)\mathbf g_b$ with
$\mathbf g_b=-\nabla\Phi_b$, its gradient consistency requires
\begin{equation}
\label{p3f:eq:v2210-lit-envelope-integrability}
 \nabla\times\mathbf g_{\rm trial}
 =\nabla\nu(y)\times\mathbf g_b=0.
\end{equation}
On a simply connected regular local domain this closed-one-form
condition is sufficient for a potential.  It is automatic on a
one-dimensional symmetry reduction, but need not hold for a general
nonspherical source.  This is a diagnostic of an already proposed
scalar representation; it adds no interpolation profile and does not
replace the native fixed-environment action/Schur readout.
\end{proposition}
\begin{proof}
The curl of a gradient vanishes and the product rule gives the equation.
The local potential statement follows from exactness of closed one-forms
on the declared domain.  The distinction between an algebraic
acceleration envelope and a conservative field is also explicit in
the primary variational MOND construction \cite{Milgrom2010QUMOND}.
\end{proof}

\subparagraph{Narrowed active obligations.}
The regular analytic compiler, canonical geodesic pump, full contact
normalization, closed-curvature kinetic identity, finite boundary-chain
derivatives and supplied-kernel causal composition are now explicit
tests on their stated native domains.  Their unresolved inputs are
the actual common-action projection and all consumed source/metric
jets; a surviving physical scalar and its populated source state;
the actual moving boundary incidence and regular inverses; and the
nonanalytic microscopic map at K16.  A quantitative local approximation
still requires its actual higher-variation error bound.  Terminal
decoupling still requires the stated endpoint solution and controlled
mixed-memory limit, rather than only its divergent algebraic mass.
Native nonspherical RAR and occupied finite-source histories still require
execution of the actual operator/state/readout.  Exact interacting
quantum reconstruction and null/corner boundary charges keep their
separate domains.  The inherited clocked characteristic equality and
local Noether/Cauchy results are not reopened.  Neither the geometric
pump nor the finite determinant identity supplies capture, cold stress,
an abundance, a new species or a globally equal action by itself.


\subsection{Exact conformal heat loads and the first raw tidal Euler response}
\label{p3f:sec:v2212-conformal-tidal-response}

The following results concern the existing pure-metric common heat
constituent and its supplied finite local real form.  Their mathematical
domains do not require identifying that constituent with the reconstructed
Born--Legendre scalar action.  In particular they do not consume the
scalar-response certificates \texttt{ACD\_H} or \texttt{EXJET7} as a
substitute for evaluating the heat operator.  A conformal metric tangent,
an exact Euclidean load, and a matched physical scalaron force remain
different operations.  The whole-action overlap of
Proposition~\ref{p3f:prop:v1164-scalar-overlap-condition} is retained.

\subsubsection{Fixed-pairing conformal insertions and their actual localizer}

\begin{proposition}[Exact conformal Dirac and heat loads]
\label{p3f:prop:v2212-conformal-heat-load}
Let $g_J=e^{-q}g_E$, with $F=e^q>0$, on a common smooth four-dimensional
spin domain.  Use the pure-metric Dirac constituent, the canonical
spin-bundle identification and the unitary $L^2$ transport to a fixed
Hilbert pairing.  Assume the common differentiable trace domain of
Proposition~\ref{p3f:prop:v1164-heat-hessian-no-insertion} and
Proposition~\ref{p3f:prop:v2206-mixed-heat-jet}.  For a fixed cutoff put
$X=D^2/\Lambda_g^2$ and let $Q=\delta q$ be a smooth multiplication
tangent.  Finite matrices are a special case only when their actual
family realizes these conformal jets.  If $Q$ leaves a retained spectral
window, its transported embedding and complementary contacts remain;
the formula is not asserted for an arbitrary compressed $D$.  Then
\begin{equation}
\label{p3f:eq:v2212-conformal-dirac-insertion}
 \boxed{\overline D(q)=e^{q/4}D_Ee^{q/4},\qquad
 D_q[Q]=\tfrac14\{Q,D\},\qquad
 X_q[Q]=\frac{\{D,\{Q,D\}\}}{4\Lambda_g^2}.}
\end{equation}
On the unlocalized closed trace domain, with unit insertion, cyclicity
gives the exact load, also for nonconstant $Q$,
\begin{equation}
\label{p3f:eq:v2212-unlocalized-conformal-load}
 \boxed{\delta_q\operatorname{Tr}(e^{-X})[Q]
              =-\operatorname{Tr}(QXe^{-X}).}
\end{equation}
For the actual transported insertion $W$, including its localizer,
pairing, weight and projector contacts, define
\[
 R_W=\int_0^1e^{-(1-u)X}We^{-uX}\,du,\qquad
 S_W=\operatorname{Tr}(We^{-X}).
\]
Its complete first conformal load is
\begin{equation}
\label{p3f:eq:v2212-localized-conformal-load}
 \boxed{(S_W)_q[Q]=\operatorname{Tr}(W_qe^{-X})
 -\frac{\operatorname{Tr}\bigl(Q\{D,\{D,R_W\}\}\bigr)}
 {4\Lambda_g^2}.}
\end{equation}
An independently held mass, odd field, gauge constituent or source
preparation retains its own actual conformal jets; the pure-metric
identity does not transform those data automatically.
\end{proposition}
\begin{proof}
Contraction of the conformally transformed spin connection gives
$D_J=e^{5q/4}D_Ee^{-3q/4}$ under the pointwise spin identification.
Since $d\mu_J=e^{-2q}d\mu_E$, the unitary pairing transport multiplies
spinors by $e^{-q}$.  Conjugating $D_J$ gives the first identity, whose
ordinary derivative gives the anticommutator insertions.  Apply the
already proved Duhamel first derivative to $S_W$ and move the trace
factors cyclically: its operator term is $-\operatorname{Tr}(R_WX_q)$.
Twice moving $D$ gives Eq.~\eqref{p3f:eq:v2212-localized-conformal-load}.
For $W=I$, $R_W=e^{-X}$ commutes with $D$; the four terms in the nested
anticommutator then give
$\operatorname{Tr}(Q\{D,\{D,e^{-X}\}\})
=4\operatorname{Tr}(QD^2e^{-X})$.  This proves the unlocalized load
without assuming that $Q$ commutes with $D$.  The trace hypotheses
justify these cyclic moves and differentiations.
\end{proof}

\begin{corollary}[Complete conformal--shape contacts]
\label{p3f:cor:v2212-conformal-shape-contacts}
Hold $Q$ in the transported reference chart and let $\alpha$ be an
actual metric-shape coordinate.  At fixed cutoff,
\[
 D_{q\alpha}[Q]=\tfrac14\{Q,D_\alpha\},\qquad
 X_{q\alpha}=\Lambda_g^{-2}
       (\{D_q,D_\alpha\}+\{D,D_{q\alpha}\}).
\]
Writing $Z=e^{-X}$, differentiation of the preceding load gives
\begin{equation}
\label{p3f:eq:v2212-localized-mixed-load}
\begin{split}
 (S_W)_{q\alpha}[Q]={}&\operatorname{Tr}(W_{q\alpha}Z+W_qZ_\alpha)\\
 &-\frac1{4\Lambda_g^2}\operatorname{Tr}Q\bigl(
 \{D_\alpha,\{D,R_W\}\}+\{D,\{D_\alpha,R_W\}\}\\
 &\hspace{39mm}+\{D,\{D,(R_W)_\alpha\}\}\bigr).
\end{split}
\end{equation}
Here $(R_W)_\alpha$ retains $W_\alpha$ and both ordered exponential
derivatives.  Equivalently, for $\mathsf E(A)=e^A$,
\[
 (R_W)_\alpha=D^2\mathsf E_{-X}[-X_\alpha,W]
                         +D\mathsf E_{-X}[W_\alpha].
\]
Changing $Q$, the cutoff or the transported domain adds its actual
product-rule contacts.  A moving finite projector must be transported
with its embedding; $Pe^{-X}P=e^{-PXP}$ is not asserted for a
noninvariant compression.
\end{corollary}
\begin{proof}
Differentiate Eq.~\eqref{p3f:eq:v2212-localized-conformal-load} and
$D_q[Q]$ in the fixed reference chart.  These are the product rule and
the already established ordered mixed heat derivative, not a second
response prescription.  Since $R_W=D\mathsf E_{-X}[W]$, its displayed
derivative follows directly.  These identities retain the seagull
$X_{q\alpha}$ in
Eq.~\eqref{p3f:eq:v2206-mixed-heat-jet}; neither a positive heat kernel
nor a nonzero mixed block determines its signed physical work.
\end{proof}

\begin{proposition}[Algebraic Weyl weights under the same conformal map]
\label{p3f:prop:v2212-weyl-conformal-weight}
For a derivative-free scalar contraction $W_n$ of $n$ Weyl tensors,
with all indices contracted by the metric, the four-dimensional map
$g_E=Fg_J$, $F>0$, obeys
\begin{equation}
\label{p3f:eq:v2212-weyl-conformal-weight}
 \boxed{\sqrt{-g_J}\,W_n(g_J)
      =\sqrt{-g_E}\,F^{n-2}W_n(g_E).}
\end{equation}
Thus a constant-coefficient $C^2$ density has no direct $q$ vertex at
fixed $g_E$, while an actual cubic Weyl constituent carries a factor
$F$.  Terms with curvature derivatives or $RC^2$ retain the additional
conformal derivative contacts and cannot be replaced by this weight.
\end{proposition}
\begin{proof}
The mixed Weyl tensor $C^a{}_{bcd}$ is conformally invariant.  Its
metric contractions give $W_n(g_J)=F^nW_n(g_E)$; the volume gives
$F^{-2}$.  Weyl covariance is pointwise also for nonconstant $F$, so
no constant-field assumption is used.  This kinematic identity neither
selects a scalaron state nor proves the scalar-action projection.
\end{proof}

\subsubsection{A fixed raw six-derivative response on Einstein metrics}

\begin{proposition}[Ordinary-Dirac sixth-jet Euler representative and coefficient]
\label{p3f:prop:v2212-dirac-a6-einstein}
Use the unprojected ordinary torsion-free pure-metric Dirac heat
coefficient, its fixed exponential profile and the finite-local clock
domain of Theorem~\ref{p3f:thm:finite-local-bosonic-clock-real-form}.
Let $P_6$ have the compact normalization of equations (3.6)--(3.7) of
\cite{LitDYNMistryPinzulRachwal2020}, retaining their factor 1008.
In the native Riemann convention, its interior first variation on any
four-dimensional constant-scalar Einstein jet is represented by
\begin{equation}
\label{p3f:eq:v2212-a6-einstein-representative}
 \boxed{P_{6,E}\simeq_\delta
       \frac{31}{2160}R^3-\frac16RC^2+\frac1{15}J_C,\qquad
 J_C=C_{ab}{}^{cd}C_{cd}{}^{ef}C_{ef}{}^{ab}.}
\end{equation}
Here $\simeq_\delta$ means equality of the evaluated interior Euler
tensors, not permission to impose the background before varying.
The raw Lorentz constituent and its inherited coefficient are
\begin{equation}
\label{p3f:eq:v2212-a6-raw-coefficient}
 \boxed{S_{6,L}^{\rm raw}=-z\int\sqrt{-g}\,P_6,\qquad
 z=\frac{\kappa_J}{84\Lambda_g^4}
   =\frac1{64512\pi^4U}>0.}
\end{equation}
This evaluates the raw coefficient placeholder of
Proposition~\ref{p3f:prop:v1164-dirac-a4-a6-firewall}: on this same
Lorentz constituent and polynomial normalization, $\zeta_6=-z$.
It is not an independently added completion of $f_{2\mid7}$.
\end{proposition}
\begin{proof}
Write $S_{\mu\nu}=\operatorname{Ric}_{\mu\nu}-Rg_{\mu\nu}/4$.
The full Weyl-basis coefficient, with its Riemann sign converted to the
native convention, contains
\[
\begin{aligned}
 P_6={}&-\tfrac1{10}R\Box R-\tfrac32C\Box C-\tfrac2{135}R^3\\
 &+\tfrac13R\operatorname{Ric}^2-\tfrac{13}{15}
             \operatorname{tr}(\operatorname{Ric}^3)+\tfrac7{12}RC^2\\
 &-\tfrac{23}{5}\operatorname{Ric}_{\mu\nu}
          \operatorname{Ric}_{\rho\sigma}C^{\mu\rho\nu\sigma}
          -\tfrac{133}{30}J_C.
\end{aligned}
\]
Terms at least quadratic in $S$ or $\nabla\operatorname{Ric}$ have zero
first variation on the stated jet; $\int R\Box R$ does too.  The
four-dimensional identity
$C_{\mu\alpha\beta\gamma}C_\nu{}^{\alpha\beta\gamma}
=g_{\mu\nu}C^2/4$ removes the possible linear $SC C$ term.  The
integrated generalized Gauss--Bonnet identity then gives, at first
variation on this same Einstein jet,
\[
 \int C\Box C\simeq_\delta
     \int(\tfrac12RC^2-3J_C).
\]
The remaining scalar coefficient is
$-2/135+1/12-13/240=31/2160$, the $RC^2$ coefficient is
$-3/4+7/12=-1/6$, and the cubic Weyl coefficient is
$9/2-133/30=1/15$.  This proves the Euler representative while
retaining the derivative contacts before specialization.  All integrations
by parts use compact interior variations or the actual transported
boundary problem, not an artificial reflecting surface.

For a common traced multiplicity $N$, the exponential moments are
$f_2=f_6=1$.  The original heat coefficient and the compact polynomial
therefore give
$\kappa_J=N\Lambda_g^2/(192\pi^2)$ and
$z=N/(16\pi^2\,1008\Lambda_g^2)$: the factor 1008 is in the
denominator when converting back from that compact normalization.
Their ratio is $1/(84\Lambda_g^4)$.  Substitute the already derived
$\kappa_J=12U$, $\Lambda_g^2=96\pi^2U$ for the final value.
Under $I_E\mapsto-iS_L$, the finite-local clock theorem gives
$S_L=-\sum_a c_a\int P_{a,L}\,d\mu_L$ and fixes the displayed minus
sign.  It does not supply an exact global nonlocal Lorentzian
continuation.
\end{proof}

\begin{theorem}[Exact raw Kottler Euler ports and their competing signs]
\label{p3f:thm:v2212-a6-kottler-ports}
On the preceding domain take any constant $R$, $m>0$, $r>0$, and
\[
 H=1-2m/r-Rr^2/12,\qquad
 ds^2=-Hdt^2+H^{-1}dr^2+r^2d\Omega^2.
\]
No native stationary-root or physical-source identification is assumed.
For comparison units $M^2=U/4$ define
\[
 \widehat E_{\mu\nu}
 =\frac2{M^2\sqrt{-g}}\frac{\delta S_{6,L}^{\rm raw}}
                              {\delta g^{\mu\nu}},\qquad
 \mathcal C_6=\frac{16zm^2}{5M^2r^9},\quad
 Z_6=-\frac{31zR^3}{4320M^2}.
\]
This tensor is a raw-action comparison, not an already matched physical
residual.  Its exact mixed components are
\begin{equation}
\label{p3f:eq:v2212-a6-kottler-components}
\boxed{\begin{aligned}
 \widehat E^t{}_t&=\mathcal C_6
                  (135r-298m-\tfrac{31}4Rr^3)+Z_6,\\
 \widehat E^r{}_r&=\mathcal C_6
                  (98m-63r+\tfrac{35}4Rr^3)+Z_6,\\
 \widehat E^\vartheta{}_\vartheta
 =\widehat E^\varphi{}_\varphi
       &=\mathcal C_6(189r-442m-\tfrac{37}4Rr^3)+Z_6.
\end{aligned}}
\end{equation}
They obey the exact covariant Ward identity.  In ingoing coordinates
$ds^2=-Hdv^2+2dv\,dr+r^2d\Omega^2$, the affine tangent
$k=-\partial_r$ has the regular null projection
\begin{equation}
\label{p3f:eq:v2212-a6-kottler-null}
 \boxed{\widehat E_{kk}
   =\frac{\widehat E^r{}_r-\widehat E^t{}_t}{H}
   =-198\mathcal C_6r
   =-\frac{3168zm^2}{5M^2r^8}<0.}
\end{equation}
For $R\ge0$, $0<r<2m$, its trace and physical interior anisotropy satisfy
\begin{equation}
\label{p3f:eq:v2212-a6-kottler-trace-anisotropy}
\begin{gathered}
 \widehat E^\mu{}_\mu
 =\mathcal C_6(450r-1084m-\tfrac{35}2Rr^3)+4Z_6<0,\\
 \widehat E_{\rm axial}-\widehat E_{\rm angular}
 =\widehat E^t{}_t-\widehat E^\vartheta{}_\vartheta
 =\mathcal C_6(144m-54r+\tfrac32Rr^3)>0,\\
 \partial_r(\widehat E_{\rm axial}-\widehat E_{\rm angular})
 =\frac{16zm^2}{5M^2r^{10}}
                   (432r-1296m-9Rr^3)<0.
\end{gathered}
\end{equation}
Thus the anisotropy increases along future homogeneous infall
$\dot r<0$.  The static coordinate $t$ is the axial spatial direction
in this interior; it must not be used as its proper-time label.
\end{theorem}
\begin{proof}
Let $\psi=m/r^3$.  The metric gives
\[
 C^2=48\psi^2,\quad J_C=96\psi^3,\quad
 \Box C^2=1440m^2/r^8-3456m^3/r^9-72Rm^2/r^6.
\]
For the inverse-metric Euler derivative
$\mathscr H_{\mu\nu}[L]=(\sqrt{-g})^{-1}
\delta\int\sqrt{-g}L/\delta g^{\mu\nu}$, evaluated after variation,
\[
\begin{split}
 \mathscr H_{\mu\nu}[RC^2]
   &=\tfrac R4 C^2g_{\mu\nu}
       +(g_{\mu\nu}\Box-\nabla_\mu\nabla_\nu)C^2,\\
 \mathscr H_{\mu\nu}[J_C]
   &=24\psi^3g_{\mu\nu}
       -12C_{\mu\alpha\beta\nu}\nabla^\alpha\nabla^\beta\psi,\\
 \mathscr H_{\mu\nu}[R^3]&=R^3g_{\mu\nu}/4.
\end{split}
\]
The first identity uses the vanishing Einstein Bach tensor.  For the
second, the type-D Weyl algebra gives
$\partial J_C/\partial\operatorname{Rm}=6\psi C$; the Einstein
divergence-free Weyl identity removes derivatives of $C$ in the
double-divergence term.  Its remaining algebraic term is
$24\psi^3g_{\mu\nu}$.  For example, where $H>0$ the orthonormal
Hessian entries of $\psi$ are
$3H'\psi/(2r)$, $12H\psi/r^2-3H'\psi/(2r)$, and
$-3H\psi/r^2$ twice.  Substitution into
Eq.~\eqref{p3f:eq:v2212-a6-einstein-representative}, followed by
$\widehat E=-2z\mathscr H[P_6]/M^2$, gives the displayed components.
This also reproduces the Schwarzschild components of
\cite{LitDYNMistryPinzulRachwal2020} after converting its
covariant-metric derivative and Riemann sign conventions.

Direct differentiation verifies
\[
 (\widehat E^r{}_r)'
 +\frac{H'}{2H}(\widehat E^r{}_r-\widehat E^t{}_t)
 +\frac2r(\widehat E^r{}_r-\widehat E^\vartheta{}_\vartheta)=0;
\]
its form multiplied by $H$ is regular also at a horizon.  More precisely
$\widehat E^r{}_r-\widehat E^t{}_t=-198\mathcal C_6rH$.
The ingoing components are
$\widehat E_{vv}=-H\widehat E^t{}_t$,
$\widehat E_{vr}=\widehat E^t{}_t$, and
$\widehat E_{rr}=-198\mathcal C_6r$.  They continue the calculation
through $H=0$ without a singular static frame.  Finally
$\Gamma^\mu{}_{rr}=0$ in that chart makes $k=-\partial_r$ affine.
Summing and subtracting the mixed components gives the trace and
anisotropy formulas.  Their signs follow from $r<2m$, $R\ge0$;
differentiation gives the last inequality.
\end{proof}

\begin{corollary}[Scalar overlap does not erase the evaluated tidal ports]
\label{p3f:cor:v2212-scalar-overlap-tidal-ports}
On a constant-scalar Einstein metric, the interior Euler tensor of any
smooth scalar-only overlap $\int\sqrt{-g}\,h(R)$ is
\[
 \mathscr H_{\mu\nu}[h(R)]
       =(Rh'(R)/4-h(R)/2)g_{\mu\nu}.
\]
Subtracting this overlap therefore changes neither the evaluated null
projection nor the axial--angular anisotropy.  It can change the scalar
trace allocation.  A general mixed quotient is not asserted to be a
scalar-only subtraction.
\end{corollary}
\begin{proof}
Constancy of $R$ removes $(g\Box-\nabla\nabla)h'(R)$ and
$R_{\mu\nu}=Rg_{\mu\nu}/4$ gives the expression.  A multiple of $g$
has zero null contraction and equal orthonormal spatial components.
\end{proof}

\subparagraph{Connection to the evaluated metric and source responses.}
Paper~VI-A,
Theorem~\ref{p8:thm:v2212-constrained-shape-hessian}, evaluates a
nonzero constrained rolling scalaron--shape block on the admitted local
scalar restriction. Paper~IV,
Corollary~\ref{p4:cor:v2212-massive-conformal-shape-kernel}, evaluates
the mixed causal source block on its admitted free massive member.
Their action, clock and preparation domains remain distinct from the
raw heat calculation above. Theorem~\ref{p8:thm:v2212-residual-flow-budget}
and Corollary~\ref{p8:cor:v2212-null-focusing-budget} give the exact
complete-curvature and null budgets into which actually matched
constituents must enter. These evaluated blocks specialize the existing
Schur theorem; they do not provide a second physical action.

\subparagraph{Narrowed physical obligations after the exact raw evaluation.}
The preceding conformal loads, Weyl weights and raw Kottler Euler ports
are now evaluated identities on their printed domains.  Normalizing
$\widehat E$ by the raw Einstein action instead multiplies it by
$M^2/(2\kappa_J)=1/96$; this preserves signs but does not construct the
physical action projection.  Only if the actual matched residual
consumes this constituent with the displayed sign does
$-\widehat E_{kk}$ supply a focusing contribution to its effective null
source, while its negative trace initially drives the bare scalaron
downward at a frozen native trace root.  Its increasing interior
anisotropy can contribute an opposing pressure-rate slot to the
complete-curvature flow.  These are competing ports, not a net
regulation theorem.  The calculation is an evaluation on a test jet,
not a solution of the corrected metric constraints or higher equations.
The actual whole-action/source projection, scalar overlap, depth-active
operator/localizer/domain, causal preparation, full higher-order or
retarded continuation, and signed invariant-range certificate remain
to be evaluated on the same physical history.  Neither a Euclidean
load nor the fixed raw coefficient establishes a universal
$\mathfrak I_{\rm curv}\le U^2$ range.  The exact nonlocal and physical
source frontiers of
Remark~\ref{p3f:rem:v118-parent-first-full-action-frontier} are unchanged.


\section{Complete source-jet and tensor-action identity tests}
\label{p3f:sec:v289-native-cap-discriminants}
The following results distinguish equality of an eliminated scalar
density from equality of the complete source-marked and tensor action.
They consume existing native/cap ports and do not select a new
microscopic gravitational mechanism.

\begin{theorem}[A finite complete-source two-jet reconstruction]
\label{p3f:thm:v289-complete-source-twojet}
Fix one real $n$-dimensional source port and its comparison frame, with a
common invertible $n\times n$ probe-incidence matrix $Q$. For $a=1,2$
let $T_a=T_a^{\mathsf T}>0$, $j_a\in\mathbb R^n$ and
$\ell_a\in\mathbb R$. Consider the quadratic source parents and their
classical stationary generating functions
\[
\begin{aligned}
 W_a(y;q)&=y^{\mathsf T}T_ay+2(j_a+Qq)^{\mathsf T}y+\ell_a,\\
 F_a(q)&=\min_yW_a(y;q)\\
 &=\ell_a-(j_a+Qq)^{\mathsf T}T_a^{-1}(j_a+Qq).
\end{aligned}
\]
Then equality of the values, all first derivatives and all second
derivatives of $F_1,F_2$ at $q=0$ is \emph{equivalent} to
\[
 \boxed{T_1=T_2,\quad j_1=j_2,\quad \ell_1=\ell_2.}
\]
For the full real $4\times4$ cap source ($n=16$), this is a
finite $1+16+16\cdot17/2=153$ scalar-jet certificate, once a
source-faithful set of $16$ actual independent probes and their
common normalization has been supplied. With incomplete or compressed
probes the implication need not hold. The statement neither creates
those physical probes nor eliminates extra source or tensor ports.
\end{theorem}
\begin{proof}
The exact source jets are
\[
 F_a(0)=\ell_a-j_a^{\mathsf T}T_a^{-1}j_a,\quad
 (DF_a)_0=-2Q^{\mathsf T}T_a^{-1}j_a,\quad
 (D^2F_a)_0=-2Q^{\mathsf T}T_a^{-1}Q.
\]
Equality of Hessians and invertibility of $Q$ first imply
$T_1^{-1}=T_2^{-1}$, hence $T_1=T_2$. Equality of gradients then
implies $j_1=j_2$, and equality of values gives
$\ell_1=\ell_2$. The converse is direct. Symmetry gives
$n(n+1)/2$ independent Hessian entries. For a nonspanning
probe $Q=(1,0)^{\mathsf T}$, the two distinct positive kernels
$T_1=I_2$ and $T_2=\operatorname{diag}(1,2)$, with zero sources,
both give $F(q)=-q^2$; this proves the stated limitation.
\end{proof}

\begin{corollary}[One scalar source contraction distinguishes a running cap from exact freezing]
\label{p3f:cor:v3028-r2-source-discriminator}
On the \emph{same identified positive source port and action owner},
let $B(R)=B_Ib(R)>0$ with
\[
 b(R)=1-d\,t(R/R_s)^2,\qquad
 t(x)=\frac{x}{1+\sqrt{1+x^2}},\qquad 0<d<1,
\]
and let $s(R)=\langle j(R),T(R)^{-1}j(R)\rangle$ be the
\emph{actual complete scalar source contraction}, including the
prescribed columns and normalization.  Suppose $s(0)=R_s^2$.
The existing one-anchor transport certificate,
Proposition~\ref{p3e:prop:v289-cap-source-ward},
identifies the inverse-square running return on a connected interval
if and only if
\[
 \boxed{\frac{d}{dR}\bigl(B(R)^2s(R)\bigr)=0.}
\]
For the fixed-reference cap $s_{\rm run}=R_s^2b^{-2}$;
for each exactly frozen scalar alternative $s_{\rm frozen}=R_s^2$.
At their common reference they obey the distinct second jets
\[
 \boxed{(B^2s_{\rm run})''(0)=0,\qquad
 (B^2s_{\rm frozen})''(0)=-B_I^2d.}
\]
Thus \emph{once physical source incidence has been proved}, one
scalar curvature two-jet separates running from exact freezing.
It does not distinguish distinct marked currents with the same
scalar contraction; this requires the common invertible physical
probe-incidence hypothesis of
Theorem~\ref{p3f:thm:v289-complete-source-twojet}, or an independent
native operator/source identity.  The same scalar result alone does
not prove equality of the full metric/tensor, boundary, connection
or quantum-measure actions.
\end{corollary}
\begin{proof}
The transport identity is equivalent to
$\frac{d}{dR}(B^2s)=0$ by
Proposition~\ref{p3e:prop:v289-cap-source-ward}.
Since $b(0)=1$, $b'(0)=0$ and
$b''(0)=-d/(2R_s^2)$, the running product is
$B_I^2R_s^2$ identically, while the frozen product is
$B_I^2R_s^2b^2$, whose second derivative at zero is
$-B_I^2d$.  The remaining qualifiers follow from the
complete-source reconstruction theorem and
Proposition~\ref{p3f:prop:v289-flrw-tensor-obstruction}.
\end{proof}


\begin{proposition}[Closed-$S^3$ FLRW blindness and the native tensor Hessian]
\label{p3f:prop:v289-flrw-tensor-obstruction}
On a smooth four-dimensional Lorentzian curvature chart with
compactly supported interior variations, let
\[
 I_f[g]=\int\sqrt{-g}\,f(R[g]),\qquad
 I_{f,\beta}[g]=I_f[g]+\beta\int\sqrt{-g}\,
 C_{\mu\nu\rho\sigma}C^{\mu\nu\rho\sigma},\qquad \beta\ne0.
\]
Every FLRW metric, including a closed-$S^3$ FLRW metric, is
conformally flat. Thus the two functionals have identical
restrictions to the \emph{entire} FLRW metric family, irrespective
of its time derivatives. At a conformally flat background $g_0$,
the Weyl-square first variation vanishes for every interior metric
tangent, while its second variation is
\[
 \boxed{D^2(I_{f,\beta}-I_f)_{g_0}[h,h]
      =2\beta\int\sqrt{-g_0}\,
          C^{(1)}[h]_{\mu\nu\rho\sigma}
          C^{(1)}[h]^{\mu\nu\rho\sigma}.}
\]
On a constant-curvature Einstein background, for transverse-traceless
$h$, $DR_{g_0}[h]=0$. Therefore the four-derivative principal Hessian
of a pure $f(R)$ density vanishes on that TT slot, while the displayed
Weyl-square Hessian is nonzero for suitable TT tangents.
At a locally flat static spatial symbol with
$h_{12}=h_{21}=u(z)$, all other components zero,
$DR[h]=0$ and $C^{(1)}[h]^2=(u''(z))^2$.
Consequently the ordinary native Dirac tensor packet of
Proposition~\ref{p3f:prop:v224-native-tensor}, with its nonzero
$\beta=3/(160\pi^2)$, cannot be identified with a \emph{pure}
metric-$f(R)$ action on the same complete off-shell metric chart
unless its tensor contribution is canceled by other full-action
terms or a complete source/field/boundary transformation is supplied.
The Euler term retains its separate topological and boundary
completion; this proposition neither adds the native tensor packet
again nor claims a physical extra pole from a local expansion.
\end{proposition}
\begin{proof}
FLRW metrics have $C=0$, which proves equality on the entire
restricted family. Since $C[g_0]=0$, every first variation of
$\sqrt{-g}C^2$ vanishes and the quadratic term is precisely the
square of the linearized Weyl tensor. The first curvature variation
is $DR[h]=\nabla_\mu\nabla_\nu h^{\mu\nu}-\Box h
-R_{\mu\nu}h^{\mu\nu}$, hence zero for TT tangents on an
Einstein background. A pure $f(R)$ density has four-derivative
quadratic principal term $f''(R_0)(DR[h])^2$, which therefore
vanishes in this slot. In the displayed flat static spatial test,
$R^{(1)}_{12}=-u''/2$, $R^{(1)}=0$ and
$(\mathrm{Riem}^{(1)})^2=2(u'')^2$,
$(\mathrm{Ric}^{(1)})^2=(u'')^2/2$.
The four-dimensional identity
$C^2=\mathrm{Riem}^2-2\mathrm{Ric}^2+R^2/3$
gives $C^{(1)2}=(u'')^2$. These are local principal-symbol
tests; an admitted compactly supported wave packet yields a
nonzero tensor quadratic form. The coefficient in
Proposition~\ref{p3f:prop:v224-native-tensor} is nonzero,
proving the conditional same-chart obstruction.
\end{proof}
\begin{remark}[What a native cap identity must now test]
The capped $f(R)$ remains the declared scalar working action and the
existing source reduction is exact on its stated comparison domain.
Equality of the scalar density, its full one-variable curvature jets,
or every closed-FLRW equation does not certify equality of the native
four-dimensional tensor action. The actual observable-sensitive
extra check is the complete TT/source-marked action two-jet, including
the native Weyl term, other sector contacts and any legitimate
field/boundary redefinition. A conditional cap plus a separately retained
tensor remainder is a possible common action \emph{only after} the
actual overlap and incidence are evaluated.
\end{remark}


\section{Petersen mirror source: an exact contact test for the complete action}
\label{p3f:sec:v319-petersen-complete-action}
On an occurring five-core $2I$ branch let $T_{ij}=J_i+J_j$ and let
$\mathcal B(Q_{ij}):\mathbf3\to\mathbf3'$ denote the core-quadrupole
mirror operator of Paper~I-B.  The present test concerns the same
real three-dimensional source/target representations.  It does not
identify them with an existing physical fermionic multiplet.

\begin{theorem}[Vanishing scalar contact, nonvanishing tensor round trip]
\label{p3f:thm:v319-petersen-mirror-contact}
For each of the ten pair axes $u=u_{ij}$ with
$Q_{ij}=uu^{\mathsf T}-I_3/3$, the exact composition is
\[
 C_{ij}:=T_{ij}^{\mathsf T}\mathcal B(Q_{ij})
       =\frac13Q_{ij},
\]
and consequently
\[
 \operatorname{Tr}C_{ij}=0,\qquad
 \operatorname{Tr}C_{ij}^2=\frac{2}{27},\qquad
 \operatorname{Tr}C_{ij}^3=\frac{2}{243}.
\]
The complementary three-core subset changes $C_{ij}$ to $-C_{ij}$.
Thus a scalar linear trace contact of these two candidate source
families vanishes identically; a quadratic tensor contraction is
nonzero and complement-even, while the cubic contraction is
complement-odd.  These are \emph{conditional action diagnostics},
not action terms selected by the geometry.
\end{theorem}
\begin{proof}
The first identity is the exact core/quadrupole interface of Paper~I-A,
Theorem~\ref{p1:thm:petersen-mirror-envelope}.  The eigenvalues of
$Q_{ij}/3$ are $2/9,-1/9,-1/9$, giving the three traces.
Complementation sends $T_{ij}\mapsto-T_{ij}$ and leaves the
fixed quadrupole $Q_{ij}$ unchanged.
\end{proof}

\begin{corollary}[Exact requirement for a physical Petersen response]
\label{p3f:cor:v319-petersen-response-owner}
Suppose an actual same-owner finite source action contains a
specified Petersen-marked operator family with licensed source ports.
Then the physical first and second response must include all varying
operator entries, markings, cross-sector blocks, kinetic metrics and
contact terms of that same action.  The finite trace and quadratic
identities above can then be evaluated in its actual Hessian;
without this premise no new mass pole, chiral splitting, hadron
spectrum, color-singlet contraction, Newton correction or running
coupling is obtained.
\end{corollary}
\begin{proof}
Apply the generated action co-consumption theorem
\ref{p3a:thm:v2216-common-action-consumption} and the complete
source-jet conditions of Paper~III-D,
Theorem~\ref{p3d:thm:v2216-complete-action-jets}.  Algebraic
existence of a finite operator is logically weaker than physical
incidence of that operator in the action.  The two hypotheses cannot
be interchanged.
\end{proof}


\subsection{Native two-core response: exact capacity and action-selection gap}
\label{p3f:sec:v338-two-core-source-owner}
Paper~I-B now constructs the intrinsic oriented two-core Hodge
projectors $P_4,P_\pm$, and Paper~V-A constructs a bilinear
Petersen--Clifford source map.  These are finite geometric source
\emph{capacities}.  Neither the tensor factors nor their spectral
weights determine the second derivative of the actual complete
action in a physical two-constituent channel.

\begin{proposition}[No unique energy functional from source eligibility]
\label{p3f:prop:v338-no-unique-two-core-mass-map}
On $U=\Lambda^2\mathbb R^5$ the unbroken $A_5$ symmetry permits
independently scaled self-adjoint quadratic forms
\[
 H_{a,b,c}=aP_4+bP_++cP_-.
\]
For example the identity $P_4+P_++P_-=I$ and the signed Hodge
response $P_+-P_-$ are both native but have different eigenvalues.
Therefore the five-core incidence and oriented-pair symmetries alone
do not fix a unique two-body rest-energy operator, nor any numerical
pion or kaon mass.  A same-action kinetic normalization and source
Hessian are necessary.
\end{proposition}
\begin{proof}
The three pairwise orthogonal equivariant idempotents have nonzero
ranks $4,3,3$.  Their real linear span consists of distinct invariant
Hermitian matrices, and symmetry places no relation on their three
scalar coefficients.  In particular, two coefficient triples can
produce different spectra while preserving every incidence and
exterior-algebra identity.  A physical operator must be derived by
applying the fixed action to the typed source rather than choosing
one such triple afterward.
\end{proof}


\section{Physical weak-action boundary for quadratic Petersen connection tangents}
\noindent\emph{Scope.} The statements below retain their declared finite geometric, source-chart, and action-relative hypotheses. They do not assert physical particle decay or unconditional source occupation.

\begin{proposition}[Gauge-transverse source family and action-owner firewall]
\label{p3f:prop:v368-petersen-quadratic-action-scope}
On the already generated relative round $S^3$ with midpoint Born
connection $A_B=\theta/2$, import the geometric first-surviving
quadratic Cartan source family of
Theorem~\ref{p1c:thm:v368-first-born-petersen-harmonic}.
For $\xi=P_uKe$ the normalized restricted static source functional is
\[
 \mathcal E_j(\xi)=\frac14+\frac5{324}\|\xi\|^2
 +\frac1{4320}\|\xi\|^4-j\cdot\xi.
\]
For each independently specified frozen real source $j$ in the
rank-two image, the stationary equation
\[
 \left(\frac5{162}+\frac{\|\xi\|^2}{1080}\right)\xi=j
\]
has exactly one solution, since the Hessian is positive on that
image. At zero source its unique solution is $\xi=0$.
This result does not infer that the complete chiral weak action
physically populates $\xi$, that a connection variation is an
occupied bosonic mode, or that its proper kinetic normalization
is the restricted static one. Those are source and state obligations
for the same local bosonic Euler--Noether equation and full
field--matter action.
\end{proposition}
\begin{proof}
Differentiate the imported quartic static Yang--Mills energy.
Its derivative is a strictly radial increasing map on the
rank-two plane; its Hessian has a positive definite isotropic term
and a positive semidefinite rank-one term.
\end{proof}

\section{Complete-response exports, action certificates, and conditional scope}
\label{paper3f:sec:v280-action-response-certificates}
\subsection{Constituted response theorems and exact scope}
\label{p3:sec:summary}
Paper~III-F begins after the physical momentum seed has already been fixed by
Paper~II-B,
\[
 \mathcal H_p\cong\mathcal H_1,\qquad m_p=1,
\]
and after Paper~III-A has completed the premise-minimal Paper-II-C/GP2 smooth
reciprocal-role domain, including the derived Born involution.  Its exact local
physical operation/generator polar is available on the derived
\((\mathrm{CAUS})\) clocked-Hopf domain, with numerical units supplied only by
the displayed calibrations.  Metric-curry is not an
existence premise for the reciprocal role.
It does not reopen the former \(\mathsf P_{\rm sli}\) selector.  The
one-profile additive class \(\Dcl\) is loop-derived in Paper~III-B.  Within
the unmarked minimal real scalar Laplace-type resolution signature, Paper~III-E now
derives the specific gravitational law \(F_g(u)=e^{-u}\); the historical
\((\mathrm{GRAV}_{\circ})\) label is its certification rather than an
occurrence datum.  A different profile would enlarge the process root.
On the separate sharp-frontier occurrence \(\mathsf{SF}_{\Lambda_F}\), homogeneous
metric scaling gives the exact cumulative reciprocal-position resolution
\[
 m_x(\widetilde t)
 =\sup_{\widetilde s\in[\widetilde t_0,\widetilde t]}
 \max\!\left\{1,
 \left\lfloor
 \sqrt{1+\bigl(\Lambda_Fa(\widetilde s)\bigr)^2}-1
 \right\rfloor\right\}.
\]
When the scale supremum \(A_*\) is attained, this is the same floor formula
with \(a\) replaced by \(A_*\); an unattained exact entry threshold has the
one-level correction proved in
Theorem~\ref{p3:thm:exact-spectral-resolution-staircase}.
Here \(\ell_F=\Lambda_F^{-1}\).  Identifying this scale with the smooth action
scale requires \((\mathrm{ID}_{F/A})\) at the point of use.  This is a spectral
theorem once the operational frontier occurs.  Monotone or
unbounded growth additionally requires the corresponding behavior of the
actual solution \(a(\widetilde t)\).
For every continuous nondecreasing realized scale history,
Theorem~\ref{p3f:thm:spectral-frontier-event-process} now closes the
frontier's event-indexed dynamics: consecutive entries obey
\(m\mapsto m+1\) and add exactly \(E_{m+1}^x\).  This is a discrete
continuation on crossing order, not a physical-time rate law.  At a strict
crossing, \(m_x\) alone fails the Paper~III-A projectability test because two
points in the same level cell can have different next-crossing outcomes after
the same time increment.  The generated residual
\(r_F=a_{m_x+1}-a\) restores the discarded phase and transports an already
generated scale process, but does not construct the scale history.

The local cotangent target, tautological form, canonical symplectic form, and
cotangent-lift gluing are derived from scale-free evaluation duality, and
Paper~III-A's generator-germ theorem makes their reciprocal-generator
realization generated/co-realized rather than conditional on a G0
\((\mathrm{GL}_{xp})\) occurrence.  The calibrated positive polar additionally
requires \((\mathrm{CAUS})\), and macroscopic metric variation follows from
the occurring minimal action's derived
\((\mathrm{VDOM}_x^{\min})\).  The legacy
\((\mathrm{GL}_{xp})\) certification is not the finite-band comparison
\((\mathrm{PAIR}_{xp})\): the latter still concerns specific position
harmonics/observables versus compact momentum operators and does not fix the
reciprocal scale calibration \((\mathrm{CAL}_{xp})\).
The protected pure-metric branch used in the local spectral and backreaction
theorems imports
\[
 (\mathrm{GRAV\!-\!ASYM})+\mathsf C_{\rm IR}^{\rm der}
\]
from Paper~III-E.  Its natural scalar-density Born theorem preserves one full
copy of the vacuum, Einstein, and higher pure-metric coefficients when only
the position metric is varied.  The legacy \((\mathrm{EH}_x)\) label records
the Einstein specialization and is not a further premise.  On this
branch, \((\mathrm{CORN})\) and \((\mathrm{BD}_{\rm EH})\) are derived
certifications of, respectively, the ambient rank-one constituent and the
Born-orbit diagonal pullback; neither is a further premise.  Nonmetric bundle,
matter, and fermionic position terms retain the III-B role weight \(1/2\), and local equations
are varied only inside the unit-localizer domain \(U_0\Subset U\).  Local bosonic matter backreaction requires
\((\mathrm{ACT}_{{\rm J,EH+YM+\Phi},U_0}^{\rm loc,der})
 +(\mathrm{SOLHIST}_{U}^{\rm phys})\).  Its Newton coupling coordinate
\(G_{J,\mathrm{raw}}\) is derived from the generated joint-EH coefficient.
Its identification with an observed Newton readout consumes the actual
physical action/source projection; numerical calibration follows that
identification.  Complete matter backreaction additionally
consumes only the actually used constituents of
\((\mathrm{ACT}_{C}^{\rm res})\); quasi-local infrared locality
requires \((\mathrm{QLOC})\).  Lorentzian interpretation uses either a direct
Lorentzian occurrence, the classical continuation bridge, or the stronger
claim-specific reconstruction tag actually asserted.
The derived protected support coefficients are
\(w_{\rm vac}=w_{\rm EH}=1/8\).  On the generic local \(f(R)\) atlas,
\((\mathrm{REN}_0(\Lambda_{\rm ren}))\) supplies the renormalized subtraction,
sign, and magnitude of the zero-derivative term, so the support weight alone
does not determine a generic vacuum background.  On the native cyclic terminal
branch this generic freedom is not consumed: the same-action inverse-Euler
endpoint fixes \(-C_7\) and
Corollary~\ref{p3f:cor:v2118-native-cyclic-ren0-firewall} forbids an additional
independent subtraction.
Likewise \(w_{\rm nl}(R_+)=1\) is a full-octet ultraviolet response weight, not
a proof of inflation, slow roll, exit, or reheating.

For the nonlinear gravitational branch, III-F imports the smooth
\(\mathfrak f_{\rm sp}\) response of Paper~III-E only after freezing one
\(\mathrm{REG}_{\sigma}\in\mathfrak A_{\rm REG}^{+}\) member, supplying
the certified \((\mathrm{CRP4})\) and \((\mathrm{PQTI})\) owner relations
of Definition~\ref{p3f:def:v2211-certified-response-domain}, and one
\((F_C,Z_C)\in\mathfrak A_{\rm TERM}\).  The
visibility/state/entropy
constituents are generated/co-realized; the scale map is their unique
owner-typed inheritance, the covariant Hessian pullback and entropy coupling
are unique in the minimal signature, and no separate active-state,
transport, or coupling selector remains.  On
the generated-owner branch,
\[
 R_+=U,\qquad
 C=U\alpha_{\rm hier}(\sigma_{\rm fin})^{-7},
\]
and its Hessian contains the identified active fraction
\(w_{\rm nl}=(1/8)8^{n_{\rm act}}\).  This identification removes a separate
active-state routing premise.  Here \(R_+=U\) follows from identity-stage inheritance and the III-E momentum-curvature actualization, while the neutral dual determinant-depth
theorem generates
\(\alpha_{\rm hier}(\sigma_{\rm fin})^{-q}\).  The physical identities
\(C/R_+=\alpha_{\rm hier}(\sigma_{\rm fin})^{-7}\) and the intermediate scale
ratios are its forced inheritance through the CRP4 response ancestry.
The generated Einstein action theorem still cannot be reversed to derive
CRP4 without circularity.  Its noncircular test is the actual
action/contact certificate, beyond construction of the action-density
dual; the v1.75 joint-Dirac theorem directly
discharges the historical \(\mathrm{ACT}_{\rm EH}^{\rm cl}\) bridge on its
local clock-adapted domain.
Paper~III-E retains \((F_C,Z_C)\) as independent raw
response endpoint data on the generic reconstruction atlas.  When that atlas is
composed with \((\mathrm{REN}_0(\Lambda_{\rm ren}))\) in the generic III-F
action, Theorem~\ref{p3f:thm:term-ren-zero-order-redundancy} shows that only
\(\overline F_C=F_C-2\Lambda_{\rm ren}\) is the physical zero-order endpoint
coordinate before positive normalization.  The native cyclic terminal branch
is the stricter specialization of
Corollary~\ref{p3f:cor:v2118-native-cyclic-ren0-firewall}: its zero-order
intercept is already generated as \(-C_7\), so no independent \(\mathrm{REN}_0\)
coordinate remains to be selected there.  On \(Z_C>0\),
Theorem~\ref{p3f:thm:term-ren-multiplicative-action-quotient} further removes
the simultaneous action scale, leaving
\(\widehat{\overline F}_C\) and \(G_{\rm term}\), while
Corollary~\ref{p3f:cor:pqti-same-action-invariance} fixes the formerly free
\(\widehat\gamma_\sigma\) to
\(16\sqrt{\alpha_\sigma}\).  An externally fixed Newton constant
turns the quotient normalization into calibration.  The response
normalization \(\kappa=2\) is only the adopted
representative, and the classical field equation additionally requires
\((\mathrm{MRESP}_{\rm cl})\).  The seven scalar scale
levels are ordered without \((\mathrm{FLAG})\); only a rank-one projector
chronology requires that refinement.  The four active support anchors are the
canonical spectral pullback imported from Paper~III-E and require no
\((\mathrm{ACT}_{\rm Cl})\); they do not assert actual transitions.  The
complete reconstructed response graph inherits the joint input witness, while
its action-role adoption remains explicit.  No obsolete fixed-\(k\) response
family is used.

The gauge layer contains two distinct constructions.  The seven-dimensional
principal carrier over the local position chart and its unique nonflat
Born-invariant homogeneous connection are generated by exact pullback and
symmetry.  The residual round vector-fibre isometry-connection reduction gives
the geometric equality \(g_L=g_R\).  Its general conditional absolute
coefficient is
\[
 \frac1{g_L^2}=\frac1{g_R^2}
 =\frac{\kappa f_2}{48}(r_{\rm fib}\Lambda_{\rm cut})^2
\]
under \((\mathrm{GRED}_7)+(\mathrm{CAL}_{\rm GRED}(\kappa))\), including the
same protected Einstein coefficient, connection, fibre average, and
Yang--Mills convention. The shorter \(\kappa f_2/48\) coefficient and its
numerical calibration examples additionally require the explicit compact-cell
condition \(r_{\rm fib}\Lambda_{\rm cut}=1\).  Independently, \((\mathrm{BULK}_8)\) gives the
ordinary eight-dimensional Born-even factor \(1/2\); importing that factor
into the seven-dimensional coefficient additionally requires
\((\mathrm{MATCH}_{8\to7})\).  The value
\(f_2/48\) is specifically the \(\kappa=1\) quaternion-generator convention;
the Paper-II spin-\(\tfrac12\) convention gives \(f_2/192\).  The protected scalar has rank
\(N_{\rm sc}^{\rm rank}=1\), while the separate corner action coefficient is
\(w_{\rm sc}^{\rm corner}=1/8\).  A physical scalar-circle metric and
\(U(1)_{\rm fib}\) kinetic coefficient are not part of the main
constitutional branch.

\subparagraph{Residual-status summary.}
The local cotangent geometry, sharp-frontier band, pure-metric Born routing,
the nonzero joint connection anchor and its same-datum descendants, hierarchy
ratios after a frozen numerical scheme boundary, the local GRED carrier, its
unique homogeneous Born connection, both faithful fibre-spin extensions, and
the geometric equality \(g_L=g_R\) inside the round diagonal-fibre branch are
derived on their stated domains.

The TERM--REN combination is an exact same-action quotient,
\[
 \mathfrak A_{\rm TERM\text{-}REN}^{\rm act}
 :=(\mathfrak A_{\rm TERM}\times\mathfrak A_{\rm REN})
   /\{\text{additive zero-order shifts and positive action scalings}\}.
\]
It is not an additional structural selector family.  On the terminal
constant-curvature vacuum branch,
Corollary~\ref{p3f:cor:normalized-terminal-vacuum-data} fixes its normalized
zero-order endpoint to \(\widehat{\overline F}_C=C/2\); actual attainment of
that endpoint is a history/solution condition.

The former premise-minimal constitutive gauge occurrence is now the theorem
certificate
\[
 \boxed{(\mathrm{JOINT}_x^{\rm der}).}
\]
It is not an assumption.  Its canonical member is nonflat, while the full
affine gauge/odd family is retained without choosing a fluctuation member.
The REG family is one exact covariant object over its affine torsor; only a
requested nonhorizontal numerical evaluation adds a boundary condition.  The
generic \(\mathfrak A_{\rm GRED}^{\rm phys}\) member is optional and is
consumed only by the enlarged diagonal-fibre normalization/fluctuation branch,
not by the premise-minimal common-heat gauge chain.  Stationarity, Lorentzian
continuation, physical-time frontier histories, and interacting reconstruction
remain theorem/solution domains rather than structural selectors.  Numerical
placements remain calibrations.  Thus no local Paper-III-F structural root is
added beyond the framework-wide roots recorded in Paper~III-E.

\begingroup
\small
\setlength{\tabcolsep}{3pt}
\renewcommand{\arraystretch}{1.15}
\begin{longtable}{@{}L{.22\linewidth}L{.44\linewidth}L{.27\linewidth}@{}}
\caption[Paper III-F export ledger]{Paper III-F export ledger.}
\label{tab:p3-export-ledger}\\
\toprule
\textbf{Export} & \textbf{Paper-III-F status} & \textbf{Downstream use} \\
\midrule
\endfirsthead
\multicolumn{3}{@{}l}{\small\emph{Export ledger (continued)}}\\
\toprule
\textbf{Export} & \textbf{Paper-III-F status} & \textbf{Downstream use} \\
\midrule
\endhead
\midrule
\multicolumn{3}{r@{}}{\small\emph{Continued on next page}}\\
\endfoot
\bottomrule
\endlastfoot
Fixed momentum seed
& \(\mathcal H_p\cong\mathcal H_1\), \(m_p=1\), imported from Paper~II-B
& finite quantum/field carrier input \\
Derived local reciprocal target
& \(T^*U\), \(\vartheta\), \(d\vartheta\), cotangent-lift gluing, and reciprocal generator-germ physical co-realization are derived; calibrated polar under \((\mathrm{CAUS})\)
& local phase-space theorem output \\
Sharp scalar resolution staircase
& exact under \(\mathsf{SF}_{\Lambda_F}\), with \(\ell_F=\Lambda_F^{-1}\);
the complete retained frontier gives a \textsc{derived-later} discharge of the
II-C scalar generated-boundary condition and generates the branchwise
\(\mathrm{HIST}_x^{\rm fin}\) record; cumulative retention determines the level
from the exact integer supremum (or \(A_*\) with its attainment qualification),
while a continuous nondecreasing scale parent derives the exact event law
\(m\mapsto m+1\) and band inclusions; \(m_x\) alone has no homogeneous
physical-time law at a strict crossing, and the \(\Dcl\) action profile alone
supplies no finite frontier
& reciprocal-position growth and Paper~IV scale bridge \\
Reciprocal-position locality
& conditional on
\((\mathrm{VDOM}_x^{\min})+(\mathrm{QLOC})\) and the
protected local spectral/EH branch; the minimal variational domain is derived
from the occurring action's configuration tangent; claim-specific
continuation only for a Lorentzian interpretation
& macroscopic/local-field limits \\
Common-action backreaction
& conditional on
\(\Dcl+(\mathrm{JOINT}_x)+(\mathrm{ACT}_{\rm cl})
+(\mathrm{SOLHIST}_{U}^{\rm phys})\), with the derived joint-Dirac Einstein--Yang--Mills--scalar subfunctional and its derived coupling coordinate \(G_{J,\mathrm{raw}}\), the III-B half weight on nonmetric position
terms, derived same-datum connection/curvature/profile relations, and
compactly supported variation in \(U_0\)
& sourced metric dynamics \\
Smooth nonlinear gravity
& current Paper~III-E integrable \(\mathfrak f_{\rm sp}\) after a frozen
\(\mathrm{REG}_{\sigma}\), \((\mathrm{CRP4})+(\mathrm{PQTI})\),
\((\mathrm{TERM}_f)\), and \((\mathrm{MRESP}_{\min})\) occurrence;
transport/coupling and the reconstructed graph are derived/co-realized; the
active state/scale map and four canonical support anchors co-realize;
classical
equation under \((\mathrm{MRESP}_{\rm cl})\)
& scalaron/cosmological response handoff \\
Round vector-fibre gauge slot
& canonical local principal carrier and nonflat homogeneous Born connection
are derived; both \(SU(2)^a_L\times SU(2)^a_R\) extensions, their faithful
\(W_L\oplus W_R\) representation, and the joint Dirac anchor are derived;
optional independent fluctuations retain geometric \(g_L=g_R\); absolute
\(\kappa f_2(r_{\rm fib}\Lambda_{\rm cut})^2/48\) under
\((\mathrm{GRED}_7)+(\mathrm{CAL}_{\rm GRED}(\kappa))\), reducing to
\(\kappa f_2/48\) only with \(r_{\rm fib}\Lambda_{\rm cut}=1\)
& later gauge-sector identifications \\
Eight-dimensional Born diagnostic
& optional \((\mathrm{BULK}_8)\) gives \(\beta_{\rm B}^{\rm bulk}=1/2\);
use in \(X_7\) additionally requires \((\mathrm{MATCH}_{8\to7})\)
& auxiliary comparison/matching only \\
Scalar-circle slot
& dimensionless representation slot only in the main branch; physical metric normalization deferred
& optional later \(U(1)\) realization \\
\end{longtable}

\endgroup

\begingroup
\small
\setlength{\tabcolsep}{3pt}
\renewcommand{\arraystretch}{1.12}
\begin{longtable}{@{}L{.23\linewidth}L{.22\linewidth}L{.48\linewidth}@{}}
\caption[Paper III-F G0 interface ledger]{G0 classification of the physical
interfaces used in Paper III-F.  Every constitutive row opens a
result-independent dependency frontier, receives relative GP1--GP3 closure,
and preserves the frozen upstream presentation.}
\label{tab:p3f-g0-interface-ledger}\\
\toprule
\textbf{Interface} & \textbf{G0 type} & \textbf{Frozen record} \\
\midrule
\endfirsthead
\multicolumn{3}{@{}l}{\small\emph{Interface ledger (continued)}}\\
\toprule
\textbf{Interface} & \textbf{G0 type} & \textbf{Frozen record} \\
\midrule
\endhead
\midrule
\multicolumn{3}{r@{}}{\small\emph{Continued on next page}}\\
\endfoot
\bottomrule
\endlastfoot
Generated gravitational heat role
& derived certification, not a new G0 root
& \(F_g=e^{-u}\) and legacy \((\mathrm{GRAV}_{\circ})\) reuse the upstream
witness; \((\mathrm{GRAV\!-!TR})\) and
\((\mathrm{GRAV\!-!ASYM})\) remain analytic domains \\
\((\mathrm{LOC}^{\rm ev}_{xp})\), legacy \((\mathrm{GL}_{xp})\)
& derived certification; no G0 root
& canonical reciprocal generator-germ localization on \(T^*U\) from
Paper~III-A; no independent \(\Delta_{\rm LOC}\) witness remains \\
\((\mathrm{VDOM}_x^{\min})\), legacy \((\mathrm{POL})\)
& derived configuration-tangent certificate
& all and only local action arguments vary; compact momentum representation
data are not action fields; enlarged variational theories are distinct parents \\
\((\mathrm{JOINT}_x^{\rm der})\)
& derived rooted connection/operator certificate; no G0 root
& canonical nonflat \(\Spin(4)_{\rm fib}\) connection represented faithfully
in one Dirac family; common associated connection, curvature, profile, and
parity-even ordinary-trace kinetic normalization are theorem outputs \\
\((\mathrm{STATFAM}_x^{\rm der})\), legacy \((\mathrm{STAT}_x)\)
& generated exact-solution-family role after \((\mathrm{ACT}_{\rm cl})\)
& Euler--Lagrange map and complete stationary moduli are generated; one actual
member/history remains \((\mathrm{SOLHIST}_{U}^{\rm phys})\), with no
nonemptiness, stability, vacuum, or global-extension claim \\
Generated MRESP constituents
& derived/co-realized; scale entries also require owner inheritance
& \(V_\tau,Z,\rho_\tau,S,q_{\rm sp}\) are generated from the active
composite; \(\scal,\sigma_{\rm act}\) use the CRP4 response ancestry and
\(\mathrm{REG}_{\sigma}\in
\mathfrak A_{\rm REG}^{+}\); the unique reconstructed response graph inherits
the joint input witness \\
Canonical active spectral profile
& derived; physical placement inherits the response owner
& four exact scale/support anchors; complete orbit family unmarked, particular
\(\Pi_3(u)\) marked; no continuous rank-changing path or actual transition \\
\((\mathrm{REG}_{\sigma})\)
& exact scheme-covariant role; numerical boundary conditional
& the affine family carries canonical determinant-weight transport; horizontal
or total-weight-zero descendants require no member; the generated \(2O\)
interface anchors \(\mathrm{REG}_0\), while unanchored branches retain the
family \\
\((\mathrm{CRP4})\), historical alias
\((\mathrm{OWN}_{\rm resp/det})\)
& derived scalar-action coefficient ancestry
& its response owner is the fundamental determinant-dual port; the rank-four
positive root and \(R_+=\Lambda_{\rm EH}^2\) are generated, with no second
alias count \\
\((\mathrm{TERM}_f)\)
& raw reconstruction boundary; quotiented in the same action
& typed endpoint data \((F_C,Z_C)\) before REN; afterward only normalized
same-action invariants remain, and the terminal-vacuum branch fixes
\(\widehat{\overline F}_C=C/2\) \\
\((\mathrm{MRESP}_{\min})+(\mathrm{PQTI})\)
& scalar-action response domain plus derived second-variation owner
& minimal response with derived tensor/entropy laws; PQTI fixes
\(\widehat\gamma_\sigma=16\sqrt{\alpha_\sigma}\), while
\(\kappa_U=1\) follows from identity-stage inheritance \\
\((\mathrm{REN}_0(\Lambda_{\rm ren}))\)
& raw subtraction boundary; paired-action quotient
& frozen zero-derivative representative before composition; additive and
positive multiplicative kernels combine it with TERM, so it is not an
independent structural root \\
\((\mathrm{GRED}_7^{\rm phys})\)
& optional physical extension/calibration on a derived carrier
& local base, principal carrier, gluing, and homogeneous Born connection are
derived; enlarged isometry-connection fluctuations, physical reduction
adoption, fibre/action normalization, and diagonal-fibre coefficient matching
are retained only when that optional branch is asserted \\
\((\mathrm{BULK}_8)\)
& optional constitutive comparison
& \(Y_8,L_8,\mathcal B,P_+\), normalized measures, fixed-set domain \\
\((\mathrm{MATCH}_{8\to7})\)
& constitutive/\allowbreak calibration
& result-independent normalized-volume coefficient-routing map \\
\(\mathsf{SF}_{\Lambda_F}\)
& constitutive
& \(\Lambda_F\), projector/test family, comparator, retention convention;
continuous nondecreasing scale histories derive the event successor, not the
scale process or its physical-time rate \\
\((\mathrm{CAL}_{\rm GRED})\), \((\mathrm{CAL}_{xp})\), \((\mathrm{CAL}_t)\)
& calibration specialization
& gauge standard, reciprocal readout map, and primitive-turn duration fixed before numerical use \\
\((\mathrm{LOR}_{\rm dir})\) or \((\mathrm{CONT}_{\rm cl})\)
& occurrence/analytic bridge
& field identity, sign, domain, and reconstruction scope \\
\end{longtable}

\endgroup

Every physical localization, same-process construction, gauge
reduction, and sharp frontier in this export ledger is typed under G0.  The
remaining occurrence/relation rows are not interchangeable:
\(\mathfrak A_{\rm TERM\text{-}REN}^{\rm act}\), and
\(\mathfrak A_{\rm GRED}^{\rm phys}\)
remain explicit realization/relation families.  The same-datum JOINT
connection/operator anchor and relations, action-kernel quotients, and
canonical GRED carrier/background are not additional selectors.  The minimal variational domain and horizontal
scheme transport are derived; a nonhorizontal numerical scheme value is a
boundary condition rather than a structural selector.  Derived consequences inside a selected branch---for
example the common trace normalization or \(g_L=g_R\)---do not select a
fluctuation member or discharge the action interface itself.  The analytic
\((\mathrm{QLOC})\) and claim-specific continuation or reconstruction clauses
remain separate hypotheses and are not consequences of geometric primacy;
\((\mathrm{CAL}_{\rm GRED})\), \((\mathrm{CAL}_{xp})\),
and \((\mathrm{CAL}_t)\) are calibrations rather than
structural selectors.

The paper makes no claim about particle names, cosmological epochs,
the Standard Model, observational fits, or downstream homogeneous realizations.  Such
identifications may consume these exports only after supplying their own
additional interfaces.


\subsection{Geometric consequences and physical-response conditions}
\label{p3f:sec:v2201-geometry-first-consequences}
\subparagraph{Geometry ancestry.}
Round curvature owner and rooted hierarchy $\to$ full running $f(R)$ and the scoped gravity/free-source action $\to$ Schur residues, exact source work and coefficient jets.
\subparagraph{Status.}
These are consequences on the domains stated below.
A mathematical certificate or an executed restriction retains every
physical source, state, clock and readout hypothesis of its inputs.
The existing complete action on the declared gravity/free-Dirac
domain is retained. The results below use that action and its inherited
boundary terms. They introduce no requirement for an unrestricted
microscopic Lorentzian action. A full interacting source claim still
consumes its own complete source and metric variations.

For the gravitational calculations in this section use $\hbar=c=1$ and
$U=1$ only for dimensionless numerical evaluation. Define
\begin{gather*}
 L_\alpha=4\pi^2+\frac{1}{6\sqrt2},\quad Q=8e^{L_\alpha/2},\quad
 C_j=Ue^{-jL_\alpha},\\
 q_R=-\frac{\log(R/U)}{L_\alpha},\quad
 \chi=f_{RR},\quad F=f_R,\quad P=RF-f,\quad M^2=U/4.
\end{gather*}
The physical curvature coordinate is $K_R=16q_R$. The customary
structural K16 coefficient freeze occurs at $q_R=2$, $R=C_2$ and
$K_R=32$. The structural successor-cell clock remains a separate owner.
Its $\eta_{8\to16}$ is not obtained by identifying that clock with
$8q_R(t)$. The running factors $D_2$ and $n_7^+$ below are exactly those
of the inherited Hessian, with no fixed-$R^2$ replacement.
See \ref{p8:def:curvature-shell-coordinate}, \ref{p3f:eq:v169-inherited-hessian}.

\subsubsection{Quantitative residue certificate for a guarded pole}
\label{p3f:sec:v2201-derived-d03}
\noindent\textit{Status: proved result on stated domain.}

For an admitted Hermitian block
\[
 H=\begin{pmatrix}A&B\\B^\dagger&C\end{pmatrix},\qquad
 K(z)=A-z-B(C-z)^{-1}B^\dagger,
\]
let $\lambda$ be a simple full eigenvalue outside $\operatorname{spec}C$,
and let the retained eigenvector $u$ have norm one. Its omitted component is
$v=-(C-\lambda)^{-1}B^\dagger u$. For
$\delta=\operatorname{dist}(\lambda,\operatorname{spec}C)>0$,
\[
 Z_\lambda=I+B(C-\lambda)^{-2}B^\dagger=-K'(\lambda),\qquad
 w_{\rm ret}=\frac1{u^\dagger Z_\lambda u}
 \ge\frac1{1+\|B\|^2/\delta^2}.
\]
\begin{proof}
Normalize the full vector $(u,v)$ by
$(1+\|v\|^2)^{-1/2}$. Its spectral projector has retained weight
$(1+\|v\|^2)^{-1}$. Differentiation of the Schur operator gives the
displayed positive matrix, and the resolvent norm gives the bound.
\end{proof}
The same omitted image that shifts a pole changes its residue. A zero of
the reduced determinant alone does not give a canonically unit source
weight. The numerical witness checks the formula against full
diagonalization; it is not an environment selection.
\par\noindent\textit{Inherited inputs:} \ref{p3f:thm:v2168-left-return-test}, \ref{p3f:cor:v2167-family-equivariant-retention}, \ref{p1b:lem:v2167-compression-multiplication}.

\subsubsection{Exact Bloch ports from the complete free Dirac source}
\label{p3f:sec:v2201-derived-d09}
\noindent\textit{Status: proved result on stated domain.}

On each inherited paired $\pm\lambda$ block set $p=|\lambda|/a$,
$p_i=|\lambda|/a_i$, $E_i=\sqrt{p_i^2+m^2}$ and
\[
 A=\frac{pp_i+m^2}{E_i},\quad B=\frac{m(p-p_i)}{E_i},\quad
 \widehat H_{\rm even}=A\sigma_3+B\sigma_1+\varepsilon_0I.
\]
For the possibly unnormalized even-state Bloch vector
$r_j=\operatorname{Tr}(\sigma_{\rm even}\sigma_j)$,
\[
 \boxed{\dot r=2(B,0,A)\times r,\qquad |r|^2=\text{constant}.}
\]
The odd-sector weights are constant and the complete pair energy, including
the central allocation, is $\mathcal E=Ar_z+Br_x+\varepsilon_0$.
Writing $\mathcal V=2\pi^2a^3$, the homogeneous ports are
\begin{align*}
 \rho&=\mathcal E/\mathcal V,\\
 3p_s\mathcal V&=\frac p{E_i}(p_i r_z+m r_x)-a\varepsilon_0'(a),\\
 (\rho-3p_s)\mathcal V&=
 \frac{m^2r_z-mp_i r_x}{E_i}+\varepsilon_0+a\varepsilon_0'.
\end{align*}
\begin{proof}
Use the Pauli commutator for the state equation and
$p_s=-\langle\partial_a\widehat H\rangle/(3\mathcal Va^{-1})$.
Since $a\partial_ap=-p$, the kinetic terms cancel in the trace and the
displayed coherence remains. The identity
$d\mathcal E/dt=-3Hp_s\mathcal V$ gives exact source conservation.
\end{proof}

This closes the finite source equations for every supplied compatible
state; no Wick approximation is used. Coupled to
\[
 \dot a=aH,\quad \dot H=R/6-2H^2-a^{-2},\quad
 \dot F=v,\quad
 \dot v=-3Hv-\frac{2f-RF}{3}+\frac{\rho-3p_s}{3M^2},
\]
and the inherited lapse constraint, the equations are an explicit finite
Cauchy system. Full invariant multiplets and a spatial-isometry-invariant
state certify exact FLRW stress. The central derivative cannot be dropped.
No claim identifies this free block with the complete K16 direct packet.
\par\noindent\textit{Inherited inputs:} \ref{p3f:thm:v169-explicit-free-dirac-source}, \ref{p3f:thm:v169-finite-source-work}, \ref{p3f:thm:v169-flrw-lapse-and-phase-action}.

\subsubsection{Scalaron work fixed by the shared source trace}
\label{p3f:sec:v2201-derived-d10}
\noindent\textit{Status: proved result on stated domain.}

For a covariant, curvature-independent Jordan matter action, after its
consistent Einstein-frame field transport,
\[
 J_\phi=-\frac{\rho_E-3p_E}{\sqrt6M},\qquad
 Q_s=J_\phi\dot\phi.
\]
Thus
\[
 \dot\rho_s+3H_E(\rho_s+p_s)=Q_s,\qquad
 \dot\rho_\phi+3H_E(\rho_\phi+p_\phi)=-\sum_sQ_s,
\]
where $\rho_\phi=\dot\phi^2/2+V_E$ and
$V_E=M^2(RF-f)/(2F^2)$ is the full running potential.
\begin{proof}
At fixed Einstein metric, $g_J=F^{-1}g_E$ and
$d\log F/d\phi=\sqrt{2/3}/M$. The Jordan metric variation is a Weyl
variation; its on-equation source is the complete trace. Comparing
$\delta S_s=-\int\sqrt{-g_E}\,J_\phi\delta\phi$ gives the first identity.
Multiply the scalar equation by $\dot\phi$ and differentiate its canonical
energy to obtain the two cancellation laws.
\end{proof}
There is one scalaron energy carrier. Massive Dirac channels exchange
energy through the derived exponential mass; a classical massless conformal
channel has zero trace exchange. An anomaly or nonminimal curvature source
must use its own complete variation. Hierarchy landmarks alone are not
additional conserved energy packets. This is the derivable shared-reservoir
law on the stated action domain; numerical branching needs the actual states
and source kernels. \par\noindent\textit{Inherited inputs:} \ref{p3f:thm:v169-explicit-free-dirac-source}, \ref{p3f:eq:v169-einstein-potential-and-mass}, \ref{p8:rem:v2134-retained-scalar}.

\subsubsection{Exact third-jet jump and geometric consequences}
\label{p3f:sec:v2201-derived-d14}
\noindent\textit{Status: proved result on stated domain.}

At $R=C_2$ the positive-part selector $n_2^+$ has zero value. From the
literal secant,
\[
 n_{2,R}(C_2^+)=
 \frac{U-C_2}{2L_\alpha C_2(U+C_2)}.
\]
Hence, defining the jump as the above-curvature value minus the
below-curvature value,
\[
 \boxed{[f_{RRR}]_{\rm above-below}
 =\chi_7(C_2)\log Q\,
 \frac{U-C_2}{2L_\alpha C_2(U+C_2)}.}
\]
Numerically $U^2[f_{RRR}]\simeq1.36530\times10^{34}$, while
$U\chi(C_2)\simeq2$ is continuous.
\begin{proof}
Differentiate $\log g_2/\log(R/U)$ at its zero numerator.
Above $C_2$, $D_{2,R}=\log Q\,n_{2,R}$; below $C_2$ it is zero.
The remaining factor is smooth, so only this term jumps.
\end{proof}
If the source trace is continuous and there is no surface impulse,
$F,\dot F,\ddot F$ and $R,\dot R$ remain continuous. In the time direction
of the descending crossing,
\[
 [\ddot R]_{\rm post-pre}
 =\frac{[\chi_R]_{\rm above-below}}{\chi(C_2)}\dot R_{16}^2.
\]
The first corresponding jump in the FLRW $H$ derivatives is in
$H^{(3)}$, through $\ddot R/6$. This is a finite-jet junction, not a
delta-function squeeze. It also states exactly which nonanalytic selector
a microscopic parent map must reproduce.
\par\noindent\textit{Inherited inputs:} \ref{p3f:thm:v266-gp-exact-k16-junction-condition}, \ref{p3f:eq:v169-inherited-hessian}, \ref{p3f:prop:v1165-analytic-k16-obstruction}.

\subsubsection{Cusp in the terminal Einstein-frame potential}
\label{p3f:sec:v2201-derived-a4}
\noindent\textit{Status: proved corollary on stated domain.}

Use \(\hbar=c=1\), \(C=C_7\), \(M^2=U/4\), \(\delta=R-C\downarrow0\)
and \(A_0=C(U+C)>0\). The restored lower jet gives
\[
\chi=\frac{2\delta}{A_0}+O(\delta^2),\quad
F=1+\frac{\delta^2}{A_0}+O(\delta^3),\quad
f=\delta+\frac{\delta^3}{3A_0}+O(\delta^4).
\]
Put \(y=F-1\). Since \(R-C=\sqrt{A_0y}+O(y)\), the Legendre potential
\(P=FR-f\) satisfies
\[
\boxed{P(F)=C+Cy+\frac23\sqrt{A_0}\,y^{3/2}+O(y^2).}
\]
The exact conformal presentation then gives
\[
V_E=\frac{M^2}{2}\frac{P}{F^2}
=\frac{M^2}{2}\left[C-Cy+\frac23\sqrt{A_0}\,y^{3/2}+O(y^2)\right],
\qquad \phi=\sqrt{\tfrac32}M\log(1+y).
\]
In particular,
\[
\boxed{V_{E,\phi}(0^+)=-\frac{MC}{\sqrt6}\ne0,\qquad
 V_{E,\phi\phi}\sim\frac{A_0}{6\delta}\longrightarrow+\infty.}
\]
\textbf{Proof.} Insert the displayed jets into \(FR-f\); the nonlinear
remainder is \(2\delta^3/(3A_0)\). Invert the leading relation for \(y\),
expand \(F^{-2}\), and differentiate with
\(d\phi/dF=\sqrt{3/2}M/F\). Equivalently,
\(V_{E,\phi}=M(2f-RF)/(\sqrt6F^2)\), which gives the endpoint slope
directly. \par\noindent\textit{Inherited inputs:} \ref{p3d:thm:v173-fixed-depth-laurent-excess-jet}, \ref{p3f:thm:v169-complete-gravity-action}, \ref{p3f:thm:v169-explicit-free-dirac-source}.

Susceptibility decoupling is therefore not a potential minimum on this
one-sided active chart. The equilibrium \(2f-RF=0\), the ideal terminal
linear comparison, and the nonlinear susceptibility endpoint have to remain
distinct. This result does not assert that an actual history reaches the
endpoint or invalidate a separately generated cap/reclosure map.


\section{Conclusions and physical-response boundaries}
\label{paper3f:sec:organized-conclusion}
\label{p3:sec:conclusion}
The doubled local geometry, joint spectral and finite-source action,
Schur reduction, metric and causal variation, and their source-response
identities are assembled on the common domains stated in the preceding
chapters. The exact constructions include native finite-source tests,
scalaron and boundary certificates, photon and transverse responses, and
conditional Petersen and homogeneous-shear comparisons. Their derivations
are mathematical consequences of the typed inputs actually used; their
physical realization is not licensed merely by the existence of an
operator, an invariant family, or a finite stationary extremum.

The unprojected action normalization retains
\(U_J^{\rm raw}=96U\) and
\(G_{J,\rm raw}=G_{\rm scal}/96\). Identification with the
observed Newton coupling requires a realized complete-action metric,
measure, source, clock, and boundary projection. Similarly, a certified
finite source and a causal Green function do not by themselves establish
its occupied physical state, mass, decay channel, or detector coupling.
All conditional source-actualization, positive Schur, finite kinetic,
Lorentzian, quantum and regularity hypotheses remain local to the
claims that consume them. The substantive theorem-level scope and the
explicit retained/open obligations are recorded in
Section~\ref{p3:sec:summary}.
\printbibliography
\end{document}
